\pdfoutput=1
\documentclass[11pt,reqno]{amsart}

\usepackage[utf8]{inputenc}
\usepackage[T1]{fontenc}

\usepackage[margin=1in]{geometry}
\usepackage{letltxmacro}
\emergencystretch=3em
\hbadness=10000
\usepackage{mathtools}
\usepackage{amssymb}
\usepackage{tikz-cd}
\usepackage{booktabs}
\usepackage{tabularx}
\usepackage{array}
\usepackage{amsthm}
\usepackage{longtable}
\usepackage{graphicx}
\usepackage{enumitem}
\usepackage{doi}
\usepackage[backend=biber,style=numeric,sorting=nyt,giveninits=true,doi=true,url=true,isbn=false]{biblatex}
\addbibresource{src/bibliography.bib}
\AtEveryBibitem{\iffieldundef{doi}{}{\clearfield{url}}}
\usepackage{xcolor}\definecolor{linkred}{rgb}{0.78,0.08,0.12}
\usepackage{hyperref}
\hypersetup{colorlinks=true,urlcolor=blue,linkcolor=linkred,citecolor=blue}
\usepackage[capitalize,nameinlink]{cleveref}
\usepackage{colortbl}

% ---- ported from F1_v168 preamble (bookmarks, Unicode, TOC depth) ----
\renewcommand{\thepart}{\Roman{part}}
\pdfstringdefDisableCommands{%
  \def\golden{φ}\def\Zphi{Z[φ]}\def\Fone{F₁}\def\chiv{χ₅}\def\OmH{Ω̂}%
  \def\Lambda{Λ}\def\mathbb#1{#1}\def\mathrm#1{#1}\def\Z{Z}\def\Q{Q}\def\R{R}\def\C{C}%
  \def\log{log }\def\zeta{ζ}\def\pi{π}\def\mu{μ}\def\lambda{λ}\def\to{→}\def\sqrt#1{√#1}\def\tfrac#1#2{#1/#2}}
% Greek letters inside Lean identifiers
\DeclareUnicodeCharacter{03C6}{\ensuremath{\varphi}}
\DeclareUnicodeCharacter{03C9}{\ensuremath{\omega}}
\DeclareUnicodeCharacter{03B9}{\ensuremath{\iota}}
\DeclareUnicodeCharacter{03C3}{\ensuremath{\sigma}}
\DeclareUnicodeCharacter{03C8}{\ensuremath{\psi}}
\DeclareUnicodeCharacter{03B4}{\ensuremath{\delta}}
\DeclareUnicodeCharacter{03B5}{\ensuremath{\varepsilon}}
\DeclareUnicodeCharacter{03BB}{\ensuremath{\lambda}}
\DeclareUnicodeCharacter{039B}{\ensuremath{\Lambda}}
\DeclareUnicodeCharacter{03BC}{\ensuremath{\mu}}
\DeclareUnicodeCharacter{03C4}{\ensuremath{\tau}}
\DeclareUnicodeCharacter{03C7}{\ensuremath{\chi}}
\DeclareUnicodeCharacter{2081}{\ensuremath{_1}}
\DeclareUnicodeCharacter{2082}{\ensuremath{_2}}
\DeclareUnicodeCharacter{2083}{\ensuremath{_3}}
\DeclareUnicodeCharacter{2085}{\ensuremath{_5}}
% The corpus's own symbols in the status ledger: the open statements and the
% classical inputs are transcribed from the declarations' doc comments, so the
% text carries the mathematics they are written in.  Each maps to the LaTeX
% it denotes; \ensuremath keeps it valid both in prose and inside a formula.
\DeclareUnicodeCharacter{039E}{\ensuremath{\Xi}}
\DeclareUnicodeCharacter{03A3}{\ensuremath{\Sigma}}
\DeclareUnicodeCharacter{03B2}{\ensuremath{\beta}}
\DeclareUnicodeCharacter{03B3}{\ensuremath{\gamma}}
\DeclareUnicodeCharacter{03B6}{\ensuremath{\zeta}}
\DeclareUnicodeCharacter{03B7}{\ensuremath{\eta}}
\DeclareUnicodeCharacter{03B8}{\ensuremath{\theta}}
\DeclareUnicodeCharacter{03C0}{\ensuremath{\pi}}
\DeclareUnicodeCharacter{03C1}{\ensuremath{\rho}}
\DeclareUnicodeCharacter{2074}{\ensuremath{^4}}
\DeclareUnicodeCharacter{2075}{\ensuremath{^5}}
\DeclareUnicodeCharacter{2077}{\ensuremath{^7}}
\DeclareUnicodeCharacter{207A}{\ensuremath{^+}}
\DeclareUnicodeCharacter{207B}{\ensuremath{^-}}
\DeclareUnicodeCharacter{2080}{\ensuremath{_0}}
\DeclareUnicodeCharacter{211A}{\ensuremath{\Q}}
\DeclareUnicodeCharacter{211D}{\ensuremath{\R}}
\DeclareUnicodeCharacter{2194}{\ensuremath{\leftrightarrow}}
\DeclareUnicodeCharacter{21D2}{\ensuremath{\Rightarrow}}
\DeclareUnicodeCharacter{2200}{\ensuremath{\forall}}
\DeclareUnicodeCharacter{2203}{\ensuremath{\exists}}
\DeclareUnicodeCharacter{2208}{\ensuremath{\in}}
\DeclareUnicodeCharacter{2212}{\ensuremath{-}}
\DeclareUnicodeCharacter{2218}{\ensuremath{\circ}}
\DeclareUnicodeCharacter{221A}{\ensuremath{\sqrt{\ }}}
\DeclareUnicodeCharacter{2248}{\ensuremath{\approx}}
\DeclareUnicodeCharacter{2260}{\ensuremath{\ne}}
\DeclareUnicodeCharacter{2264}{\ensuremath{\le}}
\DeclareUnicodeCharacter{2265}{\ensuremath{\ge}}
\DeclareUnicodeCharacter{27EA}{\ensuremath{\langle\langle}}
\DeclareUnicodeCharacter{27EB}{\ensuremath{\rangle\rangle}}
\DeclareUnicodeCharacter{27FA}{\ensuremath{\iff}}
% Transcribed ledger prose also mixes text symbols into formulas: these fall
% back to their TS1 text glyphs (\textonehalf & co.), which are undefined in
% math mode and drop the character.  The conditional keeps the text glyph in
% prose and the math symbol inside $...$.
\DeclareUnicodeCharacter{00B2}{\ifmmode^2\else\texttwosuperior\fi}
\DeclareUnicodeCharacter{00B3}{\ifmmode^3\else\textthreesuperior\fi}
\DeclareUnicodeCharacter{00B7}{\ifmmode\cdot\else\textperiodcentered\fi}
\DeclareUnicodeCharacter{00B9}{\ifmmode^1\else\textonesuperior\fi}
\DeclareUnicodeCharacter{00BD}{\ifmmode\frac12\else\textonehalf\fi}
\DeclareUnicodeCharacter{2016}{\ifmmode\Vert\else\textbardbl\fi}
\DeclareUnicodeCharacter{2192}{\ifmmode\to\else\textrightarrow\fi}
\hyphenation{gen-er-al-ised Gel-fond Schnei-der}
\raggedbottom
% Table of contents to sub-subsections; run-in paragraph headings stay out of it
\setcounter{tocdepth}{3}
\makeatletter\def\l@paragraph#1#2{}\makeatother

\crefname{section}{\S}{\S\S}
\Crefname{section}{\S}{\S\S}
\crefname{subsection}{\S}{\S\S}
\Crefname{subsection}{\S}{\S\S}
\crefname{equation}{Eq.}{Eqs.}
\Crefname{equation}{Eq.}{Eqs.}
\crefname{Theorem}{Theorem}{Theorems}
\Crefname{Theorem}{Theorem}{Theorems}
\crefname{Proposition}{Proposition}{Propositions}
\Crefname{Proposition}{Proposition}{Propositions}

\numberwithin{equation}{section}

\newtheorem{theorem}{Theorem}[section]
\newtheorem*{theorem*}{Theorem}
\newtheorem{Theorem}[theorem]{Theorem}
\newtheorem*{Theorem*}{Theorem}

\newtheorem{corollary}[theorem]{Corollary}
\newtheorem{Corollary}[theorem]{Corollary}

\newtheorem{lemma}[theorem]{Lemma}
\newtheorem{Lemma}[theorem]{Lemma}

\newtheorem{proposition}[theorem]{Proposition}
\newtheorem{Proposition}[theorem]{Proposition}

\newtheorem{conjecture}[theorem]{Conjecture}
\newtheorem{Conjecture}[theorem]{Conjecture}

\newtheorem{assumption}[theorem]{Assumption}
\newtheorem{Assumption}[theorem]{Assumption}

\newtheorem{observation}[theorem]{Observation}
\newtheorem{Observation}[theorem]{Observation}

\theoremstyle{definition}
\newtheorem{definition}[theorem]{Definition}
\newtheorem{Definition}[theorem]{Definition}

\newtheorem{note}[theorem]{Note}
\newtheorem{Note}[theorem]{Note}

\newtheorem{example}[theorem]{Example}
\newtheorem{Example}[theorem]{Example}

\newtheorem{remark}[theorem]{Remark}
\newtheorem{Remark}[theorem]{Remark}

\newtheorem{openproblem}{Open Problem}
\newtheorem{Openproblem}{Open Problem}

\crefname{openproblem}{Open Problem}{Open Problems}
\Crefname{openproblem}{Open Problem}{Open Problems}

%%% ====================================================================
%%% CUSTOM COMMANDS
%%% ====================================================================
\newcommand{\cat}[1]{\mathcal{#1}}
\newcommand{\Set}{\mathbf{Set}}
\newcommand{\Hilb}{\mathbf{Hilb}}
\newcommand{\Met}{\mathbf{Met}}
\newcommand{\ob}[1]{\mathrm{ob}(#1)}
\newcommand{\norm}[1]{\|#1\|}
\newcommand{\abs}[1]{|#1|}
\newcommand{\Zmod}[1]{\mathbb{Z}/{#1}\mathbb{Z}}
\newcommand{\Ztwenty}{(\mathbb{Z}/20\mathbb{Z})^\times}
\newcommand{\Rpcf}{R_{\mathrm{PCF}}}
\newcommand{\Ecube}{E^3}
\newcommand{\golden}{\varphi}
\newcommand{\OmH}{\widehat{\Omega}}
\newcommand{\Lchi}{L(s,\chi_5)}
\newcommand{\pcfrepo}{https://github.com/omega-pcf/02-odd-zeta}

%%% ====================================================================
%%%  PENTAGONAL F1-DESCENT MACROS — ported from F1_v168.tex (F1).
%%%  Only macros the shell did NOT already define are added; the shell
%%%  definitions of \golden \Rpcf \OmH \Set \Hilb \Met \norm
%%%  \Ztwenty and \pcfrepo (D1: omega-pcf/02-odd-zeta) are kept as-is.
%%% ====================================================================
% The mark of this manuscript is \witness[ids]{X}, defined below with its
% grammar (docs/witness-contract.md §6.5), its sizes (the 05 rule 13/15, CVR-001
% TYP-4: mark 13, ids 15, prose 16) and its one colour branch.  The class
% palette follows; colour is redundant with the letter (WCAG 1.4.1), so the
% letter alone still carries the meaning.
\definecolor{tierP}{rgb}{0.15,0.35,0.65}   % blue   -- proved in Lean
\definecolor{tierPs}{rgb}{0.70,0.45,0.05}  % ochre  -- written, not yet
% verified against the corpus (docs/witness-contract.md §6.6): the star's
% one denotation, intermediate layers only, zero stars at release
\definecolor{tierN}{rgb}{0.05,0.40,0.20}   % green  -- verified numerically
\definecolor{tierW}{rgb}{0.00,0.35,0.45}   % teal   -- witness item
\makeatletter
\@namedef{markcol@P}{\color{tierP}}
\@namedef{markcol@P*}{\color{tierPs}}
\@namedef{markcol@N}{\color{tierN}}
\@namedef{markcol@W}{\color{tierW}}
% unknown key: no colour, no error (the 05-kkgt form of \markcol)
\newcommand{\markcol}[1]{\expandafter\ifx\csname markcol@#1\endcsname\relax\else\csname markcol@#1\endcsname\fi}
% the family's breakable typewriter underscore (H2/H3: the break is a
% property of the character, and in T1 the three forms of the underscore
% measure the same, so \_ and a raw _ both print through this one constant).
% Line breaking of an id list is comma-first (contract §6.12 S8): the break
% after ", " is the natural interword break (penalty 0) and the break inside an
% identifier is DISCOURAGED (penalty 500), so a list splits at its commas and
% only an identifier longer than the remaining measure is split at all --
% never \allowbreak (= penalty 0), which made the mid-identifier break equally
% attractive and produced twenty_divisor_ / lattice_splits in the PDF.
% The trailing \relax is load-bearing: after a <dimen> TeX scans the keywords
% `plus' and `minus' of \hskip, and an identifier that begins -- right after
% the underscore -- with the letters `minus' or `plus' would be eaten as that
% keyword ("Missing number ... Illegal unit of measure").  \relax is not a
% letter, so both keyword scans push it back and the id prints whole.
\DeclareRobustCommand{\brk@us}{\textunderscore\penalty500\hskip0pt\relax}
\ExplSyntaxOn
\tl_new:N \l__tier_src_tl
\tl_new:N \l__tier_us_tl   \tl_set:Nx \l__tier_us_tl  { \char_generate:nn { 95 } { 8  } }
\tl_new:N \l__tier_usx_tl  \tl_set:Nx \l__tier_usx_tl { \char_generate:nn { 95 } { 11 } }
\tl_new:N \l__tier_uso_tl  \tl_set:Nx \l__tier_uso_tl { \char_generate:nn { 95 } { 12 } }
% \tierids{<arbitrary string>}: \_ and a raw _ (any category code) both
% become the family's breakable typewriter underscore; every other token of
% the argument passes through unchanged.  \tl_replace_all does not descend
% into brace groups, so every wrapper applies the normalisation to its own
% argument -- a group nested two levels deep is out of its reach (the only
% groups inside a bracket today are those of \testref, which normalises at
% its own level).
\cs_new_protected:Npn \tierids #1
  {
    \tl_set:Nn \l__tier_src_tl {#1}
    \tl_replace_all:Nnn \l__tier_src_tl { \_ } { \brk@us }
    \exp_args:NNV \tl_replace_all:Nnn \l__tier_src_tl \l__tier_us_tl  { \brk@us }
    \exp_args:NNV \tl_replace_all:Nnn \l__tier_src_tl \l__tier_usx_tl { \brk@us }
    \exp_args:NNV \tl_replace_all:Nnn \l__tier_src_tl \l__tier_uso_tl { \brk@us }
    \tl_use:N \l__tier_src_tl
  }
% the same, for the identifiers that may also break after a dot
\cs_new_protected:Npn \tieridsdot #1
  {
    \tl_set:Nn \l__tier_src_tl {#1}
    \tl_replace_all:Nnn \l__tier_src_tl { \_ } { \brk@us }
    \exp_args:NNV \tl_replace_all:Nnn \l__tier_src_tl \l__tier_us_tl  { \brk@us }
    \exp_args:NNV \tl_replace_all:Nnn \l__tier_src_tl \l__tier_usx_tl { \brk@us }
    \exp_args:NNV \tl_replace_all:Nnn \l__tier_src_tl \l__tier_uso_tl { \brk@us }
    \tl_replace_all:Nnn \l__tier_src_tl { . } { . \allowbreak }
    \tl_use:N \l__tier_src_tl
  }
\ExplSyntaxOff
% the test that checks a number, cited by its name in the suite that runs
\providecommand{\testref}[1]{{\upshape\small\ttfamily\tierids{#1}}}
% a code mention in prose (contract §6.4, S2): the canonical spelling, the
% same tt body as \testref at the size of the prose
\newcommand{\code}[1]{{\upshape\ttfamily\tierids{#1}}}
% a tag set inside a display must switch to text first: amsart defines
% \footnotesize as \Small, and a font-size command is invalid in math mode
% (the build warned "Command \Small invalid in math mode" at every tag of an
% equation).  In prose the wrapper is the identity, so nothing moves there.
\def\witness@text#1{\ifmmode\text{#1}\else#1\fi}
% One colour branch for the whole mark (contract §6.12 S6): the letter AND the
% ids print in the class colour, so the block and the inline mark read as one
% object -- the inline ids used to print black (#2659a5 against #000000).
\def\witness@ids#1#2{\witness@text{{\normalfont\footnotesize\bfseries\markcol{#2}#2[{\mdseries\upshape\small\ttfamily\tierids{#1}}]}}}
% ---------------------------------------------------------------------------
% THE MARK -- one name, one optional argument, placement-driven (§6.12 S7).
%
%   \witness[ids]{X}     X[ids]  -- the tag first, its ids in the brackets
%
% The shape follows the position of the call in the source, never the class:
% own line in the source (\ifvmode)  -> the block, flush right, \medskip above
% and below, indent 0, no terminal punctuation (S2-S5);
% mid-line in the source (\ifhmode) -> the mark inside the sentence, the
% paragraph does not break and no space is eaten.
% The optional argument is read as an ARBITRARY STRING (§6.5): a \testref, or
% the raw names of our own files, several of them separated by ','.  The macro
% does not judge the shape of what it holds -- a character test (\_, ',' or
% \testref present, otherwise an error) rejects valid internal ids such as
% F1PaperTorus.chi5WeilOrbit and cannot recognise a citation key -- so the only
% error it raises at run time is a literal \cite token inside the bracket: an
% external source stays a \cite set next to the tag, never fused into the mark.
% Whether the bracket holds a citation key is measured where the key set
% exists, by tests/test_witness_census.py against citation.csl.json.
% Sizes are the 05 rule 13/15 (05 CVR-001 TYP-4): the letter in \footnotesize,
% the ids in \small tt, upright.  Colour is redundant with the letter (WCAG
% 1.4.1), so the letter alone still carries the meaning.
% A mark at the head of a paragraph stays inline when the criterion keeps it
% there (§6.10): \leavevmode opens the paragraph, keeps its \parindent and
% leaves the mode \hmode, which is the branch this macro then takes.
\newcommand{\witness}[2][]{%
  \ifvmode
    \par\medskip\noindent\hfill
    \parbox{\dimexpr\linewidth-1em\relax}{\raggedleft
      \normalfont\footnotesize\bfseries\markcol{#2}#2%
      \ifx\relax#1\relax\else\mdseries{\upshape\small\ttfamily[\tierids{#1}]}\fi}%
    \par\medskip
  \else
    \ifx\relax#1\relax
      \witness@text{{\normalfont\footnotesize\bfseries\markcol{#2}[#2]}}%
    \else
      \in@{\cite}{#1}\ifin@
        \PackageError{witness}{the optional argument of \string\witness\space holds internal ids only}
        {An external source is a \string\cite\space set next to the tag, never its key;}%
      \else
        \witness@ids{#1}{#2}%
      \fi
    \fi
  \fi}
\makeatother
\newcommand{\PCF}{\mathrm{PCF}}
\newcommand{\C}{\mathbb{C}}
\newcommand{\R}{\mathbb{R}}
\newcommand{\Z}{\mathbb{Z}}
\newcommand{\Q}{\mathbb{Q}}
\newcommand{\N}{\mathbb{N}}
\newcommand{\F}{\mathbb{F}}
\renewcommand{\Re}{\operatorname{Re}}
\newcommand{\OK}{\mathcal{O}_K}
\newcommand{\Mpcf}{M_{\mathrm{PCF}}}
\newcommand{\GN}{G_{N}}
\newcommand{\epsz}{\varepsilon_{0}}
\newcommand{\lng}{\lambda_{\log}}
% Lean names and file names may break after an underscore or a dot, so that long
% identifiers never run into the margin.
\LetLtxMacro\oldtexttt\texttt
\ExplSyntaxOn
\tl_new:N \l__tt_text_tl
\DeclareRobustCommand{\texttt}[1]
  {
    \tl_set:Nn \l__tt_text_tl {#1}
    \tl_replace_all:Nnn \l__tt_text_tl { . } { . \allowbreak }
    \oldtexttt { \def\_{\textunderscore\allowbreak} \tl_use:N \l__tt_text_tl }
  }
\ExplSyntaxOff
\ExplSyntaxOn
% the same boundary normalisation as \tierids, plus the break after a dot
\DeclareRobustCommand{\LeanInline}[1]
  { \textnormal { \small \ttfamily [ \tieridsdot {#1} ] } }
\ExplSyntaxOff
% an underscore may break a line anywhere in running text
\AtBeginDocument{\DeclareRobustCommand{\_}{\ifmmode\mathunderscore\else\textunderscore\allowbreak\fi}}
%%% paper-specific
\newcommand{\Fone}{\mathbb{F}_1}
\newcommand{\Zphi}{\Z[\golden]}
\newcommand{\Zst}{(\Z/20)^{\times}}
\newcommand{\chiv}{\chi_5}
\newcommand{\Frob}{\mathrm{Fr}}
\newcommand{\LPCF}{\Lambda_{\PCF}}
\newcommand{\Oone}{\mathsf{O}}
\newcommand{\Rloc}{\Z[\golden,\golden^{-1},\tfrac15]}
\newcommand{\Ecore}{\ensuremath{\mathrm{(E)}}}
% A Lean declaration of this version's corpus, marked with a star: the star
% records the pre-compilation provenance of the name (F9-001: process artifact
% of the build, since the modules compile in the green build), not an evidence tag --
% the [P*] class it was defined for was retired in the F11 census (P* = 0).
% FLAGGED, NOT DELETED (W5a): every use of this mark in the manuscript was
% converted to \LeanInline -- the star carried no evidence class, and all 194
% names it marked exist in the green build (18248 jobs, sorry = 0).  The
% definition stays as provenance of what the mark was; the rule that keeps
% it uninvoked is the contract's, not the suite's: docs/witness-contract.md
% §6.5 gives the star one denotation (a name written, not yet verified
% against the corpus), confines it to intermediate layers, and sets the
% release gate at zero stars, so a released text carries none.
\newcommand{\LeanStar}[1]{\LeanInline{#1}$^{*}$}
% a tally: a classification made in the text, recorded as an enumerated type whose
% count is checked by decision; it certifies the count, not the classification
\newcommand{\LeanTally}[1]{\LeanInline{#1}$^{\circ}$}
% A Lean declaration of the companion development of \cite{gonzalez26cw}, not part of
% this paper's build.
\newcommand{\LeanCW}[1]{\LeanInline{#1}$^{\dagger}$}
% Lean accounting of a statement, set after it and outside the statement
\newcommand{\leanline}[1]{\par\noindent{\footnotesize\raggedright\textsc{Lean:} #1\par}\smallskip}
% provenance notes are kept in the source and not printed; see the source map
\newcommand{\src}[1]{}

\begin{document}

\title[Odd Zeta Values from Pentagonal $\Fone$-Descent]
{Odd Zeta Values from Pentagonal $\Fone$-Descent:\\
    a $\Lambda$-structure on $\Zphi$, the character $\chiv$,\\
    and the zeta function}
\thanks{Zenodo All Versions DOI: \href{https://doi.org/10.5281/zenodo.19492791}{10.5281/zenodo.19492791}}

\author[J. A. González García et al.]{Jorge Armando GONZÁLEZ GARCÍA}
\address{TTAMAYO PUNTO COM, S.A.P.I. de C.V.; Research \& Development Division; Naucalpan de Juárez, State of Mexico, Mexico}
\email{jorge@ttamayo.com, victor@ttamayo.com, luz@ttamayo.com}

\author[]{Víctor Manuel GONZÁLEZ GARCÍA}
\author[]{Itzel Marion DRESSLER PÉREZ}
\address{Independent Researcher}
\email{itzel@omega-pcf.com}

\author[]{Luz María GARCÍA ORDÓÑEZ}

\renewcommand{\datename}{\textbf{Received:}}
% Received date, last-modified date and version, with the labels in bold.
% amsart prints \date as the footnote at the foot of page 1 (\@adminfootnotes
% -> \@setdate, which prefixes \datename), so the first label lives in
% \datename and the line carries the other two: without this the class would
% prefix its own "Date:" and the footnote would read "Date: Received: ...".
% \today is reproducible: scripts/tasks/compile.ts compiles with
% SOURCE_DATE_EPOCH set to the HEAD commit's timestamp (scripts/utils/git.ts).
% The version token mirrors the "version" field of package.json.
% selfcite:begin id="version"
\date{April 10, 2026; \textbf{Last modified:} \today; \textbf{Version:} \texttt{v2.0.0}}
% selfcite:end id="version"

\keywords{Field with one element ($\mathbb{F}_1$); $\Lambda$-rings; golden ratio;
Frobenius lifts; Dirichlet $L$-functions; Weil positivity; explicit formula;
formal verification}

\subjclass[2020]{Primary 11M26, 11R04, 14G99; Secondary 11M20, 68V20}

\begin{abstract}
We demonstrate that the odd values $\zeta(2k+1)$ of the Riemann zeta
function are structurally determined by $\varphi=(1+\sqrt{5})/2$ and~$\pi$.
The Euler product of $\zeta$ is built from local factors $f_p(s)$, one
for each prime.  We prove that the Frobenius lift at each prime $p\ne5$,
congruent modulo $p$ to $\golden^p = F_p\golden + F_{p-1}$---where $F_n$ is
the $n$-th Fibonacci number---determines the splitting type of $p$ in
$\mathbb{Z}[\golden]$, and hence, with the ramified prime $5$, which admits
no lift, every $f_p(s)$. These lifts constitute $\mathbb{F}_1$-descent data
in the sense of Borger away from $5$. Through the Dedekind factorisation
\[
  \zeta(s) = \frac{\zeta_{\mathbb{Q}(\sqrt{5})}(s)}{L(s,\chi_5)},
\]
this determines $\zeta(s)$ completely: the Frobenius structure implies
the Euler product at every level (primes, splitting types, local factors,
$L$-function), each implication proved, anchored by the base case
$L(1,\chi_5) = 2\log\golden/\sqrt{5}$
(the class-number formula for $\mathbb{Q}(\sqrt{5})$),
which expresses the $L$-function value entirely in terms of~$\golden$;
and the factorisation is a natural isomorphism over $\mathbb{Q}$ and the
quadratic fields of discriminant dividing~$20$.
The determination fixes the odd values through the pentagonal arithmetic
without evaluating them, and by itself it does not decide their
arithmetic nature.  The arithmetic is formalised in Lean 4 with Mathlib, and each
statement carries its tier: proved, assumed from the literature, numerical, or open.

\end{abstract}

\maketitle

\tableofcontents

\pagebreak
\section{Introduction}\label{sec:intro}
% =====================================================================

\subsection{Objective}\label{sec:objective}

What determines the values $\zeta(3)$, $\zeta(5)$, $\zeta(7)$, \ldots?
Euler showed that the even values $\zeta(2k) = r_k\pi^{2k}$
($r_k \in \mathbb{Q}$) are governed by $\pi$ alone.
The odd values have resisted a comparable structural account for
over two centuries: Ap\'ery~\cite{apery79} proved
$\zeta(3)\notin\mathbb{Q}$,
Ball and Rivoal~\cite{ballrivoal01} showed that infinitely many
$\zeta(2k+1)$ are irrational, but no organising principle has
emerged.  The current study identifies such a principle: the pentagonal
arithmetic of $\mathbb{Q}(\sqrt{5})$, generated by
$\golden^2 = \golden + 1$.  At $s=1$ the pentagonal value $L(1,\chi_5)=2\log\golden/\sqrt{5}$ is transcendental, because $\log\golden$ is; at the odd $s\ge3$ the arithmetic nature of $L(s,\chi_5)$ is open, and the determination does not decide it (\cref{rmk:nature-L}).

The key identity is the Dedekind factorisation (\cref{eq:dedekind})
\begin{equation}\label{eq:main-factorization}
  \zeta(s) = \frac{\zeta_{\mathbb{Q}(\sqrt{5})}(s)}{L(s,\chi_5)},
\end{equation}
where the Artin character $\chi_5$ is dictated by the pentagon (\cref{thm:pentagonal-values}):
\begin{equation}\label{eq:chi5-pentagon}
  \chi_5(a) = \frac{\log|2\cos(\pi a/5)|}{\log\golden}\in\{+1,-1\}.
\end{equation}
This character classifies every prime $p$ as split or inert in
$\mathbb{Z}[\golden]$, determining the local Euler factor $f_p(s)$
and hence the entire product $\zeta_{\mathbb{Q}(\sqrt{5})}(s) = \prod_p f_p(s)$, with base case
$L(1,\chi_5) = 2\log\golden/\sqrt{5}$.  The factorisation~\eqref{eq:main-factorization} holds for every quadratic
field with its own character; what belongs to $\mathbb{Q}(\sqrt5)$ is that each of its components is fixed by
the pentagonal arithmetic and by nothing else, the character through~\eqref{eq:chi5-pentagon}.

For every $k\ge1$, $\zeta(2k+1)=\zeta_{\mathbb{Q}(\sqrt5)}(2k+1)/L(2k+1,\chi_5)$, with both sides
built from the pentagon: $\zeta_{\mathbb{Q}(\sqrt5)}$ is the Euler product whose every factor is
fixed by the splitting type, and $L(s,\chi_5)$ is a combination of the Hurwitz values $\zeta(s,a/5)$
at the pentagonal points (\cref{thm:hurwitz}).  The determination is unconditional and is proved
with the analytic objects of Mathlib (\cref{thm:odd-zeta}).  The even values come from $\pi$ and
lie in $\mathbb{Q}\pi^{2k}$, those of $L(s,\chi_5)$ in $\sqrt5\,\pi^{2k}\mathbb{Q}$
(\cref{ssec:pentagon-odd}); and from the anchor $L(1,\chi_5)=2\log\golden/\sqrt5$ every value of
$\zeta$ and of $L(\cdot,\chi_5)$ is written in base $\golden$ (\cref{ssec:base-phi}).

The objective is summarized in the following diagram.  In it, $G\colon\Zst\to\Z_2\times\Z_2$ is the
homomorphism on the eight classes of the primes modulo $20$ whose two components are the character $\chi_4$ of
$\Q(i)$ and $\chiv$; its kernel $\{1,9\}$ has two elements.  The number $\mu_3=\tfrac12$ is $1/|\ker G|$, and it is
also the modulus of the three eigenvalues $\tfrac12\zeta_3^{\,k}$ ($\zeta_3=e^{2\pi i/3}$, $k=0,1,2$), which the
symmetric group $S_3$ permutes; the subscript records the three eigenvalues, and $\mu_3$ is a number, not the group
of cube roots of unity.  These are the eigenvalues of the operator of our companion paper in its spectral form
$\OmH=\tfrac12\operatorname{diag}(1,\zeta_3,\zeta_3^2)$, which is normal; the same operator has self-adjoint
forms, on the tower and on the ordinates of the zeros (\cref{thm:operator-forms}).
\begin{equation}\label{eq:objective}
  \begin{tikzcd}[column sep=3.4em, row sep=2.6em]
  & & & {\textstyle\prod f_p}
    \ar[r, "{\text{\tiny colimit}}"', "{\text{\tiny(\ref{prop:euler-colimit})}}"]
    & \zeta_{\mathbb{Q}(\sqrt{5})}
    \ar[r, "{\text{\tiny factorise}}"', "{\text{\tiny(\ref{eq:dedekind})}}"]
    & \zeta(2k{+}1) \\
  \golden^2{=}\golden{+}1
    \ar[r, "{\text{\tiny Frob.}}"', "{\text{\tiny(\ref{thm:frobenius-congruence})}}"]
    & \Frob_p
    \ar[r, "{\text{\tiny sign}}"', "{\text{\tiny(\ref{thm:sign-is-chi5})}}"]
    & \chi_5
    \ar[r, "{\text{\tiny class}}"', "{\text{\tiny(\ref{prop:G})}}"]
    & G \ar[u] \ar[d] & & \\
  & & & \mu_3{=}\tfrac{1}{2}
    \ar[r, "{\text{\tiny triad}}"', "{\text{\tiny(\ref{thm:S3-from-euler})}}"]
    & S_3
    \ar[r, "{\text{\tiny eig.}}"', "{\text{\tiny(\ref{prop:omega-spectrum})}}"]
    & \zeta_3^{\,k}/2
\end{tikzcd}
\end{equation}
Both branches originate from the same map $G$ (\cref{prop:span}).
The upper branch produces the odd zeta values: $G$ classifies
primes, determining every $f_p$, assembling the Euler product,
and delivering $\zeta(2k+1) = \zeta_{\mathbb{Q}(\sqrt{5})}(2k+1)/L(2k+1,\chi_5)$.
The lower branch records the spectral structure: $1/|\ker G|=\tfrac12=\mu_3$, the modulus of the eigenvalues
$\zeta_3^{\,k}/2$ that $S_3$ permutes (\cref{thm:S3-from-euler,prop:omega-spectrum}).  Its arrows record the
order in which the construction meets these numbers: they are separate identities, each proved, and not
implications between them (\cref{thm:S3-from-euler}).
The regulator $R_K = \log\golden$ is the arithmetic size of the field,
$L(1,\chi_5) = 2R_K/\sqrt{5}$ (\cref{thm:L1}).
\Cref{sec:colimit} makes the categorical content
explicit: $\{\Frob_p\}$ is a functor, the Euler product is a colimit, and
the Dedekind factorisation is a natural isomorphism across the four
quadratic fields classified by~$G$; the two branches form a span
with common source~$G$ (\cref{prop:span}).

This development grows out of our companion paper \emph{The Hilbert--P\'olya Operator and the
Primitive Structure of the Complex Plane}~\cite{gonzalez25hp}, hereafter \emph{the Primitive
Structure}, where we attempted to join the framework of the P\'olya operator with the $\Fone$
programme of Borger and Manin --- rings with commuting Frobenius lifts~\cite{borger09}, and the
noncommutative tori that Manin proposes and that string theory produces~\cite{manin04,connes98} ---
together with the classification of the primes by Grisales Herrera~\cite{grisales25}, made with
the pentagonal cyclotomic polynomials and $\golden$: the classification of the eight classes modulo $20$
that $G$ divides into four fibres; evaluated at $\golden$, the cyclotomic polynomials $\Phi_n$ later
read the order of $\golden$ modulo each prime (\cref{ssec:why-phi}).  From that classification alone we
constructed a self-adjoint operator, $\tilde H_{\PCF}$, whose eigenvalues predict the heights of the zeros of
$\zeta$ with relative error below $1.7\%$ across twelve orders of magnitude (its form is recalled in
\cref{thm:operator-forms}).  The primes enter that construction only through their classes: neither their
distribution, nor the zeros, nor the Euler product is an input.  \emph{PCF} stands for
Precedent--Current--Forthcoming, the three consecutive terms that the Fibonacci recurrence binds
(\cref{ssec:tripartite}).

\tier{P}\ marks a statement proved in Lean~4~\cite{demoura21} with Mathlib~\cite{mathlib20};
\tier{N}\ marks a statement verified numerically, with the test named in the mark.  The protocol
they follow is \cref{ssec:protocol}.

The Primitive Structure formalises the forgetful functor chain
$C^*\text{-}\mathbf{Alg}\to\mathbf{Hilb}\to\mathbf{Met}\to\mathbf{Set}$ and the
spectral constraint that $\tfrac12$ is the only positive real part in the spectrum of
$\OmH$ (here \LeanInline{positive\_spectral\_half}).  Its proof that the real part of the zeros is $\tfrac12$
rested on one hypothesis, \texttt{bridge\_axiom}: the transport of that constraint to the zeros, which identifies
the real part of each zero with a spectral position of $\OmH$ (\cref{sec:context}).  The transport required a
common source for the Euler product and for the spectrum, and the span of the diagram is that source: in it
(\cref{prop:span}), $G$ produces $\zeta(s)$ through the Euler product~\cite{euler1737} and analytic
continuation~\cite{riemann1859}, and the spectral number $\mu_3=\tfrac12$ through $|\ker G|=2$
(\LeanInline{G\_hom}, \LeanInline{G\_surj}, \LeanInline{G\_kernel} \tier{P}).  The span does not by itself
discharge the axiom: none of its pieces speaks of a zero (\cref{rmk:markers}).  The bridge is
carried to its last link, which is equivalent to the statement \Ecore\ that every zero of $L(s,\chiv)$ in the
critical strip has real part $\tfrac12$ (\cref{ssec:core-forms}), and which a conjectured law of the primes by
classes, stated in the Results, would give (\cref{cor:closure}).  The determination of the odd values
(\cref{thm:odd-zeta}) does not depend on the \texttt{bridge\_axiom}.

The Primitive Structure stated its bridge for the zeros of $\zeta$; the statements about zeros in this paper are
about those of $L(s,\chiv)$, and the reason is the descent.  On $\Z$ every Frobenius lift is the identity, so the
descent carries no datum there; on $\Zphi$ the lift at each prime is the identity or the conjugation, its sign is $\chiv$ (\cref{thm:sign-is-chi5}), and $\chiv$ is the character that splits the zeta function of the field,
$\zeta_{\Q(\sqrt5)}=\zeta\cdot L(\cdot,\chiv)$.  The factor that the pentagon contributes, and that the
construction of this paper reaches, is $L(s,\chiv)$: for its zeros the first bound of the Primitive Structure,
$\mathrm{Re}(\rho)\le\mu_3$, is equivalent to GRH for $L(s,\chiv)$ (\LeanInline{premise\_A\_iff\_GRH\_chi5} \tier{P},
\cref{ssec:core-forms}).  For $\zeta$ itself the paper determines the Euler product and the odd values (\cref{thm:odd-zeta}); its
statements about the position of zeros are about those of $L(s,\chiv)$.

On the side of the $\Fone$ programme, the Primitive Structure claimed a $\Lambda$-ring structure on
$\Z[\golden,\golden^{-1},\tfrac12]$.  This paper makes it exact (\cref{rmk:monoid-not-ring}) on the ring
of the pentagon $\Zphi$, where every step can be checked; here we develop the $\Fone$ side that we left
undeveloped there, as our proposal for an $\Fone$ framework.  The descent datum has two halves: the
$\Lambda$-monoid of the golden tower, the powers $\golden^n$, $n\in\Z$, which are the positive units of
$\Zphi$, and the Artin lifts on $\Zphi$ away from $5$, obstructed at the ramified prime (\cref{thm:lambda-monoid,thm:frobenius-congruence,thm:obstruction-five}), which together are the one $\Lambda$-structure that $\Zphi$ carries at the primes other than $5$ (\cref{thm:descent-datum}); the
character $\chiv$ is read off the Frobenius, and each Euler factor is the characteristic
polynomial of a Frobenius (\cref{thm:chi5-from-frobenius,thm:frob-eigen}).  The same arithmetic
reaches the zeros of $L(s,\chiv)$ by a single thread from $\golden$: a chain of eight stages, each using only
what the previous ones produced, which never takes the zeros or the Euler products as a starting point; the primes
enter it only as an index set, through their classes, and no property of their distribution is used
(\cref{sec:topological-proof}).  Each zero $\rho$ has a coordinate $r(\rho)=|\rho-1|/|\rho|$, and \Ecore\ is $r(\rho)=1$; this splits
into a multiplicative half, $r(\rho)\,r(1-\rho)=1$, and an additive half, $r(\rho)+r(1-\rho)=2$.  For every zero the
thread proves the symmetries of its orbit $\{\rho,1-\rho,\bar\rho,1-\bar\rho\}$, that this orbit is centred at $\tfrac12$, and the
multiplicative half; what remains is the
additive half, one statement proved equivalent to GRH$(\chiv)$ in ten forms
(\cref{prop:two-halves}, \cref{ssec:core-forms}).  Along Weil's route the argument reaches a closing theorem whose
one open input is the positivity of Weil's quadratic form on the translates of a test function by every power of the
Frobenius (\cref{thm:weil-closure}); that positivity is proved for test functions supported on a first interval and
certified numerically on two longer ones (\cref{thm:first-window,thm:second-window,thm:third-window}), and the conjectured law
of the primes by classes would give it in full (\cref{conj:closure}).

Read together, these results form a dictionary in the sense of Weil's Rosetta stone (\cref{sec:weil}), built by
classification under equivalence~\cite{gonzalez26cw}.  It has four columns --- the arithmetic of a one-dimensional
vector, the line of the exponents of the golden tower (\cref{ssec:dimensional-chain}); the geometry of the pentagon and
its cyclotomic polygons; the algebra of $\Zphi$; and a Hermitian operator, the self-adjoint operator of the companion paper (\cref{thm:one-operator,thm:operator-forms}) --- crossed by a law read at each prime.  At a prime $\mathfrak p$ of $\Zphi$ the unit $\golden$ has a finite order in
the multiplicative group of the residue field, which is the order of its angular part in the completion at
$\mathfrak p$; this order, the \emph{clock} of the prime, is the law, and the cyclotomic polynomials evaluated at
$\golden$ read it: for $\mathfrak p$ prime to $n$, the clock is $n$ exactly when $\mathfrak p$ divides
$\Phi_n(\golden)$ (\cref{prop:clock-cyclotomic}).  The dictionary shows what an operator built from Fourier analysis does not reach in this structure, which we
propose as the reason why such an operator has been elusive.  Over each level $n$ of the tower, the power $\golden^n$, the structure carries a finite
group of six elements, a colour of order three and an orientation of order two.  A decomposition in abelian characters, Fourier's among them, is indexed by the level alone, and any
self-adjoint operator that depends on the level alone, with its whole functional calculus, reads only the level.
Neither sees the six elements over a level: only an operator that does not commute with the symmetry of those six
elements separates them (\cref{prop:level-operators}).  The operator of the fourth column depends on the level alone, so this
holds for it too, and it is why one self-adjoint operator is not enough here (\cref{rmk:one-operator}).  What is irrational in the structure, $\golden$ and the logarithms of the primes, the characters and the
operators of the level see only as statistics.  For
$\golden$, multiplication by $\golden$ on $\Zphi\cong\Z^2$ is the Fibonacci matrix
$\bigl(\begin{smallmatrix}0&1\\1&1\end{smallmatrix}\bigr)$: on the real torus $\R^2/\Z^2$ it is mixing, so the
characters see it only through the decay of correlations, while on the points of order $p$ it is exactly periodic,
with period the order of $\golden$ modulo $p$.  For the logarithms of the primes, it is the fluctuation of the zeros
(\cref{prop:two-faces,prop:torus-line,prop:error-partition}).

A previous version of this paper~\cite{gonzalez26oddzeta} also addressed physical aspects of the
construction --- the torus $T^2_{\PCF}$ and its period, the entropy of one bit, the Brown--Henneaux
central charge, and the Wilson coefficients of the $\mathcal{N}{=}4$ effective theory of Elvang
et al.~\cite{elvang26}.  These are now developed in two companion papers on physics, the Crystalline
Worldsheet~\cite{gonzalez26cw} and the de Sitter companion~\cite{gonzalez26ds}, which puts one object in place of an
ensemble for the observer in de Sitter space.  Here we develop the mathematical side, and \cref{sec:disc-cw}
records where the mathematics of this paper meets them.

The noncommutative tori that Manin takes as the geometry of real
multiplication are those of string theory~\cite{manin04,connes98}; the Crystalline Worldsheet places a
string, a ring, inside such a torus, much as holography places a theory on a boundary, and develops the
connection with physics, which meets the Riemann hypothesis both through the Veneziano amplitudes of
strings and through the link between the Gaussian Unitary Ensemble, the zeros of $\zeta$ and holography.
This paper develops the $\Fone$ programme of that construction, and its objects come from it: holography's
missing operator, the bridge of a ring inside a torus, the dimensional emergence of the vanishing point of the
paper on the simultaneous identities of the golden ratio~\cite{gonzalez26mersenne}, and the structure that
emergence gives the operator.  In that paper the vanishing point is the Euler rotor, read as the vanishing point of
classical perspective and formalised with $\golden$.  Here the same point is unfolded from the line of the tower to
the plane of the Euler product and of the zeros, and to a third dimension over that plane.  That third dimension is
geometric and not algebraic: it is reached by the symmetry $S_3$ of three directions, and no algebra of dimension
three is involved.

\subsection{Context}\label{sec:context}\label{ssec:intro-circ}

The proof of the Primitive Structure that $\mathrm{Re}(\rho)=1/2$ is a squeeze (Theorem~10
of~\cite{gonzalez25hp}): two independent bounds, both derived from the arithmetic of $\Ztwenty$ and of
$G$ (Definition~5 of~\cite{gonzalez25hp}), force equality at~$1/2$, and the squeeze itself is verified
in Lean~4.  The first bound, Premise~A (Proposition~3 of~\cite{gonzalez25hp}), is
$\mathrm{Re}(\rho)\le\mu_3=1/2$.  It is the hypothesis \texttt{bridge\_axiom} of \cref{sec:objective}, and in
that paper it remains an axiom (the field \texttt{bounded\_by\_mu} of the structure \texttt{TernaryLZero},
Definition~7 of~\cite{gonzalez25hp}).  The second bound, Premise~B, is $1-\mathrm{Re}(\rho)\le\mu_3$: once
Premise~A is provided, the functional equation, there the classical input \texttt{hecke\_1920}, gives it
(Theorem~10 of~\cite{gonzalez25hp}).  The two bounds give $\mathrm{Re}(\rho)=1/2$.

What the Primitive Structure offers for Premise~A is the Spectral Embedding (Proposition~5
of~\cite{gonzalez25hp}): the Euler product enters~$\Ztwenty$, $G$ classifies the primes into four
characters, and the $S_3$ symmetry of the tripartite object restricts
$\mathrm{spec}(\OmH)=\tfrac{1}{2}\{1,\zeta_3,\zeta_3^2\}$, with real parts $\{1/2,\,-1/4\}$, both
${\le}\,\mu_3$.  As justification it cites the Euler product identity, Gelfand's theory of Banach
algebras~\cite{gelfand41} and the spectral mapping theorem of Dunford and
Schwartz~\cite{dunfordschwartz58}, but it does not show that the three act on a common structure: the
Euler product organised by~$G$ through the same map that produces $\mathrm{spec}(\OmH)$.  That common
structure is the span of \cref{sec:objective}.

This paper supplies it from the side of the $\mathbb{F}_1$-geometric programme (Borger~\cite{borger09},
Manin~\cite{manin95}, L\'opez Pe\~na--Lorscheid~\cite{lopezpena09mapping}), in relation to the Euler
product.  The descent datum on $\Zphi$ is at once the $\mathbb{F}_1$-descent datum in the sense of Borger
and the source of every local Euler factor $f_p(s)$, through the Fibonacci splitting criterion
$\chi_5(p)\equiv F_p\pmod{p}$ (\cref{thm:s1}); on $\Z[\golden,\golden^{-1},\tfrac12]$, the ring of the companion paper, only its first half survives
(\cref{cor:localisation}).  In building the categorical infrastructure that carries it --- the
Frobenius system as a functor (\cref{ssec:frobenius-functor}), the Euler product as a colimit
(\cref{ssec:euler-colimit}), the Dedekind factorisation as a natural isomorphism
(\cref{ssec:dedekind-natural}), and the span (\cref{ssec:span}) --- we found that the same chain also
determines the odd zeta values: not through new closed forms, but through the
identity~\eqref{eq:main-factorization} together with the class-number formula for $\mathbb{Q}(\sqrt{5})$.
As \cref{sec:objective} states, the identity holds for every quadratic field, and what belongs to
$\mathbb{Q}(\sqrt{5})$ is that its components are fixed by the pentagonal arithmetic: by the same classification
from which the operator of the Primitive Structure was built.
This answers, structurally, the question with which this introduction opens, a problem older than the
Riemann Hypothesis.

\subsection{Organisation}\label{ssec:organisation}

\Cref{sec:background} states the problem and the $\Fone$ programme.  \Cref{sec:methods} builds the construction and carries it to the zeros.  \Crefrange{sec:choice}{sec:grisales} build the pentagon, the ring $\Zphi$, the tower with its two registers, the Artin $\Lambda$-structure with the $\delta$-ring and the polar
form of $\golden$ at each prime, and the character with the classes, the angular part of $\golden$ at each prime, its order modulo the prime and its index.  \Crefrange{sec:dedekind}{ssec:base-phi} carry it to the zeta functions: Euler factors, the colimit, Dedekind, the odd zeta values, the counting
function, the Mersenne anchor, and the values in base $\golden$.  \Crefrange{sec:operator}{sec:topological-proof} go from $\golden$ to the zeros: the operator in all its forms, then the structure that places the zeros against it, and finally the thread from $\golden$ to the zeros, in eight stages, assembled from everything built before it.  \Cref{sec:results} collects the results: what is proved on
the zero set, the Weil analogy, the windows of Weil positivity, and the verdict that
remains, which ends with Weil's route carried to a closing theorem with one open input, the positivity of
Weil's form for every power of the Frobenius (\cref{ssec:weil-closure}).  \Cref{sec:disc} discusses what has been demonstrated, against the $\Fone$ programme, Borger, Manin and
Weil, and the connection with the companion papers of the series: the paper on the simultaneous
identities of the golden ratio and Euler's identity~\cite{gonzalez26mersenne}, henceforth CR, whose module
--- origin, radius, angle and magnitude --- it reaches on the zeros, the radius being what remains open;
the Primitive Structure, whose inputs it settles and whose operator has to be taken to the prismatic and
perfectoid setting for the argument to close; and the two papers on physics, the Crystalline Worldsheet, which continues Manin's
programme on that side, and the de Sitter companion, which reads the bulk--boundary map as a code with a
one-dimensional code subspace and the closed string as the cut of a tensor network.  \Cref{sec:conclusions} concludes.

\subsection{Protocol}\label{ssec:protocol}

Two evidence classes are marked where a statement is stated: \tier{P}\ for a statement proved in
Lean~4 with Mathlib, \tier{N}\ for a statement verified numerically, with the test named in the
mark; a statement with no mark is standard, and carries its citation where it enters.

\emph{Notation.}  Several symbols are overloaded --- $\psi$, $\Lambda$, $\sigma$, $\tau$, $\ell$, $L$, $S$, $N$, $R$, $T$ and $\mu$ ---
and the index of notation below gives each reading with the place where it is used.  In particular
$\mu_3=\tfrac12$ is a number, while $\mu_n$ for $n=4,5,6,20$ is the group of $n$-th roots of unity and
$\mu(\cdot)$ is the M\"obius function.

All Lean names of the paper are in one file,
\texttt{F1\_pentagonal\_descent\_UNIFIED.lean}, whose last block is the paper itself: a
namespace for each section, with an alias for every declaration the section cites.
Where a result also appears in the companion papers it is cited for
co-attribution rather than as authority (\cref{sec:companions}).  Every claim of
the corpus, with its tier and its declaration, is listed in \cref{app:ledger};
the corrections to earlier versions of the programme are in
\cref{app:corrections}.

\subsubsection{Index of notation}\label{app:notation}
Each symbol with its meaning in one line and the place where it is defined.

{\small\setlength{\tabcolsep}{4pt}
\begin{longtable}{>{\raggedright\arraybackslash}p{0.2\textwidth}>{\raggedright\arraybackslash}p{0.56\textwidth}>{\raggedright\arraybackslash}p{0.18\textwidth}}
\toprule Symbol & Meaning & Defined in\\\midrule\endhead
\multicolumn{3}{l}{\emph{The ring $\OK=\Zphi$ and the golden tower $T_\golden$}}\\
$\golden$ & $(1+\sqrt5)/2=2\cos(\pi/5)$, the root of $X^2-X-1$ & \cref{sec:choice}\\
$\phi$ & Euler's function & \cref{ssec:protocol}\\
$K$, $\OK=\Zphi$ & $\Q(\sqrt5)$ and its ring of integers, $\Z[X]/(X^2-X-1)$ & \cref{prop:OK}\\
$\omega$ & the class of $X$ in $\Zphi$ & \cref{ssec:protocol}\\
$F_n$, $L_n$ & the Fibonacci and Lucas numbers; $\golden^n=F_n\golden+F_{n-1}$, $L_n=\golden^n+(-\golden)^{-n}$ & \cref{def:fib-PCF-labels,prop:binet}\\
$P_n,C_n,F_n^{*}$; $\Omega$; PCF & $F_{n-1},F_n,F_{n+1}$ (Precedent, Current, Forthcoming), so that $F_n^{*}=C_n+P_n$; the tripartite generator $\Omega=P\cdot C\cdot F$; PCF names both & \cref{def:fib-PCF-labels,def:tripartite-omega}\\
$T_\golden$, $T_\Z$ & the golden tower $\{\golden^x:x\in\R\}$ and its lattice $\{\golden^n:n\in\Z\}$, the positive units of $\OK$; $n$ is the \emph{level} & \cref{def:golden-tower,thm:units}\\
$X_{\PCF}$ & the exponent monoid $(\Z,+)$ of $T_\Z$ & \cref{thm:lambda-monoid}\\
$\Lambda$ & a $\Lambda$-structure in Borger's sense: a commuting family of Frobenius lifts & \cref{sec:lambda-ring}\\
$\sigma$ & the conjugation of $K$, $\golden\mapsto1-\golden$, the second endomorphism of $\Zphi$ & \cref{prop:endo}\\
$\sigma(s)$, $\tau(s)$ & the involutions $s\mapsto1-s$ and $s\mapsto\bar s$ of $\C$ generating the Klein group of the zeros; written with their argument, to distinguish $\sigma(s)$ from the conjugation $\sigma$ of $K$ & \cref{ssec:klein}\\
$\Frob_p$ & the Frobenius lift at $p$: on the tower, $\golden^n\mapsto\golden^{pn}$ (on $X_{\PCF}$, $n\mapsto pn$), the multiplicative half of the descent datum; on $\Zphi$, $p\ne5$, the identity or the conjugation by $\chiv(p)$, the additive half.  In Lean names the lifts are written \texttt{psi}, $\psi$ & \cref{thm:lambda-monoid,def:artin-frob}\\
$\mathcal V$ & the normalised triad of $\OmH$, a column $\mathbb{C}\to\mathbb{C}^3$ with $\mathcal V^\dagger\mathcal V=1$ & \cref{ssec:why-phi}\\
$M_p$ & the matrix of $\Frob_p$ on the basis $(1,\sqrt5)$, $\operatorname{diag}(1,\chiv(p))$ & \cref{thm:frob-eigen}\\
$\delta_p$ & the $p$-derivation $(\Frob_p(a)-a^p)/p$ & \cref{ssec:delta-polar}\\
clock of $\mathfrak p$, $e_p$ & the order of $\golden$ in $(\Zphi/\mathfrak p)^\times$; $e_p$, its order in $(\Zphi/p)^\times$, the least common multiple of the clocks above $p$ & \cref{ssec:classes-angular}\\
$i_p$, $N_p$ & the index of $\golden$ at $p$, with $e_pi_p=N_p$, $N_p=p-1$ ($p$ split) or $2(p+1)$ ($p$ inert) & \cref{prop:phi-layers,prop:index-ngon}\\
$\Phi_n$, $\eta_d$ & the $n$-th cyclotomic polynomial; $\eta_d=\log|N(\golden^d-1)|-d\log\golden$ & \cref{prop:clock-cyclotomic}\\
$M(x)$, $\tau(N)$ & the Mertens function $\sum_{k\le x}\mu(k)$; the number of divisors of $N$ & \cref{prop:clock-cyclotomic,prop:index-ngon}\\
$\lng$ & $\log_\golden2$, the anchor: $\golden^{\lng}=2$ & \cref{sec:mersenne}\\
$\sigma(p)$ & the exponent of the Mersenne clause, $3\golden^{\sigma(p)}=2^p$ & \cref{sec:mersenne}\\
$R_K$ & $\log\golden$, the regulator of $K$ (Lean \texttt{regulatorK}) & \cref{ssec:fund-formula}\\
$R_{\PCF}$ & $\Z[\golden,\golden^{-1},\tfrac12]$; it carries the $\Lambda$-monoid of the tower and not the Artin $\Lambda$-structure & \cref{cor:localisation}\\
$\sigma_p$ & $\log_\golden p$, the magnitude of the prime $p$ & \cref{prop:euler-tower}\\
$v(x)$ & $\log x/\log\golden$, the golden valuation & \cref{prop:tower-disc-val}\\
\multicolumn{3}{l}{\emph{The classes $\Zst$ and the characters}}\\
$\Zst$ & the eight classes of the primes modulo $20$ & \cref{sec:grisales}\\
$G$ & the homomorphism on $\Zst$ whose components give $\chiv$ and $\chi_4$; $|\operatorname{im}G|=4$, $|\ker G|=2$ & \cref{prop:G}\\
$\chiv$, $\chi_4$ & the characters of $\Q(\sqrt5)$ and of $\Q(i)$, of conductors $5$ and $4$ & \cref{thm:s1,prop:G}\\
\multicolumn{3}{l}{\emph{The zeta and $L$-functions}}\\
$\zeta_K=\zeta_{\Q(\sqrt5)}$ & the Dedekind zeta function of $K$; $\zeta_K=\zeta\cdot L(s,\chiv)$ & \cref{thm:odd-zeta}\\
$f_p(s)$ & the local Euler factor of $\zeta_K$ at $p$, a function of $\chiv(p)$ & \cref{sec:dedekind}\\
$\Lambda_K$ & the von Mangoldt function of $K$ & \cref{sec:counting}\\
$L(s,\chiv)$, $\Lambda(s,\chiv)$ & the $L$-function and its completion & \cref{sec:analytic-inputs}\\
$S(X)$, $E(X)$ & a sum of a power-residue symbol over the primes and its energy & \cref{prop:sqrt-x}\\
$\mathcal K$, $k_s$ & Cram\'er's constant and its kernel & \cref{prop:energy-scale}\\
$V_\chi$, $\kappa$, $\gamma_{\mathrm E}$ & the amplitude $2\Re\,\xi'/\xi(1,\chi^*)$, the parity of $\chi^*$, Euler's constant & \cref{prop:closure-amplitude}\\
$\theta(x;20,a)$ & $\sum_{p\le x,\,p\equiv a}\log p$ & \cref{prop:closure-amplitude}\\
$\zeta(s,a/5)$ & the Hurwitz values at the pentagonal points & \cref{thm:odd-zeta}\\
\multicolumn{3}{l}{\emph{The operators $\OmH$ and $\Oone$}}\\
$\zeta_3$ & $e^{2\pi i/3}$ & \cref{ssec:protocol}\\
$\OmH$ & the colour operator, the power $\tfrac12Y^4$ of the phase of $\Oone$, with exact spectrum $\OmH_k=\tfrac12\zeta_3^{\,k}=\golden^{-\lng}\zeta_3^{\,k}$, $k=0,1,2$ & \cref{prop:omega-spectrum}\\
$\mu=\mu_3,\sigma,\alpha_0,\beta_0$ & the parameters $\tfrac12,\tfrac32,\tfrac{59}{40},\tfrac18$, read from the cardinals of $G$ and the conductor & \cref{prop:spectrum-free-torsion,prop:pair-from-G}\\
$T^*(n)$ & the ordinate operator of \cite{gonzalez25hp} & \cref{thm:operator-forms}\\
$\mathcal J$, $D$ & the joint index of level, colour and orientation, $\N\times\Z/3\times\Z/2$, and the deck step on it & \cref{thm:fibre-mu6,def:joint-index}\\
$\Oone$, $H$, $Y$, $W$ & the operator $\Oone=HY$; its modulus $H=\operatorname{diag}(m_0\golden^\sigma)$; its phase $Y=2\OmH W$, of order six; the orientation $W=Y^3$ & \cref{def:operators-J,thm:one-operator}\\
\multicolumn{3}{l}{\emph{The zeros of $L(s,\chiv)$}}\\
$\rho=\beta+i\gamma$ & a zero in the critical strip & \cref{ssec:topo-setting}\\
\Ecore & every zero has $\Re\rho=\tfrac12$ & \cref{ssec:topo-setting,ssec:core-forms}\\
$r(\rho)$ & $|\rho-1|/|\rho|$, the (Apollonius) coordinate of the zero & \cref{ssec:coordinate}\\
$t(\rho)$ & $r+1/r$, the trace of the zero; $t=L_k$ when $r=\golden^k$ & \cref{ssec:coordinate,ssec:two-halves}\\
$\operatorname{disc}(r)$ & $t^2-4$, the discriminant of the zero & \cref{prop:tower-disc-val}\\
$\ell(\rho)$ & $1-1/\rho$, the Li map; $|\ell(\rho)|=r(\rho)$ & \cref{prop:li-map}\\
\multicolumn{3}{l}{\emph{Counting the zeros, and the closure}}\\
$E(T)$, $E^{-1}$ & the envelope: the main term of the counting law, and its inverse & \cref{prop:positions}\\
$S(T)$ & the correction, or fluctuation: $t_n=E^{-1}(n-S(t_n))$ & \cref{prop:positions,prop:what-fluctuation-takes}\\
$T$ & a height on the critical line, $s=\sigma+iT$ & \cref{prop:positions}\\
$\psi(x)$, $\psi(x;20,a)$ & Chebyshev's function, and that of the class $a$ & \cref{eq:explicit-formula,conj:closure}\\
$R_a(x)$ & $\psi(x;20,a)-x/8$, the residual of the class $a$ around its structural share $x/8$ & \cref{conj:closure}\\
\multicolumn{3}{l}{\emph{The Weil form}}\\
$Q(f)$ & the Weil form of $L(s,\chiv)$ & \cref{ssec:weil-route}\\
$L$ & the length of the window, the support of the test functions & \cref{ssec:weil-route}\\
$\psi(z)$ & the digamma function $\Gamma'/\Gamma$ & \cref{ssec:pentagon-odd,ssec:weil-route}\\
$W(r)$ & $\log\frac5\pi+\Re\psi(\tfrac14+\tfrac{ir}2)$, the archimedean term & \cref{thm:second-window}\\
\multicolumn{3}{l}{\emph{Physics (\cref{sec:disc-cw}); see \cite{gonzalez26cw,gonzalez26ds}}}\\
$T^2_{\PCF}$, $T_K$ & the PCF torus $\C/\LPCF$ and its period $T_K=2\pi R_K=2\pi\log\golden$ (Lean \texttt{periodK}) & \cref{rmk:companion-interpretation,ssec:fund-formula}\\
$\Lambda_{\PCF}$, $\Mpcf$, $\epsz$ & the lattice $\Mpcf\Z[i]$ of the PCF torus, its scale $\Mpcf=\pi/\epsz$, and $\epsz=\log\golden/(6\sqrt3)$ & \cref{rmk:companion-interpretation}\\
$\tau_{\PCF}$ & the modulus of the torus, $\tau_{\PCF}=i$, the self-dual point & \cref{sec:disc-cw}\\
$\ell$, $G_N$, $c$ & the AdS radius, Newton's constant and the central charge $c=3\ell/(2G_N)$ (not the Li map $\ell(\rho)$) & \cref{sec:disc-cw}\\
$S_{\mathrm{BH}}$, $k_{\mathrm B}$ & the Bekenstein--Hawking entropy and Boltzmann's constant & \cref{sec:disc-cw}\\
$\alpha'$, $a_{2k-1,0}$ & the Regge slope and the Wilson coefficients of \cite{elvang26} & \cref{sec:disc-cw}\\
$\Omega(\tau)$ & the microstate $\tfrac12e^{i\tau\log\golden}$ of the ring (not the colour operator $\OmH$) & \cref{sec:microstate}\\
$S(\sigma)$, $N(\sigma)$ & the saturation entropy and the mode count of the level $\sigma$ (not $S(T)$, $N(T)$) & \cref{sec:tower-info}\\
$\varepsilon(\sigma)$, $z(\sigma)$, $\tau(\sigma)$ & the energy, the throat coordinate and the weld time of the level $\sigma$ (not $\tau(s)$) & \cref{sec:tower-info}\\
\bottomrule
\end{longtable}}


% =====================================================================

\section{Background: what $\Fone$ is and where it comes from}\label{sec:background}

\subsection{The odd zeta values: a problem older than the Riemann hypothesis}\label{sec:odd-history}

The sum of the reciprocals of the squares was posed by Mengoli in 1650~\cite{mengoli1650}, and it withstood the
attempts of the Bernoullis.  Euler found it in 1734 and read the result to the St Petersburg Academy in 1735:
$\zeta(2)=\pi^2/6$ \tier[\testref{test\_zeta\_two\_pi\_squared\_over\_six}]{N}.  In the same memoir he gave $\zeta(4),\dots,\zeta(12)$ as rational multiples of the
corresponding powers of $\pi$~\cite{euler1740}.  The general law came later:
$\zeta(2k)=(-1)^{k+1}B_{2k}(2\pi)^{2k}/\bigl(2\,(2k)!\bigr)$, with $B_{2k}$ the Bernoulli
numbers~\cite{lagarias13}.  That memoir already lists the sum of the reciprocal cubes among the series he had not
been able to sum, and the odd values never yielded to his methods.  In 1772 he reduced $\zeta(3)$ to an integral,
\[
  \tfrac78\,\zeta(3)=\tfrac{\pi^2}{4}\log2+2\int_0^{\pi/2}x\log\sin x\,dx,
\]
which he could not evaluate~\cite{euler1773,lagarias13}.

The zeta function grew out of these sums.  In 1737 Euler showed that the sum over the integers is a product over
the primes, $\sum_n n^{-s}=\prod_p(1-p^{-s})^{-1}$~\cite{euler1737}, and this tied the values to the distribution
of the primes.  In 1859 Riemann took $s$ complex, continued the function to the whole plane, proved its functional
equation, and stated the hypothesis that bears his name~\cite{riemann1859}.  The
question of the odd values is therefore more than a century older than the Riemann hypothesis, and it is still
open.  Ap\'ery proved that $\zeta(3)$ is irrational~\cite{apery79}.  Ball and Rivoal proved that infinitely many
odd values are irrational~\cite{ballrivoal01}, and Zudilin that one of $\zeta(5),\zeta(7),\zeta(9),\zeta(11)$
is~\cite{zudilin01}.  No closed form is known for any of them.

In physics, a free parameter is a quantity that a theory does not fix and that has to be supplied from outside it,
by measurement or by hand.  The odd values are expected to be free in the corresponding arithmetic sense.  The
prevailing expectation, a consequence of Grothendieck's period conjecture~\cite{grothendieck66,kontsevichzagier01},
is that $\pi,\zeta(3),\zeta(5),\dots$ are algebraically independent over $\Q$.  No polynomial relation with rational
coefficients would then hold among them, so no odd value would be fixed by $\pi$ or by the other odd values.  For
their motivic versions this is a theorem.  The motivic Galois group of mixed Tate motives over $\Z$ has a unipotent
part whose Lie algebra is free on one generator in each odd degree $2n+1\ge3$~\cite{delignegoncharov05}.  The
motivic odd zeta values are algebraically independent~\cite{brown12}, so each one is a new constant.  The same
freedom appears in physics: in the bootstrap of Elvang, Herderschee and Morales, maximal supersymmetry and
factorisation leave unfixed exactly the Wilson coefficients that carry $\zeta(2k+1)$~\cite{elvang26}.

Number theory and physics share a long history around $\zeta$, and the odd values are part of it.  It begins with
the values that the continuation assigns at the negative integers.  Where the series $1+2+3+\cdots$ diverges, the
continued function takes the value $\zeta(-1)=-\tfrac1{12}$, and physics uses that value: in the bosonic string the
zero-point energy of the transverse oscillators is $(D-2)\cdot\tfrac12\zeta(-1)$, and requiring it to be $-1$ fixes
the critical dimension $D=26$~\cite{polchinski98}.  Hawking introduced the regularisation of divergent sums and
determinants by zeta functions for path integrals in curved spacetime~\cite{hawking77}.  That regularisation is
also a standard method for computing Casimir energies~\cite{elizalde94}.  The odd values appear in at least three
places:
\begin{itemize}
  \item In quantum electrodynamics, $\zeta(3)$ enters the two-loop coefficient of the anomalous magnetic moment of
    the electron~\cite{petermann58,sommerfield57}, and $\zeta(3)$ and $\zeta(5)$ enter the three-loop
    coefficient~\cite{laporta96}.
  \item In string theory, they appear in the coefficients of the low-energy ($\alpha'$) expansion of the Veneziano
    amplitude~\cite{veneziano68}.  These are the same coefficients that the bootstrap of Elvang, Herderschee and
    Morales leaves unfixed.
  \item In the path integrals of quantum field theory, they arise as periods of Feynman diagrams: the counterterms
    of $\phi^4$ theory yield multiple zeta values, beginning with $\zeta(3)$, at high loop
    orders~\cite{broadhurstkreimer97}.
\end{itemize}

Through its zeros, $\zeta$ also meets the physics of random matrices: the zeros and the spectra of black holes in
low-dimensional holography share the same statistics, a meeting that \cref{sec:polya-hist,ssec:four-requirements}
take up.

The rest of this section examines two frameworks that stand where these fields meet.  The first is the
construction of the Hilbert--P\'olya operator, a self-adjoint operator whose eigenvalues would be the ordinates
$\gamma$ of the non-trivial zeros $\tfrac12+i\gamma$~\cite{polya1982letter,berry99,connes99}.  The second is
Manin's.  It includes the noncommutative tori of M-theory~\cite{connes98}, which he proposed as the geometry of real
multiplication~\cite{manin04}.  It also includes the frameworks for a geometry below $\mathrm{Spec}\,\Z$ with which
he approached the Riemann hypothesis~\cite{manin95}: one with Borger's $\Lambda$-rings~\cite{borger09,manin13}, and
one with cyclotomic polynomials~\cite{manin08}.  The first framework is where the circularity of the problem
becomes visible, and the next subsection begins there.

\subsection{The circularity problem and Hilbert--P\'olya}\label{sec:circularity}
\src{HP \S1.1 \texttt{ssec:polya-hist} (L90--L217), rewritten}

\subsubsection{Why the zeros cannot be referred to}

The Riemann hypothesis asks where the non-trivial zeros of $\zeta(s)$ lie.  The classical analytic theory obtains
information about the zeros from $\zeta$ itself~\cite{titchmarsh86}, and there is no circularity in that.  The
circularity meant in this paper concerns a construction: a structure meant to force the zeros onto the critical
line --- an operator whose spectrum they would be, or a geometry of whose Frobenius they would be the eigenvalues
--- that is defined from the zeros, or from the primes that the Euler product encodes, presupposes what it has to
explain.  Riemann's 1859 memoir~\cite{riemann1859} makes the duality between zeros and primes exact.  Through von
Mangoldt's explicit formula,
\begin{equation}\label{eq:explicit-formula}
  \psi(x) = x - \sum_\rho \frac{x^\rho}{\rho}
    - \log(2\pi) - \tfrac12\log(1-x^{-2}),
\end{equation}
the primes and the zeros are two descriptions of the same arithmetic content:
if $\mathrm{Re}(\rho)=\tfrac12$ for every $\rho$, each oscillatory term has
amplitude $x^{1/2}$ and $|\psi(x)-x| = O(x^{1/2}\log^2 x)$; a zero with
$\mathrm{Re}(\rho)>\tfrac12$ would create a dominant mode~\cite{titchmarsh86}.
A structure defined from the distribution of the primes is therefore defined from the zeros.  The two most
influential proposals for approaching the problem ask for a structure \emph{external} to $\zeta$, one whose own
properties force the zeros onto the critical line: a self-adjoint operator (Hilbert--P\'olya), treated in the rest
of this subsection, and a geometric Frobenius with a positivity (Weil, \cref{sec:weil}).  This paper follows the
second route through the $\Fone$ programme, and records at each step what is proved and what is assumed.

\subsubsection{The P\'olya operator}\label{sec:polya-hist}

In correspondence with Odlyzko in 1982, P\'olya recalled his conjecture of
1914: the Riemann hypothesis would follow if the non-trivial zeros were the
eigenvalues of a self-adjoint operator~\cite{polya1982letter,odlyzko87}.
Montgomery's pair correlation (1973) and Dyson's observation identified the
statistics of the spacings with those of the Gaussian Unitary
Ensemble~\cite{dyson62,montgomery73}, and Odlyzko's computations near the $10^{12}$-th zero~\cite{odlyzko87} and, later, near the $10^{22}$-nd~\cite{odlyzko01} confirmed the agreement.

Random matrices entered physics as a model for spectra whose Hamiltonian is not known: Wigner
replaced the unknown Hamiltonian of a heavy nucleus by a random symmetric matrix~\cite{wigner51},
and Bohigas, Giannoni and Schmit conjectured that a quantum system whose classical limit is chaotic
has the level statistics of the ensemble of its symmetry class --- GOE with time-reversal symmetry,
GUE without~\cite{bgs84}.  The zeros of $\zeta$ show the GUE statistics with no Hamiltonian known.
The broken time-reversal symmetry has an arithmetic reading: the phases $e^{-is\log p}$ contributed
by the Euler factors are mutually incoherent, because the logarithms of the primes are linearly
independent over $\mathbb{Q}$.  Berry and Keating read the zeros as the spectrum of the quantisation
of a chaotic system without time-reversal symmetry~\cite{berry99}; Keating and Snaith modelled $\zeta$
by characteristic polynomials of random unitary matrices~\cite{keating00}, and the Euler product
enters that model through an arithmetic factor, which the hybrid Euler--Hadamard product of Gonek,
Hughes and Keating incorporates~\cite{gonekhugheskeating07}.

The programme has resisted realisation for more than a century.  Classical
analytic methods address local properties of $\zeta$~\cite{hardy14,vinogradov47,titchmarsh86};
random matrix theory supplies the statistics but, as above, not the
operator~\cite{keating00,katz99}; Connes' noncommutative geometry reduces the
hypothesis to the validity of a trace formula~\cite{connes99,connes95};
Berry--Keating and Bender--Brody--M\"uller identify classes of candidate
operators~\cite{berry99,bender17}; Yakaboylu constructs a Hamiltonian on the half-line whose eigenvalues are fixed by the zeros, and whose self-adjointness is the contested point~\cite{yakaboylu24}.
Each approach meets either \emph{constructional circularity}---the operator
needs its own spectrum to be defined---or a limitation of \emph{scope}.
In the tower of this paper --- the powers $\golden^n$, with a group of six elements, the fibre, over each level $n$
(\cref{sec:objective}) --- a third limit can be stated exactly, and it is structural:
an operator that depends on the level alone, with everything in its functional calculus, cannot
separate the points of a fibre, and the operators that separate them do not commute with the
symmetry of the fibre (\cref{prop:level-operators,rmk:one-operator}).

\subsubsection{The requirements of an operator}\label{ssec:four-requirements}

An operator for the zeros has to meet a number of requirements, set out since P\'olya and added to by later authors;
the corpus makes four of them explicit: the operator must be non-circular, be self-adjoint, have a spectrum given
exactly, and not asymptotically or by calibration on data, and be unconditional (\LeanTally{four\_requirements}).  A name marked $^{\circ}$ is a tally: a
classification of the text recorded in Lean as an enumerated type whose count is checked by decision,
which certifies the count and not the classification.
The Primitive Structure~\cite{gonzalez25hp} distinguishes three distinct problems behind them.  The first is
circularity: an operator defined from the zeros, or from the primes, presupposes what it has to explain;
\cite{gonzalez25hp} sets beside it, as a term of comparison, a circularity of another kind, the diagonal argument of
Lawvere and Yanofsky.  The second is the binary against the third dimension: the real division algebras have
dimension $1$, $2$, $4$ or $8$, and none has dimension three.  The third is the statistics of the zeros: the unitary
class is one class of a classification that is itself binary, it is the class of quantum chaos without
time-reversal symmetry, and it is the class of the black holes of AdS/CFT~\cite{cotler17}, which is where the
statistics and the holographic duality meet.  The three are set out in \cref{ssec:intro-binary}, once the $\Fone$
programme and Manin's two programmes have been presented, together with the question they leave, how to pass
between dimensions; \cref{ssec:why-phi} then states by what means this paper meets each.

Of the framework's two candidates, the operator $\tilde H_{\mathrm{PCF}}$
meets non-circularity and self-adjointness, the latter proved in the text of \cite{gonzalez25hp}
(\LeanTally{hpps\_meets\_two}), and the operator $\OmH$ meets three: non-circularity, self-adjointness, in its
tower form (\cref{thm:operator-forms}), and exact spectrum, fixed by $\golden$ and not calibrated on data
(\cref{prop:omega-spectrum}; \LeanTally{omega\_meets\_three}).  The requirement that remains, unconditionality, is
the one \cref{ssec:core} identifies as the core, which the closure conjecture would supply (\cref{conj:closure});
\cref{ssec:disc-hp-ledger} returns to the three problems with the results of the paper.

\subsection{Weil's letter and the Rosetta stone}\label{sec:weil}
\src{HP \S1.1 (L165--L185), \texttt{ssec:weil} (L1457--L1530), rewritten}

\subsubsection{The letter of 1940 and the three columns}\label{ssec:weil-letter}

Writing to his sister Simone from prison in Rouen in 1940, Andr\'e Weil
described mathematics as a Rosetta stone with three columns: number fields,
function fields over finite fields, and Riemann surfaces~\cite{weil_letter40}.
A statement proved in one column suggests its translation into the others, and
the translation between domains is itself the method.  The Riemann hypothesis
is the case where the translation is known in one direction and missing in the
other: it is a theorem in the middle column and open in the first.

\subsubsection{RH for curves over finite fields}\label{ssec:weil-curves}

For a smooth projective curve $C$ of genus $g$ over $\mathbb{F}_q$, the
Frobenius $x\mapsto x^q$ acts on $H^1_{\text{\'et}}(C,\mathbb{Q}_\ell)$ with
eigenvalues $\alpha_1,\dots,\alpha_{2g}$, and
$\#C(\mathbb{F}_{q^m}) = 1 - \sum_i\alpha_i^m + q^m$.  The analogue of the
Riemann hypothesis is $|\alpha_i| = q^{1/2}$.  Hasse proved it for elliptic
curves~\cite{hasse34}, Weil for all curves~\cite{weil40,weil48}, and Deligne for
arbitrary varieties~\cite{deligne74}.  Weil's argument works on the surface
$C\times C$: the intersection form, with Riemann--Roch and Serre duality, gives
an inequality that, applied to the graph of the iterated Frobenius, yields
$|\alpha_i|\le q^{1/2}$; the functional equation gives the reverse
bound~\cite{milne15,oort14}.

\subsubsection{The Weil template}\label{ssec:weil-template}

Weil's proof has three ingredients, and they fix what any analogue over
$\mathbb{Z}$ must supply:
\begin{enumerate}
  \item[(a)] \textbf{A space on which Frobenius acts}: the cohomology of $C$, or
    the surface $C\times C$ with its intersection theory.
  \item[(b)] \textbf{Zeros as eigenvalues}: the zeros of the zeta function of $C$
    are read from the Frobenius eigenvalues.
  \item[(c)] \textbf{Positivity}: an inequality (Castelnuovo--Severi) that bounds
    the eigenvalues, combined with the functional equation.
\end{enumerate}
The aspiration to transport this template from $\mathbb{F}_q$ to $\mathbb{Z}$
is the $\Fone$ programme (\cref{sec:f1-programme}).  The Euler product
$\prod_p(1-p^{-s})^{-1}$, with one factor per prime, is where the transport
starts: L\'opez Pe\~na and Lorscheid write each local factor as the zeta
function of the residue field $\mathbb{F}_p$~\cite{lopezpena09torified}, and
Connes reads the zeros as the absorption spectrum of a trace formula on the
ad\`ele class space~\cite{connes99,connes16}.

% ---------------------------------------------------------------------
\subsection{The $\Fone$ programme}\label{sec:f1-programme}
\src{HP \texttt{ssec:convergence} (L395--L460), rewritten; literature checked in Consensus}

\subsubsection{Tits and the limit $q\to1$}

Tits observed in 1956 (published 1957) that combinatorial analogues of algebraic groups over
$\mathbb{F}_q$ persist as $q\to1$, where they become Weyl groups and
buildings~\cite{tits56}: a ``field with one element'' would be the base over
which these limits live.  The programme that grew from this remark has one
driving motivation: to read $\mathrm{Spec}\,\mathbb{Z}$ as a curve over $\Fone$,
with a completion analogous to a curve over a finite field, so that Weil's proof
(\cref{sec:weil}) could be translated~\cite{lorscheid12blueprints}.

\subsubsection{The obstruction the programme addresses}\label{ssec:f1-obstruction}

Weil's proof lives on the surface $C\times C$.  Over $\mathbb{Z}$ the analogous
product collapses: $\mathrm{Spec}\,\mathbb{Z}\times_{\mathbb{Z}}\mathrm{Spec}\,\mathbb{Z}
=\mathrm{Spec}\,\mathbb{Z}$ is one-dimensional, which is why the direct
translation of Weil's and Bombieri's proofs has so far been
unsuccessful~\cite{vanfrankenhuijsen08}.  The collapse has a precise source:
$\mathbb{F}_p[T]$ has a coefficient field, $\mathbb{Z}$ has none that is
visible~\cite{kurokawa03}.  The $\Fone$ programme is the search for that missing
base---an object \emph{below} $\mathrm{Spec}\,\mathbb{Z}$ over which the product
$\mathrm{Spec}\,\mathbb{Z}\times_{\Fone}\mathrm{Spec}\,\mathbb{Z}$ is a genuine
surface~\cite{haran07,toen09}.

\paragraph{Hilbert--P\'olya and $\Fone$.}
Over $\mathbb{F}_q$, Frobenius acting on $H^1$ plays the role of a
Hilbert--P\'olya operator: the zeros of $Z(C,q^{-s})$ are the $s$ with
$q^{s}=\alpha_i$, where $\alpha_1,\dots,\alpha_{2g}$ are its eigenvalues
(\cref{ssec:weil-curves}).  Deninger proposed an analogue over $\mathbb{Z}$: a
cohomological formalism in which the archimedean and the local factors are
regularised determinants~\cite{deninger91,deninger92}, and in which the zeros
would be the spectrum of an operator $\Theta$ on an $H^1$ of
$\overline{\mathrm{Spec}\,\mathbb{Z}}$~\cite{deninger94}, later sought as the
generator of a flow on a foliated space~\cite{deninger98}; the cohomology itself
has not been constructed.  Manin read that formalism as a cohomology over
$\Fone$~\cite{manin95}, so that the Frobenius of the sought base would fulfil, for
$\mathbb{Z}$, the function that Frobenius fulfils over $\mathbb{F}_q$.  Connes
realised the zeros spectrally on the ad\`ele class space~\cite{connes99}, and the
arithmetic site of Connes and Consani (\cref{ssec:f1-strategies}) relates that
realisation to $\Fone$.

\subsubsection{The strategies of the $\Fone$ programme}\label{ssec:f1-strategies}

\paragraph{Non-additive geometry and monoid schemes.}
Soul\'e gave the first framework for algebraic geometry over $\Fone$~\cite{soule03};
Deitmar's monoid schemes~\cite{deitmar05}, Haran's non-additive
geometry~\cite{haran07}, To\"en--Vaqui\'e's schemes below
$\mathrm{Spec}\,\mathbb{Z}$~\cite{toen09} and Lorscheid's
blueprints~\cite{lorscheid12blueprints} drop or weaken addition.  On a monoid
the $p$-th power $u\mapsto u^p$ is always an endomorphism, so Frobenius comes
for free at the multiplicative level.

\begin{definition}[$\Lambda$-monoid]\label{def:lambda-monoid}
A commutative monoid $M$ with commuting endomorphisms $\Frob_p$, one for each
prime, satisfying $\Frob_p\circ\Frob_q=\Frob_{pq}$; the basic example is
$\Frob_p(u)=u^p$.
\end{definition}

\paragraph{Borger: $\Lambda$-rings as descent data.}
Borger identified the structure that plays the role of Frobenius in
characteristic zero~\cite{borger09,borger11}.

\begin{definition}[$\Lambda$-ring, torsion-free case]\label{def:lambda-ring}
A torsion-free commutative ring $A$ with commuting \emph{ring} endomorphisms
$\Frob_p$, one for each prime $p$, such that $\Frob_p(a)\equiv a^p\pmod{pA}$
for all $a\in A$.
\end{definition}

A $\Lambda$-ring structure is descent data from $\mathrm{Spec}\,A$ to $\Fone$ and
globalises to a big topos over $\Fone$~\cite{borger09,borger11}; the Euler
product, with one factor per prime, is its analytic shadow.  The difference
between \cref{def:lambda-monoid} and \cref{def:lambda-ring} is additivity:
$\Frob_p$ must respect sums.  That difference is the subject of
\cref{sec:lambda-monoid,sec:lambda-ring}.  When the lifts and the congruence are required only at the primes of
a set $E$, the structure is a $\Lambda_{\Z,E}$-structure in the sense of Borger~\cite[\S1]{borger11}.  That is the sense of \emph{away
from $5$} in this paper, with $E$ the set of primes other than $5$: commuting ring endomorphisms $\Frob_p$ of $\Zphi$ for every prime $p\ne5$, each congruent
to the $p$-th power modulo $p$ (\cref{thm:frobenius-congruence,thm:descent-datum}), with no lift at $5$
(\cref{thm:obstruction-five}); the comparison of \cref{ssec:non-toric} writes it as a $\Lambda$-structure on
$\Zphi[\tfrac15]$, where $5$ is invertible and the congruence at $5$ is empty.
\begin{remark}
Ramification obstructs $\Lambda$-structures on rings of integers.  Borger and de Smit
show that a reduced $\Lambda$-ring finite flat over $\Z$ is a sub-$\Lambda$-ring of a
product of group rings $\Z[C]$ of finite cyclic groups, so that its rational algebra
lies in a product of cyclotomic fields and its Frobenius lifts are Frobenius
elements \cite[Thm.~0.1 and Cor.~0.3]{borgerdesmit08}; a Frobenius element exists only where the
extension is unramified.  For $\Zphi$ the failure at the ramified prime is proved
directly, by a nilpotent, in \cref{thm:obstruction-five}.
\end{remark}

\paragraph{Connes--Consani: the arithmetic site and three strategies.}
Connes and Consani connected $\Fone$ to noncommutative geometry and built the
arithmetic site~\cite{connes08,connes14,connes16}.  Connes distinguishes three
strategies at the level of the explicit formulas: an analytic one through the
trace formula, an algebro-geometric one that transposes Weil's proof through a
Riemann--Roch theorem in characteristic one, and a cohomological one through a
suitable Weil cohomology~\cite{connes15essay}.

\paragraph{Cyclotomic towers as $\bar{\mathbb{F}}_1$.}
Kapranov and Smirnov, and later Manin, take the cyclotomic tower in place of the algebraic closure of $\Fone$.
In the analogy between $\mathbb{Z}$ and $\mathbb{F}_q[T]$, the analogue of
$\bar{\mathbb{F}}_q[T]$ would be $\mathbb{Z}\otimes_{\Fone}\bar{\mathbb{F}}_1$, and the cyclotomic tower $\bigcup_r\mathbb{Z}[\zeta_{p^r}]$ of Iwasawa theory is taken in its place~\cite{kapranovsmirnov,manin08}; roots of unity play the role of the extensions of
$\Fone$~\cite{kapranovsmirnov}.

\subsubsection{Desiderata of the $\mathbb{F}_1$ programme}
\label{ssec:desiderata}

The literature has converged on a list of properties that any
candidate theory of $\mathbb{F}_1$-geometry should satisfy.  The
list is implicit in Tits~\cite{tits56}, articulated in
Manin~\cite{manin95}, formalised in Soul\'e~\cite{soule03},
and surveyed comparatively in
L\'opez Pe\~na--Lorscheid~\cite{lopezpena09mapping} and
Lorscheid~\cite{lorscheid18}.  We adopt the following nine-item
list as a working specification.

\begin{enumerate}[label=(D\arabic*),leftmargin=*,nosep]
\item $\mathrm{Spec}\,\mathbb{F}_1$ is a single point.
\label{D1}
\item $\mathrm{GL}_n(\mathbb{F}_1) = \Sigma_n$, the symmetric
group on~$n$ letters
(Tits~\cite{tits56}: combinatorial limit $q\to 1$ of
$\mathrm{GL}_n(\mathbb{F}_q)$).\label{D2}
\item $\mathbb{P}^n(\mathbb{F}_1)$ has exactly $n+1$ points
(buildings: combinatorial limit of $\mathbb{P}^n(\mathbb{F}_q)$).
\label{D3}
\item Counting function: a candidate $\mathbb{F}_1$-object~$X$ has
a counting function $N_X(q)$ which extends to $q\to 1$ either
polynomially (Soul\'e) or distributionally (Connes--Consani).
\label{D4}
\item $\mathrm{Spec}\,\mathbb{Z}$ admits a curve structure
over~$\mathbb{F}_1$, whose compactification
$\overline{\mathrm{Spec}\,\mathbb{Z}}$ is the
$\mathbb{F}_1$-arithmetic analogue of a smooth projective
curve.\label{D5}
\item There is a base-change functor
$-\otimes_{\mathbb{F}_1}\mathbb{Z}$ from
$\mathbb{F}_1$-schemes to ordinary schemes.\label{D6}
\item Toric varieties (split torus actions on themselves) admit
$\mathbb{F}_1$-models.  This is the minimal test case: any
candidate theory must reproduce the toric case.\label{D7}
\item The Riemann zeta function emerges as the
$\mathbb{F}_1$-counting zeta of
$\overline{\mathrm{Spec}\,\mathbb{Z}}$ (Manin~\cite{manin95}).
\label{D8}
\item There is a Weil-style proof of the Riemann Hypothesis
following from absolute geometry over~$\mathbb{F}_1$
(Deninger--Manin programme,~\cite{deninger91,deninger92,manin95}).
\label{D9}
\end{enumerate}

These nine items serve as criteria of the field, against which
any candidate construction is to be assessed.  We state them here
abstractly; the verification of which items the present
construction satisfies is given in \cref{ssec:position}, after
the relevant material has been developed, and the items each approach to $\Fone$ engages are
tabulated in \cref{app:comparison}.

\subsubsection{Soul\'e's varieties and the polynomial counting
condition}\label{ssec:soule}

The first formal proposal for an algebraic geometry over
$\mathbb{F}_1$ was made by Soul\'e~\cite{soule03}, refining
ideas of Kapranov--Smirnov~\cite{kapranovsmirnov}.  An
$\mathbb{F}_1$-variety is a finite-type scheme $X$
over~$\mathbb{Z}$ together with a polynomial $N_X\in\mathbb{Z}[x]$
such that
\begin{equation}\label{eq:soule-counting}
  \#X(\mathbb{F}_{p^n}) \;=\; N_X(p^n)
  \qquad\text{for all primes } p \text{ and all } n\ge 1.
\end{equation}
The polynomial $N_X$ is called the \emph{counting polynomial}
of~$X$, and~\eqref{eq:soule-counting} is referred to as
\emph{Soul\'e's condition}.  When the condition holds, the
Hasse--Weil generating function
$Z(X,T) = \exp\!\left(\sum_{r\ge 1} N_X(q^r)T^r/r\right)$
is rational, and the limit
\begin{equation}\label{eq:soule-zeta-2}
  \zeta_X(s) \;:=\; \lim_{q\to 1}\,Z(X,q^{-s})\,(q-1)^{N_X(1)}
\end{equation}
defines the $\mathbb{F}_1$-zeta of~$X$~\cite[\S6]{soule03}, in the form corrected in \cite[Eq.~(3)]{connesconsani10}.

The toric varieties satisfy~\eqref{eq:soule-counting} cleanly: for
$X = \mathbb{P}^n_\mathbb{Z}$, $N_X(q) = q^n + q^{n-1} +
\cdots + 1$, and $\zeta_X(s) = \prod_{k=0}^n (s-k)^{-1}$ is
the elementary $\mathbb{F}_1$-zeta of projective space.  However,
$\mathrm{Spec}\,\mathbb{Z}$ does not satisfy
Soul\'e's condition: $\zeta$ has infinitely many zeros, so it is not $\zeta_N$ for any
polynomial $N$, and the candidate counting function for
$\overline{\mathrm{Spec}\,\mathbb{Z}}$ falls outside~\eqref{eq:soule-counting}.  This is the limitation
that motivates the extension of \cref{ssec:cc}.

\begin{remark}[Algebraic extensions
$\mathbb{F}_{1^n}$]\label{rmk:F1n}
Following Kapranov--Smirnov~\cite{kapranovsmirnov} and
Soul\'e~\cite{soule03}, the algebraic extensions of
$\mathbb{F}_1$ are defined by base change as
\begin{equation}\label{eq:F1n}
  \mathbb{F}_{1^n}\otimes_{\mathbb{F}_1}\mathbb{Z}
  \;:=\; \mathbb{Z}[T]/(T^n - 1).
\end{equation}
This realises the cyclotomic side of $\mathbb{F}_1$-geometry:
adjoining roots of unity is the analogue of taking a finite
field extension of~$\mathbb{F}_q$.  The Connes--Consani
refinement~\cite{connesconsani11} singles out
$\mathbb{F}_{1^2}$ as a fundamental sub-case, with
generator $-1$ \textup{(}the unique non-trivial second root of
unity\textup{)} and balanced ternary multiplicative monoid
$\{-1,0,+1\}$.  In the present construction $-1$ appears as the norm of the
fundamental unit, $N(\golden)=-1$ \tier[\testref{test\_phi\_unit\_norm\_negative\_one}]{N} \textup{(\cref{sec:ring})}.
\end{remark}

\subsubsection{Connes--Consani's distributional extension}
\label{ssec:cc}

Connes and Consani~\cite{connesconsani10} extend Soul\'e's
framework to objects whose counting function is not polynomial
but distributional.  The technical step is the replacement of the
limit $q\to 1$ in~\eqref{eq:soule-zeta-2} by an integral against
a counting distribution.

The integral formula that makes this precise is \cref{lem:cc-integral}.  Applied to
$\overline{\mathrm{Spec}\,\mathbb{Z}}$, it gives the complete Riemann zeta function for the counting distribution
\[
  N_\mathbb{Q}(u) \;=\; u
  \;-\; \frac{d}{du}\!\left[\sum_{\rho}\mathrm{ord}(\rho)\,\frac{u^{\rho+1}}{\rho+1}\right]
  \;+\; 1,
\]
the sum running over the non-trivial zeros of $\zeta$ (\cref{thm:cc-zeta-Q}, equation~\eqref{eq:N-Q}).

The structural content of this formula is
the same shape as the Lefschetz formula for curves
over~$\mathbb{F}_q$:
\begin{equation}\label{eq:weil-shape}
  \#C(\mathbb{F}_{q^n}) \;=\; q^n
  \;-\; \sum_{i=1}^{2g}\alpha_i^n
  \;+\; 1,
\end{equation}
with the non-trivial zeros~$\rho$ playing the role of Frobenius
eigenvalues~$\alpha_i$.  Connes and Consani thereby identify
the zeros of~$\zeta(s)$ with the would-be eigenvalues of a
Frobenius operator on $\overline{\mathrm{Spec}\,\mathbb{Z}}/\mathbb{F}_1$,
and the genus appears as $N_\mathbb{Q}(1) = -\infty$ \textup{(}an
infinite-genus curve\textup{)}.

The Connes--Consani extension specifically addresses
desideratum~\ref{D8}.  It does not, by itself, exhibit the
descent data structurally: \eqref{eq:N-Q} requires the zeros of
$\zeta(s)$ as analytic input.  Combining a structural source for
the descent data (Borger's $\Lambda$-rings) with the
distributional output (Connes--Consani's $N(u)$) is the
combination the present paper provides for the case
$K=\mathbb{Q}(\sqrt{5})$.

\subsubsection{Other approaches and the position of this paper}
\label{ssec:landscape}

Several other frameworks for $\mathbb{F}_1$-geometry exist; we
summarise them briefly.  L\'opez Pe\~na and Lorscheid survey them~\cite{lopezpena09mapping}, and Manin and
Marcolli survey the homotopical approaches below $\mathrm{Spec}\,\mathbb{Z}$~\cite{maninmarcolli20}.  \Cref{app:comparison} compares them, with the present one, in a single
table: the objects each defines, how it counts, whether it has Frobenius lifts, its strategy for
the Riemann hypothesis, Soul\'e's condition, and the desiderata it engages.

\paragraph{Monoid schemes (Deitmar~\cite{deitmar05}).}
$\mathbb{F}_1$-schemes are defined by replacing rings with
commutative monoids and then taking $\mathrm{Spec}$ of monoids.
The functor $\mathrm{Spec}\,M\mapsto\mathrm{Spec}\,\mathbb{Z}[M]$
gives the base change to $\mathbb{Z}$.  Monoid schemes recover
toric varieties cleanly but do not handle non-toric arithmetic
structure.

\paragraph{Relative schemes
(To\"en--Vaqui\'e~\cite{toen09}).}
$\mathbb{F}_1$-schemes are defined as schemes relative to a
symmetric monoidal category~$\mathcal{C}$.  Choosing
$\mathcal{C}=\mathrm{CommMon}$ recovers Deitmar; choosing
$\mathcal{C}=\mathrm{Set}_*$ recovers a more rigid theory.

\paragraph{Blueprints (Lorscheid~\cite{lorscheid12blueprints}).}
Blueprints unify the previous frameworks: a blueprint is a
commutative semiring~$B$ together with a multiplicative
sub-monoid~$A\subseteq B$ generating~$B$ as a semiring.  The
resulting category subsumes Deitmar monoid schemes, Soul\'e
varieties, and To\"en--Vaqui\'e schemes.  Semiring schemes in the sense of To\"en--Vaqui\'e connect $\mathbb{F}_1$-geometry to tropical geometry: Giansiracusa and Giansiracusa use them to give tropical varieties a scheme structure~\cite{giansiracusa16}, a construction later recast with blueprints.

\paragraph{Generalised rings (Durov~\cite{durov07}).}
A generalised ring is an algebraic monad on the category of
sets.  Durov's framework treats Arakelov geometry uniformly with
$\mathbb{F}_1$-geometry: the place at infinity of~$\mathbb{Q}$
becomes an $\mathbb{F}_1$-object alongside the finite primes.

\paragraph{Torified varieties
(L\'opez Pe\~na--Lorscheid~\cite{lopezpena09torified}).}
A torified variety is a scheme $X$ over~$\mathbb{Z}$ together with
a decomposition into split tori.  These provide a class of
schemes admitting both Deitmar and Soul\'e structures.

\paragraph{$M_o$-schemes
(Connes--Consani~\cite{connesconsani10}).}
A refinement combining Soul\'e and Deitmar:
$\mathrm{Spec}$ of a multiplicative monoid endowed with a
distinguished system of homomorphisms to fields.  This is the
framework in which the distributional extension of
\cref{ssec:cc} is most naturally stated.

\paragraph{Position of this paper.}
Of these geometries, the one this paper uses is Borger's: a $\Lambda$-structure as
descent datum to $\Fone$, carried out on the cyclotomic strand, on the ring $\Zphi$ of
the pentagon.  It does not use monoid schemes, blueprints, torified varieties or the
homotopical approaches; where it meets them is recorded in
\cref{sec:discussion,ssec:non-toric}.
The construction is a Borger-style $\Lambda$-monoid on the tower
$\langle\golden\rangle$ of units, together with an Artin $\Lambda$-structure on
$\Zphi$ away from the ramified prime (\cref{sec:lambda-ring}); the ring
$\Rpcf=\mathbb{Z}[\golden,\golden^{-1},\tfrac{1}{2}]$ carries the first but not the
second (\cref{cor:localisation}).  The
counting function we attach to it
\textup{(\cref{thm:pcf-counting})} is distributional in the
Connes--Consani sense.  We do not develop a blueprint or
$M_o$-scheme structure for~$\langle\golden\rangle$; whether the
construction admits such a structure is left as an open problem
\textup{(\cref{ssec:future})}.  The toric machinery of
Deitmar applies cleanly to the trivial case $K=\mathbb{Q}$
\textup{(}corresponding to $\zeta_\mathbb{Q}=\zeta$\textup{)} but
not to the pentagonal case $K=\mathbb{Q}(\sqrt{5})$, which
requires the non-toric Borger descent.  This dichotomy is
articulated in \cref{ssec:non-toric}.




% ---------------------------------------------------------------------
\subsection{Manin's two programmes}\label{sec:manin}
\src{HP \texttt{ssec:convergence} (L438--L500), rewritten}

\subsubsection{\emph{Numbers as Functions}: Buium's $p$-derivations}

In \emph{Numbers as Functions}~\cite{manin13}, Manin surveys Buium's theory of
differential equations in the arithmetic direction~\cite{buium05}.  The prime
$p$ plays the role of a variable, and the $p$-derivation
\begin{equation}\label{eq:buium}
  \delta_p(a) = \frac{\Frob_p(a) - a^p}{p}
\end{equation}
is well defined exactly when $\Frob_p$ is a Frobenius lift; for $A=\mathbb{Z}$,
$\Frob_p=\mathrm{id}$ and $\delta_p(a)=(a-a^p)/p$ by Fermat.  Arithmetic jet
spaces follow.  Borger's $\Lambda$-rings are the global form of the same data; on $\Zphi$ these $p$-derivations exist at every $p\neq5$ (\cref{prop:delta-ring}).
Buium's route, through the $p$-adic completions, is
$p$-adic; its formalisation in Lean belongs to the
perfectoid line discussed in \cref{sec:discussion},
except for the order of $\golden$ modulo each prime, the clock of the prime (\cref{sec:objective}), which the
cyclotomic polynomials already read in finite arithmetic.

Scholze and Weinstein justify the name of the diamonds of $p$-adic geometry by an image: a geometric point is
like an impurity inside a gem, which gives it a colour; it is not seen directly, only through the reflections it
produces on the surface~\cite[\S9.2]{scholzeweinstein20}.  In that image the $p$-adic reading is the one from
inside the gem, prime by prime, and the reading by the regular polygons is the one on its surface.  This paper
takes both, the inside one as far as finite arithmetic reads it: the $p$-derivations, the polar form of $\golden$
at the primes and the prism at $5$ (\cref{prop:delta-ring,prop:polar,rmk:prism-five}); on the surface, the cyclotomic polynomials evaluated at
$\golden$; and it proves that at each prime the two read the same order (\cref{prop:clock-cyclotomic}).  In the
order of the construction the surface comes first, since the clock is read before any completion is taken.

\subsubsection{Real multiplication and noncommutative tori}

In \emph{Real Multiplication and Noncommutative Geometry}~\cite{manin04}, Manin
proposes the noncommutative tori $T_\theta$ with $\theta$ a real quadratic
irrationality as the geometric arena for the arithmetic of real quadratic
fields, with $\theta=\golden$ as the canonical case.  Connes, Douglas and
Schwarz showed that toroidal compactification in M-theory produces the same
tori~\cite{connes98}.  The arena is given; what it lacks is dynamics: Manin
proposes the space, not the flow---no operator whose spectrum is the zeros and
whose construction does not refer to the zeros (\cref{sec:circularity}).

\subsubsection{A synthesis and the problems it has to confront}\label{ssec:intro-binary}

Manin's two programmes give a requirement and an arena, and the Hilbert--P\'olya proposal asks for an operator
(\cref{sec:polya-hist}).  How can they be synthesised in one coherent framework?  A synthesis has to confront the
three problems that the Primitive Structure distinguishes (\cref{ssec:four-requirements}) --- the circularity of
the operator, the binary limit on the algebras over which it can be built, and the correspondence of that limit
with the statistics of the zeros --- and the question they leave: how to pass between dimensions.

\medskip
\paragraph{Circularity.}
The first problem is the one with which the account of the noncommutative tori ends: there is no operator whose spectrum is the zeros and
whose construction does not refer to the zeros (\cref{sec:polya-hist,sec:circularity}).  The Primitive
Structure~\cite{gonzalez25hp} compared this circularity with the self-reference of the fixed-point theorem, and
proposed that a setting in which that self-reference cannot be formed could exclude it.
Lawvere's fixed-point theorem~\cite{lawvere69} states that in a cartesian closed
category, if a map $A\to Y^{A}$ is weakly point-surjective, every endomorphism of
$Y$ has a fixed point; the diagonal arguments of Cantor, Russell, G\"odel and
Tarski's undefinability of truth are instances of its
contrapositive~\cite{lawvere69,yanofsky03}.  The theorem and its contrapositive
are formalised for the category of types, where point-surjectivity is
surjectivity (\tier[F1Lawvere.lawvere\_fixed\_point, F1Lawvere.no\_surjection\_of\_fixedPointFree]{P}).  There are
settings in which the self-reference the paradoxes rest on cannot be formed.  In a
stratified type theory, self-application $x(x)$ and self-membership $x\in x$ are
not well-typed, since an object of a given type does not take arguments of that
same type, so Russell's paradox and the liar cannot be stated.  In the theory of
real closed fields, which is complete and decidable~\cite{tarski1951},
$\mathbb{N}$ is not definable, so the theory cannot represent its own formulas
and G\"odel's diagonal construction does not arise.

\medskip
\paragraph{The binary limit and the statistics of the zeros.}
Beyond the fixed-point paradox of Lawvere and Yanofsky~\cite{lawvere69,yanofsky03}, binary structure limits the algebras over which the operator sought by the
Hilbert--P\'olya proposal and the $\Fone$ programme can be built (\cref{sec:polya-hist,ssec:f1-obstruction}).
The finite-dimensional associative division algebras over $\R$ are $\R$, $\mathbb{C}$ and
$\mathbb{H}$~\cite{frobenius1878}; without associativity the dimension is still $1$, $2$, $4$ or
$8$~\cite{bottmilnor58,kervaire58}, and the normed ones are these three and the
octonions~\cite{hurwitz1898,baez02}.  There is no real division algebra of dimension $3$.

The limit bears on the zeros through their statistics, which are what such an operator has to reproduce
(\cref{sec:polya-hist}).  Dyson's three symmetry classes of Hermitian random-matrix ensembles --- real
symmetric, complex Hermitian and quaternion self-dual --- correspond to the three associative division
algebras~\cite{dyson62threefold}, and the pair correlation of the zeros of $\zeta$ agrees with that of the
complex class, the Gaussian Unitary Ensemble: an operator reproducing it lies in the complex class, Hermitian over $\mathbb{C}$, of real dimension $2$, and not in the real or the quaternionic class.
The Primitive Structure and the Crystalline Worldsheet attach to the Euler product, on the other hand, a ternary
structure~\cite{gonzalez25hp,gonzalez26cw}: the tripartite object $\Omega=P\cdot C\cdot F$, with three eigenvalues
$\tfrac12\zeta_3^{\,k}$ permuted by $S_3$, whose parameter $\sigma=|\mathrm{rot}(S_3)|^2/|S_3|=\tfrac32$ is the value
$\zeta(2)/(\pi/3)^2$ of the Euler product at $s=2$ (\tier[zeta\_two\_over\_pi\_third\_sq]{P}).  An
operator with these three eigenvalues acts on three directions, and by the limit just stated three directions
cannot be the dimensions of a real division algebra: this is the binary against the third dimension.

\medskip
\paragraph{The passage between dimensions.}
In the holographic correspondence, a gravitational theory on a $(d+1)$-dimensional anti-de~Sitter bulk
is dual to a conformal field theory on its $d$-dimensional boundary~\cite{maldacena98}, and in low dimensions the boundary is an ensemble rather than a single theory: two-dimensional Jackiw--Teitelboim gravity is dual to a double-scaled Hermitian matrix integral~\cite{saad19}, of the
unitary class, the class of the statistics of the zeros and of quantum chaos without time-reversal symmetry
(\cref{sec:polya-hist}), and approaches to three-dimensional gravity also average over ensembles of two-dimensional conformal field theories~\cite{cooper26}.  In the error-correcting formulation of the
correspondence, the bulk is embedded in the boundary by an isometry~\cite{almheiri15,pastawski15}: the
map preserves the norm, so it is injective, and the passage loses a dimension while keeping the
information, which the Primitive Structure reads as a faithful forgetful functor.  This passage is developed in more
detail in the Primitive Structure and in the Crystalline Worldsheet~\cite{gonzalez25hp,gonzalez26cw}, and \cref{sec:disc-cw}
connects it with this construction.

The holographic correspondence has an arithmetic form, and part of it comes from Manin.  Manin and Marcolli showed that the
relation between hyperbolic geometry and Arakelov geometry at arithmetic infinity involves the same geometric data
as the Euclidean holography of black holes in $\mathrm{AdS}_3$: a Riemann surface uniformised by a Schottky group is
the conformal boundary of a hyperbolic handlebody, and the Green functions that give the Arakelov intersection
indices are geodesic lengths in that bulk.  At a prime the Bruhat--Tits tree replaces hyperbolic space, with ends
$\mathbb P^1(K)$ for $K$ a finite extension of $\Q_p$, and its quotient by a Schottky group, whose boundary is a
Mumford curve, is the non-archimedean analogue of the black hole; for the modular curves the boundary is
noncommutative~\cite{maninmarcolli01holo}.  Gubser, Knaute, Parikh, Samberg and Witaszczyk constructed a
$p$-adic AdS/CFT, with an unramified extension of $\Q_p$ as the boundary and the Bruhat--Tits tree as the bulk; the
correlators of a massive scalar in the bulk, written with the local zeta functions, are strikingly close to the real
ones~\cite{gubser17padic}.  Heydeman, Marcolli, Saberi and Stoica took the bulk to be the tree of
$\mathrm{PGL}(2,\Q_p)$, whose isometries are the conformal transformations of the boundary, so that the
discretisation breaks no symmetry; they proposed the tree as the natural setting of the tensor networks that model
bulk reconstruction by error-correcting codes, read its vertical direction as a scale of renormalisation, and found
the Ryu--Takayanagi formula there~\cite{heydeman18tensor}.  Marcolli built holographic codes on the trees and
buildings of Bruhat--Tits and on the Drinfeld symmetric spaces, in the line of Manin's reading of
$\mathrm{Spec}\,\Z$ as the arithmetic coordinate of physics, where the archimedean place has non-archimedean shadows
at the primes~\cite{marcolli18codes}; and Heydeman, Marcolli, Parikh and Saberi obtained, for networks of perfect
tensors on the tree and on its Schottky quotients, an entropy of the Ryu--Takayanagi form and one of the
Bekenstein--Hawking form for the $p$-adic black hole~\cite{heydeman21entropy}.

\subsubsection{The two programmes and this paper: why $\golden$}\label{ssec:why-phi}

This paper takes from the first programme the requirement (a Frobenius lift,
additive, at every prime) and from the second the object ($\golden$ and the
quadratic field $\mathbb{Q}(\sqrt5)$), and reaches the $p$-adic data by the cyclotomic route, the strand of the
$\Fone$ programme in which the cyclotomic tower stands for $\bar{\mathbb F}_1$ and the roots of unity for the extensions
of $\Fone$ (\cref{sec:f1-programme}):
$\golden$ lives in $\mathbb{Q}(\zeta_5)$, the pentagon (\cref{sec:choice}),
and the cyclotomic polynomials evaluated at $\golden$
read at each prime the order of the angle that the
$p$-adic route reads in the completion (the paragraph
\emph{The two readings of the $n$-gon} below).

In plain terms the construction is this.  The powers of $\golden$ form a line, the tower $\golden^n$, inside
the ring $\Zphi$; on that line the operator has its self-adjoint form, $H=\operatorname{diag}(m_0\golden^i)$, with $m_0>0$ a scale that cancels in every ratio of eigenvalues: it is the modulus of the operator (\cref{thm:one-operator}).  The line winds onto the plane, where the Euler products and the zeta functions
live.  Over each level of the line there is a fibre of six elements, the hexagon $\mu_6$, which carries the
ternary structure, and a third direction over the plane enters through the geometry of $E^3$
(\cref{thm:one-two-three}).  That third dimension is geometric and not algebraic: no real division algebra has
dimension three (\cref{ssec:intro-binary}), so it is reached by symmetry, through roots of unity, and by an
isometry, which keeps the information.  It stays in step with the arithmetic: at each prime the cyclotomic
polynomials evaluated at $\golden$ read at once the regular $n$-gon and the order of $\golden$ modulo the prime,
its clock.  Through the roots of unity of the pentagon the same structure determines the Euler factors and,
with them, the odd values of $\zeta$ (\cref{thm:odd-zeta}).  It reaches the zeros of $L(s,\chiv)$ as well: every
zero $\rho$ in the critical strip comes with its companion $1-\rho$, the coordinates $r(\rho)=|1-1/\rho|$ of the
two multiply to $1$, and, given the height of the first zero, a classical input, each lies in a window around $1$ in which the only power of $\golden$ is $\golden^0$ (\cref{prop:each-zero,prop:two-halves}).  What remains open is the exact position: that the coordinate is
exactly $1$, which is $\Re\rho=\tfrac12$, the statement \Ecore.  The paragraphs that follow take these points in
turn.

\begin{figure}[ht]
\centering\small
\begin{tikzcd}[column sep=1.4em,row sep=2.6em,ampersand replacement=\&]
\begin{array}{c}\text{hexagonal fibre}\\ \mu_6\text{ over each level}\end{array} \arrow[d]
\& \begin{array}{c}\text{self-adjoint operator}\\ H=\operatorname{diag}(m_0\golden^i)\end{array} \arrow[dl,dash,dotted]
\& \begin{array}{c}\text{pentagon}\\ \zeta_5^{\,a},\ 2\cos(\pi a/5)=\golden^{\pm1}\end{array} \arrow[d]
\\
\begin{array}{c}\text{line in the ring}\\ \golden^n\in\Zphi\end{array} \arrow[r,"{\text{lift at }p}"] \arrow[d,"{\Phi_n(\golden)}"']
\& \begin{array}{c}\text{Frobenius lift}\\ \Frob_p(a)\equiv a^p\ (p)\end{array} \arrow[r,dash,"{\text{same sign}}"] \arrow[d]
\& \begin{array}{c}\text{character}\\ \chiv(a)=\log_\golden|2\cos(\pi a/5)|\end{array} \arrow[d]
\\
\begin{array}{c}\text{clock of each prime}\\ \mathfrak p\mid\Phi_n(\golden)\end{array} \arrow[d,"{\text{product formula}}"']
\& \begin{array}{c}\text{Euler factors}\\ (1-\chiv(p)p^{-s})^{-1}(1-p^{-s})^{-1}\end{array} \arrow[l,dash,dotted] \arrow[r,"{\prod_p}"]
\& \begin{array}{c}\text{Dedekind zeta}\\ \zeta_K=\zeta\cdot L(s,\chiv)\end{array} \arrow[d] \arrow[dl]
\\
\begin{array}{c}\text{$n$-gon $=$ $p$-adic order}\\ \text{archimedean}=p\text{-adic}\end{array}
\& \begin{array}{c}\text{odd values}\\ \zeta(2k+1)=\zeta_K/L\end{array}
\& \begin{array}{c}\text{zeros of $L$}\\ r(\rho)\,r(1-\rho)=1\\ \text{window where the only power of }\golden\text{ is }\golden^0\end{array} \arrow[d,dashed,"{\Ecore}"]
\\
\& \& \begin{array}{c}\text{exact position}\\ r(\rho)=1\iff\Re\rho=\tfrac12\end{array}
\end{tikzcd}
\caption{The construction, with names and formulas (\cref{ssec:why-phi}).  Solid arrows are proved; the dashed arrow is the open
statement \Ecore; dotted lines join objects that live on the same line or the same prime.  The same diagram, with
the statement behind each arrow, is \cref{fig:construction-formulas}.}
\label{fig:construction-words}
\end{figure}


The choice of $\golden$ is not incidental.  A translation between the columns of
Weil's Rosetta stone (\cref{sec:weil}) needs an object that is at home in all of
them, and $\golden$ already is.  It is at once
\begin{itemize}
  \item a \textbf{magnitude and a ratio}: the positive root of $x^2=x+1$
    (\tier[phi\_sq]{P}), and its only positive root
    (\tier[golden\_fixed\_point\_unique]{P});
  \item a \textbf{geometry}: the diagonal of the regular pentagon,
    $\golden=2\cos(\pi/5)=1+\zeta_5+\zeta_5^4$
    (\tier[phi\_eq\_two\_cos\_pi\_fifth, phi\_cyclotomic\_inscription]{P});
  \item a \textbf{line and a function}: $x\mapsto\golden^x$, the golden tower, with the line of its integer
    levels $\golden^n$, which reaches
    the binary scale through $\golden^{\lambda_{\log}}=2$
    (\tier[mersenne\_bridge]{P});
  \item an \textbf{operator}: multiplication by $\golden$ on $\Zphi$, with matrix
    $\left(\begin{smallmatrix}0&1\\1&1\end{smallmatrix}\right)$, and the
    Frobenius lifts of \cref{sec:lambda-monoid,sec:lambda-ring};
  \item a \textbf{recursion}: $\golden^{n+1}=F_{n+1}\golden+F_n$, so the powers of
    $\golden$ are written in Fibonacci numbers (\tier[phi\_pow\_binet\_integral, fib\_pow\_aux]{P}).
\end{itemize}
$\golden$ therefore translates between domains by itself: a statement about the
number is a statement about the pentagon, about the recursion and about the
operator.  What makes it arithmetic is that it also reaches the primes, and it
does so through the work of two Colombian mathematicians.

\paragraph{Grisales Herrera: $\golden$ and the cyclotomic pentagon.}
Grisales Herrera's \emph{Simetría Áurea Prima}~\cite{grisales25}, his undergraduate thesis at the University of
Tolima, shows that the
golden ratio orders the primes $p>5$ into eight infinite families according to their
class modulo $20$, which their last two digits determine, through eight rational angles
invariant under the rotations of the regular pentagon in the complex plane; the eight
classes correspond to the zeros of the cyclotomic polynomial $\Phi_{20}$, and an
analogous relation holds for the Fibonacci numbers $F_{12(5m+j)}$.  It is the direct
link between $\golden$, the pentagon and the primes that this paper needs.  What
of it is proved here: the Galois homomorphism of the classification
(\tier[G\_hom]{P}); the pentagonal values themselves,
$|2\cos(\pi a/5)|=\golden^{\chiv(a)}$ (\tier[F1CW.norm\_cosine\_phi]{P}); and the
fact that $\chiv$ selects the Fibonacci number $F_p\bmod p$ for every prime
(\tier[additive\_spectrum\_governed\_by\_chi5]{P}).  The closed formula for
the classification is cited, not formalised (\cref{sec:grisales}).

\paragraph{Acevedo Agudelo: primes through rotations and the Mersenne primes.}
Acevedo Agudelo classifies the Mersenne primes completely by their residue
classes in the multiplicative group modulo $360$, splitting the search into four
disjoint subgroups and representing the classification geometrically by a
``Mersenne trapeze''~\cite{acevedo20}.  In a second work he characterises the
\emph{ova-angular rotations} of a prime, constructed from Dirichlet's theorem, and
derives geometric properties of the primes from them~\cite{acevedo21}; the
ova-angular classification has since been used to select Mersenne
candidates~\cite{acevedo23}.  Both read the primes through angular and cyclic
structure --- residue classes as rotations --- and not
through $p$-adic analysis, and that is where this paper
starts; the next paragraph shows that the angular
reading also gives the order of the $p$-adic angle.
The meeting point is the Mersenne bridge
$\golden^{\lambda_{\log}}=2$ and the binary tower inside
the golden one (\cref{sec:mersenne}).

\paragraph{The two readings of the $n$-gon.}
The eight classes of Grisales Herrera correspond to the
primitive $20$th roots of unity, the roots of
$\Phi_{20}$, so both classifications read the primes on
the circle.  Evaluated at $\golden$, the cyclotomic
polynomials read the primes at their own places.  At a
prime $\mathfrak p$ of $\Zphi$ the angle of $\golden$ ---
its Teichm\"uller part in the completion --- is a root of
unity, and its order is the \emph{clock} of
$\mathfrak p$; for $\mathfrak p$ prime to $n$, the clock
is $n$ exactly when $\mathfrak p$ divides
$\Phi_n(\golden)$.  For $n\ge2$ the norm
$N(\Phi_n(\golden))$ is a positive integer with two
readings.  On the circle it is the product of the
distances from $\golden$ and $-1/\golden$ to the
primitive vertices of the regular $n$-gon, of logarithm
$\phi(n)\log\golden$ up to a bounded term: the weight of
the tower at the level $n$.  At the primes it is the
product of the norms of the primes of clock $n$, and of
some primes dividing $n$.  The product formula equates
the two, so on each level of the clock the archimedean
and the $p$-adic readings of $\Phi_n(\golden)$ are
already glued; the statements, their proofs and the
values at the primes of the conductor are in the Methods
(\cref{prop:clock-cyclotomic}).  The polygons of the construction combine by laws that are proved: the
square and the pentagon in the conductor $20=\operatorname{lcm}(4,5)$, through the two components $\chi_4$ and
$\chiv$ of $G$ (\cref{prop:conductor,prop:G}); the hexagon in the fibre $\mu_6$ over each level
(\cref{prop:level-operators}); the orders $3$ and $4$ carried by the angle of $\golden$ at the primes $2$ and $3$,
those of the hexagonal and the Gauss lattices, of least common multiple $12$ (\cref{prop:polar}); and the three
characters of the Gauss, the Eisenstein and the golden lattices, which depend on a prime only through its class
modulo $60$ (\cref{prop:euler-tower}).  What the clock does not
give is the size of each prime: the clock divides $p-1$
or $2(p+1)$, and the quotient, the index, is not fixed
by it.  Ordering the primes by size instead of by clock
is where the fluctuation of the zeros is expected to
enter, as the closure conjecture of the Results
proposes.

\paragraph{The descent and the circularity.}
Borger's descent has in common with the settings in which the self-reference cannot be formed (\cref{ssec:intro-binary}) the direction of
construction: the object under study is not among the inputs.  On $\Zphi$ the
descent datum is indexed by the multiplicative monoid of the integers prime
to~$5$ and generated by its values at the primes, $\Frob_m\circ\Frob_n=\Frob_{mn}$
(\cref{prop:frob-compose}), with $\Frob_p(a)\equiv a^p\pmod p$
(\cref{thm:frobenius-congruence}); the character $\chiv$ is read off the
Frobenius (\cref{thm:sign-is-chi5,thm:chi5-from-frobenius}); and for $p\ne5$ the inverse local
factor of $\zeta_{\mathbb{Q}(\sqrt5)}=\zeta\cdot L(s,\chiv)$ at $t=p^{-s}$ is
$\det(1-tM_p)$, with $M_p$ the matrix of $\Frob_p$ (\cref{thm:frob-eigen}).
Neither $\zeta$ nor its zeros enter the definition of $\Frob_p$ or of $\chiv$; the
local factors at $p\ne5$ are their output.  With them come the Dedekind factorisation and the determination
of the odd values, $\zeta(2k+1)=\zeta_K(2k+1)/L(2k+1,\chiv)$, in which $L(s,\chiv)$ is a combination of the
Hurwitz values at the pentagonal points $a/5$ (\cref{thm:odd-zeta}): the roots of unity of the pentagon reach
the zeta functions by proved steps.  The primes enter as an index set, through their classes, and no
property of their distribution is used (\cref{sec:topological-proof}).  Nor can a Frobenius lift play the part of the
fixed-point-free map that the diagonal argument needs, which on the truth values $\{0,1\}$ is $\alpha(x)=1-x$
(\cref{ssec:disc-hp-ledger}): every ring endomorphism fixes the image
of $\mathbb{Z}$, so each lift has fixed points and the whole family has a common one
(\tier[F1Lawvere.ringHom\_fixes\_intCast, F1Lawvere.family\_common\_fixed\_point, F1Lawvere.\ensuremath{\psi}\_fixes\_int]{P}).  This is the property that
\cref{sec:polya-hist,ssec:four-requirements} require of a non-circular
construction, and it is the respect in which the descent of this paper answers,
within the $\Fone$ programme, the proposal of the Primitive Structure.

\paragraph{The passage between dimensions in this paper.}
This paper also passes between dimensions by an isometry, on the arithmetic side.  The ternary structure is obtained by symmetry and not by an algebra, through roots of unity: the ring $\Zphi$ of the descent lies in $\Z[\zeta_5]$, with $\golden=1+\zeta_5+\zeta_5^4$; the eight classes of the primes correspond to the primitive $20$th roots of unity, $a\mapsto\zeta_{20}^{\,a}$, the roots of $\Phi_{20}$~\cite{grisales25,acevedo21}; and the spectrum $\tfrac12\zeta_3^{\,k}$ of the operator
is carried by the cube roots of unity, among which,
read afterwards at the prime $2$, is the angle of
$\golden$ there: the level $3$ of the clock,
$\Phi_3(\golden)=2\golden^2$, of norm $4$.  In the fibre of the golden tower it is the colour, of order three, inside the deck group
$\mu_6=\Z[\zeta_3]^{\times}$: it is torsion, so no level reads it ($\operatorname{Hom}(\mu_6,\Z)=0$), its operator
$\tfrac12\operatorname{diag}(1,\zeta_3,\zeta_3^2)$ is normal and not self-adjoint (\cref{prop:level-operators}),
and the third direction enters through the geometry of $E^3$ (\cref{thm:one-two-three}).

The descent takes neither the distribution of the primes, nor the Euler product, nor the zeros as input, and the information is kept by an isometry: each eigenvalue $\tfrac12\zeta_3^{\,k}$ has modulus $\tfrac12$, so the triad has squared norm $\tfrac34$, and normalised it is a column $\mathcal V\colon\mathbb{C}\to\mathbb{C}^3$ whose three squared moduli are $\tfrac13$, so $\mathcal V^\dagger\mathcal V=1$ (\tier[F1PaperNew.triad\_isometry]{P}); it reaches dimension $3$ without a product in dimension $3$
\cite[the lemma \emph{Isometric embedding of the triad} and the proposition \emph{Two towers, one microstate}]{gonzalez26cw}
(\LeanCW{bulkBoundary\_isometry} \tier{P}).  A name marked $^{\dagger}$ belongs to the Lean
development of that paper.  The modulus $\tfrac12$ is the same at every level of the tower, whose levels are related by the composition law of the descent, $\Frob_p(\golden^n)=(\golden^n)^p$
(\LeanCW{tower\_frobenius\_common} \tier{P}).

Read as a code, $\mathcal V$ is an encoding whose code subspace, the image
of $\mathcal V\mathcal V^\dagger$, has dimension one: a state and not a space of logical states~\cite{almheiri15,pastawski15}.  What is
proved is this isometry of the triad; what the passage needs on the zeros, that their modes are modes of an
isometry, whose eigenvalues lie on the unit circle, is open (\cref{ssec:disc-binary}).

Of the $p$-adic form of the correspondence (\cref{ssec:intro-binary}), the arithmetic side
does not appear here as the Bruhat--Tits tree: at each prime it appears as the angle of $\golden$ on the $p$-adic circle, which is the
period of the Fibonacci dynamics on the $p$-torsion of the real torus, and the product formula glues the
archimedean and the $p$-adic readings (\cref{prop:two-faces}).  How the circle of each prime sits in its tree is not
treated here.  \Cref{ssec:disc-binary} states what of the binary zone the construction reaches and what remains open.

\paragraph{The companion papers and what this paper adds.}
This perspective---$\golden$ as the object that translates between domains and, as the preceding paragraphs show,
reaches the primes through the pentagon and the Mersenne tower---was explored in
two earlier works of this programme.  \emph{Simultaneous identities of the golden
ratio and Euler's identity}~\cite{gonzalez26mersenne} proves in Lean that $\golden$
satisfies simultaneously twelve identities, inclusions and identifications across
five canonical structures --- the complex rotor of Euler's identity, the real
trigonometric and real-multiplicative lines, the Mersenne numbers modulo $20$, and
the Galois group of $\Q(\zeta_{20})/\Q$ --- with the central clause
$3\golden^{\sigma(p)}=2^p$, $\sigma(p)=p\lng-\log_\golden3$.  The clause holds for every real $p$, with $\sigma(p)$
its only solution; what that paper proves is that its constituents are fixed independently of it --- the slope
$\lng=\log_\golden2$ is transcendental by Gelfond--Schneider~\cite{gelfond34,schneider35}, and the factor $3$ is the
Mersenne number $M_2$ and the number of classes modulo $20$ that the Mersenne primes occupy --- and it specialises
the clause to the $52$ known Mersenne prime exponents~\cite{gimps}.  The Primitive Structure~\cite{gonzalez25hp} constructs the operator side and the chain from $\Zst$
to $\mathrm{Re}(\rho)=\tfrac12$, stating its dependence on one classical axiom,
\texttt{hecke\_1920} (Hecke 1920), besides the geometric constants of its
construction, and on the hypothesis \texttt{bridge\_axiom} (the spectral mapping
of Gelfand and Dunford--Schwartz, applied through the Euler factorisation), and
deferring the construction of the arithmetic--spectral span to the present paper.

This paper extends the framework of both.  Relative to~\cite{gonzalez25hp} it
\begin{itemize}
  \item \emph{discharges Hecke}: the functional equation of $\Lambda(s,\chiv)$ and
    the companion zero are theorems (\tier[functionalEquation\_chi5\_discharged, hecke\_chi5\_derived]{P}), together with the conjugation symmetry
    $\overline{L(\bar s)}=L(s)$ (\tier[hconj\_discharged]{P});
  \item \emph{constructs the arithmetic--spectral span} with common source $G$ ---
    its two branches meeting at $\log2$ (\LeanInline{branches\_meet\_at\_log\_two},
    \LeanTally{the\_span\_is\_the\_identification}) --- and proves the pieces of the
    bridge that concern $G$ and $L(1,\chiv)$ (\LeanTally{the\_bridge\_pieces\_formalized});
  \item \emph{derives membership from the bridge}
    (\tier[the\_membership\_derived]{P}), leaving the bridge as the only
    substantive input (\cref{ssec:core});
  \item \emph{replaces the $\Lambda$-ring claim} by the Artin $\Lambda$-structure on
    $\Zphi$ away from $5$ and the obstruction at $5$ (\cref{sec:lambda-ring}).
\end{itemize}
What remains of \texttt{bridge\_axiom} is the identification of the two branches of
the span (\LeanInline{TheSpanIdentifiesThem}); every form of it is proved
equivalent to \Ecore\ (\cref{ssec:core-forms})\tier{P}, and along Weil's route it would follow from the closure conjecture (\cref{cor:closure}).



\Cref{sec:methods,sec:results} build the structure, carry it to the zeta functions and to the zeros, and collect the results; \cref{sec:disc-cw} records where the torus comes from.

% =====================================================================

\section{Methods}\label{sec:methods}

This section builds the construction and carries it to the zeros.
\Crefrange{sec:choice}{sec:grisales} build the objects: the pentagon and its conductor, the ring
$\Zphi$ with its units, the tower and its two registers, the Artin $\Lambda$-structure away from
$5$ with the $\delta$-ring and the polar form of $\golden$ at each prime, and the character $\chiv$
with the Golden Prime classification, the angular part of $\golden$ by class, the clock of
$\golden$ as cyclotomic polynomials, and the index of $\golden$ by layers and on the $n$-gon of each
prime.  No zeta
function and no zero appears there; everything is finite arithmetic over $\Zphi$, the
Teichm\"uller lift of \cref{ssec:delta-polar} and elementary real estimates on the clock and
the index aside.  \Crefrange{sec:dedekind}{ssec:base-phi} carry the
construction to the analytic objects.  The Euler factors are read off the Frobenius; the Euler
product is assembled from them as a colimit; the Dedekind factorisation follows, and with it the
pentagonal determination of the odd zeta values; the counting function places the construction
against Soul\'e and Connes--Consani; the Mersenne anchor turns the tower into a function; and from
the anchor every value of $\zeta$ and $L(\cdot,\chiv)$ is written in base $\golden$.
\Crefrange{sec:operator}{sec:topological-proof} carry it to the zeros, and \cref{sec:topological-proof} assembles the whole as a single thread from $\golden$ to the zeros.

% ---------------------------------------------------------------------
\subsection{The route: from the cyclotomic to the
$p$-adic, through the clock}\label{sec:choice}

\subsubsection{$\golden$ inscribed in the pentagon}

The golden ratio is a cyclotomic number, and the two descriptions of it used in
this paper --- as an algebraic number and as a length in the pentagon --- are the
same description.  Chebyshev's identity for $\cos5\theta$ at $\theta=\pi/5$ gives
$16c^5-20c^3+5c+1=0$ for $c=\cos(\pi/5)$, whose only positive root satisfies
$4c^2-2c-1=0$; the positive root of $x^2=x+1$ satisfies the same equation after
$x\mapsto x/2$, whence
\begin{equation}\label{eq:pentagon-id}
  \golden=2\cos(\pi/5)\qquad(\LeanInline{phi\_eq\_two\_cos\_pi\_fifth}\ \tierP).
\end{equation}
With $\zeta_5=e^{2\pi i/5}$ one has $\zeta_5+\zeta_5^4=2\cos(2\pi/5)=\golden-1$,
so
\begin{equation}\label{eq:cyclotomic-inscription}
  \golden=1+\zeta_5+\zeta_5^{4}
  \qquad(\LeanInline{phi\_cyclotomic\_inscription}\ \tierP),
\end{equation}
and $\Q(\sqrt5)$ is the fixed field of complex conjugation inside
$\Q(\zeta_5)$: the real quadratic field of $\golden$ is the real subfield of the
pentagonal cyclotomic field, of conductor $5$.  The identity itself rests on the
quintuple-angle formula $\cos5\theta=16\cos^5\theta-20\cos^3\theta+5\cos\theta$
(\LeanInline{cos\_five\_mul} \tierP).

The quarter turn lives one level away.  Adjoining $i$, of conductor $4$, gives
the conductor $20=\operatorname{lcm}(4,5)$, and the two factors are coprime
(\LeanInline{conductor\_is\_derived} \tierP), and in the twentieth roots the two levels
are visible: $\zeta_{20}^5=i$ and $\zeta_{20}^4=\zeta_5$ (\LeanInline{zeta20\_pow\_five\_eq\_I},
\LeanInline{zeta20\_pow\_four\_eq\_zeta5} \tierP).  Neither factor is chosen: the $4$
is the order of $i$, the $5$ is that of the pentagon, and the group
$\Zst$ has order $8$.

\begin{proposition}[The conductor is derived, and its four is one four]\label{prop:conductor}
$20=\operatorname{lcm}(4,5)$ with $\gcd(4,5)=1$ and $i^4=1\ne i^2$.  The factor $4$ of
the conductor, the period of the binary register $2^m\bmod 20$, and the order of
$(\Z/5)^\times$ are the same four: $2^2\equiv-1\pmod5$, exactly as $i^2=-1$, because
$5=(2+i)(2-i)$ splits in $\Z[i]$ and reduction modulo $2-i$ sends $i$ to $2$; and $5$
is the only prime with $p-1=4$.
\end{proposition}
\leanline{\LeanInline{conductor\_is\_derived}, \LeanInline{conductor\_is\_joint\_period},
\LeanStar{F1CW.binary\_period\_mod\_five},
\LeanStar{F1CW.five\_is\_the\_unique\_prime\_with\_four\_units} \tierP.}  The whole of \cref{sec:grisales} lives on that group, and
\cref{ssec:two-spectra} shows that each factor governs exactly one of the two
registers of the tower.


\subsubsection{Why the tower lives on the $5$-side}\label{ssec:crystallographic}

The two factors of the conductor differ in two structural ways, and both explain
why the arithmetic of \crefrange{sec:choice}{sec:clock} runs on the $5$-side.

The name is that of the classical crystallographic restriction.  A rotation of the plane that maps a
lattice $\Lambda\cong\Z^2$ to itself has an integer matrix in a basis of $\Lambda$, so its trace
$2\cos(2\pi/n)$ is an integer, and its order $n$ is $1$, $2$, $3$, $4$ or $6$: these are the orders a
crystal can have.  The conductor $20$ contains orders on both sides of that restriction.

\begin{proposition}[Crystallographic restriction]\label{prop:crystallographic}
$2\cos(2\pi/n)\in\Z$ exactly for $n\in\{1,2,3,4,6\}$: for $n\ge7$ the trace of
$\zeta_n$ lies strictly between $1$ and $2$, and the
pentagonal cases $n=5,10$ give $\golden^{-1}$ and $\golden$.  Accordingly the divisors of
$20$ split into the crystallographic orders $\{1,2,4\}$ and the non-crystallographic
orders $\{5,10,20\}$.  The non-crystallographic orders of least degree, those for which
$2\cos(2\pi/n)$ is quadratic, are the orders with Euler's function $\phi(n)=4$, namely
$5,8,10,12$; they give the three fields $\Q(\zeta_5)=\Q(\zeta_{10})$, $\Q(\zeta_8)$,
$\Q(\zeta_{12})$.
\end{proposition}
\leanline{\LeanInline{trZeta\_crystallographic\_values},
\LeanInline{trZeta\_not\_integer\_of\_ge\_seven},
\LeanInline{twenty\_divisor\_lattice\_splits},
\LeanInline{three\_conductors\_are\_the\_minimal\_failure} \tierP, each taking the
relevant cosine values as hypotheses; the hypotheses are proved and the lemmas
instantiated in \LeanStar{F1v2.cosineWindow\_holds},
\LeanStar{F1v2.trZeta\_not\_integer\_unconditional},
\LeanStar{F1v2.trZeta\_crystallographic\_unconditional},
\LeanStar{F1v2.three\_conductors\_unconditional}.}
\begin{proof}
For $n\ge3$, $2\cos(2\pi/n)=\zeta_n+\zeta_n^{-1}$ generates the real subfield of $\Q(\zeta_n)$, of degree
$\phi(n)/2$ over $\Q$; it is an algebraic integer, so it lies in $\Z$ exactly when that degree is $1$,
that is, when $\phi(n)\le2$, which is $n\in\{3,4,6\}$; for $n=1,2$ the values are $2$ and $-2$.  For
$n\ge7$, $0<2\pi/n<\pi/3$ gives $1<2\cos(2\pi/n)<2$.  The values at $n=5,10$ are
$2\cos(2\pi/5)=\golden-1=\golden^{-1}$ and $2\cos(\pi/5)=\golden$ \tier[\testref{test\_two\_cos\_pi\_over\_five\_equals\_phi}]{N} \eqref{eq:pentagon-id}.  The value is
quadratic exactly when $\phi(n)/2=2$, and $\phi(n)=4$ exactly for $n\in\{5,8,10,12\}$.
\end{proof}

\begin{proposition}[Unit ranks]\label{prop:unit-ranks}
In the biquadratic field $\Q(i,\sqrt5)$ of the conductor, whose three quadratic subfields
(its branches) are $\Q(i)$, $\Q(\sqrt{-5})$ and $\Q(\sqrt5)$, the two imaginary branches have unit
rank $0$ --- their unit groups are torsion --- while the real branch
$\Q(\sqrt5)$ has rank $1$; $\pm1$ is the only unit common to all three.
\end{proposition}
\leanline{\LeanInline{unit\_rank\_splits\_branches}, \LeanInline{branches\_are\_the\_unit\_groups},
\LeanInline{dirichlet\_ranks\_arithmetic} \tierP/\tierC.}

The two propositions say the same thing from two sides.  The $4$-side is
crystallographic and its units are finite, so it has no tower: the quarter turn
closes after four steps.  The $5$-side is not crystallographic and its unit group
is infinite, generated by $\golden$ (\cref{thm:units}), so it has a tower and
nothing else does.  The pentagon is excluded from the lattice exactly where it is
admitted into the arithmetic.

\subsubsection{Why this route is finitely formalizable}\label{ssec:finite-route}

Buium's route (\cref{sec:manin}) works in $p$-adic
completions, and its modern forms need perfectoid
geometry.  The route of this paper starts cyclotomic and
reaches the orders of the same angles, the clock,
without leaving finite arithmetic: the objects are
residue classes in $\Z/20$ and $\Z/p$, the quadratic
ring $\Zphi$, finitely many roots of unity, and the
integers $N(\Phi_n(\golden))$, whose prime factors not
dividing $n$ are those at which the angle of $\golden$
is a primitive $n$-th root of unity
(\cref{prop:clock-cyclotomic}).  Every statement of \crefrange{sec:choice}{sec:clock} is therefore either decidable by computation in a finite structure or provable by elementary algebra over $\Zphi$ --- the exceptions being the Teichm\"uller lift of
\cref{ssec:delta-polar}, computed at finite level and
taken as classical in the limit, the product formula of
\cref{prop:clock-cyclotomic}, which is classical, and the
archimedean estimates of \cref{prop:clock-cyclotomic}(iii),(vi)
and \cref{prop:index-ngon}(i)--(ii), elementary real
inequalities proved in Lean ---, which is why it is formalised in Lean~4 with
Mathlib as it stands, and why the tiers of \cref{ssec:protocol} can be assigned
statement by statement.

% ---------------------------------------------------------------------
\subsection{The ring $\Zphi$ and its arithmetic}\label{sec:ring}
\src{GN \S3; HP \texttt{ssec:fibonacci} (L1838)}

\subsubsection{Definition, domain, endomorphisms}

Let $f_\golden=X^2-X-1$ and
\[
  \Zphi=\Z[X]/(f_\golden),\qquad \omega=\bar X,\qquad \omega^2=\omega+1,
\]
with the real embedding $\iota\colon\Zphi\to\R$, $\omega\mapsto\golden$, which
carries the abstract tower onto the real one, $\iota(\omega^n)=\golden^n$
(\LeanStar{Z\ensuremath{\varphi}}, \LeanStar{\ensuremath{\omega}\_sq},
\LeanStar{\ensuremath{\iota}\_pow} \tierPnew).  Writing $\omega$ rather than
$\golden$ inside the ring is not pedantry: the statements of
\cref{sec:lambda-ring} are about the ring and hold in every quotient
$\Zphi/(p)$, where no real number is available.

\begin{proposition}\label{prop:domain}
$\Zphi$ is an integral domain.
\end{proposition}
\leanline{\LeanStar{f\ensuremath{\varphi}\_map\_two\_irreducible},
\LeanStar{f\ensuremath{\varphi}\_prime} \tierPnew.}
\begin{proof}
Modulo $2$, $f_\golden=X^2+X+1$ has no root in $\F_2$ ($0\mapsto1$,
$1\mapsto1$), so it is irreducible there; a monic quadratic that is irreducible
modulo a prime is irreducible over $\Q$, and being primitive it is irreducible in
$\Z[X]$.  Since $\Z[X]$ is a unique factorisation domain, $f_\golden$ is prime
and the quotient is a domain.
\end{proof}

\begin{proposition}[Only two endomorphisms]\label{prop:endo}
Every ring endomorphism $f$ of $\Zphi$ satisfies $f(\omega)=\omega$ or
$f(\omega)=1-\omega$.
\end{proposition}
\leanline{\LeanInline{FrobeniusUnico.endo\_dichotomy}, \LeanInline{FrobeniusUnico.hom\_ext} \tierP; \LeanStar{endo\_on\_\ensuremath{\omega}} \tierPnew.}
\begin{proof}
$f(\omega)^2=f(\omega^2)=f(\omega)+1$, so
$(f(\omega)-\omega)\bigl(f(\omega)-(1-\omega)\bigr)=f(\omega)^2-f(\omega)-1=0$,
and $\Zphi$ is a domain by \cref{prop:domain}.
\end{proof}

\Cref{prop:endo} is small and it decides the shape of everything below: whatever
the Frobenius lift of \cref{sec:lambda-ring} turns out to be, it has only two
possible values, and the congruence that defines it chooses between them
(\cref{thm:frobenius-congruence}).

The second endomorphism is the conjugation $\golden\mapsto1-\golden=-1/\golden$, which is the map
$x\mapsto-1/x$ on the roots because $N(\golden)=\golden(1-\golden)=-1$.  Self-duality under that map
is a property of the norm, not of $\golden$ alone, as the next statement shows.

\begin{proposition}[Self-duality under $x\mapsto-1/x$]\label{prop:selfdual}
A monic quadratic $q(x)=x^2+ax+b$ with real coefficients satisfies $q(-1/x)=-q(x)/x^2$ for every
$x\ne0$ if and only if $b=-1$, whatever $a$ is.  For $a\in\Z\setminus\{0\}$ the discriminant $a^2+4$ of
$x^2+ax-1$ is not the square of an integer, so the polynomial is irreducible over $\Q$; for $a=0$ it is
$(x-1)(x+1)$.  So the smallest $|a|$ with an irreducible self-dual polynomial is $1$, and $a=-1$ gives
$x^2-x-1$, the polynomial of $\golden$.
\end{proposition}
\leanline{\LeanInline{F1OZ.selfdual\_iff}, \LeanInline{F1OZ.selfdual\_disc\_not\_square},
\LeanInline{phi\_quadratic} \tierP.}
\begin{proof}
For $x\ne0$, $x^2q(-1/x)=1-ax+bx^2$ and $-q(x)=-x^2-ax-b$; they agree for every $x\ne0$ exactly when
$b=-1$, the two linear terms being equal for every $a$.  If $m^2=a^2+4$ with $a\ne0$, then
$|a|<|m|<|a|+2$, so $|m|=|a|+1$ and $2|a|+1=4$, which has no integer solution.
\end{proof}

\subsubsection{Additive growth: the Fibonacci coordinates}\label{ssec:fib-coords}

\begin{proposition}[Binet, integrally]\label{prop:binet}
In any commutative ring containing $\alpha$ with $\alpha^2=\alpha+1$,
\begin{equation}\label{eq:binet-integral}
  \alpha^{n+1}=F_{n+1}\,\alpha+F_n\qquad(n\ge0),
\end{equation}
where $F$ is the Fibonacci sequence.
\end{proposition}
\leanline{\LeanInline{fib\_pow\_aux}, \LeanInline{phi\_pow\_binet\_integral} \tierP.}
\begin{proof}
Induction.  For $n=0$ both sides are $\alpha$.  If
$\alpha^{n+1}=F_{n+1}\alpha+F_n$ then
$\alpha^{n+2}=F_{n+1}\alpha^2+F_n\alpha=F_{n+1}(\alpha+1)+F_n\alpha
=(F_{n+1}+F_n)\alpha+F_{n+1}=F_{n+2}\alpha+F_{n+1}$.
\end{proof}

\begin{proposition}[Factorisations and parities]\label{prop:factor-parity}
The two quadratics of the pentagon factor through $\golden$:
$x^2-x-1=(x-\golden)(x+\golden^{-1})$ and $x^2+x-1=(x-\golden^{-1})(x+\golden)$.  The
Lucas and Fibonacci numbers have the same parity, $L_k\equiv F_k\pmod2$, and
$L_k^2-5F_k^2\equiv0\pmod4$.
\end{proposition}
\leanline{\LeanInline{quadratic\_factor\_upper}, \LeanInline{quadratic\_factor\_lower},
\LeanInline{two\_sides\_of\_the\_same\_pair} \tierP.}
\begin{proof}
$\golden\cdot(-\golden^{-1})=-1$ and $\golden-\golden^{-1}=1$ give the first factorisation, and
$x\mapsto-x$ the second.  $L_k=F_{k-1}+F_{k+1}=F_k+2F_{k-1}$, so $L_k\equiv F_k\pmod2$; and
$L_k^2-5F_k^2=4(-1)^k$ (the norm of $2\golden^k=L_k+F_k\sqrt5$ is $4N(\golden)^k$).
\end{proof}

Nothing of $\R$ or of $\sqrt5$ enters \eqref{eq:binet-integral}: only the
quadratic relation and the recurrence.  That is why the same coordinates govern
the tower inside $\Zphi$, inside $\OK$, and modulo every prime, and it is the
identity through which every statement about $F_p\bmod p$ below is a statement
about $\omega^p$.  Multiplication in the tower reads as addition of Fibonacci
coordinates and adjacent powers carry,
$\golden^{k+1}+\golden^{k}=\golden^{k+2}$
(\LeanInline{multiplicative\_reads\_as\_additive},
\LeanInline{carry\_in\_base\_phi} \tierP).

\subsubsection{The unit, the norm, and the inert prime $2$}\label{ssec:unit-norm}

Multiplication by $\omega$ in the basis $(1,\omega)$ has matrix
$M=\left(\begin{smallmatrix}0&1\\1&1\end{smallmatrix}\right)$
(\LeanInline{M\_is\_mul\_by\_phi} \tierP), whose determinant is the norm,
\[
  \det M=N(\omega)=\omega(1-\omega)=\omega-\omega^2=-1
  \qquad(\LeanInline{det\_M\_eq\_norm\_phi},\ \LeanInline{znorm\_phi}\ \tierP),
\]
so $\omega$ is a unit of norm $-1$ and the tower $\{\omega^n\}_{n\in\Z}$ is a
group.  The same matrix is the Fibonacci recurrence, which is
\eqref{eq:binet-integral} read as a linear map; its characteristic polynomial is
$f_\golden$.

The prime $2$ is inert: $\Zphi/(2)=\F_2[X]/(X^2+X+1)\cong\F_4$
(\LeanInline{two\_is\_inert} \tierP), and there the image of $\omega$ has order
three, being a generator of $\F_4^\times$
(\LeanInline{f4phi\_order\_three} \tierP).  This is the first instance of the dichotomy that organises the whole section: at $2$ the polynomial stays irreducible,
so the $p$-th power map moves the two roots, and \cref{sec:lambda-ring} records that as
$\Frob_2=\sigma$.  It is also the
prime whose Frobenius carries the positivity of \cref{thm:second-window}.


\subsubsection{The ring of integers and the unit group}\label{ssec:units}

Two facts about $\Zphi$ decide how much of the tower is canonical, and both are proved here
rather than taken as classical inputs.

\begin{proposition}[The ring of integers]\label{prop:OK}
$\OK=\Zphi$: an element $(a+b\sqrt5)/2$ with $a\equiv b\pmod2$ is
$c+b\golden$ with $c=(a-b)/2\in\Z$, and conversely $c+b\golden$ has
$a=2c+b\equiv b\pmod2$.
\end{proposition}
\leanline{\LeanInline{quadratic\_ring\_is\_Z\_phi}, \LeanInline{the\_rule\_is\_proved} \tierP\ for the
rewriting; the description of $\OK$ as $\{(a+b\sqrt5)/2 : a\equiv b\ (2)\}$ is classical
\tierC.  For the elements the paper uses --- roots of $x^2-tx+1$ with $t\in\Z$ --- the
inclusion in $\Zphi$ is proved directly (\LeanStar{F1v2.arrowB\_Q}).}

\begin{theorem}[The tower is the unit group, without Dirichlet]\label{thm:units}
There is no unit in the interval $(1,\golden)$; consequently every unit $u>0$ is
$u=\golden^{k}$ for some $k\in\Z$.
\end{theorem}
\leanline{\LeanInline{no\_unit\_below\_phi\_universal}, \LeanInline{unit\_is\_a\_power\_of\_phi},
\LeanInline{unit\_structure\_no\_dirichlet} \tierP, the last two for an abstract unit
predicate with three structural inputs; for the units of $\Zphi$ the inputs are
proved and the descent instantiated in \LeanStar{F1v2.floorLogφ0},
\LeanStar{F1v2.noUnitBelowφ0}, \LeanStar{F1v2.unitsZφ\_group}, \LeanStar{F1v2.arrowB}.}
\begin{proof}
Let $u\in(1,\golden)$ be a unit with conjugate $\bar u$ and trace
$a=u+\bar u\in\Z$.  If $N(u)=+1$ then $\bar u=1/u$ and
$a=u+1/u\in(2,\sqrt5)$, an interval with no integer; if $N(u)=-1$ then
$\bar u=-1/u$ and $a=u-1/u\in(0,1)$, again with no integer.  The two cases exhaust the norms $\pm1$.  For the
second statement, put $k=\lfloor\log u/\log\golden\rfloor$; then
$u/\golden^{k}\in[1,\golden)$ is a unit, hence equals $1$.
\end{proof}

\Cref{thm:units} is what makes the tower canonical rather than chosen: the
multiplicative group $\langle\golden\rangle$ of \cref{thm:lambda-monoid} is the
full group of positive units of $\OK$, and $\golden$ is its fundamental unit.  It
is proved from the minimality argument above and the group law, with no appeal to
Dirichlet's unit theorem (\LeanInline{unit\_structure\_is\_derivable} \tierP, instantiated
for $\Zphi$ in \LeanStar{F1v2.arrowB}), so the whole of \crefrange{sec:choice}{sec:clock}
rests on nothing classical at this point, other than
the exceptions named in \cref{ssec:finite-route}.

\subsubsection{The tripartite generator $\Omega = P\cdot C\cdot F$}
\label{ssec:tripartite}



The letters $P, C, F$ are used in the rest of the paper in two readings, one additive and one
multiplicative.  This subsection defines both: the first from the Fibonacci recurrence, together
with three characterisations of $\golden$ as a fixed point, and the second from the triad
$\tfrac12\zeta_3^{\,k}$ that the inert prime $2$ gives (\cref{ssec:unit-norm}): its modulus is $|2|_2=\tfrac12$ and
its angle has the order, three, of $\omega$ in $\F_4^\times$.

\begin{definition}[Fibonacci recurrence and the labelling
$P, C, F$]\label{def:fib-PCF-labels}
The Fibonacci sequence $F\colon\mathbb{N}\to\mathbb{N}$ is defined by
\begin{equation}\label{eq:fib-recurrence}
  F_0 = 0,\qquad F_1 = 1,\qquad F_{n+2} = F_{n+1} + F_n.
\end{equation}
At each index $n\ge 1$, the recurrence binds three consecutive
terms; we label them
\begin{equation}\label{eq:PCF-labels}
  P_n = F_{n-1}\;\text{(\emph{Precedent})},\quad
  C_n = F_{n}\;\text{(\emph{Current})},\quad
  F_n^{\!*} = F_{n+1}\;\text{(\emph{Forthcoming})},
\end{equation}
so that \cref{eq:fib-recurrence} becomes
$F_n^{\!*} = C_n + P_n$ at every level.  The same letters
$P, C, F$ are used in the multiplicative tripartite structure
$\Omega = P\cdot C\cdot F$ of \cref{def:tripartite-omega} below; the
acronym \emph{PCF} (Precedent--Current--Forthcoming) labels both
the additive recurrence and the multiplicative factorisation.
\end{definition}

\begin{remark}[Asymptotic ratio]\label{rmk:fib-ratio-phi}
The classical limit
\[
  \lim_{n\to\infty}\frac{C_n}{P_n}
  \;=\;\lim_{n\to\infty}\frac{F_n}{F_{n-1}}
  \;=\;\golden
\]
binds the recurrence to $\golden$.  In the present paper we use
the algebraic identity $\golden^2 = \golden + 1$ as the
ring-level shadow of the recurrence (\cref{sec:ring}), and we
record below the three formal self-references of $\golden$ that
characterise it as a fixed point.
\end{remark}

\begin{lemma}[Continued-fraction self-reference]\label{lem:phi-cf}
\begin{equation}\label{eq:phi-cf}
  \golden \;=\; 1 + \frac{1}{\golden}.
\end{equation}
The continued fraction $1+1/(1+1/(1+\cdots))$ converges to $\golden$.
\end{lemma}
\leanline{\LeanInline{phi\_continued\_fraction} \tierP\ for the convergence; the identity is \lean{φ\_self\_ref}.}

\begin{proof}
From $\golden^2 = \golden + 1$ \textup{(\cref{sec:ring})} divide both
sides by $\golden > 0$: $\golden = 1 + 1/\golden$.  The $n$-th convergent of the continued
fraction is $F_{n+1}/F_n$, by induction from $1+1/(F_{n+1}/F_n)=F_{n+2}/F_{n+1}$; by
\cref{prop:binet} with $\alpha=\golden$ and its conjugate $-\golden^{-1}$,
$F_n=(\golden^n-(-\golden)^{-n})/\sqrt5$, so $F_{n+1}/F_n\to\golden$.
\end{proof}

\begin{lemma}[Nested-radical self-reference]\label{lem:phi-nr}
\begin{equation}\label{eq:phi-nr}
  \golden \;=\; \sqrt{1+\golden}.
\end{equation}
The nested radical $\sqrt{1+\sqrt{1+\cdots}}$ converges to $\golden$.
\end{lemma}
\leanline{\LeanInline{phi\_nested\_radical} \tierP\ for the convergence; the identity is \lean{φ\_nested\_ref}.}

\begin{proof}
From $\golden^2 = \golden + 1 = 1 + \golden$ and $\golden > 0$,
take the positive square root.  The radical is the limit of $x_0=1$, $x_{n+1}=g(x_n)$ with
$g(x)=\sqrt{1+x}$; $g$ maps $[1,2]$ into itself and $|g'(x)|=1/(2\sqrt{1+x})\le1/(2\sqrt2)<1$ there, so
$g$ is a contraction of $[1,2]$ and $x_n$ converges to its unique fixed point, which is $\golden$.
\end{proof}

\begin{proposition}[Golden fixed-point uniqueness]\label{prop:golden-unique}
For $x>0$, the equation $x^2 = x+1$ has the unique solution
$x = \golden$.
\end{proposition}
\leanline{\LeanInline{golden\_fixed\_point\_unique} \tierP.}

\begin{proof}
The polynomial $x^2 - x - 1 = (x-\golden)(x+\golden^{-1})$ since
$\golden^{-1} = \golden - 1$ (from $\golden^2=\golden+1$, divide by
$\golden$).  The root $x = -\golden^{-1} < 0$ violates $x>0$, hence
$x = \golden$ is the unique positive solution.
\end{proof}

\begin{definition}[Tripartite components]\label{def:tripartite-omega}
Let $\zeta_3=e^{2\pi i/3}$ and $\lambda_k=|2|_2\,\zeta_3^{\,k}=\tfrac12\zeta_3^{\,k}$, $k=0,1,2$: the triad of the
prime $2$.  It inscribes an
equilateral triangle in the circle $|z|=R=\tfrac12$.  The three components of the tripartite
generator $\Omega=P\cdot C\cdot F$ are the three lengths of that triangle, normalised by $R$ where
the length is a radius:
\begin{equation}\label{eq:omega-tripartite}
  \norm{C}=\frac{|\lambda_0|}{R},\qquad
  \norm{F}=|\lambda_1-\lambda_0|,\qquad
  \norm{P}=\frac{R\cos(\pi/3)}{\norm{F}/2},
\end{equation}
the current radius, the chord between adjacent eigenvalues, and the ratio of the apothem to the
half-side; and $\norm{\Omega}=\norm{P}\cdot\norm{C}\cdot\norm{F}$.
\end{definition}

\begin{proposition}[The component norms, and their product]\label{thm:norm-Omega-half}
\begin{equation}\label{eq:PCF-component-norms}
  \norm{P}=\frac1{\sqrt3}=\tan\frac\pi6,\qquad \norm{C}=1,\qquad \norm{F}=\frac{\sqrt3}2=\sin\frac\pi3,
\end{equation}
\begin{equation}\label{eq:norm-Omega-half}
  \norm{P}\cdot\norm{C}\cdot\norm{F}\;=\;\tfrac{1}{2},
\end{equation}
which is also the common modulus $|\lambda_k|=|2|_2$ of the triad.
\end{proposition}
\leanline{\LeanInline{norm\_Omega\_is\_half} \tierP; \LeanCW{normP\_eq\_tan}, \LeanCW{normF\_eq\_cos},
\LeanCW{Omega\_eigenvalues} \tierP \cite{gonzalez26cw}.}

\begin{proof}
$\norm{C}=|\lambda_0|/R=1$.  $\norm{F}=\tfrac12|\zeta_3-1|=\sin(\pi/3)=\tfrac{\sqrt3}2$.  The apothem is
$R\cos(\pi/3)=\tfrac14$ and the half-side $\tfrac{\sqrt3}4$, so $\norm{P}=1/\sqrt3$.  The $\sqrt3$ of
$\norm{F}$ cancels that of $\norm{P}$, leaving $\tfrac12$; the cancellation is specific to the
equilateral triangle.
\end{proof}

\begin{remark}[The acronym \emph{PCF}, justified]\label{rmk:PCF-acronym}
\Cref{def:fib-PCF-labels,def:tripartite-omega} and \cref{thm:norm-Omega-half}
together justify the acronym.  The three letters $P, C, F$
designate simultaneously
\begin{enumerate}[label=\textup{(\roman*)},nosep]
\item the three consecutive terms $(F_{n-1}, F_n, F_{n+1})$ of the
Fibonacci recurrence \textup{(\cref{eq:fib-recurrence})}, and
\item the three components of the multiplicative tripartite
factorisation $\Omega = P\cdot C\cdot F$ with
$\norm{P}\cdot\norm{C}\cdot\norm{F} = 1/2$
\textup{(\cref{thm:norm-Omega-half})}.
\end{enumerate}
The two readings are coherent: at every Fibonacci level the
recurrence binds three consecutive terms (additively); at the
level of the triad the generator $\Omega$ factors into three
components (multiplicatively); both readings produce the same
tripartite labelling.  The framework as a whole is the
\emph{Precedent--Current--Forthcoming Framework}, abbreviated PCF.
\end{remark}

% ---------------------------------------------------------------------
\subsection{The $\Lambda$-monoid on the golden tower}\label{sec:lambda-monoid}

Borger's datum asks for one Frobenius lift per prime
(\cref{def:lambda-ring}).  On the multiplicative side the lifts are immediate,
and they are the ones the tower of rings of \cref{ssec:rel-borger} carries between
its levels.

\begin{proposition}[Monoidal descent data]\label{thm:lambda-monoid}
Let $\langle\golden\rangle=\{\golden^n:n\in\Z\}$ and let $X_{\PCF}=(\Z,+)$ be its
exponent monoid, with $\mathrm{realise}(n)=\golden^n$.  The maps
$\Frob_p(n)=pn$ satisfy
\[
  \mathrm{realise}(\Frob_pn)=\mathrm{realise}(n)^p,\qquad
  \Frob_p\circ\Frob_q=\Frob_{pq},\qquad \Frob_1=\mathrm{id},
\]
they commute, they are the unique maps with these properties, and they induce
maps $n\mapsto pn$ on $\Z/20$.  So $p\mapsto\Frob_p$ is a monoid homomorphism
$(\N_{\ge1},\times)\to\operatorname{End}\langle\golden\rangle$, which is the
$\Lambda$-structure of Borger in multiplicative form.
\end{proposition}
\leanline{\LeanInline{F1\_monoidal\_descent\_data}, \LeanInline{psi\_realise\_pow},
\LeanInline{psi\_comp}, \LeanInline{psi\_unique}, \LeanInline{X\_PCF} \tierP.}
\begin{proof}
On a generator, $\Frob_p(\Frob_q(\golden^n))=\golden^{pqn}$, and every element of
the tower is such a generator; uniqueness is the observation that a monoid
endomorphism of $\Z$ is multiplication by its value at $1$.
\end{proof}

\begin{remark}[The monoidal lift is not a ring endomorphism]\label{rmk:monoid-not-ring}
The ring $\Z[\golden,\golden^{-1},\tfrac12]$ with $\Frob_p(\golden)=\golden^p$ is not a
$\Lambda$-ring.  The levels of the tower carry rings
(\cite[\S2.8]{gonzalez26cw}), so the question is about the map, not about the
object: $\golden\mapsto\golden^p$ does not respect $\golden^2=\golden+1$, since
$\Frob_p(\golden+1)=\Frob_p(\golden^2)=\golden^{2p}\ne\golden^p+1$ for $p\ge2$.
It is an endomorphism of the multiplicative monoid, which moves between levels,
and not a ring endomorphism inside one
(\LeanInline{add\_destroyed}, \LeanStar{F1CW.psi\_not\_additive} \tierP).
The additive half lives inside the level: it is carried by the additive register
below, and \cref{sec:lambda-ring} produces the ring endomorphism the descent
datum asks for.
\end{remark}

\subsubsection{The two registers}\label{ssec:two-spectra}

The conductor $20=4\cdot5$ splits the arithmetic of the tower into two registers,
and each factor governs one of them and is blind to the other.

\begin{theorem}[The two registers of the tower]\label{thm:two-spectra}
\leavevmode
\begin{enumerate}
  \item \emph{Multiplicative register.}  The binary register $2^m\bmod20$ is
    periodic of period $4$ from $m=2$ on.
  \item \emph{Additive register.}  For a prime $p\ne5$ let $\varepsilon_p=+1$ if $x^2-x-1$ has a
    root modulo $p$ and $\varepsilon_p=-1$ if it has none.  Then
    \[
      F_p\equiv\varepsilon_p\pmod p,\qquad
      F_{p-1}\equiv\frac{1-\varepsilon_p}2\pmod p,
    \]
    and, equivalently, $\omega^p\equiv\omega\pmod{p\Zphi}$ if $\varepsilon_p=+1$ and
    $\omega^p\equiv1-\omega\pmod{p\Zphi}$ if $\varepsilon_p=-1$.
\end{enumerate}
\end{theorem}
\leanline{(1) \LeanInline{binary\_register\_has\_period\_four} \tierP; (2)
\LeanInline{fib\_split\_case}, \LeanInline{fib\_inert\_case} \tierP,
\LeanInline{FrobeniusUnico.omega\_pow\_mod}, \LeanInline{FrobeniusUnico.fib\_eq\_sign},
\LeanInline{FrobeniusUnico.fib\_pred\_eq\_sign} \tierP.  No statement in (2) uses a character.}

The word \emph{spectrum} is not used for these registers: in this paper it is kept for the
spectrum of an operator.  The two factors of the conductor govern the two registers, and
\cref{prop:two-registers-G} below shows that each is blind to the other.

\begin{proof}[Proof of (2)]
For $p=2$ the polynomial is $x^2+x+1$ over $\F_2$, with no root; $F_2=1\equiv-1$, $F_1=1$, and
$\omega^2=\omega+1\equiv1-\omega\pmod2$.  Let $p$ be odd, $p\ne5$.  By \eqref{eq:binet-integral} in
any commutative ring, $\alpha^p=F_p\alpha+F_{p-1}$ for every root $\alpha$ of $x^2-x-1$.

If $x^2-x-1$ has a root $\alpha\in\F_p$, the other root is $\alpha'=1-\alpha$, distinct because
$(\alpha-\alpha')^2=5\ne0$ in $\F_p$.  Fermat gives $\alpha^p=\alpha$ and
$\alpha'^p=\alpha'$, and subtracting the two instances of Binet,
$\alpha-\alpha'=F_p(\alpha-\alpha')$, hence $F_p=1$; substituting back, $F_{p-1}=0$.

If it has no root, it is irreducible over $\F_p$ and its roots live in
$\F_{p^2}$, where the Frobenius $x\mapsto x^p$ exchanges them: $\alpha^p=1-\alpha$.
Then Binet reads $1-\alpha=F_p\alpha+F_{p-1}$, and
$\{1,\alpha\}$ is an $\F_p$-basis of $\F_{p^2}$, so $F_p=-1$ and $F_{p-1}=1$.

Finally $\omega^p=F_p\omega+F_{p-1}$ in $\Zphi$, and reading the coefficients modulo $p$ gives the
statement about $\omega^p$.
\end{proof}

The reading that matters for \cref{sec:lambda-ring} is this.  The lift of
\cref{thm:lambda-monoid} carries the multiplicative register, between levels.
The additive register is already the ring-level Frobenius, read modulo $p$: by
\eqref{eq:binet-integral}, $\omega^p=F_p\omega+F_{p-1}$, and what
\cref{thm:two-spectra}(2) computes, from the roots of $x^2-x-1$ modulo $p$ alone, is that this is
$\omega$ or $1-\omega$.  \Cref{sec:lambda-ring} packages that as a ring endomorphism, proves the
congruence for every element rather than for the generator, and proves that the endomorphism is
unique.

\subsubsection{Two parametrisations of the same tower}\label{ssec:two-params}

The tower carries two indices.  The construction of \cite[\S2.8]{gonzalez26cw}
indexes its levels by an integer $\sigma$, scaling by $\golden$ from one level to the
next; \cref{thm:lambda-monoid} gives the monoid $\langle\golden\rangle$ the lifts
$\Frob_p(\golden^n)=\golden^{pn}$ indexed by the primes.  This subsection makes
explicit that the two are parametrisations of one object.

\begin{definition}[The golden tower]\label{def:golden-tower}
The \emph{golden tower} is the multiplicative group
\[
  T_\golden \;=\; \{\golden^x : x\in\mathbb{R}\} \;=\; \mathbb{R}_{>0}
\]
read in the exponent coordinate $x$,
with its distinguished sublattice
$T_\mathbb{Z} = \{\golden^n : n\in\mathbb{Z}\}$ and prime sublocus
$T_\mathbb{P} = \{\golden^p : p\in\mathbb{P}\}$.
The relation
\begin{equation}\label{eq:tower-extended}
  \golden^n = F_n\,\golden + F_{n-1}\qquad (n\in\mathbb{Z}),
\end{equation}
with $F_n$ the Fibonacci sequence extended to negative indices by
$F_{-n} = (-1)^{n+1} F_n$, places every element of $T_\mathbb{Z}$ in the free
$\mathbb{Z}$-module $\Zphi=\Z\oplus\Z\golden$ of rank~$2$, with coordinates
$(F_{n-1},F_n)$.
\end{definition}

\begin{proposition}[Two coordinates on $T_\golden$]\label{prop:two-coords}
The logarithm $\log_\golden\colon T_\golden\xrightarrow{\sim}\mathbb{R}$,
$\golden^x\mapsto x$, identifies the golden tower with the additive
real line and induces two distinguished embeddings:
\begin{enumerate}[label=\textup{(\alph*)}]
\item \textup{(metric, geometric)} $\mathbb{Z}\hookrightarrow T_\mathbb{Z}$,
  $\sigma\mapsto\golden^\sigma$.  Operation: $\sigma\mapsto\sigma+1$,
  realising scale auto-similarity.
\item \textup{(arithmetic, Frobenius)} $\mathbb{P}\hookrightarrow T_\mathbb{P}$, $p\mapsto\golden^p$,
  the image of $\golden$ under the lift $\Frob_p$; the monoid $(\N_{\ge1},\times)$, generated by the
  primes, acts on $T_\mathbb{Z}$ by $\golden^n\mapsto\golden^{pn}$.
  Operation: composition, realising the multiplicative composition of primes.
\end{enumerate}
The prime $p$ enters the tower in two ways, which the paper keeps apart: as the exponent of its lift,
$\golden^p=\Frob_p(\golden)$, and as a point of $T_\golden$, $p=\golden^{\log_\golden p}$, at the coordinate
$\log_\golden p$.
The two coordinates meet at the binary anchor
\textup{(\LeanInline{mersenne\_bridge})}:
\begin{equation}\label{eq:mersenne-loc}
  \golden^{\lambda_{\log}} = 2,\qquad
  \lambda_{\log} = \frac{\ln 2}{\ln\golden} \approx 1.4404.
\end{equation}
The prime $p=2$ thus occupies the irrational position
$\sigma = \lambda_{\log}$ in the integer tower, between
levels $\sigma=1$ and $\sigma=2$, and analogously
$\log_\golden p\notin\mathbb{Q}$ for every prime $p$.
\end{proposition}
\leanline{\LeanInline{mersenne\_bridge},
\LeanInline{lambda\_is\_not\_an\_integer} \tierP\ for
$p=2$; for every rational $r>0$, $r\ne1$, so for
every prime,
\LeanStar{M12F1Circle.logphi\_rat\_irrational}
\tierPnew.}

\begin{proof}
Direct from the definition of $\log_\golden$ as the inverse of the
exponential bijection $\mathbb{R}\to\mathbb{R}_{>0}$,
$x\mapsto\golden^x$.  The Mersenne bridge gives $\lambda_{\log}$
explicitly.  If $\log_\golden p=a/b$ with $b>0$, then $\golden^a=p^b$; but
$\golden^a$ is a unit of $\Zphi$ and $p^b$, of norm $p^{2b}\ne\pm1$, is not.  So
$\log_\golden p$ is irrational, and then, $\golden$ and $p$ being algebraic, it is
transcendental by the Gelfond--Schneider theorem~\cite{gelfond34,schneider35}.
\end{proof}

\begin{proposition}[Tower compatibility]\label{thm:tower-compat}
Under the metric coordinate, the Frobenius lift $\Frob_p$ acts as
multiplication by~$p$ on the exponent:
\begin{equation}\label{eq:psi-as-multiplication}
  \Frob_p(\golden^\sigma) \;=\; (\golden^\sigma)^p \;=\; \golden^{p\sigma}
  \qquad\text{for all } \sigma\in\mathbb{Z},\; p\in\mathbb{P}.
\end{equation}
Consequently, the diagram
\begin{equation}\label{eq:tower-diagram}
\begin{tikzcd}
T_\mathbb{Z} \arrow[r, "{\Frob_p}"] \arrow[d, "\log_\golden"'] &
T_\mathbb{Z} \arrow[d, "\log_\golden"] \\
\mathbb{Z} \arrow[r, "\cdot p"'] & \mathbb{Z}
\end{tikzcd}
\end{equation}
commutes, and the two parametrisations of \cref{prop:two-coords}
generate compatible monoid actions on~$T_\golden$.
\end{proposition}
\leanline{\LeanInline{psi\_add}, \LeanInline{psi\_comp} \tierP.}

\begin{proof}
Equation~\eqref{eq:psi-as-multiplication} is immediate from
$\Frob_p(\golden) = \golden^p$ and the multiplicative homomorphism
property of the Frobenius lift.  Commutativity of
\eqref{eq:tower-diagram} reduces to
$\log_\golden(\golden^{p\sigma}) = p\sigma = p\cdot\log_\golden(\golden^\sigma)$.
\end{proof}

\begin{remark}[Frobenius composition as exponent
multiplication]\label{cor:frob-as-mul}
The composition law $\Frob_p\circ\Frob_q = \Frob_{pq}$
\textup{(\cref{thm:lambda-monoid})} corresponds, under
$\log_\golden$, to the trivial identity $p\cdot q = pq$ on the
exponent.  In particular, the $\Lambda$-monoid axiom
$\Frob_p\circ\Frob_q = \Frob_{pq}$ is a tautology in the metric
coordinate; the non-trivial content lies in the
Frobenius equality $\Frob_p(\golden^n) = (\golden^n)^p$, which is
read off the lattice $T_\mathbb{Z}\subset T_\golden$ via the
exponent map $\log_\golden$ (\LeanInline{psi\_comp} \tierP).
\end{remark}

\begin{remark}[The descent perspective]\label{rmk:descent-perspective}
The map $\sigma\mapsto\golden^\sigma$ converts addition to
multiplication: the additive group $(\mathbb{Z},+)$ on the metric
side becomes the multiplicative monoid $(T_\mathbb{Z},\cdot)$ on
the arithmetic side.  This is precisely the direction of
$\mathbb{F}_1$-descent in the sense of
Borger~\cite{borger09}~--- forgetting addition while retaining
multiplication.  The golden tower~$T_\golden$ is the laboratory in
which both structures coexist before the descent: the additive
ladder $\sigma\in\mathbb{Z}$ and the multiplicative monoid
$\mathbb{P}\hookrightarrow T_\mathbb{P}$ are visible
simultaneously, and the descent functor is realised explicitly as
the choice to read $T_\mathbb{Z}$ multiplicatively.
\end{remark}

\begin{remark}[The primes as interstitial points]\label{rmk:primes-interstitial}
Under $\log_\golden$, the integer levels $\sigma\in\mathbb{Z}$ form
a uniformly spaced sublattice of~$\mathbb{R}$, while the primes
$\{\log_\golden p : p\in\mathbb{P}\}$ form a sparse irrational
sequence.  The first few values are
\[
  \log_\golden 2 \approx 1.4404,\quad
  \log_\golden 3 \approx 2.2830,\quad
  \log_\golden 5 \approx 3.3446,\quad
  \log_\golden 7 \approx 4.0438,
\]
each strictly between two consecutive integer levels.  The
arithmetic content of the Frobenius lifts thus inhabits the
\emph{interstices} of the geometric tower: prime exponents are
the irrational points where the metric ladder
$(\mathbb{Z},+)$ and the arithmetic monoid
$(\mathbb{P},\cdot)$ meet inside~$T_\golden$.
\end{remark}

\begin{remark}[The two readings of the same
tower]\label{rmk:companion-interpretation}
The metric coordinate~$\sigma$ indexes successive levels of the
lattice $\LPCF^{(\sigma)} = \golden^\sigma\LPCF$
\textup{(\cref{ssec:rel-borger})}, where $\LPCF=\Mpcf\,\Z[i]$ is the square lattice of the PCF torus
$T^2_{\PCF}=\C/\LPCF$, of modular parameter $\tau=i$, with $\Mpcf=\pi/\epsz$ and $\epsz=\log\golden/(6\sqrt3)$
\cite[\S2.4, \S3.4]{gonzalez26cw}.  The arithmetic coordinate~$p$
indexes Frobenius lifts that classify primes via
$\chi_5(p)\equiv F_p\pmod p$ and produce the local Euler factors of
$\zeta_{\mathbb{Q}(\sqrt{5})}$ (\cref{sec:colimit}).  The two
parametrisations are not independent constructions: they are two
functorial maps from~$\mathbb{R}$ into a single multiplicative
object~$T_\golden$, identified via the bijection $\log_\golden$.
\end{remark}

\begin{center}
\small
\begin{tabular}{@{}lll@{}}
\toprule
Aspect & Metric coordinate & Arithmetic coordinate \\
\midrule
Index & $\sigma\in\mathbb{Z}$ & $p\in\mathbb{P}$ \\
Operation
  & $\sigma\mapsto\sigma+1$
  & $n\mapsto pn$, composed as $p,q\mapsto pq$ \\
Generator
  & $\golden$ \;(scale step)
  & $\Frob_p$ \;(Frobenius lift) \\
Conserved invariant
  & $|\Omega| = 1/2$ at every level
  & $\Frob_p\circ\Frob_q = \Frob_{pq}$ \\
Base point
  & $\varepsilon_0 = \ln\golden/(6\sqrt{3})$ \cite{gonzalez26cw}
  & $\golden$ itself \\
Bridge to the other
  & $\sigma_p := \log_\golden p \notin\mathbb{Z}$
  & Mersenne $\sigma_2 = \lambda_{\log}$ \\
$\mathbb{F}_1$-role
  & geometric stratification
  & descent data on $\langle\golden\rangle$ \\
\bottomrule
\end{tabular}
\end{center}

The synthesis is concise.  $T_\golden$ is a single object; the two
parametrisations are two functorial maps into~$T_\golden$, related
by the bijection $\log_\golden$.  Every claim about the metric
tower translates, via this bijection, to a claim about the
arithmetic tower, and conversely; the Mersenne bridge supplies
the precise location at which a chosen prime---namely $p=2$---becomes
a distinguished interstitial point of the integer tower.



\begin{proposition}[The tower as the image of the exponents]\label{prop:tower-image}
The map $n\mapsto\golden^n$ from $\Z$ is injective and strictly increasing; the lift
$\golden^{pn}=(\golden^n)^p$ is multiplicative, $\golden^{p(m+n)}=\golden^{pm}\golden^{pn}$,
and respects inverses, $\golden^{p(-n)}=(\golden^{pn})^{-1}$; and the tower is stable
under every power, $(\golden^k)^p=\golden^{j}$ with $j=pk$.
\end{proposition}
\leanline{\LeanInline{realise\_injective}, \LeanInline{realise\_strictMono},
\LeanInline{psi\_multiplicative\_on\_the\_monoid}, \LeanInline{psi\_respects\_inverses},
\LeanInline{family\_stable\_unconditional}, \LeanInline{family\_is\_equivariant} \tierP.}

The tower's additive law characterises its base: $\golden^\sigma+\golden^{\sigma+1}=\golden^{\sigma+2}$
for every real $\sigma$, and a positive $r$ with $r^\sigma+r^{\sigma+1}=r^{\sigma+2}$ for
all real $\sigma$ is $\golden$ (\LeanInline{adjacent\_closure\_real},
\LeanInline{golden\_unique\_closure} \tierP).

The base scale of the metric tower of \cite{gonzalez26cw} is fixed by the same
generator: $\log\golden=6\sqrt3\,\varepsilon_0=6\sqrt3\,\pi/M_{\PCF}$
(\LeanInline{eps0\_relates\_log\_and\_geometry}, \LeanInline{component\_one\_denominator} \tierP).

\subsubsection{Borger's two axioms, verified}\label{ssec:borger-criteria}


\begin{remark}[Borger's two axioms on the unit monoid
$\langle\golden\rangle$]\label{rmk:borger-criteria-pcf}
The system $\{\Frob_p\}_{p\in\mathbb{P}}$ on the multiplicative monoid
$\langle\golden\rangle\subset\Rpcf^\times$ defined by
$\Frob_p(\golden^n) = \golden^{pn}$ satisfies the two canonical
Borger axioms of \cref{def:lambda-ring}:
\begin{enumerate}[label=\textup{(B\arabic*)},nosep]
\item \emph{Commuting family:}
  $\Frob_p\circ\Frob_q = \Frob_{pq}$
  \textup{(\cref{prop:Psi-functorial};
  \LeanInline{psi\_comp})}.
\item \emph{Frobenius lift property on the monoid:}
  $\Frob_p(\golden^n) = (\golden^n)^p$ holds exactly by definition
  \textup{(\cref{thm:lambda-monoid}; \LeanInline{psi\_realise\_pow})}.
\end{enumerate}
On the generator, $\Frob_p(\golden)=\golden^p=F_p\golden+F_{p-1}$ by \cref{prop:binet}:
the monoid gives the coordinates of $\golden^p$.  That $F_p\equiv\chi_5(p)\pmod p$ is a
statement about the ring $\Zphi$ reduced modulo $p$, and it is proved there
(\cref{thm:two-spectra}\textup{(2)} and \cref{thm:chi5-from-frobenius}), not by the
monoid alone.  The monoidal lift is not additive (\cref{rmk:monoid-not-ring};
\LeanInline{add\_destroyed}, \LeanStar{F1CW.psi\_not\_additive} \tierP): the
additive half of Borger's datum is the Artin $\Lambda$-structure of \cref{sec:lambda-ring},
which lives inside each level.  The collected statement for the monoid is
\LeanInline{F1\_monoidal\_descent\_data} \tierP, with the three standard axioms only.
\end{remark}

\begin{remark}[Why both halves are needed]\label{rmk:why-monoidal}
The two halves do different work.  The multiplicative half carries the tower between
levels (\cref{thm:lambda-monoid}); it is what the operator of \cref{sec:operator} and the
membership statement of \cref{ssec:core-forms} act on.  It does not by itself determine
the character: the Euler factor at $p$ needs $\chi_5(p)$, and $\chi_5(p)$ is read off
the ring Frobenius inside a level (\cref{thm:frobenius-congruence,thm:chi5-from-frobenius}).
In Lean the monoid is formalised with $\langle\golden\rangle$
identified with $(\Z,+)$ and $\Frob_p$ with multiplication by $p$
(\LeanInline{psi\_comp}, \LeanInline{psi\_unique} \tierP), and the ring lift by its congruence
(\LeanInline{FrobeniusUnico.frobenius\_exists\_unique} \tierP).
\end{remark}

\subsubsection{Cyclotomic descent and pentagonal descent}
\label{ssec:cyclotomic}

Manin~\cite{manin08} introduces cyclotomic coordinates on
$\mathbb{F}_1$ via the algebraic
extension~$\mathbb{F}_{1^n}\otimes\mathbb{Z} = \mathbb{Z}[T]/(T^n-1)$
\textup{(\cref{rmk:F1n})}.  The natural cyclotomic
structures over $\mathbb{F}_1$ are then organised by the chain of
roots-of-unity inclusions
$\mathbb{F}_{1^n} \subset \mathbb{F}_{1^{n'}}$ when $n\mid n'$.

The pentagonal descent of the present paper is the case $n=5$,
but with an additional structure that distinguishes it from the
purely cyclotomic case.  The unit
$\golden = 2\cos(\pi/5) \in \mathbb{Z}[\zeta_5]^+$
\textup{(}the maximal totally real subring of the fifth
cyclotomic field\textup{)} is the fundamental unit of the
real-quadratic field $\mathbb{Q}(\sqrt{5})$, and the descent
data of $\Rpcf$ is therefore real-quadratic, not purely
cyclotomic: it sits inside the cyclotomic fields
$\Q(\zeta_5)=\Q(\zeta_{10})\subset\Q(\zeta_{20})$ but factors through
$\mathbb{Q}(\sqrt{5})$ rather than through the full
$\mathbb{Q}(\zeta_5)$.

\begin{remark}[Manin's absolute Frobenius]\label{rmk:absolute-frobenius}
The Manin survey~\cite{manin13} of Buium's arithmetic-differential programme and its relation to $\Fone$ discusses the notion of
\emph{absolute Frobenius}: the descent data
$\mathrm{Spec}\,A\to\mathbb{F}_1$ globally encoded by the system
of $\Frob_p$ on a $\Lambda$-ring~$A$.  The monoidal lift
$\Frob_p(\golden) = \golden^p = F_p\golden + F_{p-1}$ is a concrete
realisation of this absolute Frobenius, with the additional
feature that the coefficients
$F_p, F_{p-1}$ are read directly from the pentagonal arithmetic
$\golden^2 = \golden + 1$; on the ring $\Zphi$ the lift is the identity or the
conjugation according to the roots of $x^2-x-1$ modulo $p$ (\cref{sec:lambda-ring}).
\end{remark}

% ---------------------------------------------------------------------
\subsection{The Artin $\Lambda$-structure on $\Zphi$ away from $5$}\label{sec:lambda-ring}

\subsubsection{The Frobenius lift, by its congruence}

\begin{definition}\label{def:sigma}
Let $\sigma\colon\Zphi\to\Zphi$ be the ring homomorphism with
$\sigma(\omega)=1-\omega$; it exists because $1-\omega$ is a root of
$f_\golden$, and $\sigma\circ\sigma=\mathrm{id}$ since $\sigma^2(\omega)=1-(1-\omega)=\omega$.
\end{definition}
\leanline{\LeanInline{FrobeniusUnico.Q.σ}, \LeanInline{FrobeniusUnico.Q.σ\_ω} \tierP;
\LeanStar{\ensuremath{\sigma}\_comp\_\ensuremath{\sigma}} \tierPnew.}

By \cref{prop:endo}, $\mathrm{id}$ and $\sigma$ are the only ring endomorphisms of $\Zphi$.  The
lift at $p$ is not chosen between them: the congruence chooses.

\begin{theorem}[The Frobenius lift at a prime]\label{thm:frobenius-congruence}
Let $p\ne5$ be prime.  There is exactly one ring endomorphism $\Frob_p$ of $\Zphi$ with
\[
  \Frob_p(a)\equiv a^p\pmod{p\Zphi}\qquad\text{for every }a\in\Zphi .
\]
It is the identity if $x^2-x-1$ has a root modulo $p$, and $\sigma$ if it has none.
\end{theorem}
\leanline{\LeanInline{FrobeniusUnico.frobenius\_exists\_unique},
\LeanInline{FrobeniusUnico.frob\_congruence}, \LeanInline{FrobeniusUnico.frob\_unique},
\LeanInline{FrobeniusUnico.endo\_dichotomy}, \LeanInline{FrobeniusUnico.red\_eq\_zero\_iff} \tierP;
\LeanStar{frobenius\_congruence} \tierPnew.}
\begin{proof}
\emph{At most one.}  By \cref{prop:endo} a ring endomorphism sends $\omega$ to $\omega$ or to
$1-\omega$, and is determined by that image.  If both satisfied the congruence, then
$\omega\equiv\omega^p\equiv1-\omega$, so $2\omega-1\in p\Zphi$; its coordinates $-1$ and $2$ in the
basis $1,\omega$ would both be divisible by $p$, which fails for $1$.

\emph{On the generator.}  By \cref{thm:two-spectra}(2), $\omega^p\equiv\omega$ if $x^2-x-1$ has a
root modulo $p$ and $\omega^p\equiv1-\omega$ if it has none; in the first case the identity, in the
second $\sigma$, satisfies the congruence at $a=\omega$.

\emph{On the ring.}  Write $a=u\omega+v$ with $u,v\in\Z$.  In
$\Zphi/(p)$ the $p$-th power map is additive, since the binomial coefficients
$\binom pk$ with $0<k<p$ are divisible by $p$; and it fixes $\Z/p$ by Fermat.
Hence
$a^p\equiv u^p\omega^p+v^p\equiv u\,\Frob_p(\omega)+v=\Frob_p(u\omega+v)=\Frob_p(a)$,
using that a ring endomorphism is the identity on $\Z$.
\end{proof}

\begin{definition}[The Frobenius sign]\label{def:artin-frob}
For a prime $p\ne5$ let $\varepsilon(p)=+1$ if $\Frob_p=\mathrm{id}$ and $\varepsilon(p)=-1$ if
$\Frob_p=\sigma$; we call $p$ \emph{split} in the first case and \emph{inert} in the second
(\cref{sec:dedekind} shows that these are its splitting types in $\Zphi$).  For $n$ prime to $5$,
$n=\prod p^{a_p}$, put $\Frob_n=\prod\Frob_p^{a_p}$ and $\varepsilon(n)=\prod\varepsilon(p)^{a_p}$,
and $\varepsilon(1)=1$, $\Frob_1=\mathrm{id}$.
\end{definition}

So $\varepsilon(p)=\varepsilon_p$ of \cref{thm:two-spectra}(2), and
$\Frob_n=\mathrm{id}$ if $\varepsilon(n)=+1$ and $\Frob_n=\sigma$ if $\varepsilon(n)=-1$.  That
$\varepsilon$ is the Dirichlet character $\chiv$ is \cref{thm:sign-is-chi5}; no statement before it
uses $\chiv$.  Being ring homomorphisms, the $\Frob_n$ are additive and multiplicative by
construction --- which is exactly what \cref{rmk:monoid-not-ring} says the monoidal lift is not.

\begin{proposition}\label{prop:frob-compose}
$\Frob_m\circ\Frob_n=\Frob_{mn}$ and $\varepsilon(mn)=\varepsilon(m)\varepsilon(n)$ for $m,n$
prime to $5$; in particular the $\Frob_p$ commute.
\end{proposition}
\leanline{\LeanStar{\ensuremath{\psi}\_comp}, \LeanStar{\ensuremath{\psi}\_commute} \tierPnew, stated with
$\chiv$ in place of $\varepsilon$ (\cref{thm:sign-is-chi5}).}
\begin{proof}
The $\Frob_p$ lie in the group $\{\mathrm{id},\sigma\}$ of order two, which is commutative, so the
products in \cref{def:artin-frob} do not depend on the order, and $\Frob_n=\sigma^{k}$ with $k$ the
number of inert prime factors of $n$ counted with multiplicity; $\varepsilon(n)=(-1)^k$.  Both are
multiplicative in $n$.
\end{proof}

\Cref{thm:frobenius-congruence,prop:frob-compose} give the $\Lambda$-ring condition of
\cref{def:lambda-ring} at every prime $p\ne5$: a
commuting family of ring endomorphisms, one per prime, each congruent to the
$p$-th power.  It is the additive half of the descent datum, and it holds on the
level, not between levels; \cref{thm:descent-datum} below states it with the prime $5$.

\subsubsection{The obstruction at the ramified prime}

\begin{theorem}\label{thm:obstruction-five}
No ring endomorphism $f$ of $\Zphi$ satisfies $f(\omega)\equiv\omega^5\pmod5$.
\end{theorem}
\leanline{\LeanStar{no\_frobenius\_lift\_at\_five} \tierPnew.}
\begin{proof}
By \eqref{eq:binet-integral}, $\omega^5=F_5\omega+F_4=5\omega+3\equiv3\pmod 5$,
and by \cref{prop:endo} $f(\omega)\in\{\omega,1-\omega\}$.  So the claim is that
neither $\omega-3$ nor $(1-\omega)-3=-2-\omega$ lies in $5\Zphi$.  Consider the
ring homomorphism
\[
  \delta_5\colon\Zphi\longrightarrow\F_5[\varepsilon],\qquad
  \omega\longmapsto3+\varepsilon,\qquad \varepsilon^2=0,
\]
which is well defined because $(3+\varepsilon)^2=9+6\varepsilon=4+\varepsilon
=(3+\varepsilon)+1$ in $\F_5[\varepsilon]$.  It kills $5\Zphi$, while
$\delta_5(\omega-3)=\varepsilon\ne0$ and $\delta_5(-2-\omega)=-5-\varepsilon
=-\varepsilon\ne0$.
\end{proof}

The nilpotent is the point.  In $\F_5$ the polynomial $f_\golden$ has the double
root $3$ --- $5$ is the only ramified prime, $\Delta_K=5$ --- so the reduction
cannot separate $\omega$ from $3$, and the first-order information that
distinguishes them is exactly $\varepsilon$.  The obstruction is the
ramification, exhibited.

\begin{corollary}\label{cor:localisation}
$\Z[\golden,\golden^{-1},\tfrac12]$ admits no $\Lambda$-ring structure, since the
condition at $5$ is part of the definition and fails.  Inverting $5$ instead
makes the condition at $5$ vacuous, and \cref{thm:frobenius-congruence} gives the
structure on $\Rloc$.
\end{corollary}
\leanline{The congruence at each $p\ne5$ is \tierP\ (\cref{thm:frobenius-congruence}); the
localisation itself is not formalised, \tierO.}
\begin{proof}
The nilpotent argument of \cref{thm:obstruction-five} applies to $\Z[\golden,\golden^{-1},\tfrac12]$,
since $\delta_5$ extends to it ($2$ and $3+\varepsilon$ are units of $\mathbb F_5[\varepsilon]$).
After inverting $5$ the congruence at $5$ is a congruence modulo a unit, and holds for every map; at $p\ne5$ the lifts of \cref{thm:frobenius-congruence} extend to the localisation, since they fix $5$.
\end{proof}

\begin{theorem}[The pentagonal descent datum]\label{thm:descent-datum}
Let $E$ be the set of primes other than $5$.  The ring $\Zphi$ carries exactly one
$\Lambda_{\Z,E}$-structure in the sense of Borger~\cite[\S1]{borger11}: the family
$(\Frob_p)_{p\in E}$ of \cref{thm:frobenius-congruence}, which is $\mathrm{id}$ or $\sigma$ at each $p$
according to the roots of $x^2-x-1$ modulo $p$.  It carries no $\Lambda$-structure in the sense of
\cref{def:lambda-ring}, the obstruction being the prime $5$ alone, and the $\Lambda_{\Z,E}$-structure
extends to $\Rloc$, on which the condition at $5$ is empty.
\end{theorem}
\leanline{\LeanInline{FrobeniusUnico.lambda\_structure\_away\_from\_five},
\LeanInline{FrobeniusUnico.frob\_commute}, \LeanInline{FrobeniusUnico.frobenius\_exists\_unique} \tierP;
\LeanStar{no\_frobenius\_lift\_at\_five} \tierPnew; the localisation \tierO.}
\begin{proof}
$\Zphi$ is torsion-free, and for a torsion-free ring a $\Lambda_{\Z,E}$-structure is a family of commuting
ring endomorphisms, one for each $p\in E$, each congruent to the $p$-th power modulo $p$
\cite[\S1]{borger11}.  Such a family exists, and each of its members is unique, by
\cref{thm:frobenius-congruence}; its members commute by \cref{prop:frob-compose}.  So the structure exists
and is unique.  At $5$ no ring endomorphism is congruent to the fifth power (\cref{thm:obstruction-five}),
so the condition of \cref{def:lambda-ring} fails there and only there.  The extension to $\Rloc$ is
\cref{cor:localisation}.
\end{proof}

\subsubsection{The bridge between the two directions}

\begin{corollary}\label{thm:monoid-ring-bridge}
For every prime $p\ne5$ and every $n\in\N$,
\[
  \omega^{pn}\equiv\Frob_p(\omega^{n})\pmod{p\Zphi},\qquad
  \iota(\omega^{pn})=\mathrm{realise}(\Frob_p\,n),
\]
so the lift between levels of \cref{thm:lambda-monoid} and the lift inside the
level of \cref{thm:frobenius-congruence} agree modulo every prime where the second exists.
\end{corollary}
\leanline{\LeanStar{monoid\_ring\_bridge}, \LeanStar{monoid\_side\_is\_realise} \tierPnew.}
\begin{proof}
$\omega^{pn}=(\omega^{n})^{p}\equiv\Frob_p(\omega^{n})$ by
\cref{thm:frobenius-congruence}; the second identity is the definition of
$\mathrm{realise}$ and $\iota$.
\end{proof}

% ---------------------------------------------------------------------
\begin{proposition}[The machinery of the obstruction, and the norm sign]\label{prop:artin-machinery}
\begin{enumerate}[label=\textup{(\roman*)}]
  \item A ring homomorphism out of $\Zphi$ is determined by the image of $\omega$.
    In the dual numbers over $\mathbb F_5$, $5=0$ while $\varepsilon\ne0$; the map
    $\delta_5$ sends $\omega$ to $3+\varepsilon$ and vanishes on $5\Zphi$; and
    $\omega^5=5\omega+3$.
  \item With $\bar\golden=-1/\golden$: $(\golden\bar\golden)^k=(-1)^k$, which is $1$ exactly
    for $k$ even; $\bar\golden^k=(-1)^k\golden^{-k}$; and the trace
    $\golden^k+\golden^{-k}=\golden^k+(-1)^k\bar\golden^k$ carries the sign character.
    \item An integer quadratic $x^2-tx+d$ with odd trace $t$ and unit determinant $d=\pm1$ has no root
    modulo $2$, so $2$ is inert in its ring, and its discriminant is $5$ modulo $8$.  For a prime
    $\mathfrak p$ of $\Zphi$ write $\golden_{\mathfrak p}$ for the image of $\golden$ in the residue field
    $\Zphi/\mathfrak p$, and $\golden_p$ when $p$ is inert and $\mathfrak p=(p)$.  At the inert
    primes $2$ and $3$, $\golden_2$ has order exactly $3$ in $\mathbb F_4$ and $\golden_3$ exactly $8$ in
    $\mathbb F_9$, and so has the image of $\bar\golden$, the other root.
\end{enumerate}
\end{proposition}
\leanline{\LeanStar{hom\_ext}, \LeanStar{five\_eq\_zero\_dual}, \LeanStar{eps\_ne\_zero'},
\LeanStar{δ5\_ω}, \LeanStar{δ5\_kills\_I5}, \LeanStar{ω\_pow\_five};
\LeanInline{norm\_phi\_pow}, \LeanInline{norm\_one\_iff\_even}, \LeanInline{the\_two\_pairs\_are\_one},
\LeanInline{trace\_carries\_the\_character}, \LeanInline{two\_inert\_of\_trace\_odd},
\LeanInline{disc\_mod\_eight\_of\_trace\_odd}, \LeanInline{ord\_phiBar\_mod\_two},
\LeanInline{ord\_phiBar\_mod\_three} \tierP.}
\begin{proof}
(i) Every element is $u+v\omega$, so a ring homomorphism is fixed by the image of $\omega$; the rest
is the computation in the proof of \cref{thm:obstruction-five} and $\omega^5=F_5\omega+F_4$.
(ii) $\golden\bar\golden=N(\golden)=-1$.  (iii) Modulo $2$ the polynomial is $x^2+x+1$, with no root
in $\F_2$; $t^2\equiv1\pmod8$ for $t$ odd, so $t^2-4d\equiv1\mp4\equiv5\pmod8$.  The orders are those
of a root of $x^2-x-1$ in $\F_4^\times$ and $\F_9^\times$: $\golden_2^3=1$ in $\F_4$, and in $\F_9$
$\golden_3^{4}=-1$ and $\golden_3^2\neq\pm1$; at an inert prime the image of $\bar\golden$ is $\golden_p^{\,p}$, of the
same order.
\end{proof}

\subsubsection{The $\delta$-ring at each prime, and the polar form of $\golden$}\label{ssec:delta-polar}

The Frobenius congruence says that $\Zphi$ is, at every prime $p\neq5$, a $\delta$-ring in the
sense of Bhatt and Scholze: for a $p$-torsionfree ring a lift of Frobenius is the same datum as a $p$-derivation $\delta$ with $\Frob(x)=x^p+p\,\delta(x)$, and a $\delta$-structure is the same as a
$p$-typical $\lambda$-structure in Borger's sense \cite[Remarks~2.2--2.3]{bhattscholze22}.

\begin{proposition}[The $p$-derivation of the golden ring]\label{prop:delta-ring}
\begin{enumerate}[label=\textup{(\roman*)}]
  \item For every prime $p\neq5$ and every $a\in\Zphi$ there is $\delta_p(a)\in\Zphi$ with $\Frob_p(a)=a^p+p\,\delta_p(a)$.
  \item For $a=\golden$, with $\golden^p=F_p\golden+F_{p-1}$, $\delta_p(\golden)=\bigl((1-F_p)\golden-F_{p-1}\bigr)/p$ if
    $p$ splits and $\bigl((-1-F_p)\golden+1-F_{p-1}\bigr)/p$ if $p$ is inert; for instance
    $\delta_{11}(\golden)=-5-8\golden$ and $\delta_{13}(\golden)=-11-18\golden$.
\end{enumerate}
\end{proposition}
\leanline{\LeanStar{F1PaperPadic.delta\_structure\_exists}, \LeanStar{F1PaperPadic.delta\_table\_inert}, \LeanStar{F1PaperPadic.delta\_table\_split} \tierPnew.}

At each prime $\delta_p(\golden)$ measures the difference between the multiplicative half of the
descent, the tower $\golden\mapsto\golden^p$ (\cref{sec:lambda-monoid}), and the additive half, the Artin Frobenius $\Frob_p$; it is fixed by $F_p$, $F_{p-1}$ and $\chiv(p)$ alone, and its size is that of the tower,
$|\delta_p(\golden)|/(\golden^p/p)=1.236,\ 1.146,\ 0.992,\ \ldots,\ 0.999999$ at $p=2,3,11,\ldots,31$ (\tierN).

In the completion at a prime $\mathfrak p$ the generator splits into an angular and a radial part,
$\golden=[\golden_{\mathfrak p}]\cdot\langle\golden\rangle$: the Teichm\"uller representative $[\golden_{\mathfrak p}]$ of
$\golden_{\mathfrak p}$, a root of unity,
and a principal unit $\langle\golden\rangle\equiv1\pmod p$.  No real or complex number enters.  For the golden and metallic ratios in $\Z_p$ see~\cite{akbiyik19}.

\begin{proposition}[The polar form of $\golden$ at $2$ and $3$, and the colour group]\label{prop:polar}
\begin{enumerate}[label=\textup{(\roman*)}]
  \item At $2$ the angular part of $\golden$ is a primitive cube root of unity $\zeta_3$:
        $\mathbb{Q}_2(\sqrt5)=\mathbb{Q}_2(\sqrt{-3})$, since $-15\equiv1\pmod{16}$ is a $2$-adic square; the
    reduction of $\golden$ generates $\F_4^\times$, and its Teichm\"uller lift is a primitive cube root of
    unity of $\mathbb{Q}_2(\sqrt5)$, which we denote $\zeta_3$.
  \item At $3$ the angular part has order $8$, and its square is $\pm i$ (the orders $3$ and $8$ of $\golden_2$ and $\golden_3$ in
    $\mathbb{F}_4$ and $\mathbb{F}_9$ are \cref{prop:artin-machinery}(iii); what is added is the lift and the $\pm i$).
  \item The orders $3$ and $4$ of these angular parts are those of the hexagonal and the Gauss
    lattices, of least common multiple $12$.
    \item The colour group at an inert $p$, the subgroup $\langle\golden_p\rangle$ of the residue field generated by the
    image of $\golden$, has even order for $p\geq3$, never $2$, $4$ or
    $6$, since $\golden_p^{\,p+1}=\det M=-1$; from $p=7$ on it is a proper subgroup of
    $\mathbb{F}_{p^2}^\times$ ($16$, not $48$, at $7$; $28$, not $168$, at $13$).  With determinant $+1$ (trace
    $6$) the prime $3$ gives order $4$ and the unit group $\mu_4=\Z[i]^\times$: the orientation
    $\det M=-1$ is what puts the angular part of $\golden$ in the Eisenstein and not in the
    Gauss integers.
\end{enumerate}
\end{proposition}
\leanline{\LeanStar{F1PaperPadic.polar\_at\_two}, \LeanStar{F1PaperPadic.polar\_at\_three}, \LeanStar{F1PaperPadic.angular\_parts\_give\_twelve}, \LeanStar{F1PaperPadic.minus\_fifteen\_is\_a\_two\_adic\_square} \tierPnew{}
at precision $2^8$ and $3^5$, the Teichm\"uller lift itself by Hensel's lemma \tierC;
\LeanInline{PCFGranularity.conjugate\_reduces\_to\_other\_generator}, \LeanInline{PCFGranularity.ord\_phiBar\_mod\_two}, \LeanInline{PCFGranularity.ord\_phiBar\_mod\_three}, \LeanInline{PCFGranularity.colour\_order\_even}, \LeanInline{PCFGranularity.six\_is\_the\_least\_over\_the\_colour\_order}, \LeanInline{PCFGranularity.counterexample\_order\_four}, \LeanInline{PCFGranularity.det\_plus\_one\_breaks\_the\_dichotomy} \tierP{}.}

The single exponent $n$ of $\golden^n$ carries both coordinates: reduced modulo $6$ it gives the
angular one, in $\mu_6$, and multiplied by $\log_2\golden$ the radial one; they never resynchronise,
because $\lambda_{\log}=\log2/\log\golden$ is irrational: $\lambda_{\log}=a/b$ would give $2^b=\golden^a$, whose norms are $4^b$ and $\pm1$.  The Gaussian and the Eisenstein unit groups meet
only in $\{\pm1\}$, which is the orientation.

\subsection{The Fibonacci character}\label{sec:chi5}

The sign $\varepsilon$ of \cref{def:artin-frob} was defined by the Frobenius lift.  This
subsection identifies it with the Dirichlet character $\chiv$ of \eqref{eq:chi5-pentagon}, which is $+1$ on the classes $\pm1\bmod5$ and $-1$ on $\pm2\bmod5$, since $2\cos(\pi/5)=\golden$ and $2\cos(2\pi/5)=\golden^{-1}$ (\cref{sec:choice}); it is the
first point of the section at which $\chiv$ enters.

\begin{theorem}[The Frobenius sign is $\chiv$]\label{thm:sign-is-chi5}
For every prime $p\ne5$, $\varepsilon(p)=\chiv(p)$; hence $\varepsilon(n)=\chiv(n)$ for every $n$ prime to
$5$, and $\Frob_n=\mathrm{id}$ or $\sigma$ according as $\chiv(n)=+1$ or $-1$.
\end{theorem}
\leanline{\LeanInline{FrobeniusUnico.hasRoot\_iff\_sq}, \LeanInline{FrobeniusUnico.hasRoot\_iff\_mod5},
\LeanInline{FrobeniusUnico.sign\_eq\_chi5} \tierP, the reciprocity step by Mathlib's
\texttt{ZMod.exists\_sq\_eq\_prime\_iff\_of\_mod\_four\_eq\_one}; \LeanInline{chi5\_is\_the\_mod5\_component}, \LeanInline{five\_sq\_iff\_mod5},
\LeanInline{root\_x2mx1\_iff\_five\_sq} \tierP; \LeanStar{\ensuremath{\psi}\_of\_split},
\LeanStar{\ensuremath{\psi}\_of\_inert} \tierPnew.}
\begin{proof}
For $p=2$: $x^2+x+1$ has no root in $\F_2$, so $\varepsilon(2)=-1=\chiv(2)$.  Let $p$ be odd.  The map
$\alpha\mapsto2\alpha-1$ turns a root of $x^2-x-1$ in $\F_p$ into a square root of the discriminant $5$,
and back, so $\varepsilon(p)=+1$ exactly when $5$ is a square modulo $p$.  Since $5\equiv1\pmod4$,
quadratic reciprocity gives $\left(\frac5p\right)=\left(\frac p5\right)$, which is $+1$ exactly when
$p\equiv\pm1\pmod5$, that is, when $\chiv(p)=+1$.  Both $\varepsilon$ and $\chiv$ are multiplicative
(\cref{prop:frob-compose}), so they agree on every $n$ prime to $5$.
\end{proof}

\begin{theorem}\label{thm:s1}
For every prime $p\ne5$, $\chiv(p)\equiv F_p\pmod p$.
\end{theorem}
\leanline{\LeanStar{S1Closure.chi5\_fib\_all\_primes} \tierPnew;
\LeanInline{FrobeniusUnico.fib\_eq\_sign}, \LeanInline{FrobeniusUnico.sign\_eq\_chi5} \tierP.}
\begin{proof}
$F_p\equiv\varepsilon(p)$ is \cref{thm:two-spectra}(2), and $\varepsilon(p)=\chiv(p)$ is
\cref{thm:sign-is-chi5}.
\end{proof}

\begin{remark}
The congruence is classical~\cite{lucas1878,lehmer30}.  What the chain adds is its place: it is
read off the Frobenius lift, as the next statement makes precise, and $\chiv$ enters only through
\cref{thm:sign-is-chi5}.
\end{remark}

\begin{corollary}[The Fibonacci coordinates are read off the Frobenius]\label{thm:chi5-from-frobenius}
For every prime $p\ne5$, the congruence $\Frob_p(\omega)\equiv\omega^p\pmod p$, read in the basis
$1,\omega$, gives $F_p\equiv\varepsilon(p)=\chiv(p)$ and $F_{p-1}\equiv\tfrac12(1-\chiv(p))\pmod p$.
\end{corollary}
\leanline{\LeanStar{coeffs\_mod\_p}, \LeanStar{chi5\_from\_frobenius},
\LeanStar{lambda\_structure\_gives\_chi5} \tierPnew; \LeanInline{FrobeniusUnico.fib\_eq\_sign},
\LeanInline{FrobeniusUnico.fib\_pred\_eq\_sign} \tierP.}
\begin{proof}
Coordinates can be read modulo $p$: if
$u\omega+v\in p\Zphi$ with $u,v\in\Z$ then $p\mid u$ and $p\mid v$, because $\Zphi$ is free of rank two
over $\Z$ with basis $1,\omega$.  Now
$\omega^p=F_p\omega+F_{p-1}$ by \eqref{eq:binet-integral}, while
$\Frob_p(\omega)$ is $\omega$ or $1-\omega$ according to $\varepsilon(p)$.  Comparing
coordinates gives $F_p\equiv1$, $F_{p-1}\equiv0$ in the first case and
$F_p\equiv-1$, $F_{p-1}\equiv1$ in the second; and $\varepsilon(p)=\chiv(p)$ by
\cref{thm:sign-is-chi5}.
\end{proof}

In the inert case the two coordinates give $F_{p+1}=F_p+F_{p-1}\equiv-1+1=0\pmod p$, and in the split
case $F_{p-1}\equiv0$: every prime $p\ne5$ divides $F_{p-\chiv(p)}$, and for $p\ne2$ the divisibility
$p\mid F_{p+1}$ holds exactly when $p$ is inert (\LeanStar{F1PaperPadic.fib\_succ\_zero\_iff\_inert} \tierPnew).

\begin{table}[ht]
\centering
\caption{The Fibonacci criterion for $3\le p\le31$, $p\ne5$: $F_p\equiv\chiv(p)\pmod p$, and $p$
divides $F_{p-\chiv(p)}$.  The values of $F_p$ and $F_{p-1}$ modulo $p$ are those of the $p$-derivation of
\cref{prop:delta-ring} (\LeanStar{F1PaperPadic.delta\_table\_inert}, \LeanStar{F1PaperPadic.delta\_table\_split}
\tierPnew; up to $p=19$ also \LeanInline{fib\_split\_criterion} \tierP).}
\label{tab:fib-split}
\begin{tabular}{rccccc}
\toprule
$p$ & $p\bmod5$ & $\chiv(p)$ & type & $F_p\bmod p$ & $p\mid$ \\
\midrule
3  & 3 & $-1$ & inert & $-1$ & $F_4$ \\
7  & 2 & $-1$ & inert & $-1$ & $F_8$ \\
11 & 1 & $+1$ & split & $+1$ & $F_{10}$ \\
13 & 3 & $-1$ & inert & $-1$ & $F_{14}$ \\
17 & 2 & $-1$ & inert & $-1$ & $F_{18}$ \\
19 & 4 & $+1$ & split & $+1$ & $F_{18}$ \\
23 & 3 & $-1$ & inert & $-1$ & $F_{24}$ \\
29 & 4 & $+1$ & split & $+1$ & $F_{28}$ \\
31 & 1 & $+1$ & split & $+1$ & $F_{30}$ \\
\bottomrule
\end{tabular}
\end{table}

\Cref{tab:fib-split} reads the criterion for the primes up to $31$.

The order of the chain is therefore: the roots of $x^2-x-1$ modulo $p$ decide $\omega^p$ modulo $p$
(\cref{thm:two-spectra}); that decides the lift, which is unique (\cref{thm:frobenius-congruence}); the
lift decides $\varepsilon(p)$, which is read in the coordinates $F_p$, $F_{p-1}$
(\cref{thm:chi5-from-frobenius}); and quadratic reciprocity identifies $\varepsilon$ with $\chiv$
(\cref{thm:sign-is-chi5}).  The character is reached at the last step and is used at no earlier one.

\begin{proposition}\label{prop:pisano}
The Fibonacci sequence modulo $5$ has period $20$, the conductor of \cref{sec:choice}.
\end{proposition}
\leanline{\LeanStar{pisano\_five\_is\_twenty} \tierPnew; the period is classical~\cite{wall60}.}
\begin{proof}
Modulo $5$ the sequence is $0,1,1,2,3,0,3,3,1,4,0,4,4,3,2,0,2,2,4,1,0,1$: the pair $(0,1)$ first
returns at index $20$.
\end{proof}

% ---------------------------------------------------------------------
\begin{proposition}[The two periods, and the Fibonacci identities behind them]\label{prop:two-periods}
\begin{enumerate}[label=\textup{(\roman*)}]
  \item Gauss's sequence (the values of $\chi_4$) has period $4$, the Fibonacci numbers
    modulo $5$ have period $20$, and both periods divide the conductor $20$.
    \item $\chiv$ is multiplicative, depends only on the class modulo $5$, and $\chiv(5)=0$, so the
    factor of $L(s,\chiv)$ at the ramified prime is $1$.
  \item $F_{m+n+1}=F_{m+1}F_{n+1}+F_mF_n$; and for $k$ odd, $5F_k^2-4=L_k^2$.
\end{enumerate}
\end{proposition}
\leanline{\LeanInline{the\_two\_periods\_and\_the\_conductor}, \LeanInline{chi5\_Z\_mul},
\LeanInline{chi5\_depends\_only\_on\_class}, \LeanInline{ramified\_factor\_trivial},
\LeanInline{fib\_convolution}, \LeanInline{disc\_survives\_the\_norm\_sign} \tierP.}

\subsection{The Golden Prime classification}\label{sec:grisales}
\src{HP L2636--L2680; \cite{grisales25}}

The rule that chooses between the two endomorphisms of \cref{prop:endo} is a
character, and the character is one component of a classification of the classes
modulo $20$ that Grisales Herrera introduced \cite{grisales25} and that is formalised here.

\subsubsection{The homomorphism and its fibres}

\begin{proposition}\label{prop:G}
There is a surjective group homomorphism
$G\colon\Zst\to\Z_2\times\Z_2$ whose first component is the character $\chi_4$ of
conductor $4$, trivial exactly on $\{1,9,13,17\}$, and whose second component is
$\chiv$, trivial exactly on $\{1,9,11,19\}$; its kernel is $\{1,9\}$, so
$|\ker G|=2$ and the eight classes fall into four fibres of two.
\end{proposition}
\leanline{\LeanInline{G}, \LeanInline{G\_hom}, \LeanInline{G\_surj},
\LeanInline{G\_pi1\_is\_chi4}, \LeanInline{G\_pi2\_is\_chi5},
\LeanInline{G\_kernel} \tierP.}


\begin{proposition}[The two registers on the classes]\label{prop:two-registers-G}
\leavevmode
\begin{enumerate}
  \item \emph{Separation.}  Two classes with the same reduction modulo $4$ can
    differ in $\chiv$, and two with the same reduction modulo $5$ can differ in
    $\chi_4$; together the two reductions determine the class.
  \item \emph{Descent to Galois.}  $G(pa)=G(p)+G(a)$ for $p$ prime to $20$: the monoidal lift
    $\Frob_p$, read on $\Z/20$ as $a\mapsto pa$ (\cref{thm:lambda-monoid}), is equivariant for $G$.
\end{enumerate}
\end{proposition}
\leanline{\LeanInline{the\_two\_spectra\_are\_separated}, \LeanInline{conductor\_splits},
\LeanInline{G\_psi\_equivariant}, \LeanInline{F1\_descends\_to\_Galois},
\LeanInline{F1\_spectra\_anchor\_G\_bridge} \tierP.}
\begin{proof}
(1) $1$ and $13$ agree modulo $4$ and differ in $\chiv$; $1$ and $11$ agree modulo
$5$ and differ in $\chi_4$; and $\Z/20\cong\Z/4\times\Z/5$.  (2) $G$ is a homomorphism (\cref{prop:G}).
\end{proof}

So $G$ is not a repackaging of one character but the
joint datum, and its kernel measures how much of the class is invisible to both:
exactly one bit, $|\ker G|=2$.  That number reappears in two forms below.  First as $\tfrac12=1/|\ker G|$, which is
the modulus $|2|_2$ of the triad of the inert prime $2$ (\cref{def:tripartite-omega}); the same $\tfrac12$ is later the
positive eigenvalue of the colour operator $\OmH$ and the centre of every pair of zeros $\{\rho,1-\rho\}$, the common
target called the apex in \cref{ssec:apex}, and it is the value at which the chain of operators of \cref{ssec:chain}
closes, in its seventh link.  Second as $\log2=\log|\ker G|$ in the entropy bridge of the anchor
$\golden^{\lng}=2$ (\cref{sec:mersenne}).  Two branches thus start from $G$ --- the arithmetic one, through $\chiv$ and
$L(1,\chiv)$, and the spectral one, through $|\ker G|$ --- and they meet at $\log2$; a pair of maps with a common source
is a span, and this one is \cref{prop:span}
(\LeanInline{kernel\_is\_one\_nine}, \LeanInline{log\_two\_is\_log\_kernel}, \LeanInline{branches\_meet\_at\_log\_two} \tierP).

$G$ is the quotient of $\Zst$ by its squares: $\ker G=\{1,9\}$ is exactly the set of squares, so the
characters of $\Zst$ that factor through $G$ are the quadratic ones
(\LeanStar{F1PaperNew.G\_kernel\_eq\_squares}, \LeanStar{F1PaperNew.squares\_of\_units\_twenty} \tierPnew).

\subsubsection{The component that selects the Frobenius}

For a prime $p\ne5$ the component $\chiv$ of $G$ is exactly the sign of
\cref{def:artin-frob} (\cref{thm:sign-is-chi5}): $p\equiv\pm1\pmod5$ gives $\Frob_p=\mathrm{id}$, and
$p\equiv\pm2\pmod5$ gives $\Frob_p=\sigma$
(\LeanStar{\ensuremath{\psi}\_of\_split}, \LeanStar{\ensuremath{\psi}\_of\_inert} \tierPnew).  The
component $\chi_4$ lives in $\Q(i)$ and does not act on $\Zphi$ at all: it is the
part of the classification that the ring does not see, which is why
\cref{prop:two-registers-G} can separate the two registers by the two factors.

\subsubsection{The pentagonal values}

Grisales attaches to each class a value of the golden ratio.  The values
themselves are proved.

\begin{proposition}[The pentagonal values]\label{thm:pentagonal-values}
$|2\cos(\pi a/5)|=\golden^{\chiv(a)}$ for $a=1,2,3,4$; explicitly the four values
are $\golden,\golden^{-1},\golden^{-1},\golden$.
\end{proposition}
\leanline{\LeanStar{F1CW.norm\_cosine\_phi} \tierPnew.}
\begin{proof}
$2\cos(\pi/5)=\golden$ is \eqref{eq:pentagon-id}.  The double-angle formula with
$\cos(\pi/5)=\golden/2$ gives
$2\cos(2\pi/5)=2\bigl(2(\golden/2)^2-1\bigr)=\golden^2-2=\golden-1=\golden^{-1}$,
using $\golden^2=\golden+1$ and then
$\golden(\golden-1)=\golden^2-\golden=1$.  Finally $\cos(\pi-x)=-\cos x$ carries
the remaining two cases, $\cos(3\pi/5)=-\cos(2\pi/5)$ and
$\cos(4\pi/5)=-\cos(\pi/5)$; taking absolute values gives
$\golden,\golden^{-1},\golden^{-1},\golden$ against $\chiv=+1,-1,-1,+1$.
\end{proof}

So the character is read off the pentagon itself, without any separate arithmetic
input: the four pentagonal cosines are $\pm\golden^{\pm1}$, and the exponent is
$\chiv$.  Of the four signed values $2\cos(\pi a/5)=\golden,\golden^{-1},-\golden^{-1},-\golden$, only
$\golden$ and $-\golden^{-1}$ are roots of $f_\golden$; the other two, $\golden^{-1}$ and
$-\golden$, are roots of $x^2+x-1$.

\begin{remark}[What is not verified]\label{rmk:grisales-not-verified}
The closed formula of \cite{grisales25} for the classification, in its sine form,
is cited and not formalised; in the Lean file it is a marker
(\LeanInline{GLandsInTheTower} \tierO).  What is proved is
\cref{thm:pentagonal-values}, the homomorphism of \cref{prop:G}, and the
logarithmic signature $\sum_{a=1}^4\chiv(a)\log(2\sin(\pi a/5))=-2\log\golden$.  In
Lean the proved statement is $\sin(2\pi/5)/\sin(\pi/5)=\golden$
(\LeanStar{F1CW.pentagon\_log\_signature} \tierPnew), whose whole content is
$\sin2x=2\sin x\cos x$ together with \eqref{eq:pentagon-id}; the signature itself follows
because $\sin(3\pi/5)=\sin(2\pi/5)$ and $\sin(4\pi/5)=\sin(\pi/5)$, so the sum is
$2\log\bigl(\sin(\pi/5)/\sin(2\pi/5)\bigr)$, and it is proved in full in
\LeanStar{F1v2.pentagon\_log\_signature\_full}.
\end{remark}

The Klein group $V$ generated by the two commuting involutions $a\mapsto-a$ and $a\mapsto a^{-1}$ acts on $\Zst$.
Its orbits are $\{1,19\}$ and $\{9,11\}$, of two elements each, and $\{3,7,13,17\}$, of four, on which $V$ acts
freely; this is the \emph{free orbit}, and its elements are those of order four.  The \emph{decoration} of a
digit $a$ in base $b$ is the pair $\bigl(a\bmod20,\ G(a\bmod20)\bigr)$, and a base \emph{reaches} a fibre of $G$ when
the decoration of some digit prime to $b$ lies in it.

\begin{proposition}[The classification needs base $20$]\label{prop:base-twenty}
\begin{enumerate}[label=\textup{(\roman*)}]
  \item The decoration of a unit by its digit in base $20$ is exactly $G$, and base $20$ reaches the four fibres.
    In base $10$ the decoration of the four unit digits $1,3,7,9$ agrees with $G$, and it lies on the diagonal
    $\{(0,0),(1,1)\}$: base $10$ does not reach the two off-diagonal fibres $(1,0)$ and $(0,1)$, so the last
    digit alone cannot separate them.
  \item On the classes $\{1,3,7,9\}$ the two components of $G$ agree; $\chi_4$ is odd,
    $G(-a)_1=G(a)_1+1$, while $\chiv$ is even.
  \item The free orbit $\{3,7,13,17\}$ is the union of the two fibres $(1,1)$ and $(0,1)$.
\end{enumerate}
\end{proposition}
\leanline{\LeanInline{base20\_bijection}, \LeanInline{decoration\_matches\_G},
\LeanInline{off\_diagonal\_inaccessible}, \LeanStar{off\_diagonal\_inaccessible'},
\LeanInline{diagonal\_from\_residue}, \LeanInline{chi\_four\_is\_odd},
\LeanInline{deviation\_type\_is\_chi\_four}, \LeanInline{free\_orbit\_is\_union\_of\_fibres},
\LeanInline{free\_orbit\_eq\_order\_four} \tierP.}

This is why the classification of \cite{grisales25} reads the last \emph{two} digits:
the last digit determines the class modulo $10$, which reaches only the diagonal of the
Klein group, and the parity of the penultimate digit supplies the class modulo $4$, that
is, the value of $\chi_4$.

% =====================================================================
\subsection{The clock of $\golden$}\label{sec:clock}

The classification fixes the class of a prime modulo $20$.  The clock fixes, prime by prime, the order of
the angle of $\golden$.  This subsection defines it, writes it with the cyclotomic polynomials, reads what the
class fixes of it, and adds the index, which fixes the rest of the position of $\golden$ at each prime.

\subsubsection{The clock and the cyclotomic polynomials}\label{ssec:clock-cyclotomic}

Let $e_p$ be the order of $\golden$ in $(\Zphi/p)^\times$, the order of its angular part at $p$.
Manin's coordinate $T$ of $\mathbb{F}_{1^n}\otimes\mathbb{Z}=\mathbb{Z}[T]/(T^n-1)$
(\cref{ssec:cyclotomic}), evaluated at $\golden$, reads the cyclotomic structure on the primes of $\Zphi$.
The \emph{clock} of a prime $\mathfrak p$ of $\Zphi$ is the order of $\golden$ in
$(\Zphi/\mathfrak p)^\times$, the order of the angular part of $\golden$ at
$\mathfrak p$ (\cref{ssec:delta-polar}).  The order $e_p$ just defined
is the least common multiple of the clocks of the primes above $p$; at a split prime
the two clocks can differ ($5$ and $10$ at $11$).  No prime has clock $1$ or $2$, since
$\golden-1=\golden^{-1}$ and $\golden+1=\golden^2$ are units.

\begin{proposition}[The clock as cyclotomic polynomials]\label{prop:clock-cyclotomic}
Let $n\geq1$, and let $\Phi_n$ be the $n$-th cyclotomic polynomial.
\begin{enumerate}[label=\textup{(\roman*)}]
  \item $\golden^n-1=\prod_{d\mid n}\Phi_d(\golden)$.
  \item Let $\mathfrak p$ be a prime of $\Zphi$ of residue characteristic $p\nmid n$.  Then
    $\mathfrak p$ has clock $n$ if and only if $\mathfrak p$ divides $\Phi_n(\golden)$.
  \item For $n\geq2$, $|N(\Phi_n(\golden))|=\prod_\zeta|\golden-\zeta|\,|{-1/\golden}-\zeta|$,
    the product over the primitive $n$-th roots of unity $\zeta$, the vertices of the regular
    $n$-gon; their number is $\phi(n)$, and $n\log\golden=\sum_{d\mid n}\phi(d)\log\golden$.
    With $\eta_d=\log|N(\golden^d-1)|-d\log\golden$,
    \[
      \log|N(\Phi_n(\golden))|=\phi(n)\log\golden+\sum_{d\mid n}\mu(n/d)\,\eta_d ,
    \]
    where $\eta_1=-\log\golden$, $\eta_2=-2\log\golden$, and $|\eta_d|\leq4\golden^{-d}$;
    so $\bigl|\log|N(\Phi_n(\golden))|-\phi(n)\log\golden\bigr|\leq4\golden^2$ for every
    $n\geq1$.
  \item If a prime $\mathfrak p$ of residue characteristic $q\mid n$ divides $\Phi_n(\golden)$,
    then $n$ is the clock of $\mathfrak p$ times a power $q^j$, $j\geq1$.
  \item For $n\geq2$, $N(\Phi_n(\golden))$ is a nonzero integer, and
    $\log|N(\Phi_n(\golden))|=\sum_p v_p(N(\Phi_n(\golden)))\log p$, the sum running over the
    primes of clock $n$ and the primes of \textup{(iv)}.  On each level $n$ of the clock the
    archimedean size of \textup{(iii)} and the $p$-adic weights of the primes of that level are
    therefore equal in sum.
  \item Summed over the levels, the errors of \textup{(iii)} are a combination of values of
    the Mertens function $M(x)=\sum_{k\leq x}\mu(k)$:
    \[
      \sum_{n\leq N}\bigl(\log|N(\Phi_n(\golden))|-\phi(n)\log\golden\bigr)
      =\sum_{d\leq N}\eta_d\,M(N/d),
    \]
    so that $|M(x)|\leq Cx^\theta$ for all $x\geq1$, with $C,\theta\geq0$, gives an accumulated error at most
    $4\golden^2CN^\theta$.
  \item \emph{The mass of the clock.}  Summed over the levels,
    \[
      \sum_{n\leq N}\log|N(\Phi_n(\golden))|=\frac{3}{\pi^2}\,\log\golden\;N^2+O(N\log N)
      =\frac{\log\golden}{2\zeta(2)}\,N^2+O(N\log N).
    \]
\end{enumerate}
\end{proposition}
\leanline{\LeanStar{W4Isometria.clock\_cyclotomic\_split} (i),
\LeanStar{W4Isometria.RelojCiclotomico.clock\_iff\_cyclotomic\_root} (ii) in every
residue field, \LeanStar{W4Isometria.RelojCiclotomico.mate\_pow} (the two clocks $n$ and
$2n$ of a split prime), \LeanStar{W4Isometria.totient\_is\_primitive\_roots},
\LeanStar{W4Isometria.tower\_weight\_splits\_by\_clock},
\LeanStar{W4Isometria.RelojCiclotomico.prod\_Ncyc},
\LeanStar{W4Isometria.RelojCiclotomico.err\_moebius},
\LeanStar{W4Isometria.RelojCiclotomico.log\_Npow\_bound},
\LeanStar{W4Isometria.RelojCiclotomico.log\_Ncyc\_bound} (iii),
\LeanStar{W4Isometria.RelojCiclotomico.clock\_sum\_mertens},
\LeanStar{W4Isometria.RelojCiclotomico.clock\_sum\_le\_of\_mertens} (vi) \tierPnew;
(iv) and (v) \tierC; (vii) from (vi) and $\sum_{n\leq N}\phi(n)=3N^2/\pi^2+O(N\log N)$ \tierC;
(ii) for $n\leq600$ and $p\leq20000$, $p\nmid n$ ($1\,355\,192$ pairs,
$752$ of the $2260$ primes with two clocks), (vi) up to $N=10^6$ to $3.6\cdot10^{-12}$
(\texttt{v27\_reloj\_ciclotomico.py}, \tierN), and (vii), whose ratio to the main term is % pending(D7)
$1.0038$, $1.0018$, $1.0008$ and $1.0006$ at $N=100$, $500$, $1000$ and $2000$
(\texttt{retro69.py} \tierN).} % pending(D7)

\begin{proof}
(i) Evaluate $X^n-1=\prod_{d\mid n}\Phi_d(X)$ at $\golden$.
(ii) In the residue field $k$ of $\mathfrak p$, $n\neq0$, so $X^n-1$ is separable over $k$ and its
roots of order exactly $n$ are the roots of $\Phi_n$.
(iii) $\Phi_n(X)=\prod_\zeta(X-\zeta)$ and $N(\Phi_n(\golden))=\Phi_n(\golden)\Phi_n(\bar\golden)$
with $\bar\golden=-1/\golden$; the number of the $\zeta$ is $\phi(n)$, and $\sum_{d\mid n}\phi(d)=n$.
By (i), $\log|N(\golden^n-1)|=\sum_{d\mid n}\log|N(\Phi_d(\golden))|$, and Möbius inversion with
$\sum_{d\mid n}\mu(n/d)\,d=\phi(n)$ gives the formula.  Since
$|N(\golden^d-1)|=|(-1)^d-L_d+1|$ with $L_d=\golden^d+\bar\golden^d$,
$\eta_d=\log(1+x_d)$ for $d\geq3$, with $x_d=(\bar\golden^d-1-(-1)^d)\golden^{-d}$ and
$|x_d|\leq2\golden^{-d}\leq0.48$ (for odd $d$, $x_d=(\bar\golden/\golden)^d$), so
$|\eta_d|\leq|x_d|/(1-|x_d|)\leq4\golden^{-d}$; at
$d=1,2$ the norm is $\mp1$ and $|\eta_d|=d\log\golden\leq4\golden^{-d}$.  Hence
$\sum_{d\mid n}|\eta_d|\leq4\sum_{d\geq1}\golden^{-d}=4\golden\leq4\golden^2$.
(iv) Let $m$ be the clock of $\mathfrak p$, so $q\nmid m$, and write $n=n'q^j$ with $q\nmid n'$.
Modulo $q$, $\Phi_n\equiv\Phi_{n'}^{\,\phi(q^j)}$, so $\golden$ is a root of $\Phi_{n'}$ in the
residue field and $m=n'$ by (ii); $j\geq1$ because $q\mid n$.
(v) $\golden$ is not a root of unity, so $\Phi_n(\golden)\neq0$, and its norm is an integer; the
equality is the factorisation of that integer, and its primes are described by (ii) and (iv).
(vi) Summing the formula of (iii) over $n\leq N$ and grouping the terms by $d$ gives the identity;
the bound follows from $|M(N/d)|\leq C(N/d)^\theta\leq CN^\theta$, which uses $\theta\geq0$,
and $\sum_d|\eta_d|\leq4\sum_{d\geq1}\golden^{-d}=4\golden\leq4\golden^2$.
(vii) By (vi) with $C=\theta=1$ the errors accumulate to at most $4\golden^2N$, and the main terms
sum to $\log\golden\sum_{n\leq N}\phi(n)$.
\end{proof}

The two irrationals of the construction meet in the constant of (vii): $\log\golden$, the regulator,
is the length of one turn of the clock, and $1/\zeta(2)=6/\pi^2$ is the density of the primitive
classes, the proportion of the turns that open a new level.  The mass of the levels grows with the
square of the level, whereas the primes ordered by size grow linearly in their weight,
$\sum_{p\leq x}\log p\sim x$; passing from one growth to the other is the change from the order of
the clock to the order by size.

\begin{remark}[The angles as cyclotomic norms]\label{rmk:angles-cyclotomic}
The angular parts of \cref{prop:polar} are read by \cref{prop:clock-cyclotomic} on integers,
without the completion.  For $x$ with $x^2=x+1$ and its conjugate $1-x$, the norms of
$\Phi_3$, $\Phi_4$, $\Phi_8$, $\Phi_{12}$ and $\Phi_5$ at $x$ are $4$, $5$, $9$, $4$ and $11$
(\LeanStar{W4Isometria.conductor\_cyclotomic\_norms} \tierPnew).  So $2$, inert, has clock $3$,
the angular part $\zeta_3$ of \cref{prop:polar}(i); $3$, inert, has clock $8$, the order of
\cref{prop:polar}(ii); the ramified prime $\sqrt5$ has clock $4$, since $\golden\equiv3$ modulo
$\sqrt5$; $N(\Phi_{12}(\golden))=4$ is the prime $2$ again, at the level $12=3\cdot2^2$ of
\cref{prop:clock-cyclotomic}(iv); and $11$, split, has the clocks $5$ and $10$ at its two primes.
At $5$ the levels $4\cdot5^k$ carry the ramified prime: by
\cref{prop:clock-cyclotomic}(i), $\golden^{4\cdot5^k}-1$ is the product of the $\Phi_d(\golden)$
with $d\mid4\cdot5^k$, and $\sqrt5$ divides it exactly $2k+1$ times
(\LeanStar{F1PaperSixty.phi\_clock\_identity}, \LeanStar{F1PaperSixty.phi\_clock\_exact}
\tierPnew), so $v_5(\golden^{4\cdot5^k}-1)=k+\tfrac12$; level by level, $v_5$ of
$N(\Phi_{4\cdot5^k}(\golden))$ is $1,2,2,2$ for $k=0,1,2,3$ (\tierN).  A root of unity has
absolute value $1$ in every normed field
(\LeanStar{M12F1Circle.root\_of\_unity\_norm\_one} \tierPnew), so at each prime the angular part
carries the clock and the principal unit $\langle\golden\rangle$ carries the size.
\end{remark}


\subsubsection{The classes and the angular part of $\golden$}\label{ssec:classes-angular}

Write $N_p=p-1$ if $p$ splits and $N_p=2(p+1)$ if $p$ is inert.  Then $\golden^{N_p}\equiv1\pmod p$: if $p$
splits, $\Zphi/p\cong\F_p\times\F_p$; if $p$ is inert, $\golden_p^{\,p+1}=N(\golden)=-1$.  So $e_p$ divides $N_p$.

\begin{proposition}[What the class fixes of the angular part]\label{prop:class-angular}
\begin{enumerate}[label=\textup{(\roman*)}]
  \item In the inert classes $3,7,13,17$, $\bar\golden^{\,p+1}=-1$, so the $2$-adic valuation of $e_p$ is
    that of $2(p+1)$; in $13$ and $17$, where $p\equiv1\pmod4$, it is $2$, and the angular part
    contains exactly $i$.
  \item In the split classes $11,19$, where $p\equiv3\pmod4$, the $2$-adic valuation of $e_p$ is $1$:
    the two embeddings send $\golden$ to $\golden$ and $-1/\golden$, and one of the two has even order.
  \item In the kernel $\{1,9\}$ of $G$ the class does not fix it: $41$, $61$ and $101$, all in
    the class $1$, give $3$, $2$ and $1$.
  \item A factor $3$ in $e_p$, the hexagonal component of the order, requires $p\equiv1\pmod3$ for $p$ split and
    $p\equiv2\pmod3$ for $p$ inert, a condition on the class of $p$ modulo $60$.  The condition is necessary and
    does not suffice; what decides is this: with $N_p=3^km$ and $3\nmid m$, $3\mid e_p$ if and only if
    $\golden^{m}\not\equiv1\pmod p$.
\end{enumerate}
\end{proposition}
\leanline{\LeanStar{F1PaperPadic.inert\_norm\_is\_minus\_one}, \LeanStar{F1PaperPadic.split\_three\_mod\_four\_two\_part}, \LeanStar{F1PaperPadic.kernel\_classes\_do\_not\_fix\_it} \tierPnew{}
on the primes listed; (i)--(ii) for every prime, \LeanInline{PasosDelTexto.inert\_pow\_succ}, \LeanInline{PasosDelTexto.inert\_two\_part}, \LeanInline{PasosDelTexto.two\_part\_inert\_one\_mod\_four}, \LeanInline{PasosDelTexto.split\_three\_mod\_four} \tierP; (iv)
\LeanInline{RelojTres.three\_not\_dvd\_order\_iff}, \LeanInline{RelojTres.three\_dvd\_of\_three\_dvd\_order} \tierP; over the
primes $7\le p<6000$, (i)--(ii) hold in every case, (iv) holds at every prime, and of the primes that meet the
condition of (iv) $185$ split and $197$ inert, $42$ and $52$ of them have $e_p$ prime to $3$
(\texttt{retro119.py} \tierN); the first two, $47$ inert and $139$ split, are \LeanInline{TodoEnLean.witness\_47}, % pending(D7)
\LeanInline{TodoEnLean.witness\_139} \tierP.}
\begin{proof}[Proof of \textup{(iv)}]
$e_p\mid N_p$.  If $3\mid e_p$ then $3\mid N_p$, which is $3\mid p-1$ for $p$ split and $3\mid p+1$ for $p$ inert.
In the cyclic group generated by $\golden$ modulo $p$, of order $e_p$ dividing $3^km$, the order is prime to $3$
exactly when it divides $m$, that is, when $\golden^m\equiv1\pmod p$.
\end{proof}


So the eight classes fix the $2$-adic part of the angular part of $\golden$ in six of the eight
classes, and leave it free exactly in the kernel of $G$.  That kernel is the fibre $(0,0)$ of the
four fibres of $G$, the fibre of the trivial character, and its order is the contraction factor
$1/\mu_3=|\ker G|=2$ (\LeanInline{F1PCFUnified.four\_fibres\_four\_characters}, \LeanInline{F1PCFUnified.contraction\_is\_kernel} \tierP).
What the class does not fix is the odd part of $e_p$, and, in the kernel, its $2$-adic part too.


\begin{proposition}[The fibres of $G$ on the levels of the clock]\label{prop:fibres-levels}
Let $n\geq3$ and let $\mathfrak p$ be a prime of $\Zphi$ of clock $n$ over $p\neq2,5$.
\begin{enumerate}[label=\textup{(\roman*)}]
  \item If $n$ is odd or $n\equiv2\pmod4$, $p$ splits and lies in $\{1,9\}$ or $\{11,19\}$.  If
    $n\equiv4\pmod8$, $p$ lies in $\{1,9\}$ or $\{13,17\}$.  If $8\mid n$, $p$ lies in $\{1,9\}$ or
    $\{3,7\}$.  If $5\mid n$, $p$ splits, and if $20\mid n$, $p\equiv1\pmod{20}$.
  \item For odd $m\geq3$, $|N(\Phi_m(\golden))|=|N(\Phi_{2m}(\golden))|$.  For $4\mid n$, $n\geq8$,
    $\Phi_n(\golden)$ is a unit times a rational integer, and $|N(\Phi_n(\golden))|$ is its square.
  \item Of the mass of \cref{prop:clock-cyclotomic}\textup{(vii)}, the odd levels carry $\tfrac23$, the
    levels $2\pmod4$ carry $\tfrac16$, the levels $4\pmod8$ carry $\tfrac1{12}$ and the levels
    $0\pmod8$ carry $\tfrac1{12}$, each up to $O(N\log N)$.  So the fibre $\{13,17\}$ carries at most
    $\tfrac1{12}$ of the mass, the fibre $\{3,7\}$ at most $\tfrac1{12}$, the inert primes together at
    most $\tfrac16$, and the fibre $\{11,19\}$ at most $\tfrac56$.
\end{enumerate}
\end{proposition}
\leanline{\LeanInline{FibrasPorNivel.split\_level\_four},
\LeanInline{FibrasPorNivel.split\_level\_twenty}, \LeanInline{FibrasPorNivel.inert\_level\_four},
\LeanInline{FibrasPorNivel.inert\_one\_mod\_four}, \LeanInline{FibrasPorNivel.inert\_not\_at\_five} (i)
\tierP; (ii), the first statement, \LeanInline{PasosDelTexto.palindromic\_norm}, \LeanInline{PasosDelTexto.norm\_odd\_level\_doubled} \tierP\ given that $\Phi_m$ is palindromic and $\Phi_{2m}(x)=\Phi_m(-x)$ for odd $m\ge3$ \tierC, the second \tierT; (iii) from $\sum_{m\leq x,\ m\ \mathrm{odd}}\phi(m)=2x^2/\pi^2+O(x\log x)$ \tierC;
(i)--(iii) on every prime of every level $n\leq420$; there the inert primes carry between $0.066$ and
$0.084$ of the mass and the kernel between $0.47$ and $0.56$, with no stable value
(\texttt{retro94.py} \tierN).} % pending(D7)

\begin{proof}
(i) If $p$ splits, $n\mid p-1$; if $p$ is inert, $\golden^{p+1}=N(\golden)=-1$ in the residue field, so
$n\mid2(p+1)$ and $n\nmid p+1$, which for odd $p$ forces $4\mid n$, and $8\nmid n$ when
$p\equiv1\pmod4$.  So at $4\mid n$ a split prime has $p\equiv1\pmod4$, and an inert prime has
$p\equiv1\pmod4$ exactly when $n\equiv4\pmod8$.  If $5\mid n$ and $p$ were inert, $5\mid p+1$, and $p$
would split.
(ii) For odd $m\geq3$, $\Phi_m(-X)=\Phi_{2m}(X)$ and $\Phi_{2m}$ is self-reciprocal, so
$\Phi_m(\bar\golden)=\Phi_{2m}(1/\golden)=\golden^{-\phi(m)}\Phi_{2m}(\golden)$.  For $4\mid n$,
$\Phi_n(X)=\Phi_{n/2}(X^2)$, and $\Phi_{n/2}$ is $X^{\phi(n/2)/2}$ times an integer polynomial in
$X+1/X$; at $X=\golden^2$, $X+1/X=3$.
(iii) By \cref{prop:clock-cyclotomic}\textup{(iii)} the level $n$ weighs $\phi(n)\log\golden$ up to a
bounded term; with $n=2^km$, $m$ odd, $\phi(n)=2^{k-1}\phi(m)$ for $k\geq1$, and the sums over $m$ give
$\tfrac{2}{\pi^2}N^2$, $\tfrac1{2\pi^2}N^2$, $\tfrac1{4\pi^2}N^2$ and $\tfrac1{4\pi^2}N^2$ times
$\log\golden$.  The bounds follow from (i).
\end{proof}

This is \cref{prop:class-angular} read from the level: the kernel of $G$ is the only fibre present on
every level, and each of the other three is confined to one $2$-adic type of level.  The order of the
clock and the order by size weigh the fibres differently.  By size the inert prime ideals of norm at
most $x$ are the primes $p\leq\sqrt x$, a weight $O(\sqrt x)$, and the two split fibres weigh the same
\tierC; in the order of the clock the inert primes are a positive part of the mass on the levels
$4\mid n$, and the two split fibres are not symmetric, since only the kernel reaches those levels.

\subsubsection{Residue symbols and the index of $\golden$}\label{ssec:phi-index}

\begin{proposition}[The residue symbol of $\golden$ is its angle]\label{prop:phi-residue}
Let $p$ split in $\Zphi$, let $\mathfrak p$ be a prime above $p$, so that $\golden_{\mathfrak p}\in\mathbb{F}_p$ is a root
of $x^2-x-1$, and let $n\mid p-1$.
Then $\golden$ is an $n$-th power modulo $\mathfrak p$ if and only if $\golden_{\mathfrak p}^{(p-1)/n}=1$.
This power is the image in $\mu_n$ of the angular part $[\golden_{\mathfrak p}]$ of the polar form
(\cref{prop:delta-ring}); the principal unit contributes nothing.  For $n=5$: $\golden$ is not a fifth
power at $p=11$, and is one at $p=211$, where $\golden_{\mathfrak p}=33=94^5$.
\end{proposition}
\leanline{\LeanStar{F1PaperNew.units\_pow\_iff}, \LeanStar{F1PaperNew.phi\_residue\_symbol},
\LeanStar{F1PaperNew.phi\_fifth\_power\_instances}
\tierPnew.}

Write $i_p=(p-1)/e_p$ for $p$ split and
$i_p=2(p+1)/e_p$ for $p$ inert, the index of
$\golden$ at $p$; at an inert prime it is the index
of the subgroup generated by $r$ in $\mu_{2(p+1)}$.

\begin{proposition}[The index of $\golden$ by
layers]\label{prop:phi-layers}
Let $p\ne2,5$ and $r$ a root of $x^2-x-1$ modulo $p$
(in $\mathbb{F}_p$ if $p$ splits, in $\mathbb{F}_{p^2}$
if it is inert).
\begin{enumerate}[label=\textup{(\roman*)}]
  \item For $p$ split, $d\mid p-1$ and $(p-1)/d$
    even, $d\mid i_p$ if and only if
    $r^{(p-1)/d}=1$.  For $p$ inert and
    $d\mid2(p+1)$, $d\mid i_p$ if and only if
    $r^{2(p+1)/d}=1$.
  \item For $p$ split, $i_p$ is even if and only if
    $r$ and $1-r$ are both squares, and
    \[
      4\cdot[\,2\mid i_p\,]
      = 1+\Bigl(\tfrac{r}{p}\Bigr)
        +\Bigl(\tfrac{1-r}{p}\Bigr)
        +\Bigl(\tfrac{-1}{p}\Bigr).
    \]
  \item In the kernel $\{1,9\}$ of $G$, $i_p$ is even
    if and only if $p=a^2+20b^2$, and odd if and only
    if $p=4a^2+5b^2$.
\end{enumerate}
So each layer of the index is read by a
power-residue symbol of $\golden$: the quadratic for
$d=2$, the cubic for $d=3$, the quintic for $d=5$.
\end{proposition}
\leanline{\LeanStar{W4Isometria.split\_layer\_law},
\LeanStar{W4Isometria.inert\_layer\_law},
\LeanStar{W4Isometria.even\_index\_iff\_squares},
\LeanStar{W4Isometria.even\_index\_indicator}
\tierPnew;
(iii), stated as \LeanStar{W4Isometria.KernelLaw}
\tierO, follows from (ii) since $p\equiv1\pmod4$ there: $x^4-x^2-1$ generates the ring class field of
$\Z[\sqrt{-20}]$, and $p=a^2+20b^2$ exactly when
$p$ splits completely in it
\cite[Thm.~9.2, Thm.~9.4]{cox13} \tierC; checked on
the $555$ primes below $20000$ in those classes
\tierN.}

\begin{proposition}[The golden ratio names the fifth roots of unity]\label{prop:phi-names-mu5}
Let $p\equiv1\pmod5$ and let $\mathfrak p=(p,\golden-r)$ be a prime of $\Zphi$ above $p$.
\begin{enumerate}[label=\textup{(\roman*)}]
  \item In $\Zphi$, $\zeta_5+\zeta_5^{-1}=\golden-1$ and $\zeta_5^2+\zeta_5^{-2}=-\golden$, that is
    $2\cos(2\pi/5)=1/\golden$ and $2\cos(4\pi/5)=-\golden$.  Hence in $\mathbb F_p$ there is exactly one
    pair $\{z,z^{-1}\}$ of primitive fifth roots of unity with $z+z^{-1}=r-1$: the golden ratio fixes
    the identification of $\mu_5$ at $\mathfrak p$ up to the inversion $z\mapsto z^{-1}$.
  \item With this identification the quintic symbol of $\golden$ at $\mathfrak p$ falls in one of three
    canonical classes, $\{1\}$, $\{\zeta_5^{\pm1}\}$, $\{\zeta_5^{\pm2}\}$, with densities $\tfrac15$,
    $\tfrac25$, $\tfrac25$.
  \item The Galois conjugation $\golden\mapsto\bar\golden$ exchanges the classes $\{\zeta_5^{\pm1}\}$ and
    $\{\zeta_5^{\pm2}\}$ and fixes $\{1\}$: at the conjugate prime $z+z^{-1}=\bar\golden-1=-\golden$.
    Since a prime ideal and its conjugate have the same norm, every sum over prime ideals weighted by a
    function of the norm gives the two classes equal parts, exactly.
\end{enumerate}
\end{proposition}
\leanline{(i), (iii) \LeanInline{PasosDelTexto.fifth\_root\_trace}, \LeanInline{PasosDelTexto.second\_trace}, \LeanInline{PasosDelTexto.pair\_from\_trace}, \LeanInline{PasosDelTexto.conjugation\_swaps} \tierP; (ii) Chebotarev in the layer of group $F_{20}$
(\cref{prop:phi-layers}) \tierC; over the $39\,234$ prime ideals of norm $p\equiv1\pmod5$ below $10^6$
the classes have frequencies $0.2003$, $0.3999$, $0.3999$, the identification is unambiguous at every one,
and the two exchanged parts of the prime side coincide (\texttt{retro78.py} \tierN).} % pending(D7)

This is where the two faces of $\golden$ meet at the primes: $\golden=2\cos(\pi/5)$ names the angle, and
with it the root of unity in which the symbol of $\golden$ is read.  Two involutions act here, and they are
different.  The class $9$, the element of order two of $\ker G$, is $-1$ modulo $5$ and $1$ modulo $4$, so as an
automorphism of $\Q(\zeta_{20})=\Q(i,\zeta_5)$ it sends $\zeta_5$ to $\zeta_5^{-1}$ and fixes $i$: on $\mu_5$ it is
the inversion $z\mapsto z^{-1}$ of (i), and it fixes $\golden=1+\zeta_5+\zeta_5^{-1}$.  That is why the
identification of (i) is made up to it, and why the three classes of (ii) are unions of its orbits.  The Galois
conjugation $\golden\mapsto\bar\golden$ of (iii) is the action of the classes with $\chiv=-1$, which send $\zeta_5$
to $\zeta_5^{\pm2}$; it is the exact isometry of every such sum, the prime side of the explicit formula among
them (\LeanInline{RelojTres.class\_nine}, \LeanInline{RelojTres.classes\_that\_conjugate},
\LeanInline{TodoEnLean.class\_nine\_fixes\_phi} \tierP).  The rotation of order $5$, $z\mapsto\zeta_5z$, does not: the class
$\{1\}$ and the class $\{\zeta_5^{\pm1}\}$ are different primes, and their parts differ as much as a
random partition of the same primes does.

The clock fixes the angle of $\golden$ at each prime and the index fixes the rest of its
position; on the $n$-gon of each prime both are read exactly, and $N_p=e_p\,i_p$.

\begin{proposition}[The index on the $n$-gon of each prime]\label{prop:index-ngon}
Let $p\neq2,3,5$.
\begin{enumerate}[label=\textup{(\roman*)}]
  \item \emph{The columns.}  For $m\geq1$ the primes $p$ with $\golden^{m+1}\equiv1\pmod p$
    satisfy $\sum\log p\leq(m+1)\log\golden+4\golden^{-(m+1)}$; so the primes $p\leq X$ with
    $e_p\leq y+1$ weigh at most $\tfrac{y(y+3)}2\log\golden+4\golden^2$.
  \item \emph{The hyperbola.}  If $p\leq X$ and $i_p>j$, then $e_p\,j<2(X+1)$: a large index
    forces a small clock (with $N_p\leq2(p+1)$).  With $j=X^{3/4}$ and \textup{(i)}, the primes $p\leq X$ of index
    above $X^{3/4}$ weigh $O(\sqrt X)$, with no hypothesis.
  \item \emph{The $n$-gon.}  In a cyclic group of order $N$ the elements of index above $j$
    number $\sum_{c\mid N,\ jc<N}\phi(c)$, a fraction at most $\tau(N)/j$, with $\tau(N)$ the
    number of divisors of $N$; and for $i\mid N$ the elements whose index is a multiple of $i$
    are those with $x^{N/i}=1$, exactly $N/i$ of them, the fraction $1/i$.
  \item \emph{The quadratic layer as a sum.}  At a split prime $p\equiv1\pmod4$, the sum
    $a_p=(\tfrac rp)+(\tfrac{1-r}p)$ is $2$ if $2\mid i_p$ and $-2$ otherwise: $a_p=4\bigl([2\mid i_p]-\tfrac12\bigr)$.
\end{enumerate}
\end{proposition}
\leanline{\LeanStar{W4Isometria.HiperbolaRelojTamano.sum\_log\_le\_log\_phi},
\LeanStar{W4Isometria.HiperbolaRelojTamano.tail\_bound\_phi\_explicit} (i);
\LeanStar{W4Isometria.HiperbolaRelojTamano.large\_index\_clock} (ii);
\LeanStar{W4Isometria.GeometriaDelPoligono.card\_large\_index},
\LeanStar{W4Isometria.GeometriaDelPoligono.card\_large\_index\_le},
\LeanStar{W4Isometria.CapasDelPoligono.card\_layer},
\LeanStar{W4Isometria.CapasDelPoligono.layer\_iff} (iii);
\LeanStar{W4Isometria.ReduccionAlReloj.phi\_pow\_one\_split},
\LeanStar{W4Isometria.ReduccionAlReloj.level80\_coeff\_in},
\LeanStar{W4Isometria.ReduccionAlReloj.level80\_coeff\_out} (iv) \tierPnew; the consequence
in (ii) from the two statements, \LeanInline{PasosDelTexto.hyperbola}, \LeanInline{PasosDelTexto.hyperbola\_sqrt} \tierP.  Over the $78\,495$ primes $7\leq p\leq10^6$, $N_p\leq2(p+1)$
in all, and the primes of index above $X^{3/4}$ have mass $0$ at $X=10^6$, against the
bound $1.00\sqrt X$ that the tower gives
(\texttt{v28\_hiperbola\_reloj\_tamano.py}) \tierN.} % pending(D7)

\begin{remark}[The layers of the index as $L$-functions]\label{rmk:layers-L}
The three power-residue symbols of \cref{prop:phi-layers} reach
$L$-functions of their own by the same route:
a form of weight one and level $80$, a form of level
$135$ attached to $\chi_{-15}$, and the
representation of the quintic layer, of image $F_{20}$,
cut out by $x^5+5x^3+5x-1$; all three are Hecke
$L$-functions \tierC.  By \cref{prop:index-ngon}(iv), the
coefficient of the form of level $80$ at a split prime
$p\equiv1\pmod4$ is $4([2\mid i_p]-\tfrac12)$: the form reads
the layer $2$ of the index.
\end{remark}

% ---------------------------------------------------------------------
\subsection{From the character to the Euler factors}\label{sec:dedekind}


The character now has to reach an $L$-function, and it does so by the classical
route, with each step formalised.

\subsubsection{Splitting types and local factors}

A rational prime $p$ has one of three behaviours in $\OK=\Zphi$, and the three
are exhausted by the degree formula $\sum_ie_if_i=[K:\Q]=2$
(\LeanStar{F1CW.degree\_formula\_K} \tierP): split ($p\OK=\mathfrak{p}\bar{\mathfrak{p}}$,
two primes of residue degree $1$), inert (one prime of residue degree $2$), or
ramified (one prime with $e=2$, only $p=5$ since $\Delta_K=5$).  The absolute
norms follow: $N\mathfrak{p}=p$ in the split and ramified cases, $p^2$ in the
inert one (\LeanStar{F1CW.ideal\_norm\_values\_K} \tierP).  Writing $u=p^{-s}$, the
product of the local factors over the primes above $p$ is
\begin{equation}\label{eq:local-K}
  \prod_{\mathfrak{p}\mid p}\bigl(1-u^{f(\mathfrak{p})}\bigr)^{-1}
  =\begin{cases}(1-u)^{-2},&\text{split},\\ (1-u^2)^{-1},&\text{inert},\\ (1-u)^{-1},&\text{ramified},\end{cases}
\end{equation}
which is the statement that the local factor is a function of $\chiv(p)$
(\LeanStar{F1CW.euler\_from\_ideals\_K} \tierP): the first case is $\chiv(p)=+1$,
the second $\chiv(p)=-1$, the third $\chiv(p)=0$.  \Cref{eq:local-K} is also
where the Frobenius of \cref{sec:lambda-ring} enters analytically, since by
\cref{thm:frob-eigen} the inverse local factor is the characteristic polynomial
of $\Frob_p$ acting on $\Zphi$.

Two counts follow.  Expanding the product, the number of ideals of $\OK$ of norm $n$ is
$d_K(n)=\sum_{d\mid n}\chiv(d)$, the coefficient of $n^{-s}$ in $\zeta\cdot L(\cdot,\chiv)$ \tierC\
\cite[Ch.~VII]{neukirch99}: a mass level $n$ inherits its type from the classified primes that divide
it.  And since $\chiv$ is trivial on exactly four of the eight classes of $\Zst$ (\cref{prop:G}),
Dirichlet's theorem on primes in arithmetic progressions gives the split primes density $\tfrac12$ \tierC\
\cite[Ch.~7]{apostol76}.

\begin{proposition}[The Frobenius eigenvalue]\label{thm:frob-eigen}
In $\Zphi$, $\sqrt5=2\omega-1$ satisfies $\sqrt5^{\,2}=5$ and $\sigma(\sqrt5)=-\sqrt5$;
for $n$ prime to $5$, $\Frob_n(\sqrt5)=\chiv(n)\sqrt5$ and
$\Frob_n^{k}(\sqrt5)=\chiv(n)^k\sqrt5$.  For $p\neq5$, with $M_p=\mathrm{diag}(1,\chiv(p))$ the
matrix of the ring Frobenius $\Frob_p$ on $(1,\sqrt5)$, $\det(1-tM_p)=(1-t)(1-\chiv(p)t)$, the
inverse local factor of $\zeta_{\mathbb{Q}(\sqrt5)}=\zeta\cdot L(s,\chiv)$ at
$t=p^{-s}$; at $p=5$ the inverse local factor is $1-t$, the case $\chiv(5)=0$.
\end{proposition}
\leanline{\LeanStar{frobenius\_eigenvalue\_on\_sqrt5}, \LeanStar{frobenius\_power\_eigenvalue},
\LeanStar{euler\_factor\_is\_char\_poly} \tierPnew.}
\begin{proof}
$(2\omega-1)^2=4(\omega^2-\omega)+1=5$, and $\sigma(\omega)=1-\omega$ gives $\sigma(2\omega-1)=1-2\omega$.  By
\cref{thm:sign-is-chi5}, $\Frob_n$ is the identity when $\chiv(n)=+1$ and $\sigma$ when $\chiv(n)=-1$, which gives the
eigenvalue and its powers.  In the basis $(1,\sqrt5)$ the matrix of $\Frob_p$ is therefore $M_p$, and
$\det(1-tM_p)=(1-t)(1-\chiv(p)t)$, which at $t=p^{-s}$ is $(1-p^{-s})^2$ or $(1-p^{-s})(1+p^{-s})=1-p^{-2s}$, the
split and inert cases of \eqref{eq:local-K}.
\end{proof}

The two shapes of the local factor are ring identities, valid in any commutative
ring: $(1-t)(1-t)=(1-t)^2$ at a split prime and $(1-t)(1+t)=1-t^2$ at an inert one
(\LeanStar{euler\_factor\_split}, \LeanStar{euler\_factor\_inert} \tierPnew).
The character is the trace of the Frobenius minus one: $\chiv(p)=\mathrm{tr}\,M_p-1$
(\LeanStar{trace\_minus\_one\_is\_chi5} \tierPnew).

% ---------------------------------------------------------------------
\subsection{The Euler product as a categorical colimit}\label{sec:colimit}




The homomorphism $G\colon\Zst\to\mathbb{Z}_2\times\mathbb{Z}_2$ determines the
Euler factors of $\zeta_{\mathbb{Q}(\sqrt{5})}(s)$ through its component $\chiv$ (\cref{sec:dedekind}), and its
cardinals fix the modulus $\tfrac12=1/|\ker G|$ of the triad of \cref{def:tripartite-omega}.  This section states
the categorical content: the Frobenius system is a functor, the Euler product is a colimit, and the Dedekind
factorisation is a natural isomorphism.  Each term is used in its categorical sense, with the category, the diagram
and the universal property stated.

\subsubsection{Preliminaries}\label{ssec:cat-prelim}

\begin{definition}[Coefficient systems; the category of Dirichlet series]\label{def:DirSer}
\leavevmode
\begin{enumerate}[label=\textup{(\roman*)}]
  \item A \emph{coefficient system} is a pair $(X,c)$ of a set $X$ of positive integers and a function
    $c\colon X\to\C$.  With $(X,c)\le(Y,d)$ when $X\subseteq Y$ and $d$ restricts to $c$ on $X$, the coefficient
    systems form a partially ordered set, and so a category $\mathbf{Coef}$ with at most one morphism between two
    objects; a colimit in it is a least upper bound.  The Dirichlet series of $(X,c)$ is
    $D(X,c)(s)=\sum_{n\in X}c(n)\,n^{-s}$, where it converges absolutely.
  \item The category $\mathbf{DirSer}$ has as objects the Dirichlet series $F=\sum a_nn^{-s}$ with $a_1=1$,
    absolutely convergent for $\Re s>1$; a morphism $F\to G$ is a Dirichlet series $H$ of the same kind with
    $G=F\cdot H$, the product being the Dirichlet convolution $(a*b)_n=\sum_{d\mid n}a_d\,b_{n/d}$ of the
    coefficients; composition is the product, and the identity is $1$.  Since $a_1=1$, $F$ is invertible as a formal
    Dirichlet series, so $H$, when it exists, is unique: $\mathbf{DirSer}$ has at most one morphism between two
    objects.
\end{enumerate}
\end{definition}

\begin{definition}[Spectral data]\label{def:SpecOmega}
A \emph{spectral datum} is a triple $(\sigma,\mu,(\lambda_k)_{k=0,1,2})$ with $\sigma,\mu$ real and
\[
  \sigma+\mu=\frac{|\operatorname{im}G|}{|\ker G|}=2,\qquad
  \sigma\mu=\frac{|\operatorname{im}G|-1}{|\operatorname{im}G|}=\frac34,\qquad
  \lambda_k=\mu\,\zeta_3^{\,k},
\]
with the cardinals $|\operatorname{im}G|=4$ and $|\ker G|=2$ of \cref{prop:G}.  $\mathrm{Spec}(\OmH)$ is the discrete
category on the spectral data with $0<\mu<1$; the notation anticipates the colour operator $\OmH$ of
\cref{sec:operator}, whose spectrum is the triad.  $\mathbf{B}\Zst$ is the category with one object and the eight
elements of $\Zst$ as morphisms, and $F_{\mathrm{spec}}\colon\mathbf{B}\Zst\to\mathrm{Spec}(\OmH)$ sends every
morphism to an identity.
\end{definition}

\begin{proposition}[Two spectral data, one below $1$]\label{prop:spectral-data}
The spectral data are $(\tfrac32,\tfrac12)$ and $(\tfrac12,\tfrac32)$, each with its triad, and exactly one has
$\mu<1$.  So $\mathrm{Spec}(\OmH)$ has one object and one morphism, it is the terminal category, and
$F_{\mathrm{spec}}$ is the unique functor to it.  Its triad $\tfrac12\zeta_3^{\,k}=|2|_2\,\zeta_3^{\,k}$ is the triad
of the inert prime $2$ (\cref{def:tripartite-omega}).
\end{proposition}
\leanline{\LeanInline{LadoAnalitico.spectral\_data}, \LeanInline{LadoAnalitico.spectral\_datum\_below\_one},
\LeanInline{spectral\_uniqueness} \tierP.}
\begin{proof}
With $\sigma=2-\mu$, $\sigma\mu=\tfrac34$ is $(\mu-\tfrac12)(\mu-\tfrac32)=0$.  A category with one object and one
morphism is terminal.
\end{proof}

\begin{lemma}[Diagonal contraction]\label{lem:diagonal-contraction}
For $0<t\le\mu_3=\tfrac12$ the inequality $t\le t^2$ fails.
\end{lemma}
\leanline{\LeanInline{diagonal\_blocked} \tierP.}

\begin{proof}
$t\le t^2$ with $t>0$ gives $1\le t$, against $t\le\tfrac12$.  (The same argument
gives $t^2<t$ on all of $(0,1)$, which is \LeanStar{F1v2.sq\_lt\_self\_on\_unit\_interval}.)
\end{proof}

\begin{definition}[Normed monoidal category; contractive morphism]\label{def:normed-monoidal}
A \emph{norm} on a monoidal category $(\mathcal C,\otimes)$ is a function $X\mapsto\norm{X}$ from its objects to
the positive reals with $\norm{X\otimes Y}=\norm{X}\,\norm{Y}$.  A morphism $f\colon X\to Y$ is \emph{contractive} if
$\norm{Y}\le\norm{X}$.
\end{definition}

\begin{theorem}[No contractive evaluation below norm one]\label{thm:no-diagonal}
Let $\Omega$ be an object of a normed monoidal category, with $t=\norm{\Omega}>0$.  A contractive morphism
$\Omega\otimes\Omega\to\Omega$ requires $t\le t^2$, which holds exactly when $t\ge1$.  So if $0<t<1$, no morphism
$\mathrm{eval}\colon\Omega\otimes\Omega\to\Omega$ is contractive, and in particular no diagonal
$\Delta\colon\Omega\to\Omega\otimes\Omega$ is split, $\mathrm{eval}\circ\Delta=\mathrm{id}_\Omega$, by a contractive
evaluation.  This is the No-Diagonal theorem of the Primitive Structure~\cite{gonzalez25hp}.
\end{theorem}
\leanline{\LeanInline{LadoAnalitico.contractive\_iff}, \LeanInline{LadoAnalitico.no\_contractive\_eval},
\LeanInline{diagonal\_blocked}, \LeanStar{F1v2.sq\_lt\_self\_on\_unit\_interval}, \LeanInline{contraction\_is\_kernel} \tierP;
of the category the proof uses the multiplicativity of the norm.}
\begin{proof}
$\mathrm{eval}$ is contractive when $\norm{\Omega}\le\norm{\Omega\otimes\Omega}=\norm{\Omega}^2$, that is, $t\le t^2$.
For $t>0$ this is $1\le t$; for $t\le\tfrac12$ it is \cref{lem:diagonal-contraction}.  If $\mathrm{eval}\circ\Delta=
\mathrm{id}_\Omega$ with $\mathrm{eval}$ contractive, then $t\le t^2$, which $0<t<1$ excludes.
\end{proof}

\begin{corollary}[Spectral consequence]\label{cor:no-diagonal-spectral}
Take the modulus $\mu$ of a spectral datum as the norm of $\Omega$.  The condition $\mu<1$ of $\mathrm{Spec}(\OmH)$
\textup{(\cref{def:SpecOmega})} is exactly the condition under which \cref{thm:no-diagonal} excludes a contractive
evaluation: at $\mu=\tfrac12$ none exists, and at $\mu=\tfrac32$ the inequality $\mu\le\mu^2$ holds, so the norm
gives no obstruction there.  The datum of $\mathrm{Spec}(\OmH)$ has $\norm{\OmH}=\tfrac12<1$, and its triad lies on the
circle $|z|=\tfrac12$.
\end{corollary}
\leanline{\LeanInline{LadoAnalitico.contractive\_iff}, \LeanInline{LadoAnalitico.spectral\_data},
\LeanInline{spectral\_uniqueness}, \LeanInline{diagonal\_blocked} \tierP.}
\begin{proof}
By \cref{prop:spectral-data} the moduli are $\tfrac12$ and $\tfrac32$, and by \cref{thm:no-diagonal} the inequality
$\mu\le\mu^2$ holds exactly when $\mu\ge1$.
\end{proof}

\subsubsection{The Frobenius system as a functor}
\label{ssec:frobenius-functor}

\begin{definition}[Frobenius functor]\label{def:frobenius-functor}
The \emph{Frobenius functor} of $\langle\golden\rangle$ is the
monoid homomorphism
\begin{equation}\label{eq:Psi-functor}
  \Psi\colon (\mathbb{N}_{>0},\times)
  \longrightarrow \mathrm{End}_{\mathrm{Mon}}(\langle\golden\rangle),
  \qquad
  \Psi(n)(\golden^k) = \golden^{nk}.
\end{equation}
On the generator, $\Psi(n)(\golden) = \golden^n = F_n\golden + F_{n-1}$.  A monoid homomorphism is a functor between
the corresponding one-object categories, $\mathbf B(\N_{>0},\times)\to\mathbf B\,\mathrm{End}_{\mathrm{Mon}}(\langle\golden\rangle)$,
and $\Psi(p)$ is the monoidal lift $\Frob_p$ of \cref{thm:lambda-monoid}.
\end{definition}

\begin{proposition}[Functoriality of $\Psi$]\label{prop:Psi-functorial}
$\Psi$ preserves composition and identity:
$\Psi(mn) = \Psi(m)\circ\Psi(n)$ and
$\Psi(1) = \mathrm{id}_{\langle\golden\rangle}$.
The restriction to primes generates $\Psi$ through
$\Psi(p_1^{a_1}\cdots p_k^{a_k}) = \Frob_{p_1}^{a_1}\circ\cdots\circ\Frob_{p_k}^{a_k}$, the monoidal lifts.
\end{proposition}
\leanline{\LeanInline{psi\_comp} \tierP.}

\begin{proof}
On the monoid generator,
$\Psi(mn)(\golden) = \golden^{mn} = (\golden^m)^n
= \Psi(n)(\Psi(m)(\golden))$.  Since $\langle\golden\rangle$
is generated multiplicatively by~$\golden$, the monoid
endomorphism is determined by its action on~$\golden$.
\end{proof}


\subsubsection{The Euler product as a categorical colimit}
\label{ssec:euler-colimit}

\begin{definition}[Local Euler factors]\label{prop:local-factors}
The local Euler factor of $\zeta_{\mathbb{Q}(\sqrt{5})}(s)$ at a
prime~$p$ is
\[
  f_p(s) = \begin{cases}
    (1-p^{-s})^{-1}   & \chi_5(p) = 0  \text{ (ramified; only }p=5\text{)},\\
    (1-p^{-s})^{-2}   & \chi_5(p) = +1 \text{ (split)},\\
    (1-p^{-2s})^{-1}  & \chi_5(p) = -1 \text{ (inert)}.
  \end{cases}
\]
\end{definition}

\begin{definition}[Partial Euler products]\label{def:partial-euler}
Let $\mathcal{P}$ be the category whose objects are the finite
sets~$S$ of primes and whose morphisms are the inclusions
$S\hookrightarrow S'$, and let $N_S$ be the set of positive integers whose prime factors lie in $S$.  The
product $\prod_{p\in S}f_p(s)$ is a Dirichlet series whose coefficients vanish off $N_S$; let $E(S)$ be the
coefficient system formed by $N_S$ and those coefficients, so that
\begin{equation}\label{eq:partial-euler}
  D(E(S))(s) = \prod_{p\in S} f_p(s).
\end{equation}
Each $f_p$ depends only on $\chi_5(p) = \pi_2(G(p\bmod 20))$
(\cref{prop:G}).
\end{definition}

\begin{lemma}[From the lifts to the Euler factors]\label{lem:psi-to-euler}
Let $p\neq5$.  The coordinates of $\Psi(p)(\golden)=\golden^p=F_p\golden+F_{p-1}$, read modulo $p$, decide the local
factor: $(F_p,F_{p-1})\equiv(1,0)\pmod p$ exactly when $p$ splits, and $f_p$ is the split factor;
$(F_p,F_{p-1})\equiv(-1,1)\pmod p$ exactly when $p$ is inert, and $f_p$ is the inert factor.  So $f_p$ is determined by
$\Psi(p)$ together with the reduction of its value modulo $p$ in the ring $\Zphi$, which is where the ring Frobenius
lives; the monoid alone does not determine it (\cref{rmk:why-monoidal}).
\end{lemma}
\leanline{\LeanInline{FrobeniusUnico.fib\_eq\_sign}, \LeanInline{FrobeniusUnico.fib\_pred\_eq\_sign},
\LeanInline{FrobeniusUnico.sign\_eq\_chi5}, \LeanInline{local\_factor\_identity} \tierP.}
\begin{proof}
By \cref{thm:frobenius-congruence}, $\golden^p\equiv\Frob_p(\golden)\pmod{p\Zphi}$, and $\Frob_p(\golden)$ is $\golden$
or $\sigma(\golden)=1-\golden$; since $(\golden,1)$ is a basis of $\Zphi/p$ over $\F_p$, the coordinates are
$(1,0)$ or $(-1,1)$ modulo $p$, and the two pairs differ modulo every $p$ in the second coordinate.  The first case
is $\varepsilon(p)=+1$ and the second $\varepsilon(p)=-1$ (\cref{def:artin-frob}), that is, $\chiv(p)=\pm1$
(\cref{thm:sign-is-chi5}), and $f_p$ is fixed by $\chiv(p)$ (\cref{prop:local-factors}).
\end{proof}

\begin{proposition}[Euler product as colimit]\label{prop:euler-colimit}
Let $a_n$ be the coefficients of $\zeta_{\Q(\sqrt5)}(s)$.
\begin{enumerate}[label=\textup{(\roman*)}]
  \item For $S\subseteq S'$ the coefficient of $n\in N_S$ in $E(S')$ is its coefficient in $E(S)$, and both are
    $a_n$; so $E(S)\le E(S')$, and $E$ is a functor $\mathcal P\to\mathbf{Coef}$.
  \item The coefficient system $(\N_{>0},a)$ is the colimit of $E$: it is an upper bound of every $E(S)$, and
    every upper bound $(Y,d)$ of all the $E(S)$ is an upper bound of $(\N_{>0},a)$.
  \item For $\Re s>1$ the Dirichlet series of the colimit is $\zeta_{\Q(\sqrt5)}(s)$, the limit of the
    $D(E(S))(s)$ along the finite sets of primes ordered by inclusion:
    \begin{equation}\label{eq:euler-colimit}
      \zeta_{\mathbb{Q}(\sqrt{5})}(s)
      \;=\; D\Bigl(\varinjlim_{S\in\mathcal{P}} E(S)\Bigr)(s)\;=\;\lim_{S}\prod_{p\in S}f_p(s).
    \end{equation}
\end{enumerate}
\end{proposition}
\leanline{\LeanStar{F1CW.coeff\_stabilises\_K}, \LeanStar{F1CW.colimit\_exists\_K} \tierP;
\LeanInline{LadoAnalitico.mono}, \LeanInline{LadoAnalitico.cover}, \LeanInline{LadoAnalitico.upper},
\LeanInline{LadoAnalitico.least} \tierP.}
\begin{proof}
(i) Each $f_p$ is a series in $p^{-s}$ with constant term $1$.  The coefficient of $n=\prod p^{k_p}$ in a product
over $S'$ collects the terms $p^{-k_ps}$ of the factors at the primes of $n$ and the constant term $1$ of the others,
so for $n\in N_S$ it is the same over $S$, over $S'$ and over all primes.
(ii) $\N_{>0}=\bigcup_SN_S$, since $n$ lies in $N_S$ for $S$ the set of its prime factors, and by (i) every $E(S)$
is below $(\N_{>0},a)$.  If $(Y,d)$ is above every $E(S)$, then each $n$ lies in $Y$ with $d(n)=a_n$, by taking $S$
the set of its prime factors.
(iii) For $\Re s>1$ the Euler product converges absolutely, so the partial products converge to it along the
directed set of finite sets of primes.
\end{proof}

\tier{R} The predecessor version
\cite{gonzalez26oddzeta} read this colimit beside the tower limit $\|\Omega\|_k\to0$ of
\cite{gonzalez25hp}: both are limits of diagrams fixed by $G$, the first assembled from $\chiv(p)$, the
second iterated from $\mu_3=\tfrac12$.

\subsubsection{The Dedekind factorisation as a natural isomorphism}
\label{ssec:dedekind-natural}

\begin{definition}[Category of quadratic extensions]\label{def:QuadExt}
Let $\mathbf{QuadExt}_{20}$ be the category whose objects are the
fields $\mathbb{Q}$, $\mathbb{Q}(\sqrt{-1})$, $\mathbb{Q}(\sqrt{5})$ and
$\mathbb{Q}(\sqrt{-5})$ --- $\Q$ and the quadratic fields with discriminant dividing
$20$ --- with morphisms the identities and the field inclusions
$\mathbb{Q}\hookrightarrow K$.  For each $K$ let $\chi_K$ be its quadratic character, with $\chi_\Q=1$.
\end{definition}

The three quadratic fields of the conductor have discriminants $-4$, $5$ and $-20$ \tierC: the
two imaginary ones are negative and the real one positive, and their product is
$400=20^2$, the square of the conductor (\LeanInline{three\_subfield\_discs},
\LeanInline{disc\_K\_is\_conductor\_squared} \tierP\ for the arithmetic of the classical values).

\begin{theorem}[Dedekind factorisation as natural isomorphism]
\label{thm:dedekind-natural}
\leavevmode
\begin{enumerate}[label=\textup{(\roman*)}]
  \item For each $K$ of $\mathbf{QuadExt}_{20}$, $\zeta_K(s)=\zeta(s)\cdot L(s,\chi_K)$
    \textup{(\cite{neukirch99}, Ch.~VII)}.
  \item $K\mapsto\zeta_K$, with the inclusion $\Q\hookrightarrow K$ sent to the morphism $L(\cdot,\chi_K)\colon\zeta\to\zeta_K$
    of $\mathbf{DirSer}$, is a functor $\zeta_{(-)}$; and $K\mapsto\zeta\cdot L(\cdot,\chi_K)$, with the inclusion sent
    to $L(\cdot,\chi_K)$, is a functor $\zeta_\Q\otimes L_{(-)}$.
  \item The components $\eta_K=1\colon\zeta_K\to\zeta\cdot L(\cdot,\chi_K)$ form a natural isomorphism
    $\eta\colon\zeta_{(-)}\xrightarrow{\sim}\zeta_\mathbb{Q}\otimes L_{(-)}$.  Its content is in its components: that
    $\eta_K=1$ is a morphism is \textup{(i)} at $K$.
\end{enumerate}
The four instances of \textup{(i)} are:

\begin{center}
\begin{tabular}{ccll}
\toprule
$K$ & $\chi_K$ & Dedekind factorisation & split primes \\
\midrule
$\mathbb{Q}$ & $1$ &
  $\zeta = \zeta\cdot 1$ (identity) & all \\
$\mathbb{Q}(\sqrt{-1})$ & $\chi_4$ &
  $\zeta_{\mathbb{Q}(i)} = \zeta\cdot L(s,\chi_4)$ & $p\equiv1\pmod4$ \\
$\mathbb{Q}(\sqrt{5})$ & $\chi_5$ &
  $\zeta_{\mathbb{Q}(\sqrt{5})} = \zeta\cdot L(s,\chi_5)$ \textbf{(formalised here)} & $p\equiv\pm1\pmod5$ \\
$\mathbb{Q}(\sqrt{-5})$ & $\chi_4\chi_5$ &
  $\zeta_{\mathbb{Q}(\sqrt{-5})} = \zeta\cdot L(s,\chi_4\chi_5)$ & $p\equiv1,3,7,9\pmod{20}$ \\
\bottomrule
\end{tabular}
\end{center}
The four characters $\{1,\chi_4,\chi_5,\chi_4\chi_5\}$ are exactly
the four characters of $\mathbb{Z}_2\times\mathbb{Z}_2$
\textup{(\cref{prop:G})}; as $G$ has four fibres, the image has four characters.
\end{theorem}
\leanline{\LeanInline{dedekind\_factorisation}, \LeanInline{local\_factor\_identity} \tierP\ for $\Q(\sqrt5)$;
for every quadratic character, the Euler product with the local factors of its field equals $\zeta\cdot L(\cdot,\chi_K)$,
with the split, inert and ramified shapes: \LeanInline{TodoEnLean.dedekind\_factorisation\_any},
\LeanInline{TodoEnLean.fDed\_shapes} \tierP; that these local factors are those of $\Q(i)$ and $\Q(\sqrt{-5})$ is the
splitting law of those fields \tierC, and the classes on which each character is $+1$ are
\LeanInline{three\_quadratic\_subfields} \tierP; the morphisms of $\mathbf{DirSer}$ are unique,
\LeanInline{TodoEnLean.dirser\_morphism\_unique} \tierP.}

\begin{proof}
(i) Classical \textup{(\cite{neukirch99}, Ch.~VII)}; for $\Q(\sqrt5)$ it is the local identity behind \eqref{eq:local-K}, assembled over all primes and formalised as \lean{dedekind\_factorisation}.
(ii) A morphism $\zeta\to\zeta_K$ of $\mathbf{DirSer}$ is a series $H$ with $\zeta_K=\zeta\cdot H$, and by (i)
$H=L(\cdot,\chi_K)$, which has $a_1=1$ and converges absolutely for $\Re s>1$; the identity of $\Q$ goes to $1$,
since $L(\cdot,\chi_\Q)=1$, and no two non-identity morphisms of $\mathbf{QuadExt}_{20}$ compose.
(iii) By (i), $\zeta\cdot L(\cdot,\chi_K)=\zeta_K\cdot1$, so $\eta_K=1$ is a morphism, with inverse $1$.  At the inclusion
$\Q\hookrightarrow K$ the two composites of the naturality square are $L(\cdot,\chi_K)\cdot1$ and
$1\cdot L(\cdot,\chi_K)$, which are equal.
\end{proof}

\begin{corollary}[The factorisation is uniform in the character]
\label{cor:factorisation-universal}
The identity
$\zeta(s) = \zeta_{\mathbb{Q}(\sqrt{5})}(s)/L(s,\chi_5)$ is the
$K = \mathbb{Q}(\sqrt{5})$ component of a natural isomorphism~$\eta$
defined on all of $\mathbf{QuadExt}_{20}$.  What is proved in Lean is the local
statement: for each of the four characters the local factor of $\zeta_K$ is the
product of those of $\zeta$ and $L$, uniformly in the character value $0,\pm1$.
\end{corollary}
\leanline{\LeanStar{F1v2.factorisation\_natural}, \LeanStar{F1v2.local\_dedekind\_uniform}, \LeanStar{F1v2.four\_fibres\_are\_sign\_pairs} \tierPnew.}



\subsubsection{The fundamental formula and the regulator $R_K = \log\golden$}
\label{ssec:fund-formula}




\begin{theorem}[Fundamental formula]\label{thm:L1}
\begin{equation}\label{eq:L1}
  L(1,\chi_5) \;=\; \frac{2\log\golden}{\sqrt{5}}
  \;=\; \frac{T_K}{\pi\sqrt{5}},
\end{equation}
where $T_K = 2\pi\log\golden$ is the period of $T^2_{\mathrm{PCF}}$.
\end{theorem}
\leanline{\LeanInline{anchor\_is\_a\_crossed\_identity} \tierP\ given the class-number formula \tierC.}

\begin{proof}
Dirichlet class-number formula for $\mathbb{Q}(\sqrt{5})$.  With
class number $h_K = 1$ --- the Minkowski bound of $K$ is $\tfrac12\sqrt5<2$, so every ideal class contains an ideal
of norm $1$ (\LeanStar{F1CW.class\_number\_one\_K} \tierP) ---, regulator $R_K = \log\golden$ (\cref{thm:units}), number of
roots of unity $w = 2$, and discriminant $\Delta = 5$:
\[
  \lim_{s\to 1}(s-1)\,\zeta_K(s)
  \;=\; \frac{2^{r_1}(2\pi)^{r_2}\,h_K\,R_K}{w\sqrt{\Delta}}
  \;=\; \frac{4\cdot 1\cdot 1\cdot\log\golden}{2\sqrt{5}}
  \;=\; \frac{2\log\golden}{\sqrt{5}}.
\]
Since $\zeta_K(s) = \zeta(s)\cdot L(s,\chi_5)$ and
$\lim_{s\to 1}(s-1)\zeta(s) = 1$, we obtain
$L(1,\chi_5) = 2\log\golden/\sqrt{5}$.
\end{proof}

The pentagon gives the same value by a second route, without the class number.  For an even primitive
character of conductor $5$ the Gauss--Lerch formula reads
$L(1,\chiv)=-\frac{\tau(\chiv)}{5}\sum_{a=1}^{4}\chiv(a)\log\bigl(2\sin\tfrac{\pi a}5\bigr)$ \tierC\
\cite[Ch.~4]{washington97}; with $\tau(\chiv)=\sqrt5$ and the logarithmic signature
$\sum_{a=1}^{4}\chiv(a)\log(2\sin(\pi a/5))=-2\log\golden$ of the pentagon
(\LeanStar{F1v2.pentagon\_log\_signature\_full}), it is $-\frac{\sqrt5}{5}(-2\log\golden)=2\log\golden/\sqrt5$
(\LeanInline{F1OZ.lerch\_arith} \tierP).

\begin{remark}[The unifying invariant
$R = \log\golden$]\label{rmk:regulator}
The regulator $R_K = \log\golden$ of~$K = \mathbb{Q}(\sqrt{5})$
is the unifying invariant of the construction: it appears
\textup{(}equivalently up to factors of~$\pi$\textup{)} as the
residue of $\zeta_K$ at $s=1$, as the factor in the period
$T_K = 2\pi\log\golden$ of the PCF torus of \cite[\S2.8]{gonzalez26cw}, and in the base
scale $\varepsilon_0 = \ln\golden/(6\sqrt{3})$ of the metric tower of that paper.  Three classical roles---%
class-number residue, torus period, tower generator---are
realised by the same arithmetic invariant of~$K$.
\end{remark}

\begin{remark}[The prism at $5$]\label{rmk:prism-five}
At $5$, where $\Zphi$ has no Frobenius lift (\cref{thm:obstruction-five}), the ring still has a
prism, and the regulator of \cref{rmk:regulator} has a $5$-adic twin.  Let $K_5=\Q_5(\sqrt5)$ and
$\mathcal O_{K_5}=\Z_5[\golden]=\Z_5[\sqrt5]$, with uniformiser
$\sqrt5$ and ramification index $e=2$, and let $A=\Z_5[[u]]$ with the Frobenius $\varphi_A(u)=u^5$ and
the Eisenstein polynomial $E(u)=u^2-5$, generating the ideal $I=(E)$.
\begin{enumerate}[label=\textup{(\roman*)}]
  \item $\varphi_A(E)-E^5=5D(u)$ with $D=5u^8-50u^6+250u^4-625u^2+624$, so $\delta(E)=D\equiv-1
    \pmod5$ is a unit, $E$ is distinguished, and $(A,(E))$ is a prism with $A/I=\mathcal O_{K_5}$
    \cite{bhattscholze22}.
  \item Its special fibre $A/(5,E)=\F_5[u]/(u^2)$ is the ring $\F_5[\varepsilon]$ of the proof of
    \cref{thm:obstruction-five}: $\delta_5$ is the reduction of the prism, with $u=\sqrt5\mapsto
    2\varepsilon$.
  \item At $u=\sqrt5$, $\varphi_A(E)(\sqrt5)=\sqrt5^{\,10}-5=\tau(\chiv)^2\cdot624$, with
    $\tau(\chiv)=\sqrt5$ the Gauss sum and $624$ a unit: the candidate lifts on $\mathcal O_{K_5}$
    fail by $\tau(\chiv)$, and in the prism the defect is $\tau(\chiv)^2$, absorbed as
    $5\in I+\varphi_A(I)A$.
  \item The Kummer layer of $\golden$, the $5$-torsion of the $1$-motive
    $[\Z\to\mathbb G_m,\ 1\mapsto\golden]$, is a finite flat group scheme over $\mathcal O_{K_5}$.
    With $\langle\golden\rangle$ the principal unit of the polar form (\cref{ssec:delta-polar}),
    $v_5(\langle\golden\rangle-1)=v_5(\golden^4-1)=\tfrac12$ (\cref{rmk:angles-cyclotomic}, $k=0$), so its upper
    ramification break over $K_5$ is $\tfrac32=e\,(1+\tfrac1{p-1})-1$, Fontaine's bound
    \cite{fontaine85}, which it attains; with $v_5(\langle\golden\rangle-1)=1$ it would be
    $\tfrac12$.  Over $\Q_5$ the break is $\tfrac34$ and the conductor exponent $4(1+\tfrac34)=7$.
  \item Its Kisin module modulo $5$ \cite{kisin06} has the classes $h\in\F_5u\oplus\F_5u^2$, the
    cokernel of $g\mapsto g^5-u^2g$ on $\F_5[[u]]$, of breaks $\tfrac32$ and $\tfrac12$; from
    $\golden^4=1+4\sqrt5+5(1-\golden)$, the class of $\golden$ is $h=u$, the one of Fontaine's
    bound.
  \item Its $5$-adic class is $\log_5\golden$, of valuation $\tfrac12$, and Leopoldt's formula
    gives $L_5(1,\chiv)=2\log_5\golden/\sqrt5$ \cite[Thm.~5.18]{washington97}, the twin of
    $L(1,\chiv)=2\log\golden/\sqrt5$ \tier[\testref{test\_l\_one\_chi5\_digamma\_value}]{N} (\cref{thm:L1}): the two regulators of $\Q(\sqrt5)$ at $s=1$.
\end{enumerate}
\end{remark}
\leanline{\LeanStar{F1PaperSixty.bk\_delta\_E}, \LeanStar{F1PaperSixty.bk\_E\_distinguished}
(i), \LeanStar{F1PaperSixty.bk\_special\_fiber\_is\_delta5} (ii),
\LeanStar{F1PaperSixty.bk\_frobenius\_defect\_is\_tau\_sq} (iii),
\LeanStar{F1PaperSixty.fontaine\_bound\_saturated} (iv, the arithmetic),
\LeanStar{F1PaperSixty.phi\_four\_kisin\_class}, \LeanStar{F1PaperSixty.kisin\_classes\_breaks}
(v, the arithmetic), \LeanStar{F1PaperSixty.leopoldt\_unit} (vi) \tierPnew; the prism, the
breaks, the Kisin classes and Leopoldt's formula \tierC; the cokernel of (v) computed up to degree
$33$, and $L_5(1,\chiv)$ and $2\log_5\golden/\sqrt5$ both $\equiv557\pmod{5^4}$ against the
Kubota--Leopoldt interpolation (\texttt{kisin.py}) \tierN.}

% ---------------------------------------------------------------------
\subsection{The factorisation and the odd zeta values}\label{sec:factorisation}

\subsubsection{The factorisation}

Each of the three cases factors as the product of the local factor of $\zeta$ and
the local factor of $L(\cdot,\chiv)$ --- the inert case being the identity
$1-u^2=(1-u)(1+u)$
(\LeanStar{F1CW.local\_dedekind\_ramified\_K}, \LeanStar{F1CW.local\_dedekind\_split\_K},
\LeanStar{F1CW.local\_dedekind\_inert\_K} \tierP).  Assembling over all primes, with
the Euler product taken as the colimit of the partial products over finite sets
of primes (\LeanStar{F1CW.coeff\_stabilises\_K}, \LeanStar{F1CW.colimit\_exists\_K} \tierP),
gives
\begin{equation}\label{eq:dedekind}
  \zeta_{\Q(\sqrt5)}(s)=\zeta(s)\,L(s,\chiv)
  \qquad(\LeanInline{dedekind\_factorisation}\ \tierP).
\end{equation}

\subsubsection{The pentagonal determination of the odd zeta values}\label{ssec:odd-zeta}

The factorisation is formalised with genuine analytic objects, not placeholders:
$\zeta$ is Mathlib's \texttt{riemannZeta}, $L(s,\chiv)$ is the $L$-series of the
Dirichlet character $\chiv$ with its Euler product, and $\zeta_K$ is \emph{defined}
as the product over primes of the local factors \eqref{eq:local-K}, each one fixed
by the splitting type, that is, by $\chiv(p)$ and hence by the Frobenius of
\cref{sec:lambda-ring}.

\begin{theorem}[Pentagonal determination]\label{thm:odd-zeta}
For every $k\ge1$,
\[
  \zeta(2k+1)=\frac{\zeta_K(2k+1)}{L(2k+1,\chiv)},\qquad
  \zeta_K(s)=\prod_p f_K(p,s),
\]
and the denominator is a combination of Hurwitz zeta functions at the four
pentagonal points,
\[
  L(s,\chiv)=5^{-s}\Bigl(\zeta(s,\tfrac15)-\zeta(s,\tfrac25)-\zeta(s,\tfrac35)+\zeta(s,\tfrac45)\Bigr)
  \qquad(\Re s>1).
\]
\end{theorem}
\leanline{\LeanInline{zeta\_odd\_pentagonal\_determination}, \LeanInline{dedekind\_factorisation},
\LeanInline{hurwitz\_representation\_L\_chi5} \tierP, with axioms
\texttt{[propext, Classical.choice, Quot.sound]} only, the first under the hypothesis
$L(2k+1,\chiv)\neq0$; the hypothesis is discharged from Mathlib's non-vanishing of $L(s,\chiv)$ on
$\Re s\ge1$ in \LeanStar{F1PaperNew.zeta\_odd\_pentagonal\_determination\_unconditional} \tierPnew.}

Every object on the right-hand side is built from the pentagon, and $\zeta$ itself
appears nowhere on it: the local factors of $\zeta_K$ are the products over the prime
ideals of $\Zphi$ above $p$, $\prod_{\mathfrak p\mid p}(1-N\mathfrak p^{-s})^{-1}$, with
the splitting type decided by the Frobenius of $\Zphi$
(\LeanStar{F1CW.euler\_from\_ideals\_K}, \LeanStar{F1CW.ideal\_norm\_values\_K},
\LeanStar{F1CW.degree\_formula\_K} \tierP); the
Hurwitz arguments from the points $a/5$; the signs from $\chiv$.  In this sense the
odd values are \emph{structurally determined} by the pentagonal arithmetic, and the
determination is a theorem with the analytic objects of Mathlib.  What the theorem
is not is an evaluation: it expresses $\zeta(2k+1)$ through $\zeta_K(2k+1)$ and the
Hurwitz values at $a/5$, and says nothing about their arithmetic nature: that $\zeta(3)$ is irrational~\cite{apery79}, that infinitely many odd values are~\cite{ballrivoal01}, and that one of $\zeta(5),\zeta(7),\zeta(9),\zeta(11)$ is~\cite{zudilin01} are proved by other means, and the determination neither uses nor reproves them.  The even values, by contrast, are explicit,
$\zeta(2k)\in\Q\,\pi^{2k}$ (\LeanStar{F1CW.even\_zeta\_bernoulli} \tierPnew), and so are
the even values of $L(s,\chiv)$, which lie in $\sqrt5\,\pi^{2k}\Q$ with rational
coefficients given by Bernoulli polynomials at the points $a/5$: the rationality of the
Bernoulli values is proved (\LeanStar{F1CW.bernoulli\_eval\_rational} \tierPnew), and the
formula expressing $L(2k,\chiv)$ through them is classical \tierC\ \cite[Ch.~4]{washington97}:
\[
  L(2k,\chiv)=\frac{(-1)^{k+1}\sqrt5\,(2\pi)^{2k}\,B_{2k,\chiv}}{2\,(2k)!\,5^{2k}},\qquad
  B_{n,\chiv}=5^{n-1}\sum_{a=1}^{4}\chiv(a)\,B_n\bigl(\tfrac a5\bigr),
\]
with $B_{2,\chiv}=\tfrac45$, $B_{4,\chiv}=-8$, $B_{6,\chiv}=\tfrac{804}5$, $B_{8,\chiv}=-5776$ in exact rational
arithmetic, and the formula agreeing with the Hurwitz sums at the points $a/5$ to $40$ digits
(\texttt{oz\_carried.py} \tierN); % pending(D7)
at $s=1$, $L(1,\chiv)=2\log\golden/\sqrt5$ (\cref{sec:mersenne}).

\begin{remark}[What is known of the arithmetic nature of $L(s,\chiv)$]\label{rmk:nature-L}
At the integers, $L(s,\chiv)$ is rational for $s\le0$, through the generalised Bernoulli
numbers \tierC; transcendental at $s=1$, since $L(1,\chiv)=2\log\golden/\sqrt5$ and
$\log\golden$ is transcendental --- were it algebraic and nonzero, $\golden=e^{\log\golden}$
would be transcendental by Lindemann's theorem~\cite{lindemann1882} \tierC; and
transcendental at $s=2k$, where it is a nonzero element of $\sqrt5\,\pi^{2k}\Q$ \tierC, by Lindemann's
theorem for $\pi$; the step from that form to transcendence, for $L(2k,\chiv)$ and for
$\sqrt5\,L(2k,\chiv)=\sum_{a=1}^{4}\chiv(a)\,\mathrm{Li}_{2k}(\zeta_5^{\,a})$ alike, is
\LeanInline{F1OZ.even\_value\_transcendental} \tierP, whose inputs are the transcendence of $\pi$
\cite{lindemann1882} \tierC\ and $L(2k,\chiv)\ne0$ (Mathlib's
\texttt{DirichletCharacter.LFunction\_ne\_zero\_of\_one\_le\_re}).  At
the odd $s=2k+1\ge3$ it is open \tierO.  For $s\ge2$ three statements are equivalent:
$L(s,\chiv)\in\overline{\Q}$; $\sum_{a=1}^{4}\chiv(a)\,\mathrm{Li}_s(\zeta_5^{\,a})\in\sqrt5\,\overline{\Q}$;
and $\sum_{a=1}^{4}\chiv(a)\,\zeta(s,a/5)\in5^{s}\,\overline{\Q}$.  Indeed, for $s\ge1$,
\[
  L(s,\chiv)=\frac1{\sqrt5}\sum_{a=1}^{4}\chiv(a)\,\mathrm{Li}_s(\zeta_5^{\,a})
  =\frac2{\sqrt5}\Bigl(\Re\,\mathrm{Li}_s(\zeta_5)-\Re\,\mathrm{Li}_s(\zeta_5^{2})\Bigr),
  \qquad \mathrm{Li}_1(z)=-\log(1-z),
\]
because $\sum_{a=1}^{4}\chiv(a)\zeta_5^{\,an}=2\bigl(\cos\tfrac{2\pi n}5-\cos\tfrac{4\pi n}5\bigr)=\sqrt5\,\chiv(n)$
(\LeanInline{F1OZ.fourier\_projection} \tierP; the two forms agree with the Hurwitz sums at $s=1,3,5$ to
$40$ digits, \texttt{oz\_carried.py} \tierN); and $L(s,\chiv)$ is the Hurwitz sum divided by $5^{s}$.  Both % pending(D7)
factors, $\sqrt5$ and $5^{s}$, are algebraic and nonzero (\LeanInline{PasosDelTexto.algebraic\_scale} \tierP),
which gives the equivalence; so the determination of \cref{thm:odd-zeta} does not decide whether
$\zeta(2k+1)$ is irrational.  $\mathrm{Li}_s(\zeta_5^{\,a})=\sum_n\zeta_5^{\,an}n^{-s}$ is the discrete Fourier
transform of $n^{-s}$ over $\Z/5$, and $L(s,\chiv)$ is its $\chiv$-projection.  At $s=1$ the form gives a second
proof of the transcendence of $L(1,\chiv)$, as a logarithm of distances on the circle: $2\sin(\pi a/5)=|1-\zeta_5^{\,a}|$, with
$|1-\zeta_5|^2=(5-\sqrt5)/2$, $|1-\zeta_5^2|^2=(5+\sqrt5)/2$ and ratio $\golden^2$
(\LeanInline{F1OZ.normSq\_one\_sub\_zeta5}, \LeanInline{F1OZ.normSq\_one\_sub\_zeta5\_sq},
\LeanInline{F1OZ.baker\_ratio} \tierP), so $\sqrt5\,L(1,\chiv)=\log\bigl(|1-\zeta_5^2|^2/|1-\zeta_5|^2\bigr)
=2\log\golden$, a nonzero linear form in logarithms of algebraic numbers, transcendental by Baker's theorem
\tierC\ \cite{baker66}.  If $L(2k+1,\chiv)$ is transcendental, then $\zeta(2k+1)$ and $\zeta_K(2k+1)$ are not
both algebraic, by \eqref{eq:dedekind} and \LeanInline{PasosDelTexto.algebraic\_scale} \tierP.  \tier{R} The linear form is the same for every $s$ --- the points $\zeta_5^{\,a}$ and the weights
$\chiv(a)/\sqrt5$ --- and only the kernel $\mathrm{Li}_s$ changes; at $s=1$ it reduces to logarithms and at
even $s$ to Bernoulli polynomials, and at odd $s\ge3$ to neither.  The predecessor version
\cite{gonzalez26oddzeta} read this as the reason no cancellation is expected there; it is not an argument,
and the question is the open one above.
\end{remark}

\begin{remark}[The period $T_K$ and the odd values, through $G$]\label{rmk:TK-odd}
The predecessor version \cite{gonzalez26oddzeta} asked whether $T_K=2\pi\log\golden$ appears in
$\zeta_K(2k+1)$ through the Kronecker limit formula or Siegel units.  The classes modulo $20$ say where each
factor of $T_K$ lives.  The class $19\equiv-1$ lies in the fibre $(1,0)$ of $G$
(\LeanInline{degenerate\_orbits\_are\_transversals} \tierP), so a character of $\Z_2\times\Z_2$ is even
exactly when it is trivial on the first component: of the four characters of \cref{thm:dedekind-natural},
$1$ and $\chiv$ are even, and $\chi_4$ and $\chi_4\chiv$ are odd.  At $s=1$ the even character $\chiv$ carries
the regulator, $L(1,\chiv)=2\log\golden/\sqrt5$ (\cref{thm:L1}), and the odd character $\chi_4\chiv$ of
$\Q(\sqrt{-5})$ carries the circle, $L(1,\chi_4\chiv)=\pi/\sqrt5$, by the class-number formula with $h=2$,
$w=2$ and discriminant $-20$ \tierC.  So the period is the product over the $\chiv$-fibre of $G$,
\[
  T_K=2\pi\log\golden=5\,L(1,\chiv)\,L(1,\chi_4\chiv)
\]
(\LeanInline{F1OZ.period\_from\_the\_chi5\_fibre} \tierP\ given the two values).  At $s=2k+1$ the parity
decides.  An odd character has its odd values in $\Q\,\pi^{2k+1}$ times the square root of its conductor
\tierC\ \cite[Ch.~4]{washington97}, as $L(3,\chi_4)=\pi^3/32$ and $L(3,\chi_4\chiv)=3\pi^3/(40\sqrt5)$
(\texttt{oz\_carried.py} \tierN), and with $\pi=5\arccos(\golden/2)$ and $\sqrt5=2\golden-1$ % pending(D7)
(\eqref{eq:pentagon-id}) these are written in $\golden$; while $\zeta_K(2k+1)=\zeta(2k+1)\,L(2k+1,\chiv)$ is the
product over the two even characters, the part without such a form.  The circle factor of $T_K$ continues at
odd $s$ on the odd side of $G$, outside $\zeta_K$:
$\zeta_{\Q(i,\sqrt5)}(2k+1)=\zeta_K(2k+1)\,L(2k+1,\chi_4)\,L(2k+1,\chi_4\chiv)$, the last two factors in
$\Q^\times\pi^{4k+2}/\sqrt5$ \tierC.  The regulator factor $\log\golden$ is the value of the even side at
$s=1$, and at $s=2k+1$ its place is taken by $L(2k+1,\chiv)$ itself.  The Kronecker limit formula for real
quadratic fields gives the constant term at $s=1$ \cite{zagier75} \tierC, which is where $T_K$ already
lives.  At $s=2k+1$ Borel's theorem relates $\zeta_K(2k+1)$, up to a nonzero rational factor, a power of $\pi$
and $\sqrt5$, to the regulator of $K_{4k+1}(\OK)$, of rank $r_1+r_2=2$, and not to the regulator $\log\golden$
of $K_1(\OK)=\OK^\times$ \cite{borel77} \tierC; the rank for $\Q(i,\sqrt5)$ is also $r_1+r_2=2$, so its two
odd characters add none, as their closed forms require.  The exact expression is \eqref{eq:dedekind} with the
polylogarithmic form above,
\[
  \zeta_K(2k+1)=\frac1{\sqrt5}\,\mathrm{Li}_{2k+1}(1)\sum_{a=1}^{4}\chiv(a)\,\mathrm{Li}_{2k+1}(\zeta_5^{\,a}),
\]
and $T_K$ does not appear in it.  Whether $\zeta_K(2k+1)$ satisfies another algebraic relation with
$\log\golden$ is a question of the algebraic independence of periods, and stays open \tierO.
\end{remark}

\begin{remark}[$G$ and the motivic Galois group]\label{rmk:G-motivic}
The predecessor version also asked whether $G$ is a finite-level reflection of the motivic Galois group
acting on the odd values.  It is not, in the following sense.  The odd values $\zeta(2k+1)$ are periods of
mixed Tate motives over $\Z$, and the Galois group of that Tannakian category is $\mathbb G_m\ltimes U$, with
$U$ pro-unipotent and its Lie algebra free on one generator $\sigma_{2k+1}$ in each odd degree $2k+1\ge3$
\cite{delignegoncharov05,brown12} \tierC.  That group is connected, and a connected group scheme has no
nontrivial homomorphism onto a finite group, since the image would be connected and finite; so no
nontrivial finite group, $G$ among them, is a quotient of it, and the generators $\sigma_{2k+1}$ that carry
the odd values lie where no finite quotient reaches \tierT.  $G$ is instead the Galois group of the
biquadratic field $\Q(i,\sqrt5)$, the field cut out by its kernel, the squares (\cref{prop:G}) \tierC: it acts
on the odd values through the factorisation, by deciding their local factors, and not through the motivic
action.
\end{remark}

\subsubsection{The data of the determination}
\label{ssec:pentagon-odd}

This subsection proves the Hurwitz half of \cref{thm:odd-zeta} and identifies what
generates each piece of its data.  The factorisation itself is the natural
isomorphism of \cref{thm:dedekind-natural}; what is added here is that both sides of
it, at $s=2k+1$, are $\chi_5$-projections of the Hurwitz zeta function at the
pentagonal points $\{a/5\}_{a=1}^{4}$, with every component fixed by the pentagonal
arithmetic.

\begin{theorem}[Hurwitz representation of $L(s,\chi_5)$]
\label{thm:hurwitz}
For all $s$ with $\mathrm{Re}(s)>1$, and by analytic continuation
for $s\ne 1$:
\begin{equation}\label{eq:hurwitz-rep}
  L(s,\chi_5)
    = \frac{1}{5^s}\sum_{a=1}^{4}\chi_5(a)\,\zeta\!\left(s,\tfrac{a}{5}\right),
\end{equation}
where $\zeta(s,x) = \sum_{n=0}^{\infty}(n+x)^{-s}$ is the Hurwitz
zeta function.  The four arguments $a/5$ for $a\in\{1,2,3,4\}$ are
the \emph{pentagonal points}, and $\chi_5$ is their Artin weight.

The kernels at the three classes of special values are:
\begin{align}
  s = 1\colon\quad
    & \zeta(s,a/5)-\tfrac1{s-1}\xrightarrow{s\to 1} -\psi(a/5)
    \;\Rightarrow\; L(1,\chi_5) = -\tfrac15\textstyle\sum_a\chi_5(a)\psi(\tfrac a5)=\tfrac{2\log\golden}{\sqrt{5}},
    \label{eq:kernel-s1}\\
  s = 2k\colon\quad
    & \zeta(2k,a/5)\text{ via Bernoulli polynomials}
    \;\Rightarrow\; L(2k,\chi_5) \in \sqrt{5}\cdot\pi^{2k}\cdot\mathbb{Q},
    \label{eq:kernel-even}\\
  s = 2k+1,\;k\geq 1\colon\quad
    & \zeta(2k+1,a/5) = \sum_{n=0}^{\infty}(n+a/5)^{-(2k+1)}.
    \label{eq:kernel-odd}
\end{align}
In all three cases, the organising structure is identical: the same
$\chi_5(a)$-weights over the same pentagonal points $\{a/5\}_{a=1}^{4}$.
\end{theorem}
\leanline{\LeanInline{hurwitz\_representation\_L\_chi5} \tierP\ for \eqref{eq:hurwitz-rep} with
$\Re s>1$; its continuation and the kernels \eqref{eq:kernel-s1}--\eqref{eq:kernel-even}
are classical \tierC.}

\begin{proof}
Equation~\cref{eq:hurwitz-rep} is well-established \cite[Ch.~12]{apostol76}: any periodic Dirichlet
character $\chi$ modulo $N$ admits the Hurwitz expansion
$L(s,\chi) = N^{-s}\sum_{a=1}^{N-1}\chi(a)\zeta(s,a/N)$.  For
$\chi = \chi_5$, the sum reduces to $a\in\{1,2,3,4\}$ since
$\chi_5(0) = 0$.

The kernel at $s=1$ follows from the expansion
$\zeta(s,x) = 1/(s-1) - \psi(x) + O(s-1)$ near $s=1$, with $\psi$ the
digamma function, combined with $\sum_{a=1}^{4}\chi_5(a) = 0$ (so the $1/(s-1)$ poles
cancel, and so does the term $\log 5$ from $5^{-s}$).  This yields
$L(1,\chi_5) = -\sum_a\chi_5(a)\psi(a/5)/5$, which by Gauss's digamma theorem \tierC\ produces
$2\log\golden/\sqrt{5}$, matching \cref{thm:L1}; equivalently
$\sum_{a=1}^{4}\chi_5(a)\psi(a/5)=-2\sqrt5\log\golden$ (\LeanInline{F1OZ.digamma\_sum} \tierP\ from the
two values of $L(1,\chiv)$).

The kernel at even $s=2k$ uses
$\zeta(-n,x) = -B_{n+1}(x)/(n+1)$ and the functional equation,
giving $L(2k,\chi_5) = \sqrt{5}\cdot\pi^{2k}\cdot r_k$ for some
$r_k\in\mathbb{Q}$.

The kernel at odd $s=2k+1$ is the direct Hurwitz series, convergent
for $\mathrm{Re}(s)>1$.

In the Lean formalisation, this representation is
\lean{hurwitz\_representation\_L\_chi5}; the proof is definitional in
Mathlib's \texttt{DirichletCharacter.LFunction}, which is constructed
precisely as the Hurwitz combination weighted by the character (proved by explicit
enumeration).
\end{proof}

\begin{proposition}[The Riemann zeta function at the pentagonal points]\label{prop:zeta-hurwitz}
For $\Re s>1$,
\[
  \zeta(s)=\frac1{5^s}\sum_{a=1}^{5}\zeta\!\left(s,\tfrac a5\right),\qquad
  \zeta_{\Q(\sqrt5)}(s)=\frac1{5^{2s}}\Bigl(\sum_{a=1}^{5}\zeta\!\left(s,\tfrac a5\right)\Bigr)
  \Bigl(\sum_{b=1}^{4}\chiv(b)\,\zeta\!\left(s,\tfrac b5\right)\Bigr),
\]
the term $a=5$ being $\zeta(s,1)=\zeta(s)$.  So both factors of \eqref{eq:dedekind} are sums of the Hurwitz zeta
function at the points $a/5$: $\zeta$ with weight $1$ over $a=1,\dots,5$, and $L(s,\chiv)$ with weight $\chiv(a)$ over
$a=1,\dots,4$.
\end{proposition}
\leanline{\LeanInline{ZetaPentagonal.riemannZeta\_pentagonal\_points}, \LeanInline{hurwitz\_representation\_L\_chi5},
\LeanInline{dedekind\_factorisation} \tierP.}
\begin{proof}
For $\Re s>1$ the series $\sum n^{-s}$ converges absolutely, so it can be summed by residues modulo $5$: writing
$n=5m+a$ with $1\le a\le5$ and $m\ge0$, $n^{-s}=5^{-s}(m+a/5)^{-s}$, and the sum over $m$ is $\zeta(s,a/5)$.  The
second identity is the first times \cref{thm:hurwitz}, by \eqref{eq:dedekind}.
\end{proof}

\begin{corollary}[What generates the data of \cref{thm:odd-zeta}]
\label{thm:pentagon-odd}
For $k\geq 1$, the factorisation of \cref{thm:odd-zeta}
\begin{equation}\label{eq:zeta-odd-factorisation}
  \zeta(2k+1) = \frac{\zeta_{\mathbb{Q}(\sqrt{5})}(2k+1)}{L(2k+1,\chi_5)}
\end{equation}
expresses $\zeta(2k+1)$ through the Hurwitz zeta function at the points $a/5$: $L(2k+1,\chi_5)$ is its
sum weighted by $\chi_5(a)$ over $a=1,\dots,4$, and $\zeta_{\mathbb{Q}(\sqrt{5})}(2k+1)$ is the product of that sum
with the unweighted sum over $a=1,\dots,5$ (\cref{prop:zeta-hurwitz}).  All
the data are generated by:
\begin{enumerate}[label=\textup{(\arabic*)},nosep]
\item $2\cos(\pi/5) = \golden$ \emph{(pentagonal geometry generates $\golden$)};
\item $\sum_{a=1}^{4}\chi_5(a)\log(2\sin(\pi a/5)) = -2\log\golden$
  \emph{(the logarithmic signature of $\golden$ in the pentagon)};
\item the ring Frobenius of $\Zphi$ determines $\chi_5(p)\equiv F_p\pmod p$, the Artin
  character of $\mathbb{Q}(\sqrt{5})/\mathbb{Q}$
  \emph{(Frobenius lifts connect $\golden$ to the L-function character)}.
\end{enumerate}
The case $k=0$ yields the closed form
$L(1,\chi_5) = 2\log\golden/\sqrt{5}$.  For $k\geq 1$, the
determination is structural: $L(2k+1,\chi_5)$ and
$\zeta_{\mathbb{Q}(\sqrt{5})}(2k+1)$ are built from the same values $\zeta(2k+1,a/5)$, the first with the
weights $\chi_5(a)$ and the second as the product of the weighted and the unweighted sums.  No
parameter beyond $\golden$, $\pi$, and the pentagonal arithmetic of
$\mathbb{Q}(\sqrt{5})$ is required.
\tierP/\tierPnew
\end{corollary}

\begin{proof}
Items~(1)--(3) are \cref{eq:pentagon-id}, the logarithmic signature
(\LeanStar{F1v2.pentagon\_log\_signature\_full}), and \cref{thm:chi5-from-frobenius}
respectively.  The factorisation is
\cref{thm:dedekind-natural}.  The Hurwitz representation is
\cref{thm:hurwitz}.  The $k=0$ closed form is
\cref{thm:L1}.

For $k\geq 1$: by \cref{thm:hurwitz}, both
$L(2k+1,\chi_5) = 5^{-(2k+1)}\sum_{a=1}^{4}\chi_5(a)\zeta(2k+1,a/5)$
is the Hurwitz series at the pentagonal points weighted by~$\chi_5$, and
$\zeta_{\mathbb{Q}(\sqrt{5})}(2k+1)$ is the product of the local factors fixed by
$\chi_5(p)$ (\cref{thm:odd-zeta}), and the product of the two Hurwitz sums of \cref{prop:zeta-hurwitz}.  Whether
either factor has a closed form in $\pi$ is open (cf.\ Ap\'ery's theorem and its generalisations), and
their ratio $\zeta(2k+1)$ is fixed by the arithmetic of $\Rpcf$
without any free parameter beyond the pentagonal data.

The complete cascade is formalised in Lean.  The Euler products are
\lean{euler\_product\_identity} and \lean{dirichlet\_series\_identity}---direct
invocations of \texttt{riemannZeta\_eulerProduct\_tprod} and
\texttt{DirichletCharacter.LSeries\_eulerProduct\_tprod} from
Mathlib's analytic L-series infrastructure.  The local Euler factor
classification (split, inert, ramified) is \lean{f\_K}, with the
multiplicative identity $f_K = f_\zeta \cdot f_L$ formalised as
\lean{local\_factor\_identity} (pure algebra, fully verified at all
three cases).  The structural and analytic presentations of $\chi_5$
are identified by \lean{chi5\_structural\_eq\_analytic} (the bridge
theorem), and the Dedekind factorisation
$\zeta_{\mathbb{Q}(\sqrt{5})}(s) = \zeta(s)\cdot L(s,\chi_5)$ is
\lean{dedekind\_factorisation}.  The final pentagonal determination
is \lean{zeta\_odd\_pentagonal\_determination}.
\end{proof}

\begin{table}[ht]
\centering
\caption{Zeta values from $s = 2$ to $s = 13$.  Even values (shaded)
are rational multiples of $\pi^{2k}$; odd values are ratios
$\zeta_{\mathbb{Q}(\sqrt{5})}(s)/L(s,\chi_5)$, both factors determined by
$\chi_5$ and hence by $\golden$.  The structural identity at every $s$
is the Dedekind factorisation \lean{dedekind\_factorisation}; numerical
entries are verified at 50-digit precision (see companion verification
script \texttt{verify\_odd\_zeta\_pcf\_unified.py}).}
\label{tab:zeta-values}
\begin{tabular}{rcccc}
\toprule
$s$ & $\zeta_{\mathbb{Q}(\sqrt{5})}(s)$ & $L(s,\chi_5)$ & $\zeta(s)$ & Source \\
\midrule
\rowcolor[gray]{0.93}
 2 & --- & --- & $\pi^2/6 \approx 1.6449$ & $\pi$ \\
 3 & $1.02755$ & $0.85482$ & $1.20206$ & $\golden,\pi$ \\
\rowcolor[gray]{0.93}
 4 & --- & --- & $\pi^4/90 \approx 1.0823$ & $\pi$ \\
 5 & $1.00133$ & $0.96567$ & $1.03693$ & $\golden,\pi$ \\
\rowcolor[gray]{0.93}
 6 & --- & --- & $\pi^6/945 \approx 1.0173$ & $\pi$ \\
 7 & $1.00007$ & $0.99179$ & $1.00835$ & $\golden,\pi$ \\
\rowcolor[gray]{0.93}
 8 & --- & --- & $\pi^8/9450 \approx 1.0041$ & $\pi$ \\
 9 & $1.000004$ & $0.99800$ & $1.00201$ & $\golden,\pi$ \\
\rowcolor[gray]{0.93}
10 & --- & --- & $\pi^{10}/93555 \approx 1.00099$ & $\pi$ \\
11 & $1.0000003$ & $0.99951$ & $1.00049$ & $\golden,\pi$ \\
\rowcolor[gray]{0.93}
12 & --- & --- & $691\pi^{12}/638512875 \approx 1.00025$ & $\pi$ \\
13 & $1.00000002$ & $0.99988$ & $1.00012$ & $\golden,\pi$ \\
\bottomrule
\end{tabular}
\end{table}

\begin{example}[From the pentagon to $\zeta(3)$]\label{ex:zeta3}
At $s=3$ the local factors of $\zeta_K$ are fixed by $\chiv(p)$, which is read off $F_p\bmod p$
(\cref{thm:s1}): at the inert primes $2$ and $3$, $f_K(2,3)=(1-2^{-6})^{-1}=\tfrac{64}{63}$ and
$f_K(3,3)=\tfrac{729}{728}$; at the ramified prime, $f_K(5,3)=(1-5^{-3})^{-1}=\tfrac{125}{124}$; at the split
prime $11$, $f_K(11,3)=\bigl(\tfrac{1331}{1330}\bigr)^2$ (\LeanInline{F1OZ.local\_factors\_at\_three}
\tierP).  No parameter is chosen at any prime.  The two products give $\zeta_K(3)=1.027548\ldots$ and
$L(3,\chiv)=0.85482\ldots$, and their ratio is $\zeta(3)=1.2020569\ldots$ (\cref{tab:zeta-values};
\texttt{oz\_carried.py} \tierN). % pending(D7)
\end{example}

\begin{remark}[Structural versus closed-form determination]
\label{rmk:structural-vs-closed}
\Cref{thm:pentagon-odd} establishes \emph{structural} determination:
the pentagonal lattice $\{a/5\}_{a=1}^{4}$, the Artin character
$\chi_5$, and the Hurwitz kernel $\zeta(2k+1,\cdot)$ are all fixed
by $\golden$ and $\pi$ without additional parameters.  A closed-form
expression $\zeta(2k+1) = r_k\pi^{2k+1}$ for $r_k\in\mathbb{Q}$
would be stronger; whether it exists for $k\geq 1$ is an open
problem related to Ap\'ery's theorem and its generalisations.  What
the present construction establishes is that the \emph{source}---the
organising arithmetic structure---is the pentagon generated by
$\golden$ and $\pi$, not free arithmetic data.
\end{remark}

\begin{remark}[The pentagon as primitive structure]
\label{rmk:spectral-validation}
\Cref{thm:pentagon-odd} is the sense in which the pentagon is the primitive structure
beneath the analytic apparatus: $\zeta(2k+1)$ is not added to the
construction but emerges from~$\golden$ and~$\pi$, the data already
present in the pentagonal lattice $\{a/5\}_{a=1}^{4}$ and the
$\Lambda$-monoid generator.
\end{remark}








\subsubsection{What the chain proves, and over what}

By \cref{thm:chi5-from-frobenius} the splitting type at $p$ is read off the
Frobenius of \cref{sec:lambda-ring}, so the chain
\[
  \Frob_p\ \Longrightarrow\ \chiv(p)\ \Longrightarrow\ \text{local factor}\
  \Longrightarrow\ \eqref{eq:dedekind}
\]
consists of implications, each proved.  Its scope has to be stated with the same
precision: the structure on $\Zphi$ controls $L(s,\chiv)$, and it controls
$\zeta(s)$ only through the product \eqref{eq:dedekind}, where the two factors
appear together.  Every statement of \cref{sec:topological-proof} about the position of zeros therefore
concerns the zeros of $L(s,\chiv)$, and the zero set of $\zeta_K$ is the union of those of $\zeta$ and of
$L(s,\chiv)$; what reaches $\zeta$ is the arithmetic: the Euler factors, the factorisation and the odd values.
The odd values are treated in \cref{ssec:odd-zeta}.  Since a Dirichlet series is determined on the whole
plane by its coefficients, through analytic continuation \tierC\ \cite{riemann1859}, the chain determines the
function $\zeta(s)$ everywhere; that determination is of the function, and by itself it places no zero.

% ---------------------------------------------------------------------
\subsection{The counting function: from Soul\'e to Connes--Consani}
\label{sec:counting}

The construction of \cref{sec:colimit} produced
$\zeta_{\mathbb{Q}(\sqrt{5})}(s)$ as a colimit of partial Euler
products indexed by primes, with each local factor $f_p(s)$ read
from the Artin character $\chi_5(p)$ via the Fibonacci splitting
criterion.  This output is structurally an Euler product; for it
to be legible as a zeta function in the $\mathbb{F}_1$ sense, it
must admit a counting function in one of the established
frameworks.  In this section we show that
$\zeta_{\mathbb{Q}(\sqrt{5})}(s)$ does not satisfy the polynomial
counting condition of Soul\'e~\cite{soule03} but does admit a
distributional counting function in the sense of Connes and
Consani~\cite{connesconsani10}, and we set up that counting
function explicitly.

\subsubsection{Why $\zeta_{\mathbb{Q}(\sqrt{5})}(s)$ does not satisfy
Soul\'e's condition}\label{ssec:not-soule}

\begin{definition}[Soul\'e's condition]\label{def:soule-condition}
A scheme $X$ of finite type over~$\mathbb{Z}$ is said to satisfy
\emph{Soul\'e's condition} if there exists a polynomial
$N_X\in\mathbb{Z}[x]$ such that
\begin{equation}\label{eq:soule-condition}
  \#X(\mathbb{F}_{p^n}) = N_X(p^n)
  \qquad\text{for every prime } p\text{ and every } n\ge 1.
\end{equation}
When this condition holds, the Hasse--Weil generating function
$Z(X,T) = \exp\!\left(\sum_{r\ge 1} N_X(q^r)T^r/r\right)$ is a
rational function of~$T$, and the limit
\begin{equation}\label{eq:soule-zeta}
  \zeta_X(s) \;=\; \lim_{q\to 1}\,Z(X,q^{-s})\,(q-1)^{N_X(1)}
\end{equation}
is well defined and gives the $\mathbb{F}_1$-zeta function of~$X$~\cite[\S6]{soule03}, in the form corrected in \cite[Eq.~(3)]{connesconsani10}.
\end{definition}

\begin{proposition}[Failure of Soul\'e's condition for $\zeta_K$]
\label{prop:not-soule}
The number of prime ideals of residue degree one of $\mathcal O_K$ above a prime
$p\ne5$ is $1+\chi_5(p)$, and no polynomial $N\in\Z[x]$ satisfies $N(p)=1+\chi_5(p)$ for
every prime $p\ne5$.  So the degree-one count of $\mathcal O_K$ is not a Soul\'e counting
polynomial, and $\zeta_K$ is not the zeta function of a scheme satisfying
\eqref{eq:soule-condition} through that count.
\end{proposition}
\leanline{\LeanStar{F1v2.not\_soule}, \LeanStar{F1v2.degree\_one\_count} \tierPnew.}

\begin{proof}
A split prime has two primes of degree one above it and an inert prime none, so the
count is $1+\chi_5(p)\in\{0,2\}$.  If $N(p)=1+\chi_5(p)$ for every prime $p\ne5$, then
$N(N-2)$ vanishes at infinitely many integers and is the zero polynomial, so $N=0$ or
$N=2$.  But $N(11)=1+\chi_5(11)=2$ excludes $N=0$, and $N(2)=1+\chi_5(2)=0$ excludes
$N=2$.
\end{proof}

\begin{remark}[Trivial vs.\ non-trivial cases]\label{rmk:trivial-cases}
For $K=\mathbb{Q}$ every prime has exactly one prime above it, so the
degree-one count is the constant polynomial $1$; Spec\,$\mathbb Z$ still fails Soul\'e's
condition as a whole (\cref{ssec:soule}), and $\zeta$ is reached only by the
distributional extension of Connes--Consani.  What \cref{prop:not-soule} adds is that
for $\mathbb{Q}(\sqrt5)$ already the degree-one count is not polynomial: the character
$\chi_5$ is what breaks it.
\end{remark}

\subsubsection{The integral formula and distributional extension}
\label{ssec:integral-formula}

Connes and Consani~\cite{connesconsani10} replace the limit
$q\to 1$ in~\eqref{eq:soule-zeta} by an integral against a
counting distribution.  We recall their construction in the form
needed below.

\begin{lemma}[Integral formula~\cite{connesconsani10}, Lemma~2.1]
\label{lem:cc-integral}
Let $N\colon[1,\infty)\to\mathbb{R}$ be continuous with polynomial
bound $|N(u)|\le Cu^k$, and write $d^*u=du/u$.  Define
$\zeta_N(s) := \exp\!\left(\int_1^\infty
\frac{N(u)}{\log u}\,u^{-s}\,d^*u\right)$, up to a multiplicative constant, the integral being taken as a principal
value at $u=1$ when $N(1)\neq0$ \cite[Eqs.~(9)--(10)]{connesconsani10}.  Then
\begin{equation}\label{eq:cc-integral}
  \frac{\partial_s\zeta_N(s)}{\zeta_N(s)} \;=\;
  -\int_1^\infty N(u)\,u^{-s}\,d^*u
\end{equation}
for $\mathrm{Re}(s)>k$.  When $N\in\mathbb{Z}[x]$
the polynomial counting case is recovered, with
$\zeta_N(s) = \lim_{q\to 1}Z(q,q^{-s})(q-1)^{N(1)}$.
\tierC
\end{lemma}

\Cref{lem:cc-integral} converts the limit definition into an
integral; the integral makes sense for any $N$ that is a
distribution on $[1,\infty)$ with polynomial bound, not only for
polynomial~$N$.  This is the distributional extension.

\begin{theorem}[Connes--Consani for $\zeta_\mathbb{Q}$~%
\cite{connesconsani10}, Theorem~2.2]\label{thm:cc-zeta-Q}
The complete Riemann zeta
$\zeta_\mathbb{Q}^*(s) = \pi^{-s/2}\Gamma(s/2)\zeta(s)$ satisfies
\eqref{eq:cc-integral} with
\begin{equation}\label{eq:N-Q}
  N_\mathbb{Q}(u) \;=\; u
  \;-\; \frac{d}{du}\!\left[\sum_{\rho\in Z}
  \mathrm{ord}(\rho)\,\frac{u^{\rho+1}}{\rho+1}\right]
  \;+\; 1,
\end{equation}
where $Z$ is the set of non-trivial zeros of $\zeta(s)$ counted
with multiplicity.  The distribution $N_\mathbb{Q}$ is positive on
$(1,\infty)$ but $N_\mathbb{Q}(1) = -\infty$, recovering the
predicted infinite genus of $\overline{\mathrm{Spec}\,\mathbb{Z}}$
\textup{(}Manin~\cite{manin95}\textup{)}.
\tierC
\end{theorem}

The structural content of~\eqref{eq:N-Q} is that the
$\mathbb{F}_1$-counting function over
$\overline{\mathrm{Spec}\,\mathbb{Z}}$ has the same form
$u - \sum_\alpha\alpha^{\bullet} + 1$ as the Lefschetz-type formula
for curves over $\mathbb{F}_q$, with the non-trivial zeros~$\rho$
playing the role of Frobenius eigenvalues.

\subsubsection{The PCF counting function $N_K(u)$}
\label{ssec:pcf-counting}

We now apply \cref{lem:cc-integral} to $K=\mathbb{Q}(\sqrt{5})$ and
identify the distributional counting function attached to
$\zeta_K(s)$.  Throughout, $\Lambda_K$ denotes the von Mangoldt
function of~$K$, defined on prime ideals by
$\Lambda_K(\mathfrak{p}^k) = \log N(\mathfrak{p})$ and zero on
non-prime-power ideals; $N(\mathfrak{p})$ denotes the absolute
norm.

\begin{definition}[Chebyshev function of~$K$]\label{def:chebyshev-K}
For $u>1$ define
\begin{equation}\label{eq:chebyshev-K}
  \psi_K(u) \;=\; \sum_{\substack{\mathfrak{a}\;\text{integral ideal}\\ N(\mathfrak{a})\le u}}
  \Lambda_K(\mathfrak{a}).
\end{equation}
Equivalently, $\psi_K(u) = \sum_{\mathfrak{p}^k}\log N(\mathfrak{p})$
where the sum runs over prime-power ideals of norm at most~$u$.
\end{definition}

\begin{proposition}[Logarithmic derivative
of~$\zeta_K$]\label{prop:dlog-zetaK}
For $\mathrm{Re}(s)>1$,
\begin{equation}\label{eq:dlog-zetaK}
  -\frac{\partial_s\zeta_K(s)}{\zeta_K(s)}
  \;=\; \sum_{\mathfrak{a}}\Lambda_K(\mathfrak{a})\,N(\mathfrak{a})^{-s}
  \;=\; \int_1^\infty u^{-s}\,d\psi_K(u).
\end{equation}
\end{proposition}
\leanline{\LeanStar{F1v2.local\_log\_deriv}, \LeanStar{F1v2.local\_dirichlet\_expansion} \tierPnew; the values of $\Lambda_K$: \LeanStar{F1CW.von\_mangoldt\_k\_values} \tierP.}

\begin{proof}
The first equality is the standard logarithmic derivative of the
Euler product $\zeta_K(s) = \prod_\mathfrak{p}(1-N(\mathfrak{p})^{-s})^{-1}$.
The second is integration by parts using
\cref{def:chebyshev-K}.
\end{proof}

\begin{theorem}[PCF counting function]\label{thm:pcf-counting}
The Dedekind zeta $\zeta_K(s)$ with $K=\mathbb{Q}(\sqrt{5})$
admits the distributional counting function
\begin{equation}\label{eq:N-K}
  N_K(u) \;=\; u\,\frac{d\psi_K}{du}(u)\;+\;\kappa_K(u)
  \;=\;\frac{d}{du}\Bigl(\sum_{N(\mathfrak a)<u}N(\mathfrak a)\,\Lambda_K(\mathfrak a)\Bigr)+\kappa_K(u),
\end{equation}
where $\kappa_K(u)=2u^2/(u^2-1)$, the density $u^2/(u^2-1)$ of a real place counted at the two real places
of~$K$ \cite[Appendix, `Arithmetic machinery']{gonzalez26cw}, is their contribution, regularised as a principal value at~$u=1$
(\LeanStar{F1CW.kappa\_k\_two\_real\_places}, \LeanStar{F1CW.kappa\_per\_real\_place}): $u\,d\psi_K/du$ counts $\zeta_K$, and $N_K$ counts the complete
function $\zeta_K^*(s)=\bigl(\pi^{-s/2}\Gamma(s/2)\bigr)^2\zeta_K(s)$, as $N_\Q$ counts $\zeta_\Q^*$ in
\cref{thm:cc-zeta-Q}.  Explicitly,
$N_K(u)$ is supported on the prime-power norms
$\{N(\mathfrak{p}^k)\}_{\mathfrak{p},k}$, with weights
\begin{equation}\label{eq:N-K-weights}
  N_K\big(N(\mathfrak{p}^k)\big)
  \;=\; N(\mathfrak{p}^k)\,\log N(\mathfrak{p})
  \;\text{ \textup{(}as a Dirac mass weight\textup{)}}.
\end{equation}
The support is parametrised by the Fibonacci splitting criterion
\textup{(\cref{thm:s1})}: if $\chi_5(p) = +1$, then $p$
splits in~$\mathcal{O}_K$ as $p = \mathfrak{p}\overline{\mathfrak{p}}$ and
$N(\mathfrak{p}^k) = p^k$ contributes a Dirac mass of weight
$p^k\log p$ at $u = p^k$ \textup{(}with multiplicity~$2$\textup{)};
if $\chi_5(p) = -1$, then $p$ is inert and $N(\mathfrak{p}^k) = p^{2k}$
contributes a Dirac mass of weight $p^{2k}\log p^2 = 2p^{2k}\log p$ at $u = p^{2k}$;
and the ramified prime $5$, with $N(\mathfrak p)=5$, contributes $5^k\log5$ at $u=5^k$.
The factor $N(\mathfrak p^k)$ is the normalisation of the Connes--Consani formula;
the Lean statements below concern the local values of $\Lambda_K$ and the number of
primes above $p$.
\end{theorem}
\leanline{\LeanStar{F1v2.pcf\_counting\_weights} \tierPnew\ for the weights at the prime-power norms, \LeanStar{F1v2.pcf\_counting\_local}, \LeanStar{F1v2.local\_log\_deriv} \tierPnew\ for the local assembly, the global identity being the Connes--Consani integral formula \tierC; its components are \tierP: \LeanStar{F1CW.von\_mangoldt\_k\_values} (the weights $\log N\mathfrak p$), \LeanStar{F1CW.kappa\_k\_two\_real\_places}, \LeanStar{F1CW.kappa\_k\_simple\_pole}, \LeanStar{F1CW.kappa\_per\_real\_place} (the archimedean density and its principal value).}

\begin{proof}
Substituting~\eqref{eq:dlog-zetaK} into
\eqref{eq:cc-integral}: the
right-hand side of~\eqref{eq:cc-integral} for $N_K$ equals
$-\int_1^\infty u^{-s}\,d\psi_K(u)$ provided
$N_K(u)\,du/u = d\psi_K(u)$, that is,
$N_K(u) = u\,d\psi_K/du$ as distributions on $(1,\infty)$; since
$u\,\delta_{u_0}=u_0\,\delta_{u_0}$, each Dirac mass of $d\psi_K$ at a norm $N(\mathfrak a)$
is multiplied by $N(\mathfrak a)$.  The Archimedean correction $\kappa_K(u)$
absorbs the contribution of the gamma factors at the two real
places of~$K$, regularised as a principal value at~$u=1$ in the manner of~\cite[Eq.~(12) and the definition of $\kappa$ that follows it]{connesconsani10}.
The support and weights follow from the explicit form of the
local Euler factors and the splitting criterion.
\end{proof}

\begin{remark}[The PCF counting function is structural,
not analytic]\label{rmk:counting-structural}
Theorem~\ref{thm:pcf-counting} produces $N_K(u)$ \emph{without}
reference to the zeros of~$\zeta_K(s)$.  Its support is the prime
power norms; its weights are the von Mangoldt values; both are
read off the splitting data $\chi_5(p)$, which is in turn read
off the Fibonacci congruence $\chi_5(p)\equiv F_p\pmod p$.  The
construction therefore determines the weights and the multiplicities from the pentagonal arithmetic
$\golden^2=\golden+1$ and the Frobenius lifts $\Frob_p$; the support, the prime powers themselves,
is an input the structure does not supply --- their sizes $\log p$ are what the law of
\cref{conj:closure} has to give (\cref{prop:what-fluctuation-takes,prop:euler-tower}).
\end{remark}

\begin{corollary}[What the counting function fixes, and the residue]\label{cor:residue}
\eqref{eq:N-K} fixes the logarithmic derivative of $\zeta_K$ on $\Re s>1$, so it determines $\zeta_K$ there up to a
constant factor, and the normalisation $\zeta_K(\sigma)\to1$ as $\sigma\to+\infty$ fixes the constant.  It therefore
determines the residue of $\zeta_K$ at $s=1$, whose value is given by the class-number formula:
\begin{equation}\label{eq:dedekind-residue}
  \mathrm{Res}_{s=1}\zeta_K(s) \;=\; \frac{2^{r_1}(2\pi)^{r_2}\,h_K\,R_K}{w_K\sqrt{|\Delta_K|}}
  \;=\;\frac{4\log\golden}{2\sqrt5}\;=\; \frac{2\log\golden}{\sqrt{5}},
\end{equation}
with $r_1=2$, $r_2=0$, $h_K=1$, $R_K=\log\golden$, $w_K=2$ and $|\Delta_K|=5$, the value of \cref{thm:L1}.
\end{corollary}
\leanline{\LeanInline{TodoEnLean.same\_logDeriv\_proportional}, \LeanInline{TodoEnLean.constant\_is\_one},
\LeanInline{TodoEnLean.zetaK\_tendsto\_one}, \LeanInline{anchor\_is\_a\_crossed\_identity},
\LeanInline{dedekind\_factorisation}, \LeanStar{F1CW.class\_number\_one\_K} \tierP\ given the class-number formula \tierC.}
\begin{proof}
Two functions holomorphic and without zeros on the half-plane $\Re s>1$ with the same logarithmic derivative have
a quotient of derivative zero, hence constant on that connected open set.  The Dirichlet series of $\zeta_K$ has
$a_1=1$, so $\zeta_K(\sigma)\to1$ as $\sigma\to+\infty$.  The residue is a property of $\zeta_K$, and its value is
\cref{thm:L1} with $\lim_{s\to1}(s-1)\zeta(s)=1$.
\end{proof}

\subsubsection{Comparison: descent data versus distributional counting}
\label{ssec:descent-vs-counting}

The construction yields two distinct $\mathbb{F}_1$-data attached
to the same object $\zeta_K$, with complementary roles.

\begin{remark}[Two readings of the same
zeta]\label{rmk:two-readings}
\textup{(a)} The $\Lambda$-monoid structure on
$\langle\golden\rangle\subset\Rpcf^\times$
\textup{(\cref{sec:lambda-monoid})} provides the
\emph{descent data} $\mathrm{Spec}\,\langle\golden\rangle\to\mathbb{F}_1$:
a system of commuting Frobenius lifts $\{\Frob_p\}_{p\in\mathbb{P}}$
satisfying Borger's two axioms.
This data is determined
\emph{before} any analytic identification: the lift
$\Frob_p(\golden^n) = \golden^{pn}$, equivalently
$\Frob_p(\golden) = F_p\golden + F_{p-1}$ on the generator, is read
off the pentagonal arithmetic alone, with no reference to zeros,
residues, or analytic continuation of zeta.

\textup{(b)} The distributional counting function $N_K(u)$
\textup{(\cref{thm:pcf-counting})} provides the
\emph{counting data} in the Connes--Consani sense: a distribution
on $[1,\infty)$ encoding the prime-power norms of~$K$ with
their von Mangoldt weights, and producing $\zeta_K(s)$ via
\eqref{eq:cc-integral}.  This data is the analytic output of the
arithmetic input.

The two readings are complementary: the descent data classifies
\emph{how} the arithmetic of~$K$ is encoded in commuting power
maps, while the counting data classifies \emph{what} arithmetic
information emerges as an analytic distribution on the norm
spectrum.  PCF supplies a structural source for the descent that
the Connes--Consani framework otherwise leaves analytic.
\end{remark}

\begin{remark}[Position relative to Manin's
question]\label{rmk:manin-question}
Manin~\cite{manin95} asked whether
$\overline{\mathrm{Spec}\,\mathbb{Z}}$ admits a structure of
absolute curve over~$\mathbb{F}_1$ whose zeta function is the
complete Riemann zeta.  Connes--Consani~\cite{connesconsani10}
answered affirmatively at the level of counting distributions:
\cref{thm:cc-zeta-Q} produces $N_\mathbb{Q}(u)$
explicitly.  However, their construction takes the zeros of
$\zeta(s)$ as analytic input
\textup{(}equation~\eqref{eq:N-Q} contains~$\sum_{\rho\in Z}$\textup{)},
and so does not exhibit the descent data structurally.

The pentagonal $\mathbb{F}_1$-descent of the present paper supplies
the complementary structural input at the level of
$\zeta_{\mathbb{Q}(\sqrt{5})}$: the $\Lambda_{\Z,E}$-structure on $\Zphi$ away from $5$
(\cref{thm:descent-datum}), whose Frobenius lifts produce the Euler factors directly from
$\golden^2 = \golden + 1$, without analytic input.  The two
constructions are complementary for $\mathbb{Q}(\sqrt{5})$ --- structure on the input
side, counting on the output side --- with the Dedekind factorisation
$\zeta = \zeta_K/L(\cdot,\chi_5)$ between them; neither answers Manin's question for
$\mathrm{Spec}\,\Z$.
\end{remark}

\begin{remark}[What remains analytic]\label{rmk:remains-analytic}
What is proved and what is classical here, exactly.  For $L(s,\chiv)$ the
continuation and the functional equation are proved in Lean from Mathlib's Dirichlet
$L$-functions (\LeanInline{functionalEquation\_chi5\_discharged},
\LeanInline{hecke\_chi5\_derived} \tierP); for $\Re s>1$ the Euler products of $\zeta$
and $L(s,\chiv)$ and the factorisation $\zeta_K=\zeta\cdot L(\cdot,\chiv)$ are proved
(\LeanInline{dedekind\_factorisation} \tierP).  The meromorphic continuation of
$\zeta_K$ to all of $\mathbb C$ is then the product of the two continuations; it is
not stated separately in Lean, and the classical account is Hecke's
\cite{hecke20} (see \cite[Ch.~VII]{neukirch99}).
The pentagonal $\mathbb{F}_1$-descent does not bypass these
analytic inputs; it provides a structural source for the
arithmetic data that they organise.
\end{remark}

% ---------------------------------------------------------------------
\subsection{The Mersenne bridge}\label{sec:mersenne}
\src{HP \texttt{ssec:mersenne} (L2187); \cite{acevedo20,acevedo21}}

The golden tower and the binary tower meet at one exponent, and the meeting fixes
the scale of the construction.  The binary tower $\{2^n\}$ carries the Mersenne numbers $2^e-1$
(\cref{prop:mersenne-exponent}), which gives the bridge its name.

\begin{proposition}\label{prop:mersenne-bridge}
With $\lng=\log2/\log\golden$ one has $\golden^{\lng}=2$ \tier[\testref{test\_mersenne\_bridge\_phi\_lambda\_log\_two}]{N}, so the binary tower
$\{2^n\}$ sits inside the golden one; and $2$ is not an integer power of
$\golden$.
\end{proposition}
\leanline{\LeanInline{mersenne\_bridge}, \LeanInline{binary\_tower\_inside\_golden},
\LeanInline{lambda\_is\_not\_an\_integer} \tierP.}
\begin{proof}
The first identity is the definition of $\lng$ together with
$\golden^{x}=e^{x\log\golden}$.  For the second: $n\mapsto\golden^n$ is strictly
increasing since $\golden>1$, and $\golden<2<\golden^2=\golden+1$, so $2$ falls
strictly between two consecutive powers and equals none.
\end{proof}

The second half is what gives the bridge its role.  The exponent $\lng$ is not
reachable by descent on the exponent lattice --- it is irrational, and the value
$2$ lies in a gap of the tower --- so it \emph{anchors} the scale instead of being
derived from it.  The
tower is discrete, and a real number is in it only if its logarithm is an integer
multiple of $\log\golden$.  The bit itself is the kernel: $\log2=\log|\ker G|$, the
kernel being exactly $\{1,9\}$ (\LeanInline{kernel\_is\_exactly\_one\_nine},
\LeanInline{log\_two\_is\_log\_kernel} \tierP).

The anchor is not special to $\golden$.  For every unit $\varepsilon>1$ of a real
quadratic field, $\varepsilon^{-\log2/\log\varepsilon}=\tfrac12$; for $\golden$ the
exponent is $\lng$, the half-turn is $\golden^{\lng/2}=\sqrt2$, and the counting in base
$\golden$ is the counting in base $2$ rescaled, $\log_\golden T=\lng\log_2T$
(\LeanInline{modulus\_transports}, \LeanInline{half\_three\_roles},
\LeanInline{density\_exact} \tierP).  What singles out $\golden$ is not the anchor but the
conductor: the unit whose field has conductor $5$.

The anchor also connects the arithmetic of $K$ to the one bit, $\log2$.  By \cref{thm:L1},
$L(1,\chiv)=2\log\golden/\sqrt5$, and then
\begin{equation}\label{eq:log-two}
  \lng\cdot\frac{\sqrt5}2\,L(1,\chiv)=\lng\cdot R_K=\log2
\end{equation}
(\LeanInline{branches\_meet\_at\_log\_two} \tierP),
the regulator cancelling exactly.  \Cref{eq:log-two} is the arithmetic branch of
the span of \cref{prop:span}, and $\log2=\log|\ker G|$ is where it meets the
spectral branch.

\begin{remark}[Where the Mersenne work meets this]
Acevedo Agudelo's classification of the Mersenne primes by residue classes modulo
$360$ and his study of the primes through angular rotations
\cite{acevedo20,acevedo21} read the primes by cyclic and angular structure, which
is the same register in which $\chiv$ is read off the pentagon
(\cref{thm:pentagonal-values}).  The point of contact is
\cref{prop:mersenne-bridge}: the binary tower of the Mersenne numbers inside the
golden one.  The companion paper \cite{gonzalez26mersenne} develops that thread.
\end{remark}

% ---------------------------------------------------------------------
\begin{proposition}[What the exponent of a Mersenne number fixes]\label{prop:mersenne-exponent}
Let $M=2^e-1$.
\begin{enumerate}[label=\textup{(\roman*)}]
  \item If $e\equiv1\pmod4$ then $M\equiv1\pmod5$, split, and $M\equiv11\pmod{20}$, $M\equiv31\pmod{60}$
    for $e\geq5$; if $e\equiv3\pmod4$ then $M\equiv2\pmod5$, inert, and $M\equiv7\pmod{20}$ and modulo
    $60$; and $M_2=3$ is inert.
  \item For $e$ odd, $M\equiv1\pmod3$.
  \item A Mersenne prime with $e\equiv3\pmod4$ divides $F_{M+1}$, by the inert case of the Fibonacci
    criterion (\cref{tab:fib-split} and the paragraph before it).
  \item The Mersenne primes fall in the classes $3$, $7$ and $11$ and never in the kernel $\{1,9\}$
    of $G$: always where the class fixes the $2$-adic part of the angular part of $\golden$ (\cref{prop:class-angular}).
  \item For $e\geq2$, $M\equiv3\pmod4$, and for $e$ odd $M\equiv1\pmod3$: a Mersenne prime is inert in the
    Gauss integers and split in the Eisenstein integers, and in the golden field by (i); its
    Frobenius in $\mathbb{Q}(i,\sqrt{-3},\sqrt5)$, the field of the three lattices, is $(-1,+1,\pm1)$.
\end{enumerate}
\end{proposition}
\leanline{\LeanStar{F1PaperPadic.mersenne\_mod\_five}, \LeanStar{F1PaperPadic.mersenne\_mod\_twenty}, \LeanStar{F1PaperPadic.mersenne\_mod\_sixty}, \LeanStar{F1PaperPadic.mersenne\_odd\_mod\_three}, \LeanStar{F1PaperPadic.fib\_succ\_zero\_iff\_inert}, \LeanStar{F1PaperPadic.mersenne\_inert\_rank}, \LeanStar{F1PaperPadic.mersenne\_classes\_checked}, \LeanStar{F1PaperPadic.mersenne\_frobenius\_three\_lattices} \tierPnew; \LeanInline{F1PCFUnified.mersenne\_mod\_four} \tierP;
\LeanInline{F1PCFUnified.mersenne\_concentration\_general} \tierP{} \cite{gonzalez26mersenne}.}
\begin{proof}
(i) Modulo $5$, $2$ has order four, with $2^e\equiv2,4,3,1$ for $e\equiv1,2,3,0\pmod4$; so $M\equiv1$ for $e\equiv1$
and $M\equiv2$ for $e\equiv3$.  For $e\geq2$, $M\equiv-1\equiv3\pmod4$, and with (ii) the Chinese remainder theorem
gives the classes modulo $20$ and $60$; $M_2=3$ has $\chiv(3)=-1$.
(ii) $2\equiv-1\pmod3$, so $2^e\equiv-1$ and $M\equiv-2\equiv1$ for $e$ odd.
(iii) For $q\neq2,5$, $\golden^{q+1}=\golden\cdot\golden^q\equiv\golden\,\Frob_q(\golden)\pmod q$
(\cref{thm:frobenius-congruence}).  If $q$ is inert this is $\golden(1-\golden)=N(\golden)=-1$, so
$F_{q+1}\golden+F_q\equiv-1$ and $F_{q+1}\equiv0$; if $q$ splits it is $\golden^2=\golden+1$, and $F_{q+1}\equiv1$.  A
Mersenne prime with $e\equiv3\pmod4$ is $\equiv2\pmod5$, hence inert.
(iv) A Mersenne prime has $e=2$ or $e$ an odd prime, so $M=3$ or $e\equiv1,3\pmod4$, and (i) gives the classes $3$,
$11$ and $7$, none in $\{1,9\}$.
(v) The classes modulo $4$ and $3$ are (i) and (ii); a prime $\equiv3\pmod4$ is inert in $\Z[i]$, and a prime
$\equiv1\pmod3$ splits in $\Z[\zeta_3]$.
\end{proof}

\subsection{The values in base $\golden$, from the anchor}\label{ssec:base-phi}

The classical formulas for the zeta values are written in the constants $2$, $\pi$
and $\sqrt5$.  In the present construction each of them is a function of $\golden$, and
the $L$-function of the pentagon is anchored at the two ends of its functional
equation.  The values can therefore be written in base $\golden$, starting from the
anchor, and not in the traditional base.

\begin{proposition}[The zeta values in base $\golden$]\label{thm:base-phi}
Let $\Lambda(s,\chiv)=(5/\pi)^{s/2}\,\Gamma(s/2)\,L(s,\chiv)$ be the completed $L$-function.
\begin{enumerate}[label=\textup{(\roman*)}]
  \item \emph{The anchor.}  $\Lambda(0,\chiv)=\Lambda(1,\chiv)=2\log\golden=\log(\golden+1)$,
    and $L(1,\chiv)=\dfrac{2\log\golden}{\sqrt5}=\dfrac{2\log\golden}{2\golden-1}$.
  \item \emph{The constants.}  $\sqrt5=2\golden-1$, $\pi=5\arccos(\golden/2)$,
    $2=\golden^{\lng}$ and $\tfrac12=\golden^{-\lng}$, with $\lng=\log2/\log\golden$.
  \item \emph{$s=0$.}  $\zeta(0)=-\tfrac12=-\golden^{-\lng}$.
  \item \emph{The even values.}  For $k\ge1$,
    \[
      \zeta(2k)=\frac{(-1)^{k+1}B_{2k}}{2\,(2k)!}\,
      \bigl(\golden^{\lng}\cdot5\arccos(\golden/2)\bigr)^{2k},
    \]
    where the $2\pi$ of Euler's formula has been written as
    $\golden^{\lng}\cdot5\arccos(\golden/2)$.  Likewise
    \[
      L(2k,\chiv)=(2\golden-1)\,\bigl(5\arccos(\golden/2)\bigr)^{2k}r_k,\qquad r_k\in\Q,
    \]
    the first two being
    \[
      L(2,\chiv)=\frac{4\arccos(\golden/2)^2}{2\golden-1},\qquad
      L(4,\chiv)=\frac{40\arccos(\golden/2)^4}{3\,(2\golden-1)}.
    \]
  \item \emph{The odd values.}  $\zeta(2k+1)=\zeta_K(2k+1)/L(2k+1,\chiv)$, $k\ge1$,
    (\cref{thm:odd-zeta}), with $L(\cdot,\chiv)$ the combination of Hurwitz zetas at the
    pentagonal points $a/5$, where $|2\cos(\pi a/5)|$ is $\golden$ or $\golden^{-1}$.
\end{enumerate}
\end{proposition}
\leanline{\LeanStar{F1v2.zeta\_values\_in\_base\_phi}, \LeanStar{F1v2.anchor\_in\_phi},
\LeanStar{F1v2.sqrt5\_in\_phi}, \LeanStar{F1v2.pi\_and\_two\_in\_phi},
\LeanStar{F1v2.zeta\_zero\_in\_phi}, \LeanStar{F1v2.even\_zeta\_in\_phi},
\LeanStar{F1v2.even\_L\_in\_phi}, \LeanStar{F1v2.L2\_L4\_in\_phi} \tierPnew; the
anchor \LeanInline{anchor\_is\_a\_crossed\_identity} (given the values of the anchor,
stated as \LeanInline{LOneIsTheAnchor}),
\LeanInline{sqrt\_five\_from\_phi}, \LeanStar{F1CW.pi\_eq\_five\_arccos},
\LeanInline{mersenne\_bridge}, \LeanStar{F1CW.even\_zeta\_bernoulli},
\LeanStar{F1CW.norm\_cosine\_phi} \tierPnew; the class-number formula and the closed forms of
$L(2,\chiv)$, $L(4,\chiv)$ \tierC/\tierN.}

\begin{proof}
(i) $2\log\golden=\log\golden^2=\log(\golden+1)$; $\Lambda(1,\chiv)=(5/\pi)^{1/2}\Gamma(\tfrac12)L(1,\chiv)=\sqrt5\,L(1,\chiv)$,
which is $2\log\golden$ by \cref{thm:L1}, and $\Lambda(0,\chiv)=\Lambda(1,\chiv)$ by the functional equation; $\sqrt5=2\golden-1$
rewrites $L(1,\chiv)$.  (ii) $\sqrt5=2\golden-1$ is
the definition of $\golden$; $\pi=5\arccos(\golden/2)$ is \cref{eq:pentagon-id},
$\golden=2\cos(\pi/5)$; $\golden^{\lng}=2$ is the Mersenne anchor.  (iii) $\zeta(0)=-\tfrac12$
and $\golden^{-\lng}=\tfrac12$.  (iv) Euler's formula with $2\pi$ rewritten by (ii); the
generalised Bernoulli formula for $L(2k,\chiv)$ with $\sqrt5$ and $\pi$ rewritten by (ii);
$L(2,\chiv)=4\pi^2/(25\sqrt5)$ and $L(4,\chiv)=8\pi^4/(375\sqrt5)$ rewritten the same way,
using $25\cdot\tfrac{4}{25}=4$ and $\tfrac{8\cdot625}{375}=\tfrac{40}3$.  (v) is
\cref{thm:odd-zeta} with \cref{thm:pentagonal-values}.
\end{proof}

\begin{remark}[What changes and what does not]\label{rmk:base-phi}
The values are the classical ones; what changes is the base in which they are written.
Traditionally $\zeta(2k)$ is a rational multiple of $\pi^{2k}$ and nothing more is said
about $2$ and $\pi$.  Here $2$ is the Mersenne anchor $\golden^{\lng}$, $\pi$ is the
pentagon $5\arccos(\golden/2)$, $\sqrt5$ is $2\golden-1$, and the starting point is the
anchor of the $L$-function at $s=0$ and $s=1$, which is $2\log\golden$ at both.  The
even values then carry $\golden$ in every constant, and the odd ones carry it through
$\zeta_K/L$ and the pentagonal points.  The rational coefficients $B_{2k}$ and $r_k$
are untouched: $r_1=4/5^3$ and $r_2=8/(3\cdot5^4)$, and written against
$\pi^{2k}/\sqrt5$ instead, $L(2,\chiv)=\tfrac4{25}\pi^2/\sqrt5$ and
$L(4,\chiv)=\tfrac8{375}\pi^4/\sqrt5$, whose denominators $5^2$ and $3\cdot5^3$ can be read as carrying the
conductor and the ternary (\LeanInline{denominators\_are\_conductor\_and\_ternary} \tierP\ for the factorisations).
\end{remark}

% =====================================================================

% ---------------------------------------------------------------------
\subsection{The operator}\label{sec:operator}

The operator is built from the tower and its fibre alone, before any zero is mentioned; in the thread of \cref{sec:topological-proof} it is Stage~5.  This section derives the fibre over the tower,
defines the operator on the joint index of level and fibre, proves that the forms in which it has been written are
functions of it, states the dimensional chain on which it acts, and gives its exact spectrum and the space on which
all the Frobenius lifts act at once.

\subsubsection{The fibre over the tower}\label{ssec:fibre}

The tower $\golden^{\Z}$ is indexed by the level alone.  Each level also has a reading at the inert prime $2$ and a
sign, and the two together are a point of a finite group; this subsection derives that group from $M$, the matrix of
multiplication by $\golden$ (\cref{ssec:unit-norm}).

\begin{definition}[The signature of a level]\label{def:signature}
For $n\in\Z$ let
\[
  s(n)=\bigl(M^n\bmod 2,\ \det M^n\bigr)\in GL_2(\F_2)\times\{\pm1\}.
\]
$M^n$ is multiplication by $\golden^n$ on $\Zphi$, so the first component is the class of $\golden^n$ in
$(\Zphi/2)^\times=\F_4^\times$ and the second is the norm $N(\golden^n)=(-1)^n$.
\end{definition}

\begin{theorem}[The fibre over the tower is $\mu_6$]\label{thm:fibre-mu6}
\begin{enumerate}[label=\textup{(\roman*)}]
  \item $s$ is a homomorphism from $\Z$ onto $\F_4^\times\times\{\pm1\}$, a cyclic group of order six, with kernel
    $6\Z$.  Its two components have orders three and two, so $n\mapsto(k,\epsilon)=(n\bmod3,\ n\bmod2)$ identifies
    the image with $\Z/3\times\Z/2\cong\mu_6$: $k$ is the \emph{colour} and $\epsilon$ the \emph{orientation}.
  \item $M^6-1=4M^3$ and $\det M^6=1$, while for $1\le n\le5$ either $M^n\not\equiv1\pmod2$ or $\det M^n=-1$.  So six
    is the least positive $n$ for which $M^n$ fixes every $2$-torsion point of the torus
    $\mathbb T=\R^2/\Z^2$ and preserves its orientation.
  \item Let $X$ be the mapping torus of $M$ on $\mathbb T$.  The composite $\pi_1(X)\to\Z\to\Z/6$ of the projection
    to the base circle with $s$ defines a connected cyclic cover $X_6\to X$ of degree six.  $X_6$ is the mapping
    torus of $M^6$, and its deck group is $\Z/6\cong\mu_6$, generated by the \emph{deck step} $D$, which advances one
    turn of the base circle.  By \textup{(ii)} it is the least cyclic cover of $X$ along the base on which the
    monodromy fixes $\mathbb T[2]$ pointwise and preserves the orientation of the fibre torus.
\end{enumerate}
{\raggedright \LeanInline{PCFGranularity.phiSig\_period\_six}, \LeanInline{PCFGranularity.phiSig\_six\_are\_distinct},
\LeanInline{PCFGranularity.cover\_degree\_six}, \LeanInline{PCFGranularity.M\_six\_sub\_one},
\LeanInline{PCFGranularity.colour\_factor\_has\_order\_three},
\LeanInline{PCFGranularity.orientation\_factor\_has\_order\_two}, \LeanInline{f4phi\_order\_three},
\LeanInline{toMu6\_bijective}, \LeanInline{generator\_to\_generator} \tierP; for (iii), the deck step and the two
identifications on $\Z^2\times\Z$: \LeanInline{TodoEnLean.gam\_Dstep} ($\gamma\circ D=\mathrm{id}$),
\LeanInline{TodoEnLean.gam6\_Dstep6} ($\gamma_6\circ D^6=\mathrm{id}$), \LeanInline{TodoEnLean.Dstep\_gam6}
($D\gamma_6=\gamma_6D$), \LeanInline{TodoEnLean.Dstep\_iter\_shift} (the base shift of $D^k$ is $k$) \tierP; that the
quotients by these properly discontinuous free actions are covering spaces \tierC\par}
\end{theorem}

\begin{proof}
(i) Reduction modulo $2$ and the determinant are multiplicative, so $s(m+n)=s(m)s(n)$.  The class of $\golden$
generates $\F_4^\times$, of order three, and $\det M=-1$ has order two (\cref{ssec:unit-norm}); the orders are
coprime, so $s(1)$ has order six, the image is cyclic of order six and the kernel is $6\Z$.
(ii) From $M^2=M+1$, $M^3=M^2+M=2M+1$ and $M^6=(2M+1)^2=4(M^2+M)+1=4M^3+1$; $\det M^6=(-1)^6$.  For
$n=1,2,4,5$ the class of $\golden^n$ in $\F_4^\times$ is not $1$, and $\det M^3=-1$.  $M^n$ acts on
$\mathbb T[2]=(\tfrac12\Z/\Z)^2$ as $M^n\bmod2$, and it preserves the orientation exactly when $\det M^n=1$.
(iii) Write $X=\mathbb T\times\R/\!\sim$ with $(x,t+1)\sim(Mx,t)$, and $X_6=\mathbb T\times\R/\!\sim_6$ with
$(x,t+6)\sim_6(M^6x,t)$, the mapping torus of $M^6$.  The identity of $\mathbb T\times\R$ induces a covering
$X_6\to X$ of degree six, the one attached to the subgroup of $\pi_1(X)=\Z^2\rtimes_M\Z$ over $6\Z$.
$D[x,t]=[M^{-1}x,t+1]$ covers the identity of $X$, and $D^6[x,t]=[M^{-6}x,t+6]=[x,t]$, with no smaller power the
identity.  The monodromy of the $n$-fold cyclic cover along the base is $M^n$, and \textup{(ii)} applies.
\end{proof}

The fibre is therefore a result: its six points are the six signatures, its two factors are the reading of the
level at the prime $2$ and the sign of the norm, and its group is the deck group of a cover of the mapping torus
that $M$ itself determines.

\begin{definition}[The joint index and its space]\label{def:joint-index}
$\mathcal J=\N\times\Z/3\times\Z/2$.  A point $(\sigma,k,\epsilon)$ is a level $\sigma$ of the tower together with a
point $(k,\epsilon)$ of the fibre of $X_6\to X$, identified with $\Z/3\times\Z/2$ by
\cref{thm:fibre-mu6}\textup{(i)}.  The space is $\ell^2(\mathcal J)=\ell^2(\N)\otimes\C^3\otimes\C^2$, with
orthonormal basis $e_{\sigma,k,\epsilon}$, and the deck step acts on it by
$D\,e_{\sigma,k,\epsilon}=e_{\sigma,k+1,\epsilon+1}$, a unitary of order six that preserves the level.
\end{definition}

\subsubsection{The operator on the joint index}\label{ssec:the-operator}

\begin{definition}[The operator and its parts]\label{def:operators-J}
Let $m_0>0$ be a scale and $\zeta_3=e^{2\pi i/3}$.  On $\ell^2(\mathcal J)$ the following operators are diagonal in
the basis $e_{\sigma,k,\epsilon}$:
\[
  H\,e_{\sigma,k,\epsilon}=m_0\golden^{\sigma}\,e_{\sigma,k,\epsilon},\qquad
  \OmH\,e_{\sigma,k,\epsilon}=\tfrac12\zeta_3^{\,k}\,e_{\sigma,k,\epsilon},\qquad
  W\,e_{\sigma,k,\epsilon}=(-1)^{\epsilon}\,e_{\sigma,k,\epsilon},
\]
the \emph{tower operator} (the level), the \emph{colour operator} and the \emph{orientation}; $H$ has domain
$\{v:\sum m_0^2\golden^{2\sigma}|v_{\sigma,k,\epsilon}|^2<\infty\}$, and the eigenvalues of $\OmH$ are the triad
$|2|_2\,\zeta_3^{\,k}$ of \cref{def:tripartite-omega}, so that on the colour factor
$\OmH=\tfrac12\operatorname{diag}(1,\zeta_3,\zeta_3^2)$.  The \emph{phase} is $Y=2\,\OmH\,W$, and
\[
  \Oone\;=\;H\,Y ,\qquad
  \Oone\,e_{\sigma,k,\epsilon}=m_0\,\golden^{\sigma}\,\zeta_3^{\,k}\,(-1)^{\epsilon}\,e_{\sigma,k,\epsilon},
\]
with the domain of $H$, is \emph{the operator}.
\end{definition}

\begin{theorem}[One operator]\label{thm:one-operator}
\begin{enumerate}[label=\textup{(\roman*)}]
  \item \emph{Normal, with simple spectrum.}  $\Oone$ is normal, and its eigenvalues
    $m_0\golden^{\sigma}\zeta_3^{\,k}(-1)^{\epsilon}$ are pairwise distinct.  So $\Oone$ alone separates the points
    of $\mathcal J$, and every bounded operator that commutes with $\Oone$ is a function of $\Oone$.
  \item \emph{Modulus and phase.}  $\Oone^{*}\Oone=H^2$, so $|\Oone|=H$, positive and self-adjoint with trivial
    kernel, and $\Oone=Y\,H=H\,Y$ is the polar decomposition of $\Oone$, which is unique.  $Y$ is unitary of order
    exactly six, with spectrum $\mu_6$.  The self-adjoint form of the operator is its modulus: the tower form
    $H=\operatorname{diag}(m_0\golden^{\sigma})$, whose eigenvalues have ratios
    $\lambda_i/\lambda_j=\golden^{\,i-j}$.  Its real and imaginary parts
    $\tfrac12(\Oone+\Oone^{*})$ and $\tfrac1{2i}(\Oone-\Oone^{*})$ are commuting self-adjoint operators.
  \item \emph{The forms are functions of the operator.}  $H=|\Oone|$, $Y=\Oone\,|\Oone|^{-1}$, $W=Y^3$ and
    $\OmH=\tfrac12Y^4=|2|_2\,Y^4$; and every operator of the level --- diagonal, with entries that depend on $\sigma$
    alone, as $\operatorname{diag}T^*(n)$ --- is $g(H)$ for a function $g$ on the spectrum of $H$.
  \item \emph{The deck.}  $DH=HD$ and $DYD^{-1}=\zeta_6Y$ with $\zeta_6=e^{i\pi/3}$, so
    $D\,\Oone\,D^{-1}=\zeta_6\,\Oone$: over each level the six eigenvalues are the vertices of a regular hexagon of
    radius $m_0\golden^{\sigma}$, which the deck step rotates.  A function of $\Oone$ commutes with $D$ exactly when
    it is a function of $H$.
\end{enumerate}
{\raggedright \LeanInline{UnOperador.spec\_injective}, \LeanInline{UnOperador.norm\_spec},
\LeanInline{UnOperador.spec\_polar}, \LeanInline{UnOperador.norm\_yDiag}, \LeanInline{UnOperador.forms\_from\_spec},
\LeanInline{FibraYReloj.yDiag\_pow\_three}, \LeanInline{FibraYReloj.yDiag\_pow\_four},
\LeanInline{FibraYReloj.yDiag\_injective}, \LeanInline{FibraYReloj.deck\_weyl} \tierP, on the eigenvalues; the
operator statements: \LeanInline{TodoEnLean.diagonal\_normal}, \LeanInline{TodoEnLean.commute\_diag\_offdiag},
\LeanInline{TodoEnLean.commute\_diag\_is\_function}, \LeanInline{TodoEnLean.polar\_unique},
\LeanInline{TodoEnLean.yDiag\_order\_six}, \LeanInline{TodoEnLean.spec\_deck},
\LeanInline{TodoEnLean.function\_commutes\_with\_deck\_iff}, \LeanInline{TodoEnLean.modulus\_deck\_invariant} \tierP,
the commutant on every finite set of levels, and entrywise on $\ell^2(\mathcal J)$ by the proof below\par}
\end{theorem}

\begin{proof}
All the operators are diagonal in one orthonormal basis, and the adjoint of a diagonal operator is the diagonal
operator of the conjugate entries, on the same domain.
(i) Diagonal operators with a common domain commute, so $\Oone\Oone^{*}=\Oone^{*}\Oone$.  The six numbers
$\zeta_3^{\,k}(-1)^{\epsilon}$ are the six sixth roots of unity, of modulus $1$.  If
$m_0\golden^{\sigma}y=m_0\golden^{\sigma'}y'$ with $|y|=|y'|=1$, the moduli give $\golden^{\sigma}=\golden^{\sigma'}$,
hence $\sigma=\sigma'$ because $\golden>1$, and then $y=y'$.  Let $B$ be bounded and commute with $\Oone$, and let
$\lambda_i$ be the eigenvalue at the basis vector $e_i$.  Then
$\lambda_i\langle e_j,Be_i\rangle=\langle e_j,B\Oone e_i\rangle=\langle\Oone^{*}e_j,Be_i\rangle
=\lambda_j\langle e_j,Be_i\rangle$, so $\langle e_j,Be_i\rangle=0$ for $i\ne j$: $B$ is diagonal, and it is
$f(\Oone)$ with $f(\lambda_i)=\langle e_i,Be_i\rangle$, which is well defined because the $\lambda_i$ are distinct.
(ii) $\Oone^{*}\Oone$ has entries $|m_0\golden^{\sigma}y|^2=(m_0\golden^{\sigma})^2$, those of $H^2$, and $H$ is the
positive square root.  $H$ has no kernel, so the partial isometry of the polar decomposition is determined, and it
is $Y$.  $y^6=1$ for each entry $y$ of $Y$, and the entry at $(k,\epsilon)=(2,1)$ is $-\zeta_3^{2}=\zeta_6$, of
order six.  The ratios are $m_0\golden^{i}/m_0\golden^{j}$.
(iii) For $y=\zeta_3^{\,k}(-1)^{\epsilon}$, $y^3=(-1)^{\epsilon}$ and $y^4=\zeta_3^{4k}=\zeta_3^{\,k}$.  The map
$\sigma\mapsto m_0\golden^{\sigma}$ is injective, so $g(m_0\golden^{\sigma})$ may be prescribed level by level.
(iv) $D$ preserves the level.  $D\,\OmH\,D^{-1}e_{\sigma,k,\epsilon}=\tfrac12\zeta_3^{\,k-1}e_{\sigma,k,\epsilon}$ and
$DWD^{-1}=-W$, so $DYD^{-1}=-\zeta_3^{-1}Y=\zeta_6Y$.  $f(\Oone)$ commutes with $D$ exactly when
$f(\zeta_6\lambda)=f(\lambda)$ on the spectrum, that is, when $f$ is constant on each hexagon: a function of the
modulus.
\end{proof}

This is the operator used from here on.  It is normal; what is self-adjoint in it is its modulus $H$, which reads the
level, and what its modulus leaves to be read is its phase $Y$, which reads the fibre of \cref{thm:fibre-mu6}; the
colour operator $\OmH$ and the orientation $W$ are the two powers $\tfrac12Y^4$ and $Y^3$ of that phase.

\subsubsection{The line, the plane and the third dimension}\label{ssec:dimensional-chain}

The modulus of the operator lives on a line, the line of the exponents; its phase and the Euler products live on a
plane; and its ratios $\golden^{\,i-j}$ act on a third dimension, where the triad of $\OmH$ is permuted.  This
subsection states that chain and the numbers that recur along it.

The previous version of this paper proves the chain of dimensions
$\R\xrightarrow{\times i}\C\xrightarrow{\times\golden}E^3$
\cite[\S3.6.1, Proposition `Dimensional chain from self-measurement']{gonzalez26oddzeta}, and its discussion
\cite[\S6]{gonzalez26oddzeta} reads it: the Euler product emerges from the line, the zeta functions live on the
plane, and from the two the third dimension emerges through the golden coupling.  In the dictionary of this paper,
built by classification under equivalence (\cref{sec:objective}), a dimension is not a property of one object: it is
a class under equivalence, and each class is realised in several domains.  The
following theorem states the chain with the equivalences it rests on, and places in it
the one ingredient the previous version did not have: the ratios
$\lambda_i/\lambda_j=\golden^{\,i-j}$ of the operator.

\begin{theorem}[The dimensional chain]\label{thm:one-two-three}
\begin{enumerate}[label=\textup{(\arabic*)}]
  \item \emph{The line, $d=1$.}  $\sigma\mapsto\golden^\sigma$ is an isomorphism
    $(\R,+)\cong(\R_{>0},\times)$, with the exponent lattice $\Z$ mapping onto the tower
    $\golden^\Z$.  On it the lift $\Frob_p(\golden^\sigma)=\golden^{p\sigma}$ is
    multiplication of the exponent by $p$, and $\Frob_q\circ\Frob_p=\Frob_{pq}$: the lifts
    act as the multiplicative monoid of the positive integers.  That same monoid indexes
    the Dirichlet series $\sum_{n\ge1}n^{-s}$, and unique factorisation writes it as the
    Euler product $\prod_p(1-p^{-s})^{-1}$ for $\Re s>1$.
  \item \emph{The plane, $d=2$.}  $\C=\R[x]/(x^2+1)$; the line winds onto the circle,
    $\theta\mapsto e^{i\theta}$ with period $2\pi$, and the tower closes on the torus
    $T^2_{\PCF}$ with period $T_K=2\pi\log\golden$.  The Euler products converge on the
    plane, where $\zeta$, $\zeta_K$ and $L(s,\chiv)$ live.  The second dimension is
    also the arithmetic one: $\Zphi\otimes\R\cong\R^2$, and a prime sees it exactly when
    it is inert, the Frobenius being the identity for $\chiv(p)=+1$ and the
    conjugation for $\chiv(p)=-1$ (\cref{thm:chi5-from-frobenius}).
  \item \emph{$E^3$, $d=3$.}  The golden coupling $z=\golden y$ extends $\C$ to
    $E^3=\{(x,y,\golden y)\}$, on which $\golden^2=\golden+1$ is the self-measurement
    $z^2=zy+y^2$: the square of the new coordinate folds back onto the ones that
    generated it.  The operator acts on $E^3$ through its ratios: the ratio at index
    distance $k=i-j$ sends the coupled coordinate to
    $\golden^{k}y=F_{k}\,z+F_{k-1}\,y$, an integer combination of $z$ and $y$, so
    $E^3$ is invariant under every ratio and the fold-back has Fibonacci
    coefficients; the ratio has norm $\golden^k\bar\golden^{\,k}=(-1)^k$, the parity of
    the index distance.
\end{enumerate}
\LeanStar{F1v2.dimensional\_chain}, \LeanStar{F1v2.line\_is\_the\_tower},
\LeanStar{F1v2.frobenius\_on\_the\_line}, \LeanStar{F1v2.euler\_product\_from\_the\_line},
\LeanStar{F1v2.plane\_from\_the\_line}, \LeanStar{F1v2.self\_measurement\_E3},
\LeanStar{F1v2.ratio\_folds\_back}, \LeanStar{F1v2.E3\_closed\_under\_ratios},
\LeanStar{F1v2.ratio\_norm} \tierPnew; \LeanInline{self\_measurement},
\LeanStar{F1CW.dirichlet\_is\_euler}, \LeanStar{F1CW.periodK},
\LeanInline{isometry\_with\_common\_modulus₅}, \LeanStar{F1CW.hinge\_central\_charge} \tierP
\end{theorem}

\begin{proof}
(1) $\golden^{s+t}=\golden^s\golden^t$; injectivity by taking logarithms, since
$\log\golden\ne0$; and $u>0$ is $\golden^{\log u/\log\golden}$.  The lift is
$(\golden^\sigma)^p=\golden^{p\sigma}$, and two lifts compose to the product of
exponents.  The identification of the Dirichlet series over the integer levels with
the Euler product is unique factorisation.  (2) $i^2=-1$ and
$e^{i(\theta+2\pi)}=e^{i\theta}$; the arithmetic statement is
\cref{thm:chi5-from-frobenius}, and the factorisation is \cref{eq:dedekind}.
(3) $(\golden y)^2-(\golden y)y-y^2=y^2(\golden^2-\golden-1)=0$.  By Binet,
$\golden^{k+1}=F_{k+1}\golden+F_k$ in any ring where $\golden^2=\golden+1$, which
multiplied by $y$ is the fold-back; the invariance of $E^3$ is
$\golden^k(\golden y)=\golden(\golden^k y)$; the norm is
$(\golden\bar\golden)^k=(-1)^k$.
\end{proof}

\tier{R} In the same terms the predecessor version \cite{gonzalez26oddzeta} read the factorisation
$\zeta=\zeta_K/L(\cdot,\chiv)$ as a projection: the one-dimensional $\zeta$, over the primes of $\Q$, is the
quotient of the two-dimensional $\zeta_K$, over the primes of $\Q(\sqrt5)$, by the character that records how
the second coordinate projects onto the first; the even values are the circle already closed, and the odd
values the terms where the generator, through $\chiv$, and the circle, through the points $a/5$, are coupled.

The next statement collects the numbers that recur along the chain.

\begin{theorem}[The numbers of the chain: $|\ker G|=2$, $\mu_3=\tfrac12$, $S_3$]\label{thm:S3-from-euler}
The following hold, each separately.  The Euler product converges for $\Re s>1$, its
local factors are classified by the classes of $\Zst$, and the image of $G$ is the
Klein group with $|\ker G|=2$.  The contraction factor of the spectral parameter $\mu_3=\tfrac12$ is
$1/\mu_3=2=|\ker G|$.  The \emph{spectral parameters} of dimension $n$ are $\sigma_n=n/2$ and $\mu_n=2-n/2$, so that
$\sigma_n+\mu_n=2$; at $n=3$, the dimension of $E^3$, they are $(\sigma_3,\mu_3)=(\tfrac32,\tfrac12)$, and their product
$\sigma_3\mu_3=\tfrac34=3\cdot\bigl(1\cdot(-\tfrac12)\cdot(-\tfrac12)\bigr)$ is three times the product of the real parts of
the cube roots of unity.  Conversely the constraints $\sigma+\mu=2$,
$\sigma\mu=\tfrac34$, $0<\mu<1$ force $(\sigma,\mu)=(\tfrac32,\tfrac12)$; and the three
eigenvalues $\tfrac12\zeta_3^k$ of $\OmH$ form an equilateral triangle of radius $\tfrac12$,
permuted by $S_3$, with tripartite norms $\|P\|\|C\|\|F\|=\tfrac12$ (\cref{thm:norm-Omega-half}).
\LeanInline{G\_kernel}, \LeanInline{contraction\_is\_kernel},
\LeanInline{spectral\_uniqueness}, \LeanInline{mu\_ternary\_param}, \LeanInline{sigma\_ternary},
\LeanInline{sigma\_mu\_is\_d\_times\_cos\_product}, \LeanInline{norm\_Omega\_is\_half} \tierP; they are components of
\LeanStar{F1v2.S3\_E3\_c3\_from\_the\_euler\_product} \tierPnew
\end{theorem}

\begin{proof}
$G^{-1}(0,0)=\{1,9\}$; $1/\mu_3=2$.  Substituting $\mu=2-\sigma$ gives
$\sigma^2-2\sigma+\tfrac34=0$, with roots $\tfrac32,\tfrac12$, and only $\mu=\tfrac12$
lies in $(0,1)$.  $\zeta_3^3=1$ and $|\zeta_3^k/2|=\tfrac12$;
$\tfrac1{\sqrt3}\cdot1\cdot\tfrac{\sqrt3}2=\tfrac12$.
\end{proof}

The three-dimensional class has one more realisation, internal to the arithmetic.
The lattice $\Zphi$ of the second dimension carries a space of dimension three --- its
binary quadratic forms --- on which the operator's ratios act by boosts.

\begin{proposition}[The binary forms of $\Zphi$ as a domain of $d=3$]\label{prop:binary-forms}
On the forms $ax^2+bxy+cy^2$ the discriminant $b^2-4ac=b^2+(a-c)^2-(a+c)^2$ has
signature $(2,1)$, so its isometry group is $SO(2,1)\cong PSL_2(\R)$.  Multiplication
by $\golden$ acts by $(a,b,c)\mapsto(c,\,b+2c,\,a+b+c)$, an integral isometry with
eigenvalues $\golden^2,-1,\golden^{-2}$: it reverses the norm form $x^2+xy-y^2$, of
discriminant $5=\Delta_K$, and boosts the null forms $(x+\golden y)^2$,
$(x+\bar\golden y)^2$ by $\golden^{\pm2}$, a boost of rapidity $2\log\golden$.  The ratio
$\golden^{\,i-j}$ of the operator boosts them by $\golden^{\pm2(i-j)}$.
\LeanStar{F1v2.symF\_preserves\_disc}, \LeanStar{F1v2.discQ\_signature},
\LeanStar{F1v2.norm\_form\_disc\_and\_reversed}, \LeanStar{F1v2.null\_boost\_plus},
\LeanStar{F1v2.null\_boost\_minus}, \LeanStar{F1v2.null\_forms\_and\_rapidity},
\LeanStar{F1v2.ratio\_boost} \tierPnew
\end{proposition}

\begin{proof}
The pull-back of $Q$ by $(x,y)\mapsto(y,x+y)$ is the stated map; the discriminant
scales by $\det^2=1$.  The characteristic polynomial is
$-(\lambda+1)(\lambda^2-3\lambda+1)$.  The eigenvector computations use only
$\golden^2=\golden+1$ and $\bar\golden^2=\bar\golden+1$, and $\golden^2\bar\golden^2=1$.
An eigen-direction stays one under iteration.
\end{proof}

\subsubsection{The colour operator $\OmH$ and its exact spectrum}\label{ssec:omega}

\begin{proposition}[Exact spectrum of $\OmH$]\label{prop:omega-spectrum}
On the colour factor, $\OmH=\tfrac12\operatorname{diag}(1,\zeta_3,\zeta_3^2)$ (\cref{def:operators-J}) has eigenvalues of modulus
$\tfrac12$ with real parts $\tfrac12,-\tfrac14,-\tfrac14$; the product of the moduli is
$2^{-3}$ and that of the real parts is $2^{-5}$; the trace of the real parts is $0$;
and $\tfrac12$ is the only positive value.
\LeanInline{spectral\_positions\_values}, \LeanInline{eigenvalue\_half}, \LeanStar{F1CW.triad\_real\_parts}, \LeanInline{positive\_spectral\_half} \tierP
\end{proposition}

Its spectrum is exact, not asymptotic, and it is fixed before any zero is
mentioned: this is the non-circularity requirement of
\cref{ssec:four-requirements}, met by construction
(tallied in \LeanTally{omega\_meets\_three}).

\begin{proposition}[The spectrum of $\OmH$ read at the place $2$]\label{prop:omega-two-adic}
The eigenvalues of $\OmH$ are $\tfrac12\zeta_3^{\,k}$.  The modulus $\tfrac12$ is $|2|_2$, the $2$-adic norm of the
prime $2$, and by the product formula $|2|_2=1/|2|_\infty=\golden^{-\lambda_{\log}}$; the angle $\zeta_3$ is the angular part of $\golden$ at $2$ (\cref{prop:polar}).  Under either $2$-adic embedding of $\mathbb{Q}(\zeta_3)$ --- the two exchange $\zeta_3$ and $\zeta_3^2$ and permute the eigenvalues --- the spectrum is
$\{|2|_2\,[\golden_2]^{\,k}\}_{k=0,1,2}$, the factorisation of \cref{thm:operator-forms} read $2$-adically.
{\raggedright \LeanStar{F1PaperPadic.half\_is\_two\_adic\_norm\_of\_two}, \LeanStar{F1PaperPadic.polar\_at\_two}, \LeanStar{F1PaperPadic.minus\_fifteen\_is\_a\_two\_adic\_square} \tierPnew;
\LeanInline{mu\_from\_tower} \tierP\par}
\end{proposition}

\subsubsection{One operator, several forms}\label{ssec:operator-forms}

The operator has been written from the spectrum,
from the arithmetic of the anchor, from the Klein group, from an isometry, from the
geometry of the torus, from the tower, and --- in the Primitive Structure and again
in the corpus --- as the operator whose eigenvalues follow the heights of the zeros.
\Cref{thm:one-operator} says what unites them: the spectral and arithmetic forms are the power
$\tfrac12Y^4$ of the phase of $\Oone$, the tower form is its modulus $H=|\Oone|$, the ordinate form is a function
$g(H)$ of that modulus, and the geometric form is its fibre and its deck step.  The theorem below records each form
with what is proved of it; each reaches the apex $\tfrac12$.

\begin{theorem}[The forms of the operator]\label{thm:operator-forms}
Let $\Oone=HY$ be the operator of \cref{def:operators-J}.
\begin{enumerate}[label=\textup{(\roman*)}]
  \item \emph{Spectral.}  $\OmH=\tfrac12Y^4$ is $\tfrac12\operatorname{diag}(1,\zeta_3,\zeta_3^2)$ on the colour factor, with the
    spectrum of \cref{prop:omega-spectrum}.
  \item \emph{Arithmetic.}  $\OmH_k=\golden^{-\lng}\zeta_3^k$: the spectrum is the
    Mersenne anchor $\golden^{-\lng}=\tfrac12$ times the orbit of the ternary rotation;
    the same $\tfrac12$ is the factorial face $\Gamma(\tfrac32)/\sqrt\pi$.
  \item \emph{From the group.}  The three components $(P,C,F)$ of $\OmH$ are the three
    involutions $\sigma(s)=1-s$, $\tau(s)=\bar s$ and $\sigma\tau(s)=1-\bar s$ of the Klein group $V=\operatorname{im}G$:
    the assignment $(P,C,F)\mapsto(\sigma,\tau,\sigma\tau)$ is a bijection onto the three elements of $V$ other than the identity,
    each of which is an involution (\LeanInline{ComponentesInvoluciones.components\_are\_the\_involutions} \tierP).  The involutions
    come down from the parity and the reality of $\chiv$: negation $a\mapsto-a$ and inversion $a\mapsto a^{-1}$ on $\Zst$
    preserve $\chiv$ because it is even and real (\LeanInline{both\_actions\_from\_one\_character} \tierP; the action on the zeros is \cref{prop:descent}).
  \item \emph{Isometric.}  Normalised, the spectrum is a column $\mathcal V\colon\mathbb{C}\to\mathbb{C}^3$ with $\mathcal V^\dagger\mathcal V=1$ (\cref{ssec:why-phi}); the diagonal unitary $\operatorname{diag}(i,\zeta_3)$, whose entries generate the rotations of the Gauss lattice $\Z[i]$ and of the Eisenstein lattice $\Z[\zeta_3]$, satisfies the same identity; and the norm of the operator is $\tfrac12$.
  \item \emph{Geometric.}  On the torus the Fibonacci matrix $M$ has no periodic
    frequency --- Fourier blindness, which is the irrationality of $\golden$ --- while
    on its finite quotients it has period exactly $6$; the deck group of the fibre is
    $\mu_6$ (\cref{thm:fibre-mu6}), the unit group of $\Z[\zeta_3]$, which contains the cube roots of unity, the phases of the spectral form.  The level operator $H$ is not a complete set of commuting
    observables: it has to be joined by the fibre operators, and a deck-invariant
    observable is a function of the level alone; the fibre and its dual have the same
    size.
  \item \emph{Tower.}  $H=|\Oone|=\operatorname{diag}(m_0\golden^i)$, with $m_0>0$ a scale, self-adjoint, with ratios
    $\lambda_i/\lambda_j=\golden^{\,i-j}$; its diagonal $i=j$ gives $r=1$.
  \item \emph{Ordinates.}  $\tilde H_{\PCF}=\operatorname{diag}T^*(n)$, a function of $H$ (\cref{thm:one-operator}\textup{(iii)}), with
    \[
      T^*(n)=c(n)\,\frac{\pi n}{\log(\mu n+\sigma)},\qquad
      c(n)=2+\frac{\alpha_0}{1+\beta_0\log n},
    \]
    written in \cite{gonzalez25hp}, in its regime $n>30$, with the numbers
    $\mu=\tfrac12$, $\sigma=\tfrac32$, $\alpha_0=\tfrac{59}{40}$, $\beta_0=\tfrac18$ (for
    $n\le30$ that paper interpolates $\alpha,\beta$ adaptively), and written in the corpus with the same
    parameters as powers of $\golden$: $\mu=\golden^{-\lng}$, $\sigma=3\golden^{-\lng}$,
    $\beta_0=\golden^{-3\lng}$, $\alpha_0=\golden^{-\lng}(3-\tfrac1{20})$.  The two are the
    same function, not an approximation of one by the other.  In the second form every
    parameter is structural: $\mu$ is the apex $1/|\ker G|$, $\sigma+\mu=2$ is the trace of
    the diagonal, $\sigma/\mu=3$ the ternary order, $\beta_0=1/\phi(20)$, with $\phi$ Euler's function, the inverse
    number of units of the conductor, $\alpha_0=(dN-1)/(|\ker G|N)$ with $d=3$, $N=20$,
    and $\mu n+\sigma=\golden^{-\lng}(n+3)$ puts the ternary inside the logarithm; the
    factor $\pi n$ is the angular face and $\golden^{-\lng}$ the radial one.  A third
    writing expresses every parameter through the clauses of the Mersenne
    correspondence $3\golden^{\sigma(p)}=2^p$ \cite{gonzalez26mersenne}.  Its eigenvalues
    follow the heights $\gamma_n$ of the zeros.
\end{enumerate}
Of the faces that give a value --- tower, $T^*$ in its two writings, and $\OmH$ ---
all land on the value $2$ of the trace $t=r+1/r$ of a ratio $r>0$, the only value at which $r=1$ (it is the focus of the
confocal pencil of \cref{ssec:joukowski}); the fifth face counted in the tally \LeanTally{five\_faces\_four\_land} gives no value.
\LeanInline{eigenvalue\_half}, \LeanInline{F1PCFUnified.spectrum\_is\_anchor\_times\_orbit},
\LeanInline{F1PCFUnified.operator\_is\_tied\_to\_the\_cocone},
\LeanInline{F1PCFUnified.operator\_is\_made\_of\_the\_group},
\LeanInline{F1PCFUnified.V\_unitary}, \LeanStar{F1PaperNew.triad\_isometry},
\LeanInline{PCFGranularity.chaotic\_to\_fourier\_periodic\_in\_arithmetic},
\LeanInline{PCFGranularity.blindness\_is\_irrationality},
\LeanInline{PCFGranularity.the\_deck\_group\_is\_the\_eisenstein\_unit\_group},
\LeanInline{PCFGranularity.H\_not\_complete},
\LeanInline{PCFGranularity.deck\_invariant\_is\_level\_function},
\LeanInline{PCFGranularity.dual\_and\_fibre\_have\_the\_same\_size},
\LeanStar{F1CW.tower\_ham\_is\_hermitian}, \LeanInline{diagonal\_gives\_one},
\LeanInline{hpps\_parameters\_are\_the\_cocone}, \LeanInline{four\_parameters\_in\_phi},
\LeanInline{Tstar\_rewritten\_in\_phi}, \LeanInline{eleven\_components\_anchored},
\LeanInline{Tstar\_carries\_both\_faces}, \LeanInline{Tstar\_in\_connector\_clauses},
\LeanInline{Tstar\_rewritten\_in\_phi} \tierP; the comparison of the faces is recorded as
\LeanTally{tstarPhi\_is\_the\_bridge}, \LeanTally{five\_faces\_four\_land};
the identification of the components with the involutions in (iii), which the corpus takes as the hypothesis
of \LeanInline{F1PCFUnified.operator\_is\_made\_of\_the\_group}, is proved in
\LeanInline{ComponentesInvoluciones.components\_are\_the\_involutions} \tierP, and the
descent of the involutions from $\chiv$ is \tierP; the third writing of (vii) is proved
given the Mersenne clause as hypothesis (\LeanInline{Tstar\_in\_connector\_clauses});
the agreement of $T^*(n)$ with $\gamma_n$ \tierN
\end{theorem}

\begin{remark}[Why the ordinate form is the bridge]\label{rmk:tstar-bridge}
Among the forms, only $T^*$ written in $\golden$ carries both the exponent $\lambda$ of the
anchor and the ordinates of the zeros (recorded as \LeanTally{tstarPhi\_is\_the\_bridge}; the
rewriting itself is \LeanInline{Tstar\_rewritten\_in\_phi} \tierP): the
tower has the exponent but not the ordinates, the numerical form of
\cite{gonzalez25hp} the ordinates but not the exponent.  Rewriting the operator of the Primitive Structure in base $\golden$ is therefore not a change of notation: it is what connects
the operator of \cite{gonzalez25hp} with the tower of this paper, and through the tower
with the thread of \cref{sec:topological-proof}, which closes the Methods.  The distance between $T^*(n)$ and
$\gamma_n$ is the parametrisation of the envelope plus the fluctuation, the first by far the
larger (\cref{prop:error-partition,prop:positions}), and \cref{sec:results} shows what it does and
does not carry (\cref{prop:envelope-exact}).  The same forms, read by dimension ---
the tower on the line, the torus on the plane, the triad and the ratios
$\golden^{\,i-j}$ in $E^3$ --- are the dimensional chain of \cref{thm:one-two-three}, and all are read from $\Oone$ (\cref{thm:one-operator}).
\end{remark}

\begin{proposition}[The Eisenstein and deck arithmetic of the geometric form]\label{prop:eisenstein-deck}
In the Eisenstein integers the cube root of unity $\zeta_3$ has order exactly three,
and the primitive sixth root is $-\zeta_3^2$.  The deck step of the fibre has order
exactly six, the orientation flip order exactly two, and the Fibonacci matrix $M$,
of determinant $-1$, has $\det M^2=1$ and $M^2=M+1$.
\LeanInline{w\_cubed}, \LeanInline{w\_order\_three}, \LeanInline{z6\_is\_minus\_w\_squared},
{\raggedright\LeanInline{deck\_order\_six}, \LeanInline{orientation\_order\_two}, \LeanInline{det\_M\_sq},
\LeanInline{one\_identity\_three\_readings} \tierP\par}
\end{proposition}

\begin{proposition}[The spectrum as free part times torsion;\newline the parameters of $T^*$ in $\golden$]\sloppy\label{prop:spectrum-free-torsion}
\begin{enumerate}[label=\textup{(\roman*)}]
  \item $\OmH_k=\golden^{-\lng}\,\zeta_3^{\,k}$ (the factorisation already stated in
    \cref{thm:operator-forms}): the modulus comes from the tower (the free part) and the
    phase from the cube roots of unity (the torsion), and $\zeta_3$ has order exactly three, $\zeta_3^{3k}=1$; read at the place $2$, the modulus is $|2|_2$ and the phase the angular part of $\golden$ (\cref{prop:omega-two-adic}).
  \item The parameters of the ordinate formula are powers of the anchor:
    $\mu=\golden^{-\lng}$, $\sigma=3\golden^{-\lng}$, $\beta_0=\golden^{-3\lng}=1/\phi(20)$ and
    $\alpha_0=\golden^{-\lng}(3-\tfrac1{20})=\tfrac{59}{40}$, with $59=M_2N-1$, where $M_2=2^2-1=3$ is the Mersenne number of the anchor and $N=20$ the conductor; $\sigma/\mu=3$ and
    $\sigma+\mu=2$.
  \item The coefficient $c(n)=2+\alpha_0/(1+\beta_0\log n)$ exceeds $2$ for $n\ge1$.
  \item The determinant fixes the modulus: for the triad $\mu,-\mu/2,-\mu/2$ the product is
    $\mu^3/4$, and it equals $2^{-5}$ exactly when $\mu^3=\tfrac18$.
\end{enumerate}
\LeanInline{spectrum\_is\_anchor\_times\_orbit}, \LeanInline{omega\_order\_is\_exactly\_three},
\LeanInline{omegaZ\_kernel\_three}, \LeanInline{mu\_is\_phi\_to\_minus\_lambda},
\LeanInline{sigma\_is\_d\_times\_apex}, \LeanInline{beta\_is\_phi\_to\_minus\_three\_lambda},
\LeanInline{beta\_is\_inverse\_conductor\_units}, \LeanInline{alpha\_in\_phi},
\LeanInline{fiftynine\_is\_M2\_times\_conductor\_minus\_one}, \LeanInline{sigma\_over\_mu\_is\_the\_ternary},
\LeanInline{sigma\_plus\_mu\_is\_the\_diagonal\_trace}, \LeanInline{c\_excess}, \LeanInline{c\_above\_two},
\LeanInline{det\_fixes\_mu}, \LeanInline{alpha\_from\_arity\_and\_conductor} \tierP
\end{proposition}

\begin{proposition}[The spectral pair read from $G$]\label{prop:pair-from-G}
With $|\operatorname{im}G|=4$ and $|\ker G|=2$, computed from $G$ itself: the ratio of the
pair is $\sigma/\mu=|\operatorname{im}G|-1=3$, its sum is
$\sigma+\mu=|\operatorname{im}G|/|\ker G|=2$, its product is
$\sigma\mu=(|\operatorname{im}G|-1)/|\operatorname{im}G|=\tfrac34$, the apex is
$\mu=1/|\ker G|$, the product of the three cosines is $1/|\operatorname{im}G|$, and
$|\operatorname{im}G|\cdot|\ker G|=8=\phi(20)$.  The same constants are clauses of the
Mersenne bridge: $\sigma=M_2\golden^{-\lng}$, $\golden^{\lng}N=40$ and
$\alpha_0=(M_2N-1)/(\golden^{\lng}N)$, with $M_2=3$ and $N=20$.
\LeanStar{F1PaperNew.gCardinals\_of\_G}, \LeanStar{F1PaperNew.clauseFour\_holds} \tierPnew, which
close the hypotheses of \LeanInline{three\_readings\_are\_G\_invariants},
\LeanInline{mu\_is\_kernel\_reciprocal}, \LeanInline{cos\_product\_is\_image\_reciprocal},
\LeanInline{isomorphism\_theorem}, \LeanInline{sigma\_is\_M2\_over\_two},
\LeanInline{forty\_is\_clause\_four\_times\_conductor}, \LeanInline{alpha\_from\_clauses} \tierP
\end{proposition}

None of the parameters of the operator is fitted: each is a cardinal of $G$ or a clause of
the anchor.

\begin{proposition}[Fourier on the deck; the Fibonacci dynamics]\label{prop:mu6-fourier}
\begin{enumerate}[label=\textup{(\roman*)}]
  \item The deck of the fibre is the cyclic group $\mu_6$: the identification is a
    bijective homomorphism sending generator to generator, the deck acts transitively
    with order six, and $\mu_3$ embeds in it nontrivially.
  \item The six characters of $\mu_6$ with values in the Eisenstein integers are
    orthogonal and complete; analysis and synthesis are inverse up to the factor $6$;
    Parseval holds in sesquilinear form; convolution is diagonalised; and conjugation
    commutes with the transform.
  \item An isotypic component is invariant under the dynamics, and a component of type
    zero is a function of the level alone.
  \item The Fibonacci matrix $M$ has no periodic vector: $M^n v=v$ with $n>0$ forces
    $v=0$, because $\det(M^n-1)$ is the norm form at the Fibonacci coordinates and never
    vanishes.  Its powers carry the convergents, $M^3-1=2M$, and reduced modulo $2$ it
    has order exactly three while its determinant $-1$ gives the orientation of order two.
\end{enumerate}
\LeanInline{toMu6\_bijective}, \LeanInline{toMu6\_additive}, \LeanInline{generator\_to\_generator},
\LeanInline{deck\_transitive}, \LeanInline{mu6gen\_order\_six}, \LeanInline{mu3\_embeds\_nontrivially},
\LeanInline{orthogonality}, \LeanInline{completeness}, \LeanInline{inversion},
\LeanInline{parseval\_sesquilinear}, \LeanInline{convolution\_on\_the\_basis}, \LeanInline{conj\_analysis},
\LeanInline{isotypic\_invariant\_under\_dynamics}, \LeanInline{type\_zero\_is\_a\_level\_function},
\LeanInline{no\_periodic\_frequency}, \LeanInline{det\_npow\_sub\_one}, \LeanInline{det\_npow\_sub\_one\_ne\_zero},
\LeanInline{convergents\_are\_the\_entries\_of\_M\_pow}, \LeanInline{M\_cube\_sub\_one},
\LeanInline{colour\_factor\_has\_order\_three}, \LeanInline{orientation\_factor\_has\_order\_two} \tierP
\end{proposition}

The dynamics of \cref{prop:mu6-fourier}(iv) has two faces, one on the real torus and one at each prime, and the
torus carries both.

\begin{proposition}[The two faces of the Fibonacci dynamics]\label{prop:two-faces}
Let $M$ act on the torus $\mathbb T=\R^2/\Z^2$.
\begin{enumerate}[label=\textup{(\roman*)}]
  \item \emph{The periodic points.}  The points fixed by $M^n$ are finite in number, and their number is
    $|\det(M^n-1)|=|N(\golden^n-1)|=\prod_{d\mid n}|N(\Phi_d(\golden))|$.  Every periodic point is a torsion point
    of $\mathbb T$, and every torsion point is periodic.
  \item \emph{At each prime.}  For a prime $p\neq5$, $M$ acts on the $p$-torsion $\mathbb T[p]\cong(\Z/p)^2$ with
    period $e_p$, the least common multiple of the clocks of the primes of $\Zphi$ above $p$: $M^n=1$ there exactly
    when $F_n\equiv0$ and $F_{n-1}\equiv1\pmod p$, that is, when $\golden^n\equiv1$ (\cref{prop:binet}).  It divides
    $p-1$ if $p$ splits, and $2(p+1)$ but not $p+1$ if $p$ is inert~\cite{wall60}; at $p=5$ the period is $20$
    (\cref{prop:pisano}).  So $p$ divides the number of points fixed by $M^n$ exactly as
    \cref{prop:clock-cyclotomic}(ii) says: through the factors $\Phi_d(\golden)$ of its clocks.
  \item \emph{On the characters.}  For $m\in\Z^2$ let $e_m(x)=e^{2\pi i\,m\cdot x}$.  Then $e_m\circ M^k=e_{M^km}$, and
    for $m\neq0$ the vectors $M^km$ are pairwise distinct.  So for $m,n$ not both zero
    $\langle e_m\circ M^k,e_n\rangle$ is $1$ for at most one $k$ and $0$ for every other: the correlations of
    trigonometric polynomials vanish after finitely many steps, and $M$ is mixing.  The rate of expansion is
    $\log\golden$, the regulator.
\end{enumerate}
{\raggedright \LeanInline{DosCarasFibonacci.M\_pow}, \LeanInline{DosCarasFibonacci.period\_iff},
\LeanInline{DosCarasFibonacci.torsion\_is\_periodic}, \LeanInline{DosCarasFibonacci.periodic\_is\_torsion},
\LeanInline{DosCarasFibonacci.cassini}, \LeanInline{DosCarasFibonacci.det\_pow\_sub\_one\_ne\_zero},
\LeanInline{DosCarasFibonacci.no\_periodic\_vector}, \LeanInline{DosCarasFibonacci.orbit\_injective},
\LeanInline{PeriodicasSol.lefschetz\_counts}, \LeanInline{PeriodicasSol.wall\_samples},
\LeanInline{PeriodicasSol.ramified\_five}, \LeanInline{PeriodicasSol.split\_eleven},
\LeanInline{PeriodicasSol.inert\_three}, \LeanInline{PeriodicasSol.five\_tower} \tierP; in (i), that the kernel of an
endomorphism of the torus of nonzero determinant has $|\det|$ points \tierC, counted in Lean for $n<20$ and by
enumeration for $n\le8$; in (ii), the divisibility for every prime, Wall's law \tierC, in Lean on eleven primes and
by \texttt{retro109.py} for the primes below $400$ \tierN; in (iii), the passage from trigonometric polynomials to % pending(D7)
$L^2$, by density \tierC\par}
\end{proposition}

\begin{proof}
(i) $M^n-1$ is invertible over $\Q$ because $M$ has no periodic vector, so $x\mapsto(M^n-1)x$ is a surjective
endomorphism of $\mathbb T$ whose kernel has $|\det(M^n-1)|$ points; the eigenvalues of $M$ are $\golden$ and its
conjugate, so the determinant is $N(\golden^n-1)$, and \cref{prop:clock-cyclotomic}(i) factors it.  A fixed point of
$M^n$ solves $(M^n-1)x\in\Z^2$, so it is rational; and $M$ permutes the finite set $\mathbb T[N]$ for every $N$.
(ii) For $p\neq5$ the polynomial $x^2-x-1$ has distinct roots modulo $p$, so $M$ is diagonalisable over the residue
fields of the primes above $p$, with eigenvalues the reductions of $\golden$ and of its conjugate $-1/\golden$,
whose orders are the clocks; and $M^n=\bigl(\begin{smallmatrix}F_{n-1}&F_n\\F_n&F_{n+1}\end{smallmatrix}\bigr)$, which is the identity
exactly when $F_n\equiv0$ and $F_{n-1}\equiv1$.
(iii) $M$ is symmetric, so $e_m(Mx)=e_{Mm}(x)$; if $M^km=M^lm$ with $k>l$, then $M^{k-l}m=m$ and $m=0$.  The
characters are orthonormal.
\end{proof}

The two faces are the two readings of $\Phi_n(\golden)$ of \cref{ssec:why-phi}, and the product formula that glues
them is (i).  The characters of the torus, which are what the Fourier transform reads, see the dynamics only as the
decay of correlations of (iii), a statistical law with no period; the torsion points, which no character of the
level reads (\cref{prop:level-operators}), are where the same matrix is exactly periodic, prime by prime, with the
clock as period, the order of the angle of $\golden$ on the $p$-adic circle of the prisms $(\Z_p[\golden],(p))$
(\cref{ssec:delta-polar,conj:prismatic-route}).
The Frobenius and the $p$-adic angle, the cyclotomic polygon and the root of unity, are the two faces of one
dynamics; that the pairing of a real bulk with its $p$-adic boundaries, as in the arithmetic holography of
\cref{ssec:intro-binary}, has this form is a reading \tier{R}.

\Cref{prop:mu6-fourier} gives the transform matched to the deck.  What the operators of the tower read of the
fibre, and what they cannot, is the following; it is the reason the level operator of \cref{thm:operator-forms}(v)
is not a complete set of commuting observables.

\begin{proposition}[What an operator of the level does not separate]\label{prop:level-operators}
Let $\mathcal J$ and the deck step $D$ be those of \cref{def:joint-index}, and
let $H$ be any self-adjoint operator on $\ell^2(\mathcal J)$ that depends on the level alone, as the tower operator
$|\Oone|=\operatorname{diag}(m_0\golden^{\sigma})$ and the ordinate operator $\operatorname{diag}T^*(n)$ of
\cref{thm:operator-forms} do.
\begin{enumerate}[label=\textup{(\roman*)}]
  \item Every bounded Borel function of $H$ --- its spectral projections and its unitary group $e^{itH}$ among
    them --- depends on the level alone and commutes with $D$.  Every eigenspace of $H$ contains whole fibres, so no
    operator in the algebra generated by $H$ separates two points of a fibre.
  \item No relabelling of the level repairs this: $\operatorname{Hom}(\mu_6,\Z)=0$, so every homomorphism from the deck group to
    the level group is zero.
  \item Every unitary character of the fundamental group of the mapping torus $X$ of $M$ factors through
    $H_1(X)=\Z$: a decomposition in abelian characters, Fourier's among them, is indexed by the level alone.
  \item The colour operator $\OmH=\tfrac12\operatorname{diag}(1,\zeta_3,\zeta_3^2)$ and the orientation
    $W=\operatorname{diag}((-1)^{\epsilon})$ commute with each other and with $H$, and the three separate
    $\mathcal J$, while $H$ alone has eigenspaces of dimension six.  Neither commutes with $D$:
    $D\OmH D^{-1}=\zeta_3^{-1}\OmH$ and $DWD^{-1}=-W$.  So $Y=2\OmH W$ and $D$ satisfy $DY=\zeta_6\,YD$, with
    $\zeta_6=e^{i\pi/3}$, and the words $Y^aD^b$ span the full matrix algebra of the fibre.  The colour operator is
    normal and not self-adjoint.
\end{enumerate}
{\raggedright \LeanInline{PCFGranularity.deck\_invariant\_is\_level\_function}, \LeanInline{PCFGranularity.levelHom\_trivial},
\LeanInline{PCFGranularity.mu6\_is\_not\_absorbable\_into\_the\_level}, \LeanInline{PCFGranularity.H\_not\_complete},
\LeanInline{PCFGranularity.H\_degeneracy\_six}, \LeanInline{PCFGranularity.joint\_separates},
\LeanInline{PCFGranularity.ccso\_unconditional} \tierP; \LeanInline{ObservablesDelNivel.commute\_poly},
\LeanInline{ObservablesDelNivel.weyl\_not\_commute}, \LeanInline{ObservablesDelNivel.zeta3\_not\_real},
\LeanInline{FibraYReloj.deck\_weyl}, \LeanInline{FibraYReloj.yDiag\_injective}, \LeanInline{FibraYReloj.det\_M\_sub\_one} \tierP; (i) for Borel
functions, the spectral theorem \tierC; in (iii), that a unitary character factors through the abelianisation \tierC; (iv) also \texttt{retro104.py} \tierN\par} % pending(D7)
\end{proposition}

\begin{proof}
(i) $H$ is constant on each fibre, so every function of it is, and $D$, which preserves the level, commutes with it.
(ii) Six deck steps return to the origin, so a homomorphism $f$ to $\Z$ has $6f(1)=0$, and $\Z$ has no torsion.
(iii) $\pi_1(X)=\Z^2\rtimes_M\Z$, and its abelianisation is $\Z^2/(M-1)\Z^2\times\Z$; $\det(M-1)=-1$, so the first
factor vanishes, and every unitary character of a group factors through its abelianisation.
(iv) $D$ sends the point of colour $k$ to that of colour $k+1$, so conjugation by $D$ multiplies the eigenvalue
$\tfrac12\zeta_3^{k}$ attached to colour $k+1$ by $\zeta_3^{-1}$, and the orientation by $-1$; hence
$DYD^{-1}=-\zeta_3^{-1}Y=\zeta_6Y$.  A pair of unitaries of order six with this relation and $\zeta_6$ primitive is
the clock and the shift of $\Z/6$, whose words span $M_6(\C)$.  The eigenvalues $\tfrac12\zeta_3^{\pm1}$ are not
real.
\end{proof}

\begin{remark}[Why one self-adjoint operator is not enough here]\label{rmk:one-operator}
A self-adjoint operator generates a commutative $C^*$-algebra, isomorphic to the continuous functions on its
spectrum \cite{gelfandnaimark43}, and its unitary group $e^{itH}$ is a one-parameter group of which it is the
generator \cite{stone32}: its spectral decomposition is the decomposition of that abelian action in characters.
P\'olya's conjecture asks for the zeros as the eigenvalues of one such operator (\cref{sec:polya-hist}).  In the
tower, \cref{prop:level-operators} says what an operator of the level leaves out: it is constant on the fibres
together with its whole functional calculus, and what separates them does not commute with the
deck step: $Y$ alone separates them, its six eigenvalues are the sixth roots of unity, and $DY=\zeta_6YD$, the case $\theta=\tfrac16$ of the relation $U_2U_1=e^{2\pi i\theta}U_1U_2$ between two unitaries that defines the
noncommutative tori $T_\theta$ of Manin, which take $\theta$ irrational (\cref{sec:manin}; \cite{manin04}).  The
quantised cat map shows the same in its propagator: its eigenvalues are degenerate, with the multiplicity the index
gives (\cref{prop:index-degeneracy}), and the Hecke operators, which commute with it, are what fixes an eigenbasis
\cite{kurlbergrudnick00}.  The operator $\Oone$ of \cref{thm:one-operator} joins the two: its modulus is self-adjoint, a
level operator, and its phase is $Y$, what a level operator leaves to be read.
\end{remark}

An operator of the level shows its structure as degeneracy.  What a symmetry of a self-adjoint operator shows in its
spectrum in general is the following, and two of its cases are the ones the zeros can be tested against.

\begin{proposition}[What a symmetry shows in a spectrum]\label{prop:symmetry-spectrum}
Let $H$ be a self-adjoint operator on a finite-dimensional complex Hilbert space.
\begin{enumerate}[label=\textup{(\roman*)}]
  \item If a unitary $D$ commutes with $H$, then $H$ preserves every eigenspace of $D$, and the spectrum of $H$ is the
    union of the spectra of its restrictions to them, one sector for each eigenvalue of $D$.  The commutant of the deck
    step of \cref{prop:level-operators} has dimension six, one sector for each character of $\mu_6$; the commutant of
    the pair $D,Y$ of \cref{prop:level-operators}\textup{(iv)} is the scalars, so that pair imposes no sector.
  \item If an antiunitary $T=UK$, with $K$ the complex conjugation and $U$ unitary, commutes with $H$ and $T^2=-1$,
    then $T$ fixes no ray, $\langle v,Tv\rangle=0$ for every $v$, and every eigenvalue of $H$ has even multiplicity.
    If $T^2=+1$, $T$ fixes the real vectors of a basis and forces no multiplicity.
  \item An operator of the level has every eigenvalue of multiplicity at least six.
\end{enumerate}
{\raggedright \LeanInline{EstructuraEnGUE.sector\_invariant}, \LeanInline{EstructuraEnGUE.skew\_quadratic\_zero},
\LeanInline{EstructuraEnGUE.antiunitary\_sq\_neg\_one\_no\_fixed\_ray} \tierP; the commutants of (i) and the doublets of
(ii) also \texttt{retro105.py} \tierN; the classes of (ii) are Dyson's \cite{dyson62threefold}\par} % pending(D7)
\end{proposition}

\begin{proof}
(i) $D(Hv)=HDv=\lambda Hv$ when $Dv=\lambda v$, and the eigenspaces of the unitary $D$ are orthogonal and span the space.
$D$ has the six sixth roots of unity as simple eigenvalues, so its commutant is the diagonal algebra in its eigenbasis,
of dimension six; $D$ and $Y$ generate the full matrix algebra of the fibre (\cref{prop:level-operators}), whose
commutant is the scalars.
(ii) From $T^2=UK UK=U\bar U=-1$ and $U$ unitary, $U^{\mathsf T}=-U$, so $\langle v,Tv\rangle=v^{*}U\bar v$ is the
quadratic form of an antisymmetric matrix at $\bar v$ and vanishes.  If $Tv=cv$ then $-v=T^2v=|c|^2v$, so $v=0$: no
ray is fixed.  $T$ commutes with $H$ and the eigenvalues are real, so $T$ preserves each eigenspace $E$; for a unit
$v\in E$, $v$ and $Tv$ are orthonormal in $E$, and the orthogonal complement of their span in $E$ is again
$T$-invariant, so $\dim E$ is even by induction.  If $T^2=+1$, $T$ is the conjugation in a basis of fixed vectors.
(iii) is \cref{prop:level-operators}\textup{(i)}: every eigenspace contains whole fibres, of six points.
\end{proof}

\subsubsection{One space for all the lifts}\label{ssec:primon-space}

\begin{proposition}[One space for all the Frobenius]\label{prop:primon-space}
On the algebraic core $\mathcal{D}_0\subset\ell^2(\mathbb{N})$ of the operator
$\tilde H_{\mathrm{PCF}}$ \cite{gonzalez25hp} --- the finite linear combinations of the basis vectors $e_n$ --- the multiplicative shifts
$U_p\colon e_n\mapsto e_{pn}$ satisfy $U_p\circ U_q=U_{pq}$ and commute, the law of
\cref{thm:lambda-monoid}; the twist $e_n\mapsto\chiv(n)e_n$ is covariant,
$\chiv(pn)=\chiv(p)\chiv(n)$; and for $n$ prime to $5$ the twist is the Frobenius,
$\Frob_n(\sqrt5)=\chiv(n)\sqrt5$.  The truncated twisted trace
$\sum_{n<N}\chiv(n)n^{-s}$ is the partial sum of $L(s,\chiv)$.
\LeanStar{U\_comp}, \LeanStar{U\_commute}, \LeanStar{covariance}, \LeanStar{twist\_is\_frobenius},
\LeanStar{twistedTrace} \tierPnew
\end{proposition}

\begin{remark}[What the space gives and what it does not]\label{rmk:space-gives}
Let $\hat N e_n=n\,e_n$ be the number operator.  For $\mathrm{Re}\,s>1$, $\mathrm{Tr}(\hat N^{-s}\otimes\Frob_{\hat N}\mid\ell^2(\mathbb{N})\otimes\mathbb{Q}\sqrt5)
=L(s,\chiv)$, with $\Frob_{\hat N}$ acting on $e_n\otimes\sqrt5$ as $\Frob_n$: the primon gas, the free gas labelled by the primes whose partition function is $\zeta$ \cite{julia90} (\cref{sec:primon}), twisted by the Frobenius of
$\Zphi$, and the Bost--Connes form of the Euler product~\cite{bostconnes95}.  The operator $\tilde H_{\mathrm{PCF}}=\mathrm{diag}(T^*(n))$ lives on the
same space and is self-adjoint by the standard argument for real diagonal
operators \cite{gonzalez25hp}; its Lean backing is $T^*(n)\in\mathbb{R}$ and
$T^*(n)\to\infty$, and its eigenvalues follow the zero heights with a relative error below
$1.7\%$ in the regime of \cite{gonzalez25hp}, up to $n=10^{12}$,
not exactly.  The zeros live in the analytic continuation, not in
$\mathrm{Re}\,s>1$: passing from this trace to a positive form whose spectral
side is the zeros is Connes' step to the ad\`ele class space, and its positivity
is equivalent to GRH$(\chiv)$ \cite{connes99}.  That step is the core.
\end{remark}

% ---------------------------------------------------------------------
\subsection{The structure that places the zeros}\label{sec:machinery}

This section collects the structure that places the zeros against the operator, in the order in which it is used: first
what produces the line and places each zero against it, then what carries
the operator's integrality to the zeros.  The thread of \cref{sec:topological-proof} uses the first as its Stage~6 and the second as its Stage~7.

\subsubsection{Setting}\label{ssec:topo-setting}

A zero is a structure, not a hypothesis: \LeanInline{Chi5Zero} carries $\rho$,
the vanishing $L(\rho,\chiv)=0$ and the two strip bounds $0<\Re\rho<1$, so those
bounds are available in every statement below without being assumed again.  We
write $\beta=\Re\rho$, $\gamma=\Im\rho$, and
\[
  \Ecore\colon\quad \forall z\in\text{\LeanInline{Chi5Zero}},\ \Re z.\rho=\tfrac12 .
\]

\subsubsection{The analytic inputs, discharged}\label{sec:analytic-inputs}

The inputs about $L(s,\chiv)$ used throughout are listed below: every row but the
last is a theorem; the last, the height of the first zero, is classical.

\begin{longtable}{p{4.6cm}p{6.6cm}l}
\toprule Input & Lean & Tier\\\midrule
Functional equation of $\Lambda(s,\chiv)$ & \LeanInline{functionalEquation\_chi5\_discharged} & \tierP\\
Companion zero (Hecke 1920) & \LeanInline{hecke\_chi5\_derived}, \LeanInline{zeros\_route\_is\_discharged} & \tierP\\
$\overline{L(\bar s)}=L(s)$ & \LeanStar{hconj\_discharged} & \tierPnew\\
Closure under the mate; real coefficients & \LeanStar{closedUnderMate\_chi5}, \LeanStar{realCoefficients\_chi5} & \tierPnew\\
$|\gamma|\ge2$ for the zeros used & \LeanInline{discBound\_discharged\_high} & \tierP\\
\texttt{hmate}: each zero has a zero $\rho'$ with $|1-1/\rho'|=|1-1/\rho|^{-1}$; \texttt{hcoord}: for $\Re\rho>0$, $|1-1/\rho|=1\iff\Re\rho=\tfrac12$ & \LeanInline{mate\_hypothesis\_discharged}, \LeanInline{hcoord\_discharged} & \tierP\\
First zero at $\gamma\approx6.6485$ (\LeanInline{NoLowZeros}) & classical; not proved in Lean & \tierC\\
\bottomrule\end{longtable}
The root number of the functional equation is $+1$: the Gauss sum of $\chiv$ is $\tau(\chiv)=+\sqrt5$,
by Gauss's determination of the sign of the quadratic Gauss sum (\cref{rmk:nature-L}) \tierC.

The companion deserves a word, because it is what earlier drafts took as an
axiom.  From the functional equation, if $\rho$ is a zero of $L(s,\chiv)$ in the
strip then so is $1-\rho$; the corpus derives this and attaches it to the
structure, so that \LeanInline{Chi5Zero.companion} is a function and not a
hypothesis.  Its real part is $1-\beta$
(\LeanInline{Chi5Zero.companion\_re}), which is the only property the squeeze
uses.  Only \LeanInline{NoLowZeros} remains as a classical input, and it enters
just once, in \cref{thm:reduction}.

\subsubsection{The Klein group and its orbits}\label{ssec:klein}

Let $\sigma(s)=1-s$ (the functional equation) and $\tau(s)=\bar s$ (conjugation).
Both are involutions and they commute, so they generate a Klein four-group acting
on $\C$ (\LeanInline{klein\_group\_relations} \tierP).  The composite
$\sigma\tau(s)=1-\bar s$ is the map that matters.

\begin{proposition}[The critical line is a fixed set]\label{prop:stau}
$1-\bar s=s$ if and only if $\Re s=\tfrac12$.
\LeanInline{stau\_fixed\_iff\_critical} \tierP
\end{proposition}
\begin{proof}
$1-\bar s=(1-\Re s)+i\,\Im s$, whose imaginary part always agrees with that of
$s$; the real parts agree exactly when $\Re s=1-\Re s$.
\end{proof}

So the critical line is not a locus defined by an inequality but the fixed set of
an involution --- an antiholomorphic one, which is why it is a line and not a
point: the holomorphic reflection $\sigma$ alone fixes only $s=\tfrac12$.  The
zero set of $L(\cdot,\chiv)$ is stable under the group, by the two symmetries of
\cref{sec:analytic-inputs}, and the orbits are of two kinds.

\begin{proposition}[Orbit dichotomy]\label{prop:orbits}
On the line, the orbit of a zero degenerates to the pair $\{\rho,\bar\rho\}$; off
it, the orbit is free of size four.  Consequently
\[
  \Ecore\iff\text{the Klein action on the zero set has no free orbit,}
\]
\LeanInline{orbit\_on\_line\_degenerate}, \LeanInline{orbit\_off\_line\_free},
\LeanInline{E\_iff\_no\_free\_orbit} \tierP.
\end{proposition}

This is the first reformulation: \Ecore\ has become a statement about the absence of a configuration. A second form of the same reformulation removes the group and keeps
only the mate $\rho\mapsto1-\bar\rho$.

\begin{theorem}[Shared ordinates]\label{thm:shared-ordinate}
For a set $Z\subset\C$ closed under the mate, every point of $Z$ lies on the line
if and only if no two distinct points of $Z$ share an ordinate.  For the zeros of
$L(s,\chiv)$ the hypothesis is discharged
(\LeanStar{closedUnderMate\_chi5}), so the equivalence is unconditional.
\LeanInline{E\_iff\_no\_shared\_ordinate}, \LeanStar{E\_iff\_no\_shared\_ordinate\_chi5} \tierPnew
\end{theorem}
\begin{proof}
If $\Re\rho\ne\tfrac12$ then the mate $1-\bar\rho$ is in $Z$, is distinct from
$\rho$ by \cref{prop:stau}, and has the same imaginary part; that is a shared
ordinate.  Conversely, two distinct zeros on the line with the same ordinate
would be equal, since their real parts agree.
\end{proof}

\begin{remark}[What was discarded here]
An earlier version of the programme used a measured lower bound on the spacing of
ordinates (\LeanInline{MinSpacing}) as an input.  It is a theorem on finite
windows and its global form is refuted by the very configuration it was meant to
exclude (\LeanInline{no\_uniform\_m\_of\_vanishing\_spacing} \tierP), so it is
used nowhere.  \Cref{thm:shared-ordinate} needs no metric bound: only that no two
zeros coincide in ordinate, which is strictly weaker than repulsion.
\end{remark}

\begin{proposition}[The Klein action on the zeros, in detail]\label{prop:klein-detail}
\begin{enumerate}[label=\textup{(\roman*)}]
  \item Conjugation and $\rho\mapsto1-\rho$ are involutions, and the mate
    $\rho\mapsto1-\bar\rho$ is their composite: $1-\bar\rho=\overline{1-\rho}$.
  \item The mate keeps the ordinate, and it fixes $\rho$ exactly when $\Re\rho=\tfrac12$;
    $\rho\mapsto1-\rho$ does not keep the ordinate when $\gamma\ne0$, so it cannot play the
    role of the mate.  Conjugation fixes exactly the real points, and $\rho\mapsto1-\rho$
    fixes only $\rho=\tfrac12$.
  \item For any function with the functional-equation symmetry and real coefficients,
    the zero set in the strip is closed under the mate; for $L(s,\chiv)$ both inputs
    hold, because the values of $\chiv$ and the terms of its series are fixed by
    conjugation.
  \item The orbit of a zero off the real axis has two elements exactly on the line;
    the centre of the orbit is the same for all zeros.
  \item The spectral positions are the roots of $8x^2-2x-1$, namely $\tfrac12$ and
    $-\tfrac14$, and two distinct positions are at least $\tfrac34$ apart.
\end{enumerate}
\LeanInline{klein\_acts\_on\_the\_zeros\_by\_two\_involutions}, \LeanInline{kmate\_is\_conj\_of\_companion},
\LeanInline{kmate\_keeps\_the\_ordinate}, \LeanInline{kmate\_fixed\_iff},
\LeanInline{companion\_cannot\_replace\_kmate}, \LeanInline{sigma\_fixes\_iff\_real},
\LeanInline{tau\_fixes\_iff\_half}, \LeanInline{closed\_under\_mate\_of\_fe\_rc},
\LeanInline{chi5\_values\_are\_conj\_invariant}, \LeanInline{the\_series\_terms\_are\_conj\_stable},
\LeanInline{the\_cardinal\_is\_the\_spectral\_position}, \LeanInline{the\_centre\_is\_the\_same\_for\_all\_zeros},
\LeanInline{mem\_target\_iff\_spectralPoly\_zero}, \LeanInline{minSpacing\_on\_the\_spectrum} \tierP
\end{proposition}

\begin{remark}[The mate does not use the reality of the character]\label{rmk:mate-without-G}
\Cref{prop:klein-detail}\textup{(iii)} asks for more than the mate needs.  For any primitive character
$\chi$ the functional equation carries the zeros of $L(s,\chi)$ to those of $L(s,\bar\chi)$ by
$\rho\mapsto1-\rho$, and $L(\bar s,\bar\chi)=\overline{L(s,\chi)}$ carries them back by conjugation; the
composite is the mate, so the zero set of every $L(s,\chi)$ is closed under $\rho\mapsto1-\bar\rho$
(\LeanInline{MateSinG.closed\_under\_mate\_of\_fe\_dual} \tierP; the functional equation \tierC).  The reality
of $\chi$ --- modulo $20$, that $\chi$ factors through $G$, whose kernel is the set of squares
(\cref{prop:G}) --- adds the other two involutions, and by \textup{(ii)} neither of them fixes a zero with
$\gamma\ne0$ other than through the mate.  So the orbit dichotomy in its mate-only form
(\cref{thm:shared-ordinate}) is the same statement for $\chiv$ and for the character $\omega$ of order
$4$ modulo $5$, with $\omega^2=\chiv$, which does not factor through $G$: what $G$ contributes to the
zeros is the symmetry, and the position on the line is the part of the statement that $G$ does not see.
\end{remark}

\begin{proposition}[The Klein action on the zeros is induced by $\chiv$]\label{prop:descent}
The Klein group $V=\{\pm1\}\times\{a\mapsto a^{-1}\}$ acts on $\Zst$ with the
degenerate orbits $\{1,19\},\{9,11\}$ --- the classes with $a^2=1$, all split,
$\chiv=+1$ --- and the free orbit $\{3,7,13,17\}$, exactly the inert classes,
$\chiv=-1$.  The Klein group acting on the zeros is induced by $V$ through $\chiv$:
\begin{enumerate}[label=\textup{(\roman*)}]
  \item inversion corresponds to conjugation, because $\chiv(a^{-1})=\chiv(a)$ ---
    $\chiv$ is real --- and the reality of $\chiv$ is what closes the zero set under
    $\rho\mapsto\bar\rho$;
  \item negation corresponds to $\rho\mapsto1-\rho$, because $\chiv(-a)=\chiv(a)$ ---
    $\chiv$ is even --- and the parity selects $\Gamma(s/2)$ and with it the
    functional equation that closes the zero set under $\rho\mapsto1-\rho$.
\end{enumerate}
Both generators are the two involutions of one character.  The group $V$ they generate
is a Klein four-group, isomorphic to $\operatorname{im}G\cong\mathbb Z_2\times\mathbb Z_2$;
the paper writes $V=\operatorname{im}G$ for this isomorphism.  Downstairs the orbit of a zero has size $2$ exactly on the
line and $4$ off it.
\LeanInline{both\_actions\_from\_one\_character}, \LeanInline{inversion\_gives\_sigma},
\LeanInline{negation\_gives\_tau}, \LeanInline{free\_orbit\_is\_inert},
\LeanInline{klein\_orbits\_on\_Z20\_are\_two\_and\_four},
\LeanInline{stau\_fixed\_iff\_critical} \tierP
\end{proposition}

\begin{remark}[The information of the chain]\label{rmk:descent-info}
That the chain from generator to zeros is injective --- distinct generators giving
distinct zero sets, $\golden\mapsto\gamma_1=6.6484$, $1+\sqrt2\mapsto\gamma_1=4.8999$ ---
is a numerical observation (\LeanInline{descent\_loses\_nothing} \tierP\ given the
injectivity as input \tierN), not part of \cref{prop:descent}.
\end{remark}

\subsubsection{The first contact: the orbit centre}\label{ssec:orbit-centre}

The first contact with the zeros is unconditional and costs nothing.

\begin{proposition}[Orbit embedding]\label{prop:orbit-embedding}
For every $\rho\in\C$, the centre of the pair $\{\rho,1-\rho\}$ is $\tfrac12$,
which is the positive eigenvalue of $\OmH$:
$\Re\,\mathrm{ctr}(\rho,1-\rho)=\Re\OmH_0=\tfrac12$.
\LeanInline{the\_orbit\_embedding\_is\_unconditional} \tierP
\end{proposition}

Every zero therefore sits in the spectrum \emph{through its orbit}: the centre is
always $\tfrac12$.  What is not known is the radius.

\subsubsection{The apex: cocone and colimit}\label{ssec:apex}

The value $\tfrac12$ at which every orbit is centred is not one number among others:
it is the apex of the construction, reached by independent routes and characterised
as a colimit.

\begin{proposition}[Cocone and colimit]\label{prop:cocone}
\begin{enumerate}[label=\textup{(\roman*)}]
  \item Six routes from distinct domains reach $\mu_3=\tfrac12$: the half-factorial
    $\Gamma(\tfrac32)/\Gamma(\tfrac12)$, the norm $\|\OmH_0\|$, $\cos(\pi/3)$, the
    inverse binary turn $\golden^{-\lng}$, the fixed point of $x\mapsto1-x$, and the
    common real part of $s$ and $1-s$ on the critical line.
  \item The critical line is the locus where $s$ and $1-s$ have the same real part:
    $\Re s=\Re(1-s)$ exactly when $\Re s=\tfrac12$; it is the fixed set of $s\mapsto1-\bar s$
    (\cref{prop:stau}), while $s\mapsto1-s$ alone fixes only $s=\tfrac12$.
  \item Any datum that every route projects to is the apex.
\end{enumerate}
\LeanInline{modulus\_cocone\_six\_routes}, \LeanInline{functional\_equation\_fixed\_line},
\LeanInline{apex\_colimit\_universal} \tierP
\end{proposition}

The route through the reflection says where the apex lives.  In every commutative ring in which
$2$ is invertible, $x\mapsto1-x$ has exactly one fixed point, $\tfrac12$; on $\mathbb{Z}$ it has
none, and none on $\Zphi$ either, since $2$ is not a unit there
(\LeanInline{F1Lawvere.reflection\_fixed\_iff}, \LeanInline{F1Lawvere.reflection\_fixedPointFree\_int},
\LeanInline{F1Lawvere.reflection\_fixedPointFree\_Z\ensuremath{\omega}} \tierP).  So the apex is
not an element of the ring the construction starts from.  It is reached in the golden tower, as
$\golden^{-\lng}$, through the anchor $\golden^{\lng}=2$ (\cref{sec:mersenne}).

The orbit of every zero lands on the apex (\cref{prop:orbit-embedding}); the
subsections that follow measure how far each zero itself can be from it.

\subsubsection{The Apollonius coordinate}\label{ssec:apollonius}

Write
\[
  r(\rho)=\frac{|\rho-1|}{|\rho|},\qquad t(\rho)=r(\rho)+\frac1{r(\rho)} .
\]
$r$ is the Apollonius coordinate of the pair $\{0,1\}$: the locus $r=1$ is the
perpendicular bisector of the segment, that is, the critical line
(\LeanInline{line\_is\_the\_unit\_circle}, \LeanInline{apollonius\_closes\_in\_lean} \tierP),
and the deviation $\alpha=(1-2\beta)/(\beta^2+\gamma^2)$ vanishes exactly there
(\LeanInline{alpha\_zero\_iff\_r\_one} \tierP).  Two facts about $r$ are proved
without any input beyond the strip.

\begin{proposition}\label{prop:rbound}
For $0<\beta<1$ and $\gamma\ne0$, $r(\rho)^2\le1+1/\gamma^2$.  With $|\gamma|\ge2$ this
gives $t(\rho)\le2.0125$; with $|\gamma|\ge6.6485$, the height of the first zero, it
gives $t(\rho)\le2.00013$.
\LeanInline{rSq\_le\_in\_strip} \tierP; the height of the first zero is the classical
input \LeanInline{NoLowZeros} \tierC
\end{proposition}

\begin{proposition}\label{prop:integrality-is-E}
For a zero in the strip with $|\gamma|\ge2$,
\[
  \bigl(\exists n\in\N,\ t(\rho)=n\bigr)\iff\Re\rho=\tfrac12 .
\]
\LeanInline{the\_two\_labels\_are\_one\_gap} \tierP
\end{proposition}
\begin{proof}
$t=r+1/r\ge2$ with equality exactly at $r=1$, and $r=1$ is the line by the
Apollonius description; with $|\gamma|\ge2$, \cref{prop:rbound} places $t$ in
$[2,2.0125]$, an interval containing no integer other than $2$.
\end{proof}

\Cref{prop:integrality-is-E} is the pivot of the whole construction.  It says
that \Ecore\ is exactly the \emph{integrality} of one real number attached to
each zero --- not its size, which the strip already bounds, but its being a whole
number.  Bounding and quantising are different operations, and the corpus keeps
them apart explicitly (\LeanInline{bounding\_is\_not\_quantising} \tierP).

\subsubsection{The zero as centre plus radius; its coordinate, discriminant and valuation}\label{ssec:coordinate}

Everything in this subsection holds for every zero, and none of it uses that the zero
is on the line.

\begin{proposition}[Centre plus radius]\label{prop:centre-radius}
For every $\rho\in\C$, $\rho=\mathrm{ctr}(\rho,1-\rho)+\mathrm{rad}(\rho,1-\rho)$ with centre
$\tfrac12$ and radius $\rho-\tfrac12$: the real part of the radius is $\beta-\tfrac12$ and
its imaginary part is the whole ordinate $\gamma$.  In the strip $|\Re\,\mathrm{rad}|<\tfrac12$;
the real part of the radius vanishes exactly on the line; and the radius that the
second spectral position $-\tfrac14$ would need, $-\tfrac34$, never fits.  The
symmetric product of the pair is $\rho(1-\rho)=\tfrac14-(\rho-\tfrac12)^2$; for $\gamma\ne0$
it is real exactly on the line.  And the strip itself is the condition $|1-2\beta|<1$.
\LeanInline{embedding\_is\_centre\_plus\_radius}, \LeanInline{radius\_is\_rho\_minus\_half},
\LeanInline{the\_radius\_carries\_the\_whole\_ordinate}, \LeanInline{radius\_re\_bounded\_by\_the\_apex},
\LeanInline{orbit\_radius\_vanishes\_iff\_E}, \LeanInline{the\_second\_radius\_does\_not\_fit},
\LeanInline{e\_two\_from\_the\_centre}, \LeanInline{im\_zero\_iff\_product\_zero},
\LeanInline{confinement\_is\_the\_strip} \tierP
\end{proposition}

\begin{proposition}[The Apollonius coordinate]\label{prop:apollonius-coordinate}
For $\rho\ne0$ the coordinate $r(\rho)=|1-1/\rho|$ is the ratio of distances
$|\rho-1|/|\rho|$, and its level sets are the Apollonius circles of the foci $0$ and $1$.
\begin{enumerate}[label=\textup{(\roman*)}]
  \item The Klein group acts on it: $r(\bar\rho)=r(\rho)$, $r(1-\rho)=r(\rho)^{-1}$ and
    $r(1-\bar\rho)=r(\rho)^{-1}$; the discriminant is invariant under the mate.
  \item $r(\rho)\le1+1/|\gamma|$ for $\gamma\ne0$.
  \item Its discriminant $\mathrm{disc}(r)=(r+1/r)^2-4=(r-1/r)^2$ has the closed form
    $\mathrm{disc}(r(\rho))=(1-2\beta)^2/(|\rho|^2|\rho-1|^2)\le(1-2\beta)^2/\gamma^4$,
    vanishes exactly on the line, and satisfies $\mathrm{disc}(r(\rho))\le\tfrac1{12}$ for
    every zero with $|\gamma|\ge2$.
  \item The coordinate reads the side of the line: $r(\rho)>1$ exactly when $\beta<\tfrac12$.
    Its defect is $r^2-1=(1-2\beta)/|\rho|^2$, so $|r^2-1|<1/\gamma^2$ in the strip, at most
    $\tfrac14$ for $|\gamma|\ge2$, and $(r^2-1)\gamma^2<1$: the defect times the square of
    the height stays below $1$.  For a real point of the strip, $r(\beta)=1/\beta-1$.
  \item The numerator of the trace is a quadratic form centred at $\tfrac12$:
    $|\rho-1|^2+|\rho|^2=2\bigl((\beta-\tfrac12)^2+\gamma^2\bigr)+\tfrac12$, minimal exactly
    on the line.
\end{enumerate}
\LeanInline{zeroCoord\_is\_a\_distance\_ratio}, \LeanInline{alpha\_form\_eq\_recentred},
\LeanInline{numerator\_minimal\_on\_the\_line}, \LeanInline{equator\_is\_the\_critical\_line}, \LeanInline{the\_apollonius\_circle\_is\_the\_level\_set},
\LeanInline{the\_klein\_action\_on\_the\_coordinate}, \LeanInline{disc\_invariant\_under\_kmate},
\LeanInline{coord\_le\_one\_add\_inv\_im}, \LeanInline{coord\_gt\_one\_iff\_re\_lt\_half},
\LeanInline{liMap\_sq\_defect}, \LeanInline{defect\_lt\_inv\_gamma\_sq},
\LeanInline{noLowZeros\_forces\_a\_small\_defect}, \LeanInline{height\_bounded\_by\_the\_defect},
\LeanInline{real\_coord\_closed\_form}, \LeanInline{disc\_square\_196}, \LeanInline{disc\_closed\_form},
\LeanInline{disc\_decays\_with\_height}, \LeanInline{disc\_zero\_iff\_half},
\LeanInline{discBound\_of\_high\_gamma} \tierP
\end{proposition}

\begin{proposition}[The tower seen by the discriminant and the valuation]\label{prop:tower-disc-val}
Let $v(x)=\log x/\log\golden$ be the golden valuation.
\begin{enumerate}[label=\textup{(\roman*)}]
  \item On the tower: $v(\golden^n)=n$, $v(x^p)=p\,v(x)$ (the Frobenius lift is
    multiplication by $p$ on valuations) and $v(1/x)=-v(x)$.
  \item $x>0$ lies in the tower exactly when $v(x)\in\Z$.
  \item The discriminant of a level is a norm form: $\mathrm{disc}(\golden^k)=5F_k^2$ for
    $k$ even and $L_k^2$ for $k$ odd, with $L_k^2-5F_k^2=4(-1)^k$; so
    $\mathrm{disc}(\golden^k)\ge1$ for every $k\ne0$, and $\mathrm{disc}(x)<1$ exactly
    when $\golden^{-1}<x<\golden$.  The trace of $\golden^k$ is the Lucas number $L_k$ for
    $k$ even and $\sqrt5F_k$ for $k$ odd.
  \item For every zero, $v(r(1-\rho))=-v(r(\rho))$; and for every zero with
    $|\gamma|\ge2$, $|v(r(\rho))|\le\log(2/\sqrt3)/\log\golden<0.3$.
  \item The three eigenvalues of $\OmH$ share one valuation, that of the anchor:
    $|\OmH_k|=\golden^{-\lng}$.
\end{enumerate}
\LeanInline{valPhi\_phi\_pow}, \LeanInline{valPhi\_is\_the\_bridge\_to\_psi},
\LeanInline{valPhi\_inversion\_is\_negation}, \LeanInline{monomial\_iff\_valPhi\_int},
\LeanInline{disc\_family\_is\_norm\_form}, \LeanInline{lucas\_sq\_sub\_five\_fib\_sq},
\LeanInline{family\_disc\_ge\_one\_of\_ne\_zero}, \LeanInline{disc\_lt\_one\_iff\_phi\_window},
\LeanInline{trace\_integral\_iff\_even}, \LeanInline{the\_zero\_pair\_in\_valPhi\_has\_opposite\_signs},
\LeanInline{unconditional\_valPhi\_bound}, \LeanInline{the\_spectrum\_has\_one\_valuation} \tierP
\end{proposition}

\begin{proposition}[The Li map and its fixed point]\label{prop:li-map}
The map $\ell(\rho)=1-1/\rho$ whose modulus is the coordinate has order three:
$\ell(\ell(\ell(\rho)))=\rho$ for $\rho\ne0,1$, with $\ell(\ell(\rho))=1/(1-\rho)$ and
$\ell(\rho)\,\ell(1-\rho)=1$.  Its fixed point is $e^{i\pi/3}$, the root of
$x^2-x+1$ with real part $\tfrac12$ and modulus $1$: a primitive sixth root of unity, on
the line and on the unit circle at once.  For a finite set of zeros the first Li
coefficient is $\lambda_1=\sum_\rho1/\rho$, and $\lambda_0=0$.  On the unit circle the Li
term of a point of modulus $1$ and angle $\theta$ is $4-4\cos n\theta=8\sin^2(n\theta/2)\ge0$:
on the circle, the image of the line, every such term is nonnegative.  Off the circle
the powers escape: if $r(\rho)>1$, then $|\ell(\rho)^n|$ exceeds every bound.
\LeanInline{li\_coordinate\_reads\_E}, \LeanInline{liFix\_root}, \LeanInline{liFix\_re\_is\_half},
\LeanInline{liFix\_abs\_one}, \LeanInline{li\_fixed\_point\_has\_order\_six},
\LeanInline{li\_one\_is\_reciprocal\_sum}, \LeanInline{li\_zero\_vanishes},
\LeanInline{term\_on\_circle}, \LeanInline{two\_minus\_two\_cos\_is\_square},
\LeanInline{term\_nonneg\_on\_circle}, \LeanInline{liMap\_sq}, \LeanInline{liMap\_pair\_reciprocal},
\LeanInline{off\_the\_line\_powers\_escape} \tierP
\end{proposition}

The Li map thus carries the three structures of the paper in one object: it has the
order of the ternary face, its fixed point generates the deck group $\mu_6$ of
\cref{prop:mu6-fourier}, and that point lies on the line.

\begin{remark}[Confined, and what membership adds]\label{rmk:confined}
Read together, the three propositions place every zero with $|\gamma|\ge2$ exactly.  Its
discriminant is at most $\tfrac1{12}$, while every level of the tower other than
$\golden^0$ has discriminant at least $1$.  Its valuation lies in $(-0.3,0.3)$, while the
valuations of the tower are the integers.  So the coordinate is confined to the one
stretch where the only level of the tower is $1$.  Membership adds exactly one thing:
that the valuation is an integer.  Given that, it is $0$, the coordinate is $1$, and the
zero is on the line (\LeanInline{integrality\_of\_valuation\_is\_E},
\LeanInline{monomial\_alone\_gives\_E} \tierP).  Confinement alone does not give it: the
number $\tfrac54$ has discriminant $\tfrac{81}{400}<1$, lies in the stretch, and is not a
power of $\golden$ (\LeanInline{witness\_T3}, \LeanInline{witness\_not\_in\_family} \tierP).
\end{remark}

\subsubsection{The trace spectrum is discrete, and the bound selects one value}\label{ssec:trace-spectrum}

The operator side supplies not merely an integer but a discrete set of values,
and the strip bound of \cref{prop:rbound} lands inside its first gap.

\begin{proposition}\label{prop:lucas-spectrum}
For even index the trace of the $k$-th power of the Fibonacci matrix is the Lucas
number $L_k$: the structural spectrum is $\{L_{2m}\}=\{2,3,7,18,47,\dots\}$, the
first jump is $L_2-L_0=1$, and the only value below $3$ is $L_0=2$.
\LeanInline{even\_index\_gives\_integer\_trace}, \LeanInline{target\_set\_is\_derived},
\LeanInline{spectrum\_jumps}, \LeanInline{bound\_selects\_one} \tierP\ for the first five values;
for every $m$, \LeanStar{F1PaperNew.even\_lucas\_below\_three\_iff} \tierPnew
\end{proposition}

\begin{corollary}\label{cor:value-is-kernel}
The surviving value $t=2$ is the value at which $r=1$ and $\Re\rho=\tfrac12$, and
$\tfrac12=1/|\ker G|$: the eight classes of $\Zst$ fall into the four fibres of $G$
two by two, so every fibre has the size of the kernel and the apex is its inverse
(\LeanInline{the\_apex\_is\_the\_inverse\_fibre\_size},
\LeanInline{every\_fibre\_has\_the\_kernel\_size} \tierP).  The chain of operators of
\cref{ssec:chain} closes at this value, in its seventh link, and it passes through $\tfrac12$, not through the equality of the two numbers
$L_0=2$ and $|\ker G|=2$, which by itself is only an equality of integers
(\LeanInline{the\_surviving\_value\_is\_the\_kernel}).
\end{corollary}

The same value is where the two branches of the span meet (\cref{prop:span}), and
through the Joukowski map it is the focus of the confocal pencil of the line
(\cref{prop:confocal}).

So the bound $t<3$ is not an approximate statement about a real number: between
$2$ and $3$ the structural spectrum is empty, and the earlier reading in which
the bound was $2.000001$ was an artefact of treating $\beta$ as a free parameter,
a freedom the operator does not offer (\LeanInline{spectrum\_jumps} \tierP).
With the integrality in hand the closure is exact:
$t\in\Z$ together with $t<3$ gives $t=2$ and $r=1$
(\LeanInline{trace\_integer\_plus\_bound} \tierP).

\subsubsection{The geometry: Joukowski and the confocal pencil}\label{ssec:joukowski}

The coordinate $t=r+1/r$ of \cref{ssec:apollonius} is the Joukowski map, and that
gives the critical line a second geometric description, this time in the plane of
the pentagon.

\begin{proposition}\label{prop:joukowski}
For $|z|=1$, $z+1/z=2\Re z\in[-2,2]$: the Joukowski map sends the unit circle to
the segment between the foci.  With $\zeta_{10}=e^{2\pi i/10}$ one has
$\Re\zeta_{10}=\cos(\pi/5)$, so $\zeta_{10}\mapsto\golden$: the pentagon is the
image of the tenth roots of unity under the same map.  And with
$w=(\rho-1)/\rho$, $|w|=r(\rho)=1$ is equivalent to $\Re\rho=\tfrac12$: the locus $r=1$ in
the $w$-plane is the unit circle, the circle on which the roots of unity, the
twentieth among them, lie.
\LeanInline{jouk\_circle\_to\_segment}, \LeanInline{pentagon\_is\_joukowski},
\LeanInline{line\_is\_the\_unit\_circle}, \LeanInline{the\_locus\_is\_the\_circle} \tierP
\end{proposition}

\begin{proposition}[The confocal pencil]\label{prop:confocal}
With $a=r+1/r$ and $b=r-1/r$ one has $a^2-b^2=4$ for every $r\ne0$: the images of
the circles $|z|=r$ form a confocal family with foci $\pm2$, independent of $r$.
The ellipse degenerates exactly at $r=1$, where $b=0$ and $a=2$; the focus $2$ is
the trace value $t=2$ of the line.
\LeanInline{confocal\_always}, \LeanInline{degenerates\_iff\_one} \tierP
\end{proposition}

The critical line is therefore the degenerate member of a pencil whose focus is
the trace value of the line, and the same circle carries the eight classes of
$\Zst$.  Written this way, \Ecore\ says that every zero sits on the degenerate
ellipse; the trace $t$ is its semi-major axis, and \cref{prop:lucas-spectrum} says the structural values of that axis are the Lucas numbers.



\begin{proposition}[The two pencils are one system]\label{prop:pencils}
In the $\rho$-plane, the level sets of $r$ are the Apollonius circles of the pair
$\{0,1\}$, degenerating to the critical line at $r=1$
(\LeanInline{apollonius\_circle\_equation},
\LeanInline{apollonius\_degenerates\_to\_the\_line} \tierP); the level sets of
$\arg w$ are the orthogonal coaxial pencil with the same limit points, and the
critical line is their radical axis.  Under the Joukowski map the first family
goes to the confocal ellipses of \cref{prop:confocal} and the second to the
confocal hyperbolas, with the same foci $\pm2$
(\LeanInline{orthogonal\_pencil\_maps\_to\_hyperbola},
\LeanInline{hyperbolas\_share\_the\_foci}, \LeanInline{one\_confocal\_system},
\LeanInline{two\_pencils\_one\_system} \tierP).  The numerator of $t$ is a circle
centred at $\tfrac12$ and its denominator is the Cassini oval with foci $0$ and
$1$ (\LeanInline{numerator\_is\_a\_circle},
\LeanInline{denominator\_is\_a\_cassini\_oval} \tierP).
\end{proposition}

\begin{remark}[The nodes are golden, and $\Ecore$ is one-sided]\label{rmk:nodes}
The images of the roots of unity under the Joukowski map lie in $[-2,2]$, the
zeroth node being the focus $2$, and the pentagonal nodes are
$\golden,1/\golden,-1/\golden,-\golden$
(\LeanInline{nodes\_lie\_between\_the\_foci}, \LeanInline{zeroth\_node\_is\_the\_focus} \tierP;
\LeanInline{pentagon\_nodes\_are\_golden} \tierP\ given the cosine values, which are
proved in \LeanStar{F1v2.pentagonNodes\_holds}, \LeanStar{F1v2.pentagon\_nodes\_unconditional}).  Moreover $t\ge2$ holds with no
hypothesis --- it is the triangle inequality in the recentred coordinate
(\LeanInline{trace\_ge\_two\_by\_triangle} \tierP) --- so \Ecore\ is one-sided:
$t=2$ if and only if $t\le2$ (\LeanInline{E\_is\_one\_sided} \tierP).  The floor is
geometry; only the ceiling is at stake, and \cref{prop:lucas-spectrum} is what
would supply it.
\end{remark}

\begin{remark}[Pappus' parameter]\label{rmk:pappus}
The ellipse of a zero has semi-axes $t=r+1/r$ and $|r-1/r|$ and foci $\pm2$ (\cref{prop:confocal}), so its
eccentricity is $e=2/t$: the single parameter to which Pappus reduced Apollonius' classification of
the conics.  In it the statements of this section read: $e\le1$ for every zero
(\cref{rmk:nodes}); the pair $\{\rho,1-\bar\rho\}$ shares its conic, since $r\mapsto1/r$ fixes $t$
(\cref{prop:apollonius-coordinate}); $e>0.99993$ with the height of the first zero (\cref{prop:rbound}); $e=1$ exactly on the degenerate conic, $r=1$, the line (\cref{prop:joukowski}); and the structural values are
$2/L_{2m}=1,\tfrac23,\tfrac27,\dots$, of which only $1$ lies above $\tfrac23$ (\cref{prop:lucas-spectrum}).
\LeanStar{F1PaperClosure.eccentricity\_le\_one}, \LeanStar{F1PaperClosure.klein\_pair\_shares\_the\_conic}, \LeanStar{F1PaperClosure.eccentricity\_confined}, \LeanStar{F1PaperClosure.eccentricity\_eq\_one\_iff}, \LeanStar{F1PaperClosure.zero\_conic\_degenerate\_iff\_line}, \LeanStar{F1PaperClosure.integral\_trace\_gives\_degenerate} \tierPnew; the bound with the first zero \tierC
\end{remark}

\begin{remark}[Two integralities]\label{rmk:two-integralities}
They should not be confused.  The Dirichlet coefficients of $L(s,\chiv)$ are
integers unconditionally, since $\chiv$ takes values in $\{0,\pm1\}$
(\LeanInline{chi5\_is\_integral} \tierP); the integrality of $t(\rho)$ is
conditional, and is exactly the core (\cref{prop:integrality-is-E}; the distinction is recorded as
\LeanTally{two\_integralities}).
\end{remark}

\subsubsection{The operator chain}\label{ssec:chain}

The pentagon supplies integrality on the operator side.  The chain of
the corpus has eight links, each a theorem about the operator; that none of them takes
a zero or \Ecore\ as input is recorded as a tally (\LeanTally{the\_chain\_is\_now\_all\_P},
\LeanTally{eight\_links\_none\_from\_E}), and link~6 is the inequality
$r(\rho)^2\le1+1/\gamma^2$ of \cref{prop:rbound}.

\begin{longtable}{clp{6.0cm}l}
\toprule \# & Link & Content & Tier\\\midrule
1 & \LeanStar{F1CW.tower\_ham\_is\_hermitian} Hermitian & real diagonal $\lambda_i=m_0\golden^i$ (\LeanInline{towerH}) & \tierP\\
2 & ratios & $\lambda_i/\lambda_j=\golden^{\,i-j}$; $m_0$ cancels & \tierP\\
3 & $k$ even & $t=L_k$, a Lucas number, hence an integer & \tierP\\
4 & trace & the matrices $F^{k}$, $k$ even, lie in $SL_2(\Z)$ and are constructed, not posited & \tierP\\
5 & norm & $N(r)=+1$ forces $k$ even & \tierP\\
6 & bound & $t\le2.00013$ from the strip and $|\gamma|\ge6.6485$ (\cref{prop:rbound}) & \tierP,\,\tierC\\
7 & value & $t=2$ is the value at which $\Re\rho=\tfrac12=1/|\ker G|$, the inverse of the fibre size of $G$ (\cref{cor:value-is-kernel}) & \tierP\\
8 & closure & $t=2\iff r=1\iff\Re\rho=\tfrac12$ & \tierP\\
\bottomrule\end{longtable}

Two of these were weaker in earlier versions and are not now: link~6 was a
numerical observation and is derived from the strip (\cref{prop:rbound}), and the existential
``$r(\rho)$ lies in the set of eigenvalue ratios'' was opaque and is now
constructed, the witness being explicit
(\LeanInline{existential\_is\_constructed} \tierP); on the line the exponent is
$k=0$, again proved rather than measured
(\LeanInline{k\_zero\_on\_the\_line} \tierP).  The integrality itself comes from
the golden lattice: the trace of the operator is an algebraic integer of $\Zphi$
and the traces are those of the Fibonacci matrix powers: $\mathrm{tr}\,F=1$,
$\mathrm{tr}\,F^2=3$, $\mathrm{tr}\,F^4=7$ with the determinants $-1,1,1$
(\LeanInline{the\_three\_pieces} \tierP; the assembly is \LeanTally{the\_count\_is\_the\_golden\_lattice}).

The chain therefore produces an integer on the operator side, and
\cref{prop:integrality-is-E} says that the same number on the zero side would
give \Ecore.  Whether the two numbers are one number is the question of
\cref{ssec:core}.

\subsubsection{Integrality of the trace, and membership in the tower}\label{ssec:arrow}

The two sides of \cref{prop:integrality-is-E} --- integrality of $t$ and
membership of $r$ in $\golden^{\Z}$ --- are related by two theorems, one of which makes the other unnecessary for the equivalence.

\begin{theorem}[Integrality of the trace]\label{thm:arrow}
Let $r>0$ and $t=r+1/r$.
\begin{enumerate}[label=\textup{(\alph*)}]
  \item If $t\in\Z$ and $t<3$, then $t=2$ and $r=1$.
    \LeanInline{trace\_integer\_plus\_bound} \tierP
  \item If $r\in\Q(\sqrt5)$ and $t\in\Z$, then $r$ is a unit of $\OK=\Zphi$ and
    $r=\golden^{k}$ for some $k\in\Z$; for even $k$, $t=L_k$ is the trace of the
    Fibonacci matrix $F^k$.
    \LeanStar{F1v2.arrowB\_full}, \LeanStar{F1v2.arrowB\_Q}, \LeanStar{F1v2.arrowB} \tierPnew\
    (the step into $\Zphi$ by the rational root theorem and the primality of $5$; the
    descent by \LeanInline{unit\_is\_a\_power\_of\_phi} with its three inputs proved)
  \item The hypothesis $r\in\Q(\sqrt5)$ in \textup{(b)} cannot be removed:
    $r=2+\sqrt3$ has $t=4\in\Z$ and is not a power of $\golden$
    (\LeanStar{F1v2.arrow\_needs\_the\_field} \tierPnew).
\end{enumerate}
\end{theorem}

\begin{proof}
(a)~$t\ge2$ always, with equality exactly at $r=1$; the only integer in $[2,3)$ is
$2$.  (b)~$r$ is a root of $x^2-tx+1$, a monic polynomial with integer
coefficients, so $r$ is an algebraic integer; lying in $\Q(\sqrt5)$, it lies in
$\OK=\Zphi$ (\cref{prop:OK}).  So does $1/r$, the other root, and $r\cdot(1/r)=1$,
so $r$ is a unit; being positive, $r=\golden^k$ by \cref{thm:units}, with no appeal
to Dirichlet's theorem.  The trace statement is Binet
(\cref{prop:binet}).  (c)~$(2+\sqrt3)+(2-\sqrt3)=4$, and $2+\sqrt3\notin\Q(\sqrt5)$.
\end{proof}

\begin{remark}[Where this leaves the core]\label{rmk:arrow-direction}
Part~(a) is what the equivalence uses: the implication from $t(\rho)\in\Z$ to $\Re\rho=\tfrac12$
closes from integrality and the strip bound alone, with no need for the tower.
Part~(b) identifies, inside the field of the pentagon, integrality of the trace
with membership in the tower, so that $t(\rho)\in\Z$ and $r(\rho)\in\golden^{\Z}$ are one statement there.
Both reappear as conditions~(ii) and~(iii) of \cref{thm:topological-proof}.
What is open is that the number attached to a \emph{zero} has the property the
operator's numbers have --- integrality of $t(\rho)$, or $r(\rho)\in\Q(\sqrt5)$
together with it; this is the integrality form of \cref{ssec:core-forms}.  The witness for $r(\rho)$ is constructed whenever $r$ is known
to be a spectral ratio (\LeanInline{the\_witness\_is\_constructed},
\LeanInline{is\_a\_spectral\_ratio} \tierP), and knowing that is knowing \Ecore.
\end{remark}

\begin{remark}[A correction to the corpus]\label{rmk:unitstructure}
The corpus states (b) as \LeanInline{algebraic\_integer\_gives\_MonomialHalf} with an
input \LeanInline{UnitStructure} that omits the hypothesis $r\in\Q(\sqrt5)$; as
written that input is false, by (c).  With the hypothesis added it is exactly (b),
and it is then no longer a classical input: \LeanStar{F1v2.arrowB} proves it,
instantiating the corpus's descent \LeanInline{unit\_is\_a\_power\_of\_phi} with the units
of $\Zphi$ and proving its three structural inputs --- the floor of the logarithm, the
group law, and the absence of units in $(1,\golden)$ --- rather than assuming them.
\end{remark}

\subsubsection{What is proved of each zero}\label{ssec:each-zero}

The facts established so far hold for every zero of $L(s,\chiv)$ in the strip, and
none uses that the zero lies on the line.  Collected:

\begin{proposition}[Each zero, unconditionally]\label{prop:each-zero}
Let $\rho=\beta+i\gamma$ be a zero of $L(s,\chiv)$ with $0<\beta<1$, and
$r(\rho)=|1-1/\rho|$ its coordinate.
\begin{enumerate}[label=\textup{(\roman*)}]
  \item The zero set is closed under $\rho\mapsto1-\rho$, under conjugation and under
    the mate $\rho\mapsto1-\bar\rho$.
  \item The centre of the orbit $\{\rho,1-\rho\}$ is $\tfrac12$
    (\cref{prop:orbit-embedding}); the orbit has two elements if $\beta=\tfrac12$
    and four otherwise.
  \item $r(1-\rho)=r(\rho)^{-1}$.
  \item $r(\rho)^2\le1+1/\gamma^2$ (\cref{prop:rbound}); with the height of the first
    zero, $0.989\le r(\rho)\le1.012$, an interval where the only power of $\golden$ is
    $\golden^0=1$.
  \item $t(\rho)=r+1/r\ge2$; with the height of the first zero, $t(\rho)\le2.00013$, and
    the only integer in that interval is $2$.
  \item $r(\rho)=1$ exactly when $\beta=\tfrac12$.
\end{enumerate}
{\raggedright\LeanInline{functionalEquation\_chi5\_discharged} \tierP, \LeanStar{realCoefficients\_chi5},
\LeanStar{closedUnderMate\_chi5}, \LeanStar{F1PaperNew.zIsAVSet\_chi5};
\LeanInline{the\_orbit\_embedding\_is\_unconditional}, \LeanInline{orbit\_card\_upstairs},
\LeanInline{where\_the\_chain\_breaks}, \LeanInline{rSq\_le\_in\_strip},
\LeanInline{trace\_ge\_two\_by\_triangle}, \LeanInline{liMap\_on\_unit\_circle\_iff} \tierP;
the height of the first zero \tierC\par}
\end{proposition}

What these facts do not give is that the zero is a member of the structure and not
only centred on it: that its coordinate, already confined to the stretch where the only
power of $\golden$ is $1$, \emph{is} a power of $\golden$.  That is the one open statement of the paper;
\cref{ssec:core-forms} gives its forms and \cref{ssec:route-overview} the route of closure.

\subsubsection{The squeeze}\label{ssec:squeeze}

The classical shape of the argument is a squeeze: an upper bound
$\beta\le\tfrac12$ applied to $\rho$ and to its companion gives
$1-\beta\le\tfrac12$, hence $\beta=\tfrac12$.  The corpus formalises it with the
bound as a structure field.

\begin{definition}
A \emph{ternary $L$-zero} is a real number $\beta$ with $0<\beta<1$,
$\beta\le\mu_3$ and $1-\beta\le\mu_3$, where $\mu_3=\tfrac12$ is the spectral
parameter of \cref{ssec:omega}.  \LeanInline{TernaryLZero}
\end{definition}

\begin{proposition}\label{prop:squeeze}
Every ternary $L$-zero has $\beta=\tfrac12$
(\LeanInline{rh\_squeeze} \tierP), and a $\chiv$-zero becomes one as soon as the
upper bound holds for all zeros, the companion bound being then automatic
(\LeanInline{TernaryLZero.ofChi5Zero}, \LeanInline{rh\_squeeze\_derived},
\LeanInline{rigidity\_forces\_sign} \tierP).
\end{proposition}

The step from the bound to \Ecore\ is therefore proved; the honest accounting is
what the bound costs.

\begin{proposition}[Premise A is not weaker than the conclusion]\label{prop:premiseA}
$(\forall z,\ \Re z.\rho\le\mu_3)\iff(\forall z,\ \Re z.\rho=\tfrac12)$.
\LeanInline{premise\_A\_iff\_GRH\_chi5} \tierP
\end{proposition}
\begin{proof}
One direction is \cref{prop:squeeze}.  The other is immediate.
\end{proof}

So the squeeze turns \Ecore\ into an upper bound, $\Re\rho\le\tfrac12$ for every zero, and that bound is
condition~(i) of \cref{thm:topological-proof}.  The same must be said of the functional equation by
itself: invariance of the zero set under $s\mapsto1-s$ does not place any zero,
since the pair $\{\rho,1-\rho\}$ is invariant for every $\rho$
(\LeanInline{bound\_alone\_insufficient} \tierP).  What the rest of this section
does is look for the bound elsewhere: in an operator whose spectrum is known
exactly.

\subsubsection{Membership, derived from the bridge}\label{ssec:membership}

In the corpus the identification is the hypothesis \LeanInline{BridgeAxiom}: that the
real part $\beta$ of a zero is one of the real parts $\tfrac12,-\tfrac14$ of the
spectrum of $\OmH$.  Granted that, membership in the tower is a derived statement,
not a separate assumption.

\begin{theorem}[Membership]\label{thm:membership}
Assume \LeanInline{BridgeAxiom} for $\beta$, the strip, and the companion.  Then
$\beta=\tfrac12$; equivalently $|\rho-1|=|\rho|$; equivalently
$r=1=\golden^0$ lies in the golden tower.
\LeanInline{the\_membership\_derived} \tierP
\end{theorem}

The corpus tallies eight ingredients of the derivation, five proved in Lean and three
not: the bridge, Hecke's companion and the strip
(\LeanTally{five\_proved\_three\_classical}).  The last two have since been proved
(\LeanInline{hecke\_chi5\_derived}, \LeanInline{Chi5Zero}), so the bridge is the only one
left.  The earlier statement \LeanInline{MembershipStillOpen} is therefore superseded:
membership is neither open nor unconditional; it is derivable from the bridge.

\subsubsection{The span and the bridge}\label{ssec:span}

What would identify the two numbers is a span: two branches with a common source,
the Golden Prime homomorphism $G$ of \cref{sec:grisales}.

\begin{proposition}[The span]\label{prop:span}
Both branches have source $G$: the arithmetic branch through the character
$\chiv=(-1)^{\pi_2\circ G}$, the Euler factors and $L(1,\chiv)$; the spectral
branch through $|\ker G|=2$ and the modulus $\mu_3=\tfrac12$.  They meet at
$\log2=\log|\ker G|$, which is the entropy bridge of \cref{eq:log-two}.
\LeanInline{branches\_meet\_at\_log\_two}, \LeanInline{G\_kernel}, \LeanInline{G\_pi2\_is\_chi5} \tierP;
the span as a whole is recorded as \LeanTally{the\_span\_is\_the\_identification}.
\end{proposition}

\begin{proposition}[The pieces of the bridge]\label{prop:bridge-pieces}
The corpus assembles three pieces that the earlier manuscript called on to discharge
\LeanInline{BridgeAxiom}.  What is proved of them is: the identity
$\lng\cdot\tfrac{\sqrt5}2\cdot L(1,\chiv)=\log2$ given the value of $L(1,\chiv)$
(\LeanInline{branches\_meet\_at\_log\_two} \tierP); that $G$ is a homomorphism whose second
component is $\chiv$, so $\chiv=(-1)^{\pi_2\circ G}$ (\LeanInline{G\_hom},
\LeanInline{G\_pi2\_is\_chi5} \tierP); that $\ker G=\{1,9\}$ and $\mu_3=1/|\ker G|$; and
that a third constant of the construction is positive.  The corpus collects these in
\LeanTally{the\_bridge\_pieces\_formalized}, which also checks the homomorphism on two
products and the character on two classes.
\end{proposition}

In \cref{sec:primon} the identity of the first piece is read as the entropy $S_{\mathrm{BH}}/k_B$ of one bit.

\begin{remark}[What is still a marker]\label{rmk:markers}
Neither \cref{prop:span} nor \cref{prop:bridge-pieces} identifies the two traces.
In the corpus the identification is a marker, \LeanInline{TheSpanIdentifiesThem},
defined as \texttt{True} \tierO, and so is the label \LeanInline{BridgeIsCToP}: the
spectral mapping theorem is not proved in Lean.  The three pieces speak of
$L(1,\chiv)$, of constants and of $G$; none speaks of a zero.  The claim of the
earlier manuscript that they discharge \LeanInline{BridgeAxiom} is not sustained
by the corpus, and this paper does not make it (\cref{app:superseded}).
\end{remark}

\begin{proposition}[The bridge, link by link]\label{prop:bridge-links}
The bridge of \cite{gonzalez25hp} is a chain whose links are theorems:
\begin{enumerate}[label=\textup{(\roman*)}]
  \item the Frobenius lifts realise multiplicatively, $\mathrm{realise}(\Frob_p n)=\mathrm{realise}(n)^p$;
  \item each local Euler factor splits through $\chiv$, $f_\zeta\,f_L=f_K$;
  \item the factors assemble the Dedekind factorisation $\zeta_K=\zeta\cdot L(\cdot,\chiv)$;
  \item the analytic character is the second component of $G$;
  \item every spectral position with positive real part is $\tfrac12$.
\end{enumerate}
The last link --- that the real part of each zero is a spectral position --- is the spectral
identification; given it, every zero lies on the line.
\LeanInline{bridge\_components}, \LeanInline{bridge\_assembled} \tierP
\end{proposition}

The spectral identification is one of the forms of the open statement (\cref{ssec:core-forms}), and the
closure conjecture would give it (\cref{cor:closure}).

% ---------------------------------------------------------------------
\subsection{The thread from $\golden$}\label{sec:topological-proof}
\src{HP \texttt{ssec:squeeze} (L4307); \texttt{F1\_compilado\_sin\_axiomas};
\texttt{PCF\_unificado}; \texttt{PCF\_sesion\_M12\_unificado} \S129--\S205;
\texttt{everything\_descends}}

The proof is a single thread.  It starts from one relation, $\golden^2=\golden+1$,
and at no step does it take the primes, the Euler products or the zeros as a
starting point: each of them is \emph{produced} by the thread, and each production
is a theorem.  The primes enter only as an index set, through their classes; no property of their
distribution is used.  The strategy is the one announced in the introduction --- from
arithmetic to geometry, from geometry to analysis --- and it is carried out in eight
stages.  It is assembled here, at the end of the Methods, because every stage is a result proved in the
sections before it: the stages name those results and the order in which they compose.  Every stage uses only what the previous stages produced, together with
finite arithmetic and classical facts about $\Gamma$ and the Hurwitz function; the
statement of each stage names the Lean declarations that carry it.

\subsubsection{Stage 0. The axiom: $\golden$}\label{ssec:stage0}

The only input is the golden ratio as the positive root of $x^2=x+1$, which is also
the diagonal of the regular pentagon: $\golden=2\cos(\pi/5)=1+\zeta_5+\zeta_5^4$
(\LeanInline{phi\_sq}, \LeanInline{phi\_eq\_two\_cos\_pi\_fifth},
\LeanInline{phi\_cyclotomic\_inscription} \tierP).  It is the only positive root of
$x^2=x+1$ (\LeanInline{golden\_fixed\_point\_unique} \tierP), and it is a fixed point of self-reference in two independent forms, the
continued fraction $\golden=1+1/\golden$ and the nested radical
$\golden=\sqrt{1+\golden}$ (\cref{lem:phi-cf,lem:phi-nr}).  Everything below is derived from this.

\subsubsection{Stage 1. From $\golden$ to the primes, through geometry}\label{ssec:stage1}

The pentagon fixes the factor $5$ and the quarter turn $i^4=1$ fixes the factor $4$
(\cref{prop:conductor}), and the recurrence itself returns the conductor: the
Fibonacci sequence modulo $5$ has period $20$ (\cref{prop:pisano});
the conductor $20=\operatorname{lcm}(4,5)$ is derived, not chosen
(\LeanInline{conductor\_is\_derived}, \LeanInline{conductor\_is\_joint\_period} \tierP).
On $\Zst$ the Galois functor $G\colon\Zst\to\Z_2\times\Z_2$ is a surjective
homomorphism whose components are $\chi_4$ and $\chiv$, with kernel $\{1,9\}$
(\cref{prop:G}).  The component $\chiv$ is read directly off the pentagon:
$|2\cos(\pi a/5)|=\golden^{\chiv(a)}$ (\cref{thm:pentagonal-values}).  The two
factors separate the two characters, each blind to the other
(\LeanInline{conductor\_attribution} \tierP), and the
crystallographic restriction explains why only the $5$-side can carry a tower
(\cref{prop:crystallographic,prop:unit-ranks}).

The primes enter here, and they enter as \emph{classes}: every prime $p\nmid20$
falls into one of the eight classes of $\Zst$, and the primes meet all four fibres of
$G$ (\LeanInline{primes\_hit\_all\_four\_fibres}, \LeanInline{primes\_reach\_the\_tower}
\tierP).  What a prime \emph{is}, for everything that follows, is its fibre: split or
inert, decided by the geometry of the pentagon and the turn.  No list of primes and
no property of their distribution is used.

\subsubsection{Stage 2. From geometry to rings: the two registers}\label{ssec:stage2}

The pentagon's number generates the ring $\Zphi=\OK$ (\cref{prop:OK}), a domain with
exactly two endomorphisms (\cref{prop:domain,prop:endo}), whose positive units are
exactly the powers of $\golden$, proved without Dirichlet's theorem
(\cref{thm:units}).  The rings form a tower, the levels
$\LPCF^{(\sigma)}=\golden^{\sigma}\LPCF$ \cite[\S2.8, \S3.4]{gonzalez26cw}, and the tower
preserves both registers:
\begin{itemize}
  \item the \emph{multiplicative register} --- the binary register with period $4$,
    and the lift $\golden^n\mapsto\golden^{pn}$ between levels, which composes as the
    primes multiply (\cref{thm:lambda-monoid});
  \item the \emph{additive register} --- inside each level, $\Frob_p$ is the one ring endomorphism congruent to the $p$-th power, $\Frob_p(a)\equiv a^p$, the identity or the conjugation (\cref{thm:frobenius-congruence}), the lifts
forming a commuting family that composes as the primes multiply
(\cref{prop:frob-compose}); modulo $p$
    it reads $F_p\equiv\chiv(p)$ (\cref{thm:two-spectra,thm:sign-is-chi5,thm:chi5-from-frobenius}).
\end{itemize}
The two are separated by the two factors of the conductor
(\cref{prop:two-registers-G}), they agree modulo every prime where both exist
(\cref{thm:monoid-ring-bridge}), and the only place they cannot both exist is the
ramified prime, where the obstruction is a nilpotent (\cref{thm:obstruction-five});
inverting $5$ removes it (\cref{cor:localisation}); together this is the descent datum of \cref{thm:descent-datum}.

\subsubsection{Stage 3. From rings to Euler products}\label{ssec:stage3}

The Euler factors are not an input; they are the characteristic polynomials of the
Frobenius of Stage~2: $\det(1-p^{-s}M_p)=(1-p^{-s})(1-\chiv(p)p^{-s})$, with
$\Frob_p(\sqrt5)=\chiv(p)\sqrt5$ (\cref{thm:frob-eigen,lem:psi-to-euler}).  The local factors, one per
fibre, are \eqref{eq:local-K}; the Euler product is their colimit over finite sets of
primes, each coefficient being fixed at a finite stage
(\LeanStar{F1CW.coeff\_stabilises\_K}, \LeanStar{F1CW.colimit\_exists\_K} \tierP); the Dirichlet series $\sum n^{-s}$ is the Euler product of $\zeta$
(\LeanStar{F1CW.dirichlet\_is\_euler} \tierPnew); and the
Dedekind factorisation $\zeta_K=\zeta\cdot L(\cdot,\chiv)$ holds factor by factor and
globally, with Mathlib's analytic objects (\cref{eq:dedekind}), as a natural
isomorphism that is universal among such factorisations
(\cref{thm:dedekind-natural,cor:factorisation-universal}).  Evaluated at the
odd integers it gives the pentagonal determination of $\zeta(2k+1)$, with $L$ a
combination of Hurwitz values at the points $a/5$ (\cref{thm:odd-zeta}).  So $\zeta$
and $L(s,\chiv)$ are outputs of Stages~0--2.

\subsubsection{Stage 4. Through the anchor, from rings to function}\label{ssec:stage4}

The discrete tower becomes a function through one exponent.  With
$\lng=\log2/\log\golden$, $\golden^{\lng}=2$, and $2$ is not a power of $\golden$, so
the exponent anchors the scale rather than being derived from it
(\cref{prop:mersenne-bridge}).  The function $x\mapsto\golden^{x}$ carries both
spectra of Stage~2: multiplicatively, the powers compose; additively,
$\golden^{n+1}=F_{n+1}\golden+F_n$, so multiplication in the tower reads as addition
of Fibonacci coordinates and adjacent powers carry
(\cref{ssec:fib-coords}, \cref{prop:binet}, \LeanInline{multiplicative\_reads\_as\_additive},
\LeanInline{carry\_in\_base\_phi} \tierP).  The logarithm turns the tower into the
exponent lattice $X_{\PCF}=\log\golden\cdot\Z$, the additive monoid on which the lifts
act as multiplication by $p$ (\cref{thm:lambda-monoid,cor:frob-as-mul}).  The anchor also measures
the arithmetic: $\lng\cdot R_K=\lng\cdot\frac{\sqrt5}2L(1,\chiv)=\log2$, the regulator
cancelling exactly (\cref{eq:log-two}), consistently with the residue of $\zeta_K$
at $s=1$ (\cref{cor:residue}).  From the anchor, with $\sqrt5=2\golden-1$ and
$\pi=5\arccos(\golden/2)$, every value of $\zeta$ and of $L(\cdot,\chiv)$ used later is
written in base $\golden$ (\cref{thm:base-phi}).

\subsubsection{Stage 5. From function to operator}\label{ssec:stage5}

The function yields operators with exact spectra, built before any zero is mentioned.
The eigenvalue triad $\OmH=\tfrac12\operatorname{diag}(1,\zeta_3,\zeta_3^2)$ has
real parts $\tfrac12,-\tfrac14,-\tfrac14$, only $\tfrac12$ positive, and its three
component norms multiply to $\tfrac12$ (\LeanInline{positive\_spectral\_half},
\LeanInline{norm\_Omega\_is\_half} \tierP).  The tower Hamiltonian $\mathrm{diag}(m_0\golden^{i})$, with $m_0>0$ a scale, is Hermitian, its eigenvalue ratios are
$\golden^{i-j}$, and for even index the traces of the Fibonacci matrix powers are
Lucas numbers, integral, with discrete values $2,3,7,18,47,\dots$, the value $2$
being the one at which $\Re\rho=\tfrac12=1/|\ker G|$ (\cref{thm:one-operator} for the operator and \cref{thm:operator-forms} for its forms, and \cref{cor:no-diagonal-spectral} for the bound $\|\OmH\|=\tfrac12<1$;
\LeanStar{F1CW.tower\_ham\_is\_hermitian},
\cref{prop:lucas-spectrum,cor:value-is-kernel}).  On $\ell^2(\N)$ all the lifts act at
once, by the shifts $e_n\mapsto e_{pn}$ with the Frobenius of Stage~2 as twist
(\cref{prop:primon-space}).  The operator meets the requirements of
non-circularity, self-adjointness and exact spectrum
(\cref{prop:omega-spectrum}; tallied in \LeanTally{omega\_meets\_three}): its spectrum is
fixed by $\golden$, not calibrated on data.

\subsubsection{Stage 6. The zeros meet the operator}\label{ssec:stage6}

Only now do the zeros appear, and they appear as the objects whose position the structure constrains --- the centre $\tfrac12$ of each pair, the pairing, the heights up to a correction --- not as data.  The correction is what the closure conjecture of the
Results has to supply, and its first step is already proved: on each level of the clock of $\golden$ the product
formula on the cyclotomic norms $N(\Phi_n(\golden))$ fixes the weights of the primes of that level
in sum (\cref{prop:clock-cyclotomic}); what remains is their order by size.  A zero is a structure (\LeanInline{Chi5Zero}),
and the symmetries it carries --- the functional equation, the companion, the
conjugation --- are theorems of Stage~3 (\cref{sec:analytic-inputs}).  Three
contacts are unconditional.

\emph{The centre.}  For every zero, the centre of the orbit $\{\rho,1-\rho\}$ is
$\tfrac12$, the positive eigenvalue of $\OmH$
(\LeanInline{zeros\_enter\_the\_cocone\_by\_the\_involution},
\LeanInline{centre\_descends}, \cref{prop:orbit-embedding}); the apex that receives
the primes, the operator, the monoid and the archimedean place also receives the
zeros (\LeanInline{the\_apex\_receives\_from\_five\_legs\_including\_the\_zeros} \tierP).

\emph{The line.}  The critical line is the fixed set of $s\mapsto1-\bar s$ in the
Klein group generated by the functional equation and conjugation
(\cref{prop:stau}), whose orbits on the zero set are of size two on the line and four
off it (\cref{prop:orbits}); through the Apollonius coordinate $r=|\rho-1|/|\rho|$ it is the
level $r=1$, and through the Joukowski map it is the degenerate member of a confocal
pencil whose focus is the trace value $t=2$ of the line
(\cref{prop:joukowski,prop:confocal,prop:pencils,cor:value-is-kernel}).
The line is therefore produced by the structure of Stages~1--5, not assumed.

The Klein group acting on the zeros is induced by the two involutions of $\chiv$ on $\Zst$, inversion and negation
(\cref{prop:descent}), so the orbit partition of the zeros, of size two on the line and four off it, descends from $\golden$
with the character.


\subsubsection{Stage 7. From the operator to membership}\label{ssec:stage7}

The last stage places each zero in the spectrum.  Its coordinate $t=r+1/r$ satisfies
$t\ge2$ unconditionally and $t\le2.00013$ by the strip and the height of the first zero (\cref{prop:rbound,rmk:nodes});
the operator's traces are the Lucas numbers, whose only value below $3$ is
$2$ (\cref{prop:lucas-spectrum}); integrality of the trace with that bound gives
$t=2$ and $r=1$ directly, and inside $\Q(\sqrt5)$ it is the same as membership in
the tower (\cref{thm:arrow}).  Position is then interval meets discreteness: an integer power of
$\golden$ in the window $|\log r|<\log\golden$ is $\golden^0$, so $r=1$ and
$\Re\rho=\tfrac12$ (\LeanInline{position\_is\_interval\_meets\_discreteness} \tierP).
Membership follows from the bridge with its inputs named
(\cref{thm:membership}), and the multiplicative half of \Ecore, $r\cdot r'=1$, holds on every zero, so that the additive half $r+r'=2$ alone would close it (\LeanInline{coords\_closed\_under\_inversion}, \LeanInline{union\_of\_halves\_is\_E} \tierP).

\subsubsection{The thread, assembled}\label{ssec:thread}

\begin{theorem}[The thread, and the forms of the open statement]\label{thm:topological-proof}
Starting from $\golden^2=\golden+1$ alone, Stages~1--7 produce, by proved steps, the
classes of the primes, the rings with their additive and multiplicative spectra, the
Euler products and the Dedekind factorisation, the function through the anchor, the
operator with exact spectrum, the critical line as a fixed set, and the orbit
partition of the zeros.  None of these steps takes the zeros or the Euler products as a
starting point, nor uses of the primes anything but their classes in $\Zst$; the
corpus records the eight links of Stage~7 as a tally and checks that none is tagged as
taking a zero or \Ecore\ as input (\LeanTally{the\_chain\_is\_now\_all\_P},
\LeanTally{eight\_links\_none\_from\_E}).  For a zero $\rho$ of $L(s,\chiv)$, each of
the following conditions implies $\Re\rho=\tfrac12$, and each is implied by it; each
implication is a theorem, and none of the four conditions is proved (they are four of the
forms of \cref{ssec:core-forms}):
\begin{enumerate}[label=\textup{(\roman*)}]
  \item the upper bound $\Re\rho\le\tfrac12$ (the squeeze with the companion of
    Stage~6);
  \item the coordinate $r(\rho)$ in the tower $\golden^{\Z}$ (Stage~7);
  \item the trace $t(\rho)$ integral, as the operator's traces are (Stages~5 and 7);
  \item no two zeros sharing an ordinate (the Klein orbits of Stage~6).
\end{enumerate}
\LeanInline{rh\_squeeze\_derived}, \LeanInline{position\_is\_interval\_meets\_discreteness}, \LeanInline{radius\_one\_is\_E},
\LeanInline{the\_two\_labels\_are\_one\_gap},
\LeanStar{E\_iff\_no\_shared\_ordinate\_chi5} \tierP/\tierPnew
(for (i) the bound over all zeros; for (ii) given $r(\rho)=\golden^k$; for (iii)
given $t(\rho)\in\Z$ and $|\Im\rho|\ge2$).
\end{theorem}

\begin{proof} Every result cited below is proved earlier in the Methods, in \cref{sec:operator,sec:machinery}, without appeal to this theorem.
Stages~0--7 are the propositions cited in each stage.  For the closing links, each
condition is assumed and the line derived:
(i)~if the bound holds for every zero, the companion $1-\rho$ is a zero, so the bound applied to it gives
$1-\Re\rho\le\tfrac12$, whence $\Re\rho=\tfrac12$ (\cref{prop:squeeze}; the bound is
itself equivalent to the conclusion, \cref{prop:premiseA}).
(ii)~if $r(\rho)=\golden^k$, then $|k\log\golden|=|\log r(\rho)|<0.0142<\log\golden$
forces $k=0$; $r=1$ is the line.
(iii)~if $t(\rho)\in\Z$, then $2\le t<3$ forces $t=2$, hence $r=1$.
(iv)~off the line the mate $1-\bar\rho$ is a distinct zero with the same ordinate
(\cref{thm:shared-ordinate}).
\end{proof}

\begin{figure}[ht]
\centering\footnotesize
\begin{tikzcd}[column sep=1.5em,row sep=2.6em,ampersand replacement=\&]
\mu_6\times\golden^{\Z} \arrow[d,"{\text{level}}"']
\& \ell^2\ \circlearrowleft\ H=\operatorname{diag}(m_0\golden^i) \arrow[dl,dash,dotted]
\& \zeta_5^{\,a} \arrow[d,"{2\cos(\pi a/5)=\golden^{\chiv(a)}}"]
\\
\golden^{\Z}\subset\Zphi \arrow[r,"{\Frob_p}"] \arrow[d,"{\Phi_n(\golden)}"']
\& \Frob_p\in\{\mathrm{id},\sigma\},\ \varepsilon(p) \arrow[r,equal] \arrow[d]
\& \chiv(p) \arrow[d,"{\text{Hurwitz at }a/5}"]
\\
\text{clock of }\mathfrak p \arrow[d,"{\text{product formula}}"']
\& (1-\varepsilon(p)p^{-s})^{-1}(1-p^{-s})^{-1} \arrow[r,"{\prod_p}"]
\& \zeta_K=\zeta\cdot L(s,\chiv) \arrow[dl,"{s=2k+1}"'] \arrow[d,"{\rho\mapsto r(\rho)=|1-1/\rho|}"]
\\
{|N\Phi_n(\golden)|_\infty=\prod_{\mathfrak p}|N\Phi_n(\golden)|_{\mathfrak p}^{-1}}
\& \zeta(2k+1)=\zeta_K(2k+1)/L(2k+1,\chiv)
\& r(\rho)r(1-\rho)=1,\ r(\rho)\in W,\ W\cap\golden^{\Z}=\{1\} \arrow[d,dashed,"{\Ecore}"]
\\
\& \& r(\rho)=1\iff\Re\rho=\tfrac12
\end{tikzcd}
\caption{\Cref{fig:construction-words} after the thread, with the statement behind each arrow.  Fibre and operator:
\cref{thm:fibre-mu6,thm:one-operator,thm:operator-forms,prop:level-operators}; the lift and its sign: \cref{thm:frobenius-congruence,thm:sign-is-chi5};
the pentagon: \cref{thm:pentagonal-values}; the clock and the product formula: \cref{prop:clock-cyclotomic};
Euler factors and product: \cref{thm:frob-eigen,thm:dedekind-natural}; odd values and Hurwitz: \cref{thm:odd-zeta};
the zeros: \cref{prop:each-zero,prop:two-halves}; the dashed arrow: \Ecore, \cref{thm:topological-proof}.  The
lift on the tower, $\golden^n\mapsto\golden^{pn}$, and the lift on the ring agree only modulo $p$
(\cref{thm:monoid-ring-bridge}); the diagram uses the second.}
\label{fig:construction-formulas}
\end{figure}


\begin{remark}[What is left to close]\label{rmk:left-to-close}
Everything descends from $\golden$: the classes, the rings, the registers, the Euler
products, the function, the operator, the line, and the Klein symmetry of the zeros
with its two orbit sizes, induced by $\chiv$.  What remains is the verdict on the zeros of
$L(s,\chiv)$: the descended invariant is the size of the Klein orbit of a zero, two on the line and four off it
(\cref{prop:orbits,prop:descent}), and the verdict is which of the two values each zero takes; and
the four links of \cref{thm:topological-proof} are four ways of stating that last
step (\LeanInline{everything\_descends\_but\_the\_verdict} \tierP).  The corpus has
the span that would carry it constructed with its two branches (\cref{prop:span}),
integrality of the trace shown to force $r=1$ and, inside $\Q(\sqrt5)$, to be
membership in the tower (\cref{thm:arrow}), and the multiplicative half proved on
every zero; the step still to be written is that the
coordinate of each zero carries the discreteness of the operator
(\LeanInline{MonomialHalfP}, \LeanInline{monomial\_is\_in\_image\_of\_realiseP}).  The zeros also return the conductor of Stage~1 --- their unfolded spacing has mean
closest to $1$ at $q=5$ (\cref{prop:envelope-exact}) --- and \cref{ssec:three-regimes} records the
heights where the verdict is already settled, and \cref{ssec:weil-route} the windows
of Weil positivity proved along the way.
\end{remark}

% =====================================================================

\section{Results}\label{sec:results}

This section collects what the construction and the equivalent forms establish about the zeros:
the statements settled unconditionally on the zero set, the Weil analogy with its
correspondence table, the windows of Weil positivity, and the verdict that remains.

\subsection{What is proved about the zero set}\label{sec:results-zeros}

\subsubsection{Three regimes}\label{ssec:three-regimes}

\begin{proposition}[Three regimes]\label{prop:three-regimes}
For the zeros of $L(s,\chiv)$: no zero lies off the line below $\gamma=7.3241$;
below $\gamma=70$ there are exactly $33$ zeros, all on the line; and at all heights
$0\le\sum_\rho(r_\rho-1)^2/r_\rho\le2\lambda_1-\lambda_2/2\approx3.5695\cdot10^{-4}$.  The evidence marks are
given item by item below.
\tier{N}%
\end{proposition}

Alongside the equivalent forms, the corpus proves or decides quantitative statements about
the zeros of $L(s,\chiv)$ that do not depend on the core.  They fall into three
regimes by height, and each carries its own status.

\begin{enumerate}[leftmargin=2.2em]
  \item \textbf{Below $\gamma=7.3241$: no zero off the line.}  A hierarchy of sums
    $S_n=\sum_\rho\bigl(\rho(1-\rho)\bigr)^{-n}$ over the zeros, each bounded through
    the Li coefficients, excludes an off-line pair below a height that depends on $n$;
    the cut-offs are computed numerically and enter as inputs; the cut-offs for
    $n=1,2,3,8,20$ are $5.1614$, $7.2820$, $7.3241$, $6.9615$, $6.7838$, the
    optimum is at $n=3$, and it covers the first zero $\gamma_1=6.6484$
    (\tier[optimum\_at\_three, covers\_first\_zero]{P}\ given
    the numerical cut-offs \tier{N}).
  \item \textbf{Below $\gamma=70$: decided.}  There are exactly $33$ zeros, all on
    the line, established numerically by counting the zeros on the line and in the strip
    and finding the two counts equal
    (\LeanInline{DoubleCount}, \LeanInline{invariant\_two\_on\_measured\_range} \tier{N}).  The independent computation of
    \cref{prop:envelope-exact} finds the same $33$ below $70$.
  \item \textbf{All heights: a global sandwich.}  With $r_\rho=|\rho-1|/|\rho|$,
    \[
      0\;\le\;\sum_\rho\frac{(r_\rho-1)^2}{r_\rho}\;\le\;2\lambda_1-\frac{\lambda_2}2,
    \]
    where $\lambda_1,\lambda_2$ are the first Li coefficients of $L(s,\chiv)$,
    $\lambda_n=\sum_\rho\bigl[1-(1-1/\rho)^n\bigr]$, the sum over the zeros taken in conjugate pairs
    (\LeanInline{liCoeff}); the
    lower end is an identity and vanishes term by term exactly on the line, the
    upper end comes from the Hadamard product (\tier[the\_sandwich]{P}\
    given two inputs taken as hypotheses: the value of the second Li sum\ and an
    asymptotic bound on its ratio to the excess).  Numerically the bound is
    $3.5695\cdot10^{-4}$.  The corpus also records a sharper value, $3.472359\cdot10^{-7}$,
    obtained by composing the sandwich with the decided range below $\gamma=70$; that
    composition is numerical and is not formalised, and Lean checks only that the second
    value is more than a thousand times smaller (\LeanInline{improvement\_over\_thousand}
    \tier{P}\ given both values \tier{N}).
\end{enumerate}

Four \emph{equalities} hold on the whole zero set, not merely bounds: the pairing
$r_\rho\cdot r_{\rho'}=1$ for $\rho'=1-\rho$, hence $\sum_\rho\log r_\rho=0$; the closed
form of the first Li coefficient, $\lambda_1=\sum1/\rho=\kappa/2+1-\gamma_E/2$, with $\gamma_E$ Euler's constant and
$\kappa=\log(q/\pi)+\psi(\tfrac14)$, $q=5$, $\psi=\Gamma'/\Gamma$ (the constant of \cref{thm:first-window}); and the value at $0$ of the
derivative of the phase $\vartheta(t)=\tfrac t2\log(q/\pi)+\Im\log\Gamma(\tfrac14+\tfrac{it}2)$ of the gamma factor
of $\Lambda(\tfrac12+it,\chiv)$, which gives the smooth part of the counting function: $\vartheta'(0)=\kappa/2$, the same $\kappa$
(\tier[four\_equalities\_on\_the\_set]{P}\ given a classical input).  The first is the multiplicative half of \Ecore, proved for every zero; the additive half, $r+r'=2$, is what remains --- the one-sidedness of \cref{rmk:nodes} seen in the pair $\{r,r'\}$.

\begin{remark}[An earlier Weil window, and how it compares]\label{rmk:earlier-window}
The corpus also records a Weil margin $Q(g_L)\ge0.0443$ for $2L<\log64$
(\LeanInline{extended\_window} \tier{N}), computed on a one-parameter family $g_L$ of
test functions, together with a count of the prime terms by sign: up to $64$,
twelve enter against positivity and thirteen in its favour (a count made outside
Lean \tier{N}; \LeanInline{adverse\_are\_outnumbered} checks only $12<13$), with $\chiv(4)=+1$ adverse and
$\chiv(8)=-1$ favourable (\tier[four\_is\_adverse, eight\_is\_favourable]{P}).  The window theorems of
\cref{ssec:weil-route} are of a different strength: they bound the Weil form over
\emph{all} test functions supported in the window, which is what Weil's criterion
quantifies over, while a single family gives a necessary condition only.  The sign
count is the same phenomenon as \cref{prop:semilocal}: inert primes, with
$\Frob_p=\sigma$, enter in favour of positivity.  Neither kind of statement proves
a portion of \Ecore, since a window of finite length does not constrain an
individual zero; the corpus itself records this correction to the reading of its
margin.
\end{remark}

\subsubsection{The zeros select the conductor}\label{ssec:envelope-fluctuation}

Let $N(T)$ be the number of zeros of $L(s,\chiv)$ with $0<\gamma\le T$.  The corpus writes it as
$N(T)=\mathrm{env}(T)+S(T)$ (\LeanInline{CountingLaw}), with the \emph{envelope}
$\mathrm{env}(T)=\frac T{2\pi}\log\frac{qT}{2\pi e}$ for conductor $q$ (\LeanInline{envelope}) and the
\emph{fluctuation} $S(T)$ the remainder, and splits the logarithm of the envelope into its arithmetic and universal
parts, $\log q+\log T-\log(2\pi e)$ (\tier[envelope\_splits]{P}), with only
the first distinguishing the field (\tier[field\_part\_distinguishes]{P}).  For
$\chiv$ the arithmetic part is $\log5$, fixed by the conductor
(\tier[conductor\_is\_derived]{P}).  This locates exactly what an operator
built from the structure can and cannot reproduce.

\begin{proposition}[The zeros select the conductor]\label{prop:envelope-exact}
The local spacing scale $2\pi/\log(qT/2\pi)$ is injective in the conductor $q$ at
every height (\tier[scale\_injective\_in\_conductor]{P}), and only the
arithmetic part $\log q$ of the envelope distinguishes conductors
(\tier[envelope\_splits, field\_part\_distinguishes]{P}).
Consequently ``unfolded spacings have mean one'' is a statement about the zeros,
not a choice of units: unfolding the $40$ zeros of $L(s,\chiv)$ below $80$ with the
envelope of conductor $q$ gives the mean spacings
\[
\begin{array}{c|ccccccccc}
q&1&2&3&4&\mathbf5&6&8&10&20\\\hline
\text{mean}&0.528&0.734&0.855&0.941&\mathbf{1.008}&1.062&1.148&1.215&1.422
\end{array}
\]
and the value closest to $1$ is $1.008$, at $q=5$ only (\tier{N},
\texttt{unfold\_by\_conductor.txt}).  The conductor, derived in % pending(D7)
\cref{sec:choice} as the $5$-component of $\operatorname{lcm}(4,5)$, is read back
from the zeros.
\end{proposition}

\begin{proposition}[Envelope, fluctuation, and the error of an operator]\label{prop:error-partition}
Let $N$ be the zero-counting function, $N=\mathrm{env}+S$ its splitting into the envelope
and the fluctuation, $E=\mathrm{env}$ and $E^{-1}$ its exact inverse, and $t_n$ the height of the
$n$-th zero.
\begin{enumerate}[label=\textup{(\roman*)}]
  \item $t_n=E^{-1}\bigl(n-S(t_n)\bigr)$ exactly.
  \item For any operator with eigenvalues $T(n)$,
    $t_n-T(n)=\bigl[E^{-1}(n-S(t_n))-E^{-1}(n)\bigr]+\bigl[E^{-1}(n)-T(n)\bigr]$:
    fluctuation plus parametrisation.
  \item The parametrisation term vanishes for $T=E^{-1}$.
  \item The fluctuation term is carried by no operator whose eigenvalues depend on the
    index alone: two sequences of heights that differ at $n$ cannot both be matched there.
\end{enumerate}
{\raggedright \tier[F1PaperHyp.height\_from\_envelope, F1PaperHyp.error\_decomposition, F1PaperHyp.parametrisation\_removable, F1PaperHyp.no\_index\_operator\_fits\_both]{P}, on the counting law of the corpus (\code{CountingLaw})\par}
\end{proposition}

\begin{remark}[What is removable, and what is not]\label{rmk:two-errors}
For $\zeta$, the operator of \cite{gonzalez25hp} has eigenvalues
$T^*(n)=c(n)\pi n/\log(\mu n+\sigma)$ with $c(n)=2+\alpha_0/(1+\beta_0\log n)$.  By
\cref{prop:error-partition} its error is fluctuation plus parametrisation, and only the
second is removable, and only by the exact inverse $E^{-1}$, which is obtained by solving
the counting law numerically and is not a closed form.  Among closed forms the structural
one is the better: \cite{gonzalez25hp} shows that its relative error stays below $1.7\%$
from $n=1$ to $n=10^{12}$, peaks near $n\approx1.5\cdot10^9$ and decreases monotonically
after the peak, with $T^*(n)/t_n\to1$, and that replacing $\log(\mu n+\sigma)$ by the
subleading terms of the Riemann--von Mangoldt inversion increases the error by more than an
order of magnitude.  What no operator built from the envelope carries is the fluctuation
$S$; what the envelope does carry is the conductor, and \cref{prop:envelope-exact} shows
that the zeros select $q=5$.
\end{remark}

\begin{proposition}[The envelope on and off the line; the two separations]\label{prop:envelope-edge}
\begin{enumerate}[label=\textup{(\roman*)}]
  \item The envelope splits into a conductor part and a universal part, and the conductor
    part separates conductors (\cref{prop:envelope-exact}).
  \item The second envelope is that of the Fej\'er kernel of window length $L>0$ over the Klein orbit of a zero
    $\rho$ at horizontal distance $\delta=\Re\rho-\tfrac12$ from the line,
    $\mathrm{env}(L,\delta)=2\cosh^2(L\delta)/(L|\rho|^2)$ (\LeanInline{env}; the Fej\'er form is that of
    \cref{ssec:field-positivity}).  Weighted by $L$ it is
    $L\cdot\mathrm{env}(L,\delta)=2\cosh^2(L\delta)/|\rho|^2$: on the line it is the
    constant $2/|\rho|^2$, independent of $L$, and off it it grows at least like
    $e^{2L\delta}/(2|\rho|^2)$.  The line is the edge of the envelope.
  \item The Klein partner of a zero off the line has the same ordinate and real part
    $1-\beta$: the vertical separation of the pair is always $0$, and the radial
    separation is positive exactly off the line.
  \item If the excess $(r-1)+(1/r-1)$ of a coordinate $r\le2$ is at most $B$, then
    $(r-1)^2\le2B$.
\end{enumerate}
\tier[off\_line\_L\_env, on\_line\_L\_env\_constant, on\_line\_no\_growth, off\_line\_grows, the\_line\_is\_the\_edge, free\_orbit\_doubles\_ordinate, doubled\_ordinate\_reflects, separations\_are\_complementary, deviation\_bounded]{P}%
\end{proposition}

% ---------------------------------------------------------------------
The spacings of the zeros test \cref{prop:symmetry-spectrum} directly.  Unfolded by their counting functions, the
$240$ zeros of $L(s,\chiv)$ and the $156$ zeros of $\zeta$ below height $330$ have no spacing below $0.1$: the smallest are
$0.231$ and $0.372$.  For the Gaussian unitary ensemble the expected fraction of such spacings is about $10^{-3}$; for two
independent sectors it is about $0.047$, and for six about $0.079$, which would put some eighteen of them among the
$233$ spacings of $L(s,\chiv)$.  The zeros of $\zeta_K=\zeta\cdot L(s,\chiv)$, the union of the two sets, have $11$
spacings below $0.1$ out of $388$, the smallest $0.017$, and no ordinate is shared by the two factors in this range.  So
the Galois symmetry of order two that splits $\zeta_K$ into its two factors (\cref{eq:dedekind}) shows in its zeros as
two sectors, as \cref{prop:symmetry-spectrum}\textup{(i)} says a commuting symmetry does; and neither $L(s,\chiv)$ nor
$\zeta$ shows a sector, a doublet, or the multiplicity six of an operator of the level (\texttt{chi5\_zeros.py}, % pending(D7)
\texttt{retro105.py}).  If the zeros of $L(s,\chiv)$ are the spectrum of a self-adjoint operator, the data % pending(D7)
then make it neither an operator of the level nor an operator with a commuting symmetry of order two or six whose
sectors are independent: the fibre, if it acts on it, acts through an operator that does not commute with its symmetry, as in
\cref{prop:level-operators}\textup{(iv)}.

\subsection{The Weil analogy over $\Z$}\label{sec:results-weil}

\subsubsection{The Weil analogy, and where it stops}\label{ssec:weil-analogy}

Three proved facts line up with the ingredients of Weil's proof
(\cref{ssec:weil-template}): the group $\mathbb{Z}_2\times\mathbb{Z}_2$ has
exponent two; the only eigenvalue of $\OmH$ with positive real part is
$\tfrac12$, the analogue of $|\alpha_i|=q^{1/2}$; and $|\ker G|=2$
(\tier[weil\_analogy, golden\_group\_exponent\_two, positive\_spectral\_half, G\_kernel]{P}).  The corpus reads the
last as $2=2-2g$ with $g=0$ (\LeanInline{G\_genus\_zero}); the proved content is the
cardinal\tier{P}, and the value of $g$ is supplied by the statement.

The analogy must be read with its limit.  For a curve of genus $0$ the first
cohomology vanishes and the zeta function has no zeros; $L(s,\chiv)$ has
infinitely many.  Deninger's programme expects the analogue of $H^1$ for
$\mathrm{Spec}\,\mathbb{Z}$ to be infinite dimensional --- the cohomology of a foliated
dynamical system on which the zeros are the spectrum of an infinitesimal
generator~\cite{deninger98} --- and the positivity of Weil's proof to be a statement
about it.  The facts above describe the finite part of the structure; the zeros are
what remains.  \Cref{ssec:weil-closure} uses the finite part as Weil uses $H^1$ only through the orbit of
Frobenius: the eight classes of $(\Z/20)^\times$ give a finite lattice with a Frobenius of
order dividing $4$, and the infinite-dimensional positivity is needed only on the orbits,
where it becomes one bound per class.

\subsubsection{$G$ as the $\mathbb{F}_1$-side analogue of the diagonal}
\label{ssec:G-diagonal}

The arithmetic surface $C\times C$ provides Weil with a geometric
arena where Frobenius acts and intersection theory makes sense.
Over~$\mathbb{Z}$, the analogue
$\mathrm{Spec}(\mathbb{Z})\times_{\mathbb{F}_1}\mathrm{Spec}(\mathbb{Z})$
does not exist as a conventional scheme, and its construction is
one of the central open problems of the
$\mathbb{F}_1$-programme~\cite{connesconsani10}.  In the PCF span,
the structural role of this inaccessible object is played by the
common source $\mathbf{B}\Zst$, with the homomorphism
$G\colon\Zst\to\mathbb{Z}_2\times\mathbb{Z}_2$ carrying the
arithmetic content that Weil's diagonal carries over function
fields.

\begin{theorem}[$G$ as the common source]\label{thm:G-source}
The homomorphism
$G\colon\Zst\to\mathbb{Z}_2\times\mathbb{Z}_2$ is surjective
with $|\ker G|=2$, and produces simultaneously:
\begin{enumerate}[label=\textup{(\alph*)}]
\item The local Euler factors of $\zeta_{\mathbb{Q}(\sqrt{5})}$
  via the second-coordinate projection
  $\pi_2\circ G\colon\Zst\to\mathbb{Z}_2$ identified with
  $\chi_5$ \textup{(\cref{thm:s1})};
\item The spectral parameter $\mu_3=1/|\ker G|=1/2$ \tier[\testref{test\_sigma\_mu\_polynomial\_root\_three\_halves, test\_sigma\_mu\_polynomial\_root\_one\_half}]{N}, the positive eigenvalue of
  $\OmH$, with $(\sigma,\mu)=(3/2,1/2)$ the unique admissible pair
  \textup{(\cref{cor:no-diagonal-spectral})}.
\end{enumerate}
The four characters $\{1,\chi_4,\chi_5,\chi_4\chi_5\}$ of
$\mathbb{Z}_2\times\mathbb{Z}_2$ classify the four quadratic
extensions $\mathbb{Q}, \mathbb{Q}(i), \mathbb{Q}(\sqrt{5}),
\mathbb{Q}(\sqrt{-5})$ \textup{(\cref{thm:dedekind-natural})}.
\tier[G\_hom, G\_surj, G\_pi2\_is\_chi5]{P}; recorded as
\LeanTally{the\_span\_is\_the\_identification}
\end{theorem}

The structural content of \cref{thm:G-source} is that $G$ is the
hinge.  A single finite homomorphism produces, on one side, the
arithmetic input (Euler factors via $\chi_5$), and on the other
side the spectral output (the eigenvalues of $\OmH$ on the circle $|z|=\tfrac12$).
The two outputs are not derived from each other; they share a common source.  That is
the role the diagonal plays in Weil's proof, where both the graph of Frobenius and the
two projections of $C\times C$ are measured against it; here the analogy is one of
roles, not of objects.

\begin{proposition}[Characterisation of $G$ by its projections]
\label{prop:G-characterisation}
The homomorphism $G\colon\Zst\to\mathbb{Z}_2\times\mathbb{Z}_2$
is uniquely characterised on $\Zst$ by its projection structure:
for every $a\in\Zst$,
\begin{align*}
  (G(a))_1 = 0 \quad &\Longleftrightarrow\quad
    a \in \{1,\,9,\,13,\,17\}, \\
  (G(a))_2 = 0 \quad &\Longleftrightarrow\quad
    a \in \{1,\,9,\,11,\,19\},
\end{align*}
and any map $H\colon\Zst\to\mathbb{Z}_2\times\mathbb{Z}_2$
satisfying these two biconditionals coincides with $G$ on $\Zst$.  The two sets are the kernels of
$\chi_4$ and of $\chiv$.
\code{G\_diagonal\_characterisation}, \code{G\_diagonal\_unique} \tier{P}%
\end{proposition}
\begin{proof}
The biconditionals are checked on the eight elements $1,3,7,9,11,13,17,19$ of $\Zst$ and on
the four values of $G$.  The uniqueness holds because $\mathbb{Z}_2$ has only two elements: each component of
$H(a)$ is determined by whether $a$ lies in the corresponding subset, and $G$ satisfies the same conditions, so the
components agree.
\end{proof}

The characterisation sits beside the multiplicativity (\code{G\_hom}), the surjectivity
(\code{G\_surj}) and the kernel structure $|\ker G| = 2$ (\code{G\_kernel}), which the
corpus writes as $|\ker G| = 2 - 2g$ at $g = 0$ (\code{G\_genus\_zero}), an equality of
integers rather than a genus computation, as \cref{ssec:weil-analogy} records.  It is elementary --- a map to
$\mathbb{Z}_2\times\mathbb{Z}_2$ is fixed by the zero sets of its two components --- and
its interest is that the two zero sets are the kernels of $\chi_4$ and $\chiv$.

\subsubsection{Correspondence between Weil and PCF}\label{ssec:correspondence}

The table sets Weil's proof for function fields beside the structure of this paper,
role by role.  It is an analogy: the right-hand column does not reproduce Weil's
proof, and the entry that would carry the positivity is exactly the one that is not
proved.
\begin{center}
\small
\begin{tabular}{@{}>{\raggedright\arraybackslash}p{0.30\textwidth}>{\raggedright\arraybackslash}p{0.42\textwidth}>{\raggedright\arraybackslash}p{0.22\textwidth}@{}}
\toprule
Weil over $\mathbb{F}_q$ & PCF over $\mathbb{Z}$ & Reference \\
\midrule
$C\times C$ (arithmetic surface)
  & PCF span with common source $G$
  & \cref{prop:span} \\
Horizontal $C\times\{P\}$
  & partial Euler products $E\colon\mathcal P\to\mathbf{DirSer}$
  & \cref{prop:euler-colimit} \\
Vertical $\{P\}\times C$
  & $F_{\mathrm{spec}}\to\mathrm{Spec}(\OmH)$
  & \cref{def:SpecOmega} \\
Diagonal $\Delta$
  & Common source $\mathbf{B}\Zst$ via $G$
  & \cref{thm:G-source} \\
$\mathrm{Fr}\colon x\mapsto x^q$
  & Frobenius lifts $\Frob_p(\varphi^n) = \varphi^{pn}$
  & \cref{def:frobenius-functor} \\
Intersection positivity bound
  & Premise~A: $\Re\rho\le\tfrac12$ for all zeros \emph{(open, $\Leftrightarrow$ GRH$(\chiv)$)}
  & \cref{prop:premiseA} \\
Functional equation bound
  & Premise~B: Hecke 1920 self-duality
  & \cref{sec:analytic-inputs} \\
Castelnuovo--Severi from the Hodge index
  & \textup{(W4)} from \textup{(W3)} on the planes of the orbit
  & \cref{prop:W4-from-W3} \\
An eigenvalue off the circle
  & A free Klein orbit: index $2$, one even and one odd direction
  & \cref{prop:free-orbit-index} \\
Squeeze: $|\alpha_i|=q^{1/2}$
  & Squeeze: $\mathrm{Re}(\rho)=1/2$
  & \cref{prop:squeeze} \\
Eigenvalues on $H^1_{\mathrm{\acute{e}t}}$
  & Eigenvalues $\zeta_3^k/2$ on $|z|=1/2$
  & \cref{ssec:omega} \\
\bottomrule
\end{tabular}
\end{center}

Where the analogy stops is the positivity row.  Weil's positivity comes from
intersection theory on an algebraic surface and is a theorem; Premise~A is equivalent
to GRH$(\chiv)$ (\cref{prop:premiseA}) and is not proved, and nothing in the categorical
chain or in \cref{lem:diagonal-contraction} supplies it.  What does transfer is the
other half: the functional equation gives the companion bound, as in Weil's setting,
and with it the squeeze closes the moment Premise~A is granted (\cref{prop:squeeze}).

\begin{proposition}[Split and inert as orbits of the Klein action on $\Zst$]\label{prop:split-inert-orbits}
On $\Zst$ the Klein group acts by the two commuting involutions $a\mapsto a^{-1}$ and
$a\mapsto-a$.  The fixed points of inversion are the $2$-torsion $\{1,9,11,19\}$, which
are the split classes; the free orbit $\{3,7,13,17\}$ consists of the elements of order
exactly four, which are the inert classes; and $\chiv$ is trivial on squares.
\tier[klein\_acts\_on\_Z20\_by\_two\_involutions, fixed\_points\_are\_two\_torsion, free\_orbit\_eq\_order\_four, split\_is\_the\_degenerate\_orbits, inert\_is\_the\_free\_orbit, chi5\_trivial\_on\_squares]{P}%
\end{proposition}

In the language of the analogy: the Frobenius of a split prime is fixed by the Klein
symmetry, and that of an inert prime moves in a free orbit of order four --- the same
dichotomy, degenerate orbit on the line and free orbit off it, that
\cref{prop:orbits} states for the zeros.

% ---------------------------------------------------------------------
\subsection{The Weil route: positivity in windows}\label{ssec:weil-route}

The open statement has been written pointwise, zero by zero, in the forms of
\cref{ssec:core-forms}.  The only operation on the coordinates of the zeros that the
corpus proves is the inversion $r\mapsto1/r$: for every zero $\rho$, $1-\rho$ is a zero and
$r(\rho)\,r(1-\rho)=1$ (\tier[the\_available\_mechanisms\_are\_of\_order\_two]{P}); the
descent form of \cref{ssec:core-forms} asks for the powers $r\mapsto r^p$.  The only proved bridge from the primes to
all the zeros at once is the explicit formula, and in it the primes enter
through the trace of the Frobenius of $\Zphi$.



\begin{theorem}[Explicit formula, with Frobenius traces]\label{thm:explicit}
For $g$ even, smooth, of compact support, with $h(r)=\int g(u)e^{iru}\,du$,
\[
  \sum_{\gamma}h(\gamma)=\frac1{2\pi}\int h(r)\Bigl[\log\frac5\pi
  +\mathrm{Re}\,\psi\bigl(\tfrac14+\tfrac{ir}2\bigr)\Bigr]dr
  -2\sum_{p\ne5}\sum_{k\ge1}\frac{\log p}{p^{k/2}}\,
  \mathrm{tr}\bigl(\Frob_p^{\,k}\mid\sqrt5\bigr)\,g(k\log p),
\]
the sum running over the non-trivial zeros $\rho$ of $L(s,\chiv)$ with
$\gamma=(\rho-\tfrac12)/i$.  We write the right-hand side as $A(g)-2P(g)$: $A(g)$ is the archimedean term, the
integral, and $P(g)=\sum_{p\ne5}\sum_{k\ge1}\frac{\log p}{p^{k/2}}\mathrm{tr}(\Frob_p^{\,k}\mid\sqrt5)\,g(k\log p)$ the prime
side.  Here $\gamma$ is real exactly when $\rho$ is on the line; for
complex $\gamma$, $h(\gamma)$ is the value of the entire extension of $h$.
The formula is classical (, \cite{weil52}; \cite[Thm.~5.12]{iwaniec04});
that its prime coefficients are Frobenius traces is \tier{P}%
(\LeanInline{prime\_side\_is\_frobenius\_trace}).  $\Lambda(s,\chiv)$ is entire, so
there are no pole terms.  The normalisation is checked numerically: with the
$40$ zeros of height $\le80$ (first $\gamma=6.6484533447$) and Gaussian $g$,
both sides agree to $10^{-18}$ (\texttt{verify\_weil\_chi5.py}). % pending(D7)
\end{theorem}

\begin{theorem}[The core as a prime-side bound]\label{thm:core-positivity}
For $f\in L^2(\R)$ with compact support let $\tilde f(x)=\overline{f(-x)}$, so that $g=f*\tilde f$ is even with
$h(r)=|\hat f(r)|^2\ge0$ on $\R$, and let the \emph{Weil form} be the zero side $Q(f)=\sum_\gamma h(\gamma)$ of
\cref{thm:explicit} at $g=f*\tilde f$.  Given \cref{thm:explicit}, Weil positivity for $L(s,\chiv)$ --- $Q(f)\ge0$
for every such $f$ --- is equivalent to
$2P(f*\tilde f)\le A(f*\tilde f)$ for all $f$: the Frobenius-trace side bounded
by the archimedean one.
\tier[positivity\_is\_a\_prime\_side\_bound]{P}; positivity $\Leftrightarrow$
GRH$(\chiv)$ is Weil's criterion; the bound itself is open.
\end{theorem}

\begin{remark}
If all ordinates are real the zero side is a sum of $|F(\gamma)|^2$ and is
non-negative (\tier[zero\_side\_nonneg\_of\_real\_ordinates]{P}): that is
the easy half.  Given \Ecore, \tier[position\_is\_interval\_meets\_discreteness, radius\_one\_is\_E]{P}
places every zero at $k=0$ of the additive monoid.
\end{remark}

\begin{remark}[What the numerics say about the bound]\label{rmk:saturation}
Split the prime side by the Frobenius: $P=P_{\mathrm{split}}-P_{\mathrm{inert}}$,
the primes with $\Frob_p=\mathrm{id}$ against those with $\Frob_p=\sigma$
(\tier[prime\_side\_splits\_by\_frobenius]{P}).  For Gaussian test
functions of width $w\in[0.2,1]$ the inert primes dominate ($2$ and $3$
are inert), the archimedean term becomes negative, and $2P$ follows $A$ ever
more closely: the margin $A-2P$ equals the zero side and falls from $0.63$ to
$1.3\cdot10^{-9}$, at the rate $e^{-w^2\gamma_1^2/2}$ set by the first zero
(\tier{N}).  The inequality of \cref{thm:core-positivity} is therefore not a gross
estimate but a near-cancellation between the Frobenius traces and the
archimedean place; no bound on $P$ that loses a constant can prove it.  What
could prove it is an exact positivity --- a positive form on a space carrying
all the $\Frob_p$ at once, as the intersection form does in Weil's proof.
That space is the open part of the route.
\end{remark}





\begin{proposition}[Semilocal positivity and the Frobenius at $2$, numerical]\label{prop:semilocal}
Restrict the Weil form of \cref{thm:core-positivity} to test functions
supported in $[-L/2,L/2]$ and let $\lambda_{\min}(L)$ be its smallest
eigenvalue.  Then (\texttt{weil\_window\_chi5.py}): % pending(D7)
\begin{enumerate}
  \item for $L<\log2$ no prime enters and the form is the archimedean one;
    archimedean positivity alone holds up to $L_c=0.9251$ and fails beyond;
  \item with the primes, $\lambda_{\min}(L)>0$ on every window computed, but it
    decays super-exponentially ($0.097$ at $L=1$, $5.3\cdot10^{-3}$ at $1.4$,
    $\sim10^{-6}$ at $2$; \texttt{weil\_window\_chi5.py}); % pending(D7)
  \item if the Frobenius at $2$ is replaced by the identity ($\chiv(2)=+1$),
    positivity fails from $L\approx0.9$.
\end{enumerate}
The first prime to enter, $2$, is inert: $\Frob_2=\sigma$.  Its entry at
$\log2<L_c$ is what carries positivity past the archimedean threshold.
\end{proposition}

\begin{theorem}[Weil positivity for $L(s,\chiv)$ on the first window]\label{thm:first-window}
Let $\kappa=\log\frac5\pi-\gamma_E-\frac\pi2-3\log2=-3.76274\ldots$
$(=\log\frac5\pi+\psi(\tfrac14))$ and
$K(d)=2\bigl(\operatorname{artanh}e^{-d/2}+\arctan e^{-d/2}\bigr)$.
Let $L^*=0.744725\ldots$ be the root of $\kappa+K(L/2)=0$.  For every
$f\in L^2(\mathbb{R})$ supported in an interval of length $L\le L^*$, the Weil
form of $L(s,\chiv)$ satisfies
\[
  Q(f)\;\ge\;m(L)\,\|f\|^2,\qquad
  m(L)=\min\Bigl(\kappa+K(\tfrac L2),\ \kappa+\tfrac12\bigl(K(L-\log2)+K(L)\bigr)-\tfrac{\log2}{\sqrt2}\Bigr)\;\ge0,
\]
the second term only when $L>\log2$.  At $L=\log2$, $m=0.0841$.
\tier[F1v2.first\_window\_integral, F1v2.dirichlet\_localisation, F1v2.midpoint\_bound, F1v2.potential\_integral\_bound, F1v2.absorption]{P}, given the explicit formula with Gauss's integral
for the digamma function and the numerical values of the potential
(\texttt{weil\_window\_rigorous\_check.py}); the value $\psi(\tfrac12)=-\gamma_E-2\log2$ it % pending(D7)
uses is proved (\tier[digamma\_half\_is\_negative]{P})
\end{theorem}

\begin{proof}
We take $f$ with $\int|\hat f|^2\log(2+|r|)\,dr<\infty$; otherwise $Q(f)=+\infty$ and there is nothing to prove.
\emph{The archimedean term as a Dirichlet form.}  From
$\psi(z)=-\gamma_E+\int_0^\infty\frac{e^{-t}-e^{-zt}}{1-e^{-t}}\,dt$ the archimedean
term of \cref{thm:explicit} is, for $g=f*\tilde f$,
\[
  A(g)=g(0)\bigl(\log\tfrac5\pi-\gamma_E\bigr)
  +2\!\int_0^\infty\!\frac{e^{-2u}g(0)-e^{-u/2}g(u)}{1-e^{-2u}}\,du .
\]
Since $2\int_0^\infty\frac{e^{-u/2}-e^{-2u}}{1-e^{-2u}}du=\psi(1)-\psi(\frac14)=\frac\pi2+3\log2$,
with $k(u)=\frac{2e^{-u/2}}{1-e^{-2u}}>0$,
\[
  A(g)=\kappa\,g(0)+\int_0^\infty k(u)\bigl(g(0)-g(u)\bigr)\,du
      =\kappa\|f\|^2+\tfrac14\iint_{\mathbb{R}^2}k(|x-y|)\bigl(f(x)-f(y)\bigr)^2dx\,dy .
\]
\emph{Localisation.}  Let $I$ be the support interval.  Dropping the part of the
double integral over $I\times I$ (non-negative) and keeping $I\times I^c$ and
$I^c\times I$, where $f=0$ on $I^c$,
\[
  A(g)\ \ge\ \int_I f(x)^2\,V(x)\,dx,\qquad
  V(x)=\kappa+\tfrac12\bigl(K(d_-(x))+K(d_+(x))\bigr),
\]
with $d_\pm(x)$ the distances from $x$ to the ends of $I$ and
$K(d)=\int_d^\infty k=2(\operatorname{artanh}e^{-d/2}+\arctan e^{-d/2})$.  $K'=-k$ with $k$ decreasing, so $K$ is
convex and decreasing and $V$ is minimal at the midpoint, $V\ge\kappa+K(L/2)$.
\emph{Primes.}  $g$ vanishes outside $[-L,L]$.  For $L\le\log2$ no prime term
enters and $Q=A\ge(\kappa+K(L/2))\|f\|^2$.  For $\log2<L\le L^*<\log3$ only $n=2$
enters, with $-2\Lambda(2)\chiv(2)2^{-1/2}g(\log2)=\sqrt2\log2\,g(\log2)$.  With
$\delta=L-\log2$, $g(\log2)=\int f(x)f(x-\log2)dx$ involves only the two end
strips $S$ of width $\delta$, so $|g(\log2)|\le\frac12\int_S f^2$, and on $S$,
$V\ge\kappa+\frac12(K(\delta)+K(L))$.  Hence
$Q\ge\int_{I\setminus S}f^2V+\int_S f^2\bigl(V-\frac{\log2}{\sqrt2}\bigr)\ge m(L)\|f\|^2$.
\end{proof}

\begin{theorem}[Positivity up to $L=1$, through the Frobenius at $2$]\label{thm:second-window}
For every $f\in L^2(\mathbb{R})$ supported in an interval of length $L\le1$,
the Weil form of $L(s,\chiv)$ satisfies $Q(f)\ge0.0445\,\|f\|^2$.  Computer-assisted in interval arithmetic, with outward rounding (\cref{rmk:rigour}); certificate \texttt{weil\_certificates/cert\_L1.0.json}, replay \texttt{verify\_iv.py}. % pending(D7)
\end{theorem}

\begin{proof}
\emph{Reduction to a pointwise inequality.}  With $g=f*\tilde f$ and
$h=|\hat f|^2$, $Q(f)=\frac1{2\pi}\int|\hat f(r)|^2\,\Phi(r)\,dr$ for
$\Phi=W+\sum_n a_n\cos(r\log n)$, where $W(r)=\log\frac5\pi+\mathrm{Re}\,\psi(\frac14+\frac{ir}2)$ is the archimedean
weight, $a_n=-2\Lambda(n)\chiv(n)n^{-1/2}$ and $\Lambda(n)$ is the von Mangoldt function.  Write $P_L$ for the part of
the sum with $\log n<L$.  Since
$g$ is continuous and vanishes outside $(-L,L)$, adding any
$E(r)=\sum_i b_i\cos(\lambda_i r)$ with all $\lambda_i\ge L$ does not change
$Q$, and neither do the terms with $\log n\ge L$.  Hence, if
$W+P_L+E\ge t$ on $\mathbb{R}$, then $Q(f)\ge t\|f\|^2$ by Plancherel.
\emph{The certificate.}  A linear programme gives $41$ rational pairs
$(\lambda_i,b_i)$ with $\lambda_i\in[1,16]$.  $W$
is even and increasing on $[0,\infty)$, because
$\partial_y\mathrm{Re}\,\psi(x+iy)=\sum_{k\ge0}\frac{2y(k+x)}{((k+x)^2+y^2)^2}>0$.  The
trigonometric part $T=P_L+E$ is Lipschitz with constant
$M_T=\sum|a_n|\log n+\sum|b_i|\lambda_i=11.0$.  \emph{The bound, in interval arithmetic.}  Below $W$, every partial sum of the series is a lower bound, and the tail
$\sum_{j\ge K}f_j$, with $f_j=\frac{y^2}{(j+x)((j+x)^2+y^2)}$ decreasing in $j$, is at least $\int_K^\infty=\frac12\log\bigl(1+\frac{y^2}{(K+x)^2}\bigr)$;
$\psi(\tfrac14)=-\gamma_E-\frac\pi2-3\log2$ (Gauss's digamma theorem~\cite{dlmf}).  Since $|f_j''|\le2/(j+x)^3$, $|W''|\le M_W=\tfrac12\bigl(64+1.25^{-3}+\tfrac12\,1.25^{-2}\bigr)$, and
$|T''|\le M_2=\sum|a_n|\log^2n+\sum|b_i|\lambda_i^2$; with $M=M_W+M_2=67.92$, on each cell of a grid of step
$h=\tfrac1{100}$ over $[0,600]$, $\Phi(r)\ge\min(\Phi(r_k),\Phi(r_{k+1}))-Mh^2/8$.  Evaluated in interval arithmetic with
outward rounding, the minimum is $\ge0.04458$, and for $r\ge600$, $\Phi(r)\ge W(600)-\sum|a_n|-\sum|b_i|\ge0.04999$.
The earlier double-precision replay (\texttt{verify\_rig.py}, Lipschitz cells) gave $0.0431$. % pending(D7)
\end{proof}

\begin{theorem}[Positivity up to $L=1.3$, operator form]\label{thm:third-window}
For every $f\in L^2(\mathbb{R})$ supported in an interval of length $L\le1.3$,
$Q(f)\ge0.0027\,\|f\|^2$; at $L\le1.2$ the margin is $0.0077$ and at $L\le1.1$,
$0.0261$.  Computer-assisted with validated numerics: every rounding error bounded explicitly and the top eigenvalue bounded by a verified Cholesky factorisation (\cref{rmk:rigour}); \texttt{weil\_certificates/op\_rig\_validated.py}. % pending(D7)
\end{theorem}

\begin{proof}
Fix $t>0$ and write $\Psi=W+P_L$, so $Q(f)=\frac1{2\pi}\int|\hat f|^2\Psi$.
Since $W$ is increasing on $[0,\infty)$ and $P_L$ is Lipschitz with constant
$M=\sum|a_n|\log n$, a grid of step $h$ gives a piecewise-constant
$N\ge(t-\Psi)_+$: on the cell at $r_k$, $N_k=\bigl(t-W(r_k)-P_L(r_k)+Mh\bigr)_+$,
and $N=0$ beyond the radius $R$ where $W(R)\ge t+\sum|a_n|$.  Then
$\Psi\ge t-N$ and $Q(f)\ge t\|f\|^2-\frac1{2\pi}\int|\hat f|^2N$.
For the last term, let $\varphi_j$ be the sine basis of $L^2(I)$ and split
$f=f_J+f_\perp$.  With $e_r(x)=e^{irx}\mathbf{1}_I$, $|\hat f_\perp(r)|\le
\varepsilon_J(r)\|f_\perp\|$ where $\varepsilon_J(r)^2=L-\sum_{j\le J}|\hat\varphi_j(r)|^2$,
and $\frac1{2\pi}\int N|\hat f_J|^2\le\lambda_J\|f_J\|^2$ with $\lambda_J$ the
largest eigenvalue of the $J\times J$ matrix
$G_{jk}=\frac1{2\pi}\int N\hat\varphi_j\overline{\hat\varphi_k}$.  By
Cauchy--Schwarz, $\frac1{2\pi}\int N|\hat f|^2\le(\lambda_J+\eta_J)\|f\|^2$ with
$\eta_J=\frac1{2\pi}\int N\varepsilon_J^2$, so $Q\ge(t-\lambda_J-\eta_J)\|f\|^2$.
Computation at $L=1.3$, $t=2$, $J=800$, $h=2.5\cdot10^{-4}$ gives
$\lambda_J\le1.993719$, $\eta_J=3.45\cdot10^{-3}$ and quadrature errors below $2.5\cdot10^{-5}$, hence the margin $0.00277$ (\cref{rmk:rigour}); $R=88$.  The bound on $\lambda_J$ is verified: a Cholesky factorisation of $\mu I-G$ with $\mu=\tilde\lambda_J+10^{-8}$ completes in floating point, so $\mu I-G+\Delta$ is positive semidefinite with $\|\Delta\|_2\le\gamma_{J+1}(1+\gamma_{J+1})\operatorname{tr}(\mu I-G)$ \cite{higham02}, below $1.5\cdot10^{-10}$; at $L=1.2$ and $1.1$ the validated margins are $0.0077$ and $0.0261$.  The $\hat\varphi_j$ are elementary,
$\hat\varphi_j(r)=\sqrt{2/L}\,e^{-irL/2}k_j\bigl(1-(-1)^je^{irL}\bigr)/(k_j^2-r^2)$,
$k_j=j\pi/L$.
\end{proof}

\begin{remark}[How the certificates are made rigorous]\label{rmk:rigour}
Three sources of error are controlled explicitly.  \emph{The archimedean weight}:
from $\Re\psi(x+iy)=\psi(x)+\sum_{k\ge0}\frac{y^2}{(k+x)((k+x)^2+y^2)}$, all of whose
terms are positive, every partial sum is a lower bound for $W$, with
$\psi(\tfrac14)=-\gamma_E-\frac\pi2-3\log2$; so $W$ is never evaluated by a
library routine, only bounded below (\texttt{rigW.py}), and its monotonicity lets % pending(D7)
a coarse grid of such bounds serve every fine cell.  \emph{Discretisation}: the
trigonometric part enters through its Lipschitz constant, and the midpoint quadrature of $G$ and $\eta_J$ through the bound $L^3$ on the second derivatives of $\hat\varphi_j\overline{\hat\varphi_k}$ and $|\hat\varphi_j|^2$ --- each is $\iint\varphi_j(x)\varphi_k(y)e^{-ir(x-y)}$, with $|x-y|\le L$ and $\|\varphi_j\|_1^2\le L$ --- which give errors below $4\cdot10^{-5}$.  \emph{Floating point}: an explicit margin of $10^{-9}$ per
evaluated quantity is subtracted.  The pointwise certificate at $L=1$ is replayed in interval arithmetic with outward rounding (\texttt{verify\_iv.py}), which gives $0.0445$; no floating-point step remains in it.  The operator certificates are replayed with validated numerics (\texttt{op\_rig\_validated.py}): the accumulation of $G$, the evaluation of the $\hat\varphi_j$ (relative error below $10^{-9}$, the grid staying at least $2.8\cdot10^{-6}$ from every $k_j$) and the symmetrisation are bounded explicitly, and the top eigenvalue by a verified Cholesky factorisation; the margins are $0.0261$, $0.0077$, $0.0027$ at $L=1.1,1.2,1.3$, all positive, and no step depends on unbounded rounding. % pending(D7)
\end{remark}

\begin{remark}[What the window contains]
$L=1.3$ contains the primes $2$ and $3$, both inert: on both, $\Frob_p=\sigma$ and
$a_p>0$.  The first term with the opposite sign is $n=4$ at $\log4=1.386$, where
$\chiv(4)=+1$ and $a_4=-\log2$.  The next window therefore has to absorb a
contribution that works against positivity, and it is the first place where the
split/inert alternation of the Frobenius, not just its sign at one prime, has to
be used.
\end{remark}

\begin{remark}[The sign of the Frobenius at $2$]
The window $L\le1$ passes the archimedean threshold $L_c=0.925$
(\cref{prop:semilocal}).  In the certificate the prime $2$ enters with
$a_2=-2\log2\cdot\chiv(2)/\sqrt2=+0.980$.  Its sign is $-\chiv(2)$, and $\chiv(2)$
is the eigenvalue of $\Frob_2$ on $\sqrt5$ (\cref{thm:frob-eigen}): the statement is
that $2$ is inert.  Replacing that sign, the same linear programme has optimum
$-0.573$ (\texttt{cert.py}).  The replaced sequence is not the character of any $L$-function, so this % pending(D7)
is a statement about the certificate, not about another $L$-function: it shows
which term of the prime side carries the positivity past $L_c$.  \Cref{thm:second-window} is the
first statement in which the Frobenius of $\Zphi$ does work inside Weil
positivity.  The certificate method reaches $L=1.08$ (margin $0.013$, \texttt{cert.py}); beyond, % pending(D7)
the pointwise criterion is too lossy and the operator form of the next step is
needed.
\end{remark}

\begin{remark}[Scope of \cref{thm:first-window}]
This is the $\chiv$-analogue of the archimedean positivity of Connes and
Consani, extended to include the first prime.  It constrains no individual
zero: a window of length $L$ sees the zeros only through $|\hat f|^2$ with
$\hat f$ of exponential type $L/2$.  Its proof does not use the sign of
$\chiv(2)$; \cref{prop:semilocal} shows that beyond $L\approx0.9$ that sign
becomes necessary, which is where the next theorem has to use the Frobenius.
\end{remark}

\begin{remark}[Relation to the literature on small-support positivity]\label{rmk:yoshida}
For $\zeta$, Weil positivity on test functions of small support is classical:
Yoshida proved it for autocorrelations supported in $[-\log2,\log2]$
\cite{yoshida92}, Bombieri gave the standard analysis of the quadratic functional
\cite{bombieri00}, and Connes--Consani re-proved the archimedean range by trace-formula
methods \cite{connesconsani21}.  In the convention of this paper (support of $f$ of
length $L$, autocorrelation in $[-L,L]$) that range is $L\le\log2$.  Zhu has recently certified positivity for $\zeta$ up to autocorrelation support $1.6$, i.e.\ $L=1.6$ here (his $L$ is the half-width, $0.8$), and measured the window floor, which for $\zeta$ is extremely small ($\sim10^{-6}$ at $L=1$, an upper bound; $\sim10^{-17}$ at $L=1.6$, certified on both sides) \cite{zhu26}; among his open problems he proposes repeating the reduction for Dirichlet $L$-functions with the comb of the corresponding conductor, which is the case treated here.  The theorems of this section are the
$L(s,\chiv)$ case, and two differences with $\zeta$ are visible in the numbers.  The
floor for $\chiv$ is far larger at the same window ($\lambda_{\min}\approx0.097$ at
$L=1$, stable under basis refinement and confirmed from the zero side,
\texttt{weil\_window\_chi5.py}, \texttt{zeroside\_check.py}), because the % pending(D7)
first zero sits at $6.65$ rather than $14.13$ and a band-limited function cannot
avoid it; and $\Lambda(s,\chiv)$ has no pole term.  The first window,
$L\le0.7447$, already passes the classical range, and the certificates reach
$L=1.3$.
\end{remark}

\begin{remark}[The programme this defines]
Positivity on every window is GRH$(\chiv)$.  Positivity on a given window is a
finite statement: for $L<\log2$ it is archimedean (the analogue, for
$\chiv$, of the Connes--Consani archimedean positivity
\cite{connesconsani21}), and each new prime adds one term whose sign is
the Frobenius.  The first windows are within reach of a rigorous proof by
explicit kernel estimates; the super-exponential decay of $\lambda_{\min}$ says
that no fixed-precision argument reaches all of them, and that the full
statement needs an exact positivity (\cref{rmk:saturation}).
\end{remark}

\begin{remark}[The two of $\ker G$]
Since $\Lambda(s,\chiv)$ is entire, its explicit formula has no pole terms.
The two poles of the completed $\zeta_{\mathbb{Q}(\sqrt5)}$ come from $\zeta$ and
are present for every number field, whereas $|\ker G|=8/4=2$ is specific to
conductor~$20$.  The equality of the two numbers is not structural, and the
paper does not use it.
\end{remark}

The windows of this section use test functions of compact support; the Fej\'er family
$g_L(u)=(1-|u|/2L)^+$ gives the same positivity in a form measurable from primes alone, and
\cref{ssec:field-positivity,ssec:weil-closure} show that, followed along the Frobenius orbits, it writes the open statement as one bound per class, the hypothesis of the closing theorem on
Weil's side (\cref{thm:weil-closure}), which is the Weil face of \cref{conj:closure}.

% ---------------------------------------------------------------------
\subsection{The verdict that remains}\label{sec:results-verdict}

\subsubsection{The core, and its equivalent forms}\label{ssec:core}\label{ssec:core-forms}

Every step above is proved or carries its witness mark.  What remains is one step: that each
zero is a member of the structure, and not only centred on it
(\cref{prop:each-zero}).  The corpus states that step in several forms.  Each is proved
equivalent to \Ecore, directly or through another form, so they are one step written
in different languages, not separate open problems.  How the paper proposes to take that step
--- Weil's route, and the closure conjecture it leads to --- is the subject of the last subsections,
ending in \cref{conj:closure,cor:closure}.

\begin{longtable}{p{5.2cm}p{5.6cm}p{2.6cm}l}
\toprule Form of the step & Equivalence & Condition & Tier\\\midrule\endhead
The real part of each zero is a spectral position of $\OmH$ (the last link of the bridge) & \LeanInline{GDS\_iff\_E}, \LeanInline{landing\_iff\_E} & none & \tier{P}\\
The coordinate of each zero is $1$ & \LeanInline{face\_coordinate\_iff\_GRH}, \LeanInline{liMap\_on\_unit\_circle\_iff} & none & \tier{P}\\
No two zeros share an ordinate & \LeanInline{E\_iff\_no\_shared\_ordinate\_chi5}, \LeanInline{E\_iff\_no\_shared\_ordinate\_of\_fe\_rc} & none & \tier{P}\\
The symmetric product $\rho(1-\rho)$ of every zero is real & \LeanInline{no\_shared\_ordinate\_iff\_e2\_real} with the previous row & no zero on the real axis & \tier{P}\\
Premise A: $\Re\rho\le\mu_3$ for all zeros & \LeanInline{premise\_A\_iff\_GRH\_chi5} & none & \tier{P}\\
The coordinate is a power of $\golden$ (membership in the tower) & \LeanInline{the\_zeros\_must\_sit\_at\_level\_zero}, \LeanInline{descending\_to\_F1\_is\_E} & discriminant bound of the zero & \tier{P}\\
$r(\rho)\in\golden^{\Z}$ with the strip bound & \LeanInline{position\_is\_interval\_meets\_discreteness}, \LeanInline{radius\_one\_is\_E} & Prop.~\ref{prop:rbound} & \tier{P}\\
The trace $t(\rho)$ is an integer & \LeanInline{the\_two\_labels\_are\_one\_gap}, \LeanInline{route\_B\_equals\_GRH} & $|\gamma|\ge2$; discriminant bound & \tier{P}\\
The real degree $\log r/\log\golden$ is an integer & \LeanInline{E\_iff\_realDegree\_integral} & height of the first zero & \tier{P}\\
The primes act by powers on the zero coordinates & \LeanInline{B\_iff\_E\_one\_hypothesis} & height of the first zero & \tier{P}\\
\bottomrule\end{longtable}

The conditions are three, and each is settled: the discriminant bound is proved for
every zero with $|\gamma|\ge2$ (\tier[discBound\_discharged\_high]{P}); the
height of the first zero is the classical input \LeanInline{NoLowZeros}; and the
unit structure used by the membership form is proved in its corrected version, with
$r\in\Q(\sqrt5)$ (\LeanInline{F1v2.arrowB\_full}).  The forms are also linked to one
another without passing through \Ecore: an integer degree gives membership, the
spectral form gives membership and integrality, and integrality gives membership given
the unit structure (\tier[hasDegree\_gives\_monomial, the\_membership\_derived, GDS\_gives\_routeB]{P}).  The full
graph of these implications, with the theorem behind each edge, is in the ontology of
the corpus (\cref{app:lean}).

\begin{proposition}[The step as a statement about the coordinate alone]\label{prop:step-coordinate}
\begin{enumerate}[label=\textup{(\roman*)}]
  \item For a zero in the strip, the golden valuation of its coordinate is $0$
    exactly when the zero is on the line.
  \item For $r>0$ with $\mathrm{disc}(r)\le\tfrac1{12}$, $r$ is a power of $\golden$
    exactly when $r=1$.
  \item Membership in the tower is closed under powers and inverses: it is preserved
    by every Frobenius lift and by the Klein involution $r\mapsto1/r$.
  \item For any set of positive coordinates with a discriminant bound below $1$, the
    traces are all integers exactly when every coordinate is $1$.
  \item The spectral form of the step gives Premise A.
  \item The spectral parameters are the only pair: $x^2-2x+\tfrac34=0$ exactly for
    $x=\tfrac12$ and $x=\tfrac32$.
\end{enumerate}
\tier[degenerate\_iff\_valuation\_zero, monomial\_iff\_one\_under\_confinement, step\_six\_membership\_would\_give\_the\_power, isMonomial\_inv, integrality\_iff\_E, trace\_integral\_iff\_E, GDS\_gives\_premiseA, apex\_pair\_is\_unique]{P}%
\end{proposition}

\begin{proposition}[The $\golden$ side of the step: from an integral trace to the line]\label{prop:phi-side}
Let $r>0$ and $t=r+1/r=2\cosh(\log r)$.
\begin{enumerate}[label=\textup{(\roman*)}]
  \item $t$ is an integer exactly when $r$ is an eigenvalue of an integer matrix of
    determinant $1$: if $r^2-mr+1=0$ for such a matrix of trace $m$, then $t=m$, and
    conversely an integral $t$ makes $r$ an eigenvalue of one.
  \item For a zero $\rho\ne0,1$, the trace of its coordinate is the integer $n$ as soon as
    the Apollonius relation $|\rho-1|^2+|\rho|^2=n\,|\rho-1|\,|\rho|$ holds.
  \item In the tower: $\golden(\golden-1)=1$, so $\golden$ is a unit; no $u$ with $1<u<\golden$ has
    trace reaching that of $\golden$, $2<u+1/u<\golden+1/\golden$; and a level $\golden^k$ with
    $|\log\golden^k|<\log\golden$ is $\golden^0$.
  \item If the excesses $t-2$ of finitely many coordinates sum to $0$, every one of those coordinates is $1$.
  \item Every $r>1$ with integral trace $t$ is a unit of a real quadratic order,
    $r^2=tr-1$, as $\golden$ is a unit of $\Zphi$, a root of $x^2-x-1$; with any discriminant bound below $1$, an
    integral trace gives $r=1$, and it gives the integrality of the discriminant.  The
    zeros and the tower lie on the same locus of norm one, $r\cdot r'=1$, and every level
    $\golden^k$ has integral discriminant.  An integral trace is stronger than membership:
    $\golden$ itself has trace $\sqrt5$, so integrality of the trace holds on the even
    levels only --- which contain the one level, $k=0$, that the zeros can occupy.
\end{enumerate}
\tier[trace\_of\_integer\_matrix, integer\_trace\_gives\_SL2\_eigenvalue, the\_quadratic\_form\_condition\_gives\_the\_trace, cosh\_is\_the\_mechanism, phi\_is\_the\_fundamental\_unit, the\_trace\_closes\_the\_quadrants, lattice\_meets\_small\_ball\_only\_at\_zero, sum\_zero\_forces\_each, zeroModulusUnit, phiUnit, trace\_integral\_and\_bound\_give\_E, trace\_integral\_gives\_integrality, both\_sides\_on\_norm\_one\_locus, family\_membership, trace\_not\_integral\_at\_phi]{P}, the lemmas with a parameter $\golden$ instantiated
at the golden ratio through $\golden^2=\golden+1$ (\tier[phi\_sq]{P})
\end{proposition}

With the confinement of \cref{prop:each-zero} ($t\le2.00013<3$), \cref{thm:arrow}(a) closes
the step the moment the trace is an integer; the proposition adds the two equivalent
readings of that integrality and places the unit $\golden$ inside it.  So the $\golden$ side of the corpus turns membership into one arithmetic property of one number per zero: that the trace of its coordinate is an
integer --- equivalently, that the coordinate is an eigenvalue of a unimodular integer
matrix, or that the zero satisfies the Apollonius relation with an integer $n$.  Every
arrow from there to the line is proved.

So the open statement, in every form, is about one real number per zero: that the
coordinate, already confined, is the unit of the tower.  \Cref{prop:step-coordinate}(ii)
and (iii) say this from two sides: the tower is a group on which the Frobenius lifts and
the Klein involution act, and inside the confinement its only element is $1$.  In the
language of $\Fone$ it is the statement that the zero descends: writing $r=\golden^k$, the
bound of \cref{prop:rbound} gives $|\log r|<0.0142<\log\golden$, an integer $k$ with
$|k\log\golden|<\log\golden$ is $0$, and so $r=1$ (\tier[k\_zero\_by\_discreteness, position\_is\_interval\_meets\_discreteness]{P}).  The interval alone does
not place the zero and the tower alone does not see its coordinate; what is missing is
that the coordinate lies in the image of the exponent monoid $X_{\PCF}$ of
\cref{thm:lambda-monoid}.

\begin{theorem}[Equivalence with the power action]\label{thm:reduction}
Assume \LeanInline{NoLowZeros}.  Then \Ecore\ holds if and only if the
zero coordinates carry the power action of the golden monoid --- for every prime $p$ and every zero $\rho$ there is a zero
$\rho'$ with $r(\rho')=r(\rho)^p$, as the lift $\Frob_p$ sends $\golden^n$ to $\golden^{pn}$ on the tower
(\LeanInline{ZeroCoordsCarryPowerAction}); all the analytic
inputs of \cref{sec:analytic-inputs} are discharged and no other hypothesis is
used.
\tier[B\_iff\_E\_one\_hypothesis, route\_C\_needs\_only\_noLowZeros\_and\_the\_action]{P}%
\end{theorem}

\begin{remark}[Reading the equivalence]
\Cref{thm:reduction} states in one line what \cref{thm:topological-proof} develops
in eight stages: the arithmetic of \crefrange{sec:choice}{ssec:base-phi}, with one classical input about the
height of the first zero, carries the argument to the verdict of
\cref{rmk:left-to-close}.  What is proved about the zeros unconditionally along the
way --- the orbit centre, the line as a fixed set, the Klein symmetry induced by
$\chiv$, the
conductor read back from the zeros, and the windows of Weil positivity --- is
recorded in \cref{ssec:stage6,prop:envelope-exact,ssec:weil-route}.
\end{remark}

The forms above say what would close the argument; none says how.  The structure already
built toward them is considerable: the Frobenius lifts act on the index of the operator and
give the corpus action on the coordinates (\cref{prop:primon-space}); the bridge from the
tower to the operator is proved in every link but the last; the trace of every coordinate is
confined below $2.00013$, so that integrality alone would force $t=2$ (\cref{thm:arrow}); the
orbit of every zero is centred at $\tfrac12$ (\cref{prop:each-zero}); and Weil's positivity
is proved in one window and certified in two more (\cref{ssec:weil-route}).  What is missing in each is the same thing,
seen from a different side.

\subsubsection{The two halves of \Ecore}\label{ssec:two-halves}

The forms of the table concern one quantity, the coordinate of a zero, and the statement they make splits into two halves, of which the corpus proves one on every zero.  The coordinate of a zero and that of its companion, $r=r(\rho)$ and $r'=r(1-\rho)$, satisfy two conditions on the line: their product is $1$ and their sum is $2$.  Together they give $r=r'=1$.

\begin{proposition}[The two halves]\label{prop:two-halves}
Let $\rho$ be a zero of $L(s,\chiv)$ in the strip, $r=r(\rho)$ and $r'=r(1-\rho)$.
\begin{enumerate}[label=\textup{(\roman*)}]
  \item \emph{The multiplicative half holds on every zero.}  The coordinate is equivariant,
    $r(1-\rho)=r(\rho)^{-1}$, and the zero set is closed under $\rho\mapsto1-\rho$, so $r\cdot r'=1$;
    if $r=\golden^k$ then $r'=\golden^{-k}$, so the tower passes to the companion.
  \item \emph{From the multiplicative half to the additive.}  Given (i), $r+r'$ is the trace
    $t(\rho)=r+1/r$, and the line at $\rho$, $\Re\rho=\tfrac12$, is the additive half $r+r'=2$: $r\cdot r'=1$ and $r+r'=2$ give
    $r=1$, and $t=2\iff r=1$.
  \item \emph{The additive half, confined.}  On every zero $t(\rho)\ge2$; with the height of the first zero
   , $t(\rho)^2\le4+\tfrac1{12}$, and in fact $t\le2.00013$ (\cref{prop:rbound}); and
    $t=\bigl(|\rho-1|^2+|\rho|^2\bigr)/\bigl(|\rho-1|\,|\rho|\bigr)$.  The next trace values of the tower,
    $3=L_2$ at even index and $\sqrt5$ at odd index, lie outside that window
    (\cref{prop:lucas-spectrum}), so the only value of the tower inside it is $2$.
  \item \emph{What is open is the additive half alone.}  Besides the forms of the table --- the trace an integer, the coordinate a power of $\golden$ (\cref{thm:arrow}) --- it reads $t(\rho)=2$, and the quadratic-form condition $|\rho-1|^2+|\rho|^2=n\,|\rho-1|\,|\rho|$ for an integer $n$; each is proved equivalent to $\Re\rho=\tfrac12$.
\end{enumerate}
(i) \LeanInline{coordinate\_is\_equivariant}, \LeanInline{coords\_closed\_under\_inversion},
\LeanInline{the\_companion\_plays\_the\_conjugate}, \LeanInline{monomial\_transports\_to\_the\_mate};
(ii) \LeanInline{union\_of\_halves\_is\_E}, \LeanInline{two\_halves\_one\_statement};
(iii) \LeanInline{trace\_ge\_two}, \LeanInline{the\_trace\_is\_confined\_to\_a\_tiny\_interval},
\LeanInline{the\_trace\_as\_a\_quotient}, \LeanInline{the\_exclusion\_is\_complete\_over\_both\_parities};
(iv) \tier[what\_remains\_of\_the\_gap, the\_condition\_is\_exactly\_E, the\_ladder\_is\_complete]{P}, with the
height of the first zero\ in (iii) and (iv)
\end{proposition}

So the chain is closed on the multiplicative side, and on the additive side it is closed up to one
value: the trace of every zero lies in $[2,\,2.00013]$, the tower has no value there but $2$, and
what is not proved is that the trace takes that value.  No intermediate link is missing; the open statement of \cref{ssec:core-forms} is exactly the additive half, and the route that follows is the paper's proposal for supplying it.

\subsubsection{The route of closure: Weil's proof, transported}\label{ssec:route-overview}

The paper proposes one route, the one the $\Fone$ programme was made for: Weil's proof for
curves, carried along the chain arithmetic, geometry, analysis, operator.  Weil's proof has five steps: (W1) Frobenius acts on a lattice and the zeta function
is its characteristic polynomial; (W2) the intersection numbers of the graphs of its powers
count points; (W3) the intersection form has the Hodge signature, one ample direction and the
rest definite; (W4) Castelnuovo--Severi bounds the point counts on the Frobenius orbit;
(W5) the bound for every power gives the line.  Transported along the chain:

\begin{center}\small
\begin{tabular}{@{}lp{7.4cm}p{3.6cm}@{}}
\toprule step & here & status\\\midrule
W1 & $\Frob_p$ on $\Zphi$, Euler factor = characteristic polynomial; the eight classes & \tier{P}, \tier{P}\\
W2 & Weil form; translation by $\log p$ is the Frobenius; the orbit matrix is Toeplitz & \tier{P}\\
W3 & one negative direction (the pole, the ample
class); primitive part definite on orbits; on the angles
of $\golden$, roots of unity, a sum of squares; the
modulus normalised by the tower; positivity on the
windows $L\le0.7447$, $1$ and $1.3$
& \tier{N}; angles and normalisation \tier{P}; windows
\tier{P} with declared inputs, computer-assisted\\
W4 & $|t_k|\le t_0$ on the orbit, class by class, with $t_k$ the Weil form between a test function and its translate by $k\log p$, the $k$-th power of the Frobenius at $p$: the Weil face of \cref{conj:closure}; it is \textup{(W3)} on the planes of the orbit (\cref{prop:W4-from-W3}) & on the planes of the windows \tier{P}; certified for $k\le25$ (computer-assisted); open for every $k$\\
W5 & bounded orbit $\Rightarrow$ no zero off the line & \tier{P}\\
\bottomrule\end{tabular}
\end{center}

The subsections that follow take these steps in order.  The fluctuation and the two sides
of the structure give W1 and W2 (\cref{ssec:fluct-geometry,ssec:two-sides}); the Fej\'er form
on the fibres of $G$, the intersection form in finite dimension, the rotations and the golden
field give the evidence for W3 (\cref{ssec:field-positivity,ssec:product-surface,ssec:bridge-boosts,ssec:golden-field});
and \cref{ssec:weil-closure} gives the closure: the positions, the belonging of the primes
and the entry of the zeros, the conjecture, and W4 and W5 as its face on Weil's side.

\subsubsection{W1--W2. The fluctuation as geometry}\label{ssec:fluct-geometry}

On the line $\Re s=\sigma$ the Euler factor of $L(s,\chiv)$ at $p$ is
$(1-\chiv(p)p^{-\sigma}e^{-i\theta_p})^{-1}$ with $\theta_p=T\log p$.  It has three
ingredients: the sign $\chiv(p)$, which is the Frobenius and is fixed by the class of $p$
modulo $20$; the modulus $p^{-\sigma}$, which on the critical line is $p^{-1/2}$, the $\tfrac12$
at every prime; and the angle $\theta_p$, a point of a circle, one circle per prime.  The
fluctuation is therefore one function on the torus of all these circles, read along the
line $T\mapsto(T\log p)_p$.

\begin{proposition}[The torus and the line]\label{prop:torus-line}
\begin{enumerate}[label=\textup{(\roman*)}]
  \item On the line, the factor of the torus is the local factor of the corpus:
    at $\theta_p=T\log p$ it equals $f_L(p,\sigma+iT)$.
  \item The modulus of the term of each prime does not depend on the angle:
    $|\chiv(p)p^{-\sigma}e^{-i\theta}|=p^{-\sigma}$.
  \item If $a\log p=b\log q$ with $p\ne q$ primes and $a,b\in\N$, then $a=b=0$ (unique
    factorisation): the line is equidistributed in every torus of finitely many primes.
  \item The smoothed fluctuation is a function of the torus, and on the line it is exactly
    the prime side of the explicit formula.
  \item If $L(s,\chiv)\ne0$ for $\tfrac12<\Re s<1$ (the open statement written as a zero-free region), every zero of the strip lies on the line;
    the left half is given by the functional equation.
\end{enumerate}
{\raggedright \LeanInline{F1PaperTorus.torusFactor\_on\_line}, \LeanInline{F1PaperTorus.term\_modulus\_is\_height\_free}, \LeanInline{F1PaperTorus.log\_primes\_independent}, \LeanInline{F1PaperTorus.flucTorus\_on\_line}, \LeanInline{F1PaperTorus.typical\_line\_gives\_the\_line}, and Kronecker's density for two primes, \LeanInline{F1PaperTorus.kronecker\_two\_primes} (from the irrationality of $\log p/\log q$), \tier{P}; Weyl, Bohr--Jessen and Bagchi as named inputs\par}
\end{proposition}

The law of the fluctuation is geometric.  Along the true line, at random points of the
torus, along a line with surrogate primes of the same classes, and on the zero side (the
$397$ zeros below height $500$), the smoothed fluctuation has the same distribution:
standard deviation $\approx0.62$, skewness $\approx-0.33$, kurtosis $\approx-0.55$, with
Kolmogorov--Smirnov $p$-values between $0.14$ and $0.69$ (\texttt{testGeom.py} \tier{N}).  Its
individual values depend on the point of the line, and to reproduce them at height $T$
requires $\log p$ to precision $\approx1/T$ (\texttt{testB1b.py} \tier{N}); its law is fixed by the
Frobenius signs, the modulus $\tfrac12$ and the independence of the $\log p$.  This does not mean the operator's error grows with height: by
\cref{rmk:two-errors} its relative error stays bounded, and the real part
is read with no loss of resolution in the Fej\'er coordinate (\cref{prop:envelope-edge}).

\subsubsection{W1--W2. The two sides of the structure}\label{ssec:two-sides}

\begin{proposition}[The additive side controls every integer at once]\label{prop:additive-side}
The partial sums of $\chiv$ depend only on $N\bmod5$ (they are $0,0,1,0,-1$) and are bounded
by $1$.  Hence, by Abel summation, $|\sum_{n<N}f(n)\chiv(n)|\le f(0)$ for every decreasing
$f\ge0$ and every $N$, and the Dirichlet series of $L(s,\chiv)$ converges in the whole region
$\Re s>0$.
{\raggedright \tier[F1PaperTorus.partialChi\_eq, F1PaperTorus.partialChi\_bounded, F1PaperTorus.abel\_bound, F1PaperTorus.integer\_block\_bounded]{P}; the convergence is Dirichlet's test\par}
\end{proposition}

\begin{proposition}[The multiplicative side, and the inverse of the structure]\label{prop:multiplicative-side}
\begin{enumerate}[label=\textup{(\roman*)}]
  \item $\chiv$ is completely multiplicative; hence the coefficients of $1/L(s,\chiv)$
    are $\mu(n)\chiv(n)$, that is,
    \[\textstyle\sum_{d\mid n}\mu(d)\chiv(d)\chiv(n/d)=[n=1].\]
  \item On the circle of each prime, $3+4c\cos x+\cos2x=2(1+c\cos x)^2\ge0$ for $c=\chiv(p)^k=\pm1$:
    the Frobenius sign is a half-turn of the circle, and the positivity survives it.
  \item If $L(s,\chiv)$ has an inverse at every point of $\tfrac12<\Re s<1$ --- that is, does not vanish
    there, the open statement itself in another form, as in \cref{prop:torus-line}(v) --- every zero of the
    strip lies on the line.
\end{enumerate}
{\raggedright \tier[F1PaperTorus.chi5\_mul, F1PaperTorus.inverse\_of\_the\_structure, F1PaperTorus.three\_four\_one\_signed, F1PaperTorus.circle\_positivity, F1PaperTorus.inverse\_gives\_the\_line]{P}\par}
\end{proposition}

Neither side is enough alone.  The additive side is not: the Davenport--Heilbronn
sequence, periodic modulo $5$ with values $1,\kappa_{\mathrm{DH}},-\kappa_{\mathrm{DH}},-1,0$
($\kappa_{\mathrm{DH}}=(\sqrt{10-2\sqrt5}-2)/(\sqrt5-1)$), has bounded partial sums and the functional
equation of conductor $5$, and its series has zeros off the line; the partial sums of
its inverse grow like $y^{0.788}$, those of $\mu\chiv$ like $y^{0.443}$ (\texttt{dh.py} \tier{N}).
The multiplicative side is not: summing the positivity of all circles requires absolute
convergence --- for $\sigma>1$ it gives $\|L(\chi^0,\sigma)^3L(\chiv,\sigma+it)^4L(\chiv^2,\sigma+2it)\|\ge1$
(\tier[F1PaperTorus.mertens\_positivity\_mathlib]{P}, through Mathlib), and hence no zero on
$\Re s\ge1$ (\tier[F1PaperTorus.nonvanishing\_from\_one]{P}) --- and the product stays $\ge1$ only down to
$\sigma\approx0.9$ (\texttt{mertens341.py} \tier{N}).  What would join them is the bound
$\sum_{n\le y}\mu(n)\chiv(n)=O(y^{1/2+\varepsilon})$ on the inverse of the structure
(\LeanInline{F1PaperTorus.InverseStructureBound}, open; measured $\le1.23\sqrt y$ up to $3\cdot10^7$,
\texttt{mobius.py} \tier{N}), while the direct side is bounded by a constant.

The same sequence separates the equivalent forms of the open statement from the routes to it.  It has
the functional equation, real coefficients and bounded partial sums, so its zero set is closed under the
Klein group, and \cref{prop:orbits}, \cref{thm:shared-ordinate} and \cref{prop:additive-side} hold for
it word for word; its zeros off the line, at $0.8085+85.699i$ and $0.6508+114.163i$, are free orbits
(\texttt{retro99.py}).  What it lacks is on the prime side: the coefficients of $-L'/L$ of an Euler % pending(D7)
product are supported on prime powers, and for this series $130$ of the coefficients with $n\le400$ fall
outside them, the first at $n=6$ (\texttt{dh\_logderiv.py}).  An argument that reaches the line must % pending(D7)
therefore use the support on prime powers, as the windows of \cref{ssec:weil-route} and the orbit of
\cref{thm:weil-closure} do; a statement that holds equally for this sequence is a form of the open
statement, not a step towards it.  The exponent $0.788$ measured for its inverse is close to the real part
$0.8085$ of its first zero off the line, as the classical relation between the partial sums of $1/L$ and
the supremum of the real parts of the zeros gives; this calibrates the value $y^{0.443}$ measured
for $\mu\chiv$.

\subsubsection{W3. Positivity on the fibres of $G$: the Fej\'er form of $\Q(i,\sqrt5)$}\label{ssec:field-positivity}

With the Fej\'er window $g_L(u)=(1-|u|/2L)^+$, the prime side splits into the block with
character, $k_1=-2\sum_p\chiv(p)\frac{\log p}{\sqrt p}g_L(\log p)$, the block of squares, which
carries no character because $\chiv(p^2)=1$, and the higher powers.  The corpus shows that the open statement follows from the convergence of $k_1-L$ as $L\to\infty$ (\LeanInline{CharacterBlockConverges}); its limit
is forced, $-\kappa-c_2-c_3=-0.0802$, from the conductor in the archimedean kernel
($-\kappa=\gamma_E+\log8+\pi/2+\log\pi-\log5$, with $\kappa$ the constant of \cref{thm:first-window}), the
contribution $c_2$ of the block of squares, which carries Mertens' constant, and the contribution $c_3$ of the higher powers
(\LeanInline{F1PaperTorus.forcedLimit}, \LeanInline{F1PaperTorus.character\_block\_gives\_the\_line} \tier{P}; \texttt{blocks.py} \tier{N}).  Equivalently, the Fej\'er sum over the
zeros $Z(L)=\sum_\rho h_L(\gamma_\rho)$, with $h_L(r)=2L(\sin Lr/Lr)^2\ge0$, computed from primes
alone as prime side plus archimedean term, must be $\ge0$ and tend to $0$
(\tier[F1PaperTorus.fejerKernel\_nonneg, F1PaperTorus.fejer\_sum\_nonneg]{P}; checked against the $397$ zeros to
$3\cdot10^{-4}$, \texttt{fejer\_check.py} \tier{N}).

\begin{proposition}[The Klein orbits and the four fibres of $G$]\label{prop:four-fibres}
\begin{enumerate}[label=\textup{(\roman*)}]
  \item Let $\chi_{20}=\chi_4\chiv$.  For $n$ odd and prime to $5$, each combination $1\pm\chi_4\pm\chiv\pm\chi_{20}$ with the
    signs of a character of $\Z_2\times\Z_2$ equals
    $4$ on one fibre of $G$ and $0$ on the others: $(0,0)=\{1,9\}$, $(0,1)=\{13,17\}$,
    $(1,0)=\{11,19\}$, $(1,1)=\{3,7\}$ modulo $20$.
  \item Every square falls in the kernel fibre: $n^2\equiv1,9\pmod{20}$ ($G$ has image of
    exponent $2$).
  \item If each of the four Fej\'er sums of $\Q(i,\sqrt5)$ is $\ge0$ and their sum tends to
    $0$, each tends to $0$.
\end{enumerate}
{\raggedright \tier[F1PaperTorus.four\_fibre\_indicators, F1PaperTorus.squares\_in\_kernel\_fibre, F1PaperTorus.each\_decays\_of\_field\_decays]{P}\par}
\end{proposition}

The positivity is a balance between the two kinds of Klein orbit.  At $2L=18$ (primes below
$e^{18}$) the split primes, the degenerate orbits, contribute $-1787.59$, the inert primes, the
free orbit, $+1796.72$, the archimedean term $-3.29$ and the higher powers $-5.83$, leaving
$Z=0.0113$: the two orbit types cancel to $0.5\%$, and what remains of their race is the
conductor plus the squares (\texttt{balance.py} \tier{N}).  For the four fibres, the deviation
of each fibre's Fej\'er sum from a quarter of the total is predicted by the three conductors,
the squares and the primes $2$ and $5$ alone, up to exactly $-\frac18\sum_\psi\psi(g)Z_\psi$:
measured $-2.7954,0.6963,0.7488,1.3504$, predicted $-2.7815,0.6871,0.7376,1.3568$, remainders
equal to $-\frac18\sum_\psi\psi(g)Z_\psi$ with $Z_4=0.019$, $Z_5=0.0113$, $Z_{20}=0.0814$
(\texttt{fibres.py} \tier{N}).  The four characters are those of $\Q(i,\sqrt5)$, the field of
conductor $20$ with Galois group $\Z_2\times\Z_2$, whose Artin map is $G$.  Its Fej\'er
sum, the sum of the four, falls from $0.398$ ($2L=6$) to $0.118$ ($2L=18$), with every part
positive (\texttt{dedekind.py} \tier{N}).

Positivity alone suffices.  A zero off the line, $\rho=\tfrac12+\delta+i\gamma$, enters $Z(L)$ with
a term of size $\approx e^{2L\delta}/(2L\gamma^2)$ whose sign oscillates, and it makes $Z$
negative once $L$ is of order $\log\gamma/\delta$: adding such a quadruple at $\gamma=100$ to the
$397$ zeros makes $Z<0$ at $2L\approx63$ for $\delta=0.1$ and at $2L\approx127$ for $\delta=0.05$, while
without it the minimum over $2L\in[2,240]$ is $+6.3\cdot10^{-4}$ (\texttt{offline.py} \tier{N}).
Read as positivity, the open statement is one inequality, and three objects carry it, in
increasing strength.  For the line of $\chiv$ alone it is enough that the Fej\'er sum of
$\chiv$ be $\ge0$ for every $L$ --- one of the four.  The golden field $\Q(\sqrt5)$, whose zeta is
$\zeta\cdot L(s,\chiv)$, adds the sum of $\zeta$: it brings the pole, the ample class of the Hodge
picture (\cref{ssec:golden-field}), and it also asserts the Riemann hypothesis for $\zeta$.  The
field $\Q(i,\sqrt5)$ of the four fibres of $G$ adds $\chi_4$ and $\chi_{20}$: \emph{the four Fej\'er sums of
$\Q(i,\sqrt5)$ are $\ge0$ for every $L$} (\LeanInline{F1PaperTorus.FieldFejerPositivity}, open), which gives the
line of $\chiv$ and more.  It
is the analogue, for $\mathrm{Spec}\,\Z$ and this field, of the positivity that Weil obtains on
$C\times C$ from the Hodge index; the Castelnuovo--Severi bound itself is (W4) of \cref{ssec:weil-closure}.  Here, as there, the positivity is not local --- the split primes
subtract, the inert primes add --- and the object that supplies it in Weil's proof, the product
surface, is the object the $\Fone$ programme seeks.

\subsubsection{W3. The intersection form in finite dimension}\label{ssec:product-surface}

This subsection and \cref{ssec:golden-field} study one object, the Weil form as an intersection
form: here in general and for $\zeta$ and $\chiv$ separately, there for the golden field, where the
pole of $\zeta$ becomes the ample class.

In Weil's proof for curves the positivity comes from the intersection form of $C\times C$: it
has exactly one positive direction, the ample class, and is negative on its orthogonal
complement (the Hodge index), and the Castelnuovo--Severi inequality
$|N_m-q^m-1|\le2g\,q^{m/2}$ follows.  Both have finite-dimensional analogues here.

The analogue of the point counts is $\psi(x,\chi)=\sum_{n\le x}\Lambda(n)\chi(n)$.  Up to $e^{18}$,
$|\psi(x,\chi)|/\sqrt x$ stays below $1.27$, $1.28$ and $2.22$ for $x>10^6$ ($\chi_4$, $\chiv$,
$\chi_{20}$), with no growth --- its maximum over $[3\cdot10^7,e^{18}]$ is $1.13$, $1.07$ and $1.76$ ---,
and each fibre of $G$ deviates from a quarter of the total by at most $0.79\sqrt x$
(\texttt{castelnuovo.py} \tier[\testref{test\_castelnuovo\_golden\_values, test\_castelnuovo\_matches\_manuscript\_claim, test\_castelnuovo\_window\_maxima}]{N}).

The analogue of the intersection form is the Weil quadratic form on a space of test functions $g_1,\dots,g_n$: the matrix $\bigl(\text{archimedean term}-\text{prime side}\bigr)$ of the products $g_j*g_k$.  Finite compressions of this form have recently carried an unconditional result: with a rank--trace inequality and Sylvester's law of inertia for the off-line pairs, more than two thirds of the zeros of $\zeta$ are simple and on the line, the result extends to primitive Dirichlet $L$-functions, and it is formally verified in Lean~4 \cite{alpogefurman26}.

\begin{proposition}[The intersection form]\label{prop:intersection-form}
\begin{enumerate}[label=\textup{(\roman*)}]
  \item If the heights of the zeros are real, the Weil form is the Gram matrix of the vectors
    $(\hat g_j(\gamma_\rho))_\rho$, hence nonnegative:
    $\sum_{j,k}x_jM_{jk}x_k=\sum_\rho\bigl(\sum_jx_j\hat g_j(\gamma_\rho)\bigr)^2\ge0$.
  \item (Hodge index, primitive direction.)  If the total form $M=vv^{\mathsf T}+M_{\mathrm{loc}}$ --- the
    pole, the ample direction, plus the local part --- is nonnegative, then $M_{\mathrm{loc}}$ is nonnegative
    on the orthogonal complement of $v$.
\end{enumerate}
{\raggedright\tier[F1PaperTorus.gram\_form\_nonneg, F1PaperTorus.primitive\_positivity]{P}\par}
\end{proposition}

With eight even test functions (triangles centred at $0,\pm1,\dots,\pm7$, support up to $16$)
and primes below $e^{16.2}$: for $\chiv$ the form computed from the primes equals the Gram
matrix of the $397$ zeros (entries agree to $3.5\cdot10^{-6}$) and its eight eigenvalues are
positive, from $3\cdot10^{-5}$ to $0.048$; for $\zeta$ the local part has exactly one negative
eigenvalue, $-3646$, in the direction of the pole, and seven positive ones, and the total form
is nonnegative (\texttt{gram.py} \tier{N}).  That is the signature of the Hodge index, in
dimension eight.  What remains, in the language of the surface, is that this holds in every
finite dimension for each character of $\Q(i,\sqrt5)$
(open; recorded in Lean only schematically, as \LeanInline{F1PaperTorus.ArithmeticHodgeIndex}; the statement Weil's route needs at this point is (W4) of \cref{thm:weil-closure}, the Weil face of
\cref{conj:closure}): Weil obtains it because the
intersection form of a projective surface has that signature, and the object whose geometry
would impose it here is the one the $\Fone$ programme seeks.  The closing theorem needs less:
the signature only on the Frobenius orbits, class by
class (\cref{ssec:weil-closure}).

\subsubsection{W3. Rotations on the line, boosts off it}\label{ssec:bridge-boosts}

The dimensional chain of \cref{thm:one-two-three} already carries forms with a single exceptional
direction: the golden norm form $x^2+xy-y^2=(x+\golden y)(x+\bar\golden y)$, $\bar\golden=-1/\golden$, has signature $(1,1)$, the
discriminant of binary forms has signature $(2,1)$, and the Fibonacci matrix acts on the
latter as an isometry whose null directions are its eigenvectors (\cref{prop:binary-forms}).
The bridge to the Weil form goes through the Frobenius.  In the Weil form the prime $p$ enters
as the translation by $\log p$, which on the zero side multiplies by $p^{i\gamma}=e^{i\gamma\log p}$.

\begin{proposition}[Rotations on the line, boosts off it]\label{prop:bridge-boosts}
\begin{enumerate}[label=\textup{(\roman*)}]
  \item $|e^{ia\gamma}|=e^{-a\,\Im\gamma}$.  On the line ($\gamma$ real) the translation is a rotation,
    of modulus $1$; off the line ($\gamma=t+is$) the factors on $\gamma$ and on $\bar\gamma$ are
    $p^{-s}$ and $p^{s}$, with product $1$: a boost of rapidity $s\log p$.
  \item The contribution of a quadruple of zeros off the line to the Weil form is
    $\Re(w^2)=(\Re w)^2-(\Im w)^2$: a hyperbolic plane, of signature $(1,1)$.
  \item In the chain, $\golden^2$ scales the null coordinates $x+\golden y$ and $x+\bar\golden y$ of the golden
    form by $\golden^2$ and $\golden^{-2}$ and preserves their product: a boost of the same kind.
\end{enumerate}
{\raggedright\tier[F1PaperTorus.translation\_factor, F1PaperTorus.translation\_is\_rotation\_on\_the\_line, F1PaperTorus.translation\_is\_boost\_off\_the\_line, F1PaperTorus.offline\_contribution\_is\_hyperbolic, F1PaperTorus.hyperbolic\_plane\_is\_indefinite, F1PaperTorus.golden\_boost]{P}\par}
\end{proposition}

A quadruple placed at $\gamma=20$ with $\delta=0.1$ contributes the eigenvalues
$-1.4\cdot10^{-4}$ and $+5.3\cdot10^{-4}$, a hyperbolic plane; the translation by $\log7$ scales
its two members by $1.2148$ and $0.8232$, with product $1$ (\texttt{bridge.py} \tier[\testref{test\_mersenne\_bridge\_phi\_lambda\_log\_two, test\_mersenne\_bridge\_s\_bh\_is\_log\_two, test\_regulator\_form\_of\_bridge\_equals\_l\_one\_form}]{N}).  With
eight test functions the total form stays positive: the negative direction is covered by
the positive contribution of the zeros on the line, and resolving it needs a window of
length of order $1/\delta$, the scale already met in \cref{ssec:field-positivity}.

\begin{proposition}[Free Klein orbits and the inertia index]\label{prop:free-orbit-index}
For real test functions $f$ the zero side of the Weil form (\cref{thm:core-positivity}) is
$Q(f)=\sum_\rho\hat f(\gamma_\rho)\hat f(-\gamma_\rho)$ with
$\gamma_\rho=(\rho-\tfrac12)/i$, since $h(\gamma)=\hat f(\gamma)\hat f(-\gamma)$ is the entire extension of $|\hat f|^2$; let $f^\vee(x)=f(-x)$.
\begin{enumerate}[label=\textup{(\roman*)}]
  \item A zero on the line contributes $2|\hat f(\gamma)|^2\ge0$.  A free orbit
    $\{\rho,\bar\rho,1-\rho,1-\bar\rho\}$ contributes $4\Re\bigl(\hat f(\gamma)\hat f(-\gamma)\bigr)$, of signature
    at most $(2,2)$; on even functions it is $4\Re(\hat f(\gamma)^2)$ and on odd functions
    $-4\Re(\hat f(\gamma)^2)$, each of signature at most $(1,1)$.  The hyperbolic plane of
    \cref{prop:bridge-boosts}\textup{(ii)} is the even one.
  \item $Q$ commutes with $f\mapsto f^\vee$, so its index of inertia is the sum of the indices on even and
    on odd functions, and each is at most the number of free orbits.  In particular \Ecore\ implies
    $Q\ge0$ on every window, and the dichotomy $2/4$ of \cref{prop:orbits} is the dichotomy $0/2$ of the index.
\end{enumerate}
{\raggedright\tier[OrbitaLibre.quadruple\_form, OrbitaLibre.even\_part, OrbitaLibre.odd\_part, OrbitaLibre.on\_line\_nonneg]{P}\par}
\end{proposition}
\begin{proof}
The algebra of $4\Re(AB)$ with $A=\hat f(\gamma)$, $B=\hat f(-\gamma)$: $B=A$ on even functions and $B=-A$ on
odd ones, and $4\Re(A^2)=4(\Re A)^2-4(\Im A)^2$.
\end{proof}
On the zeros of the Davenport--Heilbronn sequence of \cref{ssec:two-sides} below height $100$, which carry one
free orbit at $\gamma=85.70$, $\delta=0.3085$, the truncated zero side has index $0$ up to window $2$ and one even and
one odd negative direction from window $3$, against $1/\delta=3.24$; with the second orbit, at height $114.16$,
two of each (\texttt{retro101.py}).  In the form of $L(s,\chiv)$ with two zeros replaced by a quadruple, the index is $0$ or $2$ on every window tried, never $1$ (\texttt{retro93.py}). % pending(D7)

So the dictionary is exact.  A zero off the line would turn every Frobenius translation into
a boost on a hyperbolic plane of the Weil form, of the same kind as the golden boost of the
chain; on the line every translation is a rotation and the form has no hyperbolic plane.  The
chain has the boost --- the Frobenius of the tower,
$\golden\mapsto\golden^2$, on the golden form --- and
at each prime it also has the rotation.  In the polar
form $\golden=[\golden_{\mathfrak p}]\cdot\langle\golden\rangle$
(\cref{ssec:delta-polar}) the angular part is a root of
unity whose order is the clock of the prime
(\cref{prop:clock-cyclotomic}), and on directions that
are roots of unity the Toeplitz form of an orbit is a
sum of squares and the orbit is bounded by $t_0$, with
no zeros used
(\tier[F1PaperClosureClock.clock\_angles\_W3\_W4]{P}).  The radial part is not left free either.  On the
arithmetic side the tower normalises it exactly, as
$q^{-1/2}$ does in Weil's proof: the Fibonacci matrix
$M$ has eigenvalues $\golden$ and $-1/\golden$, its
traces $L_k$ fail the bound ($L_2=3>L_0=2$), and with the
modulus removed the eigenvalues are $\pm1$ and the bound
holds
(\tier[W4Isometria.golden\_torus\_fails\_W4, W4Isometria.normalized\_M\_satisfies\_W4]{P}); at $2$ the modulus is exactly
$|2|_2=\tfrac12=\golden^{-\lng}$
(\cref{prop:omega-two-adic}); at $5$ the principal unit
has exactly the half valuation
(\cref{rmk:angles-cyclotomic}); and on each level of the
clock the sizes are exact in sum
(\cref{prop:clock-cyclotomic}).  On the zero side the
radius of every zero is confined, $t(\rho)\in[2,2.00013]$
(\cref{prop:rbound}); the multiplicative half
$r(\rho)\,r(1-\rho)=1$ is proved
(\cref{prop:two-halves}), so that a bound on one side
suffices, since a pair in which neither member expands
lies on the line
(\tier[W4Isometria.line\_from\_bounded\_pair]{P}).  What stays open in the radial part is the
last margin: the size of each prime one by one, and on
the zero side $t=2$ exactly.
What the line requires, in this language, is that the boost does not reach the zero side:
that every Frobenius translation acts there by rotations
(\LeanInline{F1PaperTorus.FrobeniusActsByRotations}, open).  On one Frobenius orbit this is exactly
the bound $|t_k|\le t_0$ of \cref{ssec:weil-closure}: rotations keep the orbit bounded, a boost
makes it grow.

\subsubsection{W3. The golden field: Hodge index and the first windows}\label{ssec:golden-field}

The golden field $\Q(\sqrt5)$, whose zeta function is $\zeta_K=\zeta\cdot L(s,\chiv)$, is the middle
object of \cref{ssec:field-positivity}: its positivity asserts more than the line of $\chiv$ (it
includes the Riemann hypothesis for $\zeta$), and in exchange it carries the pole of $\zeta$, the ample
class, and with it the Hodge picture of \cref{ssec:product-surface}.  Its places are read off the weights $1+\chiv(p^k)$: a split prime
weighs $2$ at every power (two places of norm $p$); an inert prime weighs $0$ at odd powers and $2$
at even ones (one place of norm $p^2$).  Since $2$ and $3$ are inert, $\Q(\sqrt5)$ has no place of norm
$2$ or $3$; its first places have norms $4,5,9,11,19,29,\dots$
(\tier[F1PaperTorus.place\_weights, F1PaperTorus.no\_places\_of\_norm\_two\_or\_three]{P}).

\paragraph{The pole is the ample direction.}  Split into the two kinds of Klein orbit, the
Weil form of $\chiv$ on a space of test functions is (archimedean $-$ split primes) $+$
($-$ inert primes).  The first part has exactly one negative direction and is positive on its
complement; the second has exactly one positive direction, the same one, and is negative on its
complement; the Weil form is their sum.  That direction is the pole of $\zeta$: the prime
number theorem turns the main term of the prime sum into the pole functional, which is of
rank one, $2vv^{\mathsf T}$ with $v_j=\hat g_j(i/2)$, and each orbit type carries half of it
(ratio $0.4994$).  Pole and main term cancel exactly, and the total form is
$\mathsf K-\mathsf R$, with $\mathsf K$ the archimedean term plus a bounded explicit part (no primes) and $\mathsf R$ the
remainder $\psi_K(x)-x$ of the primes
(\LeanInline{F1PaperTorus.pole\_form\_rank\_one}, \LeanInline{F1PaperTorus.total\_is\_structural\_minus\_remainder} \tier{P}; the measurements \tier{N}).  For
$\zeta_K$ the local part (archimedean minus places) has exactly one negative eigenvalue, in the
direction of the pole, and the total is nonnegative, with $8$, $12$, $16$, $20$ and $24$ test
functions: the signature of the Hodge index, with the pole as the ample class
(\texttt{field\_signature.py} \tier{N}).  The open statement, in this language, is that the
signature persists in every finite dimension, equivalently $\mathsf R\le\mathsf K$ as quadratic forms
(open; recorded in Lean only schematically, as \LeanInline{F1PaperTorus.GoldenFieldHodgeIndex} and
\LeanInline{F1PaperTorus.RemainderDominated}; the statement Weil's route needs at this point is (W4) of \cref{thm:weil-closure}, the Weil face of
\cref{conj:closure}).

\paragraph{The archimedean form, explicitly.}  For an even test function $g$ with $h=g*g$, Gauss's
integral for $\psi$ gives the prime-free part $Q_\infty$ of the Weil form of $\zeta_K$ as
\[
  Q_\infty(g)=2\,\hat g(i/2)^2+C_0\|g\|^2+D(g),\qquad D(g)=4\int_0^\infty w(u)\bigl(h(0)-h(u)\bigr)\,du,
\]
with $w(u)=e^{-u/2}/(1-e^{-2u})$ and $C_0=\log5-2\log\pi-2\gamma_E+4\int_0^\infty\frac{e^{-2u}-e^{-u/2}}{1-e^{-2u}}\,du=-9.13493$,
where $D\ge0$ because $h(0)-h(u)=\frac12\int(g(x)-g(x+u))^2\,dx$ (\tier[F1PaperTorus.autocorrelation\_le\_energy]{P}).  The
conductor enters as $\log5$ in $C_0$; the pole is the ample term.  Without the pole the form is
negative from support length $\approx0.45$; with it, the pole removes the ground mode of $D$, and
positivity asks that the second eigenvalue of $D$ exceed $|C_0|$, which holds up to support length
$\approx1.65$, beyond the prime-free range $[0,\log4]$ (\texttt{arch\_window.py}).

\paragraph{The first windows, over the whole function space.}  For $g$ supported in an interval
of length $S$ (the convention of \cref{thm:first-window}), the Weil form is
$(1/2\pi)\int|\hat g(r)|^2m_K(r)\,dr+2\hat g(i/2)^2$ with symbol
$m_K(r)=W_K(r)-2\sum_n c_n\cos(r\log n)$, where $W_K(r)=\log(5/\pi^2)+2\Re\psi(\frac14+\frac{ir}2)$ is
increasing, $c_n=\Lambda(n)\bigl(1+\chiv(n)\bigr)n^{-1/2}$ are the weights of the places, and the prime part is bounded by $2\sum_nc_n$.  A function orthogonal to the
polynomials of degree $<N$ has at most the fraction
$(2/\pi)(RL)^{2N+1}/((2N+1)^2(N!)^2)$ of its energy below frequency $R$ ($L=S/2$); by min--max this
gives a proved lower bound $\varrho\le\lambda_{N+1}$ (\tier[F1PaperTorus.lowFrequencyFraction\_nonneg]{P}; the analytic
estimate; $\varrho$ enclosed in interval arithmetic).  Lehmann's method on the full space of
polynomials of degree $<N$, with that $\varrho$ (\LeanInline{F1PaperTorus.LehmannBounds}), then gives
lower bounds for the smallest eigenvalue with no assumption on any eigenvalue:

\begin{center}\small
\begin{tabular}{llrrrr}
\toprule support $S$ & places & $N$ & $\varrho\le\lambda_{N+1}$ & Ritz $\lambda_1$ & lower bound of $\lambda_1$\\\midrule
$\log4$ & --- & $28$ & $3.4808$ & $6.8081\cdot10^{-3}$ & $+6.8052\cdot10^{-3}$\\
$1.5$ & $4$ & $50$ & $3.0425$ & $2.4904\cdot10^{-3}$ & $+2.4876\cdot10^{-3}$\\
$1.7$ & $4,5$ & $60$ & $1.7039$ & $3.7152\cdot10^{-4}$ & $+3.7002\cdot10^{-4}$\\
$2.0$ & $4,5$ & $80$ & $1.9369$ & $1.2171\cdot10^{-5}$ & $+1.2108\cdot10^{-5}$\\
$2.3$ & $4,5,9$ & $150$ & $1.4234$ & $1.8873\cdot10^{-7}$ & $+1.8771\cdot10^{-7}$\\
\bottomrule\end{tabular}
\end{center}

\noindent (\texttt{rigorous\_chain.py}, \texttt{rho\_interval.py}).  So the Weil form of $\Q(\sqrt5)$ is positive on
every test function supported in an interval of length $\le2.3$, with the places of norm $4$
(the inert Frobenius at $2$), $5$ (ramified) and $9$ (the inert prime $3$) inside.  The status is
that of the windows of \cref{thm:second-window,thm:third-window}: the analytic steps and $\varrho$ are
rigorous, while the Lehmann matrices are computed in double precision with the operator applied
exactly (the Rayleigh quotient agrees with the Galerkin value to seven digits); a validated
quadrature of those matrices would make the windows theorems.

\paragraph{The ceiling of windows.}  The margin falls super-exponentially: the lowest eigenvalue
is $6\cdot10^{-9}$ at $S=2.5$ and $3\cdot10^{-13}$ at $S=3.0$, and the residual Lehmann needs falls with
it, while the edge behaviour of the minimiser, of type $(\log(1/d))^{-1/2}$, limits how fast any
basis reaches it (at $S=3.0$, with five trial functions in fifty digits, the bound for $\lambda_1$ is
not yet positive).  Each new window costs exponentially more than the last, as with the minimal
eigenvalue of \cref{prop:semilocal}.  Positivity for every support must come from a structural
argument --- the Hodge index of the golden field in every dimension --- not from windows; along
Weil's route it is needed only on the Frobenius orbits (\cref{ssec:weil-closure}).

\subsubsection{W4--W5. The closure}\label{ssec:weil-closure}

The route ends where the thread of \cref{sec:topological-proof} ends: at the operator of \cite{gonzalez25hp}
and at the zeros it describes.  What the structure already gives of the zeros, what it approximates, and what is left,
can be stated exactly, and every piece of it is already in the corpus.

\begin{proposition}[The positions, from the structure]\label{prop:positions}
\begin{enumerate}[label=\textup{(\roman*)}]
  \item With the counting law $N=\mathrm{env}+S$ and the exact inverse $E^{-1}$ of the envelope,
    the height of the $n$-th zero is $t_n=E^{-1}(n-S(t_n))$, and its distance to the structural
    position $E^{-1}(n)$ is the correction alone; with no correction the height is the
    structural position.
  \item The envelope is written with the $\Gamma$-factor, $\pi$ and the conductor, with no prime
    and no zero, and splits into the part of the conductor and the universal part.
  \item The parameters of the operator $T^*$ are the cocone (\cref{prop:pair-from-G}), and its error is
    the correction plus the parametrisation of the envelope (\cref{prop:error-partition}).
  \item The real part: the centre of the orbit of every zero is $\tfrac12$, which is also
    $\Gamma(3/2)/\sqrt\pi$ (the face of $\pi$), $\golden^{-\lambda_{\log}}$ (the face of $\golden$), the
    fixed set of $s\mapsto1-\bar s$ and the positive eigenvalue of $\hat\Omega$.
\end{enumerate}
{\raggedright \tier[F1PaperClosure.heights\_from\_structure, F1PaperClosure.the\_height\_is\_the\_structure\_up\_to\_the\_correction, envelope\_splits, hpps\_parameters\_are\_the\_cocone, F1PaperClosure.operator\_error\_is\_correction\_plus\_parametrisation, F1PaperClosure.the\_half\_is\_given]{P}, on declarations of the corpus \tier{P}\par}
\end{proposition}

The positions are given, and closely.  For $\zeta$, with the parameters of the operator of
\cite{gonzalez25hp} and the zeros used only to measure: the mean error of $T^*(n)$ is about
$4.9$ mean spacings; with the exact envelope, which uses no prime and no zero, it is $0.21$
spacings ($n\le1000$; the median is $0.19$), a relative error from $3\cdot10^{-3}$ at low height to $5\cdot10^{-5}$ at $n=2000$
(\texttt{corrected\_operator.py} \tier[\testref{test\_corrected\_operator\_golden\_values, test\_corrected\_operator\_matches\_manuscript\_claim}]{N}).  What separates the structural positions from the
zeros is a fraction of a spacing.

\begin{proposition}[Membership of the primes, entry of the zeros, proved]\label{prop:membership-proved}
\begin{enumerate}[label=\textup{(\roman*)}]
  \item \emph{The primes.}  Every prime $p>5$ lies in one of the eight classes of $(\Z/20)^\times$,
    on which $G$ is a homomorphism; $\chiv$ separates the inert classes, the free Klein orbit
    $\{3,7,13,17\}$, from the split ones; the primes reach $\tfrac12$ through the conductor,
    $|\operatorname{im}G|/\phi(20)=\tfrac48=\tfrac12$; and for every zero the arithmetic reaches it in one statement.
  \item \emph{The zeros.}  For every zero the centre of $\{\rho,1-\rho\}$ is exactly $\tfrac12$,
    the positive eigenvalue of $\hat\Omega$; the zero enters by three doors --- the centre in the
    spectrum, opposite golden valuations, and $r(\rho)\,r(1-\rho)=1$, the multiplicative half of \Ecore\ ---
    and the apex $\tfrac12$ is reached at once by five legs: the primes, the zeros, the operator,
    the monoid and the archimedean place.
  \item \emph{What is discharged.}  Hecke's companion is a theorem, and membership is derived from the bridge, the strip and the
    companion rather than assumed; of these, only the last link of the bridge is open.
  \item \emph{What remains, in the corpus's own words.}  The apex is the centre of the zero, not
    the membership of the zero itself; everything descends but the verdict; and what remains is the additive half of \Ecore.
\end{enumerate}
{\raggedright \tier[F1PaperClosure.the\_primes\_belong, F1PaperClosure.the\_zeros\_belong, F1PaperClosure.what\_is\_discharged, F1PaperClosure.what\_remains\_is\_the\_verdict]{P}, on declarations of the corpus \tier{P}\par}
\end{proposition}

What remains is therefore precision, of two kinds: of the real part, the radius $0$ that the additive half of \Ecore\ would give --- the open statement of \cref{ssec:core-forms}, and the only one the line
needs; and of the height, the correction, which makes the operator exact.  They are one fluctuation read from
two sides: by the explicit formula the residual of the primes of each class around its
structural share and the correction of the heights carry the same data.

The scale $\tfrac12$ that each prime carries as its modulus $p^{-1/2}$ (\cref{prop:torus-line})
is also the scale of the sums over the primes.  What that scale is can be said on the side of the
primes and on the side of the zeros, for the symbols that read the index of $\golden$
(\cref{prop:phi-layers}).

\begin{proposition}[What $\sqrt X$ is]\label{prop:sqrt-x}
Let $a_p$ be the trace of Frobenius at $p$ of the representation of one of the layers $d=2,3,5$
of \cref{prop:phi-layers} --- the forms of levels $80$ and $135$ and the representation of the
quintic layer.  At a prime split in $\Q(\sqrt5)$ the first is $(\tfrac rp)+(\tfrac{1-r}p)$, with $r$ and $1-r$ the two roots of $x^2-x-1$ modulo $p$; at a
prime split in $\Q(\sqrt5,\sqrt{-3})$ the second is $2\Re(\golden/\mathfrak p)_3$; at $p\neq5$ the
third is the number of roots of $x^5+5x^3+5x-1$ modulo $p$ minus $1$.  Let $|a_p|\leq A$, and let $\Sigma(X)=\sum_{p\leq X}a_p\log p$ and
$\mathcal E(X)=\sum_{p\leq X}a_p^2\log^2p$.
\begin{enumerate}[label=\textup{(\roman*)}]
  \item The values of each symbol are roots of unity, and each layer of the index is their
    average in $\mathbb F_p$: for $p$ split, $d\mid p-1$ with $(p-1)/d$ even, and $u=r^{(p-1)/d}$,
    $\sum_{k<d}u^k=d$ if $d\mid i_p$ and $0$ otherwise.
  \item The energy is of the size of $X$: $\mathcal E(X)\leq A^2\log4\cdot X\log X$, and
    $\tfrac1X\sum_{p\leq X}a_p^2\log p\to1$ for the three symbols.
  \item If the cross terms of $\Sigma(X)^2$ do not accumulate --- for every $\varepsilon>0$,
    $\Sigma(X)^2\leq C_\varepsilon X^{2\varepsilon}\mathcal E(X)$ --- then
    $|\Sigma(X)|\leq C'_\varepsilon X^{1/2+\varepsilon}\sqrt{\log X}$: the scale $\tfrac12$ is the
    square root of the energy.
  \item On the side of the zeros, a zero $\rho=\beta+i\gamma$ contributes $X^\rho/\rho$ to the
    explicit formula, with $X^\rho=X^\beta e^{i\gamma\log X}$: a rotation of modulus $X^\beta$.  For a prime $p$, its
    value $\lambda(\rho)=p^{\rho-1/2}$ lies on the unit circle if and only if $\beta=\tfrac12$.  So
    the scale of \textup{(iii)}, read from the primes, and the modulus of the terms of the zeros
    agree exactly on the line.
\end{enumerate}
{\raggedright \LeanInline{W4Isometria.geom\_sum\_root\_of\_unity},
\LeanInline{W4Isometria.split\_layer\_indicator} (i);
\LeanInline{W4Isometria.symEnergy\_le} (ii, the bound);
\LeanInline{W4Isometria.sqrt\_law\_of\_pythagoras} (iii);
\LeanInline{W4Isometria.specLam\_on\_circle\_iff} (iv) \tier{P}; the hypothesis of (iii),
\LeanInline{W4Isometria.PythagorasLaw}, open; the limit in (ii) by Schur orthogonality in the
Galois group and Chebotarev, measured $0.993$, $0.998$ and $1.000$ up to $10^6$, and
$\Sigma(X)^2<\mathcal E(X)$ at every scale measured (\texttt{v12\_pitagoras\_raiz\_X.py}) \tier{N}.  From the primes up to $10^6$, the % pending(D7)
exponent $\nu$ that best fits $\Sigma(X)\approx X^\nu R(X)$, with $R$ written with the computed
zeros of the forms of levels $80$ and $135$, is $0.500$ and $0.5065$, within $\pm0.01$; the first
zero alone, with its term multiplied by $X^{\delta}$, gives $\delta=0.004\pm0.03$; that zero carries
$31\%$ and $48\%$ of $\sum_\rho1/|\rho|^2$ (\texttt{escala\_ceros\_bajos.py}) \tier{N}\par}
\end{proposition}

So the scale $\tfrac12$ has two readings that meet: from the primes, the square root of an
energy that the structure fixes; from the zeros, the modulus of a rotation on the circle.  What is
open is the hypothesis of (iii), that the phases, roots of unity, do not accumulate; since the
energy is of size $X\log X$, that hypothesis is the bound itself up to $X^\varepsilon$, so (iii)
names where the scale sits and does not reduce it to something smaller.

The passage from the clock to the size takes place inside the $n$-gon of each prime
(\cref{prop:index-ngon}): in the column of
clock $e_p$ the size of $p$ is its index, and the primes of large index weigh $O(\sqrt X)$ with no
hypothesis.  The model in which $\golden$ falls on a uniform vertex of the $n$-gon of each prime
gives the measured mass of index above $j$ to within $8\%$ for $3\leq j\leq1000$, and the layer
densities $0.1237$ and $0.1666$ against $0.1243$ and $0.1667$; the forms of level $135$ and of the
quintic layer are $3([3\mid i_p]-\tfrac13)$ and $5([5\mid i_p]-\tfrac15)$ in the sums up to $10^9$,
to $2\cdot10^{-10}$ and $6\cdot10^{-12}$ (\texttt{v30\_vertice\_al\_azar.py}, % pending(D7)
\texttt{v33\_capas\_1e9.py}).  What the geometry does not give is that $\golden$ spreads over % pending(D7)
the vertices uniformly as $p$ varies, with an error of size $\sqrt X$: that is the
equidistribution of Chebotarev in the Kummer fields $\Q(\zeta_i,\golden^{1/i})$, the piece that
Hooley obtains from the generalised Riemann hypothesis~\cite{hooley67}.


That hypothesis is about the energy at every scale, and for finitely many frequencies it is a
theorem in both directions; the constant that survives infinitely many zeros depends only on
how many zeros there are per unit height, not on their gaps.

\begin{proposition}[Energy at every scale]\label{prop:energy-scale}
Let $Z(u)=\sum_jc_je^{i\gamma_ju}$, $u=\log X$, with finitely many distinct real $\gamma_j$.
\begin{enumerate}[label=\textup{(\roman*)}]
  \item In every window $[u_0,u_1]$ of length $L\geq1$,
    $\bigl|\tfrac1L\int_{u_0}^{u_1}|Z|^2-\sum_j|c_j|^2\bigr|\leq3\mathcal K$, with
    $\mathcal K=\sum_{j\neq k}|c_j||c_k|/(1+|\gamma_j-\gamma_k|)$; and the mean over $[0,U]$ tends to
    $\sum_j|c_j|^2$.
  \item If $|c_j|\leq C_0/(1+|b_j|)$ for integers $b_j$ with $|\gamma_j-b_j|\leq1$, and each
    integer carries at most $B(1+|n|)^{1-s}$ of the $b_j$, with $s>\tfrac12$, then
    $\mathcal K\leq3C_0^2B^2\sum k_s$, with $k_s(n,m)=(1+|n|)^{-s}(1+|m|)^{-s}/(1+|n-m|)$ summed over
    $\Z^2$: the same bound for every finite family.
  \item If moreover $\sum_j|c_j|^2\leq V$ and a function $F$ is within $\epsilon$ of $Z$ in mean
    square on the window, then for every $\lambda>0$ the mean energy of $F$ there is at most
    $(1+\lambda)\,(V+9C_0^2B^2\sum k_s)+(1+1/\lambda)\,\epsilon^2$.
  \item Conversely, a term $ce^{\delta u}e^{i\gamma u}$ with $\delta>0$, $c\neq0$, plus a
    continuous bounded remainder has unbounded mean energy; and the factor $e^{2\delta u}$ that a pair with
    $\beta-\tfrac12=\delta>0$ puts on its own term has mean above $1$ on every window
    $[u_0,u_1]\subset[0,\infty)$.
\end{enumerate}
{\raggedright \LeanInline{W4Isometria.LandauPlancherel.cramer\_uniform},
\LeanInline{W4Isometria.LandauPlancherel.mean\_tendsto\_on\_line} (i);
\LeanInline{W4Isometria.CramerConteo.cramer\_K\_bounded} (ii);
\LeanInline{W4Isometria.FormulaExplicita.energy\_bounded\_of\_explicit\_formula} (iii);
\LeanInline{W4Isometria.LandauPlancherel.offline\_energy\_unbounded},
\LeanInline{W4Isometria.off\_line\_excess\_pos} (iv) \tier{P}.  With $c_j=1/\rho_j$ and
$V=2\Re\,(\Lambda'/\Lambda)(1,\chiv)$, the logarithmic derivative of the completed $L$-function at $1$, (iii) bounds the energy of the normalised error of the explicit
formula given the line, the counting law of Riemann--von Mangoldt and the truncated explicit
formula; the last two are classical\ and not formalised here, and the line is the
open hypothesis.  For $\zeta$, the mean of $|(\psi(x)-x)/\sqrt x|^2$ in $\log x$ is bounded at every scale
if and only if $\zeta$ has no zeros with $\Re s>\tfrac12$: under the line by Cram\'er's
mean-square theorem \cite[\S13.1]{montgomeryvaughan07}, and conversely by Mellin--Plancherel;
likewise for the $L$-functions here.  With $2000$
zeros of $\zeta$, $\sum|c_j|^2=0.04531$ against the value $2+\gamma_{\mathrm E}-\log4\pi=0.04619$ of
$2\Re\,\xi'/\xi(1)$ for $\zeta$, with $\xi(s)=\tfrac12s(s-1)\pi^{-s/2}\Gamma(\tfrac s2)\zeta(s)$ the completed zeta function and $\gamma_{\mathrm E}$ Euler's constant, and $\mathcal K=0.158$ ($\approx0.19$ extrapolated); for the form of level $80$ with the primes up to $10^9$, the mean energy of the
error is $0.626$ and the part beyond its $66$ computed zeros $0.066$, which does not grow from
decade to decade; the $16$ twists of the three families stay within $10\%$ of their $V$ up to
$10^9$, and their mutual correlations are those of unrelated $L$-functions; from the primes alone, for
the form of level $80$, the energy up to $10^9$ leaves to a single pair off the line at height
$\gamma$ at most $\beta-\tfrac12\leq0.0013$ for $\gamma=1$ and $0.046$ for $\gamma=10$, if the rest of
the spectrum is unchanged (\texttt{v20}--\texttt{v26}) \tier{N}\par}
\end{proposition}

\begin{conjecture}[Closure conjecture]\label{conj:closure}
There is a law of the primes by classes, derived from
the structure --- the tower, the eight classes of
$(\Z/20)^\times$ and the operator --- without taking the
primes one by one.  It is obtained in two steps.  On
each level $n$ of the clock of $\golden$, the product
formula on $N(\Phi_n(\golden))$ equates the archimedean
size of the level, $\phi(n)\log\golden$ up to a bounded
term, with the $p$-adic weights
of the primes of clock $n$
(\cref{prop:clock-cyclotomic}).  From the clock to the
size, the law of the index orders those primes by size:
inside each column of the clock the size of a prime is
its index, the primes of large index weigh $O(\sqrt x)$,
and the rest is the distribution of $\golden$ among the
vertices of the $n$-gon of each prime
(\cref{prop:index-ngon}), read layer by layer by the
power-residue symbols of $\golden$, with values in roots
of unity and energy of size $x$
(\cref{prop:phi-layers,prop:sqrt-x}), its deviation in
each layer having bounded energy at every scale, which
for finitely many frequencies is \cref{prop:energy-scale}.  The law
\begin{enumerate}[label=\textup{(\roman*)}]
  \item predicts the residual of each class $a$ around its structural share,
    $R_a(x)=\psi(x;20,a)-x/8$, with $\psi(x;20,a)=\sum_{n\le x,\ n\equiv a\ (20)}\Lambda(n)$;
  \item predicts it at scale $\tfrac12$, $R_a(x)=O_\varepsilon(x^{1/2+\varepsilon})$ for every
    $\varepsilon>0$ --- the $\tfrac12$ that each prime already carries as its modulus $p^{-1/2}$
    (\cref{prop:torus-line}); and
  \item gives, window by window --- for every $L>0$ and every test function $g=h*\tilde h$ with
    $h$ supported in $[-L/2,L/2]$, Weil's class --- the correction with which the heights are
    exactly $t_n=E^{-1}(n-\text{correction})$, smoothed at resolution of order $1/L$ in $t$; the
    heights near $T$, whose spacing is $2\pi/\log(T/2\pi)$, are separated once $L$ is of order
    $\log T$.  The correction is assembled from the residuals of the classes, each prime power
    entering through its character, its modulus and its angle on the circle
    (\cref{prop:what-fluctuation-takes} below).
\end{enumerate}
\end{conjecture}
{\raggedright\LeanInline{F1PaperClosure.ExactClosure}, open, with the residual of each class defined
from $\Lambda$ (\LeanInline{F1PaperClosure.trueResidual});
with clause~(iii)'s assembly as the limit of its
partial sums,
\LeanInline{F1PaperClosure.ExactClosureAssembledLim},
\LeanInline{F1PaperClosure.CorrectionIsLimit}, open;
together with the scale $\tfrac12$ of a finite family
of symbols, given as a parameter,
\LeanInline{F1PaperClosureClock.ExactClosureThroughClock}
open, a hypothesis that implies the limit form; the
second step in the form of the clock,
\LeanInline{F1PaperClosureClock.ClockLayerEnergyBounded}
open, uniform over a finite set of layers, which
bounds the energy of every real combination of the
layers
(\tier[F1PaperClosureClock.family\_energy\_of\_clock\_layers]{P}; that the families and their twists are such
combinations is \tier{N}\ beyond level $80$); each form implies the one before,
\tier[F1PaperClosureClock.assembledLim\_of\_throughClock, F1PaperClosure.exactClosure\_of\_assembledLim]{P}; the earlier form with an unconditional sum,
\LeanInline{F1PaperClosure.ExactClosureAssembled}, is
vacuous in (iii)\par}


\begin{remark}[The status of each clause]\label{rmk:closure-status}
For the true residuals, clause~(ii) is equivalent to GRH for the characters modulo $20$
(the classical equivalence, \LeanInline{F1PaperClosure.ClassicalHalfScale}), and given it the
explicit formula writes the correction from the primes by classes.  What the conjecture asserts
beyond that classical content is the source: that the residuals of~(i) and the correction of~(iii)
are produced by the structure, at the resolution~(iii) names, without enumerating primes.  The Lean
predicate
\LeanInline{F1PaperClosure.ExactClosureAssembledLim}
formalises clauses~(i)--(iii) --- the law's residual
is the true one, it is at scale $\tfrac12$, and the
correction that makes the heights exact is the limit
of the partial assemblies of the residuals
(\LeanInline{F1PaperClosure.CorrectionIsLimit}).  The
first form, \LeanInline{F1PaperClosure.ExactClosureAssembled},
takes the assembly as an unconditional sum; with the
true residual that series is not absolutely
summable for $T\ne0$, and its terms vanish for
$T=0$, so the sum is $0$ for every $T$ and its
clause~(iii) reduces to $\Im\rho_n=E^{-1}(n)$.  The form with the
limit implies the first statement of the closure
(\tier[F1PaperClosure.exactClosure\_of\_assembledLim]{P}).  Neither formalises the derivation from
the structure, nor the windows of
clause~(iii), so both are weaker than the conjecture
as stated.
  The
mathematical form proposed for the source is \cref{conj:prismatic-route}.  On the real part, the
conjecture has only half to supply: the multiplicative half $r(\rho)\,r(1-\rho)=1$ is proved on every
zero (\cref{prop:two-halves}(i)), so what it must give there is the additive half, $t(\rho)=2$, and
it does so through \cref{cor:closure}; clause~(iii) supplies the other precision, that of the
heights.  The law is obtained in the two steps of the
statement.  The first is proved
(\cref{prop:clock-cyclotomic}).  In the second the
energy is proved (\cref{prop:sqrt-x}) and the scale is
open: if each power-residue symbol of a finite family is
at scale $\tfrac12$ (\LeanInline{W4Isometria.UnifiedLaw60})
and the error of a class of indices is a combination of
their sums plus a remainder at scale $\tfrac12$, the
error of the class is at scale $\tfrac12$, in count and
in proportion
(\tier[W4Isometria.class\_error\_of\_unified, W4Isometria.proportion\_error\_bound]{P}).  The decomposition and the remainder are
hypotheses, not constructed.  In the form of the clock
the open statement is one: uniformly over a finite set of
layers $i$, the deviation of the index from the model of the uniform
vertex has bounded energy at every scale
(\LeanInline{F1PaperClosureClock.ClockLayerEnergyBounded});
the part of large index is unconditional
(\cref{prop:index-ngon}(ii)), the layers $2$, $3$ and $5$
are the forms of levels $80$ and $135$ and of the quintic
layer (\cref{prop:index-ngon}(iv) for level $80$;
numerically for $135$ and the quintic layer), and up to $10^9$ the
twenty layers $2\leq i\leq60$ measured show no systematic
term and a fluctuation of mean exponent $0.520$
(range $0.42$--$0.59$) (\texttt{v33\_capas\_1e9.py}).  With the family taken to % pending(D7)
be the characters modulo $20$ and the power-residue
symbols of $\golden$, the hypothesis is clause~(ii) and
GRH for the $L$-functions of the forms of levels $80$ and
$135$ and of the quintic layer.  The evidence is the scale $\tfrac12$ fixed at the low
zeros (\cref{prop:sqrt-x}).  This remains open.
\end{remark}

\begin{remark}[The window and the clock]\label{rmk:closure-window}
In a window $L$ the prime side of the explicit formula for $\zeta_K$ is finite: only the prime
ideals of norm at most $e^L$ enter, and the clock gives them exactly, level by level.  A prime
ideal of level $n$ has norm $n\,i+1$ if it lies over a split prime and $(n\,i/2-1)^2$ if it lies
over an inert one, with $i$ its index; norm at most $e^L$ forces $i\le 2(e^L+1)/n$, and the sieve of
level $n$ with that bound returns all of them and nothing else
(\LeanInline{CierrePorAltura.closure\_support}; at $e^L=10^6$ the sieve returns the $78\,508$ prime
ideals and the prime side summed level by level equals the direct sum, \texttt{retro66.py}). % pending(D7)
So clause~(iii), in this form, asks of the primes nothing that the clock does not compute.
What \cref{conj:closure} asserts beyond it is the source of the law in (i)--(ii), and, for the
line, the positivity of the window functional.  That positivity is a global cancellation: the
archimedean and polar part alone is indefinite, its negative direction being that of the pole,
and in the window $e^L=10^3$ removing any one of $108$ of the $127$ levels breaks it; along the
direction of the pole a single prime ideal contributes with the sign of its size, negative below
$\log N\approx0.38\,L$ and positive above, whatever its level (\texttt{retro67.py}). % pending(D7)
The norms and the bound on the index, \tier[PasosDelTexto.window\_split, PasosDelTexto.window\_inert]{P}.
\end{remark}


\begin{proposition}[The amplitude of clause~(ii)]\label{prop:closure-amplitude}
For a character $\chi$ modulo $20$ let $\chi^*$ be the primitive character inducing it, of
conductor $q$ and parity $\mathfrak e\in\{0,1\}$ ($\chi^*(-1)=(-1)^{\mathfrak e}$), let
$\xi(s,\chi^*)=(q/\pi)^{(s+\mathfrak e)/2}\,\Gamma(\tfrac{s+\mathfrak e}2)\,L(s,\chi^*)$, and let
$V_\chi=2\Re\,\xi'/\xi(1,\chi^*)$; for the trivial character, $V=2+\gamma_{\mathrm E}-\log4\pi$, with $\gamma_{\mathrm E}$
Euler's constant.
\begin{enumerate}[label=\textup{(\roman*)}]
  \item $V_\chi=\log(q/\pi)+\Re\,\tfrac{\Gamma'}{\Gamma}(\tfrac{1+\mathfrak e}2)
    +2\Re\,\tfrac{L'}{L}(1,\chi^*)$, and
    \[
      \tfrac{L'}{L}(1,\chi^*)=-\log q
      -\frac{\sum_b\chi^*(b)\,\gamma_1(b/q)}{\sum_b\chi^*(b)\,\gamma_0(b/q)},
    \]
    with $\gamma_0=-\Gamma'/\Gamma$ and $\gamma_1$ the generalised Stieltjes constant
    \cite{berndt72,dilcher92}: the amplitude is fixed by the conductors $1,4,5,20$, the parities
    and Hurwitz values at the fractions $b/q$; no prime is enumerated, and only the ramified primes
    $2$ and $5$ enter, through the conductors.  Over the eight characters
    $\sum_\chi V_\chi=2.4056$.
  \item Over the nontrivial zeros $\rho=\beta+i\gamma$ of $L(s,\chi^*)$,
    \[
      \sum_\rho\frac1{|\rho|^2}-V_\chi
      =\sum_\rho\frac{2(\beta-\frac12)^2}{|\rho|^2\,|1-\bar\rho|^2}\;\geq\;0,
    \]
    so $\sum_\rho1/|\rho|^2=V_\chi$ if and only if all of them lie on the line: the amplitude law
    of the zeros, whose defect is the distance of each pair $\{\rho,1-\bar\rho\}$ to the line.
  \item If the nontrivial zeros of the eight $L$-functions lie on the line and their nonnegative
    ordinates are linearly independent over $\mathbb{Q}$, then for every class $a$, $R_a(x)/\sqrt x$ has logarithmic
    mean $0$ and logarithmic variance $\tfrac1{64}\sum_\chi V_\chi=0.0376$; with
    $\theta(x;20,a)=\sum_{p\leq x,\,p\equiv a}\log p$ in place of $\psi$ the variance is the same and the mean is $-n_2(a)/8$, with $n_2(a)$ the number of square
    roots of $a$ modulo $20$.
\end{enumerate}
{\raggedright (i); (ii) termwise
\LeanInline{W4Isometria.pair\_defect}, and for finite families
\tier[W4Isometria.amplitude\_defect, W4Isometria.amplitude\_eq\_iff\_line]{P}, the sum by Hadamard's product; (iii) \cite{rubinsteinsarnak94}.  On a
finite range the cross terms between characters make single classes deviate; in the mean over the
eight classes they cancel exactly (\tier[W4Isometria.class\_parseval]{P}, on the
centred values).  Over the primes up to $10^8$, with $\theta$, the logarithmic means are $-0.497$
and $-0.500$ in the classes $1$ and $9$ and below $0.04$ in absolute value in the others, and the
mean of the eight variances is $0.0379$ on $[10^4,10^8]$, within $3.5\%$ of $0.0376$ on five other
ranges, while single classes deviate by $-54\%$ to $+68\%$ (\texttt{amplitud\_clausula\_ii.py})
\tier{N}\par}
\end{proposition}

So the true residual of clause~(i), at the scale of
clause~(ii) and with independent ordinates, has a
definite amplitude: where \cref{prop:sqrt-x} gives the
scale $\tfrac12$ as the square root of an energy, this
gives its constant, which the structure fixes without
enumerating the primes.  In the language of W3
(\cref{prop:intersection-form}), $V_\chi$ plays the role
of the Weil form at one test function,
$h(x)=e^{-|x|/2}$: its transform
$\hat h(\gamma)=1/(\tfrac14+\gamma^2)$ is
$1/(\rho(1-\rho))$ at $\gamma=(\rho-\tfrac12)/i$, and
$\sum_\rho1/(\rho(1-\rho))=V_\chi$ by Hadamard's product
.  With $h=g*\tilde g$ and $g(x)=e^{-x/2}$ on
$x>0$, the Gram value of
\cref{prop:intersection-form}(i) is
$\sum_\rho|\hat g(\gamma_\rho)|^2=\sum_\rho1/|\rho|^2$.
So (ii) says that the Weil value and the Gram value at
$h$ agree exactly on the line, and $V_\chi>0$ for the
eight characters is the positivity at that vector,
fixed by the conductors and the Hurwitz values.  The
function $h$ sits at the edge of Weil's class: its
prime side is $L'/L(1,\chi^*)$, a conditionally
convergent sum over the primes, which (i) computes
without them~\cite{berndt72,dilcher92}.  Clause~(iii) asks for more, every
height and not their sum $\sum_\rho1/|\rho|^2$.

\begin{corollary}[The closure, conditional on the conjecture]\label{cor:closure}
\Cref{conj:closure} gives the line by (ii), through the classical equivalence between residuals
at scale $\tfrac12$ and zeros on the line.  By (iii) it gives the heights exactly; and since the
landing of every zero on the spectrum of
$\hat\Omega$ is equivalent to the line, it gives
that last link of the bridge through (ii).
  That the heights it gives be
distinct for distinct zeros would also give the line (\cref{thm:shared-ordinate}), but distinctness
is not a consequence of exactness: no two zeros sharing an ordinate is itself one of the forms of the
open statement (\cref{thm:topological-proof}~(iv)).  It is stated for any $L$ of character modulo $20$ and for $\zeta$, whose operator is
that of \cite{gonzalez25hp}.  \tier{P}%
\end{corollary}
{\raggedright \tier[F1PaperClosure.the\_line\_from\_exact\_closure, F1PaperClosure.heights\_exact\_from\_exact\_closure, F1PaperClosure.landing\_from\_exact\_closure, F1PaperClosure.the\_line\_for\_chi5, F1PaperClosure.the\_line\_from\_injective\_heights\_chi5, F1PaperClosure.the\_line\_from\_injective\_heights, F1PaperClosure.the\_line\_from\_exact\_closure\_by\_heights]{P}; the equivalence of the first way,
\LeanInline{F1PaperClosure.ClassicalHalfScale}\par}

\begin{proposition}[The height fixes the phase, the isometry the modulus]\label{prop:phase-modulus}
Under the Frobenius of $p$, translation by $\log p$, the mode of a zero $\rho=\beta+i\gamma$ is
multiplied by $p^{-(\beta-1/2)}e^{i\gamma\log p}$: the height gives the phase and the distance to the
line gives the modulus.
\begin{enumerate}[label=\textup{(\roman*)}]
  \item If the heights of the zeros are given by a formula injective in the index, no two distinct
    zeros share an ordinate, and every zero lies on the line; this holds for $L(s,\chiv)$ and for
    every $L$ with functional equation and real coefficients, $\zeta$ among them.
  \item An eigenvalue of an isometry has modulus $1$; so if the mode of each zero is an eigenvector
    of an isometry, its modulus $p^{-(\beta-1/2)}$ is $1$ and every zero lies on the line.
\end{enumerate}
{\raggedright \tier[F1PaperClosure.no\_shared\_ordinate\_of\_injective\_heights, F1PaperClosure.the\_line\_from\_injective\_heights, F1PaperClosure.eigenvalue\_of\_isometry\_norm\_one, F1PaperClosure.the\_line\_from\_isometric\_modes]{P}\par}
\end{proposition}

The two parts are one statement in two languages.  A zero off the line and its partner
$1-\bar\rho$ share the height, hence the phase, and have reciprocal moduli $R$ and $1/R$ --- the
multiplicative half $r(\rho)\,r(1-\rho)=1$ read as the pair of a boost (\cref{ssec:bridge-boosts}).
In the language of Klein that pair is two zeros with one ordinate; in that of the isometry, two
eigenvalues with one phase and moduli $R$ and $1/R$, which no isometry has; in that of the
coordinate, $r\neq1$, the image under $\ell(\rho)=1-1/\rho$ off the unit circle (\cref{prop:li-map}).
The operator of \cite{gonzalez25hp} is self-adjoint, and its isometric form is of rotation type
(\cref{thm:operator-forms}); what neither part supplies is that the modes of the zeros are the
modes of that operator, the last link of the bridge.  Exact heights are what would supply it: the
height fixes the phase, the isometry fixes the modulus, and together they give $r=1$.  What the structure has to provide is therefore precise: heights that are exact, a correction
that jumps by one at each zero; that they are distinct for distinct zeros is the form~(iv) of the open
statement (\cref{thm:topological-proof}).

\emph{What the law has to supply.}  Replacing each prime by the count of its window of width
$\delta$ in $\log p$ --- what a law that knows how many primes of each kind fall in each window,
but not the primes one by one, would give --- the median error of the corrected operator
($n=100$ to $3000$) is $0.194$ at $\delta=10^{-2}$, $0.170$ at $3\cdot10^{-3}$, $0.089$ at $10^{-3}$,
$0.041$ at $3\cdot10^{-4}$ and $0.025$ at $10^{-4}$, against $0.20$ with no correction and
$0.021$ prime by prime (\texttt{law\_resolution.py} \tier[\testref{test\_law\_resolution\_golden\_values, test\_law\_resolution\_matches\_manuscript\_claim}]{N}).  Taken as a sum over the primes up to a fixed $X$, the
correction does not converge: for $n=100$ to $3000$ the mean error has a minimum at a finite $X$ and grows past it ---
$0.069$ at $X=20$ and $0.23$ at $X=10^6$ for the plain sum, $0.037$ at $X\approx500$ and $0.12$ at $X=10^6$ with the
weight $1-\log n/\log X$ on each prime power $n$ (\code{x\_sweep.py} \tier[\testref{test\_law\_resolution\_golden\_values, test\_law\_resolution\_matches\_manuscript\_claim}]{N}).  The error falls between
$\delta=10^{-3}$ and $3\cdot10^{-4}$, the scale $1/T$ of the heights measured ($T\approx240$ to
$3700$): at % pending(D7)
height $T$ the law has to resolve a scale of order $1/T$.  Taking the primes one by one is not
the law: it is the explicit formula, and as a construction it is circular; the classes are the part of the law that is not, and they are needed --- removing any one of the eight worsens the error (\code{classes\_test.py} \tier[\testref{test\_law\_resolution\_golden\_values, test\_law\_resolution\_matches\_manuscript\_claim}]{N}).
What the law has to predict, and nothing else, is fixed by the local factors.

\begin{proposition}[What the fluctuation takes from each prime]\label{prop:what-fluctuation-takes}
\begin{enumerate}[label=\textup{(\roman*)}]
  \item At $s=\sigma+it$ the logarithm of the Euler factor at $p$ is $-\log(1-z)$ with
    $z=u\,e^{-i\theta}$, $u=\chi(p)\,p^{-\sigma}$, $\theta=t\log p$, and its imaginary part --- the
    contribution of $p$ to the argument of $L$, hence to $S=\arg L/\pi$ --- is
    $-\sum_k u^k\sin(k\theta)/k$: it depends on $\chi(p)^k$, on $p^{-k\sigma}$, on $k$, and on the angle
    $k\,t\log p$ measured on the circle, and on nothing else of $p$.
  \item The local factor sees only the class: $\det(1-tM)=(1-t)(1-\chi t)$ (\cref{thm:frob-eigen}).
  \item The angular part of $\golden$ at $p$ is not a function of the class: $7$ and $47$, in the
    class $7$, have angular orders $16$ and $32$; $13$ and $113$, in the class $13$, have $28$ and $76$.
    \item What the angular order does carry is the
    position of $p$: $e_p\,i_p=p-1$ if $p$ splits and
    $2(p+1)$ if $p$ is inert, so $p$ is recovered from
    $e_p$ and the index $i_p$ of
    \cref{prop:phi-layers}.

\end{enumerate}
{\raggedright \tier[F1PaperPadic.local\_phase\_contribution, F1FrobeniusEigen.euler\_factor\_is\_char\_poly, F1PaperPadic.angular\_order\_not\_a\_function\_of\_the\_class, F1PaperPadic.position\_from\_angular\_order]{P}; (i) term by term for $|u|<1$, the sum over $p$ on the line by the explicit formula (\cref{thm:explicit})\par}
\end{proposition}

So the fluctuation takes from each prime its character $\chi(p^k)$, its modulus $p^{-k/2}$ and its
frequency $\log p$, and measures it as an angle on the
circle; the angular and radial structure of $\golden$ at
each prime (\cref{ssec:delta-polar,ssec:classes-angular})
is structure of the ring, not an ingredient of the
correction prime by prime; through the clock and the
index it is the route by which \cref{conj:closure}
proposes to obtain the law.  The index is $1$ for $66.5\%$ of the primes below $2\cdot10^5$ and at most $6$ for
$90.3\%$, so the angular order carries the position of the prime; and once the position is
accounted for, it does not explain the use of any prime in the correction
(\texttt{t1\_index.py}, \texttt{t2\_loo.py} \tier{N}).

\begin{proposition}[The index as a degeneracy]\label{prop:index-degeneracy}
Let $M=\bigl(\begin{smallmatrix}0&1\\1&1\end{smallmatrix}\bigr)$ be multiplication by $\golden$ on $\Zphi$, let
$N\ge3$ be odd, let $W_N$ be the Weil representation of $\mathrm{SL}_2(\Z/N)$ on $\C^N$, normalised to be
multiplicative, and let $U_N=W_N(M^2)$ and $e$ the order of $M^2$ modulo $N$.
\begin{enumerate}[label=\textup{(\roman*)}]
  \item $U_N^{\,e}=1$, and the eigenvalues of $U_N$ are $e$-th roots of unity.
  \item For a prime $p\ne5$, every $e$-th root of unity but one has multiplicity $m=(p-\chiv(p))/e$, the index of
    $\golden^2$ at $p$.
  \item If $p$ splits the exceptional eigenvalue is $1$, with multiplicity $m+1$; if $p$ is inert it is $-1$, the
    half-turn, with multiplicity $m-1$.  In both cases $p=e\,m+\chiv(p)$.
  \item At $N=5^k$ the spectrum is simple and fills half the circle: $e=2\cdot5^k$, and, writing $j=5^vj'$ with
    $5\nmid j'$, the root $\zeta_e^j$ is an eigenvalue for $5^k\nmid j$ exactly when $\chiv(j')=+1$; of the two
    multiples of $5^k$, $j=0$ is an eigenvalue and $j=5^k$, the half-turn, is not.
\end{enumerate}
\end{proposition}

\begin{proof}
(ii)--(iii).  $M^2$ lies in a maximal torus $T$ of $\mathrm{SL}_2(\mathbb F_p)$, cyclic of order $q=p-\chiv(p)$:
split if $p$ splits, non-split if $p$ is inert.  Restricted to $T$, the Weil representation contains every
character of $T$ once, except the quadratic character $\eta$ of $T$, which occurs twice if $T$ is split and not at all
if it is non-split \cite{gurevichhadani}.  Restricting further to $\langle M^2\rangle$, of index $m=q/e$, each
of the $e$ characters of $\langle M^2\rangle$ receives $m$ characters of $T$, and $\eta$ restricts to the
character $M^2\mapsto\eta(\golden^2)=\golden^{q}$, which is $\golden^{p-1}=1$ if $p$ splits and
$\golden^{p+1}=N(\golden)=-1$ if $p$ is inert.  \tier{P}\
\LeanInline{IndiceComoDegeneracion.count\_residue}, \LeanInline{IndiceComoDegeneracion.sigma\_phase}.
\par
(iv).  The multiplicity of $\zeta_e^j$ is $\frac1e\sum_{n<e}\operatorname{tr}U^n\zeta_e^{-jn}$.  With
$\operatorname{tr}U^n=1$ for odd $n$ and $\operatorname{tr}U^{2m}=5^v\sqrt5\,\chiv(m')$ for $m=5^vm'$,
$5^k\nmid m$ (and $5^k$ at $m=0$), the odd terms give $\pm\frac12$ at $j=0,5^k$ and $0$ elsewhere; in the even
terms the part with $v_5(m)=v$ is a Gauss sum of $\chiv$ at level $5^{k-v}$, nonzero only when $v=k-1-v_5(j)$,
where it equals $5^k\chiv(j')$, since $\chiv$ is even and its Gauss sum is $\sqrt5$.  So the multiplicity is
$\frac12(1+\chiv(j'))$ for $5^k\nmid j$, $1$ at $j=0$ and $0$ at $j=5^k$.  Given the traces.
\par
{\raggedright \emph{The traces at $5^k$.}  If the entry $a$ of $g=\bigl(\begin{smallmatrix}a&b\\c&d\end{smallmatrix}\bigr)$ is a unit,
the Bruhat decomposition of $g$ gives
$\operatorname{tr}W(g)=\bigl(\tfrac a5\bigr)^k5^{-k}\sum_{x,\xi}e\bigl(q(x,\xi)/5^k\bigr)$ with
$q=\tfrac{b}{2a}x^2+(1-a^{-1})x\xi-\tfrac{c}{2a}\xi^2$, whose discriminant is $4AC-X^2=-(t-2)/a$, $t=\operatorname{tr}g$.
For $g=M^{2n}$, $t-2=L_n^2$ if $n$ is odd and $5F_n^2$ if $n$ is even.  If $n$ is odd the form is nondegenerate, the sum is
$\bigl(\tfrac a5\bigr)^k5^k$ and the trace is $1$.  If $n=2m$, $m=5^vm'$, then $g\equiv1\pmod{5^v}$, the three coefficients are
divisible by $5^v$, the sum is $5^{2v}$ times the sum of $q/5^v$ at level $5^{k-v}$, and completing the square there gives
$G(A_1)\,G(5w')$ with $\bigl(\tfrac{A_1w'}5\bigr)=1$; hence $\operatorname{tr}=5^v\sqrt5\,\bigl(\tfrac{A_1}5\bigr)$, and
$\bigl(\tfrac{A_1}5\bigr)=\bigl(\tfrac{2F_{4m}/5^v}5\bigr)=\chiv(m')$ because $F_{5n}/5F_n\equiv1$ and $F_{4m'}\equiv3m'\pmod5$.
The Fibonacci steps, $t-2=L_n^2$ or $5F_n^2$, $F_{5n}/5F_n\equiv1$ and $F_{4m'}\equiv3m'$, are \tier[PasosDelTexto.trace\_sub\_two\_odd, PasosDelTexto.trace\_sub\_two\_even, PasosDelTexto.fib\_five\_mul, PasosDelTexto.fib\_four\_mul\_mod\_five]{P}; the completion of the square\ given the decomposition, which is checked on $115$ elements, units and congruent to $1$ modulo $5$ and $25$; every
step of the chain for $k\le7$, and the traces themselves exactly for $k\le9$.  Their modulus is
$|\ker(M^{2n}-1)|^{1/2}$, as \cite{thomas08} predicts.  The restriction count for the $301$ primes $7\le p<2000$;
the Fourier step for $k\le7$; the whole statement against the diagonalisation
(\texttt{retro81.py}--\texttt{retro87.py}).  The systematic degeneracies of the quantised map are those of % pending(D7)
\cite{kurlbergrudnick00}.\par}
\end{proof}

This is clause~(iv) of \cref{prop:what-fluctuation-takes} read as a spectrum: $e_p\,i_p=p-1$ or $2(p+1)$ becomes
$p=e\,m+\chiv(p)$, the index is a multiplicity, and the character is one state --- one more at the eigenvalue $1$
where the Frobenius is the identity, one fewer at the half-turn where it is the conjugation.  At the ramified prime
the character does not count a state but selects the half of the circle the spectrum occupies, and the selection
repeats at every power of $5$, the holes of one level being those of the level below.

\begin{remark}[The regular representation of the torus, and one state]\label{rmk:regular-plus-character}
Parts (ii)--(iii) of \cref{prop:index-degeneracy} say that, restricted to the torus $T_p$ of $M^2$ --- cyclic of order
$p-\chiv(p)$ --- the Weil representation is the regular representation of $T_p$ plus $\chiv(p)$ times its quadratic
character $\eta$:
\[
  W_p\big|_{T_p}=\operatorname{Reg}(T_p)+\chiv(p)\,\eta .
\]
On $\operatorname{Reg}(T_p)$, the functions on a finite circle of $p-\chiv(p)$ points, multiplication by a character
$\varpi$ of $T_p$ is an isometry with $\varpi\,U\,\varpi^{-1}=\varpi(\golden^2)^{\pm1}U$: it permutes the eigenspaces of $U$
cyclically, and they are equal.  If $p$ splits, in the basis of the torus $U$ is the scaling
$f(x)\mapsto(\tfrac{a}{p})f(x/a)$ of $\mathbb F_p$ by $a=\golden^2$, $\varpi$ is a Dirichlet character modulo $p$, and the
extra state is the delta at the fixed point $0$ (\texttt{retro88.py}).  At $5^k$ the same holds layer by layer: % pending(D7)
the shift of the phases by $5^{v+1}$ preserves the layer $v_5(j)=v$ and its occupation, and one state is left over.
The local models glue by the Chinese remainder theorem: for odd $N$, $W_N=\bigotimes_{p^k\,\|\,N}W_{p^k}$, the trace of
$U_N^{\,n}$ is the product of the local traces --- at $5^k$ twisted by $(\tfrac{c}{5})$, $c=(N/5^k)^{-1}$ --- and
$|\operatorname{tr}U_N^{\,n}|^2$ is the number of fixed points of $M^{2n}$ on $(\Z/N)^2$, the part of the count of
\cref{prop:clock-cyclotomic} at the primes of $N$ whose clock divides $n$ (\texttt{retro89.py}). % pending(D7)
\end{remark}

\begin{remark}[Where the open step lies]\label{rmk:open-step}
By \cref{prop:what-fluctuation-takes}(i)--(ii), each prime enters with an exact arithmetic angle, the half-turn $\chiv(p)=\pm1$
(\cref{prop:multiplicative-side}(ii)), and with its archimedean angle $t\log p$ on the torus of
\cref{prop:torus-line}.  Nothing statistical enters prime by prime.  The open step is in the assembly
of the primes along $T\mapsto(T\log p)_p$, which is what clause~(iii) of \cref{conj:closure} asks for.
The two angles are those of the propositions cited; that the open step lies in their assembly is a reading.
\end{remark}

\begin{proposition}[Euler's product, the tower and the positions]\label{prop:euler-tower}
\begin{enumerate}[label=\textup{(\roman*)}]
  \item For $\mathrm{Re}\,s>1$, $\prod_p(1-p^{-s})^{-1}=\zeta(s)$ (\cref{prop:euler-colimit,thm:odd-zeta}), and the logarithm of each local factor is $\sum_k p^{-ks}/k$.
  \item At real $s$ the product carries magnitudes only: at $s=2$ it is $\pi^2/6$, and at $s=1$ the sum $\sum_p1/p$ diverges.
  \item The frequency of a prime is its coordinate on the tower of $\golden$, the golden valuation of
    \cref{prop:tower-disc-val} and the $\sigma_p$ of \cref{prop:two-coords} and \cref{rmk:primes-interstitial}:
    $\log p=\sigma_p\log\golden$ with $\sigma_p=\log_\golden p=-\log_\golden|p|_p$ and $\golden^{\sigma_p}=p$ (the
    product formula, as \cref{prop:omega-two-adic} has it for $p=2$); $\sigma_2=\lambda_{\log}$,
    $\sigma_{2^e}=e\,\lambda_{\log}$, and $\sigma_{2^e-1}=e\,\lambda_{\log}+\log_\golden(1-2^{-e})$, where $\sigma_n=\log_\golden n$ for every $n\ge1$.
  \item The characters of the Gauss, the Eisenstein and the golden lattices, $\chi_{-4}=\chi_4$, $\chi_{-3}$ (the character of $\Q(\sqrt{-3})$, of conductor $3$) and
    $\chiv$, depend on $p$ only modulo $60$; for the Mersenne primes they are $(-1,+1,\pm1)$
    (\cref{prop:mersenne-exponent}).
\end{enumerate}
{\raggedright \tier[euler\_product\_identity, F1PaperPadic.euler\_log\_local\_factor, F1PaperPadic.euler\_product\_at\_two, F1PaperPadic.sum\_of\_prime\_reciprocals\_diverges, F1PaperPadic.frequency\_is\_tower\_coordinate, F1PaperPadic.tower\_coord\_two, F1PaperPadic.tower\_coord\_two\_pow, F1PaperPadic.tower\_coord\_mersenne, F1PaperPadic.tower\_coordinate\_is\_padic\_valuation, F1PaperPadic.three\_lattices\_level\_sixty]{P},
(i), (ii) from Mathlib\par}
\end{proposition}

Euler's identity is where the additive enumeration of the integers meets their multiplicative
factorisation --- the line of the dimensional chain (\cref{thm:one-two-three}(1)) --- and it carries
the sizes of the integers into the product; the even values it gives lie in $\mathbb{Q}\pi^{2k}$
(\cref{ssec:odd-zeta}).  Its weights $1/k$ are those of \cref{prop:what-fluctuation-takes}.  At real $s$ the product
sees only magnitudes --- and they already give $\pi$; at $s=\tfrac12+it$ it turns each factor by the
angle $t\log p$, and that turn is the fluctuation (\cref{prop:what-fluctuation-takes}).  The descent to $\Fone$ keeps
the multiplicative side --- in the terms of the Primitive Structure, the Dirichlet series factors through the binary
level and the Euler product through the ternary one \cite{gonzalez25hp} --- and with it the primes as
independent generators, their classes and their modulus $\tfrac12$.  What the additive side supplied
and the descent does not keep is their sizes, $\log p=\sigma_p\log\golden$.  The structure gives the coordinate of $2$, the binary anchor $\lambda_{\log}$ (\cref{prop:two-coords}), and that of the Mersenne numbers, $e\,\lambda_{\log}+\log_\golden(1-2^{-e})$, a closed form in the binary anchor and the exponent alone
(the second term is $-2.5\cdot10^{-4}$ at $e=13$ and $-9.7\cdot10^{-10}$ at $e=31$, \texttt{euler\_link.py} \tier{N}); of the other primes it gives none one by one; on each
level $n$ of the clock it gives their sum,
$\sum v_p\,\sigma_p=\phi(n)+\sum_{d\mid n}
\mu(n/d)\,\eta_d/\log\golden$ over the primes of
clock $n$ and the intrinsic one, with
$\eta_d=\log|N(\golden^d-1)|-d\log\golden$ bounded and
$O(\golden^{-d})$ for $d\ge3$
(\cref{prop:clock-cyclotomic}), and what the law of
\cref{conj:closure} has to supply is their order by
size.  The
level $60$ is that of the Euler product of $\mathbb{Q}(i,\sqrt{-3},\sqrt5)$, the three lattices together;
for $L(s,\chiv)$ the fluctuation takes the character at level $5$.

The places are therefore already glued, in part, and exactly where they are.  Every prime, every
zero, the operator, the monoid and the archimedean place meet at the apex $\tfrac12$
(\cref{prop:membership-proved}); every prime has its angle, one of the eight classes of $\Zst$, the
roots of $\Phi_{20}$ \cite{grisales25}, with the archimedean place, classically, as the class $-1$, complex conjugation (\cref{prop:G}); the archimedean place and the place $2$ are one reading of the operator,
whose spectrum is $|2|_2\,[\golden_2]^{\,k}$ (\cref{prop:omega-two-adic}); and the magnitude is exact on the
binary tower inside the golden one, for the powers of $2$ and for the Mersenne numbers
(\cref{sec:mersenne}, \cite{acevedo20,acevedo21}).  What
is not yet glued exactly is the magnitude of each of the
other primes: on each level of the clock their
magnitudes are glued in sum
(\cref{prop:clock-cyclotomic}), and what the law of
\cref{conj:closure} has to supply is their order by
size.

\emph{The face on Weil's side.}  The operator is already self-adjoint (\cref{rmk:space-gives}); with
an exact spectrum, $e^{i\log p\cdot H}$ is unitary and the Frobenius orbit is bounded, which is the
Castelnuovo--Severi bound (W4).  So (W4) follows from
the line, given the classical expansion of the orbit
with non-negative weights
(\tier[W4Isometria.orbits\_match\_of\_line, W4Isometria.W4\_of\_orbits\_match]{P}, the expansion), and through
\cref{cor:closure} from the conjecture; it is not a
separate statement.
  The lattice of W1 is that of the eight classes, and the orbit is formalised to the
end.

\begin{proposition}[The eight classes carry the Frobenius periodically]\label{prop:eight-classes}
\begin{enumerate}[label=\textup{(\roman*)}]
  \item The eight characters $\chi_4^a\psi_5^b$ of $(\Z/20)^\times$, with $\psi_5$ the character of order four modulo $5$ with $\psi_5(2)=i$ (the $\omega$ of \cref{rmk:mate-without-G}, $\psi_5^2=\chiv$), with values in
    $\Z[i]$, separate the eight classes: $\sum_\chi\overline{\chi(m)}\chi(n)=8$ if $m=n$ and $0$ otherwise.
  \item On the lattice of the eight classes the Frobenius has order dividing $4$: $a^4\equiv1$
    modulo $20$; its orbit is periodic and its eigenvalues are roots of unity.
  \item Let $W(f,g)$ be the bilinear Weil form, whose quadratic form $W(f,f)$ is the form $Q(f)$ of \cref{thm:core-positivity},
    and $(\tau_ug)(x)=g(x-u)$ the translation by $u$.  For the orbit $t_k=W(g,\tau_{k\log p}g)$ of a test function: if the Toeplitz form is
    semidefinite then $|t_k|\le t_0$; a zero off the line would contribute a term growing like
    $(p^{|\beta-1/2|})^k$; a bounded sum of such terms with positive weights forces every factor
    to be $1$, and averages of unimodular cross terms tend to $0$ like $1/K$.
\end{enumerate}
{\raggedright \tier[F1PaperTorus.eight\_classes\_orthogonality, F1PaperTorus.frobenius\_on\_classes\_has\_order\_dividing\_four, F1PaperTorus.frobenius\_orbit\_is\_periodic, F1PaperTorus.toeplitz\_bound, F1PaperTorus.offline\_orbit\_unbounded, F1PaperTorus.bounded\_both\_ways\_forces\_one, F1PaperTorus.weil\_route\_assembly, F1PaperTorus.cesaro\_cross\_term]{P}\par}
\end{proposition}

The orbit decomposes into the eight character blocks, and each block is exactly a Toeplitz
sequence: the Weil form of two translates depends only on the difference of the translations
(\tier[F1PaperTorus.crossCorrPair\_eq, F1PaperTorus.orbit\_is\_toeplitz]{P}).  The evidence for
(W4) has four layers.

\emph{The explicit formula on the orbit, checked.}  For $\chiv$, width $a=0.3$ and $p=3$, the
orbit computed from the primes and computed from the $397$ zeros below height $500$ agree to
$10^{-8}$ for $k=0,\dots,12$; the zeros above $500$ can contribute at most $9.6\cdot10^{-7}$.  The
orbit is computed with the autocorrelation of the triangles in closed form and the
archimedean side in the variable $u$, so no interpolation enters (\texttt{zeroside.py} \tier{N}).

\emph{The bound, swept.}  For the eight blocks, the widths $a\in\{0.1,0.2,0.3,0.5,0.8,1.2\}$, the
primes $p\in\{2,3,7,11,13\}$ and every $k$ with $k\log p+2a\le14.2$ --- $2208$ orbit values --- the
bound holds everywhere, with $\max|t_k|/t_0$ between $0.929$ ($\zeta$) and $0.986$ ($\psi_5,\psi_5^3$), and
every Toeplitz matrix is positive definite, with smallest eigenvalue at least $1.27\cdot10^{-4}\,t_0$
(\code{sweep.py} \tier[\testref{test\_circle\_and\_detect\_golden\_values, test\_circle\_and\_detect\_matches\_manuscript\_claim}]{N}).  On the circle of $p=3$ the Fej\'er density is positive in all eight
blocks, $\zeta$ included, with minima from $1.2\cdot10^{-4}$ ($\zeta$) to $0.032$ ($\chi_{20}$)
(\texttt{circle\_and\_detect.py} \tier[\testref{test\_circle\_and\_detect\_golden\_values, test\_circle\_and\_detect\_matches\_manuscript\_claim}]{N}).

\emph{The certified orbit, and the families.}  In interval arithmetic with outward rounding,
for $\chiv$ with width $a=0.01$ and $p=2$, $|t_k|\le0.1122\,t_0$ for $k\le25$ and the Toeplitz matrix is
positive definite with smallest eigenvalue at least $0.248\,t_0$; for the representation of the
quintic layer, of conductor $5^7$, with $a=0.1$, the same holds with $0.158\,t_0$ and $0.506\,t_0$
(\texttt{orbita\_cert.py}).  The orbit of $\chiv$ computed from the primes alone is fitted by
a single oscillation of frequency $6.649$, the height $6.6485$ of the first zero of $L(s,\chiv)$
(\texttt{W3\_orbita\_lado\_primos.py} \tier{N}).  The three power-residue symbols of \cref{prop:phi-layers} have $L$-functions with % pending(D7)
functional equations of conductor $80$, $135$ and $5^7$; their zeros up to $t=60$, $50$ and $6.5$
($66$, $56$ and $8$ of them) lie on the line, and the explicit formula with those zeros
reproduces the sums of the symbols (\texttt{v09}--\texttt{v11} \tier{N}).

\emph{What the numbers can and cannot see.}  A zero of $L(s,\chiv)$ at height $\gamma$ with
$\Re\rho=\tfrac12+\delta$ would have made some Toeplitz matrix of the sweep indefinite as soon as
$\delta\ge0.10$ at $\gamma=30$ and $\delta\ge0.38$ at $\gamma=100$; above $\gamma\approx200$ no admissible
$\delta$ is detected, because $|\hat g_a(\gamma)|^2$ decays like $\gamma^{-4}$.  Longer orbits see
further: with $k\log p\le120$ in place of $14.2$ the thresholds fall to $0.001$, $0.009$, $0.022$
and $0.028$ at $\gamma=30$, $100$, $200$ and $300$ (\texttt{circle\_and\_detect.py} \tier[\testref{test\_circle\_and\_detect\_golden\_values, test\_circle\_and\_detect\_matches\_manuscript\_claim}]{N}).  The sweep is therefore a verification of the line at low height, of the same kind
as the classical ones, and not evidence for (W4) at every height.  That is intrinsic: (W4) for
every $k$ and every width is as strong as the line, and no finite computation reaches it.  It has
to come from structure, as Castelnuovo--Severi comes from the Hodge index in Weil's proof.

\begin{proposition}[(W4) is (W3) on the planes of the orbit]\label{prop:W4-from-W3}
Let $W$ be the bilinear Weil form of $L(s,\chiv)$ of \cref{prop:eight-classes}\textup{(iii)}, with $W(f)=W(f,f)$, $g$ a real test function and $t(u)=W(g,\tau_ug)$, so that
$t_k(a)=t(k\log p)$ for $g=g_a$.
\begin{enumerate}[label=\textup{(\roman*)}]
  \item $t(0)-t(u)=\tfrac12W(\tau_ug-g)$ and $t(0)+t(u)=\tfrac12W(\tau_ug+g)$.  Hence $|t(u)|\le t(0)$ if and only
    if $W\ge0$ on the plane spanned by $g$ and $\tau_ug$: \textup{(W4)} for the orbit of $p$ is \textup{(W3)} on
    the planes $\langle g_a,\tau_{k\log p}g_a\rangle$, and \textup{(W3)} on every window implies \textup{(W4)}, by
    the Cauchy--Schwarz inequality.
  \item For $g$ even, $\tau_ug\mp g$ are odd and even about the midpoint $u/2$: the two halves
    $t_k\le t_0$ and $t_k\ge-t_0$ of \textup{(W4)} are the two parities of \cref{prop:free-orbit-index}.
  \item The window theorems give \textup{(W4)} on every plane that fits in them, $2a+k\log p\le L$: for $p=2$ and
    $k=1$, with $a\le0.025$ by \cref{thm:first-window} and with $a\le0.15$ and $a\le0.30$ by
    \cref{thm:second-window,thm:third-window}; for $p=3$ and $k=1$, with $a\le0.10$ by \cref{thm:third-window}.  The
    certified orbit and the sweep of this subsection are \textup{(W3)} on the span of the orbit, of which \textup{(W4)}
    is the case of two translates.
  \item With a single free orbit at height $\gamma_0$, if $W(f)<0$ then $W(\tau_{\pi/\gamma_0}f-f)<0$: the opposite parity is
    negative on the window widened by $\pi/\gamma_0$, and the index is exactly $2$ on every window longer by $\pi/\gamma_0$ than
    one on which the orbit is detected.
\end{enumerate}
{\raggedright\tier[W4DeW3.plane\_values, W4DeW3.W4\_of\_W3\_on\_plane, W4DeW3.W3\_on\_plane\_of\_W4, ParidadesJuntas.other\_parity\_value, MedioPeriodo.line\_factor, MedioPeriodo.orbit\_factor, MedioPeriodo.half\_period\_closes]{P}; (iii) with the witness marks of the window theorems, \tier{P}\ with declared inputs;
(iv) with the explicit formula on the orbit, as in \cref{thm:weil-closure}\par}
\end{proposition}
\begin{proof}
(i) $W(\tau_ug\mp g)=W(\tau_ug)+W(g)\mp2W(g,\tau_ug)=2t(0)\mp2t(u)$, by translation invariance; the two values are
$\ge0$ exactly when $|t(u)|\le t(0)$.  (ii) The reflection about $u/2$ exchanges $g$ and $\tau_ug$ when $g$ is even.
(iii) The plane is supported in an interval of length $2a+u$.  (iv) By the explicit formula,
$W(\tau_uf-f)=\sum_\rho\hat f(\gamma)\hat f(-\gamma)(2-2\cos\gamma u)$.  With $W(f)=W_{\mathrm{line}}+W_{\mathrm{orb}}$, $W_{\mathrm{line}}\ge0$ the part of the
zeros on the line and $W_{\mathrm{orb}}$ that of the orbit: the factor lies in $[0,4]$ on the line, and at $u=\pi/\gamma_0$ the four
members $\gamma=\pm\gamma_0\pm i\delta$ have $\cos\gamma u=-\cosh(\pi\delta/\gamma_0)$, so the total is less than
$4W_{\mathrm{line}}-(2+2\cosh c)W_{\mathrm{line}}\le0$, with $c=\pi\delta/\gamma_0$.
\end{proof}

So the route has one open step, not two: \textup{(W4)} is what \textup{(W3)} says on the planes of the orbit, and
what separates what is proved and certified from \textup{(W5)} is that the orbit needs every $k$: the window theorems reach $k=1$, the certificates $k\le25$ for $p=2$.  On
the Davenport--Heilbronn zeros below height $130$, with the triangles $g_a$ and $p=7$, $|t_k|/t_0$ grows like
$7^{0.326k}$, against the real part $0.8085$ of its first zero off the line, the growth that
\cref{thm:weil-closure} assumes; without the two free orbits, $|t_k|\le t_0$ for every $k\le60$
(\texttt{retro102.py}). % pending(D7)

\begin{theorem}[The closing theorem]\label{thm:weil-closure}
Fix one prime $p$.  For the triangles $g_a(x)=(1-|x|/a)^+$ of every width $a>0$, let
$t_k(a)=W_{\chiv}(g_a,\tau_{k\log p}g_a)$, the archimedean side minus the prime side of the Weil
form of $L(s,\chiv)$, both written explicitly.  Assume \textup{(W4)} for the block of $\chiv$: for
every $a>0$ and every $k$, $|t_k(a)|\le t_0(a)$.  Assume the explicit formula on the orbit with the separation of phases:
a zero with $\Re\rho\neq\tfrac12$ and a test function not vanishing at its ordinate give
$c\,(p^{|\Re\rho-1/2|})^k\le|t_k|+C$ for infinitely many $k$.  Then every zero of $L(s,\chiv)$ in
the strip has real part $\tfrac12$.  \tier{P}%
\end{theorem}
{\raggedright\LeanInline{F1PaperTorus.the\_line\_from\_W4\_explicit}, with the orbit defined in \LeanInline{F1PaperTorus.chi5WeilOrbit} \tier{P};
(W4), \LeanInline{F1PaperTorus.W4Chi5}, open; the explicit formula on the orbit,
\LeanInline{F1PaperTorus.OrbitDetectsOffLine}, whose combinatorial core --- a dominant term growing like
$R^k$ in infinitely many $k$ plus a tail that does not reach it makes the orbit unbounded --- is
\tier[F1PaperTorus.dominant\_plus\_tail\_unbounded]{P}\par}

\begin{proof}
Let $\rho$ be a zero in the strip with $\Re\rho\neq\tfrac12$, and choose the width $a$ so that $\hat g_a$ does not vanish
at its ordinate.  The second hypothesis, which is the explicit formula of \cref{thm:explicit} read on the orbit, gives
constants $c>0$ and $C$ with $c\,R^k\le|t_k(a)|+C$ for infinitely many $k$, where $R=p^{|\Re\rho-1/2|}>1$.  The first
hypothesis gives $|t_k(a)|\le t_0(a)$ for every $k$.  So $c\,R^k\le t_0(a)+C$ for infinitely many $k$, which is false
because $R>1$.  The first hypothesis is \textup{(W3)} on the planes of the orbit (\cref{prop:W4-from-W3}).
\end{proof}

Two features of the statement matter.  The orbit is defined, not a parameter: its two sides
are the explicit integrals and the explicit sum, so the classical input is literally the
explicit formula for that sequence and cannot carry the conclusion by itself.  And one prime
suffices: the orbit of any single $p$ detects every zero off the line, provided the width $a$
is chosen so that $\hat g_a$ does not vanish at its ordinate, which the family of all widths
allows.

The theorem is not circular: it assumes (W4) and derives the line, and (W4) is what the exact
spectrum of the operator gives.

\subsubsection{Scope: from $L(s,\chiv)$ to $\zeta(s)$}\label{ssec:scope-zeta}
The equivalent forms of \cref{thm:topological-proof} are statements about $L(s,\chiv)$.  The
factorisation $\zeta_K=\zeta\cdot L(\cdot,\chiv)$ carries some of it to $\zeta(s)$ and
not the rest.

What transfers is the arithmetic.  The Euler factors of $\zeta$ and of $\zeta_K$ are read
off the same Frobenius (\cref{thm:frob-eigen}), the factorisation holds factor by
factor and globally (\cref{eq:dedekind}), the odd values of $\zeta$ are determined
through it (\cref{thm:odd-zeta}), every value of $\zeta$ is written in base $\golden$
(\cref{thm:base-phi}), and the counting function of $\zeta_K$ decomposes accordingly
(\cref{thm:pcf-counting}).  The zero set of $\zeta_K$ is the union of those of $\zeta$
and of $L(s,\chiv)$, with multiplicities adding.

What does not transfer are the equivalent forms.  Their structure --- the Klein symmetry induced
by $\chiv$, the orbit centre, the trace spectrum of the operator and the closing links
--- is built for the zeros of $L(s,\chiv)$, and the conductor that the zeros select is
$5$ (\cref{prop:envelope-exact}), not $1$.  The equivalent forms do not place the zeros of
$\zeta$ on the line: GRH$(\chiv)$ does not imply RH.  RH for $\zeta_K$ implies both, and conversely RH and
GRH$(\chiv)$ give RH for $\zeta_K$ through the factorisation, since in the strip the zeros of
$\zeta\cdot L(\cdot,\chiv)$ are those of its two factors (\tier[AlcanceZeta.dedekind\_strip\_zeros\_on\_line]{P}, with the two statements as hypotheses).  The two statements share the construction and the anchor, not
the closing step.  The closure conjecture is stated for $\zeta$ as well (\cref{cor:closure}): if it
holds, the line holds for $\zeta$ too; the equivalent forms do not carry it.


% =====================================================================

\section{Discussion}\label{sec:disc}

% ---------------------------------------------------------------------
\subsection{What has been demonstrated}\label{sec:discussion}

The paper proves one construction and one set of equivalent forms, and several results that stand
on their own.  The construction starts from $\golden^2=\golden+1$ and reaches the
zeros of $L(s,\chiv)$ in eight stages (\cref{sec:topological-proof}): the pentagon and
the turn derive the conductor and the classes of the primes; the ring $\Zphi$
carries two registers, a $\Lambda$-monoid between levels and an Artin
$\Lambda$-structure within each level away from $5$, whose Frobenius reads off
$\chiv$; the Euler factors are the characteristic polynomials of that Frobenius, the
Euler product their colimit and the Dedekind factorisation a natural isomorphism,
with the odd zeta values determined through it (\cref{thm:odd-zeta}) and every value
written in base $\golden$ from the anchor (\cref{thm:base-phi}); the Mersenne anchor
turns the tower into a function; the operator $\Oone=HY$ (\cref{thm:one-operator}),
of which every form in which it has been written is a function (\cref{thm:operator-forms}), has an exact spectrum, $m_0\golden^{\sigma}\zeta_3^{\,k}(-1)^{\epsilon}$, fixed before any zero, with parameters that are cardinals of $G$ (\cref{prop:pair-from-G}); and the
zeros meet it through the orbit centre, the line as a fixed set, and the Klein
symmetry induced by $\chiv$.  The equivalent forms are \cref{thm:topological-proof}: every zero lies on the line as soon as the coordinate $|\rho-1|/|\rho|$ of each zero $\rho$ is a power of $\golden$, or any of the forms of
\cref{ssec:core-forms} holds; the
integrality of the trace with the strip bound already forces $r=1$ (\cref{thm:arrow}), and the
route of closure, Weil's, is carried to one conjecture: a law of the primes by classes, derived from the structure, makes exact the positions of the zeros that the structure already approximates (\cref{conj:closure}), whose face on Weil's side is the Castelnuovo--Severi bound on
the Frobenius orbits class by class (\cref{thm:weil-closure}).  What stands on its own is in \cref{sec:results}: the zeros return the
conductor, the explicit formula is written with the Frobenius traces of $\Zphi$, and
Weil's positivity holds in three windows: the first proved, the other two certified by
computer.  The subsections below place this
against the programmes it belongs to.

\subsubsection{What the paper achieves, and how it leads to the closure}\label{ssec:achieved}

The paper is a proof of concept for the programme it belongs to: Weil's proof transported along
one explicit chain, arithmetic, geometry, analysis, operator, carried as far as it goes
(\cref{ssec:route-overview,ssec:weil-closure}), with every step either proved, classical, or isolated
as the one statement that remains.  What it
achieves is the following.

\begin{enumerate}[label=\textup{(\arabic*)}]
  \item \emph{The descent is complete on one ring.}  On $\Zphi$ the Frobenius lifts exist
    for every prime but $5$, are forced by $\chiv$, and fail at $5$ with an explicit
    obstruction; the character they determine splits $\zeta_{\Q(\sqrt5)}$
    (\cref{thm:frobenius-congruence,thm:obstruction-five,thm:chi5-from-frobenius}).
  \item \emph{The inputs of the Primitive Structure are discharged.}  Its eight named
    axioms are theorems, and its bridge is proved link by link up to the last one
    (\cref{sec:disc-hp}).
  \item \emph{Weil's route is transported step by step.}  The Frobenius acts on finite
    lattices whose Euler factors are its characteristic polynomials, among them the eight
    classes of $(\Z/20)^\times$ with a Frobenius of order dividing $4$; the Weil form is
    written explicitly, with the Frobenius as translation and the orbit matrix exactly
    Toeplitz; the pole is the ample class; (W4) is (W3) on the planes of the orbit, as Castelnuovo--Severi is the
    Hodge index in Weil's proof; and the closing theorem is formalised
    (\cref{ssec:route-overview,ssec:weil-closure,prop:W4-from-W3}).
  \item \emph{The open statement is one, and it is named.}  The route ends in one conjecture, the closure conjecture (\cref{conj:closure}), which would give the last link of the bridge.
    With it the line follows (\cref{cor:closure}).  Its law is obtained in two steps: the first, by
    the levels of the clock, is proved (\cref{prop:clock-cyclotomic}); the second, from the clock
    to the size, is one statement on the tower, the bounded energy of the deviation from the
    uniform vertex in each layer (\cref{prop:index-ngon,prop:energy-scale,rmk:closure-status}).
  \item \emph{What is missing is known exactly.}  The positions are approximated, the primes belong and
    the zeros enter the structure (\cref{prop:positions,prop:membership-proved}); what remains is precision, of the real part --- that each zero is a member, $r=1$, the open
    statement --- and of the height, which makes the operator exact; the exact heights tie the second to the first, since the height fixes the phase and the
    isometry the modulus (\cref{prop:phase-modulus}).  The fluctuation is a function on the torus of the circles of the primes, read along the line (\cref{prop:torus-line}), and its law is the same along the line, at random points of the torus and on the zeros \tier{N}\ (\cref{ssec:fluct-geometry});
    the classes are the part of the law already in hand, and what the law has to predict is
    fixed by the local factors (\cref{prop:what-fluctuation-takes}); the resolution it has to reach is measured, of
    order $1/T$ at height $T$.
  \item \emph{The verdict is measured where it can be.}  The explicit formula on the orbit
    agrees with the zeros to $10^{-8}$; the orbit bound holds in all eight classes at
    $2208$ values; Weil positivity is proved in a first window and certified in two more.
  \item \emph{What the structure fixes and what it does not is stated.}  It fixes the symmetry of the zeros, from the
    reality and the parity of $\chiv$, the value $\tfrac12$ and the constants of the positivity; the position on the
    line is fixed by the mate, which does not use the reality of the character (\cref{rmk:mate-without-G}), and a
    statement that holds equally for the Davenport--Heilbronn sequence is a form of the open statement and not a step
    towards it (\cref{ssec:two-sides}).
  \item \emph{Why the operator has been elusive is stated inside the tower.}  An operator of the level, with its
    whole functional calculus, does not separate the torsion fibre, and what separates it does not commute with its
    symmetry (\cref{prop:level-operators}); a commuting symmetry would show in a spectrum as sectors and an
    antiunitary one of square $-1$ as doublets (\cref{prop:symmetry-spectrum}); the zeros of $L(s,\chiv)$ show
    neither, and those of $\zeta_K$ show the sector of its Galois symmetry (\cref{ssec:envelope-fluctuation}).
  \item \emph{The archimedean and the $p$-adic faces are one dynamics.}  The Fibonacci matrix is mixing on the real
    torus, where the characters see it only as the decay of correlations, and exactly periodic on the torsion
    points, prime by prime, with the clock as period; the number of its periodic points is the product of the
    cyclotomic norms $|N(\Phi_d(\golden))|$, which is the product formula that glues the two readings
    (\cref{prop:two-faces}).
\end{enumerate}

How this leads to the closure.  The conjecture asks for no computation to absolute precision:
it asks for a law, as Weil's Castelnuovo--Severi bound comes from the Hodge index and not from
counting points.  The global part of that law is already proved --- the classes of the primes
carry the Frobenius, which on them is a rotation of finite order (\cref{prop:G,prop:eight-classes})
--- and what remains is the resolution: to predict, from the structure, the residual of each class
at the scale the height requires (\cref{conj:closure}); the line follows from it by \cref{cor:closure}, and its exact heights land every zero on the
spectrum of the operator (\cref{prop:phase-modulus}).  That is the
closure; it is not proved here.

\emph{Why the correction has to be a global law.}  The correction can be computed prime by prime:
summing the explicit formula over the primes up to $X$ takes the operator from $0.20$ spacings to
$0.020$ (\cref{ssec:weil-closure}).  That is significant, but it decides nothing.  It is a
computation of the kind that located the zeros one by one in Odlyzko's work \cite{odlyzko87}: it
confirms, and it uses the primes as data, so as a construction it is circular; and it does not
converge --- with the primes summed up to a fixed $X$, the error has a minimum at a finite $X$ and more primes past it make the
error larger, not smaller (\cref{ssec:weil-closure}) \tier{N}.  What closes has
to be a law, and the paper already has the global part of it: the classification of the primes
by their classes modulo $20$ of Grisales Herrera \cite{grisales25}, which this paper proves carries
the Frobenius, the Euler factors and the belonging of every prime
(\cref{prop:G,prop:membership-proved}), and without any one of whose classes the
correction is worse.  Absolute precision depends on understanding how the fluctuation changes with the height --- what it
needs is the sizes $\log p$ of the primes, their coordinates on the tower (\cref{prop:euler-tower}): its law is the geometry of the torus (\cref{prop:torus-line}), and the resolution
it has to reach grows like $1/T$.

\emph{What the operator gives without that precision.}  The operator lacks the absolute
precision the Riemann hypothesis asks for; holography does not ask for it.  What holography
needs from the operator it already has.  It is an individual where the ensemble stood: the
modulus $\tfrac12$ that \cite[\S2.13]{gonzalez26cw} produces, with the parameters of the operator read
from the cocone and its envelope placing the zeros to a fraction of a spacing
(\cref{prop:individual-ensemble,prop:positions}).  It carries the origin of the odd zeta values:
determined by $\golden$ and $\pi$ (\cref{thm:odd-zeta}), they are the coefficients of the
$\alpha'$-expansion of the Veneziano amplitude, whose Regge tower is the Euler product
(\cref{sec:primon}).  And it sits where \cite[\S1.7.5]{gonzalez26cw} places black holes,
between $\golden$ and the primes: the horizon geometry of the extremal Kerr--Newman family selects
$\golden$, the interior near the singularity selects the primes through the primon gas, and the
modulus $\tfrac12$ is one bit per fundamental cell, which \cite[\S1.7.5]{gonzalez26cw} reads as the Bekenstein bound saturated.  The operator
is built from $\golden$ and reaches the primes through their classes; that is the intermediate
point.

\subsubsection{Relation to the $\Fone$ programme}

The programme looks for a base below $\mathrm{Spec}\,\Z$ so that the product
$\mathrm{Spec}\,\Z\times_{\Fone}\mathrm{Spec}\,\Z$ is a surface
(\cref{ssec:f1-obstruction}).  This paper does not construct that base.  What it
does is take one of the strands the programme offers --- Borger's
$\Lambda$-structures --- and carry it out completely on one small ring, $\Zphi$,
where every step is checkable: the structure exists away from $5$
(\cref{thm:frobenius-congruence}), it does not exist at $5$
(\cref{thm:obstruction-five}), and what it determines is the character that
splits $\zeta_{\Q(\sqrt5)}$ (\cref{thm:chi5-from-frobenius}).  The choice of
$\Zphi$ is not arbitrary within the programme: the cyclotomic tower is the strand
Iwasawa theory already uses in place of $\Z\otimes_{\Fone}\bar{\mathbb{F}}_1$
(\cref{ssec:f1-strategies}), and $\golden$ lives in the pentagon.

\subsubsection{Relation to Borger}\label{ssec:rel-borger}

Borger's $\Lambda$-rings are descent data to $\Fone$, and the datum is additive:
$\Frob_p$ must be a ring endomorphism with $\Frob_p(a)\equiv a^p\pmod p$.  The
object that carries the descent in this programme is a tower of rings, not a
sequence of numbers: the levels are the lattices $\LPCF^{(\sigma)}=\golden^\sigma\LPCF$ of
\cite[\S2.8, \S3.4]{gonzalez26cw}, each one a module over the base ring with its
own additive structure, and the scaling by $\golden$ is a geometric map between
levels.  There are therefore two directions, and they carry different halves of
the datum.

\emph{Along the tower}, the lift $\Frob_p$ multiplies the level index,
$\golden^n\mapsto\golden^{pn}$, and composes as the primes multiply
(\cref{thm:lambda-monoid}): this is Borger's axiom in its multiplicative form,
and it is what makes the tower carry the Euler product.  \emph{Within a level},
the ring structure is the one that has to supply additivity, and it is here that
\cite{gonzalez25hp} took the map $\golden\mapsto\golden^p$ for the ring endomorphism: it does not respect
$\golden^2=\golden+1$ (\cref{rmk:monoid-not-ring}), so the tower carries the multiplicative half of the
datum and the additive half has to be produced inside the level.

\Cref{sec:lambda-ring} produces it: on $\Zphi$ the only ring endomorphisms are
the identity and the conjugation (\cref{prop:endo}), the Frobenius congruence
forces which one prime by prime according to $\chiv$
(\cref{thm:frobenius-congruence}), and the two directions agree modulo $p$
(\cref{thm:monoid-ring-bridge}).  The obstruction at $5$ is then not an accident
of the construction but the ramification: it is exhibited by a nilpotent
(\cref{thm:obstruction-five}), and it is the reason a $\Lambda$-structure on the
full ring of integers of a ramified extension is not available.

\subsubsection{Relation to Manin}\label{ssec:rel-manin}

Manin's two programmes supply the requirement and the object
(\cref{ssec:why-phi}).  The difference is the route.  Buium's $p$-derivations,
$\delta_p(a)=(\Frob_p(a)-a^p)/p$, live in $p$-adic completions and their modern form is perfectoid; the present route stays cyclotomic and finite, which is why it could be written in Lean with Mathlib as it stands (\cref{ssec:lean-vs-padic}).  The $p$-derivations it does use are global on $\Zphi$ (\cref{prop:delta-ring}), and its polar form of $\golden$ is computed at finite level, the completion entering only as a classical input (\cref{prop:polar}).  On the second programme, the noncommutative tori
$T_\theta$ with $\theta=\golden$ are the arena \cite{manin04}, and CW supplies a
worldsheet over one of them (\cref{sec:worldsheet}); what neither supplies is the
dynamics that would put the zeros in a spectrum.

\subsubsection{Relation to Weil}\label{ssec:rel-weil}

Weil's proof needs a space carrying Frobenius, the zeros as its eigenvalues, and
a positivity (\cref{ssec:weil-template}).  This paper has the first ingredient in
a precise finite form --- $\Frob_p$ acts on $\Zphi$ and the Euler factor is its
characteristic polynomial (\cref{thm:frob-eigen}) --- and the classical bridge to
the zeros, the explicit formula, in which the primes enter as Frobenius traces
(\cref{thm:explicit}).  The third ingredient is the open one, and it is open in
the same form as in Weil's argument: a positivity.  The window theorems prove it in the first window (\cref{thm:first-window}) and certify it numerically in the other two (\cref{thm:second-window,thm:third-window}), and \cref{prop:semilocal} shows that past the archimedean
threshold the sign of $\Frob_2$ is what carries it.  The analogy stops where
\cref{ssec:weil-analogy} says: $|\ker G|=2$ is not a genus-zero statement in the
sense that would make $H^1$ vanish.  \Cref{ssec:weil-closure} follows the route to its end:
with the eight classes of $(\Z/20)^\times$ as the finite lattice (Frobenius of order dividing $4$),
the Weil form with the Frobenius as translation, and the pole as ample class, the whole argument
is formalised except (W4), the bound on the Frobenius orbit, which is (W3) on the planes of the orbit, as Castelnuovo--Severi follows from the Hodge index in Weil's proof (\cref{prop:W4-from-W3}), and which holds in all eight classes
wherever it has been computed and which the closure conjecture gives (\cref{ssec:weil-closure}); the
explicit formula on the orbit is a classical input.

\subsubsection{Position within the $\mathbb{F}_1$ landscape}
\label{ssec:position}

We assess the construction against the desiderata
\ref{D1}--\ref{D9} stated in \cref{ssec:desiderata}.

\begin{description}[leftmargin=*,nosep]
\item[\ref{D1} ($\mathrm{Spec}\,\mathbb{F}_1$ a single point).]
  Not addressed: the paper does not construct $\mathrm{Spec}\,\mathbb{F}_1$.
\item[\ref{D2} ($\mathrm{GL}_n(\mathbb{F}_1) = \Sigma_n$).]
  Met in one instance, not as a theorem about $\mathbb F_1$: the symmetry group of the
  three eigenvalues of $\OmH$ \textup{(\cref{thm:S3-from-euler})} is
  $S_3=\Sigma_3=W(A_2)$, the group Tits' limit assigns to $\mathrm{GL}_3(\mathbb{F}_1)$.
\item[\ref{D3} ($\mathbb{P}^n(\mathbb{F}_1)$ has $n+1$ points).]
  Not addressed.
\item[\ref{D4} (counting function).]  Verified
  distributionally: $N_K(u)$ is computed explicitly in
  \cref{thm:pcf-counting} as a discrete distribution on
  prime-power norms.  The polynomial form fails
  \textup{(\cref{prop:not-soule})}, but the Connes--Consani
  extension applies.
\item[\ref{D5} ($\mathrm{Spec}\,\mathbb{Z}$ as $\mathbb{F}_1$-curve).]
  Partially.  The descent datum is constructed on $\Zphi$ away from $5$ and on the
  unit monoid $\langle\golden\rangle$, not on $\mathrm{Spec}\,\mathbb{Z}$; the connection
  to $\zeta$ passes through the Dedekind factorisation
  $\zeta = \zeta_K/L(\cdot,\chi_5)$.
\item[\ref{D6} (base-change functor).]  Not addressed: no functor
  $-\otimes_{\mathbb F_1}\mathbb Z$ is constructed.
\item[\ref{D7} (toric varieties).]
  The toric piece corresponds to $\zeta_\mathbb{Q}$ and is handled
  in the literature (Deitmar's monoid schemes); it is not treated here.  The non-toric content of pentagonal
  $\mathbb{F}_1$ is articulated in
  \cref{ssec:non-toric}.
\item[\ref{D8} (Riemann zeta from $\mathbb{F}_1$-counting).]
  Partially: the construction provides
  $\zeta_{\mathbb{Q}(\sqrt{5})}$ structurally, and via Dedekind
  factorisation $\zeta(s) = \zeta_K(s)/L(s,\chi_5)$ recovers
  $\zeta(s)$.  The full Connes--Consani realisation of
  $\zeta^*_\mathbb{Q}$ via $\overline{\mathrm{Spec}\,\mathbb{Z}}$
  is not engaged here; pentagonal descent contributes the
  $\mathbb{Q}(\sqrt{5})$-place.
\item[\ref{D9} (Weil-style proof of RH).]
  The bifurcated structure required by Weil's proof for function
  fields is realised concretely for the Dedekind zeta of
  $\mathbb{Q}(\sqrt{5})$: the arithmetic--spectral span
  \textup{(\cref{prop:span})} with common source~$G$ replaces
  the inaccessible
  $\mathrm{Spec}(\mathbb{Z})\times_{\mathbb{F}_1}\mathrm{Spec}(\mathbb{Z})$
  with the homomorphism $G\colon\Zst\to\mathbb{Z}_2\times\mathbb{Z}_2$, the
  Frobenius of $\Zphi$ (\cref{sec:lambda-ring}) plays the role of Weil's, the spectral uniqueness of $(\sigma,\mu)=(3/2,1/2)$
  paired with the Hecke functional equation supplies the
  two-sided constraint
  \textup{(\cref{sec:topological-proof})}, and the Dedekind factorisation
  $\zeta = \zeta_K/L(\cdot,\chi_5)$ relates this structural data
  on $\zeta_K$ to data on $\zeta(s)$ at primes split or inert
  in $\mathbb{Z}[\golden]$.  The realisation is not a proof: it writes GRH$(\chiv)$ as one open statement in equivalent forms, and carries the route of closure to one conjecture
  (\cref{conj:closure}),
  and it does not reach RH for $\zeta$ (\cref{ssec:scope-zeta}).
\end{description}

In sum: \ref{D4} is met distributionally; \ref{D2} is met in one instance;
\ref{D5}, \ref{D8} are addressed partially; \ref{D9} is realised as a bifurcated
structure with the constituents named above, which writes GRH$(\chiv)$ as one closing step and does not prove it; \ref{D1}, \ref{D3}, \ref{D6}, \ref{D7} are not addressed.  The
detailed comparison with other $\mathbb{F}_1$-approaches is
collected in \cref{app:comparison}.

\subsubsection{The non-toric character of pentagonal $\mathbb{F}_1$}
\label{ssec:non-toric}

The natural toric examples in $\mathbb{F}_1$-geometry have clean
counting polynomials and admit Deitmar monoid-scheme structures
directly.  We make precise what
``non-toric'' means in the present setting by examining the four
quadratic extensions of $\mathbb{Q}$ classified by~$G$
(\cref{thm:dedekind-natural}).

\begin{center}
\small
\begin{tabular}{@{}llll@{}}
\toprule
$K$ & Character & Toric? & $\mathbb{F}_1$-realisation \\
\midrule
$\mathbb{Q}$ & $1$ (trivial) & yes
  & $\mathrm{Spec}\,\mathbb{F}_1[T]$ (Deitmar) \\
$\mathbb{Q}(\sqrt{-1})$ & $\chi_4$ & cyclotomic
  & $\mathbb{F}_{1^4}$-extension (\cref{rmk:F1n}) \\
$\mathbb{Q}(\sqrt{5})$ & $\chi_5$ & \textbf{non-toric}
  & Artin $\Lambda$-structure on $\Zphi[\tfrac15]$, $\Lambda$-monoid on $\langle\golden\rangle$ \\
$\mathbb{Q}(\sqrt{-5})$ & $\chi_4\chi_5$ & non-toric mixed
  & combination of the above \\
\bottomrule
\end{tabular}
\end{center}

The trivial case $K=\mathbb{Q}$ admits the toric realisation
$\mathrm{Spec}\,\mathbb{F}_1[T]\otimes_{\mathbb{F}_1}\mathbb{Z}
=\mathrm{Spec}\,\mathbb{Z}[T]$.
The cyclotomic case $K=\mathbb{Q}(i)$ admits the
$\mathbb{F}_{1^4}$-extension
$\mathbb{Z}[T]/(T^4-1)$, whose quotient by $T^2+1$ is $\mathbb{Z}[i]$ (after inverting
$2$ it splits as $\mathbb{Z}[\tfrac12,i]\times\mathbb{Z}[\tfrac12][T]/(T^2-1)$), and which
fits within Manin's cyclotomic framework~\cite{manin08}.
The pentagonal case $K = \mathbb{Q}(\sqrt{5})$ does not fit
either: the splitting of primes through~$\chi_5$ has no toric
model, because~$\chi_5$ is not given by a polynomial relation
in~$p$, and the unit~$\golden$ is real-quadratic rather than a
root of unity.  The construction therefore uses Borger's $\Lambda$-structures,
which are defined without reference to a torus.

The dichotomy has a clean spectral expression on the side of~$G$ (\cref{prop:G,prop:two-registers-G}).
The four characters $\{1,\chi_4,\chi_5,\chi_4\chi_5\}$ of
$\mathbb{Z}_2\times\mathbb{Z}_2$ split into a toric/cyclotomic
pair $\{1,\chi_4\}$ and a non-toric pair $\{\chi_5,\chi_4\chi_5\}$,
with the second coordinate of $G$ \textup{(}the projection
$\pi_2\colon\mathbb{Z}_2\times\mathbb{Z}_2\to\mathbb{Z}_2$ that
extracts $\chi_5\textup{)}$ identifying the non-toric
content.  Pentagonal descent is exactly the content not visible
to monoid-scheme machinery.

\subsubsection{Formalization: Lean versus the $p$-adic route}\label{ssec:lean-vs-padic}

Everything in \crefrange{sec:choice}{sec:clock},
except the Teichm\"uller lift of \cref{ssec:delta-polar}
(classical in the limit), the product formula of
\cref{prop:clock-cyclotomic} (classical), and the
elementary real estimates on the clock and the index
(proved in Lean), is either a computation in a finite ring
(\texttt{ZMod 20}, \texttt{ZMod p}) or elementary algebra in
$\Zphi=\Z[X]/(X^2-X-1)$, which is why Mathlib suffices.  The $p$-adic route would
require perfectoid spaces, of which only the definition has been formalised
\cite{buzzard20}, or the Liquid Tensor Experiment's machinery \cite{scholze22}.
This is a claim about what is formalisable today, not about mathematical depth:
the cyclotomic route is small enough to check
completely, and that is its advantage for the present
purpose; it also reaches the order of the $p$-adic
angle: at a split prime $p\nmid n$ the residue of
$\golden$ has order $n$ exactly when it is a root of
$\Phi_n$ in $\mathbb{F}_p$
(\LeanInline{W4Isometria.clock\_iff\_cyclotomic\_root}
\tier{P}), while the lift itself stays classical.  The limits are visible too: unbounded
self-adjointness is not in Mathlib \cite{gonzalez25hp}, and the analytic steps of
\cref{thm:explicit} are not formalised here.  The first window of Weil positivity is
formalised in its four steps, with the two inputs the paper declares --- Gauss's
integral for the digamma function and the numerical values of the potential --- as
named hypotheses (\cref{thm:first-window}).

% ---------------------------------------------------------------------
\subsection{Connection with CR: the unit and the module}\label{sec:disc-cr}

The first piece of the series is \cite{gonzalez26mersenne}, CR below.  In CR the Euler rotor is read
as the Vitruvian module --- an origin, a radius, an angle and a magnitude
--- and formalised with $\golden$, the
proportion that Pacioli made canonical \cite{pacioli1509} after Vitruvius
\cite{vitruvius} and that recurs through the Renaissance and Neoclassicism; that paper
proves, among its simultaneous identities, the clause $3\golden^{\sigma(p)}=2^p$ on the
known Mersenne exponents.  \cite[\S1.7.2]{gonzalez26cw} reads the same module as the vanishing point of classical
perspective, the datum from which space is disposed on the picture.  The complex plane inherits the same
triple --- origin, modulus, angle, which organise every $z=|z|e^{i\theta}$ --- from that geometric line and
from an algebraic one, through Cardano's \emph{Ars Magna} (1545), Gauss and Argand \cite{argand1806}; and the
predecessor version \cite{gonzalez26oddzeta} compared what Brunelleschi formalised for perspective, three
dimensions encoded in two with the ratios preserved, with what $\golden^2=\golden+1$ does for the descent to
$\Fone$, the multiplicative relations surviving the loss of addition.

The module has a second line of descent, the one \cite{gonzalez25hp} traces: Apollonius classified
the conics by their parameters; Vitruvius took that parametric logic for the module --- origin,
radius, angle and magnitude, formalised in \cite{gonzalez26mersenne}; Pappus reduced the family of
conics to one parameter, the eccentricity; and the line continues to the moduli spaces, the
classification under equivalence \cite{gonzalez26cw}.  The zeros are classified here in exactly that
way.  The Apollonius coordinate puts each on a confocal ellipse of eccentricity $e=2/t$
(\cref{rmk:pappus}), and $e=1$ --- the degenerate conic --- is the radius $1$ of the module, the image
of the zero on the unit circle and the isometry (\cref{prop:li-map,prop:phase-modulus}).  For every
zero the centre is $\tfrac12$, the pair shares its conic, $e\le1$ and $e>0.99993$; what remains open is
the last margin, $e=1$.

This paper reaches each datum of the module on the zeros of $L(s,\chiv)$.  The origin is the apex
$\tfrac12$, the centre of the orbit of every zero (\cref{prop:orbit-embedding}), at which the primes, the zeros,
the operator, the monoid and the archimedean place meet (\cref{prop:membership-proved}); the angle is carried
by the classes of the primes and by the heights of the zeros (\cref{prop:G,prop:phase-modulus}); the
magnitude is the tower of $\golden$, on which every prime has its coordinate (\cref{prop:euler-tower}); and
the radius is $r(\rho)=|\rho-1|/|\rho|$, the one datum the classification above leaves open.  Its unit
is the one of Euler's identity: $e^{i\pi}+1=0$ is a rotor of radius $1$.

In the language of CR, then, what is open is the radius alone: that the module has radius exactly
$1$ at every zero.  A radius is fixed by an isometry, and the isometry is the operator's, which the
next section takes up.

\subsection{Connection with the Primitive Structure}\label{sec:disc-hp}

\subsubsection{What the Primitive Structure declared, and what is settled here}\label{ssec:disc-hp-ledger}

The Primitive Structure \cite{gonzalez25hp} constructed the operator
$\tilde H_{\PCF}$ and closed its argument conditionally on named inputs.  This paper
was written to settle them.  The account, input by input:

\begin{center}
\begin{tabular}{>{\raggedright\arraybackslash}p{4.2cm}>{\raggedright\arraybackslash}p{3.2cm}>{\raggedright\arraybackslash}p{5.2cm}}
\toprule
Input of \cite{gonzalez25hp} & Status here & Where \\
\midrule
\texttt{hecke\_1920} (Hecke 1920) & discharged \tier{P}& \LeanInline{hecke\_chi5\_derived}; functional equation, companion zero, conjugation: \cref{sec:analytic-inputs} \\
Seven geometric axioms ($\golden^2=\golden+1$, $i^2=-1$, $z=\golden y$, basis of $E^3$, continued fraction, nested radical, $\|\Omega\|=\tfrac12$) & theorems \tier{P}& Block~1 of the unified Lean file, \texttt{F1\_compilado\_sin\_axiomas}, with $0$ axioms (\cref{app:lean}) \\
Arithmetic--spectral span & constructed \tier{P}& \cref{prop:span}, with $G$ as common source (\cref{prop:G}) \\
Membership of the zeros & derived from the bridge \tier{P}& \cref{thm:membership} \\
$\Lambda$-ring on $\Zphi[\golden^{-1},\tfrac12]$ & corrected & monoid between levels, Artin structure within: \cref{rmk:monoid-not-ring,sec:lambda-ring} \\
\texttt{bridge\_axiom} & carried to its last link, which the closure conjecture would give & \cref{ssec:core-forms,cor:closure} \\
The operator's spectrum vs.\ the zeros & located exactly & the envelope places the zeros to a fraction of a spacing; the precision is \cref{conj:closure} \\
\bottomrule
\end{tabular}
\end{center}

Two of these change the reading of \cite{gonzalez25hp}.  The Hecke input is no longer
an axiom, so the operator's argument now rests on theorems up to the last link of the bridge.
And the bridge, which \cite{gonzalez25hp} took as a single axiom, is now a chain whose
every link but the last is proved (\cref{sec:topological-proof}): the span exists, the
arrow from the operator's integrality to the tower is closed (\cref{thm:arrow}), and
what remains is that the coordinate $|\rho-1|/|\rho|$ of each zero $\rho$ is a power of $\golden$, in its forms of
\cref{ssec:core-forms}, and, along Weil's route, the bound (W4) of \cref{thm:weil-closure}, which is \textup{(W3)} on the
planes of the orbit (\cref{prop:W4-from-W3}).

The triple obstruction of \cite{gonzalez25hp} distinguishes problems that are not the same, as \cref{sec:polya-hist}
set out.  One is the circularity
of an operator built from the zeros or from the primes it is meant to explain, which \cite{gonzalez25hp} compares with a
circularity of another kind, the diagonal argument of Lawvere and Yanofsky.  Another is the binary against the third
dimension, bounded by Frobenius and Hurwitz.  Another is the relation between the symplectic and the unitary classes,
since the Wigner--Dyson classification is itself binary.  With \cref{prop:level-operators,prop:symmetry-spectrum}
each can be seen inside the tower.  The diagonal argument needs an endomorphism without fixed point: $\alpha(x)=1-x$ is one on the truth values
$\{0,1\}$, and on a commutative ring it has a fixed point exactly when $2$ is invertible, the point $\tfrac12$,
none on $\Z$ or on $\Zphi$ (the paragraph after \cref{prop:cocone}); the mate $\rho\mapsto\alpha(\bar\rho)$, its continuation
to $\C$ composed with the conjugation, fixes exactly the line (\cref{prop:klein-detail};
\LeanInline{EstructuraEnGUE.lawvere\_alpha\_no\_fixed\_point}, \LeanInline{F1Lawvere.reflection\_fixed\_iff}
\tier{P}).  The binary problem is a symmetry of order two in a spectrum: an antiunitary symmetry with $T^2=-1$ fixes no
ray and forces every level into a doublet, the
symplectic class; a commuting symmetry of order two splits the spectrum into two sectors with no repulsion between
them, as the zeros of $\zeta_K$ show.  The dimensional obstruction is not met by a division algebra either: what separates
the torsion fibre does not commute with the deck step, and with it generates the matrix algebra of order six,
$DY=\zeta_6\,YD$, and not $\mathbb H$ or $\mathbb O$; the colour operator $\hat\Omega$ is normal and not self-adjoint
(\cref{prop:level-operators}).  A zero off the line would be a doublet of ordinates, two zeros at separation zero
(\cref{thm:shared-ordinate}; the configuration of arity two of \cite{gonzalez26cw}), the configuration of the symplectic
class at one level; the zeros computed show none (\cref{ssec:envelope-fluctuation}) \tier{N}, and that no zero shows it is form~(iv) of \cref{thm:topological-proof}.  In the tower the binary is the orientation, and it
enters with the colour in an operator that does not commute with the deck step, which is why no sector and no doublet
appear.  In the tower the circularity of the diagonal argument and the binary of a spectrum meet at one point, the passage of $\alpha$ from $\{0,1\}$ to
a ring in which $2$ is invertible, where the missing fixed point appears and, through the mate, becomes the line.  This last reading joins facts that are \tier{P}\ and \tier{N}.

\subsubsection{The operator and the envelope}\label{ssec:disc-hp-operator}

\cite{gonzalez25hp} builds a self-adjoint operator whose eigenvalues follow the
zero heights.  Its discrepancy splits into two terms: the fluctuation $S$, which no
operator built from the envelope can carry, and the closed-form parametrisation of
the envelope inverse, which accounts for most of the error below $1.7\%$ reported in
\cite{gonzalez25hp} (about $1\%$ at $n=500$) and is removable only by the exact inverse,
not by any closed form, the structural one being the best of these
(\cref{prop:error-partition,rmk:two-errors}).  What the envelope carries is the conductor, and the zeros select it
(\cref{prop:envelope-exact}).  So the operator side and the arithmetic side meet at
a sharp statement: the structure fixes the envelope, conductor included, and leaves
$S$, and the core (\cref{thm:core-positivity}) is what would control $S$.  The closure conjecture says what would control it (\cref{conj:closure}): the operator of \cite{gonzalez25hp} is not an approximation to
be discarded; its envelope, inverted exactly, already places the zeros to a fraction of a spacing
(\cref{prop:positions}), and what remains is the precision a law of the primes by classes has to
supply.

The statistics of the zeros can be read against \cref{prop:level-operators}.  Bohigas, Giannoni and Schmit
attach to a quantum system without further symmetry the statistics of the ensemble of its class \cite{bgs84}, and
systems with arithmetic symmetries, in which further commuting operators exist, depart from those statistics
\cite{bggs92,keating91}.  In the tower, an operator of the level fixes the envelope, conductor included
(\cref{prop:envelope-exact}), and cannot separate the fibre; the GUE statistics of the zeros is then what such an
operator shows of a structure whose fibre it does not resolve.  This is a reading, not a theorem: it does
not say that the zeros carry the fibre, only where, in this framework, an operator built on the level stops.

\subsubsection{Continuing the framework}\label{ssec:disc-hp-framework}

What \cite{gonzalez25hp} built as an operator, this paper builds as a thread: the
operator of Stage~5 is the same object, but it is now reached from $\golden$ through
the primes, the rings and the Euler products, and it is followed by the passage to
membership.  The Hilbert--P\'olya programme and the $\Fone$ programme meet in that
thread, which is the sense in which this paper continues the framework rather than
replacing the operator.

\subsubsection{How far the operator has to be taken, and how the argument closes}\label{ssec:disc-hp-closure}

The law of \cref{conj:closure} has to give
    magnitudes, and nothing else (\cref{prop:what-fluctuation-takes,prop:euler-tower}).  Two parts of this paper
    already meet it from both sides.  The descent to $\Fone$, which answers Manin's requirement
    (\cref{sec:manin,ssec:rel-manin}), gives the input at every prime: $\Zphi$ is a $\delta$-ring at every
    $p\neq5$, with global $p$-derivations (\cref{prop:delta-ring}), the Euler factor is the characteristic
    polynomial of its Frobenius (\cref{thm:frob-eigen}), and the magnitude of each prime is its
    valuation (\cref{prop:euler-tower}).  The operator gives the exact output where it already exists: it
    separates magnitude and angle, $\hat\Omega_k=\golden^{-\lambda_{\log}}\zeta_3^{\,k}$, joins $\pi$, $\golden$ and the binary place
    in $\tfrac12$, carries the classes, joins the rotations of the Gauss and the Eisenstein lattices,
    is the $2$-adic polar form of $\golden$, and places the zeros up to the correction
    (\cref{thm:operator-forms,prop:omega-two-adic,prop:positions}).  Through it the places are already glued exactly: at the apex $\tfrac12$, in the angles of the classes and in the binary magnitude (the paragraph that follows \cref{prop:euler-tower}).  The binary output sits at the level $3$ of the clock:
the angle $\zeta_3$ of $\hat\Omega$ is that of $\golden$
at $2$, and $N(\Phi_3(\golden))=4$.  Every prime has its
level, and on each level the product formula on
$N(\Phi_n(\golden))$ already glues the archimedean and
the $p$-adic readings, in sum over the primes of that
level (\cref{prop:clock-cyclotomic}).  What is missing is
the output at every other prime one by one: their order
by size, which the clock does not fix, and which is the
second step of \cref{conj:closure}.  Read by levels, the
accumulated clock carries the Mertens function
(\cref{prop:clock-cyclotomic}(vi)), and through it,
classically, the zeros of $\zeta$; read by size, inside
each column, it is the deviation of $\golden$ from a
uniform vertex of the $n$-gon of each prime, layer by
layer, whose layers $2$, $3$ and $5$ are the forms of
levels $80$ and $135$ and of the quintic layer
(\cref{prop:index-ngon}(iv) for level $80$, numerically
for the other two).  The setting of that
output is the third of Connes'
    strategies, a suitable cohomology \cite{connes15essay}, in the $p$-adic form, which the construction enters
from the cyclotomic side through the clock
(\cref{sec:manin,sec:choice}), now with what the
cyclotomic route gained: prisms take $\delta$-rings as input \cite{bhattscholze22}; the perfection $\golden^{1/p^\infty}$ of the tower
carries the magnitude and the cyclotomic towers
$\mu_{p^\infty}$ the angle --- on the clock, the
divisors $n_0p^k$ over a prime of clock $n_0$ ---, $\mu_{20}=\mu_4\times\mu_5$,
    where the eight classes live, being a finite level of those at $2$ and $5$; the tilt \cite{scholze12} carries each $p$-adic reading, exactly, to characteristic $p$, where the analogue of Weil's exactness is a theorem (\cref{ssec:weil-curves}); and the product formula joins it with the archimedean reading, as the
    operator already does at $2$ (\cref{prop:omega-two-adic}).  This is the modern form of the $p$-adic route
    Manin asked for, now with the global $\delta_p$ and with the exact template of the operator; the prime
    $5$, with no Frobenius lift (\cref{thm:obstruction-five}), needs a prism of its own.  Connes and
    Consani have recently built an absolute arithmetic curve over $\Fone$ as a common origin of
    $p$-adic Hodge theory, complex analytic geometry and the ad\`elic scaling site, in which, at each
    prime, a real locus $\mathbb{R}_+^\times/p^{\mathbb{Z}}$ and a phase space independent of $p$ appear separately
    \cite{connesconsani26}: the same separation of magnitude and angle, and the nearest contemporary
    setting for this step.  The formalisation of this geometry stands where \cref{ssec:lean-vs-padic} records, and the finite part of the route is formalised here.

The route can be stated as a conjecture, with each clause tied to what is already proved.

\begin{conjecture}[The prismatic route to the closure]\label{conj:prismatic-route}
There is a cohomology $H$ attached to $\Zphi$, assembled from the prisms $(\Z_p[\golden],(p))$ at the
primes $p\neq5$, where $\Zphi$ is a $\delta$-ring (\cref{prop:delta-ring}), and from the Breuil--Kisin prism $(\Z_5[[u]],(u^2-5))$ at $5$, where no Frobenius lift exists (\cref{thm:obstruction-five}), with Frobenius
operators $F_p$, such that:
\begin{enumerate}[label=\textup{(\arabic*)}]
  \item \emph{local factors:} $\det(1-p^{-s}F_p\mid H_p)$ is the Euler factor of $L(s,\chiv)$ at $p$,
    as it already is for $\Zphi$ (\cref{thm:frob-eigen});
  \item \emph{the binary place as template:} at $p=2$ its component is exactly the operator's ---
    angular part $\zeta_3$, the level $3$ of the clock,
modulus $|2|_2$, anchor $\golden^{\lng}=2$ (\cref{prop:omega-two-adic,prop:euler-tower});
    \item \emph{gluing:} the archimedean and $p$-adic
    components are glued by the product formula, as the
    operator already does at $2$, so that the magnitude
    of each prime is $\sigma_p=\log_\golden p$
    (\cref{prop:euler-tower}); on each level $n$ of the
    clock this gluing is exact in sum, by the product
    formula on $N(\Phi_n(\golden))$, the primes of
    clock $n$ together with those of
    \cref{prop:clock-cyclotomic}(iv),
    and what it leaves is the order of the primes of each
    level by size;

  \item \emph{assembly:} over the apex $\tfrac12$ and the eight classes, the magnitudes of all the
    primes give the correction of \cref{conj:closure} at resolution $1/T$, and the zeros land on the
    spectrum of $\hat\Omega$;
  \item \emph{positivity:} $H$ carries a positive pairing for which each $F_p$, normalised by $p^{-1/2}$,
    is an isometry.
\end{enumerate}
Then the spectrum is exact, the Frobenius is unitary, and every zero lies on the line, its mode being
an eigenvector of an isometry (\cref{prop:phase-modulus}); on Weil's side this is (W4)
(\cref{thm:weil-closure}).  The two halves of \Ecore\ are reached by different clauses:
the multiplicative half is already proved (\cref{prop:two-halves}(i)); clause~(5) gives the
additive half, since the eigenvalue of an isometry has modulus $1$, that is $r=1$, which with the
multiplicative half is $t=2$; and clauses~(1)--(4) give the magnitudes and the correction of the
heights.
\end{conjecture}

\begin{center}\small
\begin{tabular}{p{2.6cm}p{4.1cm}p{3.1cm}p{3.2cm}}
\toprule
Clause & Already proved, at the level of $\Zphi$ & Open & Refuted if \\
\midrule
Prisms at $p\ne5$ & $\Zphi$ is a $\delta$-ring (\cref{prop:delta-ring}) \tier{P}; a $p$-complete, $p$-torsion-free $\delta$-ring with $(p)$ is a crystalline prism \cite{bhattscholze22} & the formalisation & --- \\
(1) local factors & characteristic polynomial (\cref{thm:frob-eigen}) \tier{P}& that it comes from $H$ & the factor of $H$ differs from Euler's \\
(2) template at $2$ & $2$-adic polar form, anchor (\cref{prop:omega-two-adic,prop:euler-tower}) \tier{P}& its prismatic reading & the component at $2$ does not give $\zeta_3$, $|2|_2$ \\
(3) gluing & product formula at $2$; $\sigma_p=\log_\golden p$ (\cref{prop:euler-tower}) \tier{P}; exact on each level of the clock (\cref{prop:clock-cyclotomic}), with the clock as cyclotomic polynomials \tier{P}& from the clock to the size & some exact magnitude is not $\log_\golden p$ \\
the prime $5$ & the obstruction (\cref{thm:obstruction-five}) \tier{P}; the prism, distinguished, $\delta(E)\equiv-1$; special fibre $=\delta_5$; defect $\tau(\chiv)^2$ \tier{P}; conductor $5^7$; Kisin class $h=u$ \tier{N}\ (\cref{rmk:prism-five}) & the gluing & --- \\
(4) assembly & the classes and the fluctuation's intake (\cref{prop:what-fluctuation-takes}) \tier{P}; the layers of the index (\cref{prop:phi-layers}) and the energy of their phases (\cref{prop:sqrt-x}) \tier{P}& the law at resolution $1/T$; the scale $\tfrac12$ of the power-residue symbols & the assembled correction does not give exact heights \\
(5) positivity & modulus $1$ $\Leftrightarrow$ the line (\cref{prop:phase-modulus}) \tier{P}; windows (\cref{thm:first-window}); (W4) certified for $k\le25$ (\cref{ssec:weil-closure}) & the pairing & (W4) fails \\
\bottomrule
\end{tabular}
\end{center}

At $5$ the row of the table is filled by the Breuil--Kisin prism of \cref{rmk:prism-five}; its
Kummer representation is mixed, so it contributes the regulator and the ramification, not the
zeros.

\paragraph{A reading, not a theorem.}  The descent datum also has two halves: the multiplicative,
between the levels of the tower (\cref{thm:lambda-monoid}), and the additive, inside each level modulo
$p$ (\cref{thm:frobenius-congruence}), and $\delta_p(\golden)$ measures their difference at each prime
(\cref{ssec:delta-polar}).  A prism is a $\delta$-ring.  So the route of \cref{conj:prismatic-route}
works with the object that measures how the two halves of the descent are glued, as the open
statement asks how the two halves of the coordinate are: the multiplicative half proved, the
additive half to be supplied.

\begin{remark}[A model in the pentagonal field]\label{rmk:fermat-model}
The Fermat curve $x^5+y^5=1$ has good reduction at every $p\ne5$ and bad reduction at $5$.  At
$p\equiv1\pmod5$ the eigenvalues of Frobenius on $H^1$ are Jacobi sums of order $5$~\cite{weil49}, of
modulus $\sqrt p$: normalised by $p^{-1/2}$, Frobenius is an isometry --- there the conclusion of Weil's
proof, through the Hodge index on $C\times C$ --- and each local factor is its
characteristic polynomial.  So clauses~(1) and~(5) hold locally away from $5$ and fail at $5$, as
for $\Zphi$ (\cref{thm:obstruction-five}).  Their $p$-adic values are given through the Teichm\"uller
character~\cite{grosskoblitz79}, the object of the angular part (\cref{prop:delta-ring}).  For $p=11,31,41,61,71$ the projective points of $x^5+y^5=z^5$ number exactly
$p+1+\sum J(\chi^a,\chi^b)$ over the twelve pairs with $a+b\not\equiv0\pmod5$, and all twelve sums have
modulus $\sqrt p$ \tier{N}.  The model
is not $L(s,\chiv)$, and the local isometry does not give the line.
\end{remark}

That is how the argument closes.  The construction and the operator give each zero the centre $\tfrac12$ of its pair, its pairing and its height up to the correction; what remains is the radius, and with it the exact
height.  The radius is $1$ exactly when the magnitudes of the primes are known at the resolution the
height requires; those magnitudes are what the law of \cref{conj:closure}
has to supply.  On each level of the clock they are
already exact in sum, by the product formula on the
cyclotomic norms (\cref{prop:clock-cyclotomic}); what
the law adds is each one in its place, and the
prismatic and perfectoid setting is where, prime by
prime, they can be made exact
(\cref{conj:prismatic-route}).  Neither conjecture is
proved here.

\subsubsection{The binary zone: from obstruction to route}\label{ssec:disc-binary}

Each open clause of \cref{conj:prismatic-route} concerns one of three objects: the primes $2$ and $5$ of the conductor and the eight classes of $\mu_{20}$ (clauses (1)--(4)), or the isometry of \cref{ssec:intro-binary} (clause (5)).

\paragraph{What the construction reaches.}  A Frobenius lift at every prime other than $5$
(\cref{sec:lambda-ring,thm:frobenius-congruence}); the exact component of the operator at the place $2$
(\cref{prop:omega-two-adic}); the multiplicative half of \Ecore\ on every zero (\cref{prop:two-halves});
and the isometry $\mathcal V\colon\mathbb{C}\to\mathbb{C}^3$ of the triad of the operator, with modulus $\tfrac12$ at
every level of the tower (\cref{ssec:why-phi}).  The construction starts from $\golden$ and the pentagon; the place $2$ enters through the anchor $\golden^{\lng}=2$ (\cref{sec:mersenne}) and through the conductor $4$ of $\chi_4$ (\cref{prop:G}).

\paragraph{The two primes of the conductor.}  The eight classes correspond to the roots of $\Phi_{20}$, in
$\mu_{20}=\mu_4\times\mu_5$.  At $5$, the ramified prime of $\Q(\sqrt5)$, $\Zphi$ has no Frobenius lift
(\cref{thm:obstruction-five}), so clause~(1) needs there a prism other than $(\Z_5[\golden],(5))$.  At $2$
the component of $\OmH$ is exact, with angular part $\zeta_3$ and modulus $|2|_2=\golden^{-\lng}$
(\cref{prop:omega-two-adic}), and the magnitude $\sigma_2=\lng$ is glued by the product formula; for
$p\ne2$, $\sigma_p=\log_\golden p$ is not a level of the tower, and its gluing is clause~(3)
(\cref{prop:euler-tower}).

\paragraph{The classes that $G$ identifies.}  $(\Z/20)^\times=\langle3\rangle\times\langle11\rangle$ and
$\ker G=\{1,9\}=\{1,3^2\}$; the fibres of $G$ are $\{1,9\}$, $\{3,7\}$, $\{11,19\}$ and $\{13,17\}$
(\cref{prop:G}).  A character of $(\Z/20)^\times$ factors through $G$ exactly when it is trivial on $9$,
that is, when $\chi(3)^2=1$, which with $\chi(11)^2=1$ means that $\chi$ is quadratic; the other four
have $\chi(3)=\pm i$ and $\chi(9)=-1$.  Reduction modulo $5$ separates the classes of the kernel,
$1\mapsto1$ and $9\mapsto4$ (\LeanInline{the\_four\_fibres}, \LeanInline{four\_fibres\_four\_characters},
\LeanInline{carriers\_have\_order\_four}, \LeanInline{Z20star\_pi2\_fiber\_is\_mod5} \tier{P}).  The descent
reads only $\chiv$, a quadratic character (\cref{thm:chi5-from-frobenius}).  Clause~(ii) of
\cref{conj:closure} bounds the residual of each class separately, and by the orthogonality of characters
that is a statement about all eight characters modulo $20$, the four of order $4$ included, which do not
factor through $G$.  So the law of \cref{conj:closure} needs data that $G$ does not retain.

\paragraph{The isometry.}  Clause~(5) is the isometric step.  The eigenvalues of an isometry of a
finite-dimensional inner-product space lie on the unit circle, so with it the mode of each zero has
modulus $1$, and with the multiplicative half this is the additive half of \Ecore\
(\cref{prop:phase-modulus,prop:two-halves}).  The isometry of the triad exists
(\cref{ssec:why-phi}); that the modes of the zeros are modes of an isometry is not proved.

\paragraph{The circle of each prime.}  The arithmetic angle with which a prime enters an $L$-function or a Gauss sum forms a ladder.
The half-turn $\chiv(p)=\pm1$ is fixed by a rule, and so is the sign of the quadratic Gauss sum,
$\tau(\chiv)=+\sqrt5$.  The quarter-turn is the order-$4$ part that $G$ collapses (\cref{prop:G}).
The fifth-turns are the residue symbols of order $5$ (\cref{prop:phi-residue}); their Jacobi sums
are Hecke characters, whose angles equidistribute~\cite{hecke20}, and the Gauss sums are fifth roots
of them, $g(\chi)^5=\chi(-1)\,p\prod_{k=1}^{3}J(\chi,\chi^k)$.  Which root is Kummer's question: no rule is known~\cite{kummer1846,patterson78}; the angles equidistribute, for the cubic sums~\cite{heathbrownpatterson79} and for every order~\cite{patterson87}; and in the cubic case, under GRH, they carry a bias~\cite{dunnradziwill24}.  The binary angle is exact; the higher ones
are statistical.  To see $\mu_{20}=\mu_4\times\mu_5$ is to see the quarter-turn and the fifth-turns.

\paragraph{Digits, prisms and Weil.}  The ladder is read in three places of the paper, and each reads a
different part of it.

\emph{Digits.}  By \cref{prop:base-twenty} the class of $p$ in base $20$ is its Frobenius in
$\mathrm{Gal}(\Q(\zeta_{20})/\Q)\cong\Zst$.  The decoration is its image under $G$, the last decimal digit
is the reduction modulo $5$ above, and the two digits are the two primes of the conductor.  The digit
sees the half-turns and the quarter-turn.

\emph{Weil.}  Weil's template runs on the same level (\cref{ssec:rel-weil}): on the lattice of the eight
classes the Frobenius has order dividing $4$ (\cref{prop:eight-classes}), so its eigenvalues $\pm1$, $\pm i$
are the half- and quarter-turns the digit sees, and on the windows it is the half-turn at $2$ that carries
the positivity (\cref{prop:semilocal}).  Like the law of \cref{conj:closure}, it takes only the class
(\cref{prop:what-fluctuation-takes}).

\emph{Prisms.}  No digit sees the fifth-turns: $11$ and $211$ lie in the same class modulo $20$, and
$\golden$ is a fifth power at $211$ and not at $11$, for either root of $x^2-x-1$ (\cref{prop:phi-residue};
\LeanInline{F1PaperNew.phi\_residue\_not\_class\_function} \tier{P}).  They are read on the $p$-adic circle
of the prisms $(\Z_p[\golden],(p))$ from which the cohomology of \cref{conj:prismatic-route} is assembled;
the order of the angle of $\golden$ on that circle is the period of the Fibonacci dynamics on the $p$-torsion of the
real torus
(\cref{prop:two-faces}).
In the model of \cref{rmk:fermat-model} the two sides meet: the isometry is the conclusion of Weil's
proof, and its eigenvalues are computed on the $p$-adic circle.

So the class is what Weil's template and the law need, and the angle is what the prisms carry.  The
cohomology of \cref{conj:prismatic-route} is where the two would have to be one object, and Weil's proof
says what that object must have: a positivity, which for a curve comes from the Hodge index on
$C\times C$, classical; open for the cohomology.

\paragraph{From obstruction to route.}  $G$ is not injective on the classes: its fibres have two elements.
A cohomology that yields clause~(ii) must therefore separate the classes in each fibre of $G$, so it must
see $\mu_{20}$, a finite level of the cyclotomic towers at $2$ and $5$; with a prism at $5$ and a positive
pairing, this is the cohomology of \cref{conj:prismatic-route}.  That it exists is not proved; the paper
establishes where the open step lies, what the construction must satisfy clause by clause, and what would
refute each clause.

\subsection{Connection with physics}\label{sec:disc-cw}

Two companion papers carry this construction to physics, on nearly the same subject: the Crystalline
Worldsheet~\cite{gonzalez26cw}, which places a string in the torus and reads the tower as a bulk--boundary
system, and the de Sitter companion~\cite{gonzalez26ds}, which puts one object in place of an ensemble for
the observer in de Sitter space.  This section records where the mathematics of this paper meets them.

The isometry of \cref{ssec:why-phi,ssec:disc-binary} is the one the Crystalline
Worldsheet~\cite{gonzalez26cw} uses on the physical side: there it is the bulk--boundary map, with modulus
$\tfrac12$ at every level, and the unfolded spacings of the zeros favour the Gaussian Unitary Ensemble
over Poisson by a factor $7.6$ \tier{N}\ \cite[\S2.13]{gonzalez26cw}.  That paper places a ring inside a torus in order to carry a
worldsheet, and uses it to attempt a bridge --- a hybrid --- between holography and string theory.  This section follows that thread.  Holography lacks an individual
operator: a Hermitian operator with a ground state from which everything else emerges;
in its place it writes the Gaussian unitary ensemble, the same statistics that govern
the zeros of $\zeta$ (\cref{ssec:holography-problem}).  The bridge is the ring in the
torus (\cref{sec:worldsheet}); it meets strings through the quantum tori that Manin
proposes for the $\Fone$ programme, and the operator through the construction of the Primitive Structure.  Its aim is to recover the operator from the bridge.  For that, this
paper and the Primitive Structure bear on the problem mathematically, since both
address the structure the ensemble describes (\cref{ssec:individual-ensemble}), and
\cite{gonzalez26cw} addresses it physically (\cref{sec:microstate}).

The framework extends that of CR (\cref{sec:disc-cr}).  There the Euler rotor is a
vanishing point, formalised with $\golden$; here the same point is made to emerge
dimensionally --- from the line to the plane of the Euler product and of the zeros, and
from there to a third dimension over the plane, obtained by symmetry
(\cref{ssec:one-two-three}).  That structure is the base of the operator
(\cref{ssec:operator-by-dimension}), and it is the structure \cite{gonzalez26cw} uses to
build a bridge between strings and holography with a single operator
(\cref{ssec:one-operator}).

\subsubsection{The problem: holography without an operator}\label{ssec:holography-problem}

In the holographic dictionary the boundary theory supplies the Hamiltonian.  In low
dimensions, however, a gravitational bulk has no single boundary dual: two-dimensional JT gravity is recovered only as an average over an ensemble of boundary theories, a matrix integral \cite{saad19}, and in three dimensions ensembles of two-dimensional conformal field theories, averaged through the equidistribution of Hecke operators, play the same role \cite{cooper26}, although quantum gravity should carry its own
individual microstates.  What is missing is an individual --- a Hermitian operator with
a ground state, from which the spectrum and the states would follow.  The ensemble is
of Gaussian unitary type, and the same statistics govern the pair correlation of the
zeros of $\zeta$ (\cref{sec:polya-hist}).  An ensemble is what one writes for want of an
individual \cite[\S2.13]{gonzalez26cw}.  Seen from the physical side, the problem is the
Hilbert--P\'olya problem: an operator whose spectrum the ensemble describes only
statistically.

\subsubsection{The bridge: a ring inside a torus}\label{sec:worldsheet}

The objects of \crefrange{sec:choice}{ssec:base-phi} were not assembled for a problem in number theory.  They
come from a reinterpretation of the closed string as an $\Fone$ ring living on an
M-theory torus, used as the place from which an observer looks
\cite[\S1.7]{gonzalez26cw}.  The word for that place in the classical theory of
perspective is the \emph{module}: an origin, a radius, an angle and a magnitude,
the datum that fixes how space is disposed on the picture plane.  The proposal is
to make the module exact by making it a ring,
\[
  \Rpcf=\Z[\golden,\golden^{-1},\tfrac12],
\]
carried on the torus $T^2_{\PCF}$ of \cref{rmk:companion-interpretation}; and the
relation that seeds it is $\golden^2=\golden+1$ in the form
$\golden=1+1/\golden$, a self-reference with a fixed point rather than a
contradiction.

Two features of that construction meet objects of \cref{sec:methods}.

\emph{The self-dual point.}  The modular parameter of the torus is $\tau=i$, the
fixed point of $\tau\mapsto-1/\tau$ (\LeanInline{F1CW.s\_duality\_fixes\_i} \tier{P}),
and the same involution appears in the Schwinger variable as $t\mapsto1/t$, whose
unique fixed point is $t=1$ (\LeanInline{selfdual\_fixed\_point} \tier{P}).
The period of the torus is $T_K=2\pi\log\golden=2\pi R_K$: the regulator of
$\Q(\sqrt5)$, which is the invariant that reappears in \eqref{eq:log-two}.

\emph{The fibre.}  The module direction closes, once around the fundamental
domain, with the half-turn $e^{i\pi}=-1$: an orientation-reversing involution, so
the worldsheet is one-sided (\LeanCW{PaperMobius.fibre\_monodromy},
\LeanCW{fibre\_reflection\_order\_two} \tier{P}; in this paper's corpus the orientation flip
of the fibre has order exactly two, \LeanInline{orient\_flip\_order} \tier{P}).  The fibre datum is
independent of whether the base torus is commutative or a noncommutative
$T^2_\theta$ (\LeanCW{fibre\_independent\_of\_base} \tier{P}).

\emph{Where the bridge meets strings and the operator.}  On the side of strings, since
the fibre does not depend on the base, the construction sits on the noncommutative tori
$T_\theta$ that Manin proposes, with $\theta$ real quadratic, as the quantum tori of real
multiplication and of the $\Fone$ programme (\cref{sec:manin}), and on the
Connes--Douglas--Schwarz torus of M-theory compactification \cite{connes98}.  On the side of the operator it meets the construction of the Primitive Structure \cite{gonzalez25hp}: the operator built there from the same $\golden$
is the candidate for the individual the ensemble replaces, and \cref{ssec:operator-by-dimension}
shows that it carries the structure of the bridge.

\subsubsection{The mathematical individual: what it determines, and where the ensemble enters}\label{ssec:individual-ensemble}

This paper and the Primitive Structure address the structure of the ensemble from the
mathematical side.  There the individual is the operator of \cref{thm:operator-forms}, built
from $\golden$ in \cite{gonzalez25hp} and here, with no physics; on the physical side it
is the single microstate of modulus $\tfrac12$ of \cite{gonzalez26cw}
(\cref{sec:microstate}), with no ensemble average.  The question is then precise: what
does the individual determine, and where does the ensemble still enter?

\begin{proposition}[What the individual determines, and where the ensemble enters]\label{prop:individual-ensemble}
\begin{enumerate}[label=\textup{(\roman*)}]
  \item The envelope of the zeros is determined by the structure, conductor included:
    the zeros return $q=5$ (\cref{prop:envelope-exact}).
  \item The fluctuation $S$ around the envelope is not determined by any operator built
    from the envelope, and the pair-correlation kernel
    $K(u)=1-\bigl(\sin\pi u/\pi u\bigr)^2$ of the ensemble is a statement about $S$; it
    equals $1$ at every non-zero integer separation.
  \item The configuration the ensemble makes improbable --- two zeros at the same
    ordinate --- is, off the critical line, a concrete pair: $\rho$ and its mate
    $1-\bar\rho$ are distinct and share an ordinate (\cref{thm:shared-ordinate}).
  \item Hence, for a zero set closed under the mate, every zero lies on the line if and
    only if no two zeros share an ordinate.  This is link~(iv) of
    \cref{thm:topological-proof}: the individual, one-bit form of the level repulsion
    the ensemble describes, and strictly weaker than any metric bound on the spacings.
\end{enumerate}
\LeanInline{extreme\_case\_is\_binary}, \LeanInline{E\_iff\_no\_shared\_ordinate},
\LeanInline{repulsion\_excludes\_zero\_spacing},
\LeanInline{F1CW.sineKernel\_at\_nonzero\_integer} \tier{P}; (i) \tier{N}\ in its measurement,
\tier{P}\ in the injectivity of the scale; (ii) the kernel
\end{proposition}

\begin{proof}
(i) is \cref{prop:envelope-exact}.  (ii): the envelope is a smooth function of the height,
and $S$ is what remains; $\sin(\pi n)=0$ for $n\in\Z\setminus\{0\}$.  (iii) is
\cref{thm:shared-ordinate}.  (iv): off the line the pair of (iii) shares an ordinate, and
conversely two distinct zeros on the line with one ordinate would coincide.
\end{proof}

So where holography averages an ensemble for want of an operator, the construction
has the operator, and of everything the ensemble statistics describe beyond the
envelope, the verdict needs one bit: whether any two zeros share an ordinate.  That bit is link~(iv) of \cref{thm:topological-proof}.  \cite[Remark~2.88]{gonzalez26cw} records why the ensemble's
statistics are not reproduced by the operator: they reside in the fluctuations, and reproducing
them would require either an exact spectrum or a dynamical mechanism generating the level
repulsion.  The closure conjecture of this paper (\cref{conj:closure}) is the first of those
alternatives, made specific: the exact positions, the envelope from the archimedean side and the
correction from a law of the primes by classes.  The mathematical individual is this paper's and the
Primitive Structure's; the physical one --- one microstate, holography without an
ensemble average, the observer with a single pure state --- is developed in
\cite{gonzalez26cw} (\cref{sec:microstate,sec:adscft}).

\subsubsection{The physical individual: one microstate and its modulus}\label{sec:microstate}

The state carried by that ring is a single one, not an ensemble:
$\Omega(\tau)=\tfrac12e^{i\tau\log\golden}$, of constant modulus
$|\Omega|=\tfrac12$ (\LeanInline{modulus\_Omega},
\LeanCW{omega\_flow\_invariant} \tier{P}).  Three things about the number
$\tfrac12$ are proved rather than posited.

It is the value where the binary entropy is largest, $H(p)\le1$ with equality
only at $p=\tfrac12$ (\LeanCW{binary\_entropy\_le\_one},
\LeanCW{binary\_entropy\_half} \tier{P}), so one bit per cell is a maximum and
not a coincidence.  It is reached by eight independent routes that meet at the
same apex --- among them the half-factorial $\Gamma(3/2)/\sqrt\pi$, the product of
the three component norms of the eigenvalue triangle, $\cos(\pi/3)$, the inverse
binary turn $\golden^{-\lng}$, the midpoint of the Galois pair
$(\golden+\bar\golden)/2$, and the fixed set of $x\mapsto1-x$
(\LeanInline{modulus\_cocone\_six\_routes} \tier{P}\ for the six routes of \cref{prop:cocone}; the eight routes are \LeanCW{PaperS2.mu\_diagram\_commutes}).  And it is the value the
functional equation fixes, since $s$ and $1-s$ have the same real part exactly on
$\Re s=\tfrac12$ (\LeanInline{functional\_equation\_fixed\_line} \tier{P}; in \cite{gonzalez26cw} \LeanCW{faceLine\_apex}) --- which is the point of contact
between the physical construction and \cref{prop:stau}.

The last two routes deserve a remark, because they are the reason the number is
the same on both sides.  The midpoint of $\{\golden,\bar\golden\}$ and the fixed
point of $x\mapsto1-x$ are the same computation: Galois conjugation in
$\Q(\sqrt5)$ and the apex involution are one map on the orbit
(\LeanInline{phi\_bar\_eq} \tier{P}).  The arithmetic $\tfrac12$ of \crefrange{sec:choice}{sec:clock} and the holographic $\tfrac12$ of the construction are therefore not two
numbers that agree; they are one number read twice.

\subsubsection{From the vanishing point to the dimensional chain}\label{ssec:one-two-three}

Here the vanishing point of CR (\cref{sec:disc-cr}) is not held fixed but made to emerge dimensionally.  A one-dimensional vector, the line of the tower, becomes
two-dimensional: the plane on which the Euler product and the zeta functions live and
from which the analysis reaches the zeros.  It also becomes three-dimensional, but the
third dimension is geometric, over the plane, and not algebraic: it is reached by
symmetry --- the three eigenvalue directions permuted by $S_3$, on which the ratios
$\golden^{\,i-j}$ act --- a holographic re-encoding of the plane.




The chain is \cref{thm:one-two-three}, proved in \cref{ssec:dimensional-chain}, where the operator needs it: the
line of the exponents, on which the lifts act and from which the Euler product comes; the plane, on which the zeta
functions live; and $E^3$, on which the ratios $\golden^{\,i-j}$ of the operator act.

\begin{remark}[Reading the chain]\label{rmk:reading-chain}
The theorem states identities; the reading is this.  The Euler product comes from the
line, where the lifts act: its index set is the monoid of the positive integers, which acts there
(\cref{thm:one-two-three}(1)); it converges on the plane, in the half-plane $\Re s>1$, where the zeta
functions live (\cref{thm:one-two-three}(2)).  The factorisation
$\zeta=\zeta_K/L(\cdot,\chiv)$ can be read as the projection of the plane onto the
line, with the odd values $\zeta(2k+1)$ where the generator (through $\chiv$) and the
circle (through $\pi$ and the points $a/5$) meet (\cref{thm:odd-zeta}).  The same
three-dimensional class is realised by the three eigenvalue directions of $\OmH$
(\cref{thm:operator-forms}) and by the invariants $\ell=1$, $G_N=\tfrac12$, $c=3$, which
\cite[\S3.5]{gonzalez26cw} shows to be the signature the AdS/CFT corner shares with the microstate.  None of this reading is
used elsewhere.
\end{remark}

The numbers that recur along the chain on the mathematical side --- $|\ker G|=2$, $\mu_3=\tfrac12$, the triangle
that $S_3$ permutes --- are \cref{thm:S3-from-euler}.  The physical reading adds three.

\begin{proposition}[The numbers of the chain in the physical reading: $\ell=1$, $G_N=\tfrac12$, $c=3$]\label{prop:chain-physics}
$S$-duality $\tau\mapsto-1/\tau$ fixes $\tau=i$, and $\ell=1/\ell$ forces $\ell=1$; and the
Brown--Henneaux charge $3\ell/(2G_N)$ equals $3$ exactly when $G_N=\tfrac12=\mu_3$.
These and the identities of \cref{thm:S3-from-euler} are independent; the arrows of the display record the order in
which the construction meets them, not implications between them:
\[
  \underbrace{\prod_p f_p(s)}_{\text{Euler product}}\;\to\;
  \underbrace{|\ker G|=2}_{\text{contraction}}\;\to\;
  \underbrace{\mu_3=\tfrac12}_{\text{norm}}\;\to\;
  \underbrace{S_3}_{\text{symmetry}}\;\to\;
  \underbrace{E^3,\ c=3.}_{\text{geometry}}
\]
\LeanInline{F1v2.S3\_E3\_c3\_from\_the\_euler\_product} \tier{P}, which states the whole display; the three of this
statement are \LeanInline{F1CW.s\_duality\_fixes\_i}, \LeanInline{selfdual\_fixed\_point},
\LeanInline{F1CW.hinge\_central\_charge} \tier{P}%
\end{proposition}

\begin{proof}
$-1/i=i$, and $1/\ell=\ell$ with $\ell>0$ gives $\ell=1$.  Finally $3/(2G)=3\iff G=\tfrac12$.
\end{proof}

The three-dimensional class has one more realisation, internal to the arithmetic: the binary quadratic forms of
$\Zphi$, a space of dimension three on which the ratios of the operator act by boosts (\cref{prop:binary-forms}).

\begin{remark}[The three forms of the conductor, and the rotor]\label{rmk:three-forms-rotor}
The binary forms of the three faces have discriminants $-4$ (Gauss, $x^2+y^2$), $-3$
(Eisenstein) and $5$ (the norm form of $\Zphi$); the two imaginary ones are negative and
the golden one positive, and $|{-4}|\cdot5=20$ is the conductor
(\LeanInline{trichotomy\_locates\_reality}, \LeanInline{signs\_separate\_the\_faces},
\LeanInline{identity\_disc\_is\_positive} \tier{P}).  The Euler rotor of
\cite{gonzalez26mersenne}, $s\mapsto e^{i\pi s}$, restricted to the integers is the sign
character, $k\mapsto(-1)^k$, with imaginary part vanishing on the integers; at
$s=\tfrac12$ it is $i$, the self-dual point of the torus, where $t\mapsto1/t$ fixes $t=1$
(\LeanInline{rotor\_on\_integers}, \LeanInline{rotor\_imaginary\_vanishes\_on\_integers},
\LeanInline{rotor\_at\_half}, \LeanInline{rotor\_half\_is\_the\_selfdual\_point} \tier{P}).  The
vanishing point of that paper sits, in this chain, at the self-dual point.
\end{remark}

\begin{proposition}[The regulator measures the chain]\label{prop:regulator-chain}
$R_K=\log\golden$ is the same invariant at every step: the circumference of the second
dimension, $T_K=2\pi R_K$; the size of the plane over the line,
$L(1,\chiv)=2R_K/\sqrt5$; the information of the projection,
$S_{\mathrm{BH}}/k_{\mathrm B}=\lambda_{\log}R_K=\log2$; and the rapidity of the boost
in the third dimension, $\golden^2=e^{2R_K}$.
\LeanInline{F1v2.regulator\_measures\_the\_chain} \tier{P}; \LeanInline{F1CW.period\_eq\_two\_pi\_regulator},
\LeanInline{anchor\_is\_a\_crossed\_identity}, \LeanInline{F1CW.lambda\_log\_mul\_regulator}
\tier{P}\ given the class-number formula
\end{proposition}

\begin{proof}
The period is $2\pi\log\golden$ by definition of the torus; the $L$-value is the
class-number formula with $h_K=1$, $R_K=\log\golden$, $w_K=2$, $\Delta_K=5$;
$\lambda_{\log}\log\golden=\log2$ by the definition of $\lambda_{\log}$; and
$e^{2\log\golden}=\golden^2$.
\end{proof}

\subsubsection{The structure as the base of the operator}\label{ssec:operator-by-dimension}

The third dimension is the structure of the operator, as the forms of
\cref{thm:operator-forms} show and as \cite{gonzalez25hp} built it, and the forms sort by dimension.  On the line
are the tower Hamiltonian $\operatorname{diag}(m_0\golden^i)$, whose eigenvalues and
ratios live on the exponent lattice, and the ordinates $T^*(n)$, a real sequence.  On
the plane are the triad $\tfrac12\zeta_3^k$, the torus whose deck group $\mu_6$ is the
unit group of the Eisenstein lattice, and the diagonal unitary $\operatorname{diag}(i,\zeta_3)$, an isometry between the
Gauss and the Eisenstein lattices.  In three dimensions are the three eigenvalue
directions $P,C,F$ of $\OmH$, permuted by $S_3$, on which the ratios
$\golden^{\,i-j}$ act.  The operator $\Oone$ carries the three as its modulus, its phase and its ratios
(\cref{thm:one-operator,thm:one-two-three}); among its forms, one carries the three in a single formula, the operator of
\cite{gonzalez25hp}.

\begin{proposition}[The operator of the Primitive Structure carries the three dimensions]\label{prop:hp-three-dimensions}
Written in $\golden$ alone,
\[
  T^*(n)=\Bigl(2+\frac{\golden^{-\lng}\,(3-\tfrac1{20})}{1+\golden^{-3\lng}\log n}\Bigr)
  \,\frac{5\arccos(\golden/2)\;n}{\log\bigl(\golden^{-\lng}(n+3)\bigr)} ,
\]
and its factors are the three dimensions of the chain: the radial factor
$\golden^{-\lng}=\tfrac12$, with its cube $\beta_0$, is the tower, the line; the angular
factor $5\arccos(\golden/2)\,n=\pi n$ is the circle, the plane; and the ternary $d=3$
enters inside the logarithm, $\mu n+\sigma=\golden^{-\lng}(n+3)$, and in $\sigma=3\mu$,
the three directions of $E^3$.  The conductor enters in
$\alpha_0=(dN-1)/(|\ker G|\,N)$ with $N=20$.
\LeanInline{F1v2.hp\_operator\_in\_phi\_alone},
\LeanInline{F1v2.hp\_operator\_carries\_three\_dimensions} \tier{P};
\LeanInline{Tstar\_rewritten\_in\_phi}, \LeanInline{Tstar\_carries\_both\_faces},
\LeanInline{mu\_n\_plus\_sigma\_factors}, \LeanInline{eleven\_components\_anchored} \tier{P}%
\end{proposition}

\begin{proof}
The rewriting of the four parameters in $\golden$ is exact
(\cref{thm:operator-forms}\,(vii)); $\pi=5\arccos(\golden/2)$ is the pentagon
(\cref{thm:base-phi}); and $\tfrac12n+\tfrac32=\tfrac12(n+3)$.
\end{proof}

This is where the Primitive Structure and the Crystalline Worldsheet meet.  The
operator that \cite{gonzalez25hp} built to follow the zeros already carries, once it
is written in $\golden$, the line, the plane and $E^3$.  \Cref{ssec:dimensional-chain} proves the
chain of those three dimensions and the numbers that recur along it; what the operator of the Primitive Structure carries in its formula, the construction of
\cite{gonzalez26cw} realises as a worldsheet and a bulk.

% ---------------------------------------------------------------------

\subsubsection{One operator for strings and holography}\label{ssec:one-operator}

This structure is what \cite{gonzalez26cw} uses to attempt a bridge between strings and
holography with a single operator.  The chain places the Euler product, the worldsheet
and a three-dimensional bulk in one sequence generated by $\golden$; the subsections
below record how the construction of \cite{gonzalez26cw} reads each of them.

\begin{remark}[Why the chain matters, and where the physics goes]\label{rmk:one-two-three}
Each class is where the odd zeta values and the primes meet physics.  On the line, the
Euler product is the partition function of a free gas labelled by the primes
\cite{julia90}, the form in which automorphic $L$-functions appear in the
Wheeler--DeWitt wavefunctions near a black-hole singularity
\cite{hartnollyang25,declerck25}: primes and black holes.  On the plane lives the
worldsheet: the odd values are the coefficients of the $\alpha'$-expansion of the
Veneziano amplitude \cite{veneziano68}, whose coefficients $a_{2m-1,0}=g^2\zeta(2m+1)$ are exactly the ones the algebraic analysis of the bootstrap of \cite{elvang26} leaves unfixed, because the odd values are, conjecturally, algebraically independent of the other zeta values (positivity then bounds them tightly). In this paper they are not free: \cref{thm:odd-zeta} fixes them through the pentagonal arithmetic, though without evaluating them, and the reading of that difference as the effective theory not reaching the torus while the two-dimensional worldsheet fixes them is that of \cite{gonzalez26cw}; the predecessor version \cite{gonzalez26oddzeta} read it through Lawvere's
theorem \cite{lawvere69}: a system with the diagonal sees the odd values undetermined, as a G\"odel sentence is
undecidable inside its system, while the multiplicative fragment, where addition is destroyed and the diagonal
cannot be formed (\cref{thm:no-diagonal}), sees them determined.  They are also the periods that Feynman
path integrals produce, as in the three-loop coefficient of the electron's anomalous
moment \cite{laporta96}.  In the third dimension the three eigenvalue directions at
modulus $\tfrac12$ give $\ell=1$, $G_N=\tfrac12$ and $c=3$, the Brown--Henneaux central
charge of three-dimensional gravity \cite{brownhenneaux86}
(\cref{prop:chain-physics}), and the binary forms of $\Zphi$ carry the isometry group
$SO(2,1)\cong PSL_2(\R)$ of three-dimensional anti-de~Sitter space, on which the
operator's ratios act as boosts of rapidity $2R_K$ (\cref{prop:binary-forms,prop:regulator-chain}).  The chain therefore places the
Euler product, the worldsheet and a three-dimensional bulk in one chain generated by
$\golden$, with the operator's ratios acting on the third dimension: a candidate
bridge between strings and holography.  What is proved here is the chain and the
equivalences it rests on; the physical development --- the bulk, the boundary theory
and their dictionary --- belongs to \cite{gonzalez26cw}.
\end{remark}

\subsubsection{Holography from one microstate}\label{sec:adscft}

In \cite[\S3.4]{gonzalez26cw} the AdS/CFT correspondence is obtained from that
single state, level by level rather than on average.  The throat is the $AdS_5$
geometry whose curvature invariants are fixed by the modulus, the scalar curvature $R=-20$,
$\Lambda_5=-6$, $G_{AB}=6g_{AB}$ (\LeanCW{R\_scalar\_pcf},
\LeanCW{lambda5\_pcf}, \LeanCW{G\_AB\_pcf} \tier{P}); with $\ell=1$ and
$\GN=\tfrac12$ the Brown--Henneaux charge is $c=3$
(\LeanInline{F1CW.hinge\_central\_charge} \tier{P}), the arity of
\cref{sec:ring}.

The per-level statement is the one that matters structurally.  The bulk
Hamiltonian and the boundary modular generator are exponentials of the same
dilation generator with the same rate $\log\golden$, so their ratio does not
depend on the level (\LeanCW{tower\_e\_eq\_exp\_regulator},
\LeanCW{modular\_bulk\_ratio\_level\_independent} \tier{P}); an isometry acting
level by level therefore intertwines them exactly
(\LeanCW{bulk\_boundary\_intertwine},
\LeanInline{isometry\_with\_common\_modulus₅} \tier{P}).  Because no level is privileged,
the tower is an index over one flow and not a stack of different structures
(\LeanCW{no\_privileged\_level} \tier{P}).

On the de Sitter side the same modulus fixes the Gibbons--Hawking temperature and
the first law (\LeanCW{ds\_temperature\_pos},
\LeanCW{GravitySectorPCF.ds\_first\_law}, \LeanCW{komar\_first\_law}
\tier{P}), and the modular generator of the static patch is the horizon entropy and
the tower entropy at once (\LeanCW{modular\_generator\_is\_entropy},
\LeanCW{modular\_generator\_is\_tower} \tier{P}).  The duality web --- S-duality
as the Galois conjugation $\golden\mapsto-1/\golden$, T-duality as the conjugate
pair of lengths, the Kaluza--Klein bridge --- is a projection of the same tower,
and the shared signature of its two corners is the list
$|\Omega|=\tfrac12$, $\ell=1$, $\GN=\tfrac12$, $c=3$
(\LeanInline{apex\_colimit\_universal} \tier{P}).

The de Sitter companion develops three points of this holographic reading that this paper uses only through the
isometry $\mathcal V$.  First, the bulk--boundary map read as a quantum error-correcting code: an isometry $\mathcal V$ with
$\mathcal V^\dagger\mathcal V=1$ is an encoding and $P=\mathcal V\mathcal V^\dagger$ the projection onto its code subspace~\cite{almheiri15,pastawski15},
and for the triad of \cref{ssec:why-phi} that subspace has dimension one, a state and not a space of logical
states.  Second, the tensor and the string: a network on a regular tiling cut at a level has as many open legs as
edges the boundary of the level crosses, so its entropy is the area law only on average and only for a bond
dimension chosen to that end; in the throat of that paper the cut is the closed string itself, the ring of each
level, with weight $\tfrac12$ per unit of length.  The tiling of the pentagon code of~\cite{pastawski15}, $\{5,4\}$,
comes closest, with edge $\cosh e=\golden$ and layers $5F_{2k}$ growing by $\golden^2$, but a count of tiles grows
from level to level by a fraction and $\golden$ is not one.  Third, the passage from the plane to the third
direction: seen along its diagonal, the cube $\Z^3$ projects onto the hexagonal lattice $\Z[\zeta_3]$ of the
eigenvalue triangle with the diagonal as kernel, whereas the golden coupling $z=\golden y$ carries the third
direction without one.  None of these points is used in the proofs of this paper.

\subsubsection{The tower of information, and the two lengths}\label{sec:tower-info}

The tower has an information reading that connects it to \crefrange{sec:choice}{sec:clock}.  Saturation
entropy at level $\sigma$ is $S(\sigma)=\pi\golden^{\sigma}$, the mode count is
$N(\sigma)=\lfloor\pi\golden^\sigma\rfloor$ --- $3,5,8,13,21,34,56$ for
$\sigma=0,\dots,6$ (\LeanCW{tower\_start} \tier{P}) --- and the recurrence is
one multiplication by $\golden$ per level
(\LeanCW{S\_tower\_recurrence} \tier{P}).  The energy per bit is constant along
the tower, $\varepsilon(\sigma)/S(\sigma)=\epsz/\pi=1/\Mpcf$
(\LeanCW{tower\_first\_law} \tier{P}), which is the first law in this setting
and the reason the scale relation $\epsz\Mpcf=\pi$ is called a certainty
principle (\LeanCW{certainty\_principle} \tier{P}).

Two lengths run in opposite directions: the throat coordinate
$z(\sigma)=\golden^{\sigma}$ grows and the weld time
$\tau(\sigma)=\Mpcf\golden^{-\sigma}$ shrinks, and their product is constant,
$z(\sigma)\tau(\sigma)=\Mpcf$ (\LeanCW{lattice\_radius\_conjugate} \tier{P}).
That is $R\cdot(\alpha'/R)=\alpha'$ with $\alpha'=\Mpcf$: short and long distance
are one family exchanged by duality, which is the construction's reason for
having no ultraviolet regime to remove.

\subsubsection{Primes, black holes and amplitudes}\label{sec:primon}

Three points of \cite{gonzalez26cw} bear directly on this paper.

\emph{The entropy bridge.}  With $\lng=\log2/\log\golden$,
\[
  \frac{S_{\mathrm{BH}}}{k_B}=\lng\,R_K=\lng\cdot\frac{\sqrt5}{2}L(1,\chiv)=\log2
  \qquad(\LeanInline{branches\_meet\_at\_log\_two}\ \tier{P}),
\]
so the holographic bit is the regulator of $\Q(\sqrt5)$ measured in the unit
$\lng$; this is \eqref{eq:log-two} again, and it is the arithmetic branch of the
span of \cref{prop:span}.  The inputs are $h_K=1$, proved from the Minkowski
bound $\sqrt5/2<2$ (\LeanInline{F1CW.minkowski\_bound\_lt\_two},
\LeanInline{F1CW.class\_number\_one\_K} \tier{P}), and the class-number formula, which
enters as a classical input.

\emph{The Regge tower is the Euler product.}  The poles of the Veneziano
amplitude sit at the integers, every level is populated and carries spin
$\le n-1$ (\LeanCW{regge\_residue\_ne\_zero},
\LeanCW{regge\_residue\_degree} \tier{P}), and unique factorisation turns the
spectral Dirichlet series into the Euler product
(\LeanInline{F1CW.dirichlet\_is\_euler} \tier{P}).  The product over all
places is the completed zeta (\LeanCW{framework\_partition\_eq\_completed\_zeta},
\LeanCW{places\_assembly} \tier{P}), the archimedean place being fixed by the
self-dual Gaussian, which is normalised, equal to its own Fourier transform, and
of the same width in both --- three independent conditions all giving $a=\pi$
(\LeanCW{selfdual\_gaussian\_unique} \tier{P}).  That completed object is the
one whose archimedean factor supplies the weight $W$ of \cref{thm:explicit}, and its two factors
are the two sides of \cref{prop:positions,conj:closure}: the archimedean factor, fixed by the
self-dual Gaussian, carries the envelope, and the Euler product, which is the Regge tower
\cite{gonzalez26cw}, carries the residual of the primes that the law has to predict.  The predecessor
version \cite{gonzalez26oddzeta} read this tower in two further ways, which \cite{gonzalez26cw} does not take
up: graded by $\chiv$ through \eqref{eq:local-K}, a split prime contributes two single-particle modes,
$(1-p^{-s})^{-2}$, and an inert prime one mode at doubled energy, $(1-p^{-2s})^{-1}$; and the discrete Fourier
transform of $\{n^{-s}\}$ over $\Z/5$ of \cref{rmk:nature-L} is the Fourier analysis of the tower, whose
$\chiv$-projection is $L(s,\chiv)$.  The
window theorems of \cref{ssec:weil-route} are computations with the same
$\Gamma$-factor this construction fixes.

\emph{The individual and the ensemble.}  The replacement of the statistical statement
about zero spacings by a concrete pair, and what it leaves to the ensemble, is
\cref{prop:individual-ensemble}.

\emph{What the predecessor version derived in physics, and where it is now.}  The predecessor
\cite{gonzalez26oddzeta} derived, in its section on applications, the invariants of the torus at
$\tau_{\PCF}=i$.  They are developed in the two companion papers, and this paper uses none of them: the
worldline of modulus $\tfrac12$ \cite[\S3.1]{gonzalez26cw}, which is also the state at every level of
\cite[\S4]{gonzalez26ds}; the one bit of binary entropy at that modulus and the area law
\cite[\S3.2]{gonzalez26cw}; the one-loop partition function $Z_{\PCF}(i)$ \cite[\S\S2.9, 3.3]{gonzalez26cw};
and the AdS radius, Newton's constant and the Brown--Henneaux central charge \cite[\S3.4]{gonzalez26cw},
whose arithmetic is \cref{prop:chain-physics}.  Its entropy bridge is the one above.  Its common origin of
$|\eta(i)|$ and $\log\golden$ at $\tau=i$ rested on their algebraic independence, which is not proved, and is
not carried over (\cref{app:superseded}).

\subsubsection{Degrees of freedom: the three faces}\label{sec:dof}

\cite[\S5]{gonzalez26cw} reads one datum in three registers --- amplitude,
particle, fluid --- through the positive Grassmannian: a point of
$\mathrm{Gr}(k,n)$ is a projector, its trace is the rank, its normalisation
$\rho=P/k$ is a density matrix, and total positivity is a separate condition, not
implied by $\rho\ge0$ (\LeanCW{four\_faces\_one\_datum},
\LeanCW{rho\_is\_state}, \LeanCW{two\_positivities\_distinct} \tier{P});
none of the four readings depends on the frame chosen to represent the point
(\LeanCW{projector\_frame\_invariant} \tier{P}).  The seed of the ladder is the
$A_2$ lattice, with exactly six minimal vectors and Gram determinant $3$
(\LeanCW{a2\_minimal\_vectors\_exactly\_six}, \LeanCW{a2\_gram\_det}
\tier{P}), and the Frenkel--Kac--Segal count
$\dim\mathfrak g=\text{kissing}+\text{rank}$ holds at every rung
$A_2\to D_4\to E_6\to E_8$ (\LeanCW{ladder\_dim\_eq\_kissing\_plus\_rank}
\tier{P}).  On the particle face the Weinberg angle at unification is a ratio of
mode counts, $N(0)/N(2)=3/8$ (\LeanCW{weinberg\_angle\_gut} \tier{P}), and the
one-loop coefficients of the assumed content are the MSSM ones
(\LeanCW{beta\_is\_mssm} \tier{P}).

\subsubsection{What the Crystalline Worldsheet defers to this paper, and scope}\label{sec:companions}

The corpus divides as the companion paper states \cite[\S1.9]{gonzalez26cw}.  The
Crystalline Worldsheet builds the string construction and proves, for its own use,
the pentagonal identity $\golden=2\cos(\pi/5)$, the golden monoid with its
Frobenius lifts, the Euler colimit, the completed zeta with its functional
equation, the odd zeta values and the entropy bridge; it defers to the present
paper the $\Fone$ programme proper and the span aimed at the Riemann hypothesis.
We follow the same principle in reverse: what this paper's derivations use is
proved here, results shared with \cite{gonzalez26cw} are cited for
co-attribution rather than as authority, and the physical content described in
\cref{sec:disc-cw} is used nowhere.

Relative to \cite{gonzalez25hp} the account is given in
\cref{ssec:disc-hp-ledger}: Hecke discharged, the span constructed, membership derived,
the $\Lambda$-ring claim corrected, and the bridge left as one identification.
What is new here relative to \cite{gonzalez26cw}: the passage from the
$\Lambda$-monoid to the Artin $\Lambda$-structure on $\Zphi$ with its obstruction
at $5$ (\cref{sec:lambda-ring}); the separation of the two registers by the two
factors of the conductor (\cref{ssec:two-spectra}); $\chiv$ derived from the
Frobenius (\cref{thm:chi5-from-frobenius}); GRH$(\chiv)$ as one open statement in equivalent forms (\cref{thm:reduction,ssec:core-forms}); the Weil route with its three
window theorems (\cref{ssec:weil-route}); and the closing theorem with the closure conjecture
(\cref{thm:weil-closure,conj:closure}).

\phantomsection\label{sec:cw-scope}
Nothing in this section is used by \cref{sec:methods,sec:results}.  Its role is
to record where the objects came from: the torus and its self-dual point, the
modulus $\tfrac12$, the tower $\golden^{\sigma}$ and the anchor
$\golden^{\lng}=2$ were built for a question about the de Sitter observer, and
the arithmetic of \crefrange{sec:choice}{ssec:base-phi} was extracted from them afterwards.  The physical
consequences of the construction are developed in \cite{gonzalez26cw} and are not
discussed here.


% ---------------------------------------------------------------------
\subsection{Limitations and future work}\label{sec:limits}

\subsubsection{Limitations}\label{sec:open}

\Cref{ssec:achieved} separates what is proved from the one step that is open, and
the route of closure ends in one conjecture (\cref{conj:closure}), whose face on Weil's side, (W4),
is \textup{(W3)} on the planes of the orbit (\cref{prop:W4-from-W3}), is certified on the first powers of $\Frob_2$, and is not proved for all of them.  What follows limits the evidence for what
is proved; the list is derived from \cref{app:ledger}.

\begin{enumerate}
  \item \emph{Compilation.}  Twenty-seven modules of this version are compiled (\cref{app:lean}); the other declarations written for it are
    written against the other blocks of the unified file, \texttt{F1\_cw\_equivalents} among them (\cref{app:lean}), with no \texttt{sorry}
    or \texttt{axiom}, and compiled in the development.  They include the Artin $\Lambda$-structure
    and the obstruction at $5$, the closure of $\chi_5(p)\equiv F_p$ for every $p$, the
    arrow $r=\golden^k$, the first window, the failure of Soul\'e's condition, and the
    discharges of \S O of the block \texttt{F1\_v2\_formalizations}, and the closure modules of
    \cref{app:lean}: the closing theorem (\cref{thm:weil-closure}), the corollary of the conjecture (\cref{cor:closure}) and the two ways of \cref{prop:phase-modulus}; the eccentricity theorems of \cref{rmk:pappus}; and the module F1PaperPadic: the $\delta$-ring and the polar form of $\golden$, the classes and the angular
    part, the Mersenne numbers, the spectrum of $\OmH$ at $2$, and
    \cref{prop:what-fluctuation-takes,prop:euler-tower}.  The way of
    Klein for $\chiv$ rests on \texttt{E\_iff\_no\_shared\_ordinate\_chi5}, which is itself
    compiled in the development.
  Five steps of the text remain proved only in the text after the formalisations of this
    version: the second statement of \cref{prop:fibres-levels}(ii), two computations with Gauss
    sums in the proof of \cref{prop:index-degeneracy}, the gluing of the local models in the
    remark after it, and the connectedness step of \cref{rmk:G-motivic}.
  \item \emph{Tallies.}  Several classifications of the text are recorded in Lean as
    counts over enumerated types (marked $^\circ$); they certify the count, not the
    classification.
  \item \emph{Hypotheses of compiled lemmas.}  Some compiled lemmas take their inputs as
    hypotheses: numerical constants (cut-offs, margins, the two values of the
    sandwich), the value of $L(1,\chiv)$, and the identification of the operator's
    components.  Where the input is provable it is now proved and the lemma
    instantiated; the rest carry their tier.
  \item \emph{Certificates.}  The window $L\le1$ is certified in interval arithmetic; the operator windows up to $L=1.3$ are validated, with every rounding error bounded and the eigenvalue bound verified by Cholesky (\cref{rmk:rigour}); the windows of the golden field up to length $2.3$ are not yet validated: their
    Lehmann matrices are computed in double precision (\cref{ssec:golden-field}).
  \item \emph{The first window} is formalised with two declared inputs, Gauss's
    integral\ and the values of the potential \tier{N}\ (\cref{thm:first-window}).
  \item \emph{The counting density.}  The function
    \texttt{kappa\_K} of Block~1 of the unified file is not the archimedean density $\kappa_K$ of \cref{thm:pcf-counting}, and is not used.
  \item \emph{The measurements of the closure.}  The positions from the structure, the
    resolution the law has to reach and the need for every class are measured for $\zeta$, for $n\le3000$ and primes up to $10^6$; the orbit bound (W4) at $2208$ values, at low
    height (\cref{ssec:weil-closure}); the index of the angular order over the primes below $2\cdot10^5$, the
    use of each prime for $p\le300$ on $50$ zeros, and the coordinates of the Mersenne numbers for the
    exponents computed (\cref{prop:what-fluctuation-takes,prop:euler-tower}).
  \item \emph{What Lean checks of the conjecture.}  Lean states the form of the law of
    \cref{conj:closure} --- a residual for each class and the correction of the heights assembled from them --- and proves
    what follows from it.  That the law is derived from the structure without
    enumerating primes is the content of the conjecture, not a condition Lean checks.
  \item \emph{What the structure does not fix.}  The mate, the one element of the Klein group that places a zero,
    does not use the reality of the character, so what $G$ contributes to the zeros is their symmetry and not
    their position (\cref{rmk:mate-without-G}); and a statement that holds equally for the Davenport--Heilbronn
    sequence, which has the symmetry and zeros off the line, is a form of the open statement and not a step towards
    it (\cref{ssec:two-sides}).  The window theorems and the certificates improve the constants of the positivity;
    the open content is (W3) for every power of the Frobenius.
  \item \emph{The statistics of the zeros.}  The spacings are measured below height $330$; the exclusion of sectors of
    order two and six assumes that the sectors are independent, and an operator whose sectors are coupled is not
    excluded by them; that the fibre enters the zeros through an operator that does not commute with its symmetry is a reading.
\end{enumerate}

\subsubsection{Future work}\label{ssec:future}

\begin{enumerate}
  \item Interval-arithmetic replay of \cref{thm:second-window,thm:third-window}.
  \item The window at $\log4$, where the first term with $\chiv=+1$ enters.
  \item A bound for $\lambda_{\min}(L)$, which would say in advance how far any
    method of this kind can reach.
  \item Formalising the explicit formula for $L(s,\chiv)$ in Lean, with Gauss's integral
    for the digamma function, and certifying the potential values of
    \cref{thm:first-window} inside Lean with interval bounds.
  \item Formalising the localisation $\Rloc$ and the $\Lambda$-ring axioms as
    such.
  \item The same construction for other real quadratic fields: which conductor,
    which obstruction, which windows.
  \item Higher cyclotomic levels: $\Q(\zeta_5)$ itself rather than its real
    subfield.
  \item The pentagonal case of Kummer's problem beyond equidistribution, which holds for every order~\cite{patterson87}: the bias of the quintic Gauss sums at primes, and the coefficients of Kubota's quintic theta function~\cite{kubota69} with their bound, which are not available.  The same input would carry to the pentagonal field the method claimed
    in~\cite{openai26}, which bounds the partial sums of $\mu\chi$ as the inverse of the structure
    does (\cref{prop:multiplicative-side}), with exponent $\tfrac78$.
  \item The closure conjecture (\cref{conj:closure}): a law of the primes by classes, derived from
    the structure, predicting the residual of each class at scale $\tfrac12$ and the correction
    of the heights, at the resolution $1/T$ that \cref{ssec:weil-closure} measures.  With it the line follows (\cref{cor:closure}), and the exact heights, a correction that jumps
    by one at each zero, land every zero on the spectrum of the operator (\cref{prop:phase-modulus}).
  \item Compiling the closure modules and formalising the explicit formula on the orbit, with
    its tail, which the Weil face of the closure uses as a classical input.
  \item A validated quadrature of the Lehmann matrices, which would turn the windows of the
    golden field (\cref{ssec:golden-field}) into theorems.
  \item The golden field as a whole: the conjecture for $\zeta_{\Q(\sqrt5)}=\zeta\cdot L(s,\chiv)$, with the pole as ample class;
    and, with the Gauss and the Eisenstein lattices, for the Euler product of
    $\mathbb{Q}(i,\sqrt{-3},\sqrt5)$, of level $60$ (\cref{prop:euler-tower}).
    \item The way to the closure described in
    \cref{ssec:disc-hp-closure}: prisms on $\Zphi$ at each
    $p\neq5$, the binary exactness of the operator read
    prismatically at $2$, its analogue at every prime, the
    Breuil--Kisin prism at $5$ (\cref{rmk:prism-five}),
    the gluing on each level of the clock
    (\cref{prop:clock-cyclotomic}), and the passage from
    the clock to the size over the apex and the classes.
  \item The distribution of the index $i_p$ of the order
    of $\golden$ modulo $p$ --- a problem of Artin's type
    for $\golden$, and the second step of the law of
    \cref{conj:closure}: on each level of the clock the
    primes are fixed by the product formula
    (\cref{prop:clock-cyclotomic}), and their order by
    size is the law of the index, read at the layers $2$,
    $3$ and $5$ by the power-residue symbols of $\golden$
    (\cref{prop:phi-layers}); its error at scale $\tfrac12$
    would follow from GRH for their $L$-functions under
    the hypotheses of \cref{rmk:closure-status}, which are
    not constructed; in the form of the clock all the
    layers are one statement, the bounded energy of the
    deviation from the uniform vertex
    (\cref{rmk:closure-status}), measured up to $10^9$ in
    twenty layers and open.
  \item Whether the tower $\langle\golden\rangle$ admits a blueprint or $M_o$-scheme structure
    (\cref{ssec:landscape}, the paragraph `Position of this paper').
  \item The open problems on the odd values that the predecessor version \cite{gonzalez26oddzeta} stated,
    on the mathematical side: a closed form of $L(2k+1,\chiv)$, $k\ge1$, in $\log\golden$ and $\pi$, which
    with \cref{thm:odd-zeta} would give one for $\zeta(2k+1)$ (\cref{rmk:structural-vs-closed}); whether the
    period $T_K=2\pi\log\golden$ appears in $\zeta_K(2k+1)$ through the Kronecker limit formula or Siegel
    units, answered in the negative for the Kronecker route in \cref{rmk:TK-odd}; whether
    $\sum_{a=1}^4\chiv(a)\mathrm{Li}_s(\zeta_5^{\,a})$ is transcendental for every odd $s\ge3$, the cases $s=1$
    (Baker's theorem) and $s$ even (Lindemann's) being settled in \cref{rmk:nature-L}; whether the pentagonal structure of the
    motive of $L(s,\chiv)$ gives constraints toward Grothendieck's period conjecture \cite{grothendieck66}, for
    which $L(s,\chiv)$ is a period in the sense of Kontsevich and Zagier \cite{kontsevichzagier01}; whether
    multiple zeta values admit a pentagonal factorisation compatible with their stuffle relations; whether
    $\zeta(3)=\zeta_K(3)/L(3,\chiv)$ has an Ap\'ery-type rapidly convergent presentation through the pentagonal
    Hurwitz data; the question whether $G$ is a finite-level reflection of the motivic Galois group acting on
    the odd-degree generators is answered in the negative in \cref{rmk:G-motivic}.
  \item The open problems that the predecessor version stated on the physical side belong to
    \cite{gonzalez26cw}.  The relation between strings and holography through the Virasoro algebra that the
    last of them asked for is developed there: the central charge of the Virasoro tower and the
    Brown--Henneaux charge are one component of the same vertex \cite[\S3.5]{gonzalez26cw}.  Open, and not
    taken up there: whether the terminal spectral category and the span (\cref{def:SpecOmega,prop:span})
    extend that relation to $\Fone$-geometry and string compactification; whether the ratios
    $\zeta_K(2k+1)/L(2k+1,\chiv)$ can be recognised in specific Feynman topologies of the anomalous magnetic
    moment of the electron, where $\zeta(3)$ enters at two loops \cite{petermann58,sommerfield57} and higher odd
    values at higher loops \cite{laporta96}; and whether the pentagonal determination extends across bosonic,
    heterotic and superstring amplitudes beyond the $\mathcal N=4$ case of \cite{elvang26}.
\end{enumerate}

% ---------------------------------------------------------------------

\section{Conclusions}\label{sec:conclusions}

The paper builds a dictionary, by classification under equivalence, between four columns --- the arithmetic of a
one-dimensional vector, the geometry of the pentagon and its cyclotomic polygons, the algebra of $\Zphi$ with its
$\Lambda$-structure, and a Hermitian operator --- crossed by a law read at each prime, and it descends to geometry so
as not to refer to the zeros (\cref{sec:objective}).  What each column gives is the following.

The pentagon carries an $\Fone$-descent that can be written down completely
(\cref{thm:frobenius-congruence,thm:obstruction-five}).  On
$\Zphi$ the Frobenius lifts exist for every prime except the ramified one, they are
forced to be the identity or the conjugation according to $\chiv$, and the
obstruction at $5$ is exhibited by a nilpotent.  The character they determine is the
one that splits the Dedekind zeta of $\Q(\sqrt5)$, so the arithmetic of the tower
reaches $L(s,\chiv)$ by proved steps, and the odd zeta values are determined through
the pentagonal Euler product and the Hurwitz values at $a/5$ --- all of them, even
and odd, written in base $\golden$ from the anchor $2\log\golden$.

The same arithmetic reaches the zeros by a single thread from $\golden^2=\golden+1$,
which never takes the zeros or the Euler products as a starting point and uses of the
primes only their classes.  It produces the line as a fixed set, the Klein symmetry of
the zeros from the reality and parity of $\chiv$, and the operator $\Oone=HY$, with exact spectrum,
of which every form of the operator is a function, the ordinate operator of the Primitive Structure among
them (\cref{thm:one-operator,thm:operator-forms}).  Inside the tower an operator of the level, with its whole functional calculus, is blind to the
torsion fibre, which only an operator that does not commute with its symmetry separates; the zeros of $L(s,\chiv)$ show
neither the sectors nor the doublets a symmetry would produce, and those of $\zeta_K$ show the sector of its Galois
symmetry (\cref{prop:level-operators,prop:symmetry-spectrum}); and the same Fibonacci dynamics that the characters of the
real torus see only as statistics is exactly periodic on its torsion points, prime by prime, with the clock as
period (\cref{prop:two-faces}).  GRH$(\chiv)$ is shown equivalent to one statement, that the coordinate $|\rho-1|/|\rho|$ of each zero $\rho$ is a power of $\golden$, which the corpus
proves equivalent to the line in ten forms.  The route of closure is Weil's, transported along the chain (\cref{ssec:route-overview}), and it is
formalised to the end with one conjecture left (\cref{conj:closure}).  The structure approximates the positions of the zeros, the primes belong through their classes and the zeros enter it through their centre; what remains for the line is that each zero is a member, and the conjecture asks for that together with the exact heights (\cref{conj:closure}; the argument is in \cref{ssec:weil-closure,ssec:disc-hp-closure}).  Its law is obtained in two steps, the first proved, by the levels of the clock, and the second
written as one statement on the tower (\cref{rmk:closure-status}).  Its face on Weil's side is the Castelnuovo--Severi bound on
the Frobenius orbits, class by class (\cref{thm:weil-closure}); it is the positivity itself, on the planes of those orbits (\cref{prop:W4-from-W3}), certified on the first powers of $\Frob_2$ and not proved for all of them.  The chain
arithmetic, geometry, analysis, operator ends where it began, at the operator, whose positions the structure approximates and whose precision the law of the classes has to supply.
Integrality of the trace with the strip bound already forces the line, so one of the forms
asks only that one real number per zero be an integer.  The inputs of the Primitive Structure are thereby settled: its eight named axioms
are theorems, its hypothesis \texttt{bridge\_axiom} is carried to its last link, which the closure conjecture would give,
and the problems its triple obstruction distinguishes, the circularity of an operator built from the zeros, the
binary against the third dimension and the binary of the Wigner--Dyson classes, can each be seen inside the tower
(\cref{ssec:disc-hp-ledger}).

What the paper proves about the zeros beyond the equivalent forms is of a different kind.
The zeros return the conductor the pentagon derives (\cref{prop:envelope-exact}).  In Weil's positivity, written
with the Frobenius traces of $\Zphi$, three windows hold --- up to
$L^*=0.7447$ by an explicit estimate, formalised with two declared inputs, and up to
$L=1$ and $L=1.3$ by computer-assisted certificates --- and
past the archimedean threshold positivity is carried by $\Frob_2=\sigma$, because $2$
is inert.  That is the first place where the $\Lambda$-structure of the pentagon does
work inside an analytic statement about the zeros, and it is the direction in which the rest of the core, if it is reachable, would be reached. The Frobenius that does that work is the one the $\delta$-ring of $\Zphi$ carries to prisms; the places are already glued exactly at the apex $\tfrac12$, in the angles of the classes and in the binary magnitude (the paragraph after \cref{prop:euler-tower}), and the way to the magnitudes of the other primes is the one of \cref{ssec:disc-hp-closure}, listed in \cref{ssec:future}.  Finally, the construction
reads as a chain of dimensions --- the line, the plane, $E^3$ --- along which the same
numbers recur, placing the Euler product, the worldsheet and a three-dimensional bulk
in one arithmetic chain; its physical development belongs to the Crystalline Worldsheet.



\pagebreak
% =====================================================================
\appendix
% =====================================================================

\section{Status ledger}\label{app:ledger}

The ledger is segmented by purpose: a registry of the claims the text relies on, with the
declaration that carries each of them, the open statements of the corpus, and the classical
inputs it takes as given.  The declarations are those of
\texttt{F1\_pentagonal\_descent\_UNIFIED.lean} and of the twenty-seven modules compiled for this version
(\cref{app:lean}).
The entries written for this version are compiled in the development as well, and corrected
when they are found wrong.
The registry keeps, in the order of the text, every statement of the body that cites Lean,
labelled or in a paragraph of the text, so that no statement the text relies on is left out;
its last row is the closing step.  The two lists that follow name what is not proved, each
with the block and line of the unified file that declares it.  A ledger row that none of
that evidence settles is kept apart at the end, unclassified rather than given a class.
\emph{Kind} records the nature of the evidence (a Lean declaration, a numerical test, a
statement in the text), not its status; the four tables below carry no status column.

\subsection{Proved claims registry}

{\scriptsize\setlength{\tabcolsep}{3pt}
\begin{longtable}{>{\raggedright\arraybackslash}p{0.17\textwidth}>{\raggedright\arraybackslash}p{0.07\textwidth}>{\raggedright\arraybackslash}p{0.43\textwidth}>{\raggedright\arraybackslash}p{0.27\textwidth}}
\toprule Claim & Kind & Declaration & Note\\\midrule\endhead
$\chi_5(p)\equiv F_p\pmod p$, all $p\ne5$ & Lean & \texttt{S1Closure.chi5\_fib\_all\_primes}; \texttt{lambda\_structure\_gives\_chi5}; \texttt{FrobeniusUnico.fib\_eq\_sign} & Supersedes the tabulation over six primes (\texttt{fib\_split\_criterion}).\\
$F_{p-1}\bmod p$ & Lean & \texttt{F1Pendientes.fib\_pred\_from\_mod5} & New.\\
Pisano period $\pi(5)=20$ & Lean & \texttt{pisano\_five\_is\_twenty} & Closes a gap of the compiled file (modular step).\\
$\zeta(2k+1)=\zeta_K(2k+1)/L(2k+1,\chi_5)$, every $k\ge1$ & Lean & \texttt{F1PaperNew.zeta\_odd\_pentagonal\_determination\_unconditional} & Discharges the hypothesis $L(2k+1,\chi_5)\ne0$ of \texttt{zeta\_odd\_pentagonal\_determination} \witness{P}\ from Mathlib's non-vanishing on $\Re s\ge1$.\\
$L_{2m}<3\iff m=0$, every $m$ & Lean & \texttt{F1PaperNew.even\_lucas\_below\_three\_iff} & Extends \texttt{bound\_selects\_one}, which checks the first five values.\\
The clock as cyclotomic polynomials & Lean & \texttt{W4Isometria.RelojCiclotomico.clock\_iff\_cyclotomic\_root}, \texttt{prod\_Ncyc}, \texttt{log\_Ncyc\_bound}, \texttt{clock\_sum\_mertens} & New (\cref{prop:clock-cyclotomic}); product formula.\\
The index by layers & Lean & \texttt{W4Isometria.split\_layer\_law}, \texttt{even\_index\_indicator} & New (\cref{prop:phi-layers}); \texttt{KernelLaw} (open).\\
The index on the $n$-gon & Lean & \texttt{PasosDelTexto.hyperbola}, \texttt{hyperbola\_sqrt}, \texttt{W4Isometria.HiperbolaRelojTamano.large\_index\_clock}, \texttt{GeometriaDelPoligono.card\_large\_index}, \texttt{ReduccionAlReloj.level80\_coeff\_in} & New (\cref{prop:index-ngon}).\\
The prism at $5$ & Lean & \texttt{F1PaperSixty.bk\_delta\_E}, \texttt{bk\_special\_fiber\_is\_delta5}, \texttt{leopoldt\_unit} & New (\cref{rmk:prism-five}); arithmetic only.\\
$\sqrt X$ as the root of the energy & Lean; Stmt & \texttt{W4Isometria.symEnergy\_le}, \texttt{sqrt\_law\_of\_pythagoras}; \texttt{PythagorasLaw} (open) & New (\cref{prop:sqrt-x}).\\
Energy at every scale & Lean & \texttt{W4Isometria.LandauPlancherel.cramer\_uniform}, \texttt{CramerConteo.cramer\_K\_bounded}, \texttt{FormulaExplicita.energy\_bounded\_of\_explicit\_formula} & New (\cref{prop:energy-scale}); counting law and truncated explicit formula.\\
Amplitude of clause (ii) & Lean & \texttt{W4Isometria.pair\_defect}, \texttt{amplitude\_eq\_iff\_line} & New (\cref{prop:closure-amplitude}).\\
Second step of the closure, in the form of the clock & Lean; Stmt & \texttt{F1PaperClosureClock.ClockLayerEnergyBounded}; \texttt{family\_energy\_of\_clock\_layers} \witness{P}& New (\cref{conj:closure,rmk:closure-status}).\\
W3 on the angles of $\golden$ & Lean & \texttt{F1PaperClosureClock.clock\_angles\_W3\_W4} & New (\cref{ssec:bridge-boosts}).\\
Window $L\le1$, $Q\ge0.0445\|f\|^2$ & Num & \code{weil\_certificates/verify\_iv.py} & Interval arithmetic; supersedes the double-precision $0.0431$.\\
Windows $L\le1.1,1.2,1.3$ & Num & \code{weil\_certificates/op\_rig\_validated.py} & Validated numerics, eigenvalue verified by Cholesky; margins $0.0261$, $0.0077$, $0.0027$.\\
$G$ is the quotient of $\Zst$ by its squares: $\ker G=\{1,9\}$ & Lean & \texttt{F1PaperNew.G\_kernel\_eq\_squares}, \texttt{squares\_of\_units\_twenty} & \cref{prop:G}\\
The $n$-th power residue symbol of $\golden$ at a split $p$ is the image of its angular part in $\mu_n$ & Lean & \texttt{F1PaperNew.units\_pow\_iff}, \texttt{phi\_residue\_symbol} & \cref{prop:phi-residue}\\
The fifth-power residue symbol of $\golden$ is not a function of the class mod $20$ ($p=11$, $211$) & Lean & \texttt{F1PaperNew.phi\_residue\_not\_class\_function} & \cref{ssec:disc-binary}\\
Point counts of the Fermat quintic equal $p+1+\sum J$, all $|J|=\sqrt p$, $p\le71$ & Num & \code{fermat\_quintic\_points.py} & \cref{rmk:fermat-model}\\
The normalised triad of $\OmH$ is an isometry, $\mathcal V\colon\mathbb{C}\to\mathbb{C}^3$, $\mathcal V^\dagger\mathcal V=1$ & Lean & \texttt{F1PaperNew.triad\_isometry}; in \cite{gonzalez26cw} \texttt{bulkBoundary\_isometry}$^\dagger$ & \cref{ssec:why-phi}, \cref{thm:operator-forms}(iv).\\
The characters of $\Zst$ that factor through $G$ are the quadratic ones; the four of order $4$ take $-1$ at $9$ & Lean & \texttt{four\_fibres\_four\_characters}, \texttt{carriers\_have\_order\_four}, \texttt{PasosDelTexto.character\_at\_nine} & \cref{ssec:disc-binary}.\\
Lawvere's fixed-point theorem and the diagonal argument, for types & Lean & \texttt{F1Lawvere.lawvere\_fixed\_point}, \texttt{no\_surjection\_of\_fixedPointFree} & Cited in \cref{ssec:intro-binary}.\\
No Frobenius lift is fixed-point-free & Lean & \texttt{F1Lawvere.ringHom\_fixes\_intCast}, \texttt{family\_common\_fixed\_point}, \texttt{\ensuremath{\psi}\_fixes\_int} & \cref{ssec:why-phi}.\\
The reflection $x\mapsto1-x$: fixed point $\tfrac12$ where $2$ is invertible, none on $\Z$ or $\Zphi$ & Lean & \texttt{F1Lawvere.reflection\_fixed\_iff}, \texttt{reflection\_fixedPointFree\_int}, \texttt{reflection\_fixedPointFree\_Z\ensuremath{\omega}} & \cref{ssec:apex}.\\
Functional equation of $\Lambda(s,\chi_5)$ & Lean & \texttt{functionalEquation\_chi5\_discharged} & \\
Hecke's companion zero & Lean & \texttt{hecke\_chi5\_derived} & The axiom \texttt{hecke\_1920} of \cite{gonzalez25hp} is discharged.\\
$L(\bar s)=\overline{L(s)}$ & Lean & \texttt{hconj\_discharged} & Previously conditional on four pieces; all four proved.\\
Closure under the mate; real coefficients & Lean & \texttt{closedUnderMate\_chi5}, \texttt{realCoefficients\_chi5} & The conjugation is no longer a hypothesis.\\
Disc bound for $|\gamma|\ge2$ & Lean & \texttt{discBound\_discharged\_high} & \\
First zero at $\gamma\approx6.6485$ & Stmt & \texttt{NoLowZeros} (definition) & Classical; not proved in Lean.  Not equivalent to \Ecore.\\
\texttt{hmate}, \texttt{hcoord} & Lean & \texttt{mate\_hypothesis\_discharged}, \texttt{hcoord\_discharged} & \\
Strip bound on $r$ (link 6) & Lean & \texttt{rSq\_le\_in\_strip} & Previously a numerical observation.\\
The eight links of the operator chain & Lean & \texttt{the\_chain\_is\_now\_all\_P}$^\circ$, \texttt{eight\_links\_none\_from\_E}$^\circ$ & They describe the operator; that one speaks of a zero is the core.\\
Tower Hamiltonian Hermitian & Lean & \texttt{F1CW.tower\_ham\_is\_hermitian} & Real diagonal $m_0\varphi^i$; says nothing about zeros.\\
Traces of the Fibonacci powers & Lean & \texttt{the\_three\_pieces}; \texttt{the\_count\_is\_the\_golden\_lattice}$^\circ$ & \\
Orbit centre & Lean & \texttt{the\_orbit\_embedding\_is\_unconditional} & \\
Integrality of $t(\rho)$ $\Leftrightarrow$ \Ecore & Lean & \texttt{the\_two\_labels\_are\_one\_gap} & For $|\gamma|\ge2$.\\
Membership, given the bridge & Lean; Stmt & \texttt{the\_membership\_derived}; \texttt{five\_proved\_three\_classical}$^\circ$; \texttt{BridgeAxiom} & Supersedes \texttt{MembershipStillOpen}.\\
Identification of the two branches of the span & Stmt & \texttt{TheSpanIdentifiesThem} (marker) & The branches meet at $\log2$ (\texttt{branches\_meet\_at\_log\_two}, \witness{P}).\\
Pieces of the bridge & Lean & \texttt{branches\_meet\_at\_log\_two}, \texttt{G\_hom}, \texttt{G\_pi2\_is\_chi5}; \texttt{the\_bridge\_pieces\_formalized}$^\circ$ & None speaks of a zero; \texttt{BridgeIsCToP} is a marker.\\
Premise A $\Leftrightarrow$ \Ecore & Lean & \texttt{premise\_A\_iff\_GRH\_chi5} & \\
\texttt{GDS}/\texttt{bridge\_axiom} $\Leftrightarrow$ \Ecore & Lean & \texttt{GDS\_iff\_E} & The only substantive input.\\
Power action on zero coordinates $\Leftrightarrow$ \Ecore & Lean & \texttt{B\_iff\_E\_one\_hypothesis} & Given \texttt{NoLowZeros}.\\
Descent to $\mathbb F_1$ $\Leftrightarrow$ \Ecore & Lean & \texttt{descending\_to\_F1\_is\_E} & \\
Spectral embedding on zeros $\Rightarrow$ \Ecore & Lean & \texttt{the\_zero\_restricted\_form\_gives\_E} & The version with $\Re$ is false (\texttt{the\_embedding\_with\_Re\_is\_false}).\\
\Ecore\ $\Leftrightarrow$ no shared ordinate & Lean & \texttt{E\_iff\_no\_shared\_ordinate\_chi5} & Previously conditional on closure under the mate.\\
Minimum spacing as input & Stmt & \texttt{MinSpacing} (definition) & Its global form is refuted (\texttt{no\_uniform\_m\_of\_vanishing\_spacing}).\\
Radius in the tower & Lean & \texttt{position\_is\_interval\_meets\_discreteness}, \texttt{radius\_one\_is\_E} & $r=\varphi^k$ with the strip bound forces $k=0$.\\
$\Lambda$-monoid on $\langle\varphi\rangle$ & Lean & \texttt{F1\_monoidal\_descent\_data} & \\
Artin $\Lambda$-structure on $\mathbb Z[\varphi]$ away from $5$: the lift at each $p\ne5$ exists, is unique, and is fixed by the roots of $x^2-x-1$ modulo $p$ & Lean & \texttt{FrobeniusUnico.frobenius\_exists\_unique}, \texttt{F1Pendientes.frobenius\_congruence} & The monoid is not additive (\texttt{add\_destroyed}, \texttt{F1CW.psi\_not\_additive}).\\
Obstruction at $5$ & Lean & \texttt{no\_frobenius\_lift\_at\_five} & \\
Golden Prime classification & Lean & \texttt{G\_hom}, \texttt{F1CW.norm\_cosine\_phi} & The closed sine formula of \cite{grisales25} is a marker (\texttt{GLandsInTheTower}, open).\\
$\Zphi$ a $\delta$-ring at each $p\ne5$; polar form of $\golden$ at $2$ and $3$; colour group & Lean & \texttt{F1PaperPadic.delta\_structure\_exists}, \texttt{polar\_at\_two}, \texttt{PCFGranularity.colour\_order\_even} & \cref{prop:delta-ring,prop:polar}.\\
The classes fix the $2$-adic part of the angular part of $\golden$, except on $\ker G$ & Lean & \texttt{PasosDelTexto.inert\_two\_part}, \texttt{split\_three\_mod\_four}, \texttt{F1PaperPadic.inert\_norm\_is\_minus\_one}, \texttt{kernel\_classes\_do\_not\_fix\_it} & \cref{prop:class-angular}.\\
What the exponent of a Mersenne number fixes; rank of apparition & Lean & \texttt{F1PaperPadic.fib\_succ\_zero\_iff\_inert}, \texttt{mersenne\_mod\_sixty} & \cref{prop:mersenne-exponent}.\\
The spectrum of $\OmH$ read at the place $2$ & Lean & \texttt{F1PaperPadic.half\_is\_two\_adic\_norm\_of\_two} & \cref{prop:omega-two-adic}.\\
What the fluctuation takes from each prime & Lean & \texttt{F1PaperPadic.local\_phase\_contribution}, \texttt{angular\_order\_not\_a\_function\_of\_the\_class} & \cref{prop:what-fluctuation-takes}.\\
Euler's product, the tower of $\golden$ and the level $60$ & Lean & \texttt{F1PaperPadic.euler\_product\_at\_two}, \texttt{frequency\_is\_tower\_coordinate}, \texttt{three\_lattices\_level\_sixty} & \cref{prop:euler-tower}.\\
Pappus' parameter: the zeros classified by the eccentricity $e=2/t$ of their conic & Lean & \texttt{F1PaperClosure.eccentricity\_eq\_one\_iff}, \texttt{zero\_conic\_degenerate\_iff\_line} & \cref{rmk:pappus}.\\
Logarithmic signature of the pentagon & Lean & \texttt{F1v2.pentagon\_log\_signature\_full} & \\
Crystallographic restriction, pentagonal nodes & Lean & \texttt{trZeta\_*}, \texttt{pentagon\_nodes\_are\_golden}; \texttt{F1v2.cosineWindow\_holds}, \texttt{F1v2.pentagon\_nodes\_unconditional} & The hypotheses are proved in \texttt{F1\_v2}, \S O.\\
Units of $\mathbb Z[\varphi]$ are $\pm\varphi^k$ & Lean & \texttt{unit\_is\_a\_power\_of\_phi}; \texttt{F1v2.arrowB} & Without Dirichlet.\\
Frobenius eigenvalue on $\sqrt5$ & Lean & \texttt{F1FrobeniusEigen.frobenius\_eigenvalue\_on\_sqrt5}, \texttt{euler\_factor\_is\_char\_poly} & The Euler factor is $\det(1-tM_p)$.\\
Explicit formula for $L(s,\chi_5)$ & Lean; Stmt & \texttt{F1WeilFunctional.ExplicitFormulaChi5}, \texttt{prime\_side\_is\_frobenius\_trace} & Checked numerically to $10^{-18}$ with 40 zeros \witness{N}.\\
Weil positivity (full) & Lean & \texttt{positivity\_is\_a\_prime\_side\_bound} & A bound on the Frobenius-trace side by the archimedean one.\\
Semilocal positivity & Num & \code{weil\_window\_chi5.py} & Archimedean alone up to $L=0.9251$; with $\chi_5(2)=+1$ it fails from $L\approx0.9$.\\
First window, $L\le0.7447$ & Lean & \texttt{F1v2.first\_window\_integral} and its steps & $m(\log2)=0.0841$.\\
Windows $L\le1$ and $L\le1.3$ & Num & \texttt{weil\_certificates/} & Margins $0.0445$ (interval arithmetic) and $0.0027$ (validated numerics, verified Cholesky).\\
Envelope and fluctuation & Num & \code{envelope.py}, \code{tstar.py} & $\gamma_n=E^{-1}(n-S(\gamma_n))$, 40 zeros.\\
Integrality of the trace with $t<3$ & Lean & \texttt{trace\_integer\_plus\_bound} & Closes link (iii) without the tower.\\
$r\in\mathbb Q(\sqrt5)$, $t\in\mathbb Z\Rightarrow r=\varphi^k$ & Lean & \texttt{F1v2.arrowB\_full} & The hypothesis on the field is needed: $2+\sqrt3$ has trace $4$.\\
\texttt{UnitStructure} (M12) & Lean & \texttt{F1v2.arrowB\_full} & It omitted $r\in\mathbb Q(\sqrt5)$.\\
Klein action induced by $\chi_5$ & Lean & \texttt{both\_actions\_from\_one\_character}, \texttt{free\_orbit\_is\_inert} & The injectivity of the chain is \witness{N}.\\
Link 7: the value $2$ & Lean & \texttt{the\_apex\_is\_the\_inverse\_fibre\_size} & Through $\tfrac12=1/|\ker G|$.\\
Diagonal contraction & Lean & \texttt{diagonal\_blocked}; \texttt{F1v2.sq\_lt\_self\_on\_unit\_interval} & The categorical reading is a remark.\\
Local factorisation, four characters & Lean & \texttt{F1v2.factorisation\_natural} & \\
Failure of Soul\'e's condition & Lean & \texttt{F1v2.not\_soule} & No $N\in\mathbb Z[x]$ with $N(p)=1+\chi_5(p)$.\\
Counting function: weights, local assembly & Lean & \texttt{F1v2.pcf\_counting\_weights}, \texttt{local\_log\_deriv} & Global identity: Connes--Consani integral formula.\\
Residue of $\zeta_K$ at $s=1$ & Lean & \texttt{L1\_from\_class\_number\_formula}, \texttt{class\_number\_one\_K} & Given the class-number formula.\\
$\kappa_K$ & Lean & \texttt{F1CW.kappaK} & $\kappa_K=2u^2/(u^2-1)$; the \texttt{kappa\_K} of the first block is a different function and is not used.\\
Cocone and colimit of the apex & Lean & \texttt{modulus\_cocone\_six\_routes}, \texttt{apex\_colimit\_universal} & \\
The bridge, link by link & Lean & \texttt{bridge\_components}, \texttt{bridge\_assembled} & The last link is open; the closure conjecture would give it.\\
Reconnections of compiled lemmas & Lean & \texttt{F1PaperNew.*\_holds}, \texttt{F1v2} \S O & 21 hypotheses on abstract objects proved for the actual ones.\\
Connections of separately defined objects & Lean & \texttt{chi5\_ideals\_is\_chi5\_structural}, \texttt{conductor20\_characters\_are\_G}, \texttt{norm\_form\_factors\_through\_phi} & Every central object linked to $\varphi$.\\
Partition of an operator's error; only the parametrisation is removable & Lean & \texttt{error\_decomposition}, \texttt{no\_index\_operator\_fits\_both} & \cref{prop:error-partition}.\\
Fluctuation as a function of the torus; independence of the $\log p$; law of the fluctuation & Lean & \texttt{torusFactor\_on\_line}, \texttt{log\_primes\_independent} & \cref{prop:torus-line}.\\
Additive side bounded; multiplicativity and the inverse of the structure & Lean & \texttt{abel\_bound}, \texttt{chi5\_mul}, \texttt{inverse\_of\_the\_structure} & Davenport--Heilbronn shows the additive side alone is not enough.\\
Four fibres of $G$; squares in the kernel; field squeeze & Lean & \texttt{four\_fibre\_indicators}, \texttt{squares\_in\_kernel\_fibre} & \cref{prop:four-fibres}.\\
Finite-dimensional intersection form; Hodge index, primitive direction & Lean & \texttt{gram\_form\_nonneg}, \texttt{primitive\_positivity} & \cref{prop:intersection-form}; one negative direction (the pole) for $\zeta$.\\
Frobenius translation: rotation on the line, boost off it; the golden boost & Lean & \texttt{translation\_is\_boost\_off\_the\_line}, \texttt{golden\_boost} & \cref{prop:bridge-boosts}; an off-line quadruple is a hyperbolic plane.\\
Golden field: place weights, rank-one pole, total $=K-R$; first windows up to length $2.3$ & Lean & \texttt{place\_weights}, \texttt{pole\_form\_rank\_one}, \texttt{total\_is\_structural\_minus\_remainder} & \cref{ssec:golden-field}; $\rho\le\lambda_{N+1}$ proved, matrices in double precision.\\
The fibres of $G$ on the levels of the clock & Lean & \texttt{FibrasPorNivel}, \texttt{PasosDelTexto.palindromic\_norm} & New (\cref{prop:fibres-levels}).\\
Free Klein orbits and the inertia index & Lean & \texttt{OrbitaLibre.quadruple\_form}, \texttt{even\_part}, \texttt{odd\_part} & New (\cref{prop:free-orbit-index}); completes \cref{prop:bridge-boosts}(ii).\\
(W4) is (W3) on the planes of the orbit; half period & Lean & \texttt{W4DeW3}, \texttt{ParidadesJuntas}, \texttt{MedioPeriodo} & New (\cref{prop:W4-from-W3}); (W4) for every $k$ remains open.\\
The mate without the reality of the character & Lean & \texttt{MateSinG.closed\_under\_mate\_of\_fe\_dual} & New (\cref{rmk:mate-without-G}).\\
Davenport--Heilbronn as a test of routes & Num & \code{retro99.py}, \code{dh\_logderiv.py} & New (\cref{ssec:two-sides}).\\
What an operator of the level does not separate & Lean & \texttt{PCFGranularity}, \texttt{ObservablesDelNivel}, \texttt{FibraYReloj} & New (\cref{prop:level-operators}); brings in the content of the granularity module.\\
The two faces of the Fibonacci dynamics: periodic on the torsion, mixing on the characters & Lean & \texttt{DosCarasFibonacci}, \texttt{PeriodicasSol} & New (\cref{prop:two-faces}).\\
What a symmetry shows in a spectrum; the spacings of the zeros & Lean; Num & \texttt{EstructuraEnGUE}, \code{retro105.py} & New (\cref{prop:symmetry-spectrum}, \cref{ssec:envelope-fluctuation}).\\
\midrule\multicolumn{4}{l}{\emph{Statements of the body not listed above, with the declarations the body gives}}\\
\midrule
The constraint of the Primitive Structure, $\tfrac12$ the only positive real part of $\OmH$; $G$ a surjective homomorphism & Lean & \texttt{positive\_spectral\_half}, \texttt{G\_surj} & \cref{sec:objective}.\\
The four requirements of an operator & Lean & \texttt{four\_requirements} & \cref{ssec:four-requirements}.\\
The requirements that each candidate operator meets & Lean & \texttt{hpps\_meets\_two}, \texttt{omega\_meets\_three} & \cref{ssec:four-requirements}.\\
The parameter $\sigma=\tfrac32$ of the tripartite object as $\zeta(2)/(\pi/3)^2$ & Lean & \texttt{zeta\_two\_over\_pi\_third\_sq} & \cref{ssec:intro-binary}.\\
$\golden$ as root of $x^2=x+1$ and as diagonal of the pentagon & Lean & \texttt{phi\_sq}, \texttt{phi\_eq\_two\_cos\_pi\_fifth}, \texttt{phi\_cyclotomic\_inscription} & \cref{ssec:why-phi}.\\
$\chiv$ selects $F_p\bmod p$ for every prime & Lean & \texttt{additive\_spectrum\_governed\_by\_chi5} & \cref{ssec:why-phi}.\\
The span constructed: its two branches identified & Lean & \texttt{the\_span\_is\_the\_identification} & \cref{ssec:why-phi}.\\
$\golden=2\cos(\pi/5)$ from Chebyshev's identity for $\cos5\theta$ & Lean & \texttt{cos\_five\_mul} & \cref{sec:choice}.\\
The two levels of the conductor in the twentieth roots of unity & Lean & \texttt{zeta20\_pow\_five\_eq\_I}, \texttt{zeta20\_pow\_four\_eq\_zeta5} & \cref{sec:choice}.\\
The conductor is derived, and its four is one four & Lean & \texttt{conductor\_is\_derived}, \texttt{conductor\_is\_joint\_period}, \texttt{F1CW.binary\_period\_mod\_five}, \dots & \cref{prop:conductor}.\\
Crystallographic restriction & Lean & \texttt{trZeta\_crystallographic\_values}, \texttt{F1v2.three\_conductors\_unconditional}, \dots & \cref{prop:crystallographic}.\\
Unit ranks & Lean & \texttt{unit\_rank\_splits\_branches}, \texttt{branches\_are\_the\_unit\_groups}, \texttt{dirichlet\_ranks\_arithmetic} & \cref{prop:unit-ranks}.\\
$\Zphi=\Z[X]/(X^2-X-1)$ and its real embedding onto the tower & Lean & \texttt{Z\ensuremath{\varphi}}, \texttt{\ensuremath{\omega}\_sq}, \texttt{\ensuremath{\iota}\_pow} & \cref{sec:ring}.\\
$\Zphi$ is an integral domain & Lean & \texttt{f\ensuremath{\varphi}\_map\_two\_irreducible}, \texttt{f\ensuremath{\varphi}\_prime} & \cref{prop:domain}.\\
Self-duality under $x\mapsto-1/x$ & Lean & \texttt{F1OZ.selfdual\_iff}, \texttt{F1OZ.selfdual\_disc\_not\_square}, \texttt{phi\_quadratic} & \cref{prop:selfdual}; corrects \cite{gonzalez26oddzeta}.\\
Binet, integrally & Lean & \texttt{fib\_pow\_aux}, \texttt{phi\_pow\_binet\_integral} & \cref{prop:binet}.\\
Factorisations and parities & Lean & \texttt{quadratic\_factor\_upper}, \texttt{quadratic\_factor\_lower}, \texttt{two\_sides\_of\_the\_same\_pair}, \dots & \cref{prop:factor-parity}.\\
Multiplication in the tower read as addition of Fibonacci coordinates & Lean & \texttt{multiplicative\_reads\_as\_additive}, \texttt{carry\_in\_base\_phi} & After \cref{prop:factor-parity}.\\
The matrix of multiplication by $\omega$; its determinant is the norm $-1$ & Lean & \texttt{M\_is\_mul\_by\_phi}, \texttt{det\_M\_eq\_norm\_phi}, \texttt{znorm\_phi} & \cref{ssec:unit-norm}.\\
The prime $2$ is inert, and $\omega$ has order three in $\F_4^\times$ & Lean & \texttt{two\_is\_inert}, \texttt{f4phi\_order\_three} & \cref{ssec:unit-norm}.\\
The ring of integers & Lean & \texttt{quadratic\_ring\_is\_Z\_phi}, \texttt{the\_rule\_is\_proved}, \texttt{F1v2.arrowB\_Q} & \cref{prop:OK}.\\
The unit structure derived without Dirichlet's theorem & Lean & \texttt{unit\_structure\_is\_derivable} & After \cref{thm:units}.\\
Continued-fraction self-reference & Lean & \texttt{phi\_continued\_fraction} & \cref{lem:phi-cf}.\\
Nested-radical self-reference & Lean & \texttt{phi\_nested\_radical} & \cref{lem:phi-nr}.\\
Golden fixed-point uniqueness & Lean & \texttt{golden\_fixed\_point\_unique} & \cref{prop:golden-unique}.\\
The component norms, and their product $\tfrac12$ & Lean & \texttt{norm\_Omega\_is\_half}; \texttt{normP\_eq\_tan}, \texttt{normF\_eq\_cos} (CW) & \cref{thm:norm-Omega-half}; the norms are derived from the triangle of the triad.\\
Two coordinates on $T_\golden$ & Lean & \texttt{mersenne\_bridge}, \texttt{lambda\_is\_not\_an\_integer}, \texttt{M12F1Circle.logphi\_rat\_irrational} & \cref{prop:two-coords}.\\
Tower compatibility & Lean & \texttt{psi\_add}, \texttt{psi\_comp} & \cref{thm:tower-compat}.\\
Frobenius composition as exponent multiplication (remark) & Lean & \texttt{psi\_comp} & \cref{cor:frob-as-mul}.\\
The tower as the image of the exponents & Lean & \texttt{realise\_injective}, \texttt{realise\_strictMono}, \texttt{psi\_multiplicative\_on\_the\_monoid}, \dots & \cref{prop:tower-image}.\\
The additive law of the tower characterises its base & Lean & \texttt{adjacent\_closure\_real}, \texttt{golden\_unique\_closure} & After \cref{prop:tower-image}.\\
$\log\golden=6\sqrt3\,\pi/M_{\PCF}$, the base scale of the metric tower & Lean & \texttt{eps0\_relates\_log\_and\_geometry}, \texttt{component\_one\_denominator} & After \cref{prop:tower-image}.\\
The conjugation $\sigma(\omega)=1-\omega$, an involution & Lean & \texttt{FrobeniusUnico.Q.σ}, \texttt{FrobeniusUnico.Q.σ\_ω}, \texttt{\ensuremath{\sigma}\_comp\_\ensuremath{\sigma}} & \cref{def:sigma}.\\
$\Frob_m\circ\Frob_n=\Frob_{mn}$; the lifts commute & Lean & \texttt{\ensuremath{\psi}\_comp}, \texttt{\ensuremath{\psi}\_commute} & \cref{prop:frob-compose}.\\
The pentagonal descent datum: exactly one $\Lambda_{\Z,E}$-structure on $\Zphi$, $E$ the primes $\ne5$; none at $5$ & Lean & \texttt{FrobeniusUnico.lambda\_structure\_away\_from\_five}, \texttt{frob\_commute} & \cref{thm:descent-datum}.\\
The bridge between the two directions & Lean & \texttt{monoid\_ring\_bridge}, \texttt{monoid\_side\_is\_realise} & \cref{thm:monoid-ring-bridge}.\\
The machinery of the obstruction, and the norm sign & Lean & \texttt{hom\_ext}, \texttt{five\_eq\_zero\_dual}, \texttt{eps\_ne\_zero'}, \dots & \cref{prop:artin-machinery}.\\
The Frobenius sign is $\chiv$ & Lean & \texttt{five\_sq\_iff\_mod5}, \texttt{root\_x2mx1\_iff\_five\_sq}, \texttt{chi5\_is\_the\_mod5\_component} & \cref{thm:sign-is-chi5}; the first point at which $\chiv$ enters the chain.\\
$p\mid F_{p-\chiv(p)}$; the table for $p\le31$ & Lean & \texttt{F1PaperPadic.fib\_succ\_zero\_iff\_inert}, \texttt{delta\_table\_inert}, \texttt{delta\_table\_split}, \texttt{fib\_split\_criterion} & \cref{tab:fib-split}.\\
The two periods, and the Fibonacci identities behind them & Lean & \texttt{the\_two\_periods\_and\_the\_conductor}, \texttt{chi5\_Z\_mul}, \texttt{chi5\_depends\_only\_on\_class}, \dots & \cref{prop:two-periods}.\\
The two registers on the classes & Lean & \texttt{the\_two\_spectra\_are\_separated}, \texttt{G\_psi\_equivariant} & \cref{prop:two-registers-G}.\\
$|\ker G|=2$: the bit as the kernel & Lean & \texttt{kernel\_is\_one\_nine}, \texttt{log\_two\_is\_log\_kernel} & After \cref{prop:two-registers-G}.\\
$\Frob_p$ is the identity at split primes and $\sigma$ at inert ones & Lean & \texttt{\ensuremath{\psi}\_of\_split}, \texttt{\ensuremath{\psi}\_of\_inert} & After \cref{prop:two-registers-G}.\\
The classification needs base $20$ & Lean & \texttt{base20\_bijection}, \texttt{decoration\_matches\_G}, \texttt{off\_diagonal\_inaccessible}, \dots & \cref{prop:base-twenty}.\\
The angles as cyclotomic norms & Lean & \texttt{W4Isometria.conductor\_cyclotomic\_norms}, \texttt{F1PaperSixty.phi\_clock\_identity}, \texttt{F1PaperSixty.phi\_clock\_exact}, \dots & \cref{rmk:angles-cyclotomic}.\\
The kernel of $G$ leaves the $2$-adic angular part free; its order is the contraction factor & Lean & \texttt{F1PCFUnified.contraction\_is\_kernel} & After \cref{prop:class-angular}.\\
The golden ratio names the fifth roots of unity & Lean; Num & \texttt{PasosDelTexto.fifth\_root\_trace}, \texttt{PasosDelTexto.second\_trace}, \texttt{PasosDelTexto.pair\_from\_trace}, \dots, \code{retro78.py} & \cref{prop:phi-names-mu5}.\\
The two involutions at the primes: the class $9$ and the conjugation by the classes with $\chiv=-1$ & Lean & \texttt{RelojTres.class\_nine}, \texttt{RelojTres.classes\_that\_conjugate}, \texttt{TodoEnLean.class\_nine\_fixes\_phi} & After \cref{prop:phi-names-mu5}.\\
Ideals of norm $n$; split primes of density $\tfrac12$ & Lean & \texttt{G\_pi2\_is\_chi5} & \cref{sec:dedekind}.\\
The degree formula and the norms of the prime ideals of $\OK$ & Lean & \texttt{F1CW.degree\_formula\_K}, \texttt{F1CW.ideal\_norm\_values\_K}, \texttt{F1CW.euler\_from\_ideals\_K} & \cref{sec:dedekind}.\\
The two shapes of the local factor; $\chiv(p)=\mathrm{tr}\,M_p-1$ & Lean & \texttt{euler\_factor\_split}, \texttt{euler\_factor\_inert}, \texttt{trace\_minus\_one\_is\_chi5} & After \cref{thm:frob-eigen}.\\
Two spectral data, one below $1$ & Lean & \texttt{LadoAnalitico.spectral\_data}, \texttt{LadoAnalitico.spectral\_datum\_below\_one}, \texttt{spectral\_uniqueness} & \cref{prop:spectral-data}.\\
Functoriality of $\Psi$ & Lean & \texttt{psi\_comp} & \cref{prop:Psi-functorial}.\\
From Frobenius to Euler factors & Lean & \texttt{local\_factor\_identity} & \cref{lem:psi-to-euler}.\\
Euler product as colimit & Lean & \texttt{F1CW.coeff\_stabilises\_K}, \texttt{F1CW.colimit\_exists\_K} & \cref{prop:euler-colimit}.\\
The discriminants of the three quadratic fields of the conductor & Lean & \texttt{three\_subfield\_discs}, \texttt{disc\_K\_is\_conductor\_squared} & After \cref{def:QuadExt}.\\
Split laws of the four fields & Lean & \texttt{three\_quadratic\_subfields} & \cref{thm:dedekind-natural}.\\
Dedekind factorisation as natural isomorphism & Lean & \texttt{dedekind\_factorisation}, \texttt{local\_factor\_identity} & \cref{thm:dedekind-natural}.\\
Fundamental formula & Lean & \texttt{anchor\_is\_a\_crossed\_identity} & \cref{thm:L1}.\\
$L(1,\chiv)$ by Gauss--Lerch & Lean & \texttt{F1OZ.lerch\_arith} & After \cref{thm:L1}.\\
The local Dedekind factorisation at the three kinds of prime & Lean & \texttt{F1CW.local\_dedekind\_ramified\_K}, \texttt{F1CW.local\_dedekind\_split\_K}, \texttt{F1CW.local\_dedekind\_inert\_K} & \cref{sec:factorisation}.\\
Generalised Bernoulli values $B_{2k,\chiv}$, $k\le4$ & Num & \code{oz\_carried.py} & \cref{ssec:odd-zeta}.\\
The right-hand side of the determination built from the pentagon; the even values in $\Q\,\pi^{2k}$ & Lean & \texttt{F1CW.even\_zeta\_bernoulli}, \texttt{F1CW.bernoulli\_eval\_rational} & After \cref{thm:odd-zeta}.\\
Polylogarithmic form; distances on the circle; the conditional on algebraicity & Lean; Num & \texttt{F1OZ.fourier\_projection}, \texttt{F1OZ.normSq\_one\_sub\_zeta5}, \texttt{F1OZ.normSq\_one\_sub\_zeta5\_sq}, \texttt{F1OZ.baker\_ratio}, \texttt{PasosDelTexto.algebraic\_scale}, \code{oz\_carried.py} & \cref{rmk:nature-L}.\\
$L(2k,\chiv)$ and $\sum\chiv(a)\mathrm{Li}_{2k}(\zeta_5^a)$ transcendental & Lean & \texttt{F1OZ.even\_value\_transcendental} & \cref{rmk:nature-L}; Lindemann for $\pi$ as input.\\
$T_K=5L(1,\chiv)L(1,\chi_4\chiv)$; $T_K$ and $\zeta_K(2k+1)$ & Lean; Num & \texttt{F1OZ.period\_from\_the\_chi5\_fibre}, \texttt{degenerate\_orbits\_are\_transversals}, \code{oz\_carried.py} & \cref{rmk:TK-odd}.\\
Hurwitz representation of $L(s,\chi_5)$ & Lean & \texttt{hurwitz\_representation\_L\_chi5} & \cref{thm:hurwitz}.\\
$\sum\chiv(a)\psi(a/5)=-2\sqrt5\log\golden$ & Lean & \texttt{F1OZ.digamma\_sum} & \cref{thm:hurwitz}.\\
Local factors at $s=3$; $\zeta_K(3)$ and $L(3,\chiv)$ & Lean; Num & \texttt{F1OZ.local\_factors\_at\_three}, \code{oz\_carried.py} & \cref{ex:zeta3}.\\
The binary tower inside the golden one; $2$ is not a power of $\golden$ & Lean & \texttt{mersenne\_bridge}, \texttt{binary\_tower\_inside\_golden}, \texttt{lambda\_is\_not\_an\_integer} & \cref{prop:mersenne-bridge}.\\
$\log2=\log|\ker G|$, the kernel exactly $\{1,9\}$ & Lean & \texttt{kernel\_is\_exactly\_one\_nine} & After \cref{prop:mersenne-bridge}.\\
The anchor is not special to $\golden$; counting in base $\golden$ is counting in base $2$ rescaled & Lean & \texttt{modulus\_transports}, \texttt{half\_three\_roles}, \texttt{density\_exact} & After \cref{prop:mersenne-bridge}.\\
The denominators of $L(2,\chiv)$ and $L(4,\chiv)$ against $\pi^{2k}/\sqrt5$ & Lean & \texttt{denominators\_are\_conductor\_and\_ternary} & \cref{rmk:base-phi}.\\
The fibre over the tower is $\mu_6$ & Lean & \texttt{PCFGranularity.phiSig\_period\_six}, \texttt{PCFGranularity.phiSig\_six\_are\_distinct}, \texttt{PCFGranularity.cover\_degree\_six}, \dots & \cref{thm:fibre-mu6}.\\
One operator: simple spectrum, modulus and phase, the forms as functions & Lean & \texttt{UnOperador.spec\_injective}, \texttt{UnOperador.norm\_spec}, \texttt{UnOperador.spec\_polar}, \texttt{UnOperador.forms\_from\_spec} & \cref{thm:one-operator}.\\
The dimensional chain & Lean & \texttt{F1v2.dimensional\_chain}, \texttt{F1v2.line\_is\_the\_tower}, \texttt{F1v2.frobenius\_on\_the\_line}, \dots & \cref{thm:one-two-three}.\\
The binary forms of $\Zphi$ as a domain of $d=3$ & Lean & \texttt{F1v2.symF\_preserves\_disc}, \texttt{F1v2.discQ\_signature}, \texttt{F1v2.norm\_form\_disc\_and\_reversed}, \dots & \cref{prop:binary-forms}.\\
Exact spectrum of $\OmH$ & Lean & \texttt{spectral\_positions\_values}, \texttt{eigenvalue\_half}, \texttt{F1CW.triad\_real\_parts}, \dots & \cref{prop:omega-spectrum}.\\
The Eisenstein and deck arithmetic of the geometric form & Lean & \texttt{w\_cubed}, \texttt{w\_order\_three}, \texttt{z6\_is\_minus\_w\_squared}, \dots & \cref{prop:eisenstein-deck}.\\
The spectrum as free part times torsion; the parameters of $T^*$ in $\golden$ & Lean & \texttt{spectrum\_is\_anchor\_times\_orbit}, \texttt{omega\_order\_is\_exactly\_three}, \texttt{omegaZ\_kernel\_three}, \dots & \cref{prop:spectrum-free-torsion}.\\
The spectral pair read from $G$ & Lean & \texttt{F1PaperNew.gCardinals\_of\_G}, \texttt{F1PaperNew.clauseFour\_holds}, \texttt{three\_readings\_are\_G\_invariants}, \dots & \cref{prop:pair-from-G}.\\
Fourier on the deck; the Fibonacci dynamics & Lean & \texttt{toMu6\_bijective}, \texttt{toMu6\_additive}, \texttt{generator\_to\_generator}, \dots & \cref{prop:mu6-fourier}.\\
One space for all the Frobenius & Lean & \texttt{U\_comp}, \texttt{U\_commute}, \texttt{covariance}, \dots & \cref{prop:primon-space}.\\
The analytic inputs, discharged (table) & Lean & \texttt{zeros\_route\_is\_discharged} & \cref{sec:analytic-inputs}.\\
The companion zero is a function, with real part $1-\beta$ & Lean & \texttt{Chi5Zero.companion\_re} & \cref{sec:analytic-inputs}.\\
$s\mapsto1-s$ and $s\mapsto\bar s$ generate a Klein four-group & Lean & \texttt{klein\_group\_relations} & \cref{ssec:klein}.\\
The critical line is a fixed set & Lean & \texttt{stau\_fixed\_iff\_critical} & \cref{prop:stau}.\\
Orbit dichotomy & Lean & \texttt{orbit\_on\_line\_degenerate}, \texttt{orbit\_off\_line\_free}, \texttt{E\_iff\_no\_free\_orbit} & \cref{prop:orbits}.\\
The Klein action on the zeros, in detail & Lean & \texttt{klein\_acts\_on\_the\_zeros\_by\_two\_involutions}, \texttt{kmate\_is\_conj\_of\_companion}, \texttt{kmate\_keeps\_the\_ordinate}, \dots & \cref{prop:klein-detail}.\\
The chain from the generator to the zeros loses nothing & Lean & \texttt{descent\_loses\_nothing} & \cref{rmk:descent-info}.\\
The Apollonius coordinate: its level set $r=1$ is the critical line & Lean & \texttt{apollonius\_closes\_in\_lean}, \texttt{alpha\_zero\_iff\_r\_one} & \cref{ssec:apollonius}.\\
Bounding and quantising kept apart & Lean & \texttt{bounding\_is\_not\_quantising} & After \cref{prop:integrality-is-E}.\\
Centre plus radius & Lean & \texttt{embedding\_is\_centre\_plus\_radius}, \texttt{radius\_is\_rho\_minus\_half}, \texttt{the\_radius\_carries\_the\_whole\_ordinate}, \dots & \cref{prop:centre-radius}.\\
The Apollonius coordinate & Lean & \texttt{zeroCoord\_is\_a\_distance\_ratio}, \texttt{alpha\_form\_eq\_recentred}, \texttt{numerator\_minimal\_on\_the\_line}, \dots & \cref{prop:apollonius-coordinate}.\\
The tower seen by the discriminant and the valuation & Lean & \texttt{valPhi\_phi\_pow}, \texttt{valPhi\_is\_the\_bridge\_to\_psi}, \texttt{valPhi\_inversion\_is\_negation}, \dots & \cref{prop:tower-disc-val}.\\
The Li map and its fixed point & Lean & \texttt{li\_coordinate\_reads\_E}, \texttt{liFix\_root}, \texttt{liFix\_re\_is\_half}, \dots & \cref{prop:li-map}.\\
Confined, and what membership adds & Lean & \texttt{integrality\_of\_valuation\_is\_E}, \texttt{monomial\_alone\_gives\_E}, \texttt{witness\_T3}, \dots & \cref{rmk:confined}.\\
Between $2$ and $3$ the structural spectrum is empty & Lean & \texttt{spectrum\_jumps} & After \cref{cor:value-is-kernel}.\\
The Joukowski map: the circle to the segment between the foci, the pentagon from the tenth roots of unity & Lean & \texttt{jouk\_circle\_to\_segment}, \texttt{pentagon\_is\_joukowski}, \texttt{line\_is\_the\_unit\_circle}, \dots & \cref{prop:joukowski}.\\
The confocal pencil & Lean & \texttt{confocal\_always}, \texttt{degenerates\_iff\_one} & \cref{prop:confocal}.\\
The two pencils are one system & Lean & \texttt{apollonius\_circle\_equation}, \texttt{apollonius\_degenerates\_to\_the\_line}, \texttt{orthogonal\_pencil\_maps\_to\_hyperbola}, \dots & \cref{prop:pencils}.\\
Two integralities: of the coefficients of $L(s,\chiv)$, and of $t(\rho)$ & Lean & \texttt{chi5\_is\_integral}, \texttt{two\_integralities} & \cref{rmk:two-integralities}.\\
The operator chain, link by link (table) & Lean & \texttt{towerH} & \cref{ssec:chain}.\\
The existential of the chain constructed; $k=0$ on the line & Lean & \texttt{existential\_is\_constructed}, \texttt{k\_zero\_on\_the\_line} & \cref{ssec:chain}.\\
Integrality of the trace and the spectral ratio & Lean & \texttt{the\_witness\_is\_constructed}, \texttt{is\_a\_spectral\_ratio} & \cref{rmk:arrow-direction}.\\
The squeeze: every ternary $L$-zero has $\beta=\tfrac12$ & Lean & \texttt{rh\_squeeze}, \texttt{TernaryLZero.ofChi5Zero}, \texttt{rh\_squeeze\_derived}, \dots & \cref{prop:squeeze}.\\
The functional equation alone places no zero & Lean & \texttt{bound\_alone\_insufficient} & After \cref{prop:premiseA}.\\
The conductor returned by the recurrence and attributed to its two factors & Lean & \texttt{conductor\_attribution} & \cref{ssec:stage1}.\\
The primes meet all four fibres of $G$ & Lean & \texttt{primes\_hit\_all\_four\_fibres}, \texttt{primes\_reach\_the\_tower} & \cref{ssec:stage1}.\\
The Euler product of Stage~3 as the Dirichlet series of the tower & Lean & \texttt{F1CW.dirichlet\_is\_euler} & \cref{ssec:stage3}.\\
The centre of every orbit of zeros is $\tfrac12$; the apex receives the zeros & Lean & \texttt{zeros\_enter\_the\_cocone\_by\_the\_involution}, \texttt{centre\_descends}, \texttt{the\_apex\_receives\_from\_five\_legs\_including\_the\_zeros} & \cref{ssec:stage6}.\\
The last stage: membership as the union of the two halves & Lean & \texttt{union\_of\_halves\_is\_E} & \cref{ssec:stage7}.\\
Everything descends from $\golden$ but the verdict & Lean; Stmt & \texttt{everything\_descends\_but\_the\_verdict}, \texttt{MonomialHalfP}, \texttt{monomial\_is\_in\_image\_of\_realiseP} & \cref{rmk:left-to-close}.\\
The three regimes of the zeros: below $\gamma=7.3241$, below $70$, and the sandwich & Lean; Stmt & \texttt{optimum\_at\_three}, \texttt{covers\_first\_zero}, \texttt{DoubleCount}, \texttt{invariant\_two\_on\_measured\_range}, \dots & After \cref{prop:three-regimes}.\\
Four equalities on the whole zero set & Lean & \texttt{four\_equalities\_on\_the\_set} & After \cref{prop:three-regimes}.\\
An earlier Weil window, and the prime terms by sign & Lean & \texttt{extended\_window}, \texttt{adverse\_are\_outnumbered}, \texttt{four\_is\_adverse}, \texttt{eight\_is\_favourable} & \cref{rmk:earlier-window}.\\
The counting law: envelope and fluctuation & Stmt & \texttt{CountingLaw} & \cref{ssec:envelope-fluctuation}.\\
The zeros select the conductor & Lean & \texttt{scale\_injective\_in\_conductor}, \texttt{envelope\_splits}, \texttt{field\_part\_distinguishes} & \cref{prop:envelope-exact}.\\
The envelope on and off the line; the two separations & Lean & \texttt{off\_line\_L\_env}, \texttt{on\_line\_L\_env\_constant}, \texttt{on\_line\_no\_growth}, \dots & \cref{prop:envelope-edge}.\\
Three proved facts aligned with Weil's proof & Lean & \texttt{weil\_analogy}, \texttt{golden\_group\_exponent\_two}, \texttt{G\_genus\_zero} & \cref{ssec:weil-analogy}.\\
Split and inert as orbits of the Klein action on $\Zst$ & Lean & \texttt{klein\_acts\_on\_Z20\_by\_two\_involutions}, \texttt{fixed\_points\_are\_two\_torsion}, \texttt{free\_orbit\_eq\_order\_four}, \dots & \cref{prop:split-inert-orbits}.\\
The only operation on the coordinates of the zeros is the inversion & Lean & \texttt{the\_available\_mechanisms\_are\_of\_order\_two} & \cref{ssec:weil-route}.\\
The prime side split by the Frobenius, and the numerics of the bound & Lean; Num & \texttt{prime\_side\_splits\_by\_frobenius}, \code{weil\_window\_chi5.py} & \cref{rmk:saturation}.\\
The forms of the open statement (table) & Lean & \texttt{landing\_iff\_E}, \texttt{face\_coordinate\_iff\_GRH}, \texttt{liMap\_on\_unit\_circle\_iff}, \texttt{E\_iff\_no\_shared\_ordinate\_of\_fe\_rc}, \dots & \cref{ssec:core-forms}.\\
The forms linked without passing through \Ecore & Lean & \texttt{hasDegree\_gives\_monomial}, \texttt{GDS\_gives\_routeB} & \cref{ssec:core-forms}.\\
The step as a statement about the coordinate alone & Lean & \texttt{degenerate\_iff\_valuation\_zero}, \texttt{monomial\_iff\_one\_under\_confinement}, \texttt{step\_six\_membership\_would\_give\_the\_power}, \dots & \cref{prop:step-coordinate}.\\
The $\golden$ side of the step: from an integral trace to the line & Lean & \texttt{trace\_of\_integer\_matrix}, \texttt{integer\_trace\_gives\_SL2\_eigenvalue}, \texttt{the\_quadratic\_form\_condition\_gives\_the\_trace}, \dots & \cref{prop:phi-side}.\\
The zero descends: $r=\golden^k$ with $k=0$ by discreteness & Lean & \texttt{k\_zero\_by\_discreteness} & After \cref{prop:phi-side}.\\
The two halves & Lean & \texttt{coordinate\_is\_equivariant}, \texttt{coords\_closed\_under\_inversion}, \texttt{the\_companion\_plays\_the\_conjugate}, \dots & \cref{prop:two-halves}.\\
The line from the convergence of the character block of the Fej\'er sum & Lean; Stmt & \texttt{CharacterBlockConverges}, \texttt{F1PaperTorus.forcedLimit}, \texttt{F1PaperTorus.character\_block\_gives\_the\_line}, \texttt{F1PaperTorus.fejerKernel\_nonneg}, \dots & \cref{ssec:field-positivity}.\\
The places of $\Q(\sqrt5)$: none of norm $2$ or $3$ & Lean & \texttt{F1PaperTorus.no\_places\_of\_norm\_two\_or\_three} & \cref{ssec:golden-field}.\\
The archimedean form of $\zeta_K$, explicitly; $D\ge0$ & Lean & \texttt{F1PaperTorus.autocorrelation\_le\_energy} & \cref{ssec:golden-field}.\\
The positions, from the structure & Lean & \texttt{F1PaperClosure.heights\_from\_structure}, \texttt{F1PaperClosure.the\_height\_is\_the\_structure\_up\_to\_the\_correction}, \texttt{envelope\_splits}, \dots & \cref{prop:positions}.\\
Membership of the primes, entry of the zeros, proved & Lean & \texttt{F1PaperClosure.the\_primes\_belong}, \texttt{F1PaperClosure.the\_zeros\_belong}, \texttt{F1PaperClosure.what\_is\_discharged}, \dots & \cref{prop:membership-proved}.\\
The prime side in a window, given by the clock & Lean; Num & \texttt{CierrePorAltura.closure\_support}, \texttt{PasosDelTexto.window\_split}, \texttt{PasosDelTexto.window\_inert}, \code{retro66.py}, \code{retro67.py} & \cref{rmk:closure-window}.\\
The height fixes the phase, the isometry the modulus & Lean & \texttt{F1PaperClosure.no\_shared\_ordinate\_of\_injective\_heights}, \texttt{F1PaperClosure.the\_line\_from\_injective\_heights}, \texttt{F1PaperClosure.eigenvalue\_of\_isometry\_norm\_one}, \dots & \cref{prop:phase-modulus}.\\
The index as a degeneracy & Lean & \texttt{IndiceComoDegeneracion.count\_residue}, \texttt{IndiceComoDegeneracion.sigma\_phase}, \texttt{PasosDelTexto.fib\_five\_mul} & \cref{prop:index-degeneracy}.\\
(W4) follows from the line, given the expansion of the orbit & Lean & \texttt{W4Isometria.orbits\_match\_of\_line}, \texttt{W4Isometria.W4\_of\_orbits\_match} & After \cref{prop:euler-tower}.\\
The orbit in eight character blocks, each a Toeplitz sequence & Lean & \texttt{F1PaperTorus.crossCorrPair\_eq}, \texttt{F1PaperTorus.orbit\_is\_toeplitz} & After \cref{prop:eight-classes}.\\
RH for $\zeta$ and GRH$(\chiv)$ give the line for the zeros of $\zeta_K$ in the strip & Lean & \texttt{AlcanceZeta.dedekind\_strip\_zeros\_on\_line} & \cref{ssec:scope-zeta}.\\
The diagonal argument of Lawvere and Yanofsky on the phase & Lean & \texttt{EstructuraEnGUE.lawvere\_alpha\_no\_fixed\_point} & \cref{ssec:disc-hp-ledger}.\\
The classes that $G$ identifies, and the characters that factor through it & Lean & \texttt{the\_four\_fibres}, \texttt{Z20star\_pi2\_fiber\_is\_mod5} & \cref{ssec:disc-binary}.\\
The self-dual point: $\tau=i$ and $t=1$ & Lean & \texttt{selfdual\_fixed\_point} & \cref{sec:worldsheet}.\\
The fibre: the orientation flip has order two & Lean & \texttt{orient\_flip\_order} & \cref{sec:worldsheet}.\\
What the individual determines, and where the ensemble enters & Lean & \texttt{extreme\_case\_is\_binary}, \texttt{E\_iff\_no\_shared\_ordinate}, \texttt{repulsion\_excludes\_zero\_spacing}, \texttt{F1CW.sineKernel\_at\_nonzero\_integer} & \cref{prop:individual-ensemble}.\\
The single state $\Omega(\tau)$, of modulus $\tfrac12$ & Lean & \texttt{modulus\_Omega} & \cref{sec:microstate}.\\
$\tfrac12$ is the real part fixed by the functional equation & Lean & \texttt{functional\_equation\_fixed\_line} & \cref{sec:microstate}.\\
Galois conjugation and the apex involution are one map & Lean & \texttt{phi\_bar\_eq} & \cref{sec:microstate}.\\
The numbers of the chain in the physical reading & Lean & \texttt{F1v2.S3\_E3\_c3\_from\_the\_euler\_product}, \texttt{F1CW.s\_duality\_fixes\_i}, \texttt{F1CW.hinge\_central\_charge} & \cref{prop:chain-physics}.\\
The three forms of the conductor, and the rotor & Lean & \texttt{trichotomy\_locates\_reality}, \texttt{signs\_separate\_the\_faces}, \texttt{identity\_disc\_is\_positive}, \dots & \cref{rmk:three-forms-rotor}.\\
The regulator measures the chain & Lean & \texttt{F1v2.regulator\_measures\_the\_chain}, \texttt{F1CW.period\_eq\_two\_pi\_regulator}, \texttt{anchor\_is\_a\_crossed\_identity}, \dots & \cref{prop:regulator-chain}.\\
The operator of the Primitive Structure carries the three dimensions & Lean & \texttt{F1v2.hp\_operator\_in\_phi\_alone}, \texttt{F1v2.hp\_operator\_carries\_three\_dimensions}, \texttt{Tstar\_rewritten\_in\_phi}, \dots & \cref{prop:hp-three-dimensions}.\\
An isometry with a common modulus intertwines bulk and boundary level by level & Lean & \texttt{isometry\_with\_common\_modulus₅} & \cref{sec:adscft}.\\
$h_K=1$ from the Minkowski bound & Lean & \texttt{F1CW.minkowski\_bound\_lt\_two} & \cref{sec:primon}.\\
\midrule
The closing step: each zero is a member of the structure (\Ecore) & Stmt & \texttt{ConjectureE} & Implied by the closure conjecture (\cref{cor:closure}), which is stronger; every form of it is listed in \cref{ssec:core-forms}.\\
\bottomrule\end{longtable}}

\subsection{Open problems}

The corpus carries $26$ open statements with a name, among them the closure
conjecture (\cref{conj:closure}) and its face on Weil's side, (W4).  The statements the
ledger already carries keep their row; the rest are named here.

{\scriptsize\setlength{\tabcolsep}{3pt}
\begin{longtable}{>{\raggedright\arraybackslash}p{0.17\textwidth}>{\raggedright\arraybackslash}p{0.07\textwidth}>{\raggedright\arraybackslash}p{0.43\textwidth}>{\raggedright\arraybackslash}p{0.27\textwidth}}
\toprule Statement & Kind & Declaration & Note\\\midrule\endhead
Closure conjecture with clause (iii) formalised
& Stmt & \texttt{F1PaperClosure.ExactClosureAssembledLim},
  \texttt{CorrectionIsLimit}
& Implies \texttt{ExactClosure}
  (\texttt{exactClosure\_of\_assembledLim} \witness{P});
  the first form, with \texttt{assembledCorrection},
  is vacuous in (iii); the derivation from the
  structure is not formalised
  (\cref{rmk:closure-status}).\\
Positivity for the field of the four fibres & Stmt & \texttt{FieldFejerPositivity} & The four Fej\'er sums of $\Q(i,\sqrt5)$ are $\ge0$ for every $L$; stronger than the line of $\chiv$, which needs only its own sum.\\
The closure: positions from the structure, belonging of the primes and entry of the zeros, the conjecture of the two precisions & Lean; Stmt & \texttt{F1PaperClosure.ExactClosure}, \texttt{the\_line\_from\_W4}, \texttt{eight\_classes\_orthogonality} & \cref{thm:weil-closure}; (W4) measured in all eight classes.\\
Neither side of the structure is enough alone & Lean; Stmt & \texttt{F1PaperTorus.mertens\_positivity\_mathlib}, \texttt{F1PaperTorus.nonvanishing\_from\_one}, \texttt{F1PaperTorus.InverseStructureBound} & After \cref{prop:multiplicative-side}.\\
The intersection form computed from the primes; the Hodge index for $\chiv$ & Stmt & \texttt{F1PaperTorus.ArithmeticHodgeIndex} & After \cref{prop:intersection-form}.\\
Frobenius translations as rotations on the line and boosts off it & Lean; Stmt & \texttt{W4Isometria.golden\_torus\_fails\_W4}, \texttt{W4Isometria.normalized\_M\_satisfies\_W4}, \texttt{W4Isometria.line\_from\_bounded\_pair}, \texttt{F1PaperTorus.FrobeniusActsByRotations} & After \cref{prop:free-orbit-index}.\\
The pole as the ample direction; the Hodge index of the golden field & Stmt & \texttt{F1PaperTorus.GoldenFieldHodgeIndex}, \texttt{F1PaperTorus.RemainderDominated} & \cref{ssec:golden-field}.\\
Given encodings \code{lev n}, \code{fib n} of each zero, the fluctuation \code{S n} depends on (level, fibre), not level alone. & Stmt & \texttt{FluctuationIsFibreData} (Hyp \textbullet{} 69862) & \\
$Re ρ$ of each zero is a function of structural data σ. & Stmt & \texttt{RealPartIsStructural} (Hyp \textbullet{} 69871) & \\
$∀s, ½<Re s<1 → L(s,χ₅)≠0$ — no zeros strictly right of ½ in the strip (= RH-content for χ₅) & Stmt & \texttt{LineIsTypical} (Torus \textbullet{} 70056) & \\
$L(s,χ₅)=exp z$ in $½<Re<1$ — log exists ⟺ nonvanishing in the strip; Euler-product error 2.7 \%→0.2 \% & Stmt & \texttt{LogTransport} (Torus \textbullet{} 70127) & \\
$k₁L − L − forcedLimit → 0$; *"Lo que falta, en una de sus formas"* & Stmt & \texttt{CharacterBlockToForcedLimit} (Torus \textbullet{} 70189) & \\
Fejér sums of the zeros \code{Tendsto 0}; *" … es lo mismo que \code{CharacterBlockToForcedLimit}"* & Stmt & \texttt{FejerDecay} (Torus \textbullet{} 70254) & \\
$Z(L)=arq+pot−2·descomp+2·inertes ≥0$ and $→0⁺$; Z=0.0113 at 2L=18 & Stmt & \texttt{KleinOrbitBalance} (Torus \textbullet{} 70467) & \\
each fibre's Fejér deviation ≥0 and →0 (contains \code{FejerDecay}×3) & Stmt & \texttt{FourFibreBalance} (Torus \textbullet{} 70535) & \\
the 4 Dedekind-zeta Fejér sums →0, each part ≥0; Z\_K 0.398→0.118 & Stmt & \texttt{FieldFejerDecay} (Torus \textbullet{} 70577) & \\
for every finite family: one negative direction (arch−split) = one positive (inert), same direction, sum ≥0; cosine 1.0 & Stmt & \texttt{MirrorSignatures} (Torus \textbullet{} 70771) & \\
the 8×8 \code{M(s)} is Hermitian positive-definite for all $Re s>½$ (Lagarias for the 8 characters at once); min Re Λ'/Λ >0 at σ=0.55 & Stmt & \texttt{EightClassGlobalCriterion} (Torus \textbullet{} 71154) & \\
$∀k, |t_k| ≤ t_0$ on a place's orbit (Weil face; dual-tagged \code{} by EF) & Stmt & \texttt{OrbitCastelnuovoSeveri} (Torus \textbullet{} 71390) & \\
$∀a∈family ∀k, |O.t a k| ≤ O.t a 0$ — W4 in the χ₅ block; max 0.803 & Stmt & \texttt{W4Chi5} (Torus \textbullet{} 71538) & \\
∃ inner-product E, isometry V, vectors g\_a with $t_k(a)=⟪g_a,V^k g_a⟫$, $t_0(a)=‖g_a‖²$ & Stmt & \texttt{IsometricRealization} (Torus \textbullet{} 71771) & \\
the same realization *via the exact Hilbert–Pólya operator* (E⁻¹ on the π side + S self-adjoint); "caso de \code{IsometricRealization}" & Stmt & \texttt{OperatorRealization} (Torus \textbullet{} 71828) & \\
clause (iii): the correction is the assembly of the class residuals; *"derivation of the law … not formalised"* & Stmt & \texttt{ExactClosureAssembled} (Closure \textbullet{} 72356) & \\
$∀S∈(0,log4)$, the archimedean form is ≥0 (primes-free window). Doc's own "falta" list: *"la aritmética de intervalos de las cuadraturas y del residuo, y una cota inferior rigurosa de λ₂ (basta λ₂ ≥ 0.05)"* & Stmt & \texttt{GoldenArchimedeanWindow} (Torus \textbullet{} 70871 `[O→N]`) & \\
same window with S ≤ log 4: λ₂(D)=9.509 > |C₀|=9.135 ⇒ min Q/‖g‖² = +0.0071; without the pole Q<0 already at S≈0.45 & Stmt & \texttt{GoldenArchimedeanHodge} (Torus \textbullet{} 70899 `[N→O]`) & \\
\bottomrule\end{longtable}}

\subsection{Classical inputs}

The corpus carries $21$ classical inputs with a name, of which three --- Kronecker's
density for two primes, the non-vanishing on $\Re s\ge1$ and Mertens' positivity --- are now
proved in the file itself.  Each is named below with the block and line that declares it.

{\scriptsize\setlength{\tabcolsep}{3pt}
\begin{longtable}{>{\raggedright\arraybackslash}p{0.17\textwidth}>{\raggedright\arraybackslash}p{0.07\textwidth}>{\raggedright\arraybackslash}p{0.43\textwidth}>{\raggedright\arraybackslash}p{0.27\textwidth}}
\toprule Claim & Kind & Declaration & Source\\\midrule\endhead
The first windows over the whole function space: the low-frequency fraction and Lehmann's bounds & Lean; Stmt & \texttt{F1PaperTorus.lowFrequencyFraction\_nonneg}, \texttt{F1PaperTorus.LehmannBounds} & \cref{ssec:golden-field}.\\
$log p/log q$ irrational ⇒ ${a·log p/log q+b}$ dense in ℝ & Stmt & \texttt{KroneckerTwoPrimes} (Torus \textbullet{} 70025) & \\
equidistribution: line average of $F(flucTorus∘line)$ = torus average; KS p=0.14 & Stmt & \texttt{WeylLaw} (Torus \textbullet{} 70033) & \\
Bohr–Jessen/Kolmogorov: torus product ≠0 for σ>½ at typical points & Stmt & \texttt{TorusNonvanishing} (Torus \textbullet{} 70042) & \\
RH for L(s,χ₅) ⟺ strong recurrence along the line & Stmt & \texttt{BagchiEquivalence} (Torus \textbullet{} 70047) & \\
Dirichlet criterion ⇒ $Σχ₅(n)n^{−s} → L(s)$ for all $Re s>0$. & Stmt & \texttt{AdditiveSeriesConverges} (Torus \textbullet{} 70108) & \\
$∀s, Re s≥1 → L(s,χ₅)≠0$ (Mathlib) & Stmt & \texttt{NonvanishingFromOne} (Torus \textbullet{} 70115) & \\
$LogTransport ↔ CharacterBlockConverges k₁$ (EF with the Fejér family) & Stmt & \texttt{TransportIsTheCharacterBlock} (Torus \textbullet{} 70162) & \\
from the inverse bound: Dirichlet series of μχ₅ converges for $Re>½$ and equals $1/L$. & Stmt & \texttt{InverseRepresentation} (Torus \textbullet{} 70324) & \\
Davenport–Heilbronn 1936: an explicit period-5 sequence, bounded partial sums, conductor-5 functional equation, with a zero at ≈$0.8085+85.699i$ & Stmt & \texttt{PeriodicBoundedIsNotEnough} (Torus \textbullet{} 70390) & \\
$|ζ₅(σ)³·L(σ+it,χ₅)⁴·L(σ+2it)| ≥ 1$ for σ>1 → nonvanishing at Re=1 & Stmt & \texttt{MertensPositivity} (Torus \textbullet{} 70439) & \\
Weil criterion: $(∀L≥0: 0≤Z L) → FejerDecay Z$ — Fejér positivity detects any off-line zero; +6.3·10⁻⁴ clean / negative with a quartet & Stmt & \texttt{PositivityDetectsOffLineZeros} (Torus \textbullet{} 70596) & \\
$η<ρ → η−ε²/(ρ−η) ≤ λ₁$ (self-adjoint operator inequality) & Stmt & \texttt{KatoTemple} (Torus \textbullet{} 70918) & \\
(a) orthogonality identity + (b) Taylor remainder $|e^{iθ}−Σ_{m<N}|≤|θ|^N/N!$ ⇒ pointwise bound (c) & Stmt & \texttt{PointwiseLowFrequencyBound} (Torus \textbullet{} 71103) & \\
RH ⟺ $Re Λ'/Λ(s)>0$ for all $Re s>½$ (Hadamard) & Stmt & \texttt{LagariasCriterion} (Torus \textbullet{} 71144) & \\
Pólya–Jensen: Ξ zeros real ⟺ Jensen polynomials hyperbolic ∀d,n & Stmt & \texttt{JensenHyperbolic} (Torus \textbullet{} 71162) & \\
Griffin–Ono–Rolen–Zagier PNAS 2019 Thm 3: asymptotic ⇒ Jensen hyperbolic for large n & Stmt & \texttt{GORZTheorem} (Torus \textbullet{} 71179) & \\
explicit formula on the orbit + phase separation: a zero with $Re≠½$ dominates the orbit at infinitely many k; tag *"presupone que el módulo máximo … se alcanza"* & Stmt & \texttt{OrbitDetectsOffLine} (Torus \textbullet{} 71551) & \\
$:= OrbitDetectsOffLine$ (the classical entry of the Weil face, with the tail) & Stmt & \texttt{ExplicitFormulaOnTheOrbit} (Pending \textbullet{} 71886) & \\
half-scale residues ⇒ all zeros in the band on the line ("sexta forma del corpus, que allí sólo está como marcador") & Stmt & \texttt{ClassicalHalfScale} (Closure \textbullet{} 72089) & \\
the closed form $(2/π)(RL)^{2N+1}/((2N+1)²(N!)²)$; doc: *"el enunciado analítico … es clásico y queda con nombre"*; nonneg already proved (\code{lowFrequencyFraction\_nonneg}, 70980) & Lean & \texttt{LowFrequencyFraction} (Torus \textbullet{} 70977) & \\
\bottomrule\end{longtable}}

\subsection{Ledger rows whose class is not settled}

These rows of the ledger carry no declaration of their own, and none of the evidence
above settles whether the statement they record is proved, open or classical; they are
recorded here rather than given a class.

{\scriptsize\setlength{\tabcolsep}{3pt}
\begin{longtable}{>{\raggedright\arraybackslash}p{0.17\textwidth}>{\raggedright\arraybackslash}p{0.07\textwidth}>{\raggedright\arraybackslash}p{0.43\textwidth}>{\raggedright\arraybackslash}p{0.27\textwidth}}
\toprule Claim & Kind & Declaration & Note\\\midrule\endhead
Prismatic route to the closure & — & --- & \cref{conj:prismatic-route}; its clauses' proved parts are listed with it.\\
Each zero, unconditionally: closure under the symmetries, orbit centred at $\tfrac12$, $r(1-\rho)=r(\rho)^{-1}$, confinement of $r$ and $t$ & — & \cref{prop:each-zero} & None uses that the zero is on the line.\\
Forms of the closing step & — & \cref{ssec:core-forms} & One step in several forms; conditions settled or classical.\\
$G$ and the motivic Galois group & — & --- & \cref{rmk:G-motivic}.\\
\bottomrule\end{longtable}}

\section{Superseded statements and corrections to earlier versions}\label{app:superseded}\label{app:corrections}

Statements of earlier versions of the programme --- of the corpus, of the predecessor
\cite{gonzalez26oddzeta} and of \cite{gonzalez25hp} --- that this version corrects or
passes, with what now carries them.  None of them is cited in this paper.

\begin{longtable}{p{0.25\textwidth}p{0.33\textwidth}p{0.33\textwidth}}
\toprule Earlier statement & What it claimed & Current status\\\midrule\endhead
\texttt{fib\_split\_criterion} & $\chi_5\equiv F_p$, by tabulation over six primes & \texttt{S1Closure.chi5\_fib\_all\_primes} (all $p\ne5$)\\
\texttt{MembershipStillOpen} & membership open & derived from the bridge (\texttt{the\_membership\_derived})\\
\texttt{heckeComp} & Hecke as a classical input & derived (\texttt{hecke\_chi5\_derived})\\
\texttt{closedUnderMate\_needs\_only\_conj\_now} & conditional on \texttt{hconj} & \texttt{closedUnderMate\_chi5}\\
\texttt{B\_iff\_E} & two hypotheses & \texttt{B\_iff\_E\_one\_hypothesis}\\
\texttt{MinSpacing} as input & measured bound & discarded (\texttt{no\_uniform\_m\_of\_vanishing\_spacing})\\
$\Lambda$-ring claim of \cite{gonzalez25hp} and of \cite{gonzalez26oddzeta} & $\Z[\varphi,\varphi^{-1},\tfrac12]$ a $\Lambda$-ring with $\Frob_p(\varphi)=\varphi^p$ & $\Lambda$-monoid on the tower and Artin $\Lambda$-structure away from $5$ (\cref{rmk:monoid-not-ring}, \cref{sec:lambda-ring})\\
Closure of the chain in \cite{gonzalez25hp} & the chain closed, no gaps & conditional on the core (\cref{thm:topological-proof})\\
\texttt{UnitStructure} & unit structure without the field hypothesis & false as written; \texttt{F1v2.arrowB\_full}\\
Uniqueness of $\golden$ in \cite{gonzalez26oddzeta} & $\golden$ the only quadratic irrationality whose polynomial is self-dual under $x\mapsto-1/x$ & false as stated: every $x^2+ax-1$, $a\ne0$, is; \cref{prop:selfdual}\\
Split law of $\Q(\sqrt{-5})$ in \cite{gonzalez26oddzeta} & $p\equiv1,9\pmod{20}$ & $p\equiv1,3,7,9\pmod{20}$ (\cref{thm:dedekind-natural})\\
Transport of \cite{gonzalez26oddzeta} & \texttt{zeta\_positive\_spectral\_half}: positive spectral values of $\zeta$ equal $\tfrac12$ & definitional: \texttt{zeta\_spectral} is defined as \texttt{PCF\_spectral}, so the statement is about $\OmH$ and says nothing about the zeros of $\zeta$\\
Axiom of \cite{gonzalez26oddzeta} & $\golden^2=\golden+1$ as the one axiom of its Lean file & a theorem, \texttt{phi\_sq}; the file has no axiom\\
Formalisation gap of \cite{gonzalez26oddzeta} & the Euler product, the Dedekind factorisation and the class-number formula not formalisable & the first two formalised with Mathlib's $L$-series (\texttt{dedekind\_factorisation}); the class-number formula\\
Irrationality in \cite{gonzalez26oddzeta} & the irrationality of the odd values traced to the transcendence of $\log\golden$ & the determination does not decide it (\cref{rmk:nature-L})\\
Common origin at $\tau=i$ in \cite{gonzalez26oddzeta} & $|\eta(i)|$ and $\log\golden$ algebraically independent & not proved; not used\\
Generator of the line in \cite{gonzalez26oddzeta} & no quadratic generator other than $\golden$ gives $\R\cong\R_{>0}$ & every base $>1$ does; what singles out $\golden$ is \cref{prop:selfdual} and the unit group (\cref{thm:units})\\
Discharge of the bridge in \cite{gonzalez26oddzeta} & the \texttt{bridge\_axiom} a theorem & the bridge proved link by link up to the last (\cref{thm:topological-proof})\\
Polylogarithmic Baker question of \cite{gonzalez26oddzeta} & transcendence of $\sum\chiv(a)\mathrm{Li}_s(\zeta_5^a)$ open for every $s\ge2$ & settled for even $s$ (\cref{rmk:nature-L}); open for odd $s\ge3$\\
Kronecker question of \cite{gonzalez26oddzeta} & whether $T_K$ appears in $\zeta_K(2k+1)$ through the Kronecker limit formula & not through that route (\cref{rmk:TK-odd})\\
Motivic question of \cite{gonzalez26oddzeta} & whether $G$ is a finite-level reflection of the motivic Galois group & no; that group is connected (\cref{rmk:G-motivic})\\
\bottomrule\end{longtable}

\section{Lean formalization audit}\label{app:lean}

\subsection{Files and roles}
The formalization is delivered as one file, \texttt{F1\_pentagonal\_descent\_UNIFIED.lean}:
twenty-seven blocks in dependency order, then the paper as a namespace tree of aliases, one for every declaration the text cites (897), all pointing at declarations of the same
file.  The same content is also a Lake project, \texttt{F1Paper}.  The blocks are of
three kinds.
\begin{description}
  \item[Compiled, unchanged] \texttt{F1\_spectra\_anchor\_G\_standalone},
    \texttt{F1\_PCF\_session\_modules}, \texttt{PCF\_sesion\_2026\_08\_18\_unificado},
    \texttt{PCF\_sesion\_M12\_unificado},
    \texttt{PCF\_granularity\_UNIFIED\_with\_mathlib\_verified},
    \texttt{everything\_descends}, \texttt{the\_anchor\_segments\_the\_circle} and
    \texttt{the\_mate\_carries\_the\_conjugate} are the compiled files of the project,
    line for line.
  \item[Compiled, with a documented edit] \texttt{F1\_compilado\_sin\_axiomas}, the
    core formalisation of this paper (its first block: $\golden$, the pentagon, $G$,
    $\chiv$, the Euler factors, the squeeze, and the eight named axioms of
    \cite{gonzalez25hp} discharged); two \texttt{private} lemmas are renamed with the
    suffix \texttt{\_priv}, which in one file would clash with public lemmas of the
    same name, and its comments cite the labels of this paper instead of the section
    numbers of the earlier version; \texttt{PCF\_unificado} (its one
    use of the Lean development of \cite{gonzalez26cw}, a spectral-uniqueness lemma,
    replaced by the corpus's own \code{spectral\_uniqueness}, and that import
    removed); \texttt{PCF\_sesion\_completa} (its two \texttt{sorry}, in
    \code{level\_two\_character\_from\_fibonacci} and
    \code{pisano\_five\_is\_twenty}, closed by \texttt{S1\_chi5\_fib} and by a
    modular-arithmetic step).
  \item[Written] \texttt{S1\_chi5\_fib},
    \texttt{F1\_pendientes}, \texttt{F1\_pendientes\_2}, \texttt{F1\_frobenius\_eigenvalue},
    \texttt{F1\_weil\_functional}, \texttt{F1\_primon\_space}, \texttt{D16\_hconj},
    \texttt{F1\_cw\_equivalents} (the facts of \cite{gonzalez26cw} that \cref{sec:background,sec:methods,sec:results} use;
    its \S4 is taken verbatim from that paper's compiled development, and its declarations are
    cited \witness{P}),
    \texttt{F1\_v2\_formalizations}, and the modules of the paper \texttt{F1PaperNew},
    \texttt{F1PaperHyp}, \texttt{F1PaperTorus}, \texttt{F1PaperPending}, \texttt{F1PaperClosure} and \texttt{F1PaperPadic}; and, for the clock, the index, the energy and the prism at $5$, \texttt{F1PaperSixty}, \texttt{F1PaperSixtyWeil}, \texttt{M12F1Circle}, \texttt{PhiCelda}, \texttt{W4Isometria} and \texttt{F1PaperClosureClock}.  Their declarations are compiled in the development; cited inside a
    group of \witness{P}\ names they carry the mark $^{*}$.
  \item[Compiled] with the
    imports it needs and with no \texttt{sorry}, \texttt{axiom} or \texttt{native\_decide}: for the window of the closure and
    the index as a degeneracy, \texttt{CierrePorAltura}, \texttt{IndiceComoDegeneracion}, \texttt{MasaDelReloj} and
    \texttt{NivelesYOrientacion}; for the inertia index and the mate, \texttt{FibrasPorNivel},
    \texttt{OrbitaLibreEIndice}, \texttt{ParidadesJuntas}, \texttt{MedioPeriodo}, \texttt{W4DeW3} and \texttt{MateSinG}; for
    the operators of the level and the symmetries of a spectrum, \texttt{ObservablesDelNivel}, \texttt{FibraYReloj}
    and \texttt{EstructuraEnGUE}; for the two faces of the Fibonacci dynamics, \texttt{DosCarasFibonacci} and
    \texttt{PeriodicasSol}; the steps that earlier versions proved only in the text, \texttt{PasosDelTexto}; for the diagonal argument and the reflection $x\mapsto1-x$, \texttt{F1Lawvere}; for the Frobenius lift defined by its congruence, \texttt{FrobeniusUnico}; for the one operator $\Oone=HY$, \texttt{UnOperador}; for the clock with a factor $3$, \texttt{RelojTres}; for the Euler product as a colimit, \texttt{LadoAnalitico}; for $\zeta$ at the points $a/5$, \texttt{ZetaPentagonal}; for the scope to $\zeta$, \texttt{AlcanceZeta}; for the steps of the text on the operator, \texttt{TodoEnLean}; for the results carried from the predecessor version, \texttt{OZTraidos}, and for the even values, \texttt{OZCierres} (both in the namespace \texttt{F1OZ}); and, for the components of the colour operator as the involutions of the Klein group, \texttt{ComponentesInvoluciones}.  These twenty-six modules follow the paper's block in the file, as its blocks S1--S27.  Their declarations are cited \witness{P}.  Six of them did not compile as first written and
    were corrected; the corrections changed proofs and no statement, except that
    \texttt{DosCarasFibonacci.periodic\_is\_torsion} is stated on the entries of the matrix.
\end{description}
Two files of the project are not included because their declarations are already in
other blocks: \texttt{the\_chain\_is\_now\_all\_P} (all $15$ in
\texttt{PCF\_sesion\_M12\_unificado}) and \texttt{PCF\_session\_unified} ($313$ of $314$
in \texttt{PCF\_sesion\_completa}; the remaining one, the auxiliary definition
\texttt{RealDigamma}, is cited nowhere).  The Lean development of \cite{gonzalez26cw}
is not part of the file: what \cref{sec:background,sec:methods,sec:results} use of it is in the corpus under another name
or in \texttt{F1\_cw\_equivalents} (the correspondence is \texttt{CW\_EQUIVALENCIAS.md})
,
and the declarations marked $^{\dagger}$, cited only in \cref{sec:disc-cw}, are that
paper's.

\subsection{Ontology of the corpus: connections and reconnections}\label{ssec:ontology}
The whole Lean file was read as an ontology: every declaration with its statement, its
proof and the declarations it uses.  Two
kinds of repair came out of it, both in the block \texttt{F1PaperNew} and all
compiled in the development.

\emph{Reconnections.}  Several compiled lemmas take as hypotheses, on an abstract
$\golden$ or an abstract map, facts that the corpus proves elsewhere.  Each is now
proved for the actual objects: $2\cos(\pi/5)=\golden$ under its four names, $\golden^{\lng}=2$,
$\golden^{-\lng}=\tfrac12$ under two names, $|1-1/\rho|=1\Leftrightarrow\Re\rho=\tfrac12$
under three names and, with the modulus $|\rho-1|/|\rho|$, for every $\rho$ under two
more, the pentagonal cosines, Binet's $\golden^{2k}-\golden^{-2k}=\sqrt5F_{2k}$, the
cardinals of $G$ computed from $G$, $\chi_4$ as Gauss's sequence, $\|\OmH\|=\tfrac12$
under two names, the positive spectral position under three, and the closure of the
zero set under $\rho\mapsto1-\rho$ and conjugation
\witness[F1PaperNew.pentagonBridge\_holds, F1PaperNew.clauseFour\_holds, F1PaperNew.apexIsPhiPow\_holds,
F1PaperNew.liActsOnZeros\_holds, F1PaperNew.modulusOnLine\_holds, F1PaperNew.pentagonCosines\_holds,
F1PaperNew.binetDifference\_holds, F1PaperNew.gCardinals\_of\_G, F1PaperNew.gaussSplit\_holds,
F1PaperNew.operatorNormIsApex\_holds, F1PaperNew.positiveSpectralPosition\_holds,
F1PaperNew.zIsAVSet\_chi5]{P}, and the discharges of \S O of \texttt{F1\_v2\_formalizations}.

\emph{Connections.}  Some objects are defined separately and no theorem related them to
the objects built from $\golden$.  Each relation is now stated: the $\chiv$ of the ideal
machinery is the $\chiv$ of the structural chain; on $\Zst$ the characters of conductor
$20$ are the two components of $G$; the norm form of $\Zphi$ factors as
$(x+\golden y)(x+\bar\golden y)$ with discriminant $5$; and the two definitions of the
eigenvalue triad coincide \witness[F1PaperNew.chi5\_ideals\_is\_chi5\_structural,
F1PaperNew.conductor20\_characters\_are\_G, F1PaperNew.norm\_form\_factors\_through\_phi,
F1PaperNew.omega\_eigenvalues\_are\_omega\_hat]{P}.  After them every central
object of the construction is linked to $\golden$ by a chain of at most two theorems.

\emph{Duplicates.}  The corpus names the same object in several ways: $80$ groups of
definitions, $277$ definitions in all, are the same object under different names (the
mate $\rho\mapsto1-\bar\rho$ under seven, the predicate ``$r$ is a power of $\golden$''
under fifteen, the trace $r+1/r$ under nine, the tower eigenvalues $m_0\golden^i$ under
seven, and the golden ratio itself under five).  Each group is now linked by
definitional equalities \witness[F1PaperNew.same\_def\_1]{P} and the $196$ that follow, and
the paper uses one name per object.  In the same way $294$ groups of theorems state the
same thing under different names; the paper cites one name per group.  Besides these,
$99$ definitions are markers defined as \texttt{True}; none is cited as evidence.

\emph{The closure modules.}  Five further blocks carry the analysis of the open statement
(\cref{ssec:route-overview,ssec:fluct-geometry,ssec:two-sides,ssec:field-positivity}):
\texttt{F1PaperHyp}, which connects the forms of the open statement to the corpus (the index form of the
Frobenius action, the partition of the operator's error), and \texttt{F1PaperTorus}, the
fluctuation as geometry of the torus, the two sides of the structure and the Fej\'er
positivity of $\Q(i,\sqrt5)$, with the finite-dimensional intersection form and the bridge to the chain
(\cref{ssec:product-surface,ssec:bridge-boosts,ssec:golden-field,ssec:weil-closure}), ending in the closing theorem; and
\texttt{F1PaperPending}, which contains only the Weil face of what remains, stated on the explicit
orbit --- (W4) and the explicit formula on the orbit --- with the closing theorem restated from them; and \texttt{F1PaperClosure}, the closure of the paper: the
positions from the structure, the belonging of the primes and the entry of the zeros, and the conjecture of the two precisions with its consequences; and F1PaperPadic, the $p$-adic
structure of $\golden$ (the $\delta$-ring, the polar form, the classes and the angular part), the
Mersenne numbers, and what the fluctuation takes from each prime. Together the five blocks add $159$ proved statements (their \texttt{theorem} declarations in
the unified file, counted with the comments stripped), $26$ open
statements with a name\ (among them the closure conjecture, \cref{conj:closure}, and its face on Weil's side, (W4)) and $21$ classical inputs with a name, of which three --- Kronecker's
density for two primes, the non-vanishing on $\Re s\ge1$ and Mertens' positivity --- are now proved; none duplicates a
statement of the corpus, and where the corpus has the piece (the counting law, the power
action, the functional equation, the block with character, the local factor) it is used, not
redefined.  The numerical tests are the scripts ported into \texttt{tests/}.

\subsection{Build order}
Each source file is wrapped in its own \texttt{section}, so that its \texttt{open} and
\texttt{set\_option} commands stay local; the order preserves, for every file, the files
it was compiled against.  Across the file no two declarations share a fully qualified
name, no private and public declarations clash, and no block uses a name that the merge
makes ambiguous.  Some unqualified names exist twice in different namespaces, so the
paper layer always writes fully qualified names.

\subsection{Axiom report}
The report is produced by \texttt{\#print axioms} on the compiled file.  For the
compiled blocks the expected output is \texttt{[propext, Classical.choice, Quot.sound]};
for Theorem~\ref{thm:odd-zeta} it is stated with the theorem.  The report for the
declarations written for this version is produced after their compilation.

\subsection{Numerical scripts}
This subsection is selective: it lists the scripts that back a statement of the
text, while the full census lives in \texttt{tests/} ($83$ scripts, ported into the
suite) with their committed outputs in \texttt{tests/fixtures/salidas/} ($35$ files).
\code{verify\_weil\_chi5.py} (explicit formula, $10^{-18}$),
\code{weil\_window\_chi5.py} (semilocal eigenvalues),
\code{weil\_window\_rigorous\_check.py} (constants of
\cref{thm:first-window}), \texttt{weil\_certificates/} (certificates of
\cref{thm:second-window,thm:third-window}: \code{cert.py}, \code{rigW.py},
\code{opcert\_rig.py}; the rigorous replays \code{verify\_iv.py}, in interval arithmetic, and \code{op\_rig\_validated.py}, validated with a verified eigenvalue bound), \code{zeroside\_check.py} (the minimum of the $L=1$ window
from the zero side), \code{envelope.py} and \code{tstar.py}
(\cref{prop:envelope-exact}); and the scripts of \cref{sec:results-verdict}, ported into \texttt{tests/}:
\code{testGeom.py}, \code{testB1b.py}, \code{fejer\_check.py}, \code{balance.py},
\code{fibres.py}, \code{castelnuovo.py}, \code{gram.py}, \code{bridge.py}, \code{offline.py},
\code{dh.py}, \code{mertens341.py}, \code{mobius.py}, \code{rigorous\_chain.py},
\code{rho\_interval.py}, \code{field\_signature.py}, \code{arch\_window.py},
\code{w4exact.py}, \code{zeroside.py}, \code{sweep.py}, \code{circle\_and\_detect.py},
\code{corrected\_operator.py}, \code{classes\_test.py}, \code{law\_resolution.py}, \code{x\_sweep.py}, \code{oz\_carried.py}, \code{teich.py}, \code{t1\_index.py}, \code{t2\_loo.py}, \code{euler\_link.py}, \code{retro93.py}, \code{retro94.py}, \code{retro99.py}, \code{dh\_logderiv.py}, \code{dh\_zeros.py}, \code{retro101.py}, \code{retro102.py}, \code{retro104.py}, \code{chi5\_zeros.py}, \code{retro105.py} and \code{retro119.py}.

\section{Comparative landscape of $\mathbb{F}_1$-approaches}
\label{app:comparison}

The following table situates the present construction within the
landscape of $\mathbb{F}_1$-approaches.  Columns: object type
(what category of objects the approach defines), counting
(how zeta functions are computed), Frobenius (whether and how
Frobenius lifts are present), RH strategy (the proposed approach
to the Riemann Hypothesis, when stated), Soul\'e's condition (yes
/ no / generalised), and Tits desiderata
\ref{D1}--\ref{D9} engaged.

\begin{center}
\footnotesize
\begin{tabular}{@{}p{2.5cm}p{2cm}p{2cm}p{1.8cm}p{2cm}p{1.5cm}p{2cm}@{}}
\toprule
Approach & Object type & Counting & Frobenius & RH strategy
& Soul\'e? & Desiderata \\
\midrule
Tits 1957~\cite{tits56}
  & buildings, Coxeter complexes
  & combinatorial limit
  & absent
  & none stated
  & N/A
  & \ref{D2}, \ref{D3} \\
Manin 1995~\cite{manin95}
  & motivic curve $\overline{\mathrm{Spec}\,\mathbb{Z}}$
  & via Deninger determinants
  & hypothetical
  & Weil-type via $C\times C$
  & N/A
  & \ref{D5}, \ref{D8}, \ref{D9} \\
Soul\'e 2004~\cite{soule03}
  & varieties with $N_X\in\mathbb{Z}[x]$
  & $\zeta_X(s) = \lim_{q\to 1}$
  & implicit
  & not addressed
  & yes
  & \ref{D4} (toric only) \\
Deitmar 2005~\cite{deitmar05}
  & monoid schemes
  & Hasse--Weil from monoid count
  & monoid endomorphism
  & not addressed
  & yes (toric)
  & \ref{D1}, \ref{D6}, \ref{D7} \\
To\"en--Vaqui\'e 2009~\cite{toen09}
  & relative schemes
  & via base change
  & inherited
  & not addressed
  & \hspace{0pt}generalised
  & \ref{D6} \\
Borger 2009~\cite{borger09}
  & $\Lambda$-rings
  & Witt-functorial
  & explicit $\Frob_p$
  & not addressed directly
  & \hspace{0pt}generalised
  & \ref{D5}, \ref{D6} \\
Connes--Consani 2010~\cite{connesconsani10}
  & $M_o$-schemes
  & distributional $N(u)$
  & implicit via zeros
  & spectral on ad\`ele classes
  & \hspace{0pt}generalised
  & \ref{D4}, \ref{D8} \\
Lorscheid 2012~\cite{lorscheid12blueprints}
  & blueprints
  & via base change
  & inherited
  & not addressed
  & \hspace{0pt}generalised
  & \ref{D6}, \ref{D7} \\
L\'opez Pe\~na--Lorscheid 2011~\cite{lopezpena09torified}
  & torified varieties
  & via torus decomposition
  & inherited
  & not addressed
  & yes (torified)
  & \ref{D7} \\
Durov 2007~\cite{durov07}
  & \hspace{0pt}generalised rings
  & Arakelov-uniform
  & via monad
  & not addressed
  & \hspace{0pt}generalised
  & \ref{D5}, \ref{D6} \\
\midrule
\textbf{PCF (this work)}
  & Artin $\Lambda$-structure on $\Zphi[\tfrac15]$; $\Lambda$-monoid on $\langle\golden\rangle$
  & $N_K(u)$ via Connes--Consani
  & $\Frob_p\in\{\mathrm{id},\sigma\}$ by $\chi_5(p)$; $\golden^n\mapsto\golden^{pn}$
  & equivalent forms of GRH$(\chiv)$ via $G$
  & \hspace{0pt}generalised
  & \ref{D4}; \ref{D2} in one instance; \ref{D9} as equivalent forms; partial \ref{D5}, \ref{D8} \\
\bottomrule
\end{tabular}
\end{center}

The PCF row situates the construction within the Borger--Manin
$\Lambda$-ring tradition --- an Artin $\Lambda$-structure on $\Zphi$ away from $5$ and a
$\Lambda$-monoid on the tower ---
with distributional counting in the
Connes--Consani sense, and with the structural realisation of
the bifurcated Weil geometry via the common source~$G$.  The
construction is thus orthogonal to the toric / monoid-scheme
strand of the field: it does not compete with Deitmar or
blueprints on toric varieties, but rather addresses the
non-toric arithmetic content that those frameworks do not see.




\section{Source map}\label{app:sources}
Each section against its sources: the companion paper and section it comes from
(Primitive Structure for \cite{gonzalez25hp}, CW for \cite{gonzalez26cw}, CR, and OZ for the predecessor version) and the corpus files that carry its Lean
names.  The line-by-line correspondence of every cited name is the paper block of
\texttt{F1\_pentagonal\_descent\_UNIFIED.lean}.

\begin{longtable}{p{0.27\textwidth}p{0.36\textwidth}p{0.30\textwidth}}
\toprule Section & Source in the companions & Corpus\\\midrule\endhead
\S\ref{sec:intro} Introduction & new in this version & ---\\
\S\S\ref{sec:circularity}--\ref{sec:weil} Hilbert--P\'olya, Weil & Primitive Structure \S\S1--2 (Weil template); OZ \S1 & ---\\
\S\S\ref{sec:f1-programme}--\ref{sec:manin} The $\Fone$ programme, Manin & OZ \S1; CW \S1.6; \S\S\ref{sec:polya-hist}--\ref{ssec:four-requirements} from Primitive Structure \S1 (history, requirements of an operator); \S\ref{ssec:intro-binary} from Primitive Structure (the triple obstruction, the forgetful functor, holography) and CW (the isometry of the triad) & \texttt{F1Lawvere}, \texttt{F1PaperNew.triad\_isometry}\\
\S\ref{sec:choice} The route & OZ v1 (pentagonal source); Primitive Structure \S3 & \texttt{F1\_compilado}, \texttt{F1\_PCF\_session\_modules}\\
\S\ref{sec:ring} The ring $\Zphi$ & OZ v1; CW \S\S2.1--2.4 & \texttt{F1\_pendientes}, M12\\
\S\ref{sec:lambda-monoid} The $\Lambda$-monoid & OZ v1 (Frobenius lifts); CW \S2.8 & \texttt{F1\_compilado}, \texttt{F1\_spectra\_anchor\_G}\\
\S\ref{sec:lambda-ring} The Artin $\Lambda$-structure & new in this version; \S\ref{ssec:delta-polar} from the module \texttt{PCFGranularity} and \cite{bhattscholze22} & \texttt{F1\_pendientes}, \texttt{F1\_pendientes\_2}, \texttt{F1\_frobenius\_eigenvalue}\\
\S\ref{sec:chi5} The Fibonacci character & OZ (Fibonacci splitting); CR \S3 & \texttt{S1\_chi5\_fib}\\
\S\ref{sec:grisales} Golden Prime classification & Primitive Structure (Golden Prime classes); \cite{grisales25}; \S\ref{ssec:classes-angular} new, with the module \texttt{PCFGranularity}; the clock, the index by layers and on the $n$-gon from the session \texttt{W4Isometria} (\S\S34--54, 68--73) & \texttt{F1\_compilado}, \texttt{F1\_cw\_equivalents}\\
\S\S\ref{sec:dedekind}--\ref{sec:factorisation} Euler factors, colimit, odd values & OZ \S\S2--4; CW \S2.9 (Def.~2.48, Prop.~2.50); the prism at $5$ (\cref{rmk:prism-five}) from the block of $60$ (\texttt{Sixty}, Kisin and Breuil--Kisin) & \texttt{F1\_compilado}, \texttt{F1\_cw\_equivalents}\\
\S\ref{sec:counting} The counting function & OZ v1 (counting); CW \S2.9, App.~A.8 & \texttt{F1\_v2\_formalizations}\\
\S\ref{sec:mersenne} The Mersenne bridge & Primitive Structure (Mersenne bridge); CR & \texttt{F1\_compilado}, M12\\
\S\ref{ssec:base-phi} Values in base $\golden$ & OZ (anchor $L(1,\chiv)$); new statement & M12, \texttt{F1\_v2}\\
\S\ref{sec:operator} The operator & Primitive Structure (operator $T^*$); OZ; the granularity module & M12, 08\_18, \texttt{PCF\_granularity}\\
\S\ref{sec:machinery} The structure that places the zeros & Primitive Structure; new; \cref{rmk:pappus} from Primitive Structure \S2.1 and CR \cite{gonzalez26mersenne} & \texttt{PCF\_sesion\_completa}, \texttt{PCF\_unificado}, M12\\
\S\ref{sec:topological-proof} The thread & Primitive Structure \S squeeze (restructured); new & \texttt{PCF\_unificado}, \texttt{PCF\_sesion\_completa}\\
\S\ref{sec:results-zeros} The zero set & new & \texttt{PCF\_sesion\_completa}, M12, scripts\\
\S\S\ref{sec:results-weil}--\ref{sec:results-verdict} Weil analogy, windows, verdict & OZ (Weil analogy); Primitive Structure (the operator); CR \cite{gonzalez26mersenne} and \cite{bhattscholze22} (\cref{prop:what-fluctuation-takes,prop:euler-tower}); new ; $\sqrt X$, the energy at every scale, the amplitude and the certified orbit from the session \texttt{W4Isometria} (\S\S55--67) and \code{orbita\_cert.py} & \texttt{F1\_weil\_functional}, \texttt{F1\_v2}, certificates, \texttt{F1PaperHyp}, \texttt{F1PaperTorus}, \texttt{F1PaperPending}, \texttt{F1PaperClosure}, \texttt{F1PaperPadic}, \texttt{tests}\\
\S\ref{sec:discussion} What has been demonstrated & new & ---\\
\S\ref{sec:disc-cr} Connection with CR & CR \cite{gonzalez26mersenne}; the genealogy of the module, Primitive Structure \S2.1, CW; OZ \S6 (the genealogy of the plane) & ---\\
\S\ref{sec:disc-hp} Connection with Primitive Structure & Primitive Structure; the way to the closure with \cite{bhattscholze22,scholze12,connesconsani26}; \S\ref{ssec:disc-binary} new & corpus (the fibres and characters of $G$)\\
\S\ref{sec:disc-cw} Crystalline Worldsheet & OZ Discussion (dimensional chain); OZ \S4 (applications) and \S6 (readings); CW \S\S1.7, 1.9, 2.8, 2.13, 3.2--3.4, 5 & \texttt{F1\_v2}\\
\S\ref{sec:limits} Limitations and future work & new; the open problems of OZ \S6 & ---\\
\S\ref{sec:conclusions} Conclusions & new & ---\\
Appendices & generated from the corpus & \texttt{fqn\_map.json}\\
\bottomrule
\end{longtable}



\section*{Disclosure}\label{sec:disclosure}\src{OZ \texttt{sec:disclosure}, carried over from the predecessor version}

\noindent \textbf{Financial Support:} This research was supported by TTAMAYO PUNTO COM, S.A.P.I. de C.V. (Mexico). The authors also wish to extend their deepest appreciation to Mrs. Kary Fajer. Her steadfast and providential patronage predates the formal existence of TTAMAYO PUNTO COM, having served as a foundational legacy of personal support to the authors that subsequently ensured the long-term resilience and institutional continuity of the organization. This enduring commitment across many years has sustained both the authors and the institution through various critical periods, providing the essential foundation upon which this project was realized. Furthermore, the authors express their profound gratitude to the faculty and the entire team at TTAMAYO PUNTO COM, as well as to our students and clients, whose continued preference and trust have been a vital catalyst in sustaining the creative passion and intellectual rigor that drive our work.

\noindent \textbf{Competing Interests:} None declared.


\section*{Acknowledgments}\src{OZ \texttt{sec:ack}, carried over from the predecessor version}

This paper is dedicated to nature and all her beauty.

The rose has served across traditions as symbol of the
Divine Mother---Tripurasundari in the \textit{Sri Vidya},
the Gothic rosette, the Sufi \textit{gul}---representing the
manifestation of perfection in material form.
The structure is $\golden^2 = \golden + 1$: the rose's five
petals trace the pentagon, her spirals follow the Fibonacci
sequence $F_n/F_{n-1} \to \golden$, and each petal is placed at
the golden angle $360^\circ/\golden^2 \approx 137.5^\circ$.
In the language of Hegel's \textit{Encyclop\"adie}~(\S247),
nature is the Idea in the form of otherness---the rational
structure externalised and asleep in matter.  The mathematics
that arises from observing her is the return: the Idea
awakening to itself.  This paper traces that return from the
pentagon to $\zeta(2k+1)$.

As a companion paper to~\cite{gonzalez25hp}, this work shares the
foundational, institutional, and technical acknowledgments
therein.



\printbibliography

\end{document}
